% Generated mechanically from frozen input p05/ko/flat.tex.
% Do not edit this generated file; edit the bound locale source instead.

% OK, start here.
%

\setcounter{chapter}{37}
\chapter{평탄성 심화}


\phantomsection
\label{flat-section-phantom}




\section{들어가며}
\label{flat-section-introduction}

\noindent
이 장에서는 평탄 가군과 스킴의 평탄 사상에 관한 몇 가지 고급 결과와
그 응용을 논한다. 평탄성에 관한 결과 대부분은 Raynaud와 Gruson의 논문
\cite{GruRay}에서 찾을 수 있다.

\medskip\noindent
이 장을 읽기 전에 다음 결과들을 살펴볼 것을 권한다(이 목록은 앞선
결과들에 대한 안내 역할도 한다).
\begin{enumerate}
\item 평탄 가군에 관한 일반적인 논의는 Algebra, Section
\ref{algebra-section-flat}에 있다.
\item $\text{Tor}$-군과 평탄성의 관계는 Algebra, Section
\ref{algebra-section-tor}을 보라.
\item 평탄성 판정법은 Algebra, Section
\ref{algebra-section-criteria-flatness}(뇌터 경우), Algebra, Section
\ref{algebra-section-flatness-artinian}(아르틴 경우), Algebra, Section
\ref{algebra-section-more-flatness-criteria}(비뇌터 경우), 마지막으로
More on Morphisms, Section
\ref{more-morphisms-section-criterion-flat-fibres}을 보라.
\item 일반 평탄성은 Algebra, Section
\ref{algebra-section-generic-flatness}와 Morphisms, Section
\ref{morphisms-section-generic-flatness}을 보라.
\item 평탄 자취의 열림성은 Algebra, Section
\ref{algebra-section-open-flat}과 More on Morphisms, Section
\ref{more-morphisms-section-open-flat}을 보라.
\item 평탄화는 More on Algebra, Sections
\ref{more-algebra-section-flattening},
\ref{more-algebra-section-flattening-artinian},
\ref{more-algebra-section-flattening-local-base},
\ref{more-algebra-section-flattening-local-source-base}, 및
\ref{more-algebra-section-flattening-Noetherian-complete-local}을 보라.
\item 추가 결과는 More on Algebra, Sections
\ref{more-algebra-section-descent-flatness-integral},
\ref{more-algebra-section-torsion-flat},
\ref{more-algebra-section-flat-finite}, 및
\ref{more-algebra-section-blowup-flat}에 있다.
\end{enumerate}
평탄성에 관한 내용의 응용으로 다음 주제를 다룬다. Grothendieck 존재
정리의 비뇌터 판본, 블로업과 평탄성, 콤팩트화에 관한 Nagata 정리,
h 위상, 블로업 정사각형과 강하, 약한 정규화, 양의 표수에서 벡터 다발의
강하, 그리고 h 위상에서 완전 복합체의 국소 구조이다.






\section{에탈 국소화에 관한 보조정리}
\label{flat-section-etale-localization}

\noindent
이 절에서는 이 장의 뒤에서 유용하게 쓰일 에탈 국소화에 관한 몇 가지
보조정리를 열거한다. 처음 읽을 때에는 이 절을 건너뛰어도 된다.

\begin{lemma}
\label{flat-lemma-lift-etale}
$i : Z \to X$를 아핀 스킴들의 닫힌 몰입이라 하자. $Z' \to Z$를
$Z'$가 아핀인 에탈 사상이라 하자. 그러면 $X'$가 아핀이고 $Z$ 위의
스킴으로서 $Z' \cong Z \times_X X'$가 되는 에탈 사상
$X' \to X$가 존재한다.
\end{lemma}

\begin{proof}
Algebra, Lemma \ref{algebra-lemma-lift-etale}를 보라.
\end{proof}

\begin{lemma}
\label{flat-lemma-etale-at-point}
다음이
$$
\xymatrix{
X \ar[d] & X' \ar[l] \ar[d] \\
S & S' \ar[l]
}
$$
$X' \to X$와 $S' \to S$가 에탈인 스킴들의 가환도식이라고 하자.
$s' \in S'$를 점이라 하자. 그러면
$$
X' \times_{S'} \Spec(\mathcal{O}_{S', s'})
\longrightarrow
X \times_S \Spec(\mathcal{O}_{S', s'})
$$
은 에탈이다.
\end{lemma}

\begin{proof}
$X' \to X_{S'}$는 $X$ 위의 에탈 스킴들 사이의 사상으로서 에탈이다.
Morphisms, Lemma \ref{morphisms-lemma-etale-permanence}를 보라. 또한 에탈
사상의 밑변환은 에탈이다. Morphisms, Lemma
\ref{morphisms-lemma-base-change-etale}를 보라.
\end{proof}

\begin{lemma}
\label{flat-lemma-etale-flat-up-down}
$X \to T \to S$를 스킴의 사상들이라 하고 $T \to S$가 에탈이라고
하자. $\mathcal{F}$를 준연접 $\mathcal{O}_X$-가군이라 하고
$x \in X$를 점이라 하자. 그러면
$$
\mathcal{F}\text{는 }S\text{ 위에서 }x\text{에서 평탄하다}
\Leftrightarrow
\mathcal{F}\text{는 }T\text{ 위에서 }x\text{에서 평탄하다}
$$
특히 $\mathcal{F}$가 $S$ 위에서 평탄일 필요충분조건은 $\mathcal{F}$가
$T$ 위에서 평탄인 것이다.
\end{lemma}

\begin{proof}
에탈 사상은 평탄 사상이므로(Morphisms, Lemma
\ref{morphisms-lemma-etale-flat} 참조), 함의 ``$\Leftarrow$''는 Algebra,
Lemma \ref{algebra-lemma-composition-flat}에서 따른다. 역으로
$\mathcal{F}$가 $S$ 위에서 $x$에서 평탄이라고 가정하자.
$X$에서 $x$ 위에 있고 $T$에서 $x$의 $T$에서의 상 위에 있는 점을
$\tilde x \in X \times_S T$라 표기하자. 그러면
$\text{pr}_2 : X \times_S T \to T$를 통해
$(X \times_S T \to X)^*\mathcal{F}$는 $T$ 위에서 $\tilde x$에서
평탄이다. Morphisms, Lemma
\ref{morphisms-lemma-base-change-module-flat}를 보라. 대각사상
$\Delta_{T/S} : T \to T \times_S T$는 열린 몰입이다. Morphisms, Lemmas
\ref{morphisms-lemma-diagonal-unramified-morphism}와
\ref{morphisms-lemma-etale-smooth-unramified}을 함께 사용하라. 따라서
$X$는 $X \times_S T$의 열린부분스킴과 동일시된다. 이 열린부분에 대한
$\text{pr}_2$의 제한은 주어진 사상 $X \to T$이고, 점 $\tilde x$는 이
열린부분의 점 $x$에 대응하며, $(X \times_S T \to X)^*\mathcal{F}$를
이 열린부분에 제한하면 $\mathcal{F}$가 된다. 그러므로 $\mathcal{F}$는
$T$ 위에서 $x$에서 평탄이다.
\end{proof}

\begin{lemma}
\label{flat-lemma-etale-flat-up-down-local-ring}
$T \to S$를 에탈 사상이라 하자. $t \in T$의 상을 $s \in S$라 하자.
$M$을 $\mathcal{O}_{T, t}$-가군이라 하자. 그러면
$$
M\text{은 }\mathcal{O}_{S, s}\text{ 위에서 평탄하다}
\Leftrightarrow
M\text{은 }\mathcal{O}_{T, t}\text{ 위에서 평탄하다}.
$$
\end{lemma}

\begin{proof}
$S$를 $s$의 아핀 근방으로 대체한 뒤 $T$를 $t$의 아핀 근방으로
대체해도 된다. $\mathcal{F} = (\Spec(\mathcal{O}_{T, t}) \to
T)_*\widetilde M$으로 놓자. 이는 $T$ 위의 준연접층이고(Schemes, Lemma
\ref{schemes-lemma-push-forward-quasi-coherent}를 보거나 직접 논증하라),
$t$에서의 줄기는 $M$이다(세부사항은 생략한다). Lemma
\ref{flat-lemma-etale-flat-up-down}을 적용하라.
\end{proof}

\begin{lemma}
\label{flat-lemma-flat-up-down-henselization}
$S$를 스킴이라 하고 $s \in S$를 점이라 하자. $\mathcal{O}_{S, s}^h$
(각각 $\mathcal{O}_{S, s}^{sh}$)로 헨젤화(각각 강한 헨젤화)를
나타내자. Algebra, Definition \ref{algebra-definition-henselization}을
보라. $M^{sh}$를 $\mathcal{O}_{S, s}^{sh}$-가군이라 하자. 다음은
서로 동치이다.
\begin{enumerate}
\item $M^{sh}$는 $\mathcal{O}_{S, s}$ 위에서 평탄이다.
\item $M^{sh}$는 $\mathcal{O}_{S, s}^h$ 위에서 평탄이다.
\item $M^{sh}$는 $\mathcal{O}_{S, s}^{sh}$ 위에서 평탄이다.
\end{enumerate}
$M^{sh} = M^h \otimes_{\mathcal{O}_{S, s}^h} \mathcal{O}_{S, s}^{sh}$이면
이들은 다음과도 동치이다.
\begin{enumerate}
\item[(4)] $M^h$는 $\mathcal{O}_{S, s}$ 위에서 평탄이다.
\item[(5)] $M^h$는 $\mathcal{O}_{S, s}^h$ 위에서 평탄이다.
\end{enumerate}
$M^h = M \otimes_{\mathcal{O}_{S, s}} \mathcal{O}_{S, s}^h$이면 이들은
다음과도 동치이다.
\begin{enumerate}
\item[(6)] $M$은 $\mathcal{O}_{S, s}$ 위에서 평탄이다.
\end{enumerate}
\end{lemma}

\begin{proof}
More on Algebra, Lemma
\ref{more-algebra-lemma-dumb-properties-henselization}에 의해 국소환 사상
$\mathcal{O}_{S, s} \to \mathcal{O}_{S, s}^h \to \mathcal{O}_{S, s}^{sh}$
은 충실 평탄이다. 따라서 (3) $\Rightarrow$ (2) $\Rightarrow$ (1)과
(5) $\Rightarrow$ (4)는 Algebra, Lemma
\ref{algebra-lemma-composition-flat}에서 따른다. 충실 평탄성에 의해
(6) $\Leftrightarrow$ (5)와 (5) $\Leftrightarrow$ (3)은 Algebra, Lemma
\ref{algebra-lemma-flatness-descends}에서 따른다. 그러므로
(1) $\Rightarrow$ (2) $\Rightarrow$ (3)과 (4) $\Rightarrow$ (5)를
보이면 충분하다. 이를 증명할 때 $S$가 아핀 스킴이라고 가정해도 된다.

\medskip\noindent
(1)을 가정하자. Lemma \ref{flat-lemma-etale-flat-up-down-local-ring}에 의해
임의의 에탈 근방 $(T, t) \to (S, s)$에 대하여 $M^{sh}$는
$\mathcal{O}_{T, t}$ 위에서 평탄이다. $\mathcal{O}_{S, s}^h$와
$\mathcal{O}_{S, s}^{sh}$는 $\mathcal{O}_{T, t}$ 꼴의 국소환들의 유향
극한이므로(Algebra, Lemmas \ref{algebra-lemma-henselization-different}와
\ref{algebra-lemma-strict-henselization-different} 참조), Algebra, Lemma
\ref{algebra-lemma-colimit-rings-flat}에 의해 $M^{sh}$는
$\mathcal{O}_{S, s}^h$와 $\mathcal{O}_{S, s}^{sh}$ 위에서 평탄이다.
따라서 (1)은 (2)와 (3)을 함의한다. $\mathcal{O}_{S, s}$를
$\mathcal{O}_{S, s}^h$로 대체하면 물론 (2) $\Rightarrow$ (3)도 따른다.
같은 논증으로 (4) $\Rightarrow$ (5)가 증명된다.
\end{proof}

\begin{lemma}
\label{flat-lemma-tor-amplitude-up-down-henselization}
$S$를 스킴이라 하고 $s \in S$를 점이라 하자. $\mathcal{O}_{S, s}^h$
(각각 $\mathcal{O}_{S, s}^{sh}$)로 헨젤화(각각 강한 헨젤화)를
나타내자. Algebra, Definition \ref{algebra-definition-henselization}을
보라. $M^{sh}$를 $D(\mathcal{O}_{S, s}^{sh})$의 대상이라 하고
$a, b \in \mathbf{Z}$라 하자. 다음은 서로 동치이다.
\begin{enumerate}
\item $M^{sh}$는 $\mathcal{O}_{S, s}$ 위에서 $[a, b]$의 Tor 진폭을 갖는다.
\item $M^{sh}$는 $\mathcal{O}_{S, s}^h$ 위에서 $[a, b]$의 Tor 진폭을 갖는다.
\item $M^{sh}$는 $\mathcal{O}_{S, s}^{sh}$ 위에서 $[a, b]$의 Tor 진폭을 갖는다.
\end{enumerate}
만일 $M^{sh} =
M^h \otimes_{\mathcal{O}_{S, s}^h}^\mathbf{L} \mathcal{O}_{S, s}^{sh}$
이고 $M^h \in D(\mathcal{O}_{S, s}^h)$이면 이들은 다음과도 동치이다.
\begin{enumerate}
\item[(4)] $M^h$는 $\mathcal{O}_{S, s}$ 위에서 $[a, b]$의 Tor 진폭을 갖는다.
\item[(5)] $M^h$는 $\mathcal{O}_{S, s}^h$ 위에서 $[a, b]$의 Tor 진폭을 갖는다.
\end{enumerate}
만일 $M^h = M \otimes_{\mathcal{O}_{S, s}}^\mathbf{L} \mathcal{O}_{S, s}^h$
이고 $M \in D(\mathcal{O}_{S, s})$이면 이들은 다음과도 동치이다.
\begin{enumerate}
\item[(6)] $M$은 $\mathcal{O}_{S, s}$ 위에서 $[a, b]$의 Tor 진폭을 갖는다.
\end{enumerate}
\end{lemma}

\begin{proof}
More on Algebra, Lemma
\ref{more-algebra-lemma-dumb-properties-henselization}에 의해 국소환 사상
$\mathcal{O}_{S, s} \to \mathcal{O}_{S, s}^h \to \mathcal{O}_{S, s}^{sh}$
은 충실 평탄이다. 따라서 (3) $\Rightarrow$ (2) $\Rightarrow$ (1)과
(5) $\Rightarrow$ (4)는 More on Algebra, Lemma
\ref{more-algebra-lemma-flat-push-tor-amplitude}에서 따른다. 충실
평탄성에 의해 (6) $\Leftrightarrow$ (5)와 (5) $\Leftrightarrow$ (3)은
More on Algebra, Lemma
\ref{more-algebra-lemma-flat-descent-tor-amplitude}에서 따른다. 그러므로
(1) $\Rightarrow$ (3), (2) $\Rightarrow$ (3), 그리고
(4) $\Rightarrow$ (5)를 보이면 충분하다.

\medskip\noindent
(1)을 가정하자. 특히 $M^{sh}$의 코호몰로지는 차수 $< a$와 $> b$에서
소멸한다. 따라서 $i > b$이면 $P^i = 0$인 자유
% P05-EN-CH38-0001
$\mathcal{O}_{X, x}^{sh}$-가군들의 복합체 $P^\bullet$로 $M^{sh}$를
표현할 수 있다(예를 들어 매우 일반적인 Derived Categories, Lemma
\ref{derived-lemma-subcategory-left-resolution}을 보라). 모든 $n$에
대하여 $P^n$은 $\mathcal{O}_{S, s}$ 위에서 평탄이다.
$\Coker(d_P^{a - 1})$를 생각하자. More on Algebra, Lemma
\ref{more-algebra-lemma-last-one-flat}에 의해 이는 평탄
$\mathcal{O}_{S, s}$-가군이다. 따라서 Lemma
\ref{flat-lemma-flat-up-down-henselization}에 의해 평탄
$\mathcal{O}_{S, s}^{sh}$-가군이다. 그러므로
$\tau_{\geq a}P^\bullet$는
% P05-EN-CH38-0002
$D(\mathcal{O}_{S, s}^{sh}$에서 $M^{sh}$를 표현하는 평탄
$\mathcal{O}_{S, s}^{sh}$-가군들의 복합체이고, $M^{sh}$가
$[a, b]$의 Tor 진폭을 가짐을 알 수 있다. More on Algebra, Lemma
\ref{more-algebra-lemma-tor-amplitude}을 보라. 따라서 (1)은 (3)을
함의한다. $\mathcal{O}_{S, s}$를 $\mathcal{O}_{S, s}^h$로 대체하면
물론 (2) $\Rightarrow$ (3)도 따른다. 같은 논증으로
(4) $\Rightarrow$ (5)가 증명된다.
\end{proof}

\begin{lemma}
\label{flat-lemma-finite-flat-weak-assassin-up-down}
$g : T \to S$를 스킴의 유한 평탄 사상이라 하자. $\mathcal{G}$를
준연접 $\mathcal{O}_S$-가군이라 하자. $t \in T$의 상을
$s \in S$라 하자. 그러면
$$
t \in \text{WeakAss}(g^*\mathcal{G})
\Leftrightarrow
s \in \text{WeakAss}(\mathcal{G})
$$
\end{lemma}

\begin{proof}
함의 ``$\Leftarrow$''는 Divisors, Lemma
\ref{divisors-lemma-weakly-ass-pullback}에서 바로 따른다.
$t \in \text{WeakAss}(g^*\mathcal{G})$라고 가정하자.
$\Spec(A) \subset S$를 $s$의 아핀 열린 근방이라 하자. $\mathcal{G}$를
$A$-가군 $M$에 대응하는 준연접층이라 하고, $\mathfrak p \subset A$를
$s$에 대응하는 소 아이디얼이라 하자. $g$가 유한 평탄이므로 어떤 유한
평탄 $A$-대수 $B$에 대하여 $g^{-1}(\Spec(A)) = \Spec(B)$이다.
$g^*\mathcal{G}$는 $B$-가군 $M \otimes_A B$에 대응하는 준연접
$\mathcal{O}_{\Spec(B)}$-가군이고, $g_*g^*\mathcal{G}$는 $A$-가군
$M \otimes_A B$에 대응하는 준연접 $\mathcal{O}_{\Spec(A)}$-가군이다.
Algebra, Lemma \ref{algebra-lemma-finite-flat-local}에 의해 어떤 정수
$n \geq 0$에 대하여 $B_{\mathfrak p} \cong A_{\mathfrak p}^{\oplus n}$이다.
$s$ 위에 놓이는 $T$의 점이 적어도 하나 존재한다고 가정했으므로
$n \geq 1$임에 유의하자. 따라서 줄기들을 살펴보면
$$
s \in \text{WeakAss}(\mathcal{G})
\Leftrightarrow
s \in \text{WeakAss}(g_*g^*\mathcal{G})
$$
이제 가정 $t \in \text{WeakAss}(g^*\mathcal{G})$는 Divisors, Lemma
\ref{divisors-lemma-weakly-associated-finite}에 의해
$s \in \text{WeakAss}(g_*g^*\mathcal{G})$를 함의한다. 따라서 위의
동치에 의해 $s \in  \text{WeakAss}(\mathcal{G})$이다.
\end{proof}

\begin{lemma}
\label{flat-lemma-etale-weak-assassin-up-down}
$h : U \to S$를 스킴의 에탈 사상이라 하자. $\mathcal{G}$를 준연접
$\mathcal{O}_S$-가군이라 하자. $u \in U$의 상을 $s \in S$라 하자.
그러면
$$
u \in \text{WeakAss}(h^*\mathcal{G})
\Leftrightarrow
s \in \text{WeakAss}(\mathcal{G})
$$
\end{lemma}

\begin{proof}
$S$와 $U$를 각각 $s$와 $u$의 아핀 근방으로 대체하면
% P05-EN-CH38-0003
$g$가 아핀 스킴들의 표준 에탈 사상이라고 가정해도 된다. Morphisms,
Lemma \ref{morphisms-lemma-etale-locally-standard-etale}를 보라. 따라서
$S = \Spec(A)$이고
% P05-EN-CH38-0004
$X = \Spec(A[x, 1/g]/(f))$라고 가정해도 된다. 여기서 $f$는 일계수이고
$f'$은 $A[x, 1/g]$에서 가역이다. $A[x, 1/g]/(f) = (A[x]/(f))_g$는
유한 자유 $A$-대수 $A[x]/(f)$의 국소화이기도 하다. 따라서 $U$를
$S$ 위에서 유한 국소 자유인 스킴 $T = \Spec(A[x]/(f))$의 열린부분스킴으로
생각할 수 있다. 이로써 위의 Lemma
\ref{flat-lemma-finite-flat-weak-assassin-up-down}로 환원된다.
\end{proof}

\begin{lemma}
\label{flat-lemma-weakly-associated-henselization}
$S$를 스킴이라 하고 $s \in S$를 점이라 하자. $\mathcal{O}_{S, s}^h$
(각각 $\mathcal{O}_{S, s}^{sh}$)로 헨젤화(각각 강한 헨젤화)를
나타내자. Algebra, Definition \ref{algebra-definition-henselization}을
보라. $\mathcal{F}$를 준연접 $\mathcal{O}_S$-가군이라 하자. 다음은
서로 동치이다.
\begin{enumerate}
\item $s$는 $\mathcal{F}$의 약연관점이다.
\item $\mathfrak m_s$는 $\mathcal{F}_s$의 약연관 소 아이디얼이다.
\item $\mathfrak m_s^h$는
$\mathcal{F}_s \otimes_{\mathcal{O}_{S, s}} \mathcal{O}_{S, s}^h$의
약연관 소 아이디얼이다.
\item $\mathfrak m_s^{sh}$는
$\mathcal{F}_s \otimes_{\mathcal{O}_{S, s}} \mathcal{O}_{S, s}^{sh}$의
약연관 소 아이디얼이다.
\end{enumerate}
\end{lemma}

\begin{proof}
(1)과 (2)의 동치는 정의이다. Divisors, Definition
\ref{divisors-definition-weakly-associated}를 보라.
% P05-EN-CH38-0005
(2) $\Rightarrow$ (3) $\Rightarrow$ (4)는 평탄 사상들(More on Algebra,
Lemma \ref{more-algebra-lemma-dumb-properties-henselization})
$$
\Spec(\mathcal{O}_{S, s}^{sh}) \to
\Spec(\mathcal{O}_{S, s}^h) \to
\Spec(\mathcal{O}_{S, s})
$$
과 닫힌점들에 Divisors, Lemma
\ref{divisors-lemma-weakly-ass-pullback}을 적용하여 얻는다.
(4) $\Rightarrow$ (2)를 증명하기 위해 $S$를 아핀 근방으로 대체해도
된다. 다음을 만족하는 원소
$x \in \mathcal{F}_s \otimes_{\mathcal{O}_{S, s}} \mathcal{O}_{S, s}^{sh}$
를 생각하자. 그 소멸자의 근기는 $\mathfrak m_s^{sh}$와 같다고 하자
(Algebra, Lemma \ref{algebra-lemma-weakly-ass-local} 참조).
% P05-EN-CH38-0006
Algebra, Lemma \ref{algebra-lemma-strict-henselization-different}에 의해
$\mathcal{O}_{S, s}^{sh}$는 에탈 근방 $f : (U, u) \to (S, s)$에 걸친
$\mathcal{O}_{U, u}$의 극한과 같으므로, $x$가 어떤
$x' \in \mathcal{F}_s \otimes_{\mathcal{O}_{S, s}} \mathcal{O}_{U, u}$
의 상이라고 가정해도 된다. 국소환 사상
$\mathcal{O}_{U, u} \to \mathcal{O}_{S, s}^{sh}$는 강한 헨젤화이므로
충실 평탄이고, 따라서 보편 단사이다(Algebra, Lemma
\ref{algebra-lemma-faithfully-flat-universally-injective}). 그러므로
$x'$의 소멸자는 $x$의 소멸자의 역상이다. 따라서 이 소멸자의 근기는
$\mathfrak m_u$와 같다. 그러므로 $u$는 $f^*\mathcal{F}$의 약연관점이다.
Lemma \ref{flat-lemma-etale-weak-assassin-up-down}에 의해 $s$는
$\mathcal{F}$의 약연관점이다.
\end{proof}





\section{유한형 가군의 국소 구조}
\label{flat-section-local-structure-module}

\noindent
이 장의 많은 논증을 작동하게 하는 핵심 기술적 보조정리는 기하학적인
Lemma \ref{flat-lemma-elementary-devissage}이다.

\begin{lemma}
\label{flat-lemma-sheaf-lives-on-subscheme}
$f : X \to S$를 아핀 스킴들의 유한형 사상이라 하자.
$\mathcal{F}$를 유한형 준연접 $\mathcal{O}_X$-가군이라 하자.
$x \in X$의 $S$에서의 상을 $s = f(x)$라 하자.
$\mathcal{F}_s = \mathcal{F}|_{X_s}$로 두자.
그러면 유한표현인 닫힌 몰입 $i : Z \to X$와 유한형 준연접
$\mathcal{O}_Z$-가군 $\mathcal{G}$가 존재하여
$i_*\mathcal{G} = \mathcal{F}$이고
$Z_s = \text{Supp}(\mathcal{F}_s)$이다.
\end{lemma}

\begin{proof}
$f : X \to S$가 환 준동형 $A \to B$로 주어지고, $\mathcal{F}$가
$B$-가군 $M$에 대응하는 준연접층이라고 하자. Morphisms, Lemma
\ref{morphisms-lemma-locally-finite-type-characterize}에 의해
$A \to B$는 유한형 환 준동형이고, Properties, Lemma
\ref{properties-lemma-finite-type-module}에 의해 $M$은 유한
$B$-가군이다. 특히 $\mathcal{F}$의 지지는 $M$의 소멸자
$I = \{x \in B \mid xm = 0\ \forall m \in M\}$로 잘라낸
$\Spec(B)$의 닫힌 부분스킴이다. Algebra, Lemma
\ref{algebra-lemma-support-closed}를 보라.
$x$에 대응하는 소 아이디얼을 $\mathfrak q \subset B$, $s$에 대응하는
소 아이디얼을 $\mathfrak p \subset A$라 하자.
$X_s = \Spec(B \otimes_A \kappa(\mathfrak p))$이고
$\mathcal{F}_s$는 $B \otimes_A \kappa(\mathfrak p)$-가군
$M \otimes_A \kappa(\mathfrak p)$에 대응하는 준연접층임에 유의하자.
Morphisms, Lemma \ref{morphisms-lemma-support-finite-type}에 의해
$\mathcal{F}_s$의 지지는
$V(I(B \otimes_A \kappa(\mathfrak p)))$와 같다.
$B \otimes_A \kappa(\mathfrak p)$가 $\kappa(\mathfrak p)$ 위 유한형이므로
다음을 만족하는 유한개의 원소 $f_1, \ldots, f_m \in I$가 존재한다.
% P05-EN-CH38-0007
$$
I(B \otimes_A \kappa(\mathfrak p)) =
(f_1, \ldots, f_n)(B \otimes_A \kappa(\mathfrak p)).
$$
$(f_1, \ldots, f_m)$로 잘라낸 닫힌 부분스킴을 $i : Z \to X$라
쓰자. 식으로는 $Z = \Spec(B/(f_1, \ldots, f_m))$이다.
$M$은 $I$에 의해 소멸하므로 $M$을
$B/(f_1, \ldots, f_m)$-가군으로 볼 수 있다. 다시 말해
$\mathcal{F}$는 $Z$ 위 유한형 가군의 직상이다.
구성상 $Z_s = \text{Supp}(\mathcal{F}_s)$이므로 보조정리가 증명된다.
\end{proof}

\begin{lemma}
\label{flat-lemma-elementary-devissage}
% P05-EN-CH38-0008
$f : X \to S$를 국소적으로 유한형인 스킴 사상이라 하자.
$\mathcal{F}$를 유한형 준연접 $\mathcal{O}_X$-가군이라 하자.
$x \in X$의 $S$에서의 상을 $s = f(x)$라 하자.
$\mathcal{F}_s = \mathcal{F}|_{X_s}$ 및
$n = \dim_x(\text{Supp}(\mathcal{F}_s))$로 두자.
그러면 다음을 구성할 수 있다.
\begin{enumerate}
\item 기본 에탈 근방들 $g : (X', x') \to (X, x)$,
$e : (S', s') \to (S, s)$,
\item 가환 그림
$$
\xymatrix{
X \ar[dd]_f & X' \ar[dd] \ar[l]^g & Z' \ar[l]^i \ar[d]^\pi \\
& & Y' \ar[d]^h \\
S & S' \ar[l]_e & S' \ar@{=}[l]
}
$$
\item $i(z') = x'$, $y' = \pi(z')$, $h(y') = s'$를 만족하는 점
$z' \in Z'$,
\item 유한형 준연접 $\mathcal{O}_{Z'}$-가군 $\mathcal{G}$.
\end{enumerate}
이들은 다음 성질들을 만족한다.
\begin{enumerate}
\item $X'$, $Z'$, $Y'$, $S'$는 아핀 스킴이다.
\item $i$는 유한표현인 닫힌 몰입이다.
\item $i_*(\mathcal{G}) \cong g^*\mathcal{F}$,
\item $\pi$는 유한이고 $\pi^{-1}(\{y'\}) = \{z'\}$이다.
\item 확대 $\kappa(y')/\kappa(s')$는 순수 초월이다.
\item $h$는 상대차원이 $n$인 매끄러운 사상이고 그 올들은
기하적으로 정역이다.
\end{enumerate}
\end{lemma}

\begin{proof}
$V \subset S$를 $s$의 아핀 근방이라 하고,
$U \subset f^{-1}(V)$를 $x$의 아핀 근방이라 하자.
그러면 $f|_U : U \to V$와 $\mathcal{F}|_U$에 대해 보조정리를
증명하면 충분하다. 따라서 이하의 증명에서는 $X$와 $S$가 아핀이라고
가정한다.

\medskip\noindent
먼저 $X_s = \text{Supp}(\mathcal{F}_s)$라고 가정하자. 특히
$n = \dim_x(X_s)$이다. More on Morphisms, Lemmas
\ref{more-morphisms-lemma-local-local-structure-finite-type} 및
\ref{more-morphisms-lemma-local-local-structure-finite-type-affine}를
적용한다. 그러면 가환 그림
$$
\xymatrix{
X \ar[dd] & X' \ar[l]^g \ar[d]^\pi \\
& Y' \ar[d]^h  \\
S & S' \ar[l]_e
}
$$
과 점 $x' \in X'$를 얻는다. $Z' = X'$, $i = \text{id}$,
$\mathcal{G} = g^*\mathcal{F}$로 두면 이 경우의 원하는 자료를 얻는다.

\medskip\noindent
일반적인 경우에는 Lemma \ref{flat-lemma-sheaf-lives-on-subscheme}와 같이
닫힌 몰입 $Z \to X$와 $Z$ 위의 층 $\mathcal{G}$를 택한다.
앞 문단의 결과를 $Z \to S$와 $\mathcal{G}$에 적용하면 그림
$$
\xymatrix{
X \ar[dd]_f & Z \ar[l] \ar[dd]_{f|_Z} & Z' \ar[l]^g \ar[d]^\pi \\
& & Y' \ar[d]^h \\
S & S \ar@{=}[l] & S' \ar[l]_e
}
$$
과 필요한 모든 성질을 만족하는 점 $z' \in Z'$를 얻는다.
Lemma \ref{flat-lemma-lift-etale}를 이용하여 $Z'$을 $X$ 위 에탈인 스킴에
매장할 것이다. $X'$가 $S'$ 위의 스킴이기를 원하므로 이 보조정리를
직접 적용할 수는 없다. 대신 다음 사상들을 생각한다.
$$
\xymatrix{
Z' \ar[r] & Z \times_S S' \ar[r] & X \times_S S'
}
$$
첫 번째 사상은 Morphisms, Lemma
\ref{morphisms-lemma-etale-permanence}에 의해 에탈이다.
두 번째 사상은 닫힌 몰입의 기저변환이므로 닫힌 몰입이다.
마지막으로 $X$, $S$, $S'$, $Z$, $Z'$가 모두 아핀이므로
Lemma \ref{flat-lemma-lift-etale}를 적용하여 다음을 만족하는 아핀 스킴들의
에탈 사상 $X' \to X \times_S S'$를 얻는다.
$$
Z' = (Z \times_S S') \times_{(X \times_S S')} X' = Z \times_X X'.
$$
$Z \to X$가 유한표현인 닫힌 몰입이므로 $Z' \to X'$도 그러하다.
$x' \in X'$를 $z' \in Z'$에 대응하는 점이라 하자.
그러면 완성된 그림
$$
\xymatrix{
X \ar[dd] & X' \ar[dd] \ar[l] & Z' \ar[l]^i \ar[d]^\pi \\
& & Y' \ar[d]^h \\
S & S' \ar[l]_e & S' \ar@{=}[l]
}
$$
은 원래 문제의 해이다.
\end{proof}

\begin{lemma}
\label{flat-lemma-devissage-finite-presentation}
가정과 표기는 Lemma \ref{flat-lemma-elementary-devissage}와 같다고 하자.
$f$가 국소 유한표현이면 $\pi$는 유한표현이다.
이 경우 다음은 서로 동치이다.
\begin{enumerate}
\item $\mathcal{F}$는 $x$의 근방에서 유한표현
$\mathcal{O}_X$-가군이다.
\item $\mathcal{G}$는 $z'$의 근방에서 유한표현
$\mathcal{O}_{Z'}$-가군이다.
\item $\pi_*\mathcal{G}$는 $y'$의 근방에서 유한표현
$\mathcal{O}_{Y'}$-가군이다.
\end{enumerate}
계속해서 $f$가 국소 유한표현이라고 가정하면 다음도 서로 동치이다.
\begin{enumerate}
\item[(a)] $\mathcal{F}_x$는 유한표현
$\mathcal{O}_{X, x}$-가군이다.
\item[(b)] $\mathcal{G}_{z'}$는 유한표현
$\mathcal{O}_{Z', z'}$-가군이다.
\item[(c)] $(\pi_*\mathcal{G})_{y'}$는 유한표현
$\mathcal{O}_{Y', y'}$-가군이다.
\end{enumerate}
\end{lemma}

\begin{proof}
$f$가 국소 유한표현이라고 가정하자. 그러면 $Z' \to S$는 그러한
사상들의 합성이므로 국소 유한표현이다. Morphisms, Lemma
\ref{morphisms-lemma-composition-finite-presentation}를 보라.
$Y' \to S$도 매끄러운 사상과 에탈 사상의 합성이므로 국소
유한표현임에 유의하자. 따라서 Morphisms, Lemma
\ref{morphisms-lemma-finite-presentation-permanence}에 의해
$\pi$는 국소 유한표현이다. $\pi$는 유한이므로 분리이고 준콤팩트하며,
따라서 $\pi$는 실제로 유한표현이다.

\medskip\noindent
(1), (2), (3)의 동치를 증명하기 위해 다음 조건도 생각하자.
(4) $g^*\mathcal{F}$는 $x'$의 근방에서 유한표현
$\mathcal{O}_{X'}$-가군이다. 유한표현 가군의 당김은 유한표현이다.
Modules, Lemma \ref{modules-lemma-pullback-finite-presentation}를 보라.
따라서 (1) $\Rightarrow$ (4)이다.
에탈 사상 $g$는 열린 사상이다. Morphisms, Lemma
\ref{morphisms-lemma-etale-open}을 보라.
따라서 $x'$의 임의의 열린 근방 $U' \subset X'$에 대해 상 $g(U')$는
$x$의 열린 근방이고, 사상족 $\{U' \to g(U')\}$은 에탈 덮개이다.
그러므로 Descent, Lemma
\ref{descent-lemma-finite-presentation-descends}에 의해
(4) $\Rightarrow$ (1)이다. Descent, Lemma
\ref{descent-lemma-finite-finitely-presented-module}와 몇 가지 쉬운
위상적 논증(More on Morphisms, Lemma
\ref{more-morphisms-lemma-finite-morphism-single-point-in-fibre} 참조)을
사용하면 (4) $\Leftrightarrow$ (2) $\Leftrightarrow$ (3)임을 알 수 있다.

\medskip\noindent
(a), (b), (c)의 동치를 증명하기 위해 환 준동형들을 생각하자.
$$
\mathcal{O}_{X, x} \to
\mathcal{O}_{X', x'} \to
\mathcal{O}_{Z', z'} \leftarrow
\mathcal{O}_{Y', y'}
$$
첫 번째 환 준동형은 충실히 평탄하다. 따라서
$\mathcal{F}_x$가 $\mathcal{O}_{X, x}$ 위 유한표현인 것과
% P05-EN-CH38-0010
$g^*\mathcal{F}_{x'}$가 $\mathcal{O}_{X', x'}$ 위 유한표현인 것은
동치이다. Algebra, Lemma
\ref{algebra-lemma-descend-properties-modules}를 보라.
두 번째 환 준동형은 전사(따라서 유한)이고 가정에 의해 유한표현이므로,
$g^*\mathcal{F}_{x'}$가 $\mathcal{O}_{X', x'}$ 위 유한표현인 것과
$\mathcal{G}_{z'}$가 $\mathcal{O}_{Z', z'}$ 위 유한표현인 것은
동치이다. Algebra, Lemma
\ref{algebra-lemma-finite-finitely-presented-extension}을 보라.
% P05-EN-CH38-0009
$\pi$는 유한이고 유한표현이며
$\pi^{-1}(\{y'\}) = \{x'\}$이므로, 환 준동형
$\mathcal{O}_{Y', y'} \leftarrow \mathcal{O}_{Z', z'}$는 유한이고
유한표현이다. More on Morphisms, Lemma
\ref{more-morphisms-lemma-finite-morphism-single-point-in-fibre}를 보라.
따라서 $\mathcal{G}_{z'}$가 $\mathcal{O}_{Z', z'}$ 위 유한표현인 것과
% P05-EN-CH38-0011
$\pi_*\mathcal{G}_{y'}$가 $\mathcal{O}_{Y', y'}$ 위 유한표현인 것은
동치이다. Algebra, Lemma
\ref{algebra-lemma-finite-finitely-presented-extension}을 보라.
\end{proof}

\begin{lemma}
\label{flat-lemma-devissage-flat}
가정과 표기는 Lemma \ref{flat-lemma-elementary-devissage}와 같다고 하자.
다음은 서로 동치이다.
\begin{enumerate}
\item $\mathcal{F}$는 $x$의 근방에서 $S$ 위 평탄하다.
\item $\mathcal{G}$는 $z'$의 근방에서 $S'$ 위 평탄하다.
\item $\pi_*\mathcal{G}$는 $y'$의 근방에서 $S'$ 위 평탄하다.
\end{enumerate}
다음도 서로 동치이다.
\begin{enumerate}
\item[(a)] $\mathcal{F}_x$는 $\mathcal{O}_{S, s}$ 위 평탄하다.
\item[(b)] $\mathcal{G}_{z'}$는 $\mathcal{O}_{S', s'}$ 위 평탄하다.
\item[(c)] $(\pi_*\mathcal{G})_{y'}$는 $\mathcal{O}_{S', s'}$ 위
평탄하다.
\end{enumerate}
\end{lemma}

\begin{proof}
(1), (2), (3)의 동치를 증명하기 위해 다음 조건도 생각하자.
(4) $g^*\mathcal{F}$는 $x'$의 근방에서 $S$ 위 평탄하다.
이하에서는 따로 말하지 않고 Lemma \ref{flat-lemma-etale-flat-up-down}을
사용하여 $S$ 위의 평탄성과 $S'$ 위의 평탄성을 동일시하겠다.
에탈 사상 $g$는 평탄하고 열린 사상이다. Morphisms, Lemma
\ref{morphisms-lemma-etale-open}을 보라.
따라서 $x'$의 임의의 열린 근방 $U' \subset X'$에 대해 상 $g(U')$는
$x$의 열린 근방이고 사상 $U' \to g(U')$은 전사이고 평탄하다.
그러므로 Morphisms, Lemma
\ref{morphisms-lemma-flat-permanence}에 의해
(4) $\Leftrightarrow$ (1)이다. 다음에 유의하자.
$$
\Gamma(X', g^*\mathcal{F}) =
\Gamma(Z', \mathcal{G}) =
\Gamma(Y', \pi_*\mathcal{G})
$$
따라서 $S'$ 위에서 $g^*\mathcal{F}$, $\mathcal{G}$,
$\pi_*\mathcal{G}$가 평탄하다는 조건들은 모두 동치이다
($X'$, $Z'$, $Y'$, $S'$가 모두 아핀이라는 사실을 사용한다).
아핀 근방에 관한 생략된 몇 가지 위상적 논증(More on Morphisms,
Lemma \ref{more-morphisms-lemma-finite-morphism-single-point-in-fibre}와
비교하라)을 적용하면 (4) $\Leftrightarrow$ (2) $\Leftrightarrow$ (3)임을
알 수 있다.

\medskip\noindent
(a), (b), (c)의 동치를 증명하기 위해 국소환 준동형들의 가환 그림을
생각하자.
$$
\xymatrix{
\mathcal{O}_{X', x'} \ar[r]_\iota &
\mathcal{O}_{Z', z'} &
\mathcal{O}_{Y', y'} \ar[l]^\alpha &
\mathcal{O}_{S', s'} \ar[l]^\beta \\
\mathcal{O}_{X, x} \ar[u]^\gamma & & &
\mathcal{O}_{S, s} \ar[lll]_\varphi \ar[u]_\epsilon
}
$$
이하에서는 따로 말하지 않고 Lemma
\ref{flat-lemma-etale-flat-up-down-local-ring}을 사용하여
$\mathcal{O}_{S, s}$ 위의 평탄성과 $\mathcal{O}_{S', s'}$ 위의
평탄성을 동일시하겠다. 사상 $\gamma$는 충실히 평탄하다. 따라서
$\mathcal{F}_x$가 $\mathcal{O}_{S, s}$ 위 평탄인 것과
% P05-EN-CH38-0010
$g^*\mathcal{F}_{x'}$가 $\mathcal{O}_{S', s'}$ 위 평탄인 것은
동치이다. Algebra, Lemma
\ref{algebra-lemma-flatness-descends-more-general}을 보라.
$\mathcal{O}_{S', s'}$-가군으로서 가군들
$g^*\mathcal{F}_{x'}$, $\mathcal{G}_{z'}$, 그리고
% P05-EN-CH38-0011
$\pi_*\mathcal{G}_{y'}$는 모두 동형이다. More on Morphisms, Lemma
\ref{more-morphisms-lemma-finite-morphism-single-point-in-fibre}를 보라.
이로써 증명이 끝난다.
\end{proof}









\section{1단계 데비사주}
\label{flat-section-one-step-devissage}

\noindent
이 절에서는 가군의 1단계 데비사주가 무엇인지 설명한다.
% P05-EN-CH38-0012
1단계 데비사주는 밑과 표적 위에서 에탈 국소적으로 존재한다.
또한 1단계 데비사주의 기저변환, Zariski 축소, 에탈 국소화를 논한다.

\begin{definition}
\label{flat-definition-one-step-devissage}
$S$를 스킴이라 하자.
$X$가 $S$ 위에서 국소적으로 유한형이라고 하자.
$\mathcal{F}$를 유한형 준연접 $\mathcal{O}_X$-가군이라 하자.
$s \in S$를 점이라 하자.
{\it $s$ 위에서 $\mathcal{F}/X/S$의 1단계 데비사주}란
$S$ 위 스킴들의 사상
$$
\xymatrix{
X & Z \ar[l]_i \ar[r]^\pi & Y
}
$$
과 유한형 준연접 $\mathcal{O}_Z$-가군 $\mathcal{G}$로 주어지며,
다음을 만족하는 자료를 말한다.
\begin{enumerate}
\item $X$, $S$, $Z$, $Y$는 아핀이다.
\item $i$는 유한표현인 닫힌 몰입이다.
\item $\mathcal{F} \cong i_*\mathcal{G}$,
\item $\pi$는 유한이다.
\item 구조 사상 $Y \to S$는 매끄럽고 그 올들은 기하적으로 기약이며
차원이 $\dim(\text{Supp}(\mathcal{F}_s))$이다.
\end{enumerate}
이때 $(Z, Y, i, \pi, \mathcal{G})$를 $s$ 위에서
$\mathcal{F}/X/S$의 1단계 데비사주라고 한다.
\end{definition}

\noindent
그러한 1단계 데비사주는 $X$와 $S$가 아핀일 때에만 존재할 수 있음에
유의하자. 위 정의에서는 $X$가 $S$ 위에서 (국소적으로) 유한형일 것만을
요구하며, 아래에서도 이 설정으로 작업한다. \cite{GruRay}의 저자들은
% P05-EN-CH38-0013
$X \to S$가 국소 유한표현이고 $\mathcal{F}$가 유한형 준연접
$\mathcal{O}_X$-가군인 설정을 일관되게 사용한다. 이 선택의 장점은
$\mathcal{F}$가 $\mathcal{O}_X$-가군으로서 유한표현인지 묻는 것이
``의미가 있는'' 반면, 우리의 설정에서는 ``의미가 없다''는 점이다.
이에 관한 논의는 More on Morphisms, Section
\ref{more-morphisms-section-finite-type-finite-presentation}을 보라.
거기의 관찰에 따르면 우리의 설정에서는 $\mathcal{F}$가
``$S$에 대해 상대적으로 국소 유한표현''이라는 조건을 생각할 수 있고,
이 개념으로 일관되게 작업할 수도 있다. 그러나 우리는 그 대신 이 문제가
나타날 때마다 Lemma \ref{flat-lemma-devissage-finite-presentation}의 결과와
Remark \ref{flat-remark-finite-presentation}의 관찰을 사용하여
임시방편적으로 처리할 것이다.

\begin{definition}
\label{flat-definition-one-step-devissage-at-x}
$S$를 스킴이라 하자.
$X$가 $S$ 위에서 국소적으로 유한형이라고 하자.
$\mathcal{F}$를 유한형 준연접 $\mathcal{O}_X$-가군이라 하자.
$x \in X$의 $S$에서의 상을 $s$라 하자.
{\it $x$에서 $\mathcal{F}/X/S$의 1단계 데비사주}란
계 $(Z, Y, i, \pi, \mathcal{G}, z, y)$를 말한다. 여기서
$(Z, Y, i, \pi, \mathcal{G})$는 $s$ 위에서
$\mathcal{F}/X/S$의 1단계 데비사주이고, 다음을 만족한다.
\begin{enumerate}
\item $\dim_x(\text{Supp}(\mathcal{F}_s)) = \dim(\text{Supp}(\mathcal{F}_s))$,
\item $z \in Z$는 $i(z) = x$ 및 $\pi(z) = y$를 만족하는 점이다.
\item $\pi^{-1}(\{y\}) = \{z\}$이다.
\item 확대 $\kappa(y)/\kappa(s)$는 순수 초월이다.
\end{enumerate}
\end{definition}

\noindent
$x$에서 $\mathcal{F}/X/S$의 1단계 데비사주는 $X$와 $S$가 아핀일
때에만 존재할 수 있다. 조건 (1)과 Definition
\ref{flat-definition-one-step-devissage}의 조건 (5)에 의해
$Y \to S$의 상대차원은 $\dim_x(\text{Supp}(\mathcal{F}_s))$와 같다.

\begin{lemma}
\label{flat-lemma-elementary-devissage-variant}
% P05-EN-CH38-0014
$f : X \to S$를 국소적으로 유한형인 스킴 사상이라 하자.
$\mathcal{F}$를 유한형 준연접 $\mathcal{O}_X$-가군이라 하자.
$x \in X$의 $S$에서의 상을 $s = f(x)$라 하자.
그러면 점을 갖춘 스킴들의 가환 그림
$$
\xymatrix{
(X, x) \ar[d]_f & (X', x') \ar[l]^g \ar[d] \\
(S, s) & (S', s') \ar[l] \\
}
$$
이 존재하여 $(S', s') \to (S, s)$와 $(X', x') \to (X, x)$는
기본 에탈 근방이고, $g^*\mathcal{F}/X'/S'$는 $x'$에서
1단계 데비사주를 갖는다.
\end{lemma}

\begin{proof}
Definition \ref{flat-definition-one-step-devissage-at-x}와
Lemma \ref{flat-lemma-elementary-devissage}에서 바로 따른다.
\end{proof}

\begin{lemma}
\label{flat-lemma-base-change-one-step}
$S$, $X$, $\mathcal{F}$, $s$가 Definition
\ref{flat-definition-one-step-devissage}에서와 같다고 하자.
$(Z, Y, i, \pi, \mathcal{G})$를 $s$ 위에서
$\mathcal{F}/X/S$의 1단계 데비사주라 하자.
$(S', s') \to (S, s)$를 점을 갖춘 스킴들의 임의의 사상이라 하자.
$X', Z', Y', i', \pi'$를 각각 $S' \to S$를 통한
$X, Z, Y, i, \pi$의 기저변환이라 하자.
$\mathcal{F}'$를 $\mathcal{F}$의 $X'$로의 당김,
$\mathcal{G}'$를 $\mathcal{G}$의 $Z'$로의 당김이라 하자.
$S'$가 아핀이면 $(Z', Y', i', \pi', \mathcal{G}')$는
$s'$ 위에서 $\mathcal{F}'/X'/S'$의 1단계 데비사주이다.
\end{lemma}

\begin{proof}
아핀 스킴들의 섬유곱은 아핀이다. Schemes, Lemma
\ref{schemes-lemma-fibre-product-affines}를 보라.
기저변환은 닫힌 몰입, 유한표현 사상, 유한 사상, 매끄러운 사상,
기하적으로 기약인 올을 갖는 사상, 상대차원이 $n$인 사상을 보존한다.
Morphisms, Lemmas \ref{morphisms-lemma-base-change-closed-immersion},
\ref{morphisms-lemma-base-change-finite-presentation},
\ref{morphisms-lemma-base-change-finite},
\ref{morphisms-lemma-base-change-smooth},
\ref{morphisms-lemma-base-change-relative-dimension-d}, 및
More on Morphisms, Lemma
\ref{more-morphisms-lemma-base-change-fibres-geometrically-irreducible}를
보라. 유한 사상 $i$를 따른 직상은 기저변환과 가환하므로
$i'_*\mathcal{G}' \cong \mathcal{F}'$이다. Cohomology of Schemes,
Lemma \ref{coherent-lemma-affine-base-change}를 보라.
또한
$\dim(\text{Supp}(\mathcal{F}_s)) = \dim(\text{Supp}(\mathcal{F}'_{s'}))$
이다. 실제로 Morphisms, Lemma
\ref{morphisms-lemma-dimension-fibre-after-base-change}에 의해
$$
\text{Supp}(\mathcal{F}_s) \times_s s' = \text{Supp}(\mathcal{F}'_{s'}).
$$
이로써 보조정리가 증명된다.
\end{proof}

\begin{lemma}
\label{flat-lemma-base-change-one-step-at-x}
$S$, $X$, $\mathcal{F}$, $x$, $s$가 Definition
\ref{flat-definition-one-step-devissage-at-x}에서와 같다고 하자.
$(Z, Y, i, \pi, \mathcal{G}, z, y)$를 $x$에서
$\mathcal{F}/X/S$의 1단계 데비사주라 하자.
$(S', s') \to (S, s)$를 동형 $\kappa(s) = \kappa(s')$을 유도하는
점을 갖춘 스킴들의 사상이라 하자.
$(Z', Y', i', \pi', \mathcal{G}')$를 Lemma
\ref{flat-lemma-base-change-one-step}에서 구성한 것이라 하고,
$x' \in X'$ (각각 $z' \in Z'$, $y' \in Y'$)를
$x \in X$ (각각 $z \in Z$, $y \in Y$)와 $s' \in S'$ 모두로
가는 유일한 점이라 하자. $S'$가 아핀이면
$(Z', Y', i', \pi', \mathcal{G}', z', y')$는 $x'$에서
$\mathcal{F}'/X'/S'$의 1단계 데비사주이다.
\end{lemma}

\begin{proof}
Lemma \ref{flat-lemma-base-change-one-step}에 의해
$(Z', Y', i', \pi', \mathcal{G}')$는 $s'$ 위에서
$\mathcal{F}'/X'/S'$의 1단계 데비사주이다.
가정 $\kappa(s) = \kappa(s')$에 의해 올들
$X'_{s'}$, $Z'_{s'}$, $Y'_{s'}$가 각각 $X_s$, $Z_s$, $Y_s$와
동형이므로, $(Z', Y', i', \pi', \mathcal{G}', z', y')$는
Definition \ref{flat-definition-one-step-devissage-at-x}의 성질 (1)--(4)를
만족한다.
\end{proof}

\begin{definition}
\label{flat-definition-shrink}
$S$, $X$, $\mathcal{F}$, $x$, $s$가 Definition
\ref{flat-definition-one-step-devissage-at-x}에서와 같다고 하자.
$(Z, Y, i, \pi, \mathcal{G}, z, y)$를 $x$에서
$\mathcal{F}/X/S$의 1단계 데비사주라 하자. 이 상황의
{\it 표준 축소}란 표준 열린집합들
$S' \subset S$, $X' \subset X$, $Z' \subset Z$, $Y' \subset Y$로
주어지고, $s \in S'$, $x \in X'$, $z \in Z'$, $y \in Y'$이며,
$$
(Z', Y', i|_{Z'}, \pi|_{Z'}, \mathcal{G}|_{Z'}, z, y)
$$
가 $x$에서 $\mathcal{F}|_{X'}/X'/S'$의 1단계 데비사주인 것을 말한다.
\end{definition}

\begin{lemma}
\label{flat-lemma-shrink}
가정과 표기는 Definition \ref{flat-definition-shrink}에서와 같다고 하자.
다음이 성립한다.
\begin{enumerate}
\item
\label{flat-item-shrink-base}
$S' \subset S$가 $s$의 표준 열린 근방이면,
$X' = X_{S'}$, $Z' = Z_{S'}$, $Y' = Y_{S'}$로 두어
표준 축소를 얻는다.
\item
\label{flat-item-shrink-on-Y}
$W \subset Y$를 $y$의 표준 열린 근방이라 하자.
그러면 $Y' = W \times_S S'$인 표준 축소가 존재한다.
\item
\label{flat-item-shrink-on-X}
$U \subset X$를 $x$의 열린 근방이라 하자.
그러면 $X' \subset U$인 표준 축소가 존재한다.
\end{enumerate}
\end{lemma}

\begin{proof}
(1)은 Lemma \ref{flat-lemma-base-change-one-step-at-x}와, 아핀 스킴들의
사상에 의한 표준 열린집합의 역상이 표준 열린집합이라는 사실에서
바로 따른다. Algebra, Lemma \ref{algebra-lemma-spec-functorial}을 보라.

\medskip\noindent
(2)에서와 같이 $W \subset Y$라 하자. $Y \to S$는 매끄러우므로
열린 사상이다. Morphisms, Lemma
\ref{morphisms-lemma-smooth-open}을 보라. 따라서 $W$의 상에 포함되는
$s$의 표준 열린 근방 $S'$를 찾을 수 있다. 그러면
$W_{S'} \to S'$의 올들은 $S'$ 위에서 $Y \to S$의 올들의 공집합이 아닌
열린 부분스킴이므로 역시 기하적으로 기약이다.
$Y' = W_{S'}$ 및 $Z' = \pi^{-1}(Y')$로 두면
$Z' \subset Z$는 $z$의 표준 열린 근방이다.
$Z' = D(\overline{h})$가 되게 하는 함수
$\overline{h} \in \Gamma(Z, \mathcal{O}_Z)$를 택하자.
$i : Z \to X$가 닫힌 몰입이므로
$i^\sharp(h) = \overline{h}$인 함수
$h \in \Gamma(X, \mathcal{O}_X)$를 찾을 수 있다.
$X' = D(h) \subset X$로 둔다. 이렇게 하여 (2)의 표준 축소를 얻는다.

\medskip\noindent
(3)에서와 같이 $U \subset X$라 하자. $U$를 축소하여 $U$가 표준
열린집합이라고 가정해도 된다. More on Morphisms, Lemma
\ref{more-morphisms-lemma-finite-morphism-single-point-in-fibre}에 의해
$\pi^{-1}(W) \subset i^{-1}(U)$를 만족하는 $y$의 표준 열린 근방
$W \subset Y$가 존재한다. (2)를 적용하여
$Y' = W_{S'}$인 표준 축소 $X', S', Z', Y'$를 얻는다.
$Z' \subset \pi^{-1}(W) \subset i^{-1}(U)$이므로, 표준 축소를
정의하는 어느 조건도 깨뜨리지 않고 $X'$를 $X' \cap U$로 바꿀 수 있다
($U$도 표준 열린집합이므로 이는 여전히 표준 열린집합이다).
따라서 원하는 결론을 얻는다.
\end{proof}

\begin{lemma}
\label{flat-lemma-elementary-etale-neighbourhood}
$S$, $X$, $\mathcal{F}$, $x$, $s$가 Definition
\ref{flat-definition-one-step-devissage-at-x}에서와 같다고 하자.
$(Z, Y, i, \pi, \mathcal{G}, z, y)$를 $x$에서
$\mathcal{F}/X/S$의 1단계 데비사주라 하자. 점을 갖춘 스킴들의 가환 그림
$$
\xymatrix{
(Y, y) \ar[d] & (Y', y') \ar[l] \ar[d] \\
(S, s) & (S', s') \ar[l]
}
$$
이 주어지고 수평 화살표들이 기본 에탈 근방이라고 하자. 그러면
다음 성질들을 갖는 점을 갖춘 스킴들의 가환 그림
$$
\xymatrix{
& & (X'', x'') \ar[lld] \ar[d] & (Z'', z'') \ar[l] \ar[lld] \ar[d] \\
(X, x) \ar[d] & (Z, z) \ar[l] \ar[d] &
(S'', s'') \ar[lld] & (Y'', y'') \ar[lld] \ar[l] \\
(S, s) & (Y, y) \ar[l]
}
$$
이 존재한다.
\begin{enumerate}
\item $(S'', s'') \to (S', s')$는 기본 에탈 근방이고,
사상 $S'' \to S$는 합성 $S'' \to S' \to S$이다.
\item $Y''$은 $Y' \times_{S'} S''$의 열린 부분스킴이다.
\item $Z'' = Z \times_Y Y''$이다.
\item $(X'', x'') \to (X, x)$는 기본 에탈 근방이다.
\item $(Z'', Y'', i'', \pi'', \mathcal{G}'', z'', y'')$는
층 $\mathcal{F}''$의 $x''$에서의 1단계 데비사주이다.
\end{enumerate}
여기서 $\mathcal{F}''$ (각각 $\mathcal{G}''$)는 사상
$X'' \to X$ (각각 $Z'' \to Z$)에 의한
$\mathcal{F}$ (각각 $\mathcal{G}$)의 당김이고,
$i'' : Z'' \to X''$와 $\pi'' : Z'' \to Y''$는 그림에서와 같다.
\end{lemma}

\begin{proof}
$(S'', s'') \to (S', s')$를 $S''$이 아핀인 임의의 기본 에탈
근방이라 하자. $Y'' \subset Y' \times_{S'} S''$을
점 $y'' = (y', s'')$을 포함하는 임의의 아핀 열린 근방이라 하자.
그러면 (3)에 의해 아핀인 $(Z'', z'')$을 얻는다. 또한
$Z_{S''} \to X_{S''}$은 닫힌 몰입이고
$Z'' \to Z_{S''}$은 에탈 사상이다. 따라서
Lemma \ref{flat-lemma-lift-etale}를 적용할 수 있고,
% P05-EN-CH38-0015
$Z'' \cong X'' \times_{X_{S'}} Z_{S'}$을 만족하는 아핀 스킴들의
에탈 사상 $X'' \to X_{S'}$을 찾을 수 있다.
대응하는 닫힌 몰입을 $i'' : Z'' \to X''$라 쓰자.
$x'' = i''(z'')$로 두면 보조정리의 가환 그림을 얻는다.
성질 (1), (2), (3), (4)는 구성상 성립한다.
따라서 $(S'', s'') \to (S', s')$와 $Y''$을 알맞게 택했을 때
(5)가 성립함을 보이면 충분하다.

\medskip\noindent
먼저 첫 문단과 같은 $(S'', s'') \to (S', s')$와 $Y''$의
임의의 선택에 대해 성립하는 성질들을 열거하자.
구성상 $Z'' = X'' \times_X Z$이므로
$i''_*\mathcal{G}'' = \mathcal{F}''$이다
(표기는 보조정리의 명제에서와 같다). Cohomology of Schemes,
Lemma \ref{coherent-lemma-affine-base-change}를 보라.
$n = \dim(\text{Supp}(\mathcal{F}_s)) =
\dim_x(\text{Supp}(\mathcal{F}_s))$로 두자.
사상 $Y'' \to S''$은 상대차원이 $n$인 매끄러운 사상이다.
(실제로 $Y' \to S'$는 상대차원이 $n$인 에탈 사상과 매끄러운 사상의
합성 $Y' \to Y_{S'} \to S'$로서 상대차원이 $n$인 매끄러운
사상이고, 기저변환은 상대차원이 $n$인 매끄러운 사상을 보존한다.)
$\kappa(y'') = \kappa(y)$이고 $\kappa(s) = \kappa(s'')$이므로
$\kappa(y'')$는 $\kappa(s'')$의 순수 초월 확대이다.
올의 사상 $X''_{s''} \to X_s$는
$\kappa(s) = \kappa(s'')$ 위 아핀 스킴들의 에탈 사상으로서
점 $x''$을 점 $x$로 보내고 $\mathcal{F}_s$를
$\mathcal{F}''_{s''}$로 당긴다. 따라서
$$
\dim(\text{Supp}(\mathcal{F}''_{s''})) =
\dim(\text{Supp}(\mathcal{F}_s)) = n =
\dim_x(\text{Supp}(\mathcal{F}_s)) =
\dim_{x''}(\text{Supp}(\mathcal{F}''_{s''}))
$$
이다. 차원은 에탈 국소화 아래 불변이기 때문이다. Descent, Lemma
\ref{descent-lemma-dimension-at-point-local}을 보라.
$\pi'' : Z'' \to Y''$은 $\pi$의 기저변환이므로 $\pi''$은 유한이고,
$\kappa(y) = \kappa(y'')$이므로
% P05-EN-CH38-0016
$\pi^{-1}(\{y''\}) = \{z''\}$이다.

\medskip\noindent
이 시점에서 Definition \ref{flat-definition-one-step-devissage}의 조건 중
$Y'' \to S''$의 올들이 기하적으로 기약이라는 것만 아직 확인하지
않았다. 물론 일반적으로 이는 참이 아니며, 바로 여기서
$\kappa(s) \subset \kappa(y)$가 순수 초월이라는 사실을 사용한다.
$T \subset Y'_{s'}$를 $y' = (y, s')$을 포함하는
$Y'_{s'}$의 기약 성분이라 하자. $Y'_{s'}$은
$\kappa(s')$ 위의 매끄러운 스킴이므로 $T$는 그 열린 부분스킴임에
유의하자. Varieties, Lemma
\ref{varieties-lemma-geometrically-connected-if-connected-and-point}에 의해
$\kappa(s') = \kappa(s)$가
$\kappa(y') = \kappa(y)$에서 대수적으로 닫혀 있으므로
$T$는 기하적으로 연결이다. $T$가 매끄러우므로 $T$는 기하적으로 기약이다.
따라서 More on Morphisms, Lemma
\ref{more-morphisms-lemma-normal-morphism-irreducible}를 적용하여,
$Y'' \to S''$의 모든 올이 기하적으로 기약이고 $T = Y''_{s''}$인
기본 에탈 사상 $(S'', s'') \to (S', s')$와
아핀 열린집합 $Y'' \subset Y'_{S''}$을 찾을 수 있다.
(먼저 $Y''$을, 그다음 $S''$을) 축소하여
$Y''$와 $S''$가 모두 아핀이라고 가정해도 된다.
이로써 보조정리의 증명이 끝난다.
\end{proof}

\begin{lemma}
\label{flat-lemma-existence-alpha}
$S$, $X$, $\mathcal{F}$, $s$가 Definition
\ref{flat-definition-one-step-devissage}에서와 같다고 하자.
$(Z, Y, i, \pi, \mathcal{G})$를 $s$ 위에서
$\mathcal{F}/X/S$의 1단계 데비사주라 하자.
$\xi \in Y_s$를 (유일한) 일반점이라 하자.
그러면 정수 $r > 0$과 $\mathcal{O}_Y$-가군 사상
$$
\alpha : \mathcal{O}_Y^{\oplus r} \longrightarrow \pi_*\mathcal{G}
$$
이 존재하여
$$
\alpha :
\kappa(\xi)^{\oplus r}
\longrightarrow
(\pi_*\mathcal{G})_\xi \otimes_{\mathcal{O}_{Y, \xi}} \kappa(\xi)
$$
는 동형이다. 더욱이 이 경우
$$
\dim(\text{Supp}(\Coker(\alpha)_s)) < \dim(\text{Supp}(\mathcal{F}_s)).
$$
이다.
\end{lemma}

\begin{proof}
가정에 의해 스킴 $S$와 $Y$는 아핀이다.
$S = \Spec(A)$, $Y = \Spec(B)$로 쓰자.
$\pi$가 유한이므로 $\mathcal{O}_Y$-가군
$\pi_*\mathcal{G}$는 유한형 준연접 $\mathcal{O}_Y$-가군이다.
따라서 어떤 유한 $B$-가군 $N$에 대해
$\pi_*\mathcal{G} = \widetilde{N}$이다.
$\mathfrak p \subset B$를 $\xi$에 대응하는 소 아이디얼이라 하자.
$\alpha$를 얻기 위해
$r = \dim_{\kappa(\mathfrak p)} N \otimes_B \kappa(\mathfrak p)$로
두고, $N \otimes_B \kappa(\mathfrak p)$의 기저를 이루는
$x_1, \ldots, x_r \in N$을 택한다.
$\alpha : B^{\oplus r} \to N$을
식 $\alpha(b_1, \ldots, b_r) = \sum b_ix_i$로 주어지는 사상으로 택하자.
그러면
$\alpha : \kappa(\mathfrak p)^{\oplus r} \to
N \otimes_B \kappa(\mathfrak p)$가 원하는 대로 동형임은 명백하다.
마지막으로 $\alpha$가 이 성질을 갖는 임의의 사상이라고 하자.
그러면 $N' = \Coker(\alpha)$는
$N' \otimes \kappa(\mathfrak p) = 0$인 유한 $B$-가군이다.
Nakayama 보조정리(Algebra, Lemma \ref{algebra-lemma-NAK})에 의해
$N'_{\mathfrak p} = 0$이다. 올 $Y_s$는 일반점 $\xi$를 갖는,
% P05-EN-CH38-0017
차원이 $n$인 기하적으로 기약인 스킴이고, 방금 $\xi$가
$\Coker(\alpha)$의 지지에 속하지 않음을 보였으므로 보조정리의
마지막 주장이 성립한다.
\end{proof}


\section{완전 데비사주}
\label{flat-section-complete-devissage}

\noindent
% P05-EN-CH38-0018: 문법상 유일하게 정해지는 보정은 부록 원장에 기록했다.
이 절에서는 가군의 완전 데비사주가 무엇인지 설명하고 그러한
데비사주가 존재함을 증명한다. 이 절의 내용은 주로 자료를 정리하는 일이다.

\begin{definition}
\label{flat-definition-complete-devissage}
$S$를 스킴이라 하자.
$X$가 $S$ 위에서 국소 유한형이라고 하자.
$\mathcal{F}$를 유한형 준연접 $\mathcal{O}_X$-가군이라 하자.
$s \in S$를 점이라 하자.
{\it $s$ 위에서 $\mathcal{F}/X/S$의 완전 데비사주}란
다음 그림을 이루는 $S$ 위의 스킴들
$$
\xymatrix{
X & Z_1 \ar[l]^{i_1} \ar[d]^{\pi_1} \\
& Y_1 & Z_2 \ar[l]^{i_2} \ar[d]^{\pi_2} \\
& & Y_2 & Z_3 \ar[l] \ar[d] \\
& & & ... & ... \ar[l] \ar[d] \\
& & & & Y_n
}
$$
과 유한형 준연접 $\mathcal{O}_{Z_k}$-가군 $\mathcal{G}_k$,
그리고 $\mathcal{O}_{Y_k}$-가군 사상
$$
\alpha_k :
\mathcal{O}_{Y_k}^{\oplus r_k}
\longrightarrow
\pi_{k, *}\mathcal{G}_k,
\quad
k = 1, \ldots, n
$$
으로 주어지며 다음 성질들을 만족하는 것을 말한다.
\begin{enumerate}
\item $(Z_1, Y_1, i_1, \pi_1, \mathcal{G}_1)$은 $s$ 위에서
$\mathcal{F}/X/S$의 1단계 데비사주이다.
\item 사상 $\alpha_k$는 동형
$$
\kappa(\xi_k)^{\oplus r_k} \longrightarrow
(\pi_{k, *}\mathcal{G}_k)_{\xi_k}
\otimes_{\mathcal{O}_{Y_k, \xi_k}} \kappa(\xi_k)
$$
을 유도한다. 여기서 $\xi_k \in (Y_k)_s$는 유일한 일반점이다.
\item $k = 2, \ldots, n$에 대하여 계
$(Z_k, Y_k, i_k, \pi_k, \mathcal{G}_k)$는 $s$ 위에서
$\Coker(\alpha_{k - 1})/Y_{k - 1}/S$의 1단계 데비사주이다.
\item $\Coker(\alpha_n) = 0$.
\end{enumerate}
이때
$(Z_k, Y_k, i_k, \pi_k, \mathcal{G}_k, \alpha_k)_{k = 1, \ldots, n}$
을 $s$ 위에서 $\mathcal{F}/X/S$의 완전 데비사주라고 한다.
\end{definition}

\begin{definition}
\label{flat-definition-complete-devissage-at-x}
$S$를 스킴이라 하자.
$X$가 $S$ 위에서 국소 유한형이라고 하자.
$\mathcal{F}$를 유한형 준연접 $\mathcal{O}_X$-가군이라 하자.
$x \in X$를 그 상이 $s \in S$인 점이라 하자.
{\it $x$에서 $\mathcal{F}/X/S$의 완전 데비사주}란 계
$$
(Z_k, Y_k, i_k, \pi_k, \mathcal{G}_k, \alpha_k, z_k, y_k)_{k = 1, \ldots, n}
$$
으로서 $(Z_k, Y_k, i_k, \pi_k, \mathcal{G}_k, \alpha_k)$가
$s$ 위에서 $\mathcal{F}/X/S$의 완전 데비사주이고 다음을 만족하는 것을 말한다.
\begin{enumerate}
\item $(Z_1, Y_1, i_1, \pi_1, \mathcal{G}_1, z_1, y_1)$은
$x$에서 $\mathcal{F}/X/S$의 1단계 데비사주이다.
\item $k = 2, \ldots, n$에 대하여 계
$(Z_k, Y_k, i_k, \pi_k, \mathcal{G}_k, z_k, y_k)$는
$y_{k - 1}$에서 $\Coker(\alpha_{k - 1})/Y_{k - 1}/S$의
1단계 데비사주이다.
\end{enumerate}
\end{definition}

\noindent
다시 말해, 완전 데비사주가 존재할 수 있는 것은 $X$와 $S$가
아핀인 경우뿐임에 유의하자.

\begin{lemma}
\label{flat-lemma-base-change-complete}
$S$, $X$, $\mathcal{F}$, $s$가
Definition \ref{flat-definition-complete-devissage}에서와 같다고 하자.
$(S', s') \to (S, s)$를 점을 갖춘 스킴들의 임의의 사상이라 하자.
$(Z_k, Y_k, i_k, \pi_k, \mathcal{G}_k, \alpha_k)_{k = 1, \ldots, n}$
을 $s$ 위에서 $\mathcal{F}/X/S$의 완전 데비사주라 하자.
이 데이터에 대하여 $X', Z'_k, Y'_k, i'_k, \pi'_k$를 각각
$S' \to S$에 의한 $X, Z_k, Y_k, i_k, \pi_k$의 기저변환이라 하자.
$\mathcal{F}'$를 $\mathcal{F}$의 $X'$로의 당김이라 하고
$\mathcal{G}'_k$를 $\mathcal{G}_k$의 $Z'_k$로의 당김이라 하자.
$\alpha'_k$를 $\alpha_k$의 $Y'_k$로의 당김이라 하자.
$S'$가 아핀이면
$(Z'_k, Y'_k, i'_k, \pi'_k, \mathcal{G}'_k, \alpha'_k)_{k = 1, \ldots, n}$
은 $s'$ 위에서 $\mathcal{F}'/X'/S'$의 완전 데비사주이다.
\end{lemma}

\begin{proof}
Lemma \ref{flat-lemma-base-change-one-step}에 의해 1단계 데비사주의
기저변환은 1단계 데비사주이다. 따라서 $\Coker(\alpha_k)$의 형성이
기저변환과 가환하고 Definition \ref{flat-definition-complete-devissage}의
조건 (2)가 기저변환으로 보존됨을 보이면 충분하다. 첫 번째 주장은
$\pi'_{k, *}\mathcal{G}'_k$가 $\pi_{k, *}\mathcal{G}_k$의
당김이고(Cohomology of Schemes, Lemma
\ref{coherent-lemma-affine-base-change}) $\otimes$가 우완전이므로 성립한다.
두 번째 주장도 같은 이유로
$$
(\pi_{k, *}\mathcal{G}_k)_{\xi_k}
\otimes_{\mathcal{O}_{Y_k, \xi_k}} \kappa(\xi_k)
\otimes_{\kappa(\xi_k)} \kappa(\xi'_k)
\cong
(\pi'_{k, *}\mathcal{G}'_k)_{\xi'_k}
\otimes_{\mathcal{O}_{Y'_k, \xi'_k}} \kappa(\xi'_k)
$$
가 성립하기 때문이다. 기호의 뜻은 명백하다.
\end{proof}

\begin{lemma}
\label{flat-lemma-base-change-complete-at-x}
$S$, $X$, $\mathcal{F}$, $x$, $s$가 Definition
\ref{flat-definition-complete-devissage-at-x}에서와 같다고 하자.
$(S', s') \to (S, s)$를 동형 $\kappa(s) = \kappa(s')$를 유도하는
점을 갖춘 스킴들의 사상이라 하자.
$(Z_k, Y_k, i_k, \pi_k, \mathcal{G}_k, \alpha_k, z_k, y_k)_{k = 1, \ldots, n}$
을 $x$에서 $\mathcal{F}/X/S$의 완전 데비사주라 하자.
$(Z'_k, Y'_k, i'_k, \pi'_k, \mathcal{G}'_k, \alpha'_k)_{k = 1, \ldots, n}$
을 Lemma \ref{flat-lemma-base-change-complete}에서 구성한 것이라 하고,
% P05-EN-CH38-0019: 첨자가 빠진 표적 스킴을 동결 원문대로 보존했다.
$x' \in X'$ (각각 $z'_k \in Z'$, $y'_k \in Y'$)를
$x \in X$ (각각 $z_k \in Z_k$, $y_k \in Y_k$)와
$s' \in S'$ 양쪽으로 가는 유일한 점이라 하자.
$S'$가 아핀이면
$(Z'_k, Y'_k, i'_k, \pi'_k, \mathcal{G}'_k, \alpha'_k,
z'_k, y'_k)_{k = 1, \ldots, n}$
은 $x'$에서 $\mathcal{F}'/X'/S'$의 완전 데비사주이다.
\end{lemma}

\begin{proof}
Lemma \ref{flat-lemma-base-change-complete}와
Lemma \ref{flat-lemma-base-change-one-step-at-x}를 함께 적용하면 된다.
\end{proof}

\begin{definition}
\label{flat-definition-shrink-complete}
$S$, $X$, $\mathcal{F}$, $x$, $s$가 Definition
\ref{flat-definition-complete-devissage-at-x}에서와 같다고 하자.
$x$에서 $\mathcal{F}/X/S$의 완전 데비사주
$(Z_k, Y_k, i_k, \pi_k, \mathcal{G}_k, \alpha_k, z_k, y_k)_{k = 1, \ldots, n}$
을 생각하자. 이 상황의 {\it 표준 축소}란 다음과 같은 것을 말한다.
% P05-EN-CH38-0020: 정의되지 않은 점 첨자와 첨자 없는 열린집합을 동결 원문대로 보존했다.
표준 열린집합 $S' \subset S$, $X' \subset X$,
$Z'_k \subset Z_k$, $Y'_k \subset Y_k$가 주어져
$s_k \in S'$, $x_k \in X'$, $z_k \in Z'$, $y_k \in Y'$이고,
$$
(Z'_k, Y'_k, i'_k, \pi'_k,
\mathcal{G}'_k, \alpha'_k, z_k, y_k)_{k = 1, \ldots, n}
$$
% P05-EN-CH38-0021: 1단계/완전 데비사주 불일치를 동결 원문대로 보존했다.
가 $x$에서 $\mathcal{F}'/X'/S'$의 1단계 데비사주인 것을 말한다. 여기서
$\mathcal{G}'_k = \mathcal{G}_k|_{Z'_k}$이고
$\mathcal{F}' = \mathcal{F}|_{X'}$이다.
\end{definition}

\begin{lemma}
\label{flat-lemma-shrink-complete}
Definition \ref{flat-definition-shrink-complete}의 가정과 기호 아래에서
다음이 성립한다.
\begin{enumerate}
\item
\label{flat-item-shrink-base-complete}
$S' \subset S$가 $s$의 표준 열린 근방이면
$X' = X_{S'}$, $Z'_k = Z_{S'}$, $Y'_k = Y_{S'}$로 두어
표준 축소를 얻는다.
\item
\label{flat-item-shrink-on-Y-complete}
% P05-EN-CH38-0022: 정의되지 않은 점 y를 동결 원문대로 보존했다.
$W \subset Y_n$을 $y$의 표준 열린 근방이라 하자.
그러면 $Y'_n = W \times_S S'$인 표준 축소가 존재한다.
\item
\label{flat-item-shrink-on-X-complete}
$U \subset X$를 $x$의 열린 근방이라 하자.
그러면 $X' \subset U$인 표준 축소가 존재한다.
\end{enumerate}
\end{lemma}

\begin{proof}
(1)은 Lemma \ref{flat-lemma-base-change-complete-at-x}와
Lemma \ref{flat-lemma-shrink}에서 바로 따른다.

\medskip\noindent
(2)의 증명. 편의를 위해 $X = Y_0$로 쓰자.
Lemma \ref{flat-lemma-shrink} (\ref{flat-item-shrink-on-Y})를 적용하여
$y_{n - 1}$에서 $\Coker(\alpha_{n - 1})/Y_{n - 1}/S$의
1단계 데비사주의 표준 축소
$S', Y'_{n - 1}, Z'_n, Y'_n$
으로서 $Y'_n = W \times_S S'$인 것을 찾는다. 이 절차를 반복하여
$y_{n - 2}$에서 $\Coker(\alpha_{n - 2})/Y_{n - 2}/S$의
1단계 데비사주의 표준 축소
$S'', Y''_{n - 2}, Z''_{n - 1}, Y''_{n - 1}$
으로서 $Y''_{n - 1} = Y'_{n - 1} \times_S S''$인 것을 찾을 수 있다.
이와 같이 계속하여
$S^{(n)}, Y^{(n)}_0, Z^{(n)}_1, Y^{(n)}_1$을 얻는다.
이때 $S^{(n)}$, $X^{(n)}$, $Z_k^{(n - k)} \times_S S^{(n)}$,
$Y_k^{(n - k)} \times_S S^{(n)}$을 취하면 원하는 성질을 갖는
표준 축소를 얻음이 명백하다.

\medskip\noindent
(3)의 증명. 완전 데비사주의 길이에 대해 귀납법을 쓴다.
먼저 Lemma \ref{flat-lemma-shrink} (\ref{flat-item-shrink-on-X})를 적용하여
$x$에서 $\mathcal{F}/X/S$의 1단계 데비사주의 표준 축소
$S', X', Z'_1, Y'_1$으로서 $X' \subset U$인 것을 찾는다.
$n = 1$이면 끝난다. $n > 1$이면 귀납가정에 의해
% P05-EN-CH38-0023: 잘못 지정된 점 x를 동결 원문대로 보존했다.
$x$에서 $\Coker(\alpha_1)/Y_1/S$의 완전 데비사주
$(Z_k, Y_k, i_k, \pi_k, \mathcal{G}_k, \alpha_k, z_k, y_k)_{k = 2, \ldots, n}$
의 표준 축소 $S''$, $Y''_1$, $Z''_k$, $Y''_k$로서
$Y''_1 \subset Y'_1$인 것을 찾을 수 있다.
Lemma \ref{flat-lemma-shrink} (\ref{flat-item-shrink-on-Y})를 쓰면
$S''' \subset S'$, $X''' \subset X'$, $Z'''_1$과
$Y'''_1 = Y''_1 \times_S S'''$으로 이루어진 표준 축소를 찾을 수 있다.
우리 문제의 해는
$$
S''', X''', Z'''_1, Y'''_1, Z''_2 \times_S S''',
Y''_2 \times_S S''', \ldots, Z''_n \times_S S''', Y''_n \times_S S'''
$$
을 취하면 된다. 이로써 보조정리의 증명이 끝난다.
\end{proof}

\begin{proposition}
\label{flat-proposition-existence-complete-at-x}
$S$를 스킴이라 하자.
$X$가 $S$ 위에서 국소 유한형이라고 하자.
$x \in X$를 그 상이 $s \in S$인 점이라 하자.
% P05-EN-CH38-0024: F의 선언이 빠진 것을 동결 원문대로 보존했다.
다음과 같은 점을 갖춘 스킴들의 가환 그림
$$
\xymatrix{
(X, x) \ar[d] & (X', x') \ar[l]^g \ar[d] \\
(S, s) & (S', s') \ar[l]
}
$$
으로서 수평 화살표들이 기본 에탈 근방이고
% P05-EN-CH38-0025: 당긴 층의 데비사주가 x'가 아닌 x에 있다고 한 것을 동결 원문대로 보존했다.
$g^*\mathcal{F}/X'/S'$가 $x$에서 완전 데비사주를 갖는 것이 존재한다.
\end{proposition}

\begin{proof}
정수 $d = \dim_x(\text{Supp}(\mathcal{F}_s))$에 대해 귀납법으로 증명한다.
Lemma \ref{flat-lemma-elementary-devissage-variant}에 의해
다음과 같은 점을 갖춘 스킴들의 그림
$$
\xymatrix{
(X, x) \ar[d] & (X', x') \ar[l]^g \ar[d] \\
(S, s) & (S', s') \ar[l]
}
$$
으로서 수평 화살표들이 기본 에탈 근방이고
$g^*\mathcal{F}/X'/S'$가 $x'$에서 1단계 데비사주를 갖는 것이 존재한다.
문제의 국소적 성질에 의해 $(X, x) \to (S, s)$를
$(X', x') \to (S', s')$로 바꾸어도 된다. 따라서 그렇게 한 뒤에는
$x$에서 $\mathcal{F}/X/S$의 1단계 데비사주
$(Z_1, Y_1, i_1, \pi_1, \mathcal{G}_1)$이 존재한다고 가정해도 된다.

\medskip\noindent
Lemma \ref{flat-lemma-existence-alpha}를 적용하여
$$
\alpha_1 :
\mathcal{O}_{Y_1}^{\oplus r_1}
\longrightarrow
\pi_{1, *}\mathcal{G}_1
$$
인 사상을 찾는다. 이 사상은 $\kappa(\xi_1)$ 위의 벡터 공간들의
동형을 유도한다. 여기서 $\xi_1 \in Y_1$은 $s$ 위의 $Y_1$의 올의
유일한 일반점이다. 더욱이
$\dim_{y_1}(\text{Supp}(\Coker(\alpha_1)_s)) < d$이다.
$\Coker(\alpha_1)_s$의 $y_1$에서의 줄기가 영일 수 있다.
이 경우 Lemma \ref{flat-lemma-shrink} (\ref{flat-item-shrink-on-Y})에 의해
$Y_1$을 축소하여 $\Coker(\alpha_1) = 0$이라고 가정할 수 있으므로
% P05-EN-CH38-0026: 길이 0/1 불일치를 동결 원문대로 보존했다.
길이가 0인 완전 데비사주를 얻는다.

\medskip\noindent
이제 $\Coker(\alpha_1)_s$의 $y_1$에서의 줄기가 영이 아니라고 하자.
이 경우 귀납가정에 의해 다음과 같은 가환 그림
\begin{equation}
\label{flat-equation-overcome-this}
\vcenter{
\xymatrix{
(Y_1, y_1) \ar[d] & (Y'_1, y'_1) \ar[l]^h \ar[d] \\
(S, s) & (S', s') \ar[l]
}
}
\end{equation}
으로서 수평 화살표들이 기본 에탈 근방이고
$h^*\Coker(\alpha_1)/Y'_1/S'$가 $y'_1$에서 완전 데비사주
$$
(Z_k, Y_k, i_k, \pi_k, \mathcal{G}_k, \alpha_k, z_k, y_k)_{k = 2, \ldots, n}
$$
를 갖는 것이 존재한다. (특히
% P05-EN-CH38-0027: 닫힌 몰입의 잘못된 표적 Y'_2를 동결 원문대로 보존했다.
$i_2 : Z_2 \to Y'_1$은 $Y'_2$로의 닫힌 몰입이다.)
이제 Lemma \ref{flat-lemma-elementary-etale-neighbourhood}를
$S, X, \mathcal{F}, x, s$, 계
$(Z_1, Y_1, i_1, \pi_1, \mathcal{G}_1)$ 및 그림
(\ref{flat-equation-overcome-this})에 적용한다. 그러면 다음 그림
$$
\xymatrix{
& & (X'', x'') \ar[lld] \ar[d] & (Z''_1, z''_1) \ar[l] \ar[lld] \ar[d] \\
(X, x) \ar[d] & (Z_1, z_1) \ar[l] \ar[d] &
(S'', s'') \ar[lld] & (Y''_1, y''_1) \ar[lld] \ar[l] \\
(S, s) & (Y_1, y_1) \ar[l]
}
$$
으로서 인용한 보조정리에 열거된 모든 성질을 갖는 것을 얻는다.
특히 $Y''_1 \subset Y'_1 \times_{S'} S''$이다.
$X_1 = Y'_1 \times_{S'} S''$로 두고
$\mathcal{F}_1$을 $\Coker(\alpha_1)$의 당김이라 하자.
Lemma \ref{flat-lemma-base-change-complete-at-x}에 의해 계
\begin{equation}
\label{flat-equation-shrink-this}
(Z_k \times_{S'} S'',
Y_k \times_{S'} S'', i''_k, \pi''_k, \mathcal{G}''_k,
\alpha''_k, z''_k, y''_k)_{k = 2, \ldots, n}
\end{equation}
% P05-EN-CH38-0028: 불완전하게 쓰인 상대 데이터를 동결 원문대로 보존했다.
은 $\mathcal{F}_1$에서 $X_1$로의 완전 데비사주이다.
다시 문제의 성질에 의해 $(X, x) \to (S, s)$를
$(X'', x'') \to (S'', s'')$로 바꾸어도 된다.
이로부터 다음을 가정해도 됨을 알 수 있다.
\begin{enumerate}
\item[(a)] $x$에서 $\mathcal{F}/X/S$의 1단계 데비사주
$(Z_1, Y_1, i_1, \pi_1, \mathcal{G}_1)$이 존재한다.
\item[(b)] $\alpha_1 : \mathcal{O}_{Y_1}^{\oplus r_1}
\to \pi_{1, *}\mathcal{G}_1$로서
% P05-EN-CH38-0029: alpha/alpha_1 불일치를 동결 원문대로 보존했다.
$\alpha \otimes \kappa(\xi_1)$가 동형인 것이 존재한다.
\item[(c)] $Y_1 \subset X_1$은 열린집합이고 $y_1 = x_1$이며
$\mathcal{F}_1|_{Y_1} \cong \Coker(\alpha_1)$이다.
\item[(d)] $x_1$에서 $\mathcal{F}_1/X_1/S$의 완전 데비사주
$(Z_k, Y_k, i_k, \pi_k, \mathcal{G}_k, \alpha_k, z_k, y_k)_{k = 2, \ldots, n}$
가 존재한다.
\end{enumerate}
증명을 끝내려면 1단계 데비사주와 완전 데비사주를 축소하여
이들이 하나의 완전 데비사주로 이어지게 하면 된다.
(이어지는 증명을 읽는 대신 Lemma \ref{flat-lemma-shrink}와
Lemma \ref{flat-lemma-shrink-complete}를 사용하여 독자가 직접 해 보기를 권한다.)
$Y_1 \subset X_1$은 $x_1$의 열린 근방이므로
Lemma \ref{flat-lemma-shrink-complete} (\ref{flat-item-shrink-on-X-complete})를
적용하여 데이터 (d)의 표준 축소
$S', X'_1, Z'_2, Y'_2, \ldots, Y'_n$으로서
$X'_1 \subset Y_1$인 것을 찾는다. $X'_1$은 아핀 스킴 $Y_1$의
표준 열린집합이기도 함에 유의하자.
다음으로 데이터 (a)를 다음과 같이 축소한다. 먼저
Lemma \ref{flat-lemma-shrink} (\ref{flat-item-shrink-base})에 의해
밑 $S$를 $S'$로 축소하고, 이어서
Lemma \ref{flat-lemma-shrink} (\ref{flat-item-shrink-on-Y})를 사용하여 그 결과를
$S''$, $X''$, $Z''_1$, $Y''_1$로 축소하되, 마침내
$Y''_1 = X'_1 \times_S S''$이고 $S'' \subset S'$이 되게 한다.
그러면
$$
Z''_1, Y''_1, Z'_2 \times_{S'} S'', Y'_2 \times_{S'} S'', \ldots,
Y'_n \times_{S'} S''
$$
이 우리가 찾던 완전 데비사주를 줌을 알 수 있다.
\end{proof}

\noindent
자료를 조금 더 정리하면 다음 결과를 얻는다.

\begin{lemma}
\label{flat-lemma-existence-complete}
$X \to S$를 스킴의 유한형 사상이라 하자.
$\mathcal{F}$를 유한형 준연접 $\mathcal{O}_X$-가군이라 하자.
$s \in S$를 점이라 하자.
기본 에탈 근방 $(S', s') \to (S, s)$와 에탈 사상들
$h_i : Y_i \to X_{S'}$, $i = 1, \ldots, n$이 존재하여, 각
$i$에 대하여 $s'$ 위에서 $\mathcal{F}_i/Y_i/S'$의 완전 데비사주가
존재한다. 여기서 $\mathcal{F}_i$는 $\mathcal{F}$의 $Y_i$로의 당김이며,
또한 $X_s = (X_{S'})_{s'} \subset \bigcup h_i(Y_i)$이다.
\end{lemma}

\begin{proof}
모든 점 $x \in X_s$에 대하여 다음과 같은 그림
$$
\xymatrix{
(X, x) \ar[d] & (X', x') \ar[l]^g \ar[d] \\
(S, s) & (S', s') \ar[l]
}
$$
으로서 수평 화살표들이 기본 에탈 근방이고
$g^*\mathcal{F}/X'/S'$가 $x'$에서 완전 데비사주를 갖는 것을 찾을 수 있다.
$X \to S$가 유한형이므로 올 $X_s$는 준콤팩트하고, 위의 각
$g : X' \to X$가 열린 사상이므로 유한 개의 $g(X'_{s'})$로
$X_s$를 덮을 수 있다. 따라서 다음과 같은 그림들의 유한족
$$
\vcenter{
\xymatrix{
% P05-EN-CH38-0030: 첨자가 빠진 구별된 점 x를 동결 원문대로 보존했다.
(X, x) \ar[d] & (X'_i, x'_i) \ar[l]^{g_i} \ar[d] \\
(S, s) & (S'_i, s'_i) \ar[l]
}
}
\quad i = 1, \ldots, n
$$
% P05-EN-CH38-0031: 올의 상이 아닌 전체 열린 상과의 등식을 동결 원문대로 보존했다.
으로서 $X_s = \bigcup g_i(X'_i)$인 것을 찾을 수 있다.
$S' = S'_1 \times_S \ldots \times_S S'_n$로 두고
% P05-EN-CH38-0032: 정의되지 않은 X_i를 동결 원문대로 보존했다.
$Y_i = X_i \times_{S'_i} S'$를 $X'_i$의 $S'$로의 기저변환이라 하자.
Lemma \ref{flat-lemma-base-change-complete}에 의해 $\mathcal{F}$의
$Y_i$로의 당김은
% P05-EN-CH38-0033: s/s' 불일치를 동결 원문대로 보존했다.
$s$ 위에서 완전 데비사주를 가지며, 이로써 증명이 끝난다.
\end{proof}



\section{대수로의 번역}
\label{flat-section-translation}

\noindent
완전 데비사주를 갖는다는 것이 무엇을 뜻하는지 대수적으로 명시해 두는 것이
유용할 수 있다. 다음 개념을 도입한다. (그다지 유용한 개념은 아니므로
터무니없이 긴 이름을 붙인다.)

\begin{definition}
\label{flat-definition-elementary-etale-neighbourhood}
$R \to S$를 환 준동형이라 하자. $\mathfrak q$를 $S$의 소 아이디얼로서
$R$의 소 아이디얼 $\mathfrak p$ 위에 놓인 것이라 하자.
{\it $\mathfrak q$에서 환 준동형 $R \to S$의 기본 에탈 국소화}란
환들의 가환 그림과 그에 딸린 소 아이디얼들
$$
\xymatrix{
S \ar[r] & S' \\
R \ar[u] \ar[r] & R' \ar[u]
}
\quad\quad
\xymatrix{
\mathfrak q \ar@{-}[r] & \mathfrak q' \\
\mathfrak p \ar@{-}[u] \ar@{-}[r] & \mathfrak p' \ar@{-}[u]
}
$$
로 주어지며, $R \to R'$와 $S \to S'$가 에탈 환 준동형이고
$\kappa(\mathfrak p) = \kappa(\mathfrak p')$ 및
$\kappa(\mathfrak q) = \kappa(\mathfrak q')$를 만족하는 것을 말한다.
\end{definition}

\begin{definition}
\label{flat-definition-complete-devissage-algebra}
$R \to S$를 유한형 환 준동형이라 하자.
$\mathfrak r$을 $R$의 소 아이디얼이라 하자.
$N$을 유한 $S$-가군이라 하자.
{\it $\mathfrak r$ 위에서 $N/S/R$의 완전 데비사주}란
$R$-대수 준동형들
$$
\xymatrix{
& A_1 & & A_2 & & ... & & A_n \\
S \ar[ru] & & B_1 \ar[lu] \ar[ru] & & ... \ar[lu] \ar[ru] & &
... \ar[lu] \ar[ru] & & B_n \ar[lu]
}
$$
과 유한 $A_i$-가군 $M_i$, 그리고 $B_i$-가군 사상
$\alpha_i : B_i^{\oplus r_i} \to M_i$로 주어지며 다음을 만족한다.
\begin{enumerate}
\item $S \to A_1$은 전사이고 유한표현이다.
% P05-EN-CH38-0034: 조건 (2), (6)의 빠진 i<n 범위를 동결 원문대로 보존했다.
\item $B_i \to A_{i + 1}$은 전사이고 유한표현이다.
\item $B_i \to A_i$는 유한이다.
\item $R \to B_i$는 매끄럽고 기하적으로 기약인 올들을 갖는다.
\item $S$-가군으로서 $N \cong M_1$이다.
\item $B_i$-가군으로서 $\Coker(\alpha_i) \cong M_{i + 1}$이다.
\item $\alpha_i : \kappa(\mathfrak p_i)^{\oplus r_i}
\to M_i \otimes_{B_i} \kappa(\mathfrak p_i)$는 동형이다. 여기서
$\mathfrak p_i = \mathfrak rB_i$이다.
\item $\Coker(\alpha_n) = 0$.
\end{enumerate}
이 상황에서
$(A_i, B_i, M_i, \alpha_i)_{i = 1, \ldots, n}$
을 $\mathfrak r$ 위에서 $N/S/R$의 완전 데비사주라고 한다.
\end{definition}

\begin{remark}
\label{flat-remark-finite-presentation}
모든 $i$에 대한 $R$-대수 $B_i$와 $i \geq 2$에 대한 $A_i$는
$R$ 위에서 유한표현임에 유의하자. $S$가 $R$ 위에서 유한표현이면
$A_1$도 $R$ 위에서 유한표현이다. 이 경우 완전 데비사주에 나오는
모든 환 준동형은 유한표현이다. Algebra, Lemma
\ref{algebra-lemma-compose-finite-type}를 보라.
계속해서 $S$가 $R$ 위에서 유한표현이라고 가정하면 다음은 서로 동치이다.
\begin{enumerate}
% P05-EN-CH38-0035: 정의되지 않은 M/N 불일치를 동결 원문대로 보존했다.
\item $M$은 $S$ 위에서 유한표현이다.
\item $M_1$은 $A_1$ 위에서 유한표현이다.
\item $M_1$은 $B_1$ 위에서 유한표현이다.
\item 각 $M_i$는 $A_i$-가군으로서도 $B_i$-가군으로서도 유한표현이다.
\end{enumerate}
동치 (1) $\Leftrightarrow$ (2)와 (2) $\Leftrightarrow$ (3)은
Algebra, Lemma \ref{algebra-lemma-finite-finitely-presented-extension}에서
따른다. $M_1$이 유한표현이면 $\Coker(\alpha_1)$도 유한표현이고
(Algebra, Lemma \ref{algebra-lemma-extension}를 보라),
따라서 $M_2$ 등도 그러하다.
\end{remark}

\begin{definition}
\label{flat-definition-complete-devissage-at-x-algebra}
$R \to S$를 유한형 환 준동형이라 하자.
$\mathfrak q$를 $S$의 소 아이디얼로서 $R$의 소 아이디얼
$\mathfrak r$ 위에 놓인 것이라 하자.
$N$을 유한 $S$-가군이라 하자.
{\it $\mathfrak q$에서 $N/S/R$의 완전 데비사주}란
$\mathfrak r$ 위에서 $N/S/R$의 완전 데비사주
$(A_i, B_i, M_i, \alpha_i)_{i = 1, \ldots, n}$과
$\mathfrak r$ 위에 놓인 소 아이디얼들 $\mathfrak q_i \subset B_i$로
주어지며 다음을 만족하는 것을 말한다.
\begin{enumerate}
\item $\kappa(\mathfrak r) \subset \kappa(\mathfrak q_i)$는 순수 초월이다.
\item $\mathfrak q_i \subset B_i$ 위에 놓인 소 아이디얼
$\mathfrak q'_i \subset A_i$가 유일하게 존재한다.
% P05-EN-CH38-0036: A_i로의 형식에 맞지 않는 축약을 동결 원문대로 보존했다.
\item $\mathfrak q = \mathfrak q'_1 \cap S$이고
$\mathfrak q_i = \mathfrak q'_{i + 1} \cap A_i$이다.
\item $R \to B_i$의 상대차원은
$\dim_{\mathfrak q_i}(\text{Supp}(M_i \otimes_R \kappa(\mathfrak r)))$이다.
\end{enumerate}
\end{definition}

\begin{remark}
\label{flat-remark-same-notion}
$A \to B$를 유한형 환 준동형이라 하고 $N$을 유한 $B$-가군이라 하자.
$\mathfrak q$를 $B$의 소 아이디얼로서 $A$의 소 아이디얼
$\mathfrak r$ 위에 놓인 것이라 하자.
$X = \Spec(B)$, $S = \Spec(A)$로 두고 $X$ 위에서
$\mathcal{F} = \widetilde{N}$로 두자. $x$를 $\mathfrak q$에 대응하는
점이라 하고, $s \in S$를
% P05-EN-CH38-0037: 이 비고 전체의 정의되지 않은 p/r 불일치를 동결 원문대로 보존했다.
$\mathfrak p$에 대응하는 점이라 하자. 그러면
\begin{enumerate}
\item $\mathcal{F}/X/S$가 $s$ 위에서 완전 데비사주를 가지면
$N/B/A$가 $\mathfrak p$ 위에서 완전 데비사주를 갖는다.
\item $\mathcal{F}/X/S$가 $x$에서 완전 데비사주를 갖는 것과
$N/B/A$가 $\mathfrak q$에서 완전 데비사주를 갖는 것은 동치이다.
\end{enumerate}
다만 ``$\mathfrak p$ 위에서 $N/B/A$의 완전 데비사주''를 정식화할 때
상대차원에 관한 조건을 생략했다는 작은 차이가 있으며, 이 때문에 (1)의
함의는 한 방향으로만 성립한다. $\mathfrak q$에서의 완전 데비사주라는
개념에는 이 조건이 포함되어 있다. 어쨌든 여기서는
$\mathcal{F}/X/S$에 대한 존재가 $N/B/A$에 대한 존재를 함의한다는
사실만 사용할 것이다.
\end{remark}

\begin{lemma}
\label{flat-lemma-existence-algebra}
$R \to S$를 유한형 환 준동형이라 하자.
$M$을 유한 $S$-가군이라 하자.
$\mathfrak q$를 $S$의 소 아이디얼이라 하자.
환 준동형 $R \to S$의 $\mathfrak q$에서의 기본 에탈 국소화
$R' \to S', \mathfrak q', \mathfrak p'$로서
$(M \otimes_S S')/S'/R'$가 $\mathfrak q'$에서
완전 데비사주를 갖는 것이 존재한다.
\end{lemma}

\begin{proof}
이는 Proposition \ref{flat-proposition-existence-complete-at-x}를
Remark \ref{flat-remark-same-notion}를 통해 다시 쓴 것이다.
\end{proof}




\section{국소화와 보편 단사 사상}
\label{flat-section-localize-universally-injective}


\begin{lemma}
\label{flat-lemma-homothety-spectrum}
$R \to S$를 환 준동형이라 하자.
$N$을 $S$-가군이라 하자. 다음을 가정한다.
\begin{enumerate}
\item $R$은 극대 아이디얼 $\mathfrak m$을 갖는 국소환이다.
\item $\overline{S} = S/\mathfrak m S$는 뇌터 환이다.
% P05-EN-CH38-0038: 정의되지 않은 m_R/m 불일치를 동결 원문대로 보존했다.
\item $\overline{N} = N/\mathfrak m_R N$은 유한 $\overline{S}$-가군이다.
\end{enumerate}
$\Sigma \subset S$를 $\overline{N}$ 위에서 영인자가 아닌 원소들로
이루어진 곱셈적 부분집합이라 하자. 그러면 $\Sigma^{-1}S$는
반국소환이고, 그 스펙트럼은 $\text{Ass}_S(\overline{N})$의 어떤
원소에 포함되는 소 아이디얼 $\mathfrak q \subset S$들로 이루어진다.
더욱이 $\Sigma^{-1}S$의 모든 극대 아이디얼은
$\overline{S}$ 위의 $\overline{N}$의 연관 소아이디얼에 대응한다.
\end{lemma}

\begin{proof}
Algebra, Lemma \ref{algebra-lemma-ass-quotient-ring}에 의해
$\text{Ass}_S(\overline{N}) = \text{Ass}_{\overline{S}}(\overline{N})$
임에 유의하자. Algebra, Lemma \ref{algebra-lemma-finite-ass}에 의해
이는 유한집합이다.
$\{\mathfrak q_1, \ldots, \mathfrak q_r\} = \text{Ass}_S(\overline{N})$라 하자.
Algebra, Lemma \ref{algebra-lemma-ass-zero-divisors}에 의해
$\Sigma = S \setminus (\bigcup \mathfrak q_i)$이다.
Algebra, Lemma \ref{algebra-lemma-spec-localization}에서
$\Spec(\Sigma^{-1}S)$를 기술한 방식과 Algebra, Lemma
\ref{algebra-lemma-silly}에 의해 $\Sigma^{-1}S$의 소 아이디얼들은
어느 $\mathfrak q_i$에 포함되는 $S$의 소 아이디얼들과 대응한다.
따라서 $\Sigma^{-1}S$의 극대 아이디얼들은 집합
$\{\mathfrak q_1, \ldots, \mathfrak q_r\}$에서 포함관계에 관해 극대인
원소들과 일대일로 대응한다. 이로써 보조정리가 증명되었다.
\end{proof}

\begin{lemma}
\label{flat-lemma-homothety-universally-injective}
가정과 기호는 Lemma \ref{flat-lemma-homothety-spectrum}에서와 같다고 하자.
더 나아가 다음을 가정한다.
\begin{enumerate}
\item $S$는 국소환이고 $R \to S$는 국소 준동형이다.
\item $S$는 $R$ 위에서 본질적 유한표현이다.
\item $N$은 $S$ 위에서 유한표현이다.
\item $N$은 $R$ 위에서 평탄하다.
\end{enumerate}
그러면 각 $s \in \Sigma$는 보편 단사인 $R$-가군 사상
$s : N \to N$을 정하며, 사상 $N \to \Sigma^{-1}N$은
$R$-보편 단사이다.
\end{lemma}

\begin{proof}
Algebra, Lemma \ref{algebra-lemma-mod-injective-general}에 의해 완전열
$0 \to N \to N \to N/sN \to 0$에서 $N/sN$은 $R$ 위에서 평탄하다.
이는 $s : N \to N$이 보편 단사임을 뜻한다. Algebra, Lemma
\ref{algebra-lemma-flat-tor-zero}를 보라.
사상 $N \to \Sigma^{-1}N$은 사상들 $s : N \to N$의
유향 여극한이므로 보편 단사이다.
\end{proof}

\begin{lemma}
\label{flat-lemma-base-change-universally-flat-local}
$R \to S$를 환 준동형이라 하자.
$N$을 $S$-가군이라 하자.
$S \to S'$를 환 준동형이라 하자.
다음을 가정한다.
\begin{enumerate}
% P05-EN-CH38-0039: 목록 문장부호의 유일하게 정해지는 보정은 부록 원장에 기록했다.
\item $R \to S$는 국소환들 사이의 국소 준동형이다.
\item $S$는 $R$ 위에서 본질적 유한표현이다.
\item $N$은 $S$ 위에서 유한표현이다.
\item $N$은 $R$ 위에서 평탄하다.
\item $S \to S'$는 평탄하다.
\item $\Spec(S') \to \Spec(S)$의 상은 $\mathfrak m_R$ 위에 놓인
$S$의 소 아이디얼 $\mathfrak q$ 중에서 $\mathfrak q$가
$N/\mathfrak m_R N$의 연관 소아이디얼인 것을 모두 포함한다.
\end{enumerate}
그러면 $N \to N \otimes_S S'$는 $R$-보편 단사이다.
\end{lemma}

\begin{proof}
% P05-EN-CH38-0040: R/S 텐서 불일치를 동결 원문대로 보존했다.
$N' = N \otimes_R S'$로 두자. 다음 가환 그림을 생각하자.
$$
\xymatrix{
N \ar[d] \ar[r] & N' \ar[d] \\
\Sigma^{-1}N \ar[r] & \Sigma^{-1}N'
}
$$
여기서 $\Sigma \subset S$는 $N/\mathfrak m_R N$ 위에서
영인자가 아닌 원소들의 집합이다. 사상
$N \to \Sigma^{-1}N'$이 보편 단사임을 보일 수 있다면
$N \to N'$도 그러하다(Algebra, Lemma
\ref{algebra-lemma-universally-injective-permanence}를 보라).

\medskip\noindent
Lemma \ref{flat-lemma-homothety-spectrum}에 의해 환 $\Sigma^{-1}S$는
그 극대 아이디얼들이 $N/\mathfrak m_R N$의 연관 소아이디얼에
대응하는 반국소환이다. 따라서 가정에 의해
$\Spec(\Sigma^{-1}S') \to \Spec(\Sigma^{-1}S)$의 상은 이 극대
아이디얼들을 모두 포함한다. Algebra, Lemma
\ref{algebra-lemma-ff-rings}에 의해 환 준동형
$\Sigma^{-1}S \to \Sigma^{-1}S'$는 충실히 평탄하다. 따라서 사상
$\Sigma^{-1}N \to \Sigma^{-1}N'$, 즉 사상
$$
N \otimes_S \Sigma^{-1}S \longrightarrow N \otimes_S \Sigma^{-1}S'
$$
는 보편 단사이다. Algebra, Lemmas
\ref{algebra-lemma-faithfully-flat-universally-injective}와
\ref{algebra-lemma-universally-injective-tensor}를 보라.
끝으로 Lemma \ref{flat-lemma-homothety-universally-injective}를 적용하면
$N \to \Sigma^{-1}N$이 보편 단사임을 알 수 있다.
보편 단사인 가군 사상들의 합성은 보편 단사이므로(Algebra, Lemma
\ref{algebra-lemma-composition-universally-injective}를 보라)
$N \to \Sigma^{-1}N'$이 보편 단사이고, 이로써 증명이 끝난다.
\end{proof}

\begin{lemma}
\label{flat-lemma-base-change-universally-flat}
$R \to S$를 환 준동형이라 하자.
$N$을 $S$-가군이라 하자.
$S \to S'$를 환 준동형이라 하자.
다음을 가정한다.
\begin{enumerate}
\item $R \to S$는 유한표현이고 $N$은 $S$ 위에서 유한표현이다.
\item $N$은 $R$ 위에서 평탄하다.
\item $S \to S'$는 평탄하다.
\item $\Spec(S') \to \Spec(S)$의 상은 다음 조건을 만족하는 모든
소 아이디얼 $\mathfrak q$를 포함한다. 즉 $\mathfrak q$는
$N \otimes_R \kappa(\mathfrak p)$의 연관 소아이디얼이며, 여기서
$\mathfrak p$는 $\mathfrak q$의 $R$에서의 역상이다.
\end{enumerate}
그러면 $N \to N \otimes_S S'$는 $R$-보편 단사이다.
\end{lemma}

\begin{proof}
Algebra, Lemma \ref{algebra-lemma-universally-injective-check-stalks}에 의해,
$R$의 $\mathfrak p$ 위에 놓인 $S$의 임의의 소 아이디얼
$\mathfrak q$에 대하여
% P05-EN-CH38-0041: R/S 텐서 불일치를 동결 원문대로 보존했다.
$N_{\mathfrak q} \to (N \otimes_R S')_{\mathfrak q}$가
% P05-EN-CH38-0042: 불필요한 부정관사 보정은 한국어 술어로 자연스럽게 나타냈다.
$R_{\mathfrak p}$-보편 단사임을 보이면 충분하다. 따라서 환 준동형들
$R_{\mathfrak p} \to S_{\mathfrak q} \to S'_{\mathfrak q}$와
가군 $N_{\mathfrak q}$에 Lemma
\ref{flat-lemma-base-change-universally-flat-local}를 적용하면 된다.
\end{proof}

\noindent
독자는 다음 보조정리를 Algebra, Lemmas
\ref{algebra-lemma-mod-injective}와
\ref{algebra-lemma-mod-injective-general}, 그리고 Section
\ref{flat-section-variants-mod-injective}의 결과들과 비교해 볼 수 있다.
각 경우 결론은 사상 $u : M \to N$이 평탄한 여핵을 갖는
보편 단사라는 것이다.

\begin{lemma}
\label{flat-lemma-universally-injective-local}
$(R, \mathfrak m)$을 국소환이라 하자.
$u : M \to N$을 $R$-가군 사상이라 하자.
$M$이 사영 $R$-가군이고 $N$이 평탄 $R$-가군이며
$\overline{u} : M/\mathfrak mM \to N/\mathfrak mN$이 단사이면
$u$는 보편 단사이다.
\end{lemma}

\begin{proof}
Algebra, Theorem \ref{algebra-theorem-projective-free-over-local-ring}에 의해
가군 $M$은 자유이다. $M$의 유한 생성 직합 인자마다 결과가 성립함을
보이면 보조정리가 따르므로, $M$이 유한 자유라고 가정해도 된다.
$N = \colim_i N_i$를 유한 자유 가군들의 유향 여극한으로 쓰자.
Algebra, Theorem \ref{algebra-theorem-lazard}를 보라.
$M$이 유한 자유이므로 $u : M \to N$은 어떤 $i$에 대해
$N_i$를 통하여 분해된다. 이에 대응하는 $R$-가군 사상을
$u_i : M \to N_i$라 쓰자. $\overline{u}$가 단사이므로
$\overline{u_i} : M/\mathfrak mM \to N_i/\mathfrak mN_i$는 단사이고,
모든 $i' \geq i$에 대해 사상
$N_i/\mathfrak mN_i \to N_{i'}/\mathfrak mN_{i'}$와 합성한 뒤에도
단사이다. $M$과 $N_{i'}$는 국소환 $R$ 위의 유한 자유 가군이므로
모든 $i' \geq i$에 대해 $M \to N_{i'}$는 분할 단사이다.
따라서 임의의 $R$-가군 $Q$에 대해
$M \otimes_R Q \to N_{i'} \otimes_R Q$는 모든 $i' \geq i$에 대해
단사이다. $- \otimes_R Q$는 여극한과 가환하므로
% P05-EN-CH38-0043: N_{i'}/N 여극한 표적 불일치를 동결 원문대로 보존했다.
$M \otimes_R Q \to N_{i'} \otimes_R Q$는 원하는 대로 단사이다.
\end{proof}

\begin{lemma}
\label{flat-lemma-invert-universally-injective}
가정과 기호는 Lemma \ref{flat-lemma-homothety-spectrum}에서와 같다고 하자.
더 나아가 $N$이 $R$-가군으로서 사영이라고 가정하자.
그러면 각 $s \in \Sigma$는 보편 단사인 $R$-가군 사상
$s : N \to N$을 정하며, 사상 $N \to \Sigma^{-1}N$은
$R$-보편 단사이다.
\end{lemma}

\begin{proof}
$s \in \Sigma$를 택하자. Lemma
\ref{flat-lemma-universally-injective-local}에 의해
사상 $s : N \to N$은 보편 단사이다.
사상 $N \to \Sigma^{-1}N$은 사상들 $s : N \to N$의
유향 여극한이므로 보편 단사이다.
\end{proof}






\section{완비화와 Mittag--Leffler 가군}
\label{flat-section-completion-ML}


\begin{lemma}
\label{flat-lemma-completed-direct-sum-ML}
$R$을 환이라 하자. $I \subset R$을 아이디얼이라 하자.
$A$를 집합이라 하자. $R$이 뇌터 환이고 $I$에 관해 완비라고 가정하자.
완비화
$(\bigoplus\nolimits_{\alpha \in A} R)^\wedge$
는 평탄하고 Mittag--Leffler이다.
\end{lemma}

\begin{proof}
More on Algebra, Lemma
\ref{more-algebra-lemma-ui-completion-direct-sum-into-product}에 의해 사상
$(\bigoplus\nolimits_{\alpha \in A} R)^\wedge
\to \prod_{\alpha \in A} R$은 보편 단사이다.
따라서 Algebra, Lemmas \ref{algebra-lemma-ui-flat-domain}와
\ref{algebra-lemma-pure-submodule-ML}에 의해
$\prod_{\alpha \in A} R$이 평탄하고 Mittag--Leffler임을 보이면 충분하다.
Algebra, Proposition \ref{algebra-proposition-characterize-coherent}
(그리고 Algebra, Lemma \ref{algebra-lemma-Noetherian-coherent})에 의해
$\prod_{\alpha \in A} R$이 평탄함을 알 수 있다.
한편 $R$의 복사본들의 곱은 Mittag--Leffler이므로 결론이 따른다.
Algebra, Lemma \ref{algebra-lemma-product-over-Noetherian-ring}를 보라.
\end{proof}

\begin{lemma}
\label{flat-lemma-lift-ML}
$R$을 환이라 하자. $I \subset R$을 아이디얼이라 하자.
$M$을 $R$-가군이라 하자.
다음을 가정한다.
\begin{enumerate}
\item $R$은 뇌터 환이고 $I$-진 완비이다.
\item $M$은 $R$ 위에서 평탄하다.
\item $M/IM$은 사영 $R/I$-가군이다.
\end{enumerate}
그러면 $I$-진 완비화 $M^\wedge$는 평탄한 Mittag--Leffler
$R$-가군이다.
\end{lemma}

\begin{proof}
$F$가 자유 $R$-가군인 전사 $F \to M$을 택하자.
Algebra, Lemma \ref{algebra-lemma-split-completed-sequence}에 의해
가군 $M^\wedge$는 가군 $F^\wedge$의 직합 인자이다.
따라서 $F$에 대해 보조정리를 증명하면 충분하다.
이 경우 결론은 Lemma \ref{flat-lemma-completed-direct-sum-ML}에서 따른다.
\end{proof}

\noindent
Lemmas \ref{flat-lemma-universally-injective-to-completion}와
\ref{flat-lemma-universally-injective-to-completion-flat}에서
$S$가 뇌터 환이라는 가정은 $R \to S$가 유한형이면 성립한다.
Algebra, Lemma \ref{algebra-lemma-Noetherian-permanence}를 보라.

\begin{lemma}
\label{flat-lemma-universally-injective-to-completion}
$R$을 환이라 하자.
$I \subset R$을 아이디얼이라 하자.
$R \to S$를 환 준동형이라 하고 $N$을 $S$-가군이라 하자.
다음을 가정한다.
\begin{enumerate}
\item $R$은 뇌터 환이다.
\item $S$는 뇌터 환이다.
\item $N$은 유한 $S$-가군이다.
\item 임의의 유한 $R$-가군 $Q$와 임의의
$\mathfrak q \in \text{Ass}_S(Q \otimes_R N)$에 대하여
$IS + \mathfrak q \not = S$이다.
\end{enumerate}
그러면 $N$에서 $N$의 $I$-진 완비화로 가는 사상
$N \to N^\wedge$는 $R$-가군 사상으로서 보편 단사이다.
\end{lemma}

\begin{proof}
Algebra, Theorem \ref{algebra-theorem-universally-exact-criteria}에 의해
임의의 유한 $R$-가군 $Q$에 대해 사상
$Q \otimes_R N \to Q \otimes_R N^\wedge$가 단사임을 보여야 한다.
표준 사상 $Q \otimes_R N^\wedge \to (Q \otimes_R N)^\wedge$가
있으므로 표준 사상
$Q \otimes_R N \to (Q \otimes_R N)^\wedge$가 단사임을 보이면 충분하다.
따라서 $N$을 $Q \otimes_R N$으로 바꾸어도 되며, 사상
$N \to N^\wedge$의 단사성을 증명하면 충분하다.

\medskip\noindent
$K = \Ker(N \to N^\wedge)$로 두자.
$N$은 $\prod_{\mathfrak q \in \text{Ass}(N)} N_{\mathfrak q}$의
부분가군이므로 모든 $\mathfrak q \in \text{Ass}(N)$에 대해
$K_{\mathfrak q} = 0$임을 보이면 충분하다. Algebra, Lemma
\ref{algebra-lemma-zero-at-ass-zero}를 보라.
$\mathfrak q \in \text{Ass}(N)$를 택하자. 마지막 가정에 의해
$\mathfrak q' \supset IS + \mathfrak q$인 소 아이디얼이 존재한다.
$K_{\mathfrak q}$는 $K_{\mathfrak q'}$의 국소화이므로
$K_{\mathfrak q'}$가 영임을 보이면 충분하다.
$K = \bigcap I^nN$이므로
$K_{\mathfrak q'} \subset \bigcap I^nN_{\mathfrak q'}$이다.
따라서 Algebra, Lemma
\ref{algebra-lemma-intersect-powers-ideal-module-zero}에 의해
$K_{\mathfrak q'} = 0$이다.
\end{proof}

\begin{lemma}
\label{flat-lemma-universally-injective-to-completion-flat}
$R$을 환이라 하자.
$I \subset R$을 아이디얼이라 하자.
$R \to S$를 환 준동형이라 하고 $N$을 $S$-가군이라 하자.
다음을 가정한다.
\begin{enumerate}
\item $R$은 뇌터 환이다.
\item $S$는 뇌터 환이다.
\item $N$은 유한 $S$-가군이다.
\item $N$은 $R$ 위에서 평탄하다.
\item $N \otimes_R \kappa(\mathfrak p)$의 연관 소아이디얼인 임의의
소 아이디얼 $\mathfrak q \subset S$에 대하여
% P05-EN-CH38-0044: 임의의 환 사상에서 형식에 맞지 않는 R∩q 축약을 동결 원문대로 보존했다.
$\mathfrak p = R \cap \mathfrak q$이면
$IS + \mathfrak q \not = S$이다.
\end{enumerate}
그러면 $N$에서 $N$의 $I$-진 완비화로 가는 사상
$N \to N^\wedge$는 $R$-가군 사상으로서 보편 단사이다.
\end{lemma}

\begin{proof}
이는 Lemma \ref{flat-lemma-universally-injective-to-completion}에서 따른다.
실제로 Algebra, Lemma
\ref{algebra-lemma-bourbaki-fibres}와
Remark \ref{algebra-remark-bourbaki}에 의해
% P05-EN-CH38-0045: 단수 주어의 수 일치 오류는 한국어로 자연스럽게 나타냈다.
텐서곱 $N \otimes_R Q$의 연관 소아이디얼들의 집합은
가군들 $N \otimes_R \kappa(\mathfrak p)$의 연관 소아이디얼들의
집합에 포함된다.
\end{proof}




\section{사영 가군}
\label{flat-section-projective}

\noindent
다음 보조정리를 사용하면 밑에 관한 뇌터 귀납법으로 사영성을 증명할 수 있다.
Lemma \ref{flat-lemma-fibres-irreducible-flat-projective}를 보라.

\begin{lemma}
\label{flat-lemma-flat-pure-over-complete-projective}
$R$을 환이라 하자.
$I \subset R$을 아이디얼이라 하자.
$R \to S$를 환 준동형이라 하고 $N$을 $S$-가군이라 하자.
다음을 가정한다.
\begin{enumerate}
\item $R$은 뇌터 환이고 $I$-진 완비이다.
\item $R \to S$는 유한형이다.
\item $N$은 유한 $S$-가군이다.
\item $N$은 $R$ 위에서 평탄하다.
\item $N/IN$은 사영 $R/I$-가군이다.
\item $N \otimes_R \kappa(\mathfrak p)$의 연관 소아이디얼인 임의의
소 아이디얼 $\mathfrak q \subset S$에 대하여
% P05-EN-CH38-0046: 임의의 환 사상에서 형식에 맞지 않는 R∩q 축약을 동결 원문대로 보존했다.
$\mathfrak p = R \cap \mathfrak q$이면
$IS + \mathfrak q \not = S$이다.
\end{enumerate}
그러면 $N$은 사영 $R$-가군이다.
\end{lemma}

\begin{proof}
Lemma \ref{flat-lemma-universally-injective-to-completion-flat}에 의해
사상 $N \to N^\wedge$는 보편 단사이다.
Lemma \ref{flat-lemma-lift-ML}에 의해 가군 $N^\wedge$는
Mittag--Leffler이다. Algebra, Lemma
\ref{algebra-lemma-pure-submodule-ML}에 의해 $N$도
Mittag--Leffler임을 얻는다. 따라서 $N$은 $R$-가군으로서
가산 생성이고 평탄하며 Mittag--Leffler이므로, Algebra, Lemma
\ref{algebra-lemma-countgen-projective}에 의해 사영이다.
\end{proof}

\begin{lemma}
\label{flat-lemma-fibres-irreducible-flat-projective}
$R$을 환이라 하자.
$R \to S$를 환 준동형이라 하자.
다음을 가정한다.
\begin{enumerate}
\item $R$은 뇌터 환이다.
\item $R \to S$는 유한형이고 평탄하다.
\item 모든 올 환 $S \otimes_R \kappa(\mathfrak p)$는
$\kappa(\mathfrak p)$ 위에서 기하적으로 정역이다.
\end{enumerate}
그러면 $S$는 사영 $R$-가군이다.
\end{lemma}

\begin{proof}
집합
$$
\{I \subset R \mid S/IS\text{가 }R/I\text{-가군으로서 사영이 아니다}\}
$$
을 생각하자. 이 집합이 공집합임을 보여야 한다. 모순을 얻기 위해
공집합이 아니라고 가정하자. 그러면 극대원 $I$를 갖는다.
그 근기를 $J = \sqrt{I}$라 하자. $I \not = J$이면
$S/JS$는 사영 $R/J$-가군이고, $S/IS$는 $R/I$ 위에서 평탄하며
$J/I$는 $R/I$의 멱영 아이디얼이다. Algebra, Lemma
\ref{algebra-lemma-lift-projective}를 적용하면 $S/IS$가 사영
$R/I$-가군임을 알 수 있는데, 이는 모순이다.
따라서 $I$가 근기 아이디얼이라고 가정해도 된다. 다시 말해
$R$이 축약환이고 $R$의 모든 영이 아닌 아이디얼 $I$에 대해
$S/IS$가 사영 $R/I$-가군인 경우에 보조정리를 증명하는 문제로 환원된다.

\medskip\noindent
$R$이 축약환이고 $R$의 모든 영이 아닌 아이디얼 $I$에 대해
$S/IS$가 사영 $R/I$-가군이라고 가정하자.
일반 평탄성, 즉 Algebra, Lemma
\ref{algebra-lemma-generic-flatness-Noetherian}
(정역인 국소화 $R_g$에 적용한다) 또는 더 일반적인 Algebra, Lemma
\ref{algebra-lemma-generic-flatness-reduced}에 의해
$S_f$가 자유 $R_f$-가군이 되는 영이 아닌 $f \in R$이 존재한다.
$R$의 $(f)$-진 완비화를
$R^\wedge = \lim R/(f^n)$라 쓰자. 환 준동형
$$
R \longrightarrow R_f \times R^\wedge
$$
은 충실히 평탄하다. Algebra, Lemma
\ref{algebra-lemma-completion-flat}를 보라.
따라서 사영성의 충실히 평탄한 강하(Algebra, Theorem
\ref{algebra-theorem-ffdescent-projectivity}를 보라)에 의해
$S \otimes_R R^\wedge$가 사영 $R^\wedge$-가군임을 보이면 충분하다.
이를 위해 Lemma \ref{flat-lemma-flat-pure-over-complete-projective}의
판정법을 사용한다. 먼저
$S/fS = (S \otimes_R R^\wedge)/f(S \otimes_R R^\wedge)$가
사영 $R/(f)$-가군이고, $S \otimes_R R^\wedge$가 그러한 사상의
기저변환으로서 $R^\wedge$ 위에서 평탄하고 유한형임에 유의하자.
다음으로 $\mathfrak p^\wedge$를 $R^\wedge$의 소 아이디얼이라 하자.
$\mathfrak p \subset R$을 $R$의 대응하는 소 아이디얼이라 하자.
$R \to S$가 기하적으로 정역인 올 환들을 가지므로 임의의 기저변환의
올 환들에 대해서도 마찬가지이다. 따라서
$\mathfrak q^\wedge = \mathfrak p^\wedge(S \otimes_R R^\wedge)$는
% P05-EN-CH38-0047: 단복수와 문장부호 오류는 한국어로 자연스럽게 나타냈다.
$\mathfrak p^\wedge$ 위에 놓인 소 아이디얼이고
$S \otimes_R \kappa(\mathfrak p^\wedge)$의 유일한 연관 소아이디얼이다.
따라서
$f(S \otimes_R R^\wedge) + \mathfrak q^\wedge \not = S \otimes_R R^\wedge$
이면 끝난다. 이는 참이다. 실제로 $f$가 $f$-진 완비환
$R^\wedge$의 Jacobson 근기에 속하므로
$\mathfrak p^\wedge + fR^\wedge \not = R^\wedge$이고,
$R^\wedge \to S \otimes_R R^\wedge$의 올들이 공집합이 아니므로
(기약 공간은 공집합이 아니다) 이 사상은 스펙트럼 위에서 전사이다.
\end{proof}

\begin{lemma}
\label{flat-lemma-fibres-irreducible-flat-projective-nonnoetherian}
$R$을 환이라 하자. $R \to S$를 환 준동형이라 하자.
다음을 가정한다.
\begin{enumerate}
\item $R \to S$는 유한표현이고 평탄하다.
\item 모든 올 환 $S \otimes_R \kappa(\mathfrak p)$는
$\kappa(\mathfrak p)$ 위에서 기하적으로 정역이다.
\end{enumerate}
그러면 $S$는 사영 $R$-가군이다.
\end{lemma}

\begin{proof}
다음과 같은 환들의 여카르테시안 그림을 찾을 수 있다.
$$
\xymatrix{
S_0 \ar[r] & S \\
R_0 \ar[u] \ar[r] & R \ar[u]
}
$$
여기서 $R_0$는 $\mathbf{Z}$ 위에서 유한형이고, 사상
$R_0 \to S_0$는 기하적으로 정역인 올들을 갖는 유한형 평탄
사상이다. More on Morphisms, Lemmas
\ref{more-morphisms-lemma-Noetherian-approximation-flat},
\ref{more-morphisms-lemma-Noetherian-approximation-geometrically-reduced},
\ref{more-morphisms-lemma-Noetherian-approximation-geometrically-irreducible},
그리고 \ref{more-morphisms-lemma-Noetherian-approximation-combine}를 보라.
Lemma \ref{flat-lemma-fibres-irreducible-flat-projective}에 의해
$S_0$는 사영 $R_0$-가군이다. 따라서
$S = S_0 \otimes_{R_0} R$은 사영 $R$-가군이다.
Algebra, Lemma \ref{algebra-lemma-ascend-properties-modules}를 보라.
\end{proof}

\begin{remark}
\label{flat-remark-how-in-RG}
Lemma \ref{flat-lemma-fibres-irreducible-flat-projective-nonnoetherian}는
이 장의 결과들을 전개하는 데 중요한 단계이다.
\cite{GruRay}에서 이 보조정리에 대응하는 것은
\cite[I Proposition 3.3.1]{GruRay}이다. 즉
$R \to S$가 기하적으로 정역인 올들을 갖는 매끄러운 사상이면
$S$는 사영 $R$-가군이다. 이는 Lemma
\ref{flat-lemma-fibres-irreducible-flat-projective-nonnoetherian}의 특수한
경우이지만, 나중에 어차피 이 보조정리를 개선할 것이므로 이 시점에서
더 강한 결과를 얻어도 큰 이득은 없다.
\cite{GruRay}에 주어진 이 증명을 간단히 개괄한다.
\begin{enumerate}
\item 먼저 위와 같이 $R$이 뇌터 환인 경우로 환원한다.
\item 사영성은 충실히 평탄한 환 준동형을 따라 강하하므로(Algebra, Theorem
\ref{algebra-theorem-ffdescent-projectivity}를 보라)
$R$ 위의 fppf 위상에서 국소적으로 다루어도 된다. 따라서
$R \to S$가 단면 $\sigma : S \to R$을 갖는다고 가정해도 된다.
(통상적인 방법대로 $S$로 기저변환하면 된다.)
$I = \Ker(S \to R)$로 두자.
\item $R$ 위에서 조금 더 국소화하면 $I/I^2$이 자유 $R$-가군이고,
$I$에 관한 $S$의 완비화 $S^\wedge$가
$R[[t_1, \ldots, t_n]]$와 동형이라고 가정해도 된다.
Morphisms, Lemma \ref{morphisms-lemma-section-smooth-morphism}를 보라.
여기서 $R \to S$가 매끄럽다는 사실을 사용한다.
\item $S$가 사영 $R$-가군임을 증명하려면 $S$가 $R$-가군으로서
평탄하고 가산 생성이며 Mittag--Leffler임을 증명하면 충분하다.
Algebra, Lemma \ref{algebra-lemma-countgen-projective}를 보라.
처음 두 성질은 명백하다. 따라서 $S$가 $R$-가군으로서
Mittag--Leffler임을 증명하면 충분하다. Algebra, Lemma
\ref{algebra-lemma-power-series-ML}에 의해 가군
$R[[t_1, \ldots, t_n]]$은 $R$ 위에서 Mittag--Leffler이다.
따라서 Algebra, Lemma \ref{algebra-lemma-pure-submodule-ML}에 의해
% P05-EN-CH38-0048: 사상 앞의 불필요한 정관사는 한국어로 자연스럽게 나타냈다.
$S \to S^\wedge$가 $R$-가군 사상으로서 보편 단사임을 보이면 충분하다.
\item Lemma \ref{flat-lemma-base-change-universally-flat}를 적용하면
$S \to S^\wedge$가 $R$-보편 단사임을 알 수 있다.
실제로 $R \to S$가 기하적으로 정역인 올들을 가지므로, 임의의 올 환의
모든 연관점은 그 올 환의 일반점일 뿐이며 이는
$\Spec(S^\wedge) \to \Spec(S)$의 상에 들어 있다.
\end{enumerate}
방금 개괄한 증명과 Lemma
\ref{flat-lemma-fibres-irreducible-flat-projective-nonnoetherian}의
증명으로 이어지는 논증 전개 사이에는 유사점이 있다.
% P05-EN-CH38-0049: 도입구 뒤 문장부호 오류는 한국어로 자연스럽게 나타냈다.
양쪽 모두에서 완비화가 본질적인 역할을 하고, 두 경우 모두
기하적으로 정역인 올들을 갖는다는 가정이 $S$에서 $S$의 완비화로
가는 사상이 $R$-보편 단사임을 보장한다.
\end{remark}












\section{유한형 평탄 가군 I}
\label{flat-section-finite-type-flat-I}

\noindent
어떤 경우에는 유한표현 환 준동형 $R \to S$와 유한 $S$-가군
$N$이 주어졌을 때 $N$이 $R$ 위에서 평탄이면 $N$이
유한표현이다. 이 절에서는 이것이 ``점별로'' 참임을 증명한다.
Proposition \ref{flat-proposition-finite-type-flat-at-point}의 첫 번째 증명은
Section \ref{flat-section-local-structure-module}의 기하학적 결과들을
사용하지만 완전 데비사주의 존재는 사용하지 않음을 밝혀 둔다.

\begin{lemma}
\label{flat-lemma-induction-step}
$(R, \mathfrak m)$을 국소환이라 하자.
$R \to S$를 기하적으로 정역인 올들을 갖는 유한표현 평탄 환 준동형이라 하자.
$\mathfrak p = \mathfrak mS$로 쓰자.
$\mathfrak q \subset S$를 $\mathfrak m$ 위에 놓인 소 아이디얼이라 하자.
$N$을 유한 $S$-가군이라 하자.
$r \geq 0$과 $S$-가군 사상
$$
\alpha : S^{\oplus r} \longrightarrow N
$$
이 존재하여
$\alpha : \kappa(\mathfrak p)^{\oplus r} \to N \otimes_S \kappa(\mathfrak p)$
는 동형이다. 이러한 임의의 $\alpha$에 대해 다음은 서로 동치이다.
\begin{enumerate}
\item $N_{\mathfrak q}$는 $R$-평탄이다.
\item $\alpha$는 $R$-보편 단사이고
$\Coker(\alpha)_{\mathfrak q}$는 $R$-평탄이다.
\item $\alpha$는 단사이고
$\Coker(\alpha)_{\mathfrak q}$는 $R$-평탄이다.
\item $\alpha_{\mathfrak p}$는 동형이고
$\Coker(\alpha)_{\mathfrak q}$는 $R$-평탄이다.
\item $\alpha_{\mathfrak q}$는 단사이고
$\Coker(\alpha)_{\mathfrak q}$는 $R$-평탄이다.
\end{enumerate}
\end{lemma}

\begin{proof}
$\alpha$를 얻기 위해
$r = \dim_{\kappa(\mathfrak p)} N \otimes_S \kappa(\mathfrak p)$로 두고,
$N \otimes_S \kappa(\mathfrak p)$의 기저를 이루는
$x_1, \ldots, x_r \in N$을 택한다.
$\alpha(s_1, \ldots, s_r) = \sum s_i x_i$로 정의하자.
이로써 존재성이 증명된다.

\medskip\noindent
$\alpha$ 하나를 고정하자. 가장 흥미로운 함의는
(1) $\Rightarrow$ (2)이며 이를 먼저 증명한다. (1)을 가정하자.
$S/\mathfrak mS$는 분수체가 $\kappa(\mathfrak p)$인 정역이므로
$(S/\mathfrak mS)^{\oplus r} \to
N_{\mathfrak p}/\mathfrak mN_{\mathfrak p} = N \otimes_S \kappa(\mathfrak p)$
가 단사임을 알 수 있다. 따라서
% P05-EN-CH38-0129: 인용 뒤의 잘못된 마침표는 한국어 문장으로 자연스럽게 나타냈다.
Lemmas \ref{flat-lemma-universally-injective-local}와
\ref{flat-lemma-fibres-irreducible-flat-projective-nonnoetherian}에 의해
사상 $S^{\oplus r} \to N_{\mathfrak p}$는 $R$-보편 단사이다.
따라서 $S^{\oplus r} \to N$도 $R$-보편 단사이다. Algebra, Lemma
\ref{algebra-lemma-universally-injective-permanence}를 보라.
그러면 국소화 $\alpha_{\mathfrak q}$도 $R$-보편 단사이다.
Algebra, Lemma \ref{algebra-lemma-universally-injective-localize}를 보라.
Algebra, Lemma \ref{algebra-lemma-ui-flat-domain}에 의해
$\Coker(\alpha)_{\mathfrak q}$가 $R$-평탄이라는 결론을 얻는다.

\medskip\noindent
함의 (2) $\Rightarrow$ (3)은 자명하다. (3)이 성립하면
$\alpha_{\mathfrak p}$는 단사 가군 사상의 국소화이므로 단사이다.
Nakayama 보조정리(Algebra, Lemma \ref{algebra-lemma-NAK})에 의해
$\alpha_{\mathfrak p}$는 전사이기도 하다. 따라서
(3) $\Rightarrow$ (4)이다. (4)가 성립하면
$\alpha_{\mathfrak p}$는 동형이므로
$S_{\mathfrak q} \to S_{\mathfrak p}$가 단사라는 사실에 의해
$\alpha$는 단사이다. 실제로 Lemmas
\ref{flat-lemma-invert-universally-injective}와
\ref{flat-lemma-fibres-irreducible-flat-projective-nonnoetherian}를
함께 적용하면 $S \setminus \mathfrak p$의 원소들은 $S$ 위에서
영인자가 아니다. 따라서 (4) $\Rightarrow$ (5)이다.
끝으로 (5)가 성립하면 $N_{\mathfrak q}$는 평탄 가군들의 확대이므로
$R$-평탄이다. Algebra, Lemma \ref{algebra-lemma-flat-ses}를 보라.
따라서 (5) $\Rightarrow$ (1)이고 증명이 끝난다.
\end{proof}

\begin{lemma}
\label{flat-lemma-complete-devissage-flat-finite-type-module}
$(R, \mathfrak m)$을 국소환이라 하자.
$R \to S$를 유한표현 환 준동형이라 하자.
$N$을 유한 $S$-가군이라 하자.
$\mathfrak q$를 $S$의 소 아이디얼로서 $\mathfrak m$ 위에 놓인 것이라 하자.
$N_{\mathfrak q}$가 $R$ 위에서 평탄이고,
$\mathfrak q$에서 $N/S/R$의 완전 데비사주가 존재한다고 가정하자.
그러면 $N$은 유한표현 $S$-가군이고 자유 $R$-가군이며, 다음 동형
$$
N \cong B_1^{\oplus r_1} \oplus \ldots \oplus B_n^{\oplus r_n}
$$
이 존재한다. 여기서 각 $B_i$는 기하적으로 정역인 올들을 갖는
매끄러운 $R$-대수이고, 이 동형은 $R$-가군의 동형이다.
\end{lemma}

\begin{proof}
$(A_i, B_i, M_i, \alpha_i, \mathfrak q_i)_{i = 1, \ldots, n}$를
주어진 완전 데비사주라 하자. $n$에 대한 귀납법으로 보조정리를 증명한다.
$N$이 유한표현 $S$-가군일 필요충분조건은
% P05-EN-CH38-0130: 잘못된 부정관사는 한국어로 자연스럽게 나타냈다.
$M_1$이 유한표현 $B_1$-가군인 것이다. Remark
\ref{flat-remark-finite-presentation}을 보라.
$N_{\mathfrak q} \cong (M_1)_{\mathfrak q_1}$는 $R$-가군의
동형임에 유의하자. 실제로 (a)
$N_{\mathfrak q} \cong (M_1)_{\mathfrak q'_1}$이고, 여기서
$\mathfrak q'_1$은 $\mathfrak q_1$ 위에 놓인 $A_1$의 유일한
소 아이디얼이며, (b) Algebra, Lemma
\ref{algebra-lemma-unique-prime-over-localize-below}에 의해
$(A_1)_{\mathfrak q'_1} = (A_1)_{\mathfrak q_1}$이고, 따라서 (c)
$(M_1)_{\mathfrak q'_1} \cong (M_1)_{\mathfrak q_1}$이다.
그러므로 $(M_1)_{\mathfrak q_1}$은 평탄 $R$-가군이다.
따라서 보조정리를 증명하기 위해 $(S, N)$을 $(B_1, M_1)$으로
바꾸어도 된다. Lemma \ref{flat-lemma-induction-step}에 의해 사상
$\alpha_1 : B_1^{\oplus r_1} \to M_1$은 $R$-보편 단사이고
% P05-EN-CH38-0131: q/q_1 불일치를 동결 원문대로 보존했다.
$\Coker(\alpha_1)_{\mathfrak q}$는 $R$-평탄이다.
$(A_i, B_i, M_i, \alpha_i, \mathfrak q_i)_{i = 2, \ldots, n}$는
$\mathfrak q_1$에서 $\Coker(\alpha_1)/B_1/R$의 완전 데비사주임에
유의하자. 따라서 귀납가정에 의해 $\Coker(\alpha_1)$은
유한표현 $B_1$-가군이고 자유 $R$-가군이며 보조정리에서와 같은
분해를 갖는다. 이는 $M_1$이 유한표현 $B_1$-가군임을 뜻한다.
Algebra, Lemma \ref{algebra-lemma-extension}을 보라.
또한
$M_1 \cong B_1^{\oplus r_1} \oplus \Coker(\alpha_1)$는
$R$-가군의 동형이므로 보조정리에서와 같은 분해를 얻는다.
끝으로 Lemma
\ref{flat-lemma-fibres-irreducible-flat-projective-nonnoetherian}에 의해
$B_1$은 사영 $R$-가군이고, 따라서 Algebra, Theorem
\ref{algebra-theorem-projective-free-over-local-ring}에 의해 자유
$R$-가군이다. 이로써 증명이 끝난다.
\end{proof}

\begin{proposition}
\label{flat-proposition-finite-type-flat-at-point}
$f : X \to S$를 스킴의 사상이라 하자.
$\mathcal{F}$를 $X$ 위의 준연접층이라 하자.
$x \in X$의 상을 $s \in S$라 하자.
다음을 가정한다.
\begin{enumerate}
\item $f$는 국소 유한표현이다.
\item $\mathcal{F}$는 유한형이다.
\item $\mathcal{F}$는 $x$에서 $S$ 위로 평탄하다.
\end{enumerate}
그러면 기본 에탈 근방 $(S', s') \to (S, s)$와
다음 열린 부분스킴
$$
V \subset X \times_S \Spec(\mathcal{O}_{S', s'})
$$
이 존재한다. 이 부분스킴은 $x$로 가는
$X \times_S \Spec(\mathcal{O}_{S', s'})$의 유일한 점을 포함하고,
$\mathcal{F}$의 $V$로의 당김은 유한표현
$\mathcal{O}_V$-가군이며 $\mathcal{O}_{S', s'}$ 위에서 평탄하다.
\end{proposition}

\begin{proof}[첫 번째 증명]
이 증명은 더 길지만 완전 데비사주의 존재를 사용하지 않는다.
문제는 $x$와 $s$의 근방에서 국소적이므로 $X$와 $S$가
아핀이라고 가정해도 된다. 증명 도중에
% P05-EN-CH38-0132: 부사의 잘못된 위치는 한국어로 자연스럽게 나타냈다.
$S$를 $(S, s)$의 기본 에탈 근방으로 유한 번 바꿀 것이다.
그러면 목표는 (그러한 교체 뒤에)
$x$를 포함하는 열린집합
$V \subset X \times_S \Spec(\mathcal{O}_{S, s})$를 찾아
$\mathcal{F}|_V$가 $S$ 위에서 평탄이고 유한표현이 되게 하는 것이다.
물론 증명의 어느 지점에서든 $S$를
$\Spec(\mathcal{O}_{S, s})$로 바꿀 수도 있다. 즉
$S$가 국소 스킴이라고 가정해도 된다.
정수 $n = \dim_x(\text{Supp}(\mathcal{F}_s))$에 대한 귀납법으로
명제를 증명할 것이다.

\medskip\noindent
다음을 택할 수 있다.
\begin{enumerate}
\item 기본 에탈 근방들 $g : (X', x') \to (X, x)$,
$e : (S', s') \to (S, s)$,
\item 가환 그림
$$
\xymatrix{
X \ar[dd]_f & X' \ar[dd] \ar[l]^g & Z' \ar[l]^i \ar[d]^\pi \\
& & Y' \ar[d]^h \\
S & S' \ar[l]_e & S' \ar@{=}[l]
}
$$
\item $i(z') = x'$, $y' = \pi(z')$, $h(y') = s'$를 만족하는
점 $z' \in Z'$,
\item 유한형 준연접 $\mathcal{O}_{Z'}$-가군 $\mathcal{G}$.
\end{enumerate}
이들은 Lemma \ref{flat-lemma-elementary-devissage}에서와 같다.
증명의 첫 문단에서 설명한 대로 $S$를
$\Spec(\mathcal{O}_{S', s'})$로 바꿀 것이다.
다음 그림을 생각하자.
$$
\xymatrix{
X_{\mathcal{O}_{S', s'}} \ar[ddr]_f &
X'_{\mathcal{O}_{S', s'}} \ar[dd] \ar[l]^g &
Z'_{\mathcal{O}_{S', s'}} \ar[l]^i \ar[d]^\pi \\
& & Y'_{\mathcal{O}_{S', s'}} \ar[dl]^h \\
& \Spec(\mathcal{O}_{S', s'})
}
$$
여기서는 $S'$ 위의 스킴 $X', Z', Y'$를
$\Spec(\mathcal{O}_{S', s'}) \to S'$를 따라 기저변환했고,
$S$ 위의 스킴 $X$를
$\Spec(\mathcal{O}_{S', s'}) \to S$를 따라 기저변환했다.
여전히 $g$는 에탈이다. Lemma \ref{flat-lemma-etale-at-point}를 보라.
$X$를 $X_{\mathcal{O}_{S', s'}}$로,
$X'$를 $X'_{\mathcal{O}_{S', s'}}$로,
$Z'$를 $Z'_{\mathcal{O}_{S', s'}}$로,
$Y'$를 $Y'_{\mathcal{O}_{S', s'}}$로 바꾼 뒤에는
Lemma \ref{flat-lemma-elementary-devissage}에서와 같은 그림을 가지며,
추가로 $S = S'$는 닫힌점이 $s$인 국소 스킴이라고 가정해도 된다.
Lemmas \ref{flat-lemma-devissage-finite-presentation}과
\ref{flat-lemma-devissage-flat}에 의해 $Y' \to S$, 층
$\pi_*\mathcal{G}$ 및 점 $y'$에 대한 결과는
$X \to S$, $\mathcal{F}$ 및 $x$에 대한 결과를 함의한다.
따라서 $S$가 국소이고 $X \to S$가 상대차원 $n$의 기하적으로
기약인 올들을 갖는 아핀 스킴들의 매끄러운 사상이라고 가정해도 된다.

\medskip\noindent
귀납법의 기초 단계: $n = 0$.
$X \to S$는 차원이 $0$인 기하적으로 기약인 올들을 갖는
매끄러운 사상이므로 $X \to S$는 열린 몰입이다. Descent, Lemma
\ref{descent-lemma-universally-injective-etale-open-immersion}를 보라.
$S$가 국소이고 그 닫힌점이 $X \to S$의 상에 들어 있으므로
$X = S$라는 결론을 얻는다. 따라서 $\mathcal{F}$는 유한 평탄
$\mathcal{O}_{S, s}$-가군에 대응한다. 이 경우 Algebra, Lemma
\ref{algebra-lemma-finite-flat-local}에 의해 $\mathcal{F}$가 실제로
유한 자유임을 알 수 있으므로 결론이 따른다.

\medskip\noindent
귀납 단계. 닫힌 올에서 지지의 차원이 $< n$일 때마다 결과가
성립한다고 가정하자. 어떤 $B$-가군 $N$에 대해
$S = \Spec(A)$, $X = \Spec(B)$, $\mathcal{F} = \widetilde{N}$로 쓰자.
$A$는 국소환이며 그 극대 아이디얼을 $\mathfrak m$이라 쓰자.
$X \to S$가 기하적으로 기약인 올들을 가지므로
$\mathfrak p = \mathfrak mB$는 $\mathfrak m$ 위에 놓인 유일한
극소 소 아이디얼이다. 끝으로 $\mathfrak q \subset B$를 $x$에
대응하는 소 아이디얼이라 하자. Lemma \ref{flat-lemma-induction-step}에 의해
다음 사상
$$
\alpha : B^{\oplus r} \to N
$$
을 택하여
$\kappa(\mathfrak p)^{\oplus r} \to N \otimes_B \kappa(\mathfrak p)$가
동형이 되게 할 수 있다. 더욱이 $N_{\mathfrak q}$가 $A$-평탄이므로
같은 보조정리에 의해 $\alpha$가 단사이고
$\Coker(\alpha)_{\mathfrak q}$가 $A$-평탄이다.
$Q = \Coker(\alpha)$로 두자. $Q/\mathfrak mQ$의 지지는
$\mathfrak p$를 포함하지 않음에 유의하자. 따라서 확실히
$\dim_{\mathfrak q}(\text{Supp}(Q/\mathfrak mQ)) < n$이다.
$Q$에 관해 아는 모든 사실을 합치면 귀납가정을 $Q$에 적용할 수 있다.
따라서
% P05-EN-CH38-0133: S' 위의 잘못된 점 s를 동결 원문대로 보존했다.
기본 에탈 사상 $(S', s) \to (S, s)$가 존재하여,
$A' = \mathcal{O}_{S', s'}$일 때
$B \otimes_A A'$ 위의 $Q \otimes_A A'$에 대해 결론이 성립한다.
$A$를 $A'$로 바꾼 뒤에는 다음 유한 $B$-가군들의 완전열
$$
0 \to B^{\oplus r} \to N \to Q \to 0
$$
을 갖고(여기서는 위에서 말했듯 $\alpha$가 단사임을 사용한다),
또 $g \in B$, $g \not \in \mathfrak q$인 원소를 얻어
$Q_g$가 $B_g$ 위에서 유한표현이고 $A$ 위에서 평탄이 되게 한다.
국소화는 완전하므로
$$
0 \to B_g^{\oplus r} \to N_g \to Q_g \to 0
$$
은 여전히 완전하다. $B_g$와 $Q_g$가 $A$ 위에서 평탄이므로
$N_g$도 $A$ 위에서 평탄이다. Algebra, Lemma
\ref{algebra-lemma-flat-ses}를 보라. 또 $B_g$와 $Q_g$가
$B_g$ 위에서 유한표현이므로 $N_g$도 그러하다. Algebra, Lemma
\ref{algebra-lemma-extension}을 보라.
\end{proof}

\begin{proof}[두 번째 증명]
Proposition \ref{flat-proposition-existence-complete-at-x}를 적용하여
다음 가환 그림을 얻는다.
$$
\xymatrix{
(X, x) \ar[d] & (X', x') \ar[l]^g \ar[d] \\
(S, s) & (S', s') \ar[l]
}
$$
% P05-EN-CH38-0134: 완전 데비사주의 점 x/x' 불일치를 동결 원문대로 보존했다.
이는 점을 표시한 스킴들의 그림이고, 수평 화살표들은 기본 에탈
근방이며, $g^*\mathcal{F}/X'/S'$는 $x$에서 완전 데비사주를 갖는다.
(특히 $S'$와 $X'$는 아핀이다.) Morphisms, Lemma
\ref{morphisms-lemma-flat-permanence}에 의해 $g^*\mathcal{F}$는
$x'$에서 $S$ 위로 평탄이고, Lemma \ref{flat-lemma-etale-flat-up-down}에
의해 $x'$에서 $S'$ 위로도 평탄이다. Remark
\ref{flat-remark-same-notion}을 통하여 다음 몫들의 열
$$
\Gamma(X', g^*\mathcal{F})/
\Gamma(X', \mathcal{O}_{X'})/
\Gamma(S', \mathcal{O}_{S'})
$$
이 $x'$에 대응하는 $\Gamma(X', \mathcal{O}_{X'})$의 소 아이디얼에서
완전 데비사주를 가짐을 얻는다. 이 완전 데비사주를 $s'$에 대응하는
$\Gamma(S', \mathcal{O}_{S'})$의 소 아이디얼에서의 국소환
$\mathcal{O}_{S', s'}$으로 기저변환할 수 있다. 따라서
Lemma \ref{flat-lemma-complete-devissage-flat-finite-type-module}에 의해
% P05-EN-CH38-0135: 정의되지 않은 \mathcal{F}' 표기를 동결 원문대로 보존했다.
$$
\Gamma(X', \mathcal{F}')
\otimes_{\Gamma(S', \mathcal{O}_{S'})}
\mathcal{O}_{S', s'}
$$
는 $\mathcal{O}_{S', s'}$ 위에서 평탄이고
$\Gamma(X', \mathcal{O}_{X'})
\otimes_{\Gamma(S', \mathcal{O}_{S'})} \mathcal{O}_{S', s'}$ 위에서
유한표현이다. 다시 말해 $\mathcal{F}$의
$X' \times_{S'} \Spec(\mathcal{O}_{S', s'})$로의 제한은
유한표현이고 $\mathcal{O}_{S', s'}$ 위에서 평탄이다. 사상
$X' \times_{S'} \Spec(\mathcal{O}_{S', s'})
\to X \times_S \Spec(\mathcal{O}_{S', s'})$
은 에탈이므로(Lemma \ref{flat-lemma-etale-at-point}), 그 상
$V \subset X \times_S \Spec(\mathcal{O}_{S', s'})$는 열린
부분스킴이다. 에탈 강하에 의해 $\mathcal{F}$의 $V$로의 제한은
유한표현이고 $\mathcal{O}_{S', s'}$ 위에서 평탄이다. (사용한 결과는
Morphisms, Lemma \ref{morphisms-lemma-etale-open}, Descent, Lemma
\ref{descent-lemma-finite-presentation-descends}, 그리고 Morphisms,
Lemma \ref{morphisms-lemma-flat-permanence}이다.)
\end{proof}

\begin{lemma}
\label{flat-lemma-open-in-fibre-where-flat}
$f : X \to S$를 국소 유한형인 스킴의 사상이라 하자.
$\mathcal{F}$를 유한형 준연접 $\mathcal{O}_X$-가군이라 하자.
$s \in S$라 하자. 그러면 집합
$$
\{x \in X_s \mid \mathcal{F}\text{가 }x\text{에서 }S\text{ 위로 평탄}\}
$$
은 올 $X_s$에서 열린집합이다.
\end{lemma}

\begin{proof}
% P05-EN-CH38-0136: 명제에서 기호 U를 정의하지 않은 채 증명에서 사용하는 결함을 동결 원문대로 보존했다.
$x \in U$라고 하자. Proposition
\ref{flat-proposition-finite-type-flat-at-point}에서와 같은 기본 에탈 근방
$(S', s') \to (S, s)$와 열린집합
$V \subset X \times_S \Spec(\mathcal{O}_{S', s'})$를 택한다.
$\kappa(s) = \kappa(s')$이므로 $X_{s'} = X_s$임에 유의하자.
$x' \in V \cap X_{s'}$이면 $\mathcal{F}$의
$X \times_S S'$로의 당김은 $x'$에서 $S'$ 위로 평탄하다.
따라서 $\mathcal{F}$는 $x'$에서 $S$ 위로 평탄하다. Morphisms,
Lemma \ref{morphisms-lemma-flat-permanence}를 보라. 다시 말해
$X_s \cap V \subset U$는 $U$에서 $x$의 열린 근방이다.
\end{proof}

\begin{lemma}
\label{flat-lemma-finite-type-flat-at-point}
$f : X \to S$를 스킴의 사상이라 하자.
$\mathcal{F}$를 $X$ 위의 준연접층이라 하자.
$x \in X$의 상을 $s \in S$라 하자.
다음을 가정한다.
\begin{enumerate}
\item $f$는 국소 유한형이다.
\item $\mathcal{F}$는 유한형이다.
\item $\mathcal{F}$는 $x$에서 $S$ 위로 평탄하다.
\end{enumerate}
그러면 기본 에탈 근방 $(S', s') \to (S, s)$와 열린 부분스킴
$$
V \subset X \times_S \Spec(\mathcal{O}_{S', s'})
$$
이 존재한다. 이 부분스킴은 $x$로 가는
$X \times_S \Spec(\mathcal{O}_{S', s'})$의 유일한 점을 포함하고,
$\mathcal{F}$의 $V$로의 당김은 $\mathcal{O}_{S', s'}$ 위에서 평탄하다.
\end{lemma}

\begin{proof}
% P05-EN-CH38-0146: 앞 명제의 ``국소 유한표현''과 이 증명의 ``유한표현'' 불일치를 동결 원문대로 보존했다.
(이 명제와 Proposition \ref{flat-proposition-finite-type-flat-at-point}의
유일한 차이는 $f$가 유한표현이라고 가정하지 않는다는 것이다.)
문제는 $X$와 $S$ 위에서 국소적이므로 $X$와 $S$가 아핀이라고
가정해도 된다. $X = \Spec(B)$, $S = \Spec(A)$로 쓰고
$B = A[x_1, \ldots, x_n]/I$로 쓰자. 다시 말해 닫힌 몰입
$i : X \to \mathbf{A}^n_S$를 얻는다.
$t = i(x) \in \mathbf{A}^n_S$로 쓰자.
$\mathbf{A}^n_S \to S$, 층 $i_*\mathcal{F}$ 및 점 $t$에
Proposition \ref{flat-proposition-finite-type-flat-at-point}를 적용할 수 있다.
그리하여 기본 에탈 근방 $(S', s') \to (S, s)$와 열린 부분스킴
$$
W \subset \mathbf{A}^n_{\mathcal{O}_{S', s'}}
$$
을 얻으며, $i_*\mathcal{F}$의 $W$로의 당김은
$\mathcal{O}_{S', s'}$ 위에서 평탄하다. 이는
$V := W \cap \big(X \times_S \Spec(\mathcal{O}_{S', s'})\big)$
가 원하는 열린 부분스킴이라는 뜻이다.
\end{proof}

\begin{lemma}
\label{flat-lemma-finite-type-flat-along-fibre}
$f : X \to S$를 스킴의 사상이라 하자.
$\mathcal{F}$를 $X$ 위의 준연접층이라 하자.
$s \in S$라 하자.
다음을 가정한다.
\begin{enumerate}
\item $f$는 유한표현이다.
\item $\mathcal{F}$는 유한형이다.
\item $\mathcal{F}$는 올 $X_s$의 모든 점에서 $S$ 위로 평탄하다.
\end{enumerate}
그러면 기본 에탈 근방 $(S', s') \to (S, s)$와 열린 부분스킴
$$
V \subset X \times_S \Spec(\mathcal{O}_{S', s'})
$$
이 존재한다. 이 부분스킴은 올 $X_s = X \times_S s'$를 포함하고,
$\mathcal{F}$의 $V$로의 당김은 유한표현 $\mathcal{O}_V$-가군이며
$\mathcal{O}_{S', s'}$ 위에서 평탄하다.
\end{lemma}

\begin{proof}
각 점 $x \in X_s$에 Proposition
\ref{flat-proposition-finite-type-flat-at-point}를 적용하면 기본 에탈 근방
$(S_x, s_x) \to (S, s)$와 열린집합
$V_x \subset X \times_S \Spec(\mathcal{O}_{S_x, s_x})$를 찾을 수 있다.
여기서 $x \in X_s = X \times_S s_x$는 $V_x$에 들어 있고,
$\mathcal{F}$의 $V_x$로의 당김은 유한표현
$\mathcal{O}_{V_x}$-가군이며 $\mathcal{O}_{S_x, s_x}$ 위에서 평탄하다.
특히 올 $(V_x)_{s_x}$를 $X_s$에서 $x$의 열린 근방으로 볼 수 있다.
$X_s$는 준콤팩트이므로 $X_s$가 $(V_{x_i})_{s_{x_i}}$들의 합집합이
되도록 하는 유한 개의 점 $x_1, \ldots, x_n \in X_s$를 찾을 수 있다.
각 근방 $(S_{x_i}, s_{x_i})$를 지배하는 기본 에탈 근방
$(S' , s') \to (S, s)$를 택한다. More on Morphisms, Lemma
\ref{more-morphisms-lemma-elementary-etale-neighbourhoods}를 보라.
% P05-EN-CH38-0137: 단수 V_i와 복수 inverse images의 문법 불일치는 한국어로 자연스럽게 나타냈다.
$V = \bigcup V_i$로 두자. 여기서 $V_i$는 다음 사상에 의한
열린집합 $V_{x_i}$의 역상이다.
$$
X \times_S \Spec(\mathcal{O}_{S', s'})
\longrightarrow
X \times_S \Spec(\mathcal{O}_{S_{x_i}, s_{x_i}})
$$
구성상 $V$는 $X_s$를 포함하고, $\mathcal{F}$의 $V$로의 당김은
유한표현 $\mathcal{O}_V$-가군이며 $\mathcal{O}_{S', s'}$ 위에서
평탄하다.
\end{proof}

\begin{lemma}
\label{flat-lemma-finite-type-flat-along-fibre-variant}
$f : X \to S$를 스킴의 사상이라 하자.
$\mathcal{F}$를 $X$ 위의 준연접층이라 하자.
$s \in S$라 하자.
다음을 가정한다.
\begin{enumerate}
\item $f$는 유한형이다.
\item $\mathcal{F}$는 유한형이다.
\item $\mathcal{F}$는 올 $X_s$의 모든 점에서 $S$ 위로 평탄하다.
\end{enumerate}
그러면 기본 에탈 근방 $(S', s') \to (S, s)$와 열린 부분스킴
$$
V \subset X \times_S \Spec(\mathcal{O}_{S', s'})
$$
이 존재한다. 이 부분스킴은 올 $X_s = X \times_S s'$를 포함하고,
$\mathcal{F}$의 $V$로의 당김은 $\mathcal{O}_{S', s'}$ 위에서 평탄하다.
\end{lemma}

\begin{proof}
(이 명제와 Lemma \ref{flat-lemma-finite-type-flat-along-fibre}의 유일한
차이는 $f$가 유한표현이라고 가정하지 않는다는 것이다.)
각 점 $x \in X_s$에 Lemma \ref{flat-lemma-finite-type-flat-at-point}를
적용하면 기본 에탈 근방 $(S_x, s_x) \to (S, s)$와 열린집합
$V_x \subset X \times_S \Spec(\mathcal{O}_{S_x, s_x})$를 찾을 수 있다.
여기서 $x \in X_s = X \times_S s_x$는 $V_x$에 들어 있고,
$\mathcal{F}$의 $V_x$로의 당김은 $\mathcal{O}_{S_x, s_x}$ 위에서
평탄하다. 특히 올 $(V_x)_{s_x}$를 $X_s$에서 $x$의 열린 근방으로
볼 수 있다. $X_s$는 준콤팩트이므로 $X_s$가
$(V_{x_i})_{s_{x_i}}$들의 합집합이 되도록 하는 유한 개의 점
$x_1, \ldots, x_n \in X_s$를 찾을 수 있다. 각 근방
$(S_{x_i}, s_{x_i})$를 지배하는 기본 에탈 근방
$(S' , s') \to (S, s)$를 택한다. More on Morphisms, Lemma
\ref{more-morphisms-lemma-elementary-etale-neighbourhoods}를 보라.
% P05-EN-CH38-0138: 단수 V_i와 복수 inverse images의 문법 불일치는 한국어로 자연스럽게 나타냈다.
$V = \bigcup V_i$로 두자. 여기서 $V_i$는 다음 사상에 의한
열린집합 $V_{x_i}$의 역상이다.
$$
X \times_S \Spec(\mathcal{O}_{S', s'})
\longrightarrow
X \times_S \Spec(\mathcal{O}_{S_{x_i}, s_{x_i}})
$$
구성상 $V$는 $X_s$를 포함하고, $\mathcal{F}$의 $V$로의 당김은
$\mathcal{O}_{S', s'}$ 위에서 평탄하다.
\end{proof}

\begin{lemma}
\label{flat-lemma-finite-type-flat-at-point-X}
$S$를 스킴이라 하고 $X$가 $S$ 위에서 국소 유한형이라고 하자.
$x \in X$의 상을 $s \in S$라 하자. $X$가 $x$에서 $S$ 위로
평탄이면, 기본 에탈 근방 $(S', s') \to (S, s)$와 열린 부분스킴
$$
V \subset X \times_S \Spec(\mathcal{O}_{S', s'})
$$
이 존재한다. 이 부분스킴은 $x$로 가는
$X \times_S \Spec(\mathcal{O}_{S', s'})$의 유일한 점을 포함하고,
$V \to \Spec(\mathcal{O}_{S', s'})$는 평탄이고 유한표현이다.
\end{lemma}

\begin{proof}
문제는 $X$와 $S$ 위에서 국소적이므로 $X$와 $S$가 아핀이라고
가정해도 된다. $X = \Spec(B)$, $S = \Spec(A)$로 쓰고
$B = A[x_1, \ldots, x_n]/I$로 쓰자. 다시 말해 닫힌 몰입
$i : X \to \mathbf{A}^n_S$를 얻는다.
$t = i(x) \in \mathbf{A}^n_S$로 쓰자.
$\mathbf{A}^n_S \to S$, 층 $\mathcal{F} = i_*\mathcal{O}_X$ 및
점 $t$에 Proposition \ref{flat-proposition-finite-type-flat-at-point}를
적용할 수 있다. 그리하여 기본 에탈 근방
$(S', s') \to (S, s)$와 열린 부분스킴
$$
W \subset \mathbf{A}^n_{\mathcal{O}_{S', s'}}
$$
을 얻으며, $i_*\mathcal{O}_X$의 당김은 평탄이고 유한표현이다. 이는
$V := W \cap \big(X \times_S \Spec(\mathcal{O}_{S', s'})\big)$
가 원하는 열린 부분스킴이라는 뜻이다.
\end{proof}

\begin{lemma}
\label{flat-lemma-finite-type-flat-at-point-local}
$f : X \to S$를 국소 유한표현인 사상이라 하자.
$\mathcal{F}$를 유한형 준연접 $\mathcal{O}_X$-가군이라 하자.
$x \in X$이고 $\mathcal{F}$가 $x$에서 $S$ 위로 평탄이면,
$\mathcal{F}_x$는 유한표현 $\mathcal{O}_{X, x}$-가군이다.
\end{lemma}

\begin{proof}
$s = f(x)$로 두자. Proposition
\ref{flat-proposition-finite-type-flat-at-point}에 의해 기본 에탈 근방
$(S', s') \to (S, s)$가 존재하여, $\mathcal{F}$의
$X \times_S \Spec(\mathcal{O}_{S', s'})$로의 당김은 $x$에 대응하는
점 $x' \in X_{s'} = X_s$의 한 근방에서 유한표현이다. 환 준동형
$$
\mathcal{O}_{X, x} \longrightarrow
\mathcal{O}_{X \times_S \Spec(\mathcal{O}_{S', s'}), x'}
=
\mathcal{O}_{X \times_S S', x'}
$$
은 에탈 환 준동형의 국소화이므로 평탄 국소 환 준동형이다. 따라서
강하에 의해 $\mathcal{F}_x$는 $\mathcal{O}_{X, x}$ 위에서
유한표현이다. Algebra, Lemma
\ref{algebra-lemma-descend-properties-modules}를 보라. (또한 평탄 국소
환 준동형은 충실히 평탄이다. Algebra, Lemma
\ref{algebra-lemma-local-flat-ff}를 보라.)
\end{proof}

\begin{lemma}
\label{flat-lemma-finite-type-flat-at-point-local-X}
$f : X \to S$를 국소 유한형인 사상이라 하자.
$x \in X$의 상을 $s \in S$라 하자. $f$가 $x$에서 $S$ 위로
평탄이면 $\mathcal{O}_{X, x}$는 $\mathcal{O}_{S, s}$ 위에서
본질적으로 유한표현이다.
\end{lemma}

\begin{proof}
$X$와 $S$가 아핀이라고 가정해도 된다. $X = \Spec(B)$,
$S = \Spec(A)$로 쓰고 $B = A[x_1, \ldots, x_n]/I$로 쓰자.
다시 말해 닫힌 몰입 $i : X \to \mathbf{A}^n_S$를 얻는다.
$t = i(x) \in \mathbf{A}^n_S$로 쓰자.
$\mathbf{A}^n_S \to S$, 층 $\mathcal{F} = i_*\mathcal{O}_X$ 및
점 $t$에 Lemma \ref{flat-lemma-finite-type-flat-at-point-local}을 적용할 수
있다. $\mathcal{O}_{X, x}$가
$\mathcal{O}_{\mathbf{A}^n_S, t}$ 위에서 유한표현이라는 결론을 얻고,
이는 원하는 바를 함의한다.
\end{proof}







\section{열린집합에서 성질의 확장}
\label{flat-section-extending-properties}

\noindent
이 절에서는 다음 형태의 결과들을 모은다. $f : X \to S$가 스킴의
평탄 사상이고 $f$가 $S$의 조밀 열린집합 위에서 어떤 성질을 만족하면,
$f$는 $S$ 전체 위에서 같은 성질을 만족한다.

\begin{lemma}
\label{flat-lemma-flat-finite-type-finitely-presented-over-dense-open}
\begin{slogan}
(좋은) 열린집합 위의 유한표현 가군을 $S$-평탄이고 유한형으로 확장한
가군은 $X$ 위에서도 유한표현이다.
\end{slogan}
$f : X \to S$를 스킴의 사상이라 하자. $\mathcal{F}$를 준연접
$\mathcal{O}_X$-가군이라 하자. $U \subset S$를 열린집합이라 하자.
다음을 가정한다.
\begin{enumerate}
\item $f$는 국소 유한표현이다.
\item $\mathcal{F}$는 유한형이고 $S$ 위에서 평탄하다.
\item $U \subset S$는 역콤팩트이고 스킴론적으로 조밀하다.
\item $\mathcal{F}|_{f^{-1}U}$는 유한표현이다.
\end{enumerate}
그러면 $\mathcal{F}$는 유한표현이다.
\end{lemma}

\begin{proof}
문제는 $X$와 $S$ 위에서 국소적이므로 $X$와 $S$가 아핀이라고
가정해도 된다. $S = \Spec(A)$와 $X = \Spec(B)$로 쓰자.
$\mathcal{F}$가 $N$에 연관된 준연접층이 되도록 유한 $B$-가군
$N$을 택하자. 어떤 $f_i \in A$에 대해
$U = D(f_1) \cup \ldots \cup D(f_n)$이다. Algebra, Lemma
\ref{algebra-lemma-qc-open}을 보라. $U$가 스킴론적으로 조밀하므로 사상
$A \to A_{f_1} \times \ldots \times A_{f_n}$은 단사이다.
$s \in S$로 가는 $x \in X$에 대응하는
$\mathfrak p \subset A$ 위의 소 아이디얼 $\mathfrak q \subset B$를
택하자. Lemma \ref{flat-lemma-finite-type-flat-at-point-local}에 의해
가군 $N_\mathfrak q$는 $B_\mathfrak q$ 위에서 유한표현이다.
$B$-가군의 전사 $\varphi : B^{\oplus m} \to N$을 택하자.
$k_1, \ldots, k_t \in \Ker(\varphi)$를 택하고
$N' = B^{\oplus m}/\sum Bk_j$로 두자. 표준적인 전사 $N' \to N$이
있으며, $N$은 이와 같이 구성한 $B$-가군 $N'$들의 여과 여극한이다.
따라서 (a) $N'_{f_i} \cong N_{f_i}$, $i = 1, \ldots, n$이고
(b) $N'_\mathfrak q \cong N_\mathfrak q$가 되도록
$k_1, \ldots, k_t$를 택할 수 있다. 특히 이는
$N'_\mathfrak q$가 $A$ 위에서 평탄임을 뜻한다. 평탄성의 개방성,
Algebra, Theorem \ref{algebra-theorem-openness-flatness}에 의해
$g \in B$, $g \not \in \mathfrak q$가 존재하여 $N'_g$가 $A$ 위에서
평탄이라는 결론을 얻는다. 다음 가환 그림을 생각하자.
$$
\xymatrix{
N'_g \ar[r] \ar[d] & N_g \ar[d] \\
\prod N'_{gf_i} \ar[r] & \prod N_{gf_i}
}
$$
아래쪽 화살표는 $k_1, \ldots, k_t$의 선택에 의해 동형이다.
$A \to \prod A_{f_i}$가 단사이고 $N'_g$가 $A$ 위에서 평탄하므로
왼쪽 수직 화살표는 단사이다. 따라서 위쪽 수평 화살표는 단사이고,
그러므로 동형이다. 이로써 $N_g$가 $B_g$ 위에서 유한표현임이 증명된다.
Algebra, Lemma \ref{algebra-lemma-cover}를 적용하여 결론을 얻는다.
\end{proof}

\begin{lemma}
\label{flat-lemma-flat-finite-type-finitely-presented-over-dense-open-X}
$f : X \to S$를 스킴의 사상이라 하자. $U \subset S$를 열린집합이라
하자. 다음을 가정한다.
\begin{enumerate}
\item $f$는 국소 유한형이고 평탄하다.
\item $U \subset S$는 역콤팩트이고 스킴론적으로 조밀하다.
\item $f|_{f^{-1}U} : f^{-1}U \to U$는 국소 유한표현이다.
\end{enumerate}
% P05-EN-CH38-0139: 중복된 전치사 of는 한국어로 자연스럽게 나타냈다.
그러면 $f$는 국소 유한표현이다.
\end{lemma}

\begin{proof}
문제는 $X$와 $S$ 위에서 국소적이므로 $X$와 $S$가 아핀이라고
가정해도 된다. 닫힌 몰입 $i : X \to \mathbf{A}^n_S$를 택하고
$i_*\mathcal{O}_X$에 Lemma
\ref{flat-lemma-flat-finite-type-finitely-presented-over-dense-open}를 적용한다.
% P05-EN-CH38-0140: 원문의 명시적 증명 생략을 보충하지 않고 그대로 보존했다.
일부 세부사항은 생략한다.
\end{proof}

\begin{lemma}
\label{flat-lemma-flat-finite-presentation-dimension-over-dense-open}
$f : X \to S$를 평탄이고 국소 유한형인 스킴의 사상이라 하자.
$U \subset S$를 조밀 열린집합으로서 $X_U \to U$가 상대차원
$\leq e$를 갖는 것이라 하자. Morphisms, Definition
\ref{morphisms-definition-relative-dimension-d}를 보라. 더 나아가 다음
중 하나가 성립한다고 하자.
\begin{enumerate}
\item $f$는 국소 유한표현이다.
\item $U \subset S$는 역콤팩트이다.
\end{enumerate}
그러면 $f$는 상대차원 $\leq e$를 갖는다.
\end{lemma}

\begin{proof}
(1)의 경우를 증명하자. More on Morphisms, Lemmas
\ref{more-morphisms-lemma-flat-finite-presentation-CM-open}과
\ref{more-morphisms-lemma-flat-finite-presentation-CM-pieces}에서 구성하고
연구한 열린 부분스킴을 $W \subset X$라 하자. 모든 올의 모든 일반점은
$W$에 들어 있으므로 $W$에 대해 결과를 증명하면 충분하다.
$W = \bigcup_{d \geq 0} U_d$이므로 $d > e$일 때
$U_d = \emptyset$임을 증명하면 충분하다. $f$는 평탄이고 국소
유한표현이므로 열린 사상이고, 따라서 $f(U_d)$는 열린집합이다.
Morphisms, Lemma \ref{morphisms-lemma-fppf-open}을 보라. 그러므로
$U_d$가 비어 있지 않으면 원하는 대로
$f(U_d) \cap U \not = \emptyset$이다.

\medskip\noindent
(2)의 경우를 증명하자. $S$를 그 축약으로 대체해도 된다. 그러면
$U$는 스킴론적으로 조밀하다. 따라서 Lemma
\ref{flat-lemma-flat-finite-type-finitely-presented-over-dense-open-X}에 의해
$f$는 국소 유한표현이다. 이로써 (1)의 경우로 환원된다.
\end{proof}

\begin{lemma}
\label{flat-lemma-proper-flat-finite-over-dense-open}
$f : X \to S$를 평탄 고유인 스킴의 사상이라 하자.
$U \subset S$를 조밀 열린집합으로서 $X_U \to U$가 유한인 것이라 하자.
더 나아가 $f$가 국소 유한표현이거나 $U \subset S$가 역콤팩트이면,
$f$는 유한이다.
\end{lemma}

\begin{proof}
Lemma \ref{flat-lemma-flat-finite-presentation-dimension-over-dense-open}에 의해
$f$의 올들은 차원 0을 갖는다. 따라서 $f$는 준유한이고
(Morphisms, Lemma
\ref{morphisms-lemma-locally-quasi-finite-rel-dimension-0}), 그러므로 유한한
올들을 갖는다(Morphisms, Lemma \ref{morphisms-lemma-quasi-finite}).
따라서 More on Morphisms, Lemma
\ref{more-morphisms-lemma-characterize-finite}에 의해 $f$는 유한이다.
\end{proof}

\begin{lemma}
\label{flat-lemma-zariski}
$f : X \to S$를 스킴의 사상이라 하고 $U \subset S$를 열린집합이라 하자.
다음이 성립한다고 하자.
\begin{enumerate}
\item $f$는 분리이고 국소 유한형이며 평탄하다.
\item $f^{-1}(U) \to U$는 동형이다.
\item $U \subset S$는 역콤팩트이고 스킴론적으로 조밀하다.
\end{enumerate}
그러면 $f$는 열린 몰입이다.
\end{lemma}

\begin{proof}
Lemma \ref{flat-lemma-flat-finite-type-finitely-presented-over-dense-open-X}에
의해 사상 $f$는 국소 유한표현이다. 상 $f(X) \subset S$는 열린집합이므로
(Morphisms, Lemma \ref{morphisms-lemma-fppf-open}), $S$를 $f(X)$로
대체해도 된다. 따라서 $f$가 동형임을 증명해야 한다. $S$가 아핀이라고
가정해도 된다. 상이 $S$인 임의의 준콤팩트 열린집합 $X' \subset X$가
$S$로 동형사상된다는 것을 보이면 충분하므로, $X$가 준콤팩트인 경우로
환원할 수 있다. 따라서 $f$가 준콤팩트라고 가정해도 된다.
Lemma \ref{flat-lemma-flat-finite-presentation-dimension-over-dense-open}에 의해
$f$의 모든 올은 차원 $0$을 갖는다. 따라서 $f$는 준유한이다.
Morphisms, Lemma
\ref{morphisms-lemma-locally-quasi-finite-rel-dimension-0}을 보라.
$s \in S$라 하자. $V \to T$가 유한이고 $W_t = \emptyset$이며
$X \times_S T = V \amalg W$가 되도록 하는 기본 에탈 근방
$g : (T, t) \to (S, s)$를 택한다. More on Morphisms, Lemma
\ref{more-morphisms-lemma-etale-splits-off-quasi-finite-part}를 보라.
주어진 사상을 $\pi : V \amalg W \to T$라 쓰자. $\pi$는 평탄이고
국소 유한표현이므로 $\pi(V)$는 $T$에서 열린집합이다(Morphisms, Lemma
\ref{morphisms-lemma-fppf-open}). $T$를 축소하여
$T = \pi(V)$라고 가정해도 된다. $f$는 $U$ 위에서 동형이므로
$\pi$는 $g^{-1}U$ 위에서 동형이다. $\pi(V) = T$이므로 이는
$\pi^{-1}g^{-1}U$가 $V$에 들어 있음을 뜻한다. Morphisms, Lemma
\ref{morphisms-lemma-flat-morphism-scheme-theoretically-dense-open}에 의해
$\pi^{-1}g^{-1}U \subset V \amalg W$는 스킴론적으로 조밀하다.
따라서 $W = \emptyset$이다. 그러므로 $X \times_S T = V$는 $T$ 위에서
유한이다. 예를 들어 Descent, Lemma
\ref{descent-lemma-descending-property-finite}에 의해 이는 ($S$를 $s$의
한 열린 근방으로 대체한 뒤) $f$가 유한임을 뜻한다. 그러면 $f$는
유한 국소 자유이다(Morphisms, Lemma
\ref{morphisms-lemma-finite-flat}). $S$를 $s$의 더 작은 열린 근방으로
축소하면 $f$가 어떤 차수 $d$의 유한 국소 자유 사상임을 얻는다
(Morphisms, Lemma \ref{morphisms-lemma-finite-locally-free}). 그러나 조밀
열린집합 $U$ 위에서 차수가 $1$이라는 사실로부터 명백히 $d = 1$이다.
따라서 $f$는 동형이다.
\end{proof}





\section{평탄 유한표현 가군}
\label{flat-section-finitely-presented-flat}

\noindent
어떤 경우에는 유한표현 환 준동형 $R \to S$와 유한표현 $S$-가군
$N$이 주어졌을 때, $N$이 $R$ 위에서 평탄이면 적어도 $S$를 에탈
확대로 대체한 뒤에는 $N$이 사영 $R$-가군이다.
% P05-EN-CH38-0141: 중복된 부정관사는 한국어로 자연스럽게 나타냈다.
이 절에서는 이러한 성질의 결과 몇 가지를 모은다.

\begin{lemma}
\label{flat-lemma-induction-step-fp}
$R$을 환이라 하자. $R \to S$를 기하적으로 정역인 올들을 갖는
유한표현 평탄 환 준동형이라 하자. $\mathfrak q \subset S$를
소 아이디얼 $\mathfrak r \subset R$ 위에 놓인 소 아이디얼이라 하자.
$\mathfrak p = \mathfrak r S$로 두자. $N$을 유한표현 $S$-가군이라 하자.
$r \geq 0$과 $S$-가군 사상
$$
\alpha : S^{\oplus r} \longrightarrow N
$$
이 존재하여
$\alpha : \kappa(\mathfrak p)^{\oplus r} \to N \otimes_S \kappa(\mathfrak p)$
는 동형이다. 이러한 임의의 $\alpha$에 대해 다음은 서로 동치이다.
\begin{enumerate}
\item $N_{\mathfrak q}$는 $R$-평탄이다.
\item $f \in R$, $f \not \in \mathfrak r$가 존재하여
$\alpha_f : S_f^{\oplus r} \to N_f$는 $R_f$-보편 단사이고,
$g \in S$, $g \not \in \mathfrak q$가 존재하여
$\Coker(\alpha)_g$는 $R$-평탄이다.
% P05-EN-CH38-0142: 세 번째 항목의 누락된 종결 문장부호는 한국어 목록으로 자연스럽게 나타냈다.
\item $\alpha_{\mathfrak r}$는 $R_{\mathfrak r}$-보편 단사이고
$\Coker(\alpha)_{\mathfrak q}$는 $R$-평탄이다.
\item $\alpha_{\mathfrak r}$는 단사이고
$\Coker(\alpha)_{\mathfrak q}$는 $R$-평탄이다.
\item $\alpha_{\mathfrak p}$는 동형이고
$\Coker(\alpha)_{\mathfrak q}$는 $R$-평탄이다.
\item $\alpha_{\mathfrak q}$는 단사이고
$\Coker(\alpha)_{\mathfrak q}$는 $R$-평탄이다.
\end{enumerate}
\end{lemma}

\begin{proof}
$\alpha$를 얻기 위해
$r = \dim_{\kappa(\mathfrak p)} N \otimes_S \kappa(\mathfrak p)$로 두고,
$N \otimes_S \kappa(\mathfrak p)$의 기저를 이루는
$x_1, \ldots, x_r \in N$을 택한다.
$\alpha(s_1, \ldots, s_r) = \sum s_i x_i$로 정의하자. 이로써
존재성이 증명된다.

\medskip\noindent
$\alpha$ 하나를 고정하자. 사상
$\alpha_{\mathfrak r} : S_{\mathfrak r}^{\oplus r} \to N_{\mathfrak r}$에
Lemma \ref{flat-lemma-induction-step}을 적용할 수 있다. 따라서 (1), (3),
(4), (5), (6)은 모두 서로 동치이다. 또한 (2)가 (3)을 함의함은
명백하므로 (1)이 (2)를 함의함만 보이면 된다.

\medskip\noindent
(1)을 가정하자. 평탄성의 개방성, Algebra, Theorem
\ref{algebra-theorem-openness-flatness}에 의해 집합
$$
U_1 = \{\mathfrak q' \subset S \mid N_{\mathfrak q'}\text{가 }R\text{ 위에서 평탄}\}
$$
은 $\Spec(S)$에서 열린집합이다. 가정에 의해 $\mathfrak q$를 포함하고,
따라서 $\mathfrak p$도 포함한다. $S^{\oplus r}$와 $N$은 유한표현
$S$-가군이므로 집합
$$
U_2 = \{\mathfrak q' \subset S \mid
\alpha_{\mathfrak q'}\text{가 동형}\}
$$
은 $\Spec(S)$에서 열린집합이다. Algebra, Lemma
\ref{algebra-lemma-map-between-finitely-presented}를 보라. (5)에 의해
이는 $\mathfrak p$를 포함한다. $R \to S$는 유한표현이고 평탄하므로
사상 $\Phi : \Spec(S) \to \Spec(R)$는 열린 사상이다. Algebra,
Proposition \ref{algebra-proposition-fppf-open}을 보라.
임의의 소 아이디얼 $\mathfrak r' \in \Phi(U_1 \cap U_2)$에 대해,
$\mathfrak r'$ 위에 놓인 소 아이디얼 $\mathfrak q'$가 존재하여
$N_{\mathfrak q'}$는 평탄이고 $\alpha_{\mathfrak q'}$는 동형이다.
이는 $\mathfrak p' = \mathfrak r' S$일 때
$\alpha \otimes \kappa(\mathfrak p')$가 동형임을 뜻한다. 따라서 함의
(1) $\Rightarrow$ (3)에 의해 $\alpha_{\mathfrak r'}$는
$R_{\mathfrak r'}$-보편 단사이다. 그러므로
$D(f) \subset \Phi(U_1 \cap U_2)$가 되도록 하는
$f \in R$, $f \not \in \mathfrak r$를 택하면 $\alpha_f$가
$R_f$-보편 단사라는 결론을 얻는다. Algebra, Lemma
\ref{algebra-lemma-universally-injective-check-stalks}를 보라.
같은 논증은 임의의
$\mathfrak q' \in U_1 \cap \Phi^{-1}(\Phi(U_1 \cap U_2))$에 대해
가군 $\Coker(\alpha)_{\mathfrak q'}$가 $R$-평탄임도 보인다.
$\mathfrak q \in U_1 \cap \Phi^{-1}(\Phi(U_1 \cap U_2))$임에 유의하자.
따라서 $D(g) \subset U_1 \cap \Phi^{-1}(\Phi(U_1 \cap U_2))$가
되도록 하는 $g \in S$, $g \not \in \mathfrak q$를 찾을 수 있고,
이로써 원하는 결론을 얻는다.
\end{proof}

\begin{lemma}
\label{flat-lemma-complete-devissage-flat-finitely-presented-module}
$R \to S$를 유한표현 환 준동형이라 하자. $N$을 $R$ 위에서 평탄인
유한표현 $S$-가군이라 하자. $\mathfrak r \subset R$을 소 아이디얼이라
하자. $\mathfrak r$ 위에서 $N/S/R$의 완전 데비사주가 존재한다고
가정하자. 그러면 $f \in R$, $f \not \in \mathfrak r$가 존재하여
$$
N_f \cong B_1^{\oplus r_1} \oplus \ldots \oplus B_n^{\oplus r_n}
$$
는 $R$-가군의 동형이다. 여기서 각 $B_i$는 기하적으로 기약인 올들을
갖는 매끄러운 $R_f$-대수이다. 더욱이 $N_f$는 사영 $R_f$-가군이다.
\end{lemma}

\begin{proof}
$(A_i, B_i, M_i, \alpha_i)_{i = 1, \ldots, n}$를 주어진 완전
데비사주라 하자. $n$에 대한 귀납법으로 보조정리를 증명한다.
보조정리의 주장들은 전적으로 $R$-가군으로서의 $N$의 구조에 관한
것임에 유의하자. 따라서 $N$을 $M_1$으로 대체하고 $M_1$을
$B_1$-가군으로 생각해도 된다. $M_1$이 유한표현 $B_1$-가군인 이유는
Remark \ref{flat-remark-finite-presentation}을 보라. Lemma
\ref{flat-lemma-induction-step-fp}에 의해 어떤
$f \in R$, $f \not \in \mathfrak r$에 대해 $R$을 $R_f$로 대체한 뒤,
사상 $\alpha_1 : B_1^{\oplus r_1} \to M_1$이 $R$-보편 단사라고
가정해도 된다. $M_1$과 $B_1^{\oplus r_1}$은 $R$-평탄이고
유한표현 $B_1$-가군이므로 $\Coker(\alpha_1)$은 $R$-평탄이고
(Algebra, Lemma \ref{algebra-lemma-ui-flat-domain}) 유한표현
$B_1$-가군이다. $(A_i, B_i, M_i, \alpha_i)_{i = 2, \ldots, n}$가
$\Coker(\alpha_1)$의 완전 데비사주임에 유의하자. 따라서 귀납가정에
의해 어떤 $f \in R$, $f \not \in \mathfrak r$에 대해 $R$을
$R_f$로 대체한 뒤, $\Coker(\alpha_1)$이 보조정리에서와 같은 분해를
가지며 사영이라고 가정해도 된다. 특히
$M_1 = B_1^{\oplus r_1} \oplus \Coker(\alpha_1)$.
이는 분해에 관한 주장을 증명한다. 사영성에 관한 주장은 Lemma
\ref{flat-lemma-fibres-irreducible-flat-projective-nonnoetherian}에 의해
$B_1$이 사영 $R$-가군이라는 사실에서 따른다.
\end{proof}

\begin{remark}
\label{flat-remark-complete-devissage-flat-finitely-presented-module}
Lemma \ref{flat-lemma-complete-devissage-flat-finitely-presented-module}에는
다음 변형이 있다. 평탄성 조건은 $\mathfrak r$ 위에 놓인 어떤 주어진
소 아이디얼 $\mathfrak q$에서만 $N$이 평탄이라고 가정하는 것으로
약화하는 한편, 데비사주 조건은 {\it $\mathfrak q$에서} 완전
데비사주가 존재한다고 가정하는 것으로 강화한다. Lemma
\ref{flat-lemma-complete-devissage-flat-finite-type-module}와 비교하라.
\end{remark}

\noindent
다음은 이 절의 주결과이다.

\begin{proposition}
\label{flat-proposition-finite-presentation-flat-at-point}
$f : X \to S$를 스킴의 사상이라 하자. $\mathcal{F}$를 $X$ 위의
준연접층이라 하자. $x \in X$의 상을 $s \in S$라 하자.
다음을 가정한다.
\begin{enumerate}
\item $f$는 국소 유한표현이다.
\item $\mathcal{F}$는 유한표현이다.
\item $\mathcal{F}$는 $x$에서 $S$ 위로 평탄하다.
\end{enumerate}
그러면 점을 표시한 스킴들의 다음 가환 그림이 존재한다.
$$
\xymatrix{
(X, x) \ar[d] & (X', x') \ar[l]^g \ar[d] \\
(S, s) & (S', s') \ar[l]
}
$$
여기서 수평 화살표들은 기본 에탈 근방이고, $X'$와 $S'$는 아핀이며,
$\Gamma(X', g^*\mathcal{F})$는 사영
$\Gamma(S', \mathcal{O}_{S'})$-가군이다.
\end{proposition}

\begin{proof}
평탄성의 개방성, More on Morphisms, Theorem
\ref{more-morphisms-theorem-openness-flatness}에 의해 $X$를 $x$의 열린
근방으로 대체하고 $\mathcal{F}$가 $S$ 위에서 평탄이라고 가정해도 된다.
이제 Proposition \ref{flat-proposition-existence-complete-at-x}를 적용하여
명제의 주장과 같은 그림으로서 $g^*\mathcal{F}/X'/S'$가 $s'$ 위에서
완전 데비사주를 갖는 것을 찾는다. (특히 $S'$와 $X'$는 아핀이다.)
Morphisms, Lemma \ref{morphisms-lemma-flat-permanence}에 의해
$g^*\mathcal{F}$는 $S$ 위에서 평탄이고, Lemma
\ref{flat-lemma-etale-flat-up-down}에 의해 $S'$ 위에서도 평탄이다.
Remark \ref{flat-remark-same-notion}을 통하여 다음 몫들의 열
$$
\Gamma(X', g^*\mathcal{F})/
\Gamma(X', \mathcal{O}_{X'})/
\Gamma(S', \mathcal{O}_{S'})
$$
이 $s'$에 대응하는 $\Gamma(S', \mathcal{O}_{S'})$의 소 아이디얼 위에서
완전 데비사주를 가짐을 얻는다. 따라서 Lemma
\ref{flat-lemma-complete-devissage-flat-finitely-presented-module}에 의해
$S'$를 $s'$의 표준 열린 근방으로 대체한 뒤 명제의 결과가 성립한다.
\end{proof}

\noindent
이 절의 나머지에서는 이 결과의 몇 가지 변형을 증명한다. 첫 번째는
``대역적'' 변형이다.

\begin{lemma}
\label{flat-lemma-finite-presentation-flat-along-fibre}
$f : X \to S$를 스킴의 사상이라 하자. $\mathcal{F}$를 $X$ 위의
준연접층이라 하자. $s \in S$라 하자. 다음을 가정한다.
\begin{enumerate}
\item $f$는 유한표현이다.
\item $\mathcal{F}$는 유한표현이다.
\item $\mathcal{F}$는 올 $X_s$의 모든 점에서 $S$ 위로 평탄하다.
\end{enumerate}
그러면 기본 에탈 근방 $(S', s') \to (S, s)$와 스킴들의 다음 가환
그림이 존재한다.
$$
\xymatrix{
X \ar[d] & X' \ar[l]^g \ar[d] \\
S & S' \ar[l]
}
$$
여기서 $g$는 에탈이고 $X_s \subset g(X')$이며, 스킴 $X'$와 $S'$는
아핀이고, $\Gamma(X', g^*\mathcal{F})$는 사영
$\Gamma(S', \mathcal{O}_{S'})$-가군이다.
\end{lemma}

\begin{proof}
각 점 $x \in X_s$에 Proposition
\ref{flat-proposition-finite-presentation-flat-at-point}를 적용하여 다음 가환
그림을 찾을 수 있다.
$$
\xymatrix{
(X, x) \ar[d] & (Y_x, y_x) \ar[l]^{g_x} \ar[d] \\
(S, s) & (S_x, s_x) \ar[l]
}
$$
여기서 수평 화살표들은 기본 에탈 근방이고, $Y_x$와 $S_x$는 아핀이며,
$\Gamma(Y_x, g_x^*\mathcal{F})$는 사영
$\Gamma(S_x, \mathcal{O}_{S_x})$-가군이다. 특히
$g_x(Y_x) \cap X_s$는 $X_s$에서 $x$의 열린 근방이다. $X_s$는
준콤팩트이므로 $X_s$가 $g_{x_i}(Y_{x_i}) \cap X_s$들의 합집합이
되도록 하는 유한 개의 점 $x_1, \ldots, x_n \in X_s$를 찾을 수 있다.
각 근방 $(S_{x_i}, s_{x_i})$를 지배하는 기본 에탈 근방
$(S' , s') \to (S, s)$를 택한다. More on Morphisms, Lemma
\ref{more-morphisms-lemma-elementary-etale-neighbourhoods}를 보라.
$S'$가 아핀이라고도 가정해도 된다.
$X' = \coprod Y_{x_i} \times_{S_{x_i}} S'$로 두고 자명한 사상
$g : X' \to X$를 부여하자. 구성상 $g(X')$는 $X_s$를 포함하고
$$
\Gamma(X', g^*\mathcal{F})
=
\bigoplus \Gamma(Y_{x_i}, g_{x_i}^*\mathcal{F})
\otimes_{\Gamma(S_{x_i}, \mathcal{O}_{S_{x_i}})}
\Gamma(S', \mathcal{O}_{S'}).
$$
이는 사영 $\Gamma(S', \mathcal{O}_{S'})$-가군이다. Algebra, Lemma
\ref{algebra-lemma-ascend-properties-modules}를 보라.
\end{proof}

\noindent
다음 두 보조정리는 $\mathcal{F} = \mathcal{O}_X$인 경우 위 결과들을
다시 서술한 것이다.

\begin{lemma}
\label{flat-lemma-finite-presentation-flat-at-point-X}
$f : X \to S$를 국소 유한표현이라 하자. $x \in X$의 상을
$s \in S$라 하자. $f$가 $x$에서 $S$ 위로 평탄이면, 점을 표시한
스킴들의 다음 가환 그림이 존재한다.
$$
\xymatrix{
(X, x) \ar[d] & (X', x') \ar[l]^g \ar[d] \\
(S, s) & (S', s') \ar[l]
}
$$
여기서 수평 화살표들은 기본 에탈 근방이고, $X'$와 $S'$는 아핀이며,
$\Gamma(X', \mathcal{O}_{X'})$는 사영
$\Gamma(S', \mathcal{O}_{S'})$-가군이다.
\end{lemma}

\begin{proof}
이는 Proposition \ref{flat-proposition-finite-presentation-flat-at-point}의
특별한 경우이다.
\end{proof}

\begin{lemma}
\label{flat-lemma-finite-presentation-flat-along-fibre-X}
$f : X \to S$를 유한표현이라 하자. $s \in S$라 하자. $X$가
$X_s$의 모든 점에서 $S$ 위로 평탄이면, 기본 에탈 근방
$(S', s') \to (S, s)$와 스킴들의 다음 가환 그림이 존재한다.
$$
\xymatrix{
X \ar[d] & X' \ar[l]^g \ar[d] \\
S & S' \ar[l]
}
$$
여기서 $g$는 에탈이고 $X_s \subset g(X')$이며, $X'$와 $S'$는
아핀이고, $\Gamma(X', \mathcal{O}_{X'})$는 사영
$\Gamma(S', \mathcal{O}_{S'})$-가군이다.
\end{lemma}

\begin{proof}
이는 Lemma \ref{flat-lemma-finite-presentation-flat-along-fibre}의 특별한
경우이다.
\end{proof}

\noindent
다음 보조정리들은 사상과 층이 유한형이라고만 (반드시 유한표현이라고는
하지 않고) 가정하는 경우 Proposition
\ref{flat-proposition-finite-presentation-flat-at-point}의 귀결을 설명한다.

\begin{lemma}
\label{flat-lemma-finite-type-flat-at-point-free}
$f : X \to S$를 스킴의 사상이라 하자. $\mathcal{F}$를 $X$ 위의
준연접층이라 하자. $x \in X$의 상을 $s \in S$라 하자.
다음을 가정한다.
\begin{enumerate}
\item $f$는 국소 유한표현이다.
\item $\mathcal{F}$는 유한형이다.
\item $\mathcal{F}$는 $x$에서 $S$ 위로 평탄하다.
\end{enumerate}
그러면 기본 에탈 근방 $(S', s') \to (S, s)$와 점을 표시한 스킴들의
다음 가환 그림이 존재한다.
$$
\xymatrix{
(X, x) \ar[d] & (X', x') \ar[l]^g \ar[d] \\
(S, s) & (\Spec(\mathcal{O}_{S', s'}), s') \ar[l]
}
$$
여기서 $X' \to X \times_S \Spec(\mathcal{O}_{S', s'})$는 에탈이고,
$\kappa(x) = \kappa(x')$이며, 스킴 $X'$는
$\mathcal{O}_{S', s'}$ 위 유한표현 아핀 스킴이고, 층
$g^*\mathcal{F}$는 $\mathcal{O}_{X'}$ 위에서 유한표현이며,
$\Gamma(X', g^*\mathcal{F})$는 자유 $\mathcal{O}_{S', s'}$-가군이다.
\end{lemma}

\begin{proof}
보조정리를 증명하기 위해 $(S, s)$를 임의의 기본 에탈 근방으로
대체해도 되고, $S$를 $\Spec(\mathcal{O}_{S, s})$로 대체해도 된다.
따라서 Proposition \ref{flat-proposition-finite-type-flat-at-point}에 의해
$x$의 한 근방에서 $\mathcal{F}$가 유한표현이고 $S$ 위에서 평탄이라고
가정해도 된다. 이 경우 결과는 Proposition
\ref{flat-proposition-finite-presentation-flat-at-point}에서 따른다. 실제로
Algebra, Theorem \ref{algebra-theorem-projective-free-over-local-ring}에
의해 국소환 위에서는 사영 $=$ 자유이다.
\end{proof}

\begin{lemma}
\label{flat-lemma-finite-type-flat-at-point-free-variant}
$f : X \to S$를 스킴의 사상이라 하자. $\mathcal{F}$를 $X$ 위의
준연접층이라 하자. $x \in X$의 상을 $s \in S$라 하자.
다음을 가정한다.
\begin{enumerate}
\item $f$는 국소 유한형이다.
\item $\mathcal{F}$는 유한형이다.
\item $\mathcal{F}$는 $x$에서 $S$ 위로 평탄하다.
\end{enumerate}
그러면 기본 에탈 근방 $(S', s') \to (S, s)$와 점을 표시한 스킴들의
다음 가환 그림이 존재한다.
$$
\xymatrix{
(X, x) \ar[d] & (X', x') \ar[l]^g \ar[d] \\
(S, s) & (\Spec(\mathcal{O}_{S', s'}), s') \ar[l]
}
$$
여기서 $X' \to X \times_S \Spec(\mathcal{O}_{S', s'})$는 에탈이고,
$\kappa(x) = \kappa(x')$이며, 스킴 $X'$는 아핀이고,
$\Gamma(X', g^*\mathcal{F})$는 자유 $\mathcal{O}_{S', s'}$-가군이다.
\end{lemma}

\begin{proof}
% P05-EN-CH38-0147: 앞 보조정리의 ``국소 유한표현''에서 국소 한정어가 빠진 원문을 그대로 보존했다.
(이 보조정리와 Lemma \ref{flat-lemma-finite-type-flat-at-point-free}의 유일한
차이는 $f$가 유한표현이라고 가정하지 않는다는 것이다.) 문제는
$X$와 $S$ 위에서 국소적이다. 따라서 $X$와 $S$가 아핀이라고
가정해도 된다. $X = \Spec(B)$와 $S = \Spec(A)$로 쓰자. $B$는
유한형 $A$-대수이므로 전사 $A[x_1, \ldots, x_n] \to B$를 찾을 수
있다. 다시 말해 닫힌 몰입 $i : X \to \mathbf{A}^n_S$를 택할 수 있다.
$t = i(x)$와 $\mathcal{G} = i_*\mathcal{F}$로 두자.
% P05-EN-CH38-0143: 단수 동형 뒤의 복수 동사 불일치는 한국어로 자연스럽게 나타냈다.
$\mathcal{G}_t \cong \mathcal{F}_x$는 $\mathcal{O}_{S, s}$-가군의
동형임에 유의하자. 따라서 $\mathcal{G}$는 $t$에서 $S$ 위로 평탄하다.
사상 $\mathbf{A}^n_S \to S$, 점 $t$, 층 $\mathcal{G}$에 Lemma
\ref{flat-lemma-finite-type-flat-at-point-free}를 적용한다. 그리하여 기본
에탈 근방 $(S', s') \to (S, s)$와 점을 표시한 스킴들의 다음 가환
그림을 찾을 수 있다.
$$
\xymatrix{
(\mathbf{A}^n_S, t) \ar[d] & (Y, y) \ar[l]^h \ar[d] \\
(S, s) & (\Spec(\mathcal{O}_{S', s'}), s') \ar[l]
}
$$
여기서 $Y \to \mathbf{A}^n_{\mathcal{O}_{S', s'}}$는 에탈이고,
$\kappa(t) = \kappa(y)$이며, 스킴 $Y$는 아핀이고,
% P05-EN-CH38-0144: 앞 보조정리의 자유 결론을 사영으로 약화한 원문을 그대로 보존했다.
$\Gamma(Y, h^*\mathcal{G})$는 사영 $\mathcal{O}_{S', s'}$-가군이다.
그러면 원래 문제의 해는 $Y$의 닫힌 부분스킴
$X' = Y \times_{\mathbf{A}^n_S} X$로 주어진다.
\end{proof}

\begin{lemma}
\label{flat-lemma-finite-type-flat-along-fibre-free}
$f : X \to S$를 스킴의 사상이라 하자. $\mathcal{F}$를 $X$ 위의
준연접층이라 하자. $s \in S$라 하자. 다음을 가정한다.
\begin{enumerate}
\item $f$는 유한표현이다.
\item $\mathcal{F}$는 유한형이다.
\item $\mathcal{F}$는 $X_s$의 모든 점에서 $S$ 위로 평탄하다.
\end{enumerate}
그러면 기본 에탈 근방 $(S', s') \to (S, s)$와 스킴들의 다음 가환
그림이 존재한다.
$$
\xymatrix{
X \ar[d] & X' \ar[l]^g \ar[d] \\
S & \Spec(\mathcal{O}_{S', s'}) \ar[l]
}
$$
여기서 $X' \to X \times_S \Spec(\mathcal{O}_{S', s'})$는 에탈이고,
$X_s = g((X')_{s'})$이며, 스킴 $X'$는 $\mathcal{O}_{S', s'}$ 위
유한표현 아핀 스킴이고, 층 $g^*\mathcal{F}$는
$\mathcal{O}_{X'}$ 위에서 유한표현이며,
$\Gamma(X', g^*\mathcal{F})$는 자유 $\mathcal{O}_{S', s'}$-가군이다.
\end{lemma}

\begin{proof}
각 점 $x \in X_s$에 Lemma \ref{flat-lemma-finite-type-flat-at-point-free}를
적용하여 기본 에탈 근방 $(S_x , s_x) \to (S, s)$와 다음 가환
그림을 찾을 수 있다.
$$
\xymatrix{
(X, x) \ar[d] & (Y_x, y_x) \ar[l]^{g_x} \ar[d] \\
(S, s) & (\Spec(\mathcal{O}_{S_x, s_x}), s_x) \ar[l]
}
$$
여기서 $Y_x \to X \times_S \Spec(\mathcal{O}_{S_x, s_x})$는 에탈이고,
$\kappa(x) = \kappa(y_x)$이며, 스킴 $Y_x$는
$\mathcal{O}_{S_x, s_x}$ 위 유한표현 아핀 스킴이고, 층
$g_x^*\mathcal{F}$는 $\mathcal{O}_{Y_x}$ 위에서 유한표현이며,
$\Gamma(Y_x, g_x^*\mathcal{F})$는 자유
$\mathcal{O}_{S_x, s_x}$-가군이다. 특히 $g_x((Y_x)_{s_x})$는
$X_s$에서 $x$의 열린 근방이다. $X_s$는 준콤팩트이므로 $X_s$가
$g_{x_i}((Y_{x_i})_{s_{x_i}})$들의 합집합이 되도록 하는 유한 개의 점
$x_1, \ldots, x_n \in X_s$를 찾을 수 있다. 각 근방
$(S_{x_i}, s_{x_i})$를 지배하는 기본 에탈 근방
$(S' , s') \to (S, s)$를 택한다. More on Morphisms, Lemma
\ref{more-morphisms-lemma-elementary-etale-neighbourhoods}를 보라. 다음과 같이 두자.
$$
X' = \coprod Y_{x_i} \times_{\Spec(\mathcal{O}_{S_{x_i}, s_{x_i}})}
\Spec(\mathcal{O}_{S', s'})
$$
그리고 자명한 사상 $g : X' \to X$를 부여하자. 구성상
$X_s = g(X'_{s'})$이고
$$
\Gamma(X', g^*\mathcal{F})
=
\bigoplus \Gamma(Y_{x_i}, g_{x_i}^*\mathcal{F})
\otimes_{\mathcal{O}_{S_{x_i}, s_{x_i}}}
\mathcal{O}_{S', s'}.
$$
이는 자유 가군들의 기저변환들의 직합이므로 자유
$\mathcal{O}_{S', s'}$-가군이다.
% P05-EN-CH38-0145: 원문의 명시적 사소한 세부사항 생략을 보충하지 않고 그대로 보존했다.
몇 가지 사소한 세부사항은 생략한다.
\end{proof}

\begin{lemma}
\label{flat-lemma-finite-type-flat-along-fibre-free-variant}
$f : X \to S$를 스킴의 사상이라 하자. $\mathcal{F}$를 $X$ 위의
준연접층이라 하자. $s \in S$라 하자. 다음을 가정한다.
\begin{enumerate}
\item $f$는 유한형이다.
\item $\mathcal{F}$는 유한형이다.
\item $\mathcal{F}$는 $X_s$의 모든 점에서 $S$ 위로 평탄하다.
\end{enumerate}
그러면 기본 에탈 근방 $(S', s') \to (S, s)$와 스킴들의 다음 가환
그림이 존재한다.
$$
\xymatrix{
X \ar[d] & X' \ar[l]^g \ar[d] \\
S & \Spec(\mathcal{O}_{S', s'}) \ar[l]
}
$$
여기서 $X' \to X \times_S \Spec(\mathcal{O}_{S', s'})$는 에탈이고,
$X_s = g((X')_{s'})$이며, 스킴 $X'$는 아핀이고,
$\Gamma(X', g^*\mathcal{F})$는 자유 $\mathcal{O}_{S', s'}$-가군이다.
\end{lemma}

\begin{proof}
(이 보조정리와 Lemma \ref{flat-lemma-finite-type-flat-along-fibre-free}의
유일한 차이는 $f$가 유한표현이라고 가정하지 않는다는 것이다.)
각 점 $x \in X_s$에 Lemma
\ref{flat-lemma-finite-type-flat-at-point-free-variant}를 적용하여 기본 에탈
근방 $(S_x , s_x) \to (S, s)$와 다음 가환 그림을 찾을 수 있다.
$$
\xymatrix{
(X, x) \ar[d] & (Y_x, y_x) \ar[l]^{g_x} \ar[d] \\
(S, s) & (\Spec(\mathcal{O}_{S_x, s_x}), s_x) \ar[l]
}
$$
여기서 $Y_x \to X \times_S \Spec(\mathcal{O}_{S_x, s_x})$는 에탈이고,
$\kappa(x) = \kappa(y_x)$이며, 스킴 $Y_x$는 아핀이고,
$\Gamma(Y_x, g_x^*\mathcal{F})$는 자유
$\mathcal{O}_{S_x, s_x}$-가군이다. 특히 $g_x((Y_x)_{s_x})$는
$X_s$에서 $x$의 열린 근방이다. $X_s$는 준콤팩트이므로 $X_s$가
$g_{x_i}((Y_{x_i})_{s_{x_i}})$들의 합집합이 되도록 하는 유한 개의 점
$x_1, \ldots, x_n \in X_s$를 찾을 수 있다. 각 근방
$(S_{x_i}, s_{x_i})$를 지배하는 기본 에탈 근방
$(S' , s') \to (S, s)$를 택한다. More on Morphisms, Lemma
\ref{more-morphisms-lemma-elementary-etale-neighbourhoods}를 보라. 다음과 같이 두자.
$$
X' = \coprod Y_{x_i} \times_{\Spec(\mathcal{O}_{S_{x_i}, s_{x_i}})}
\Spec(\mathcal{O}_{S', s'})
$$
그리고 자명한 사상 $g : X' \to X$를 부여하자. 구성상
$X_s = g(X'_{s'})$이고
$$
\Gamma(X', g^*\mathcal{F})
=
\bigoplus \Gamma(Y_{x_i}, g_{x_i}^*\mathcal{F})
\otimes_{\mathcal{O}_{S_{x_i}, s_{x_i}}}
\mathcal{O}_{S', s'}.
$$
이는 자유 가군들의 기저변환들의 직합이므로 자유
$\mathcal{O}_{S', s'}$-가군이다.
\end{proof}




\section{유한형 평탄 가군 II}
\label{flat-section-finite-type-flat-II}

\noindent
다음 보조정리가 필요하다.

\begin{lemma}
\label{flat-lemma-weak-bourbaki-pre-pre}
$R \to S$를 유한표현 환 준동형이라 하자. $N$을 유한표현
$S$-가군이라 하자. $\mathfrak q \subset S$를
$\mathfrak p \subset R$ 위에 놓인 소 아이디얼이라 하자. 다음과 같이 두자.
$\overline{S} = S \otimes_R \kappa(\mathfrak p)$,
$\overline{\mathfrak q} = \mathfrak q \overline{S}$,
$\overline{N} = N \otimes_R \kappa(\mathfrak p)$. 그러면
$g \in S$, $g \not \in \mathfrak q$를 찾아서,
$\mathfrak r \not \subset \overline{\mathfrak q}$인 모든
$\mathfrak r \in \text{Ass}_{\overline{S}}(\overline{N})$에 대해
$\overline{g} \in \mathfrak r$가 되게 할 수 있다.
\end{lemma}

\begin{proof}
실제로 $\text{Ass}_{\overline{S}}(\overline{N}) =
\{\mathfrak r_1, \ldots, \mathfrak r_n\}$라고 하자(유한성은 Algebra,
Lemma \ref{algebra-lemma-finite-ass}에 따른다). 번호를 다시 붙여 다음을
가정해도 된다.
$$
\mathfrak r_1 \subset \overline{\mathfrak q},
\ldots,
\mathfrak r_r \subset \overline{\mathfrak q}, \quad
\mathfrak r_{r + 1} \not \subset \overline{\mathfrak q},
\ldots,
\mathfrak r_n \not \subset \overline{\mathfrak q}
$$
$\overline{\mathfrak q}$는 소 아이디얼이므로 곱
$\mathfrak r_{r + 1} \ldots \mathfrak r_n$은
$\overline{\mathfrak q}$에 포함되지 않는다. 따라서
$\mathfrak r_{r + 1}, \ldots, \mathfrak r_n$ 각각에는 들어 있지만
$\overline{\mathfrak q}$에는 들어 있지 않은 $\overline{S}$의 원소
$a$를 택할 수 있다. $a$로 가는 $g \in S$가 존재하면 이 $g$가
조건을 만족한다. 일반적으로는 영이 아닌 원소
$\lambda \in \kappa(\mathfrak p)$를 찾아서 $\lambda a$가 어떤
$g \in S$의 상이 되게 할 수 있다.
\end{proof}

\noindent
다음 보조정리에는 아래의 조금 더 강한 변형 Lemma
\ref{flat-lemma-weak-bourbaki}가 있다.

\begin{lemma}
\label{flat-lemma-weak-bourbaki-pre}
$R \to S$를 유한표현 환 준동형이라 하자. $N$을 $R$-가군으로서
평탄인 유한표현 $S$-가군이라 하자. $M$을 $R$-가군이라 하자.
$\mathfrak q$를 $\mathfrak p \subset R$ 위에 놓인 $S$의 소
아이디얼이라 하자. 그러면
$$
\mathfrak q \in \text{WeakAss}_S(M \otimes_R N)
\Leftrightarrow
\Big(
\mathfrak p \in \text{WeakAss}_R(M)
\text{ 이고 }
\overline{\mathfrak q} \in \text{Ass}_{\overline{S}}(\overline{N})
\Big)
$$
이다. 여기서 $\overline{S} = S \otimes_R \kappa(\mathfrak p)$,
$\overline{\mathfrak q} = \mathfrak q \overline{S}$,
$\overline{N} = N \otimes_R \kappa(\mathfrak p)$이다.
\end{lemma}

\begin{proof}
Lemma \ref{flat-lemma-weak-bourbaki-pre-pre}에서와 같은 $g \in S$를 택하자.
스킴의 사상 $\Spec(S_g) \to \Spec(R)$, $N_g$에 연관된 준연접 가군,
그리고 소 아이디얼 $\mathfrak qS_g$와 $\mathfrak p$에 대응하는 점들에
Proposition \ref{flat-proposition-finite-presentation-flat-at-point}를 적용한다.
대수학으로 번역하면 다음 환들의 가환 그림을 얻는다.
$$
\xymatrix{
S \ar[r] & S_g \ar[r] & S' \\
& R \ar[lu] \ar[u] \ar[r] & R' \ar[u]
}
\quad\quad
\xymatrix{
\mathfrak q \ar@{-}[r] \ar@{-}[rd] &
\mathfrak qS_g \ar@{-}[d] \ar@{-}[r] & \mathfrak q' \ar@{-}[d] \\
& \mathfrak p \ar@{-}[r] & \mathfrak p'
}
$$
표시한 대로 소 아이디얼들이 부여되어 있고 수평 화살표들은 에탈이며,
$N \otimes_S S'$는 사영 $R'$-가군이다. 다음과 같이 두자.
$N' = N \otimes_S S'$, $M' = M \otimes_R R'$,
% P05-EN-CH38-0148: \overline{S}'의 잉여체에 q'를 쓴 원문을 p'로 고치지 않고 보존했다.
$\overline{S}' = S' \otimes_{R'} \kappa(\mathfrak q')$,
$\overline{\mathfrak q}' = \mathfrak q' \overline{S}'$,
그리고
$$
\overline{N}' = N' \otimes_{R'} \kappa(\mathfrak p') =
\overline{N} \otimes_{\overline{S}} \overline{S}'
$$
Lemma \ref{flat-lemma-etale-weak-assassin-up-down}에 의해 다음을 얻는다.
\begin{align*}
\text{WeakAss}_{S'}(M' \otimes_{R'} N') & =
(\Spec(S') \to \Spec(S))^{-1}\text{WeakAss}_S(M \otimes_R N) \\
\text{WeakAss}_{R'}(M') & =
(\Spec(R') \to \Spec(R))^{-1}\text{WeakAss}_R(M) \\
\text{Ass}_{\overline{S}'}(\overline{N}') & =
(\Spec(\overline{S}') \to \Spec(\overline{S}))^{-1}
\text{Ass}_{\overline{S}}(\overline{N})
\end{align*}
$\overline{N}$과 $\overline{N}'$에 Algebra, Lemma
\ref{algebra-lemma-ass-weakly-ass}를 사용하자. 특히 다음을 얻는다.
\begin{align*}
\mathfrak q \in \text{WeakAss}_S(M \otimes_R N)
& \Leftrightarrow
\mathfrak q' \in \text{WeakAss}_{S'}(M' \otimes_{R'} N') \\
\mathfrak p \in \text{WeakAss}_R(M)
& \Leftrightarrow
\mathfrak p' \in \text{WeakAss}_{R'}(M') \\
\overline{\mathfrak q} \in \text{Ass}_{\overline{S}}(\overline{N})
& \Leftrightarrow
\overline{\mathfrak q}' \in \text{WeakAss}_{\overline{S}'}(\overline{N}')
\end{align*}
$g$의 신중한 선택과 위의
$\text{Ass}_{\overline{S}'}(\overline{N}')$에 대한 공식은 다음을 보인다.
\begin{equation}
\label{flat-equation-key-observation}
\text{만약 }\mathfrak r' \in \text{Ass}_{\overline{S}'}(\overline{N}')
\text{가 }\mathfrak r \subset \overline{S}\text{ 위에 놓이면 }
\mathfrak r \subset \overline{\mathfrak q}
\end{equation}
이는 증명의 뒤에서 핵심 관찰이 된다. 이제부터 따로 언급하지 않고
Algebra, Lemma \ref{algebra-lemma-weakly-ass-local}의 약연관 소 아이디얼
특성화를 사용한다.

\medskip\noindent
$\overline{\mathfrak q} \not \in
\text{Ass}_{\overline{S}}(\overline{N})$라고 하자. 그러면
$\overline{\mathfrak q}' \not \in
\text{Ass}_{\overline{S}'}(\overline{N}')$이다. Algebra, Lemmas
\ref{algebra-lemma-ass-zero-divisors}, \ref{algebra-lemma-finite-ass},
\ref{algebra-lemma-silly}에 의해 $\overline{N}'$ 위의 영인자가 아닌 원소
$\overline{a}' \in \overline{\mathfrak q}'$가 존재한다. 어떤 영이 아닌
$\lambda \in \kappa(\mathfrak p)$에 대해 $\overline{a}'$를
$\lambda \overline{a}'$로 대체하면 $\overline{a}'$로 가는
$a' \in \mathfrak q'$를 찾을 수 있다. Lemma
\ref{flat-lemma-invert-universally-injective}에 의해 사상
$a' : N'_{\mathfrak p'} \to N'_{\mathfrak p'}$는
$R'_{\mathfrak p'}$-보편 단사이다. 특히 사상
$a' : M' \otimes_{R'} N' \to M' \otimes_{R'} N'$는
$\mathfrak p'$에서 국소화한 뒤 단사이고, 따라서 $\mathfrak q'$에서
국소화한 뒤에도 단사이다. 이는 명백히
$\mathfrak q' \not \in \text{WeakAss}_{S'}(M' \otimes_{R'} N')$를
뜻한다. 따라서 $\mathfrak q \in \text{WeakAss}_S(M \otimes_R N)$이면
$\overline{\mathfrak q} \in \text{Ass}_{\overline{S}}(\overline{N})$이다.

\medskip\noindent
$\mathfrak q \in \text{WeakAss}_S(M \otimes_R N)$이라고 가정하자.
% P05-EN-CH38-0149: R-가군 M의 약연관 소 아이디얼에 S를 쓴 원문을 그대로 보존했다.
$\mathfrak p \in \text{WeakAss}_S(M)$임을 보이고자 한다.
$\mathfrak q$가 $J = \text{Ann}_S(z)$ 위에서 극소가 되도록 하는 원소
$z \in M \otimes_R N$을 택하자. $f_i \in \mathfrak p$, $i \in I$를
아이디얼 $\mathfrak p$의 생성원 집합이라 하자. $\mathfrak q$는
$\mathfrak p$ 위에 놓이므로 각 $i$에 대해 $n_i \geq 1$과
$g_i \in S$, $g_i \not \in \mathfrak q$를 택하여
$g_i f_i^{n_i} \in J$, 즉 $g_i f_i^{n_i} z = 0$이 되게 할 수 있다.
$z$의 상을 $z' \in (M' \otimes_{R'} N')_{\mathfrak p'}$라 하자.
$z$는 $(M \otimes_R N)_\mathfrak q$에서 영이 아닌 상을 갖고
$S_\mathfrak q \to S'_{\mathfrak q'}$는 충실히 평탄이므로 $z'$는
영이 아니다. $f_i^{n_i} z' = 0$이라고 주장한다.

\medskip\noindent
주장을 증명하자. $g_i$의 상을 $g'_i \in S'$라 하자. 핵심 관찰
(\ref{flat-equation-key-observation})에 의해 상
$\overline{g}'_i \in \overline{S}'$는 $\mathfrak r'$에 들어 있지 않는다. 여기서
% P05-EN-CH38-0150: \overline{S}' 위 연관 소 아이디얼 식의 \overline{N}을 원문대로 보존했다.
$\mathfrak r' \in \text{Ass}_{\overline{S}'}(\overline{N})$는 임의이다.
따라서 Lemma \ref{flat-lemma-invert-universally-injective}에 의해
$g'_i : N'_{\mathfrak p'} \to N'_{\mathfrak p'}$는
$R'_{\mathfrak p'}$-보편 단사이다. 특히 사상
$g'_i : M' \otimes_{R'} N' \to M' \otimes_{R'} N'$는
% P05-EN-CH38-0151: localizating 철자 오류는 한국어로 자연스럽게 나타냈다.
$\mathfrak p'$에서 국소화한 뒤 단사이다. $g_i f_i^{n_i} z' = 0$이므로
주장이 따른다.

\medskip\noindent
우리의 주장은 $R'_{\mathfrak p'}$에서 $z'$의 소멸자가 원소들
$f_i^{n_i}$를 포함함을 보인다. $R \to R'$는 에탈이므로 Algebra,
Lemma \ref{algebra-lemma-etale-at-prime}에 의해
$\mathfrak p'R'_{\mathfrak p'} = \mathfrak pR'_{\mathfrak p'}$이다.
따라서 $R'_{\mathfrak p'}$에서 $z'$의 소멸자의 근기는
% P05-EN-CH38-0152: p'의 주변환에 빠진 프라임을 동결 원문대로 보존했다.
$\mathfrak p' R_{\mathfrak p'}$와 같다(여기서 $z'$가 영이 아님을
사용한다). 한편
$$
z' \in (M' \otimes_{R'} N')_{\mathfrak p'} =
M'_{\mathfrak p'} \otimes_{R'_{\mathfrak p'}} N'_{\mathfrak p'}
$$
$N'_{\mathfrak p'}$는 국소환 $R'_{\mathfrak p'}$ 위에서 사영이므로
자유이다(Algebra, Theorem
\ref{algebra-theorem-projective-free-over-local-ring}). 따라서
$z' \in M'_{\mathfrak p'} \otimes_{R'_{\mathfrak p'}} F'$가 되도록 하는
유한 자유 직합인자 $F' \subset N'_{\mathfrak p'}$를 찾을 수 있다.
$F'$의 계수가 $n$이면 다음을 얻는다.
$\mathfrak p' R'_{\mathfrak p'} \in
\text{WeakAss}_{R'_{\mathfrak p'}}({M'_{\mathfrak p'}}^{\oplus n})$.
이는
$\mathfrak p'R'_{\mathfrak p'} \in \text{WeakAss}(M'_{\mathfrak p'})$
를 함의한다. 예를 들어 Algebra, Lemma
\ref{algebra-lemma-weakly-ass}를 보라. 그러면
$\mathfrak p' \in \text{WeakAss}_{R'}(M')$이고, 이는 다시
$\mathfrak p \in \text{WeakAss}_R(M)$을 준다. 이로써 보조정리 동치의
함의 ``$\Rightarrow$''의 증명이 끝난다.

\medskip\noindent
$\mathfrak p \in \text{WeakAss}_R(M)$이고
$\overline{\mathfrak q} \in \text{Ass}_{\overline{S}}(\overline{N})$라고
가정하자. $\mathfrak q$가 $M \otimes_R N$의 약연관 소 아이디얼임을
보이고자 한다. $\overline{\mathfrak q}'$는
$\text{Ass}_{\overline{S}'}(\overline{N}')$의 극대 원소임에 유의하자.
이는 (\ref{flat-equation-key-observation})와 $\overline{\mathfrak q}$ 위에 놓인
$\overline{S}'$의 소 아이디얼들 사이에는 포함관계가 없다는 사실의
귀결이다(에탈 사상의 올은 이산이기 때문이다. Morphisms, Lemma
\ref{morphisms-lemma-etale-over-field}를 보라). 따라서
$R, S, \mathfrak p, \mathfrak q, M, N$을 각각
$R', S', \mathfrak p', \mathfrak q', M', N'$으로 대체한 뒤,
보조정리의 가정에 더하여 다음을 가정해도 된다.
\begin{enumerate}
\item $\mathfrak p \in \text{WeakAss}_R(M)$이다.
\item $\overline{\mathfrak q} \in \text{Ass}_{\overline{S}}(\overline{N})$이다.
\item $N$은 사영 $R$-가군이다.
\item $\overline{\mathfrak q}$는
$\text{Ass}_{\overline{S}}(\overline{N})$에서 극대이다.
\end{enumerate}
한 번 더 환원할 수 있다. 즉 $R, S, M, N$을 $\mathfrak p$에서의
국소화들로 대체해도 된다. 이로부터 조건 하나가 더 생긴다.
\begin{enumerate}
\item[(5)] $R$은 극대 아이디얼이 $\mathfrak p$인 국소환이다.
\end{enumerate}
(1) -- (5)가 $\mathfrak q \in \text{WeakAss}(M \otimes_R N)$을
함의함을 보여 증명을 끝내겠다.

\medskip\noindent
$R$은 국소이고 $\mathfrak p \in \text{WeakAss}_R(M)$이므로 소멸자
$I$의 근기가 $\mathfrak p$와 같은 $y \in M$을 택할 수 있다.
어떤 $\overline{g}_i \in \overline{S}$에 대해
$\overline{\mathfrak q} = (\overline{g}_1, \ldots, \overline{g}_n)$로
쓰자. $\overline{g}_i$로 가는 $g_i \in S$를 택하자. 그러면
$\mathfrak q = \mathfrak pS + g_1S + \ldots + g_nS$이다. 사상
$$
\Psi : N/IN \longrightarrow (N/IN)^{\oplus n}, \quad
z \longmapsto (g_1z, \ldots, g_nz).
$$
을 생각하자. 이는 사영 $R/I$-가군의 준동형이다. $\mathfrak p/I$는
국소 멱영이므로 국소환 $R/I$는 자기연관이다(More on Algebra,
Definition \ref{more-algebra-definition-auto-ass}).
$\overline{\mathfrak q} \in \text{Ass}_{\overline{S}}(\overline{N})$이므로
사상 $\Psi \otimes \kappa(\mathfrak p)$는 단사가 아니다. 따라서
More on Algebra, Lemma \ref{more-algebra-lemma-P-fPD-zero}에 의해
$\Psi$는 단사가 아니다. $\Psi$의 핵에서 영이 아닌
$z \in N/IN$을 택하자. 구성상 소멸자 $J = \text{Ann}_S(z)$는
$IS$와 $g_i$를 포함한다. 따라서 $\sqrt{J} \subset S$는
$\mathfrak q$를 포함한다. $\mathfrak s \subset S$를 $J$ 위의 극소 소
아이디얼이라 하자. 그러면 $\mathfrak q \subset \mathfrak s$이고,
$\mathfrak s$는 $\mathfrak p$ 위에 놓이며,
$\mathfrak s \in \text{WeakAss}_S(N/IN)$이다.
% P05-EN-CH38-0153: 동사가 빠진 원문 문장은 한국어로 자연스럽게 나타냈다.
마지막 사실은 약연관 소 아이디얼의 정의에서 따른다. 이미 증명한
보조정리의 ``$\Rightarrow$'' 부분을 환 준동형 $R \to S$와 가군
$R/I$ 및 $N$에 적용하면
$\overline{\mathfrak s} \in \text{Ass}_{\overline{S}}(\overline{N})$라는
결론을 얻는다. $\overline{\mathfrak q} \subset
\overline{\mathfrak s}$이므로 위 조건 (4)의 $\overline{\mathfrak q}$의
극대성에 의해 $\overline{\mathfrak q} = \overline{\mathfrak s}$이다.
이는 $\mathfrak q = \mathfrak s$임을 보인다.
% P05-EN-CH38-0154: conlude 철자 오류는 한국어로 자연스럽게 나타냈다.
따라서 원하는 결론을 얻는다.
\end{proof}

\begin{lemma}
\label{flat-lemma-bourbaki-finite-type-general-base-at-point}
$S$를 스킴이라 하자. $f : X \to S$를 국소 유한형이라 하자.
$x \in X$의 상을 $s \in S$라 하자. $\mathcal{F}$를 $X$ 위의
유한형 준연접층이라 하고, $\mathcal{G}$를 $S$ 위의 준연접층이라 하자.
$\mathcal{F}$가 $x$에서 $S$ 위로 평탄이면
$$
x \in \text{WeakAss}_X(\mathcal{F} \otimes_{\mathcal{O}_X} f^*\mathcal{G})
\Leftrightarrow
s \in \text{WeakAss}_S(\mathcal{G})
\text{ 이고 }
x \in \text{Ass}_{X_s}(\mathcal{F}_s).
$$
\end{lemma}

\begin{proof}
이 문단에서는 $f$가 유한표현인 경우로 환원한다. 문제는 $X$와 $S$
위에서 국소적이므로 $X$와 $S$가 아핀이라고 가정해도 된다.
$X = \Spec(B)$, $S = \Spec(A)$로 쓰고
$B = A[x_1, \ldots, x_n]/I$로 쓰자. 다시 말해 $S$ 위의 닫힌 몰입
$i : X \to \mathbf{A}^n_S$를 얻는다.
$t = i(x) \in \mathbf{A}^n_S$로 쓰자. $i_*\mathcal{F}$는
$\mathbf{A}^n_S$ 위의 유한형 준연접층이고 $t$에서 $S$ 위로
평탄하며, 다음이 성립함에 유의하자.
$$
i_*(\mathcal{F} \otimes_{\mathcal{O}_X} f^*\mathcal{G}) =
i_*\mathcal{F} \otimes_{\mathcal{O}_{\mathbf{A}^n_S}} p^*\mathcal{G}
$$
여기서 $p : \mathbf{A}^n_S \to S$는 사영이다.
$t$가 $i_*(\mathcal{F} \otimes_{\mathcal{O}_X} f^*\mathcal{G})$의
약연관점일 필요충분조건은 $x$가
$\mathcal{F} \otimes_{\mathcal{O}_X} f^*\mathcal{G}$의 약연관점인 것이다.
Divisors, Lemma \ref{divisors-lemma-weakly-associated-finite}를 보라.
마찬가지로 $x \in \text{Ass}_{X_s}(\mathcal{F}_s)$일 필요충분조건은
$t \in \text{Ass}_{\mathbf{A}^n_s}((i_*\mathcal{F})_s)$인 것이다
(Algebra, Lemma \ref{algebra-lemma-ass-quotient-ring}을 보라). 따라서
$X = \mathbf{A}^n_S$인 경우에 보조정리를 증명하면 충분하다. 그러므로
$X \to S$가 유한표현이라고 가정해도 된다.

\medskip\noindent
이 문단에서는 $\mathcal{F}$가 유한표현이고 $S$ 위에서 평탄인 경우로
환원한다. Proposition \ref{flat-proposition-finite-type-flat-at-point}에서와
같은 기본 에탈 근방 $e : (S', s') \to (S, s)$와 열린집합
$V \subset X \times_S \Spec(\mathcal{O}_{S', s'})$를 택하자.
$x$와 $s'$로 가는 유일한 점을 $x' \in X' = X \times_S S'$라 하자.
그러면 $X' \to S'$, $x'$, $s'$, $(X' \to X)^*\mathcal{F}$ 및
$e^*\mathcal{G}$에 대해 주장을 증명하면 충분하다. Lemma
\ref{flat-lemma-etale-weak-assassin-up-down}을 보라.
% P05-EN-CH38-0155: v를 도입하는 문장에 빠진 be는 한국어로 자연스럽게 나타냈다.
$x'$로 가는 유일한 점을 $v \in V$라 하고,
$s' \in \Spec(\mathcal{O}_{S', s'})$를 닫힌점이라 하자. 그러면
$\mathcal{O}_{V, v} = \mathcal{O}_{X', x'}$이고,
$\mathcal{O}_{\Spec(\mathcal{O}_{S', s'}), s'} =
\mathcal{O}_{S', s'}$이며, $\mathcal{F}$와 $\mathcal{G}$의 당김들의
줄기에도 같은 사실이 성립한다. 또한
$V_{s'} \subset X'_{s'}$는 열린 부분스킴이다.
% P05-EN-CH38-0156: 단수 주어 condition 뒤의 수 일치 오류는 한국어로 자연스럽게 나타냈다.
약연관점이라는 조건은 층의 줄기에만 의존하므로
$X' \to S'$, $x'$, $s'$, $(X' \to X)^*\mathcal{F}$ 및
$e^*\mathcal{G}$를
$V \to \Spec(\mathcal{O}_{S', s'})$, $v$, $s'$,
$(V \to X)^*\mathcal{F}$ 및
$(\Spec(\mathcal{O}_{S', s'}) \to S)^*\mathcal{G}$로 대체해도 된다.
따라서 $f$가 유한표현이고 $\mathcal{F}$가 유한표현이며 $S$ 위에서
평탄이라고 가정해도 된다.

\medskip\noindent
$f$가 유한표현이고 $\mathcal{F}$가 유한표현이며 $S$ 위에서
평탄이라고 가정하자. $X$와 $S$를 $x$와 $s$의 아핀 근방으로
축소한 뒤에는 Lemma \ref{flat-lemma-weak-bourbaki-pre}가 이 경우를 처리한다.
\end{proof}

\begin{lemma}
\label{flat-lemma-weak-bourbaki}
$R \to S$를 본질적으로 유한형인 환 준동형이라 하자. $N$을
$R$ 위에서 평탄인 유한 $S$-가군의 한 국소화라 하자. $M$을
$R$-가군이라 하자. 그러면
$$
\text{WeakAss}_S(M \otimes_R N)
=
\bigcup\nolimits_{\mathfrak p \in \text{WeakAss}_R(M)}
\text{Ass}_{S \otimes_R \kappa(\mathfrak p)}(N \otimes_R \kappa(\mathfrak p))
$$
\end{lemma}

\begin{proof}
이 보조정리는 Lemma
\ref{flat-lemma-bourbaki-finite-type-general-base-at-point}를 대수학으로
번역한 것이다. 번역의 세부사항은 생략한다.
\end{proof}

\begin{lemma}
\label{flat-lemma-bourbaki-finite-type-general-base}
$f : X \to S$를 국소 유한형인 사상이라 하자. $\mathcal{F}$를
$S$ 위에서 평탄인 $X$ 위의 유한형 준연접층이라 하고,
$\mathcal{G}$를 $S$ 위의 준연접층이라 하자. 그러면
$$
\text{WeakAss}_X(\mathcal{F} \otimes_{\mathcal{O}_X} f^*\mathcal{G}) =
\bigcup\nolimits_{s \in \text{WeakAss}_S(\mathcal{G})}
\text{Ass}_{X_s}(\mathcal{F}_s)
$$
\end{lemma}

\begin{proof}
Lemma \ref{flat-lemma-bourbaki-finite-type-general-base-at-point}의 직접적인
귀결이다.
\end{proof}

\begin{theorem}
\label{flat-theorem-finite-type-flat}
\begin{slogan}
평탄 궤적은 열린집합이다(비뇌터 버전).
\end{slogan}
$f : X \to S$를 스킴의 사상이라 하자. $\mathcal{F}$를 준연접
$\mathcal{O}_X$-가군이라 하자. 다음을 가정한다.
\begin{enumerate}
\item $X \to S$는 국소 유한표현이다.
\item $\mathcal{F}$는 유한형 $\mathcal{O}_X$-가군이다.
\item $S$의 약연관점들의 집합은 $S$에서 국소 유한이다.
\end{enumerate}
그러면 $U = \{x \in X \mid \mathcal{F}\text{가 }x\text{에서 }S\text{ 위로 평탄}\}$는
$X$에서 열린집합이고, $\mathcal{F}|_U$는 유한표현
$\mathcal{O}_U$-가군이며 $S$ 위에서 평탄하다.
\end{theorem}

\begin{proof}
$\mathcal{F}$가 $x$에서 $S$ 위로 평탄인 $x \in X$를 택하자.
$\mathcal{F}$의 제한이 $S$-평탄 유한표현 가군이 되는 $x$의 열린
근방을 찾아야 한다. 문제는 $X$와 $S$ 위에서 국소적이므로 $X$와
$S$가 아핀이라고 가정해도 된다. Lemma
\ref{flat-lemma-finite-type-flat-at-point-local}에 의해 $\mathcal{F}_x$는
유한표현 $\mathcal{O}_{X, x}$-가군이다. 따라서 Algebra, Lemma
\ref{algebra-lemma-construct-fp-module-from-stalk}에 의해 유한표현
$\mathcal{O}_X$-가군 $\mathcal{F}'$와 사상
$\varphi : \mathcal{F}' \to \mathcal{F}$가 존재하여 동형
$\varphi_x : \mathcal{F}'_x \to \mathcal{F}_x$를 유도한다. 특히
$\mathcal{F}'$는 $x$에서 $S$ 위로 평탄하다. 따라서 평탄성의 개방성,
More on Morphisms, Theorem
\ref{more-morphisms-theorem-openness-flatness}에 의해 $X$를 축소한 뒤
$\mathcal{F}'$가 $S$ 위에서 평탄이라고 가정해도 된다.
$\mathcal{F}$는 유한형이므로 $X$를 축소한 뒤 $\varphi$가 전사라고
가정해도 된다. Modules, Lemma
\ref{modules-lemma-finite-type-surjective-on-stalk}을 보라. 또는 Nakayama
보조정리(Algebra, Lemma \ref{algebra-lemma-NAK})를 사용해도 된다.
Lemma \ref{flat-lemma-bourbaki-finite-type-general-base}에 의해
$$
\text{WeakAss}_X(\mathcal{F}') \subset
\bigcup\nolimits_{s \in \text{WeakAss}(S)} \text{Ass}_{X_s}(\mathcal{F}'_s)
$$
이다. 가정에 의해 $\text{WeakAss}(S)$는 유한이고 Divisors, Lemma
\ref{divisors-lemma-finite-ass}에 의해
$\text{Ass}_{X_s}(\mathcal{F}'_s)$도 유한이므로
$\text{WeakAss}_X(\mathcal{F}')$가 유한이라는 결론을 얻는다.
Algebra, Lemma \ref{algebra-lemma-silly}를 사용하여 $X$를 한 번 더
축소한 뒤 $\text{WeakAss}_X(\mathcal{F}')$가 $x$의 일반화들에
포함된다고 가정해도 된다. 이제 $\mathcal{K} = \Ker(\varphi)$를
생각하자. 다음을 얻는다.
$\text{WeakAss}_X(\mathcal{K}) \subset \text{WeakAss}_X(\mathcal{F}')$
(Divisors, Lemma \ref{divisors-lemma-ses-weakly-ass}에 따른다.) 한편
% P05-EN-CH38-0157: 두 독립절의 잘못된 접속은 한국어로 자연스럽게 나타냈다.
$\varphi_x$는 동형이고, 모든 $x' \leadsto x$에 대해
$\varphi_{x'}$도 동형이다. 따라서
$\text{WeakAss}_X(\mathcal{K}) = \emptyset$이고, Divisors, Lemma
\ref{divisors-lemma-weakly-ass-zero}에 의해 $\mathcal{K} = 0$이다.
\end{proof}

\begin{lemma}
\label{flat-lemma-finite-type-flat-algebra}
$R \to S$를 유한표현 환 준동형이라 하자. $M$을 유한 $S$-가군이라
하자. $\text{WeakAss}_R(R)$이 유한이라고 가정하자. 그러면
$$
U = \{\mathfrak q \subset S \mid M_{\mathfrak q}\text{가 }R\text{ 위에서 평탄}\}
$$
는 $\Spec(S)$에서 열린집합이고, $D(g) \subset U$를 만족하는 모든
$g \in S$에 대해 국소화 $M_g$는 $R$ 위에서 평탄인 유한표현
$S_g$-가군이다.
\end{lemma}

\begin{proof}
Theorem \ref{flat-theorem-finite-type-flat}에서 곧바로 따른다.
\end{proof}

\begin{lemma}
\label{flat-lemma-finite-type-flat-X}
$f : X \to S$를 국소 유한형인 스킴의 사상이라 하자. $S$의
약연관점들의 집합이 $S$에서 국소 유한이라고 가정하자. 그러면 $f$가
평탄인 점 $x \in X$들의 집합은 열린 부분스킴 $U \subset X$이고,
$U \to S$는 평탄이고 국소 유한표현이다.
\end{lemma}

\begin{proof}
문제는 $X$와 $S$ 위에서 국소적이므로 $X$와 $S$가 아핀이라고
가정해도 된다. 그러면 $X \to S$는 유한형 환 준동형 $A \to B$에
대응한다. 전사 $A[x_1, \ldots, x_n] \to B$를 택하고 $B$를
$A[x_1, \ldots, x_n]$-가군으로 생각하자. Lemma
\ref{flat-lemma-finite-type-flat-algebra}를 적용하면 증명이 끝난다.
\end{proof}

\begin{lemma}
\label{flat-lemma-finite-type-flat-over-integral}
$f : X \to S$를 국소 유한형이고 평탄인 스킴의 사상이라 하자.
$S$가 정역 스킴이면 $f$는 국소 유한표현이다.
\end{lemma}

\begin{proof}
Lemma \ref{flat-lemma-finite-type-flat-X}의 특별한 경우이다.
\end{proof}

\begin{proposition}
\label{flat-proposition-flat-finite-type-finite-presentation-domain}
$R$을 정역이라 하자. $R \to S$를 유한형 환 준동형이라 하자.
$M$을 유한 $S$-가군이라 하자.
\begin{enumerate}
\item $S$가 $R$ 위에서 평탄이면 $S$는 유한표현 $R$-대수이다.
\item $M$이 $R$-가군으로서 평탄이면 $M$은 유한표현 $S$-가군이다.
\end{enumerate}
\end{proposition}

\begin{proof}
(1)은 Lemma \ref{flat-lemma-finite-type-flat-over-integral}의 특별한 경우이다.
(2)를 위해 전사 $R[x_1, \ldots, x_n] \to S$를 택하자. Lemma
\ref{flat-lemma-finite-type-flat-algebra}에 의해 $M$은 유한표현
$R[x_1, \ldots, x_n]$-가군이다. Algebra, Lemma
\ref{algebra-lemma-finitely-presented-over-subring}에 의해 결론을 얻는다.
\end{proof}

\begin{lemma}[Theorem \ref{flat-theorem-finite-type-flat}의 유한형 판]
\label{flat-lemma-finite-type-flat}
$f : X \to S$를 스킴의 사상이라 하자. $\mathcal{F}$를 준연접
$\mathcal{O}_X$-가군이라 하자. 다음을 가정한다.
\begin{enumerate}
\item $X \to S$는 국소 유한형이다.
\item $\mathcal{F}$는 유한형 $\mathcal{O}_X$-가군이다.
\item $S$의 약연관점들의 집합은 $S$에서 국소 유한이다.
\end{enumerate}
그러면 $U = \{x \in X \mid \mathcal{F}\text{가 }x\text{에서 }S\text{ 위로 평탄}\}$는
$X$에서 열린집합이고, $\mathcal{F}|_U$는 $S$ 위에서 평탄이며
$S$에 대해 상대적으로 국소 유한표현이다(More on Morphisms,
Definition
\ref{more-morphisms-definition-relatively-finitely-presented-sheaf}을 보라).
\end{lemma}

\begin{proof}
문제는 $X$와 $S$ 위에서 국소적이다. 따라서 $X$와 $S$가 아핀이라고
가정해도 된다. 그러면 닫힌 몰입 $i : X \to \mathbf{A}^n_S$를 택할 수
있다. $X' = \mathbf{A}^n_S \to S$와 유한형 준연접 가군
$\mathcal{F}' = i_*\mathcal{F}$에 Theorem
\ref{flat-theorem-finite-type-flat}을 적용하면
$$
U' = \{x' \in X' \mid \mathcal{F}'\text{가 }x'\text{에서 }S\text{ 위로 평탄}\}
$$
가 $X'$에서 열린집합이고 $\mathcal{F}'|_{U'}$가 유한표현임을 얻는다.
$\mathcal{F}'$는 $X' \setminus i(X)$에서 영으로 제한되고, 모든
$x \in X$에 대해 $\mathcal{F}'_{i(x)} \cong \mathcal{F}_x$이므로
$$
U' = i(U) \amalg (X' \setminus i(X))
$$
이다. 따라서 $U = i^{-1}(U')$는 열린집합이다. 더욱이
$\mathcal{F}'|_{U'} = (i|_U)_*(\mathcal{F}|_U)$임은 명백하다. 그러므로
More on Morphisms, Lemmas
\ref{more-morphisms-lemma-relative-finite-presentation}과
\ref{more-morphisms-lemma-finite-morphism-relative-finite-presentation}에 의해
$\mathcal{F}|_U$가 $S$에 대해 상대적으로 유한표현이라는 결론을 얻는다.
\end{proof}

\begin{lemma}
\label{flat-lemma-finite-type-flat-algebra-bis}
$R \to S$를 유한형 환 준동형이라 하자. $M$을 유한 $S$-가군이라
하자. $\text{WeakAss}_R(R)$이 유한이라고 가정하자. 그러면
$$
U = \{\mathfrak q \subset S \mid M_{\mathfrak q}\text{가 }R\text{ 위에서 평탄}\}
$$
는 $\Spec(S)$에서 열린집합이고, $D(g) \subset U$를 만족하는 모든
$g \in S$에 대해 국소화 $M_g$는 $R$ 위에서 평탄이고 $R$에 대해
상대적으로 유한표현인 $S_g$-가군이다(More on Algebra, Definition
\ref{more-algebra-definition-relatively-finitely-presented}를 보라).
\end{lemma}

\begin{proof}
이는 Lemma \ref{flat-lemma-finite-type-flat}을 대수학으로 번역한 것이다.
\end{proof}








\section{상대적으로 순수한 가군의 예}
\label{flat-section-examples-pure-modules}

\noindent
% P05-EN-CH38-0158: 현재 절을 가리키는 지시어가 빠진 원문을 한국어로 자연스럽게 나타냈다.
이 짧은 절에서는 나중에 도입할 {\it 상대적으로 순수한 가군}이라는 개념과
{\it 비순수성}이라는 개념의 동기가 되는 결과들의 예를 몇 가지 논한다.
각 예는 보조정리로 서술한다. 연관 소아이디얼에 관한 조건이 Lemmas
\ref{flat-lemma-base-change-universally-flat},
\ref{flat-lemma-universally-injective-to-completion},
\ref{flat-lemma-universally-injective-to-completion-flat},
\ref{flat-lemma-flat-pure-over-complete-projective}에 나타나는 조건들과 비슷함에
유의하자. 이에 관한 논의는 Algebra, Lemma
\ref{algebra-lemma-compare-relative-assassins}도 보라.

\begin{lemma}
\label{flat-lemma-explain-why-pure}
$R$을 극대 아이디얼이 $\mathfrak m$인 국소환이라 하자.
$R \to S$를 환 준동형이라 하고, $N$을 $S$-가군이라 하자. 다음을 가정하자.
\begin{enumerate}
\item $N$은 $R$-가군으로서 사영이고,
\item $S/\mathfrak mS$는 뇌터이며 $N/\mathfrak mN$은 유한
$S/\mathfrak mS$-가군이다.
\end{enumerate}
그러면 $\mathfrak p = R \cap \mathfrak q$라 할 때
$N \otimes_R \kappa(\mathfrak p)$의 연관 소아이디얼인 임의의
소아이디얼 $\mathfrak q \subset S$에 대하여
$\mathfrak q + \mathfrak m S \not = S$이다.
\end{lemma}

\begin{proof}
Lemmas \ref{flat-lemma-homothety-spectrum}과
\ref{flat-lemma-invert-universally-injective}의 가정들이 성립함에 유의하자.
이 보조정리들의 결론은 더 언급하지 않고 사용하겠다.
$\Sigma \subset S$를 $N/\mathfrak mN$ 위에서 영인자가 아닌 원소들의
곱셈적 집합이라 하자. 사상 $N \to \Sigma^{-1}N$은 $R$-보편 단사이다.
따라서 $N \otimes_R \kappa(\mathfrak p)$의 연관 소아이디얼인 임의의
$\mathfrak q \subset S$는
$\Sigma^{-1}N \otimes_R \kappa(\mathfrak p)$의 연관 소아이디얼이기도 하다.
이는 분명 $\mathfrak q$가 $\Sigma^{-1}S$의 한 소아이디얼에 대응함을
뜻한다. 따라서 $N/\mathfrak mN$의 한 연관 소아이디얼에 대응하는
$\mathfrak q'$에 대하여 $\mathfrak q \subset \mathfrak q'$이고,
결론을 얻는다.
\end{proof}

\noindent
다음 보조정리는 이 현상의 또 다른 (조금 어리석은) 예를 준다.

\begin{lemma}
\label{flat-lemma-explain-why-pure-complete}
$R$을 환이라 하고 $I \subset R$을 아이디얼이라 하자.
$R \to S$를 환 준동형이라 하고 $N$을 $S$-가군이라 하자.
$N$이 $I$-진 완비이면, 임의의 $R$-가군 $M$과
$N \otimes_R M$의 연관 소아이디얼인 임의의 소아이디얼
$\mathfrak q \subset S$에 대하여 $\mathfrak q + I S \not = S$이다.
\end{lemma}

\begin{proof}
$S^\wedge$로 $S$의 $I$-진 완비화를 나타내자.
$N$은 $S^\wedge$-가군이고, 따라서 $N \otimes_R M$도
$S^\wedge$-가군임에 유의하자.
$\mathfrak q = \text{Ann}_S(z)$를 만족하는 원소
$z \in N \otimes_R M$을 택하자. $z \not = 0$이므로
$\text{Ann}_{S^\wedge}(z) \not = S^\wedge$이고, 따라서
$\mathfrak q S^\wedge \not = S^\wedge$이다. 그러므로
$\mathfrak q S^\wedge \subset \mathfrak m$을 만족하는 극대 아이디얼
$\mathfrak m \subset S^\wedge$가 존재한다. Algebra, Lemma
\ref{algebra-lemma-radical-completion}에 의해
$IS^\wedge \subset \mathfrak m$이므로 결론을 얻는다.
\end{proof}

\noindent
국소환 위의 사영 가군은 자유 가군이므로 다음 보조정리는
Lemma \ref{flat-lemma-explain-why-pure}의 다른 증명을 준다는 점에 유의하자.
Algebra, Theorem \ref{algebra-theorem-projective-free-over-local-ring}을 보라.

\begin{lemma}
\label{flat-lemma-explain-why-pure-direct-sum-finite-modules}
$R$을 극대 아이디얼이 $\mathfrak m$인 국소환이라 하자.
$R \to S$를 환 준동형이라 하고 $N$을 $S$-가군이라 하자.
$N$이 $R$-가군으로서 유한 $R$-가군들의 직합과 동형이라고 가정하자.
그러면 임의의 $R$-가군 $M$과 $N \otimes_R M$의 연관 소아이디얼인
임의의 소아이디얼 $\mathfrak q \subset S$에 대하여
$\mathfrak q + \mathfrak m S \not = S$이다.
\end{lemma}

\begin{proof}
각 $M_i$가 유한 $R$-가군이 되도록
$N = \bigoplus_{i \in I} M_i$로 쓰자. $M$을 $R$-가군이라 하고,
$\mathfrak q + \mathfrak m S = S$를 만족하는 $N \otimes_R M$의
연관 소아이디얼 $\mathfrak q \subset S$를 택하자.
$\mathfrak q = \text{Ann}_S(z)$를 만족하는 원소
$z \in N \otimes_R M$을 택한다. 직합 분해를 조금 바꾸면 어떤
원소 $1 \in I$에 대하여 $z \in M_1 \otimes_R M$이라고 가정할 수 있다.
어떤 $f \in \mathfrak q$, $x_j \in \mathfrak m$, $g_j \in S$에 대하여
$1 = f + \sum x_j g_j$로 쓰자.
% P05-EN-CH38-0159: denote-by 구문에서 by가 빠진 원문을 한국어로 자연스럽게 나타냈다.
임의의 $g \in S$에 대하여 다음 $R$-선형 사상을 $g'$으로 나타내자.
$$
M_1 \to N \xrightarrow{g} N \to M_1
$$
여기서 첫째 화살표는 포함 사상이고, 둘째 화살표는 $g$에 의한 곱셈이며,
셋째 화살표는 사영 사상이다. 각 $x_j \in R$이므로 다음 등식을 얻는다.
$$
f' + \sum x_j g'_j = \text{id}_{M_1} \in \text{End}_R(M_1)
$$
나카야마 보조정리(Algebra, Lemma \ref{algebra-lemma-NAK})에 의해
$f'$은 전사이다. 따라서 Algebra, Lemma \ref{algebra-lemma-fun}에 의해
$f'$은 동형이다. 특히 사상
$$
M_1 \otimes_R M \to N \otimes_R M \xrightarrow{f} N \otimes_R M
\to M_1 \otimes_R M
$$
은 동형이다. 이는 $fz = 0$이라는 가정에 모순이다.
\end{proof}

\begin{lemma}
\label{flat-lemma-explain-why-pure-ML}
$R$을 극대 아이디얼이 $\mathfrak m$인 헨젤 국소환이라 하자.
$R \to S$를 환 준동형이라 하고 $N$을 $S$-가군이라 하자.
$N$이 $R$-가군으로서 가산 생성이고 Mittag-Leffler라고 가정하자.
그러면 임의의 $R$-가군 $M$과 $N \otimes_R M$의 연관 소아이디얼인
임의의 소아이디얼 $\mathfrak q \subset S$에 대하여
$\mathfrak q + \mathfrak m S \not = S$이다.
\end{lemma}

\begin{proof}
Algebra, Lemma \ref{algebra-lemma-split-ML-henselian}에 의해
이 보조정리는 Lemma
\ref{flat-lemma-explain-why-pure-direct-sum-finite-modules}로 귀착된다.
\end{proof}

\noindent
$f : X \to S$가 스킴의 사상이고 $\mathcal{F}$가 $X$ 위의
준연접 가군이라고 하자.
$\xi \in \text{Ass}_{X/S}(\mathcal{F})$라 하고
$Z = \overline{\{\xi\}}$라 하자. 다음 그림을 생각하자.
$$
\xymatrix{
\xi \ar@{|->}[d] & Z \ar[r] \ar[rd] & X \ar[d]^f \\
f(\xi) & & S
}
$$
$f(Z) \subset \overline{\{f(\xi)\}}$이고, 등식이 성립할 때, 즉
$f(Z) = \overline{\{f(\xi)\}}$일 때 그리고 그때에만 $f(Z)$가
닫혀 있음에 유의하자. Lemma \ref{flat-lemma-explain-why-pure}에 따르면,
$S$, $X$가 아핀이고, 올 $X_s$들이 뇌터이며, $\mathcal{F}$가 유한형이고,
$\Gamma(X, \mathcal{F})$가 사영 $\Gamma(S, \mathcal{O}_S)$-가군이면
$f(Z) = \overline{\{f(\xi)\}}$는 닫힌 부분집합이다.
% P05-EN-CH38-0160: 복수 주어와 단수 동사의 불일치를 한국어로 자연스럽게 나타냈다.
Lemmas \ref{flat-lemma-explain-why-pure-complete},
\ref{flat-lemma-explain-why-pure-direct-sum-finite-modules}, 및
\ref{flat-lemma-explain-why-pure-ML}에 서술된 경우들에도 조금씩 다른 유사한
명제들이 성립한다.




\section{비순수성}
\label{flat-section-impure}

\noindent
Section \ref{flat-section-examples-pure-modules}에서 예를 든 현상을
$\text{Ass}_{X/S}(\mathcal{F})$의 점들의 특수화로 형식화하고자 한다.
또한 점 $s \in S$ 주위에서 국소적으로 작업하고자 한다. 이를 위해
다음 정의들을 둔다.

\begin{situation}
\label{flat-situation-pre-pure}
여기서 $S$, $X$는 스킴이고 $f : X \to S$는 유한형 사상이다.
또한 $\mathcal{F}$는 유한형 준연접 $\mathcal{O}_X$-가군이고,
$s$는 $S$의 점이다.
\end{situation}

\noindent
이 상황에서 사상 $g : T \to S$, $g(t) = s$를 만족하는 점 $t \in T$,
특수화 $t' \leadsto t$, 그리고 $X$의 기저변환에서 $t'$ 위에 놓이는 점
$\xi \in X_T$를 생각하자. 그림으로는 다음과 같다.
\begin{equation}
\label{flat-equation-impurity}
\vcenter{
\xymatrix{
\xi \ar@{|->}[d] & \\
t' \ar@{~>}[r] & t \ar@{|->}[r] & s
}
}
\quad\quad
\vcenter{
\xymatrix{
X_T \ar[d] \ar[r] & X \ar[d] \\
T \ar[r]^g & S
}
}
\end{equation}
% P05-EN-CH38-0161: denote-by 구문에서 by가 빠진 원문을 한국어로 자연스럽게 나타냈다.
또한 $\mathcal{F}$의 $X_T$로의 당김을 $\mathcal{F}_T$로 나타내자.

\begin{definition}
\label{flat-definition-impurity}
Situation \ref{flat-situation-pre-pure}에서
$\xi \in \text{Ass}_{X_T/T}(\mathcal{F}_T)$이고
$\overline{\{\xi\}} \cap X_t = \emptyset$이면, 그림
(\ref{flat-equation-impurity})가 {\it $s$ 위에서 $\mathcal{F}$의 비순수성}을
정의한다고 한다. 이를 ``$(g : T \to S, t' \leadsto t, \xi)$를
$s$ 위에서 $\mathcal{F}$의 비순수성이라 하자''라고 나타내겠다.
\end{definition}

\begin{lemma}
\label{flat-lemma-impure-finite-presentation}
Situation \ref{flat-situation-pre-pure}의 상황이라 하자.
$s$ 위에서 $\mathcal{F}$의 비순수성이 존재하면, $g$가 국소 유한표현이고
% P05-EN-CH38-0162: 둘째 병렬절에 동사가 빠진 원문을 한국어로 자연스럽게 나타냈다.
$t$가 $s$ 위의 $g$의 올의 닫힌점인, $s$ 위에서 $\mathcal{F}$의 비순수성
$(g : T \to S, t' \leadsto t, \xi)$가 존재한다.
\end{lemma}

\begin{proof}
$(g : T \to S, t' \leadsto t, \xi)$를 $s$ 위에서 $\mathcal{F}$의
임의의 비순수성이라 하자. $t \in T$와 $Z = \overline{\{\xi\}}$에
Limits, Lemma \ref{limits-lemma-separate}를 적용하면 $t$의 열린 근방
$V \subset T$, 가환 그림
$$
\xymatrix{
V \ar[d] \ar[r]_a & T' \ar[d]^b \\
T \ar[r]^g & S,
}
$$
및 다음을 만족하는 닫힌 부분스킴 $Z' \subset X_{T'}$를 얻는다.
\begin{enumerate}
\item 사상 $b : T' \to S$는 국소 유한표현이고,
\item $Z' \cap X_{a(t)} = \emptyset$이며,
\item $Z \cap X_V$는 사상 $X_V \to X_{T'}$에 의해 $Z'$ 안으로 간다.
\end{enumerate}
$t'$가 $t$로 특수화하므로 $T$를 $t$의 열린 근방 $V$로 바꾸어도 된다.
따라서 $b \circ a = g$인 가환 그림
$$
\xymatrix{
X_T \ar[d] \ar[r] &
X_{T'} \ar[d] \ar[r] &
X \ar[d] \\
T \ar[r]^a & T' \ar[r]^b & S
}
$$
을 얻는다. $\xi' \in X_{T'}$로 $\xi$의 상을 나타내자.
Divisors, Lemma \ref{divisors-lemma-base-change-relative-assassin}에 의해
$\xi' \in \text{Ass}_{X_{T'}/T'}(\mathcal{F}_{T'})$이다.
% P05-EN-CH38-0163: 이미 폐포인 기호의 폐포를 다시 취하는 중복 표현을 제거했다.
더욱이 구성에 의해 $\overline{\{\xi'\}}$는 올 $X_{a(t)}$와 만나지 않는
닫힌 부분집합 $Z'$에 포함된다. 이로부터
$(T' \to S, a(t') \leadsto a(t), \xi')$가 $s$ 위에서
$\mathcal{F}$의 비순수성임을 안다.

\medskip\noindent
따라서 $g : T \to S$가 국소 유한표현이라고 가정해도 된다.
$Z = \overline{\{\xi\}}$라 하자. 가정에 의해 $Z_t = \emptyset$이다.
More on Morphisms, Lemma
\ref{more-morphisms-lemma-empty-generic-fibre}에 따르면 이는
$\overline{\{t\}}$의 어떤 열린 부분집합에 속하는 $t''$에 대하여
$Z_{t''} = \emptyset$이라는 뜻이다. $s$ 위의 $T \to S$의 올은
Jacobson 스킴이므로(Morphisms, Lemma
\ref{morphisms-lemma-ubiquity-Jacobson-schemes}을 보라),
% P05-EN-CH38-0164: 단수 주어와 복수 동사의 불일치를 한국어로 자연스럽게 나타냈다.
$Z_{t''} = \emptyset$을 만족하는 닫힌점
$t'' \in \overline{\{t\}}$가 존재한다. 그러면
$(g : T \to S, t' \leadsto t'', \xi)$가 원하는 비순수성이다.
\end{proof}

\begin{lemma}
\label{flat-lemma-impure-limit}
Situation \ref{flat-situation-pre-pure}의 상황이라 하자.
$(g : T \to S, t' \leadsto t, \xi)$를 $s$ 위에서 $\mathcal{F}$의
비순수성이라 하자. $T = \lim_{i \in I} T_i$가 $S$ 위의 아핀 스킴들의
유향 극한이라고 가정하자. 그러면 어떤 $i$에 대하여 삼중항
$(T_i \to S, t'_i \leadsto t_i, \xi_i)$는 $s$ 위에서
$\mathcal{F}$의 비순수성이다.
\end{lemma}

\begin{proof}
명제의 표기는 다음을 뜻한다. $p_i : T \to T_i$를 사영 사상이라 하고,
$t_i = p_i(t)$ 및 $t'_i = p_i(t')$라 하자. 끝으로
$\xi_i \in X_{T_i}$는 $\xi$의 상이다. Divisors, Lemma
\ref{divisors-lemma-base-change-relative-assassin}에 의해 $\xi_i$는
$T_i$ 위에서 $\mathcal{F}_{T_i}$의 상대 연관점 집합에 속한다.
따라서 남은 일은 어떤 $i$에 대하여
$\overline{\{\xi_i\}} \cap X_{t_i} = \emptyset$임을 보이는 것뿐이다.

\medskip\noindent
첫째 증명. $Z_i = \overline{\{\xi_i\}} \subset X_{T_i}$와
$Z = \overline{\{\xi\}} \subset X_T$에 축약 유도 스킴 구조를 주자.
Limits, Lemma \ref{limits-lemma-inverse-limit-irreducibles}에 의해
$Z = \lim Z_i$이다. 체 $k$와 상이 $t$인 사상
$\Spec(k) \to T$를 택하자. 그러면
$$
\emptyset =
Z \times_T \Spec(k) = (\lim Z_i) \times_{(\lim T_i)} \Spec(k)
= \lim Z_i \times_{T_i} \Spec(k)
$$
이다. 극한은 올곱과 가환하기 때문이다(극한은 극한과 가환한다).
$X_{T_i} \to T_i$가 유한형이고 따라서 $Z_i \to T_i$가 유한형이므로,
각 $Z_i \times_{T_i} \Spec(k)$는 준콤팩트이다. 그러므로 Limits,
Lemma \ref{limits-lemma-limit-nonempty}에 의해 어떤 $i$에 대하여
$Z_i \times_{T_i} \Spec(k)$가 공집합이다. 합성
$\Spec(k) \to T \to T_i$의 상은 $t_i$이므로 원하는 결론을 얻는다.

\medskip\noindent
둘째 증명. $Z = \overline{\{\xi\}}$라 하자. 이 상황에 Limits, Lemma
\ref{limits-lemma-separate}를 적용하면 $t$의 열린 근방 $V \subset T$,
가환 그림
$$
\xymatrix{
V \ar[d] \ar[r]_a & T' \ar[d]^b \\
T \ar[r]^g & S,
}
$$
및 다음을 만족하는 닫힌 부분스킴 $Z' \subset X_{T'}$를 얻는다.
\begin{enumerate}
\item 사상 $b : T' \to S$는 국소 유한표현이고,
\item $Z' \cap X_{a(t)} = \emptyset$이며,
\item $Z \cap X_V$는 사상 $X_V \to X_{T'}$에 의해 $Z'$ 안으로 간다.
\end{enumerate}
$V$가 $T$의 아핀 열린집합이라고 가정해도 된다. 따라서 Limits, Lemmas
\ref{limits-lemma-descend-opens}와 \ref{limits-lemma-limit-affine}에 의해
% P05-EN-CH38-0165: 앞에서 정의한 p_i 대신 미정의 f_i를 쓰는 원문 표기를 그대로 보존했다.
$V = f_i^{-1}(V_i)$인 어떤 $i$와 아핀 열린집합 $V_i \subset T_i$를
찾을 수 있다. Limits, Proposition
\ref{limits-proposition-characterize-locally-finite-presentation}에 의해,
필요하면 $i$를 조금 키운 뒤 $a = a_i \circ f_i|_V$를 만족하는 사상
$a_i : V_i \to T'$를 찾을 수 있다. 유도된 사상
$X_{V_i} \to X_{T'}$는 $\xi_i$를 $Z'$ 안으로 보낸다.
$Z' \cap X_{a(t)} = \emptyset$이므로
$(T_i \to S, t'_i \leadsto t_i, \xi_i)$가 $s$ 위에서
$\mathcal{F}$의 비순수성이라고 결론짓는다.
\end{proof}

\begin{lemma}
\label{flat-lemma-quasi-finite-impurity-elementary}
Situation \ref{flat-situation-pre-pure}의 상황이라 하자.
$g$가 $t$에서 준유한인, $s$ 위에서 $\mathcal{F}$의 비순수성
$(g : T \to S, t' \leadsto t, \xi)$가 존재하면,
$(T, t) \to (S, s)$가 기본 에탈 근방인 비순수성
$(g : T \to S, t' \leadsto t, \xi)$가 존재한다.
\end{lemma}

\begin{proof}
$(g : T \to S, t' \leadsto t, \xi)$를 $g$가 $t$에서 준유한인,
$s$ 위에서 $\mathcal{F}$의 비순수성이라 하자. $T$를 축소한 뒤
$g$가 국소 유한형이라고 가정해도 된다. $T \to S$와 $t \mapsto s$에
More on Morphisms, Lemma
\ref{more-morphisms-lemma-etale-makes-quasi-finite-finite-at-point}를 적용하면
다음 그림을 얻는다.
$$
\xymatrix{
T \ar[d] & T \times_S U \ar[l] \ar[d] & V \ar[l] \ar[ld] \\
S & U \ar[l]
}
$$
여기서 $(U, u) \to (S, s)$는 기본 에탈 근방이고,
$V \subset T \times_S U$는 $v = (t, u)$의 열린 근방이며,
$V \to U$는 유한이고 $v$는 $u$ 위에 놓이는 $V$의 유일한 점이다.
사상 $V \to T$는 에탈이고 따라서 평탄이므로, $v' \mapsto t'$를
만족하는 특수화 $v' \leadsto v$가 존재한다.
$\kappa(t') \subset \kappa(v')$는 유한 분리 확장임에 유의하자.
$\xi \in X_{t'}$로 가는 임의의 점 $\zeta \in X_{v'}$를 택하자.
Divisors, Lemma \ref{divisors-lemma-base-change-relative-assassin}에 의해
$\zeta \in \text{Ass}_{X_V/V}(\mathcal{F}_V)$이다. 더욱이 가정에 의해
폐포 $\overline{\{\xi\}}$가 $X_t$와 만나지 않으므로, 폐포
$\overline{\{\zeta\}}$도 올 $X_v$와 만나지 않는다. 다시 말해
% P05-EN-CH38-0166: 고정된 점 s 대신 대문자 S를 쓰는 원문을 그대로 보존했다.
$(V \to S, v' \leadsto v, \zeta)$는 $S$ 위에서
$\mathcal{F}$의 비순수성이다.

\medskip\noindent
다음으로 $u' \in U'$를 $v'$의 상이라 하고,
% P05-EN-CH38-0167: 정의되지 않은 U'를 쓰는 원문 표기를 그대로 보존했다.
$\theta \in X_U$를 $\zeta$의 상이라 하자. 그러면
$\theta \mapsto u'$이고 $u' \leadsto u$이다. Divisors, Lemma
\ref{divisors-lemma-base-change-relative-assassin}에 의해
% P05-EN-CH38-0168: 기저변환 첨자 없이 F를 쓰는 원문을 그대로 보존했다.
$\theta \in \text{Ass}_{X_U/U}(\mathcal{F})$이다. 더욱이
$\pi : X_V \to X_U$는 유한이므로
$\pi\big(\overline{\{\zeta\}}\big) = \overline{\{\pi(\zeta)\}}$이다.
$v$는 $u$ 위에 놓이는 $V$의 유일한 점이고
$X_v \cap \overline{\{\zeta\}} = \emptyset$이므로
$X_u \cap \overline{\{\pi(\zeta)\}} = \emptyset$이다.
따라서 $(U \to S, u' \leadsto u, \theta)$는 $s$ 위에서
$\mathcal{F}$의 비순수성이며, 결론을 얻는다.
\end{proof}

\begin{lemma}
\label{flat-lemma-Noetherian-impurity-quasi-finite}
Situation \ref{flat-situation-pre-pure}의 상황이라 하자.
$S$가 국소 뇌터라고 가정하자. $s$ 위에서 $\mathcal{F}$의 비순수성이
존재하면, $g$가 $t$에서 준유한인, $s$ 위에서 $\mathcal{F}$의 비순수성
$(g : T \to S, t' \leadsto t, \xi)$가 존재한다.
\end{lemma}

\begin{proof}
$S$를 $s$의 아핀 근방으로 바꾸어도 된다. Lemma
\ref{flat-lemma-impure-finite-presentation}에 의해 다음과 같이 가정해도 된다.
% P05-EN-CH38-0169: 목적어 없는 of와 such that이 이어진 손상된 절을 문맥대로 나타냈다.
비순수성 $(g : T \to S, t' \leadsto t, \xi)$가 있고, $g$는 국소 유한형이며
$t$는 $s$ 위의 $g$의 올의 닫힌점이다. $T$를
$\overline{\{t'\}}$ 위의 축약 유도 스킴 구조로 바꾸어도 된다.
$Z = \overline{\{\xi\}} \subset X_T$라 하자. 가정에 의해
$Z_t = \emptyset$이고 $Z \to T$의 상은 $t'$를 포함한다.
More on Morphisms, Lemma
\ref{more-morphisms-lemma-relative-assassin-in-neighbourhood}에 의해,
임의의 $w \in f(V)$에 대하여 $V_w$의 임의의 일반점 $\xi'$가
$\text{Ass}_{X_T/T}(\mathcal{F}_T)$에 속하도록 하는 공집합이 아닌
열린집합 $V \subset Z$가 존재한다. More on Morphisms, Lemma
\ref{more-morphisms-lemma-nonempty-generic-fibre}에 의해
$W \subset f(V)$인 공집합이 아닌 열린집합 $W \subset T$가 존재한다.
More on Morphisms, Lemma
\ref{more-morphisms-lemma-quasi-finite-quasi-section-meeting-nearby-open-X}에
의해 $t \in T'$이고 $T' \to S$가 $t$에서 준유한이며, $s$로 가지 않는
점 $z \in T' \cap W$와 특수화 $z \leadsto t$가 존재하도록 하는
닫힌 부분스킴 $T' \subset T$가 존재한다. 공집합이 아닌 스킴 $V_z$의
임의의 일반점 $\xi'$를 택하자. 그러면
$(T' \to S, z \leadsto t, \xi')$가 원하는 비순수성이다.
\end{proof}

\noindent
이하에서는 $s$에서 $S$의 헨젤화
$S^h = \Spec(\mathcal{O}_{S, s}^h)$를 사용한다. \'Etale Cohomology,
Definition \ref{etale-cohomology-definition-etale-local-rings}을 보라.
% P05-EN-CH38-0170: 헨젤화의 닫힌점을 보내는 절의 관사와 어순 오류를 자연스럽게 나타냈다.
$S^h \to S$는 $S^h$의 닫힌점을 $s$로 보내고 잉여체들의 동형을
유도하므로,
% P05-EN-CH38-0171: 표기 지정의 어순이 손상된 원문을 자연스럽게 나타냈다.
이 닫힌점도 $s \in S^h$로 나타내겠다. 따라서
$(S^h, s) \to (S, s)$는 점이 지정된 스킴들의 사상이다.

\begin{lemma}
\label{flat-lemma-impurity-on-henselization}
Situation \ref{flat-situation-pre-pure}의 상황이라 하자.
$s$ 위에서 $\mathcal{F}$의 비순수성
$(S^h \to S, s' \leadsto s, \xi)$가 존재하면,
$(T, t) \to (S, s)$가 기본 에탈 근방인, $s$ 위에서
$\mathcal{F}$의 비순수성 $(T \to S, t' \leadsto t, \xi)$가 존재한다.
\end{lemma}

\begin{proof}
$S$를 $s$의 아핀 근방으로 바꾸어도 된다. $S = \Spec(A)$라 하고,
$s$가 소아이디얼 $\mathfrak p \subset A$에 대응한다고 하자. 그러면
$\mathcal{O}_{S, s}^h = \colim_{(T, t)} \Gamma(T, \mathcal{O}_T)$
이다. 여기서 여극한은 $(S, s)$의 아핀 기본 에탈 근방 $(T, t)$들의
여과 범주의 반대 범주 위에서 취한다. More on Morphisms, Lemma
\ref{more-morphisms-lemma-describe-henselization}과 그 증명을 보라.
따라서 $S^h = \lim_i T_i$이고, Lemma \ref{flat-lemma-impure-limit}에 의해
결론을 얻는다.
\end{proof}

\begin{lemma}
\label{flat-lemma-pure-along-X-s}
Situation \ref{flat-situation-pre-pure}에서 다음은 서로 동치이다.
\begin{enumerate}
\item $S^h$가 $s$에서 $S$의 헨젤화일 때, $s$ 위에서 $\mathcal{F}$의
비순수성 $(S^h \to S, s' \leadsto s, \xi)$가 존재한다.
\item $(T, t) \to (S, s)$가 기본 에탈 근방인, $s$ 위에서
$\mathcal{F}$의 비순수성 $(T \to S, t' \leadsto t, \xi)$가 존재한다.
\item $T \to S$가 $t$에서 준유한인, $s$ 위에서 $\mathcal{F}$의
비순수성 $(T \to S, t' \leadsto t, \xi)$가 존재한다.
\end{enumerate}
\end{lemma}

\begin{proof}
에탈 사상은 국소 준유한이므로 (2)가 (3)을 함의함은 분명하다.
Lemma \ref{flat-lemma-quasi-finite-impurity-elementary}에서 (3)이 (2)를
함의함을 보았고, Lemma \ref{flat-lemma-impurity-on-henselization}에서
(1)이 (2)를 함의함을 보았다. 끝으로
$(T \to S, t' \leadsto t, \xi)$가 $s$ 위에서 $\mathcal{F}$의
비순수성이고 $(T, t) \to (S, s)$가 기본 에탈 근방이면, 구조 사상
$S^h \to S$의 분해 $S^h \to T \to S$를 택할 수 있다.
$t'$로 가는 임의의 점 $s' \in S^h$와
$\xi \in X_{t'}$로 가는 임의의 $\xi' \in X_{s'}$를 택하자. 그러면
$(S^h \to S, s' \leadsto s, \xi')$는 $s$ 위에서
$\mathcal{F}$의 비순수성이다.
% P05-EN-CH38-0172: 특수화·상대 연관점·빈 교집합 확인을 생략한 원문을 그대로 보존했다.
자세한 내용은 생략한다.
\end{proof}





\section{상대적으로 순수한 가군}
\label{flat-section-pure}

\noindent
기저에 대해 순수한 가군이라는 개념은 \cite{GruRay}에서 도입되었다.

\begin{definition}
\label{flat-definition-pure}
$f : X \to S$를 유한형인 스킴의 사상이라 하고, $\mathcal{F}$를
유한형 준연접 $\mathcal{O}_X$-가군이라 하자.
\begin{enumerate}
\item $s \in S$라 하자. $(T, t) \to (S, s)$가 기본 에탈 근방인,
$s$ 위에서 $\mathcal{F}$의 비순수성
$(g : T \to S, t' \leadsto t, \xi)$가 없으면 $\mathcal{F}$가
{\it $X_s$를 따라 순수하다}고 한다.
\item $s$ 위에서 $\mathcal{F}$의 비순수성이 전혀 존재하지 않으면
$\mathcal{F}$가 {\it $X_s$를 따라 보편적으로 순수하다}고 한다.
\item $\mathcal{O}_X$가 $X_s$를 따라 순수하면 $X$가
{\it $X_s$를 따라 순수하다}고 한다.
\item 모든 $s \in S$에 대하여 $\mathcal{F}$가 $X_s$를 따라 보편적으로
순수하면, $\mathcal{F}$가 {\it 보편적으로 $S$-순수하다}, 또는
{\it $S$에 대해 보편적으로 순수하다}고 한다.
\item 모든 $s \in S$에 대하여 $\mathcal{F}$가 $X_s$를 따라 순수하면,
$\mathcal{F}$가 {\it $S$-순수하다}, 또는 {\it $S$에 대해 순수하다}고 한다.
\item $\mathcal{O}_X$가 $S$에 대해 순수하면 $X$가
{\it $S$-순수하다}, 또는 {\it $S$에 대해 순수하다}고 한다.
\end{enumerate}
\end{definition}

\noindent
여기서는 의도적으로 국소 유한형 사상만이 아니라 유한형 사상으로
범위를 제한한다. 이에 관한 논의는 Remark
\ref{flat-remark-discuss-finite-type}를 보라. 이 정의의 상황에서 Lemma
\ref{flat-lemma-pure-along-X-s}는 다음이 서로 동치임을 말해 준다.
\begin{enumerate}
\item $\mathcal{F}$는 $X_s$를 따라 순수하다.
\item $g$가 $t$에서 준유한인 비순수성
$(g : T \to S, t' \leadsto t, \xi)$가 없다.
\item $S^h$가 $s$에서 $S$의 헨젤화일 때,
$(S^h \to S, s' \leadsto s, \xi)$ 꼴의 비순수성이 전혀 존재하지 않는다.
\end{enumerate}
$X^h = X \times_S S^h$라 쓰고 $\mathcal{F}$의 $X^h$로의 당김을
$\mathcal{F}^h$로 나타내면, 마지막 조건을 다음과 같이 더 긍정적인
방식으로 서술할 수 있다.
\begin{enumerate}
\item[(4)] $\text{Ass}_{X^h/S^h}(\mathcal{F}^h)$의 모든 점은
$X_s$의 점으로 특수화한다.
\end{enumerate}
특히 $\mathcal{F}$가 $X_s$를 따라 순수할 필요충분조건은
$\mathcal{F}$의 $X \times_S \Spec(\mathcal{O}_{S, s})$로의 당김이
$X_s$를 따라 순수한 것임이 분명하다.

\begin{remark}
\label{flat-remark-discuss-finite-type}
$f : X \to S$를 국소 유한형 사상이라 하고 $\mathcal{F}$를 유한형
준연접 $\mathcal{O}_X$-가군이라 하자. 이 경우에도 위의 (1)과 (2)는
서로 동치이다. Lemma \ref{flat-lemma-quasi-finite-impurity-elementary}의
증명은 $f$가 준콤팩트라는 사실을 사용하지 않기 때문이다.
(3)과 (4)가 서로 동치임도 분명하다. 그러나 (1)과 (3)이 서로
동치인지는 알지 못한다. 이 경우에는 \cite{GruRay}에서 하듯 동치인
조건 (3)과 (4)를 사용하여 순수성을 정의하는 편이 때로 더 편리할 수 있다.
한편 많은 응용에서는 올바른 개념이 실제로 보편적 순수성인 듯하다.
\end{remark}

\noindent
기저에 대해 순수하다는 성질이 기저변환 아래 보존되는지, 즉 순수함과
보편적으로 순수함이 같은 것인지는 자연스러운 질문이다. 뇌터 기저
스킴 위에서는 이것이 참이고(Lemma \ref{flat-lemma-Noetherian-base-change}를
보라), 층이 평탄해도 참이다(Lemmas
\ref{flat-lemma-finite-type-flat-pure-along-fibre-is-universal}와
\ref{flat-lemma-finite-type-flat-pure-is-universal}을 보라).
사상과 층이 유한표현이어도 일반적으로는 참이 아니다.
% P05-EN-CH38-0177: counter example의 비표준 띄어쓰기를 자연스럽게 나타냈다.
반례는 Examples, Section \ref{examples-section-pure-not-universally}을 보라.
먼저 여기서 쓰는 ``보편적으로''라는 말을 통상적인 개념과 맞춘다.

\begin{lemma}
\label{flat-lemma-base-change-universally}
$f : X \to S$를 유한형인 스킴의 사상이라 하고, $\mathcal{F}$를
유한형 준연접 $\mathcal{O}_X$-가군이라 하자. $s \in S$라 하자.
다음은 서로 동치이다.
\begin{enumerate}
\item $\mathcal{F}$는 $X_s$를 따라 보편적으로 순수하다.
\item 점이 지정된 스킴들의 모든 사상 $(S', s') \to (S, s)$에 대하여
당김 $\mathcal{F}_{S'}$는 $X_{s'}$를 따라 순수하다.
\end{enumerate}
특히 $\mathcal{F}$가 $S$에 대해 보편적으로 순수할 필요충분조건은
$\mathcal{F}$의 모든 기저변환 $\mathcal{F}_{S'}$가 $S'$에 대해
순수한 것이다.
\end{lemma}

\begin{proof}
이는 형식적이다.
\end{proof}

\begin{lemma}
\label{flat-lemma-quasi-finite-base-change}
$f : X \to S$를 유한형인 스킴의 사상이라 하고, $\mathcal{F}$를
유한형 준연접 $\mathcal{O}_X$-가군이라 하자. $s \in S$라 하고,
$(S', s') \to (S, s)$를 점이 지정된 스킴들의 사상이라 하자.
$S' \to S$가 $s'$에서 준유한이고 $\mathcal{F}$가 $X_s$를 따라
순수하면, $\mathcal{F}_{S'}$는 $X_{s'}$를 따라 순수하다.
\end{lemma}

\begin{proof}
% P05-EN-CH38-0173: 조건 접속사 If의 오자 It을 한국어로 자연스럽게 나타냈다.
$(T \to S', t' \leadsto t, \xi)$가 $T \to S'$가 $t$에서 준유한인,
$s'$ 위에서 $\mathcal{F}_{S'}$의 비순수성이면,
% P05-EN-CH38-0174: 특수화 화살표 대신 보통 화살표를 쓰는 원문을 그대로 보존했다.
$(T \to S, t' \to t, \xi)$는 $T \to S$가 $t$에서 준유한인,
$s$ 위에서 $\mathcal{F}$의 비순수성이다. Morphisms, Lemma
\ref{morphisms-lemma-composition-quasi-finite}을 보라. 따라서 이 보조정리는
% P05-EN-CH38-0175: 정의 뒤 논의를 가리키는 전치사와 관사가 빠진 원문을 자연스럽게 나타냈다.
Definition \ref{flat-definition-pure} 뒤의 논의에서 주어진 순수성의 특성화
(2)로부터 즉시 따른다.
\end{proof}

\begin{lemma}
\label{flat-lemma-Noetherian-base-change}
$f : X \to S$를 유한형인 스킴의 사상이라 하고, $\mathcal{F}$를
유한형 준연접 $\mathcal{O}_X$-가군이라 하자. $s \in S$라 하자.
$\mathcal{O}_{S, s}$가 뇌터이면, $\mathcal{F}$가 $X_s$를 따라
순수할 필요충분조건은 $\mathcal{F}$가 $X_s$를 따라 보편적으로
순수한 것이다.
\end{lemma}

\begin{proof}
먼저 $S$를 $\Spec(\mathcal{O}_{S, s})$로 바꾸어도 된다. 즉 $S$가
뇌터라고 가정해도 된다. 다음으로 Lemma
\ref{flat-lemma-Noetherian-impurity-quasi-finite}와
% P05-EN-CH38-0176: 특정 특성화 앞의 관사가 빠진 원문을 자연스럽게 나타냈다.
Definition \ref{flat-definition-pure} 뒤의 논의에서 주어진 순수성의 특성화
(2)를 사용하면 결론을 얻는다.
\end{proof}

\noindent
순수성은 평탄 하강을 만족한다.

\begin{lemma}
\label{flat-lemma-flat-descend-pure}
$f : X \to S$를 유한형인 스킴의 사상이라 하고, $\mathcal{F}$를
유한형 준연접 $\mathcal{O}_X$-가군이라 하자. $s \in S$라 하고,
$(S', s') \to (S, s)$를 점이 지정된 스킴들의 사상이라 하자.
$S' \to S$가 $s'$에서 평탄하다고 가정하자.
\begin{enumerate}
\item $\mathcal{F}_{S'}$가 $X_{s'}$를 따라 순수하면,
$\mathcal{F}$는 $X_s$를 따라 순수하다.
\item $\mathcal{F}_{S'}$가 $X_{s'}$를 따라 보편적으로 순수하면,
$\mathcal{F}$는 $X_s$를 따라 보편적으로 순수하다.
\end{enumerate}
\end{lemma}

\begin{proof}
$(T \to S, t' \leadsto t, \xi)$를 $s$ 위에서 $\mathcal{F}$의
비순수성이라 하자. $T_1 = T \times_S S'$로 두고,
% P05-EN-CH38-0178: 일반적으로 거짓인 유일점 주장을 원문대로 보존했다.
$t_1$을 $t$와 $s'$로 가는 $T_1$의 유일한 점이라 하자.
$T_1 \to T$는 $t_1$에서 평탄하므로(Morphisms, Lemma
\ref{morphisms-lemma-base-change-flat}을 보라), $t' \leadsto t$ 위에 놓이는
특수화 $t'_1 \leadsto t_1$가 존재한다. Algebra, Section
\ref{algebra-section-going-up}을 보라.
$\Spec(\kappa(t'_1) \otimes_{\kappa(t')} \kappa(\xi))$의 한 일반점에
대응하는 점 $\xi_1 \in X_{t'_1}$를 택하자. Schemes, Lemma
\ref{schemes-lemma-points-fibre-product}을 보라. Divisors, Lemma
\ref{divisors-lemma-base-change-relative-assassin}에 의해
$\xi_1 \in \text{Ass}_{X_{T_1}/T_1}(\mathcal{F}_{T_1})$이다.
$X_{T_1}$에서 $\{\xi_1\}$의 자리스키 폐포는 $X_T$에서 $\{\xi\}$의
자리스키 폐포 안으로 가므로, 이 폐포는 $X_{t_1}$와 서로소이다.
따라서 $(T_1 \to S', t'_1 \leadsto t_1, \xi_1)$는 $s'$ 위에서
$\mathcal{F}_{S'}$의 비순수성이다. 다시 말해 보조정리의 (2)의
% P05-EN-CH38-0179: contrapositive 뒤의 비표준 전치사를 자연스럽게 나타냈다.
대우를 증명했다. 끝으로 $(T, t) \to (S, s)$가 기본 에탈 근방이면
$(T_1, t_1) \to (S', s')$도 기본 에탈 근방이고, 이로부터 (1)을 얻는다.
\end{proof}

\begin{lemma}
\label{flat-lemma-supported-on-closed}
$i : Z \to X$를 스킴 $S$ 위의 유한형 스킴들의 닫힌 몰입이라 하자.
$s \in S$라 하고, $\mathcal{F}$를 $Z$ 위의 유한형 준연접 층이라 하자.
그러면 $\mathcal{F}$가 $Z_s$를 따라 (보편적으로) 순수할 필요충분조건은
$i_*\mathcal{F}$가 $X_s$를 따라 (보편적으로) 순수한 것이다.
\end{lemma}

\begin{proof}
이는 Divisors, Lemma
\ref{divisors-lemma-relative-weak-assassin-finite}에서 따른다.
\end{proof}








\section{상대적으로 순수한 층의 예}
\label{flat-section-examples-pure-sheaves}

\noindent
순수성이 무엇을 뜻하는지 알아볼 수 있는 몇 가지 예를 든다.

\begin{lemma}
\label{flat-lemma-proper-pure}
$f : X \to S$를 유한형인 스킴의 사상이라 하고, $\mathcal{F}$를
유한형 준연접 $\mathcal{O}_X$-가군이라 하자.
\begin{enumerate}
\item $\mathcal{F}$의 지지가 $S$ 위에서 고유이면,
$\mathcal{F}$는 $S$에 대해 보편적으로 순수하다.
\item $f$가 고유이면 $\mathcal{F}$는 $S$에 대해 보편적으로 순수하다.
\item $f$가 고유이면 $X$는 $S$에 대해 보편적으로 순수하다.
\end{enumerate}
\end{lemma}

\begin{proof}
먼저 (1)을 (2)로 귀착시킨다. $Z \subset X$를 $\mathcal{F}$의
스킴론적 지지라 하고, $i : Z \to X$를 대응하는 닫힌 몰입이라 하자.
어떤 유한형 준연접 $\mathcal{O}_Z$-가군 $\mathcal{G}$에 대하여
$\mathcal{F} = i_*\mathcal{G}$로 쓸 수 있다. Morphisms, Section
\ref{morphisms-section-support}을 보라. (1)의 경우에는 가정에 의해
$Z \to S$가 고유이다. 따라서 Lemma \ref{flat-lemma-supported-on-closed}에
의해 (1)은 (2)로 귀착된다.

\medskip\noindent
$f$가 고유라고 가정하자. $(g : T \to S, t' \leadsto t, \xi)$를
$s \in S$ 위에서 $\mathcal{F}$의 비순수성이라 하자. $f$는
고유이므로 보편적으로 닫힌 사상이다. 따라서 $f_T : X_T \to T$는
닫힌 사상이다. $f_T(\xi) = t'$이므로
% P05-EN-CH38-0180: 기저변환 사상 f_T 대신 f를 쓰는 원문을 그대로 보존했다.
$t \in f(\overline{\{\xi\}})$인데, 이는 모순이다.
\end{proof}

\begin{lemma}
\label{flat-lemma-quasi-finite-pure}
$f : X \to S$를 분리 유한형인 스킴의 사상이라 하고, $\mathcal{F}$를
유한형 준연접 $\mathcal{O}_X$-가군이라 하자. 모든 $s \in S$에 대하여
$\text{Supp}(\mathcal{F}_s)$가 유한이라고 가정하자. 다음은 서로 동치이다.
\begin{enumerate}
\item $\mathcal{F}$는 $S$에 대해 순수하다.
\item $\mathcal{F}$의 스킴론적 지지는 $S$ 위에서 유한이다.
\item $\mathcal{F}$는 $S$에 대해 보편적으로 순수하다.
\end{enumerate}
특히 준유한 분리 사상 $X \to S$가 주어졌을 때, $X$가 $S$에 대해
순수할 필요충분조건은 $X \to S$가 유한인 것이다.
\end{lemma}

\begin{proof}
$Z \subset X$를 $\mathcal{F}$의 스킴론적 지지라 하자. Morphisms,
Definition \ref{morphisms-definition-scheme-theoretic-support}을 보라.
그러면 $Z \to S$는 올들이 유한인 분리 유한형 사상이다. 따라서 이는
분리 준유한 사상이다. Morphisms, Lemma
\ref{morphisms-lemma-quasi-finite}을 보라. Lemma
\ref{flat-lemma-supported-on-closed}에 의해, $Z \to S$와 $Z$ 위의 유한형
준연접 가군으로 보는 층 $\mathcal{F}$에 대하여 보조정리를 증명하면
충분하다. 따라서 $X \to S$가 분리 준유한이고
$\text{Supp}(\mathcal{F}) = X$라고 가정해도 된다.

\medskip\noindent
Lemma \ref{flat-lemma-proper-pure}와 Morphisms, Lemma
\ref{morphisms-lemma-finite-proper}에 의해 (2)는 (3)을 함의한다.
(3)이 (1)을 함의함은 자명하다. (1)이 성립한다고 가정하고 (2)를
증명하자. $S$가 아핀이라고 가정해도 됨은 분명하다. More on Morphisms,
Lemma \ref{more-morphisms-lemma-quasi-finite-separated-pass-through-finite}에
의해 다음 그림을 찾을 수 있다.
$$
\xymatrix{
X \ar[rd]_f \ar[rr]_j & & T \ar[ld]^\pi \\
& S &
}
$$
여기서 $\pi$는 유한이고 $j$는 준콤팩트 열린 몰입이다. $j$가 닫혀
있음을 보이면 $j$는 닫힌 몰입이고, 따라서 $f = \pi \circ j$가
유한이라고 결론지을 수 있다. $j$가 닫혀 있음을 보이려면 특수화를
$j$를 따라 올릴 수 있음을 보이면 충분하다. Schemes, Lemma
\ref{schemes-lemma-quasi-compact-closed}를 보라. $x \in X$라 하고,
$t' = j(x)$로 두며, $t' \leadsto t$를 특수화라 하자.
$t \in j(X)$임을 보여야 한다. $s' = f(x)$와 $s = \pi(t)$로 두면
$s' \leadsto s$이다. More on Morphisms, Lemma
\ref{more-morphisms-lemma-etale-splits-off-quasi-finite-part-technical}에 의해
기본 에탈 근방 $(U, u) \to (S, s)$와 분해
$$
T_U = T \times_S U = V \amalg W
$$
를 찾을 수 있다. 여기서 두 부분스킴은 열리고 닫혀 있고,
$V \to U$는 유한이며, $u$로 가는 $V$의 유일한 점 $v$가 존재하고,
$v$는 $T$에서 $t$로 간다. $V \to T$가 에탈이므로 일반화를 올릴 수
있다. Morphisms, Lemmas
\ref{morphisms-lemma-generalizations-lift-flat}과
\ref{morphisms-lemma-etale-flat}을 보라. 따라서 특수화
$v' \leadsto v$가 존재하며, $v'$의 상은 $t' \in T$이다.
특히 $v' \in X_U \subset T_U$이다.
% P05-EN-CH38-0181: v'가 아니라 t'의 상으로 u'를 정의하는 원문을 그대로 보존했다.
$u' \in U$를 $t'$의 상으로 나타내자. $X_{u'}$는 유한 이산 집합이고
$X_{u'} = \text{Supp}(\mathcal{F}_{u'})$이므로
$v' \in \text{Ass}_{X_U/U}(\mathcal{F})$임에 유의하자.
$\mathcal{F}$가 $S$에 대해 순수하므로 $v'$는 $X_u$의 한 점으로
특수화해야 한다. $v$는 $u$ 위에 놓이는 $V$의 유일한 점이고
($W$의 어떤 점도 $v'$의 특수화일 수 없으므로), $v \in X_u$이다.
따라서 $t \in X$이다.
\end{proof}

\begin{lemma}
\label{flat-lemma-flat-geometrically-integral-fibres-pure}
$f : X \to S$를 올들이 기하적으로 정역인 유한형 평탄 사상이라 하자.
그러면 $X$는 $S$ 위에서 보편적으로 순수하다.
\end{lemma}

\begin{proof}
$\xi \in X$이고 $s' = f(\xi)$라 하며, $s' \leadsto s$를 $S$의
특수화라 하자. $\xi$가 $X_{s'}$의 연관점이면 $X_{s'}$가 정역 스킴이므로
$\xi$는 유일한 일반점이다. $\xi_0$를 $X_s$의 유일한 일반점이라 하자.
$X \to S$가 평탄하므로 $s' \leadsto s$를 $X$의 특수화
$\xi' \leadsto \xi_0$로 올릴 수 있다. Morphisms, Lemma
\ref{morphisms-lemma-generalizations-lift-flat}을 보라.
% P05-EN-CH38-0182: 접속사 Then의 오자 The를 한국어로 자연스럽게 나타냈다.
$\xi$는 $X_{s'}$의 일반점이므로 $\xi \leadsto \xi'$이고, 따라서
$\xi \leadsto \xi_0$이다. 이는
% P05-EN-CH38-0183: 특수화 화살표 대신 보통 화살표를 쓰는 원문을 그대로 보존했다.
$(\text{id}_S, s' \to s, \xi)$가 $s$ 위에서 $\mathcal{O}_X$의
비순수성이 아니라는 뜻이다. $f$가 올들이 기하적으로 정역인 유한형
평탄 사상이라는 가정은 기저변환 아래 보존되므로, 어떤 기저변환 뒤에도
비순수성이 존재하지 않는다. 따라서 $X$는 보편적으로 $S$-순수하다.
\end{proof}

\begin{lemma}
\label{flat-lemma-affine-locally-projective-pure}
$f : X \to S$를 유한형 아핀 사상이라 하자. $\mathcal{F}$를
$f_*\mathcal{F}$가 $S$ 위에서 국소 사영인 유한형 준연접
$\mathcal{O}_X$-가군이라 하자. Properties, Definition
\ref{properties-definition-locally-projective}를 보라. 그러면
$\mathcal{F}$는 $S$ 위에서 보편적으로 순수하다.
\end{lemma}

\begin{proof}
$S$가 헨젤 국소환의 스펙트럼인 경우로 귀착시키면 이는 Lemma
\ref{flat-lemma-explain-why-pure}에서 따른다.
\end{proof}




\section{순수성의 판정법}
\label{flat-section-criterion-purity}

\noindent
먼저 유한형 준연접 층들의 평탄한 족이 주어졌을 때, 상대 연관점 집합의
점들이 (특수화한다면) 가까운 올들의 상대 연관점 집합의 점들로
특수화함을 증명한다.

\begin{lemma}
\label{flat-lemma-associated-point-specializes}
$f : X \to S$를 유한형인 스킴의 사상이라 하고, $\mathcal{F}$를
유한형 준연접 $\mathcal{O}_X$-가군이라 하자. $s \in S$라 하자.
$\mathcal{F}$가 $X_s$의 모든 점에서 $S$ 위로 평탄하다고 가정하자.
$f(x') = s'$이고 $s' \leadsto s$가 $S$의 특수화가 되도록
$x' \in \text{Ass}_{X/S}(\mathcal{F})$를 택하자. $x'$가 $X_s$의
한 점으로 특수화하면, 어떤 $x \in \text{Ass}_{X_s}(\mathcal{F}_s)$에
대하여 $x' \leadsto x$이다.
\end{lemma}

\begin{proof}
$t \in X_s$에 대하여 $x' \leadsto t$라고 하자. 그러면 $x$가
$\overline{\{x'\}} \cap f^{-1}(\{s\})$의 한 기약 성분의 일반점에
대응하도록 하는 특수화 $x' \leadsto x \leadsto t$를 찾을 수 있다.
가정에 의해 $\mathcal{F}$는 $x$에서 $S$ 위로 평탄하다.
More on Morphisms, Lemma
\ref{more-morphisms-lemma-associated-point-specializes}에 의해 원하는 대로
$x \in \text{Ass}_{X/S}(\mathcal{F})$이다.
\end{proof}

\begin{lemma}
\label{flat-lemma-criterion}
$f : X \to S$를 유한형인 스킴의 사상이라 하고, $\mathcal{F}$를
유한형 준연접 $\mathcal{O}_X$-가군이라 하자. $s \in S$라 하자.
$(S', s') \to (S, s)$를 기본 에탈 근방이라 하고,
$$
\xymatrix{
X \ar[d] & X' \ar[l]^g \ar[d] \\
S & S' \ar[l]
}
$$
을 스킴 사상들의 가환 그림이라 하자. 다음을 가정한다.
\begin{enumerate}
\item $\mathcal{F}$는 $X_s$의 모든 점에서 $S$ 위로 평탄하다.
\item $X' \to S'$는 유한형이다.
\item $g^*\mathcal{F}$는 $X'_{s'}$를 따라 순수하다.
\item $g : X' \to X$는 에탈이다.
\item $g(X')$는 $\text{Ass}_{X_s}(\mathcal{F}_s)$를 포함한다.
\end{enumerate}
이 상황에서 $\mathcal{F}$가 $X_s$를 따라 순수할 필요충분조건은
$X' \to X \times_S S'$의 상이 $s'$로 특수화하는 $S'$의 점들 위에 놓인
$\text{Ass}_{X \times_S S'/S'}(\mathcal{F} \times_S S')$의 점들을
포함하는 것이다.
\end{lemma}

\begin{proof}
사상 $S' \to S$는 에탈이므로 $\mathcal{F}$가 $X_s$를 따라 순수하면
% P05-EN-CH38-0184: 기저변환 뒤의 올을 X_{s'}가 아니라 X_s로 쓰는 원문을 그대로 보존했다.
$\mathcal{F} \times_S S'$는 $X_s$를 따라 순수하다. Lemma
\ref{flat-lemma-quasi-finite-base-change}를 보라. 순수성은 평탄 하강을
만족하므로(Lemma \ref{flat-lemma-flat-descend-pure}를 보라),
$\mathcal{F} \times_S S'$가 $X_{s'}$를 따라 순수하면 $\mathcal{F}$는
$X_s$를 따라 순수하다. 따라서 $S$를 $S'$로 바꾸고 $S = S'$라고
가정해도 된다. 이때 $g : X' \to X$는 $S$ 위의 유한형 스킴들 사이의
에탈 사상이다. 더욱이 $S$를 $\Spec(\mathcal{O}_{S, s})$로 바꾸고
$S$가 국소라고 가정해도 된다.

\medskip\noindent
먼저 $\mathcal{F}$가 $X_s$를 따라 순수하다고 가정하자. 이 경우
순수성에 의해 $\text{Ass}_{X/S}(\mathcal{F})$의 모든 점은 $X_s$의
한 점으로 특수화한다. 따라서 Lemma
\ref{flat-lemma-associated-point-specializes}에 의해
$\text{Ass}_{X/S}(\mathcal{F})$의 모든 점은
$\text{Ass}_{X_s}(\mathcal{F}_s)$의 한 점으로 특수화한다.
그러므로 $\text{Ass}_{X/S}(\mathcal{F})$의 모든 점은 $g$의 상에 속한다
(그 상은 열린집합이고 $\text{Ass}_{X_s}(\mathcal{F}_s)$를 포함한다).

\medskip\noindent
역으로 $g(X')$가 $\text{Ass}_{X/S}(\mathcal{F})$를 포함한다고 가정하자.
$S^h = \Spec(\mathcal{O}_{S, s}^h)$를 $s$에서 $S$의 헨젤화라 하자.
% P05-EN-CH38-0185: denote-by 구문에서 by가 빠진 원문을 한국어로 자연스럽게 나타냈다.
$S^h \to S$에 의한 $g$의 기저변환을
$g^h : (X')^h \to X^h$로 나타내고, $\mathcal{F}$의 $X^h$로의 당김을
$\mathcal{F}^h$로 나타내자. Divisors, Lemma
\ref{divisors-lemma-base-change-relative-assassin}과 Remark
\ref{divisors-remark-base-change-relative-assassin}에 의해 상대 연관점 집합
$\text{Ass}_{X^h/S^h}(\mathcal{F}^h)$는 사영 $X^h \to X$에 의한
$\text{Ass}_{X/S}(\mathcal{F})$의 역상이다. $g(X')$가
$\text{Ass}_{X/S}(\mathcal{F})$를 포함한다고 가정했으므로 기저변환
$g^h((X')^h) = g(X') \times_S S^h$는
$\text{Ass}_{X^h/S^h}(\mathcal{F}^h)$를 포함한다. 이로써 $S$가
헨젤 국소환의 스펙트럼인 경우로 귀착된다.
$x \in \text{Ass}_{X/S}(\mathcal{F})$라 하자. 보조정리의 증명을
끝내려면 $x$가 $X_s$의 한 점으로 특수화함을 보여야 한다.
% P05-EN-CH38-0186: 특정 논의 앞의 관사가 빠진 원문을 자연스럽게 나타냈다.
Definition \ref{flat-definition-pure} 뒤의 논의에 있는 순수성의 판정법 (4)를
보라. 가정에 의해
% P05-EN-CH38-0187: x' 앞의 부정관사 일치 오류를 한국어로 자연스럽게 나타냈다.
$g(x') = x$를 만족하는 $x' \in X'$가 존재한다.
$g : X' \to X$는 에탈이므로
$x' \in \text{Ass}_{X'/S}(g^*\mathcal{F})$이다. Lemma
\ref{flat-lemma-etale-weak-assassin-up-down}을 보라(여기서는 $w \in S$가
$x'$의 상일 때 올들의 사상 $X'_w \to X_w$에 이를 적용한다).
$g^*\mathcal{F}$는 $X'_s$를 따라 순수하므로 어떤 $y \in X'_s$에 대하여
$x' \leadsto y$이다. 따라서 $x = g(x') \leadsto g(y)$이고,
원하는 대로 $g(y) \in X_s$이다.
\end{proof}

\begin{lemma}
\label{flat-lemma-finite-type-flat-pure-along-fibre-is-universal}
$f : X \to S$를 스킴의 사상이라 하고, $\mathcal{F}$를 준연접
$\mathcal{O}_X$-가군이라 하자. $s \in S$라 하자. 다음을 가정한다.
\begin{enumerate}
\item $f$는 유한형이다.
\item $\mathcal{F}$는 유한형이다.
\item $\mathcal{F}$는 $X_s$의 모든 점에서 $S$ 위로 평탄하다.
\item $\mathcal{F}$는 $X_s$를 따라 순수하다.
\end{enumerate}
그러면 $\mathcal{F}$는 $X_s$를 따라 보편적으로 순수하다.
\end{lemma}

\begin{proof}
먼저 예비 관찰을 하나 하자. $(S', s') \to (S, s)$가 기본 에탈
근방이라고 하자.
% P05-EN-CH38-0188: denote-by 구문에서 by가 빠진 원문을 한국어로 자연스럽게 나타냈다.
$\mathcal{F}$의 $X' = X \times_S S'$로의 당김을 $\mathcal{F}'$로 나타내자.
Definition \ref{flat-definition-pure} 뒤의 논의에 의해 $\mathcal{F}'$는
$X'_{s'}$를 따라 순수하다. 더욱이 $\mathcal{F}'$는 $X'_{s'}$를 따라
$S'$ 위로 평탄하다. 그러므로 $\mathcal{F}'$가 $X'_{s'}$를 따라
보편적으로 순수함을 증명하면 충분하다. 실제로 점이 지정된 스킴들의
임의의 사상 $(T, t) \to (S, s)$가 주어지면 올곱
$(T', t') = (T \times_S S', (t, s'))$는 $(T, t)$ 위에서 평탄하다.
따라서 $\mathcal{F}_{T'}$가 $X_{t'}$를 따라 순수하면 Lemma
\ref{flat-lemma-flat-descend-pure}에 의해 $\mathcal{F}_T$는 $X_t$를 따라
순수하다. 그러므로 증명 중 언제든 $(s, S)$를 기본 에탈 근방으로
바꾸어도 된다. 또한 문제는 국소적이므로 $S$를
$\Spec(\mathcal{O}_{S, s})$로 바꾸어도 된다.

\medskip\noindent
기본 에탈 근방 $(S', s') \to (S, s)$와 가환 그림
$$
\xymatrix{
X \ar[d] & X' \ar[l]^g \ar[d] \\
S & \Spec(\mathcal{O}_{S', s'}) \ar[l]
}
$$
을 다음과 같이 택하자. 사상
$X' \to X \times_S \Spec(\mathcal{O}_{S', s'})$는 에탈이고,
$X_s = g((X')_{s'})$이며, 스킴 $X'$는 아핀이고,
$\Gamma(X', g^*\mathcal{F})$는 자유 $\mathcal{O}_{S', s'}$-가군이다.
Lemma \ref{flat-lemma-finite-type-flat-along-fibre-free-variant}를 보라.
$X' \to \Spec(\mathcal{O}_{S', s'})$는 유한형임에 유의하자
(이는 에탈 사상과 유한형 사상의 기저변환의 합성인 준콤팩트 사상이다).
증명의 첫 문단의 예비 관찰에 의해 $S$를
$\Spec(\mathcal{O}_{S', s'})$로 바꾸어도 된다. 따라서 가환 그림
$$
\xymatrix{
X \ar[dr] & & X' \ar[ll]^g \ar[ld] \\
& S &
}
$$
이 존재한다고 가정해도 된다. 여기서 두 스킴은 $S$ 위에서 유한형이고,
$g$는 에탈이며, $X_s \subset g(X')$이고, $S$는 닫힌점이 $s$인
국소 스킴이며, $X'$는 아핀이고, $\Gamma(X', g^*\mathcal{F})$는 자유
$\Gamma(S, \mathcal{O}_S)$-가군이다. 이 경우 $g^*\mathcal{F}$는
$S$ 위에서 보편적으로 순수함에 유의하자. Lemma
\ref{flat-lemma-affine-locally-projective-pure}를 보라.

\medskip\noindent
이 상황에서 Lemma \ref{flat-lemma-criterion}를 적용한다. $\mathcal{F}$가
$X_s$를 따라 순수하고 $X_s$를 따라 $S$ 위로 평탄하다는 가정으로부터
$\text{Ass}_{X/S}(\mathcal{F}) \subset g(X')$를 얻는다. Divisors, Lemma
\ref{divisors-lemma-base-change-relative-assassin}과 Remark
\ref{divisors-remark-base-change-relative-assassin}에 의해, 점이 지정된
스킴들의 임의의 사상 $(T, t) \to (S, s)$에 대하여
$$
\text{Ass}_{X_T/T}(\mathcal{F}_T) \subset
(X_T \to X)^{-1}(\text{Ass}_{X/S}(\mathcal{F})) \subset
g(X') \times_S T = g_T(X'_T).
$$
이다. 따라서 위 그림의 $(T, t)$로의 기저변환에 Lemma
\ref{flat-lemma-criterion}를 적용하면 원하는 대로 $\mathcal{F}_T$가
$X_t$를 따라 순수하다고 결론짓는다.
\end{proof}

\begin{lemma}
\label{flat-lemma-finite-type-flat-pure-is-universal}
$f : X \to S$를 유한형인 스킴의 사상이라 하고, $\mathcal{F}$를
유한형 준연접 $\mathcal{O}_X$-가군이라 하자. $\mathcal{F}$가
$S$ 위로 평탄하다고 가정하자. 이 경우 $\mathcal{F}$가 $S$에 대해
순수할 필요충분조건은 $\mathcal{F}$가 $S$에 대해 보편적으로 순수한 것이다.
\end{lemma}

\begin{proof}
Lemma \ref{flat-lemma-finite-type-flat-pure-along-fibre-is-universal}와
정의들로부터 즉시 따른다.
\end{proof}

\begin{lemma}
\label{flat-lemma-limit-purity}
$I$를 유향 집합이라 하고, $(S_i, g_{ii'})$를 $I$ 위의 아핀 스킴들의
역계라 하자. $S = \lim_i S_i$라 두고 $s \in S$라 하자.
$g_i : S \to S_i$로 사영들을 나타내고 $s_i = g_i(s)$로 두자.
$f : X \to S$가 유한표현 사상이고 $\mathcal{F}$가 유한표현 준연접
$\mathcal{O}_X$-가군이며, $X_s$를 따라 순수하고 $X_s$의 모든 점에서
$S$ 위로 평탄하다고 하자. 그러면 어떤 $i \in I$, 유한표현 사상
$X_i \to S_i$, 그리고 유한표현 준연접 $\mathcal{O}_{X_i}$-가군
$\mathcal{F}_i$가 존재한다. 이 가군은 $(X_i)_{s_i}$를
따라 순수하고 $(X_i)_{s_i}$의 모든 점에서 $S_i$ 위로 평탄하며,
$X \cong X_i \times_{S_i} S$이고 $\mathcal{F}_i$의 $X$로의 당김은
$\mathcal{F}$와 동형이다.
\end{lemma}

\begin{proof}
$U \subset X$를 $\mathcal{F}$가 $S$ 위로 평탄한 점들의 집합이라 하자.
More on Morphisms, Theorem
\ref{more-morphisms-theorem-openness-flatness}에 의해 이는 $X$의 열린
부분스킴이다. 가정에 의해 $X_s \subset U$이다. $X_s$는 준콤팩트이므로
$X_s \subset U'$인 준콤팩트 열린집합 $U' \subset U$를 찾을 수 있다.
Limits, Lemma \ref{limits-lemma-descend-finite-presentation}에 의해 어떤
$i \in I$와, $S$로의 기저변환이
% P05-EN-CH38-0189: 내려보낼 f 대신 f_i 자신과 동형이라고 한 원문을 그대로 보존했다.
$f_i$와 동형인 유한표현 사상 $f_i : X_i \to S_i$를 찾을 수 있다.
하나를 고정하고 $X_{i'} = X_i \times_{S_i} S_{i'}$로 두자. 그러면
전이 사상들이 아핀인 $X = \lim_{i'} X_{i'}$를 얻는다. Limits, Lemma
\ref{limits-lemma-descend-modules-finite-presentation}에 의해, 필요하면
$i$를 키운 뒤 $S$로의 기저변환이 $\mathcal{F}$와 동형인 유한표현
준연접 $\mathcal{O}_{X_i}$-가군 $\mathcal{F}_i$가 존재한다고 가정해도
된다. Limits, Lemma \ref{limits-lemma-descend-opens}에 의해, 필요하면
$i$를 키운 뒤 $X$에서의 역상이 $U'$인 열린집합 $U'_i \subset X_i$가
존재한다고 가정해도 된다. 특히 $(X_i)_{s_i} \subset U'_i$임에 유의하자.
Limits, Lemma
\ref{limits-lemma-descend-module-flat-finite-presentation}에 의해($i$를
한 번 더 키운 뒤) $\mathcal{F}_i$가 $U'_i$ 위에서 평탄하다고 가정해도
된다. 특히 $\mathcal{F}_i$는 $(X_i)_{s_i}$를 따라 평탄하다.

\medskip\noindent
다음으로 Lemma \ref{flat-lemma-finite-presentation-flat-along-fibre}를 사용하여
기본 에탈 근방 $(S_i', s_i') \to (S_i, s_i)$와 스킴의 가환 그림
$$
\xymatrix{
X_i \ar[d] & X_i' \ar[l]^{g_i} \ar[d] \\
S_i & S_i' \ar[l]
}
$$
을 다음과 같이 택하자. $g_i$는 에탈이고,
$(X_i)_{s_i} \subset g_i(X_i')$이며, 스킴 $X_i'$, $S_i'$는 아핀이고,
$\Gamma(X_i', g_i^*\mathcal{F}_i)$는 사영
$\Gamma(S_i', \mathcal{O}_{S_i'})$-가군이다.
$g_i^*\mathcal{F}_i$는 $S'_i$ 위에서 보편적으로 순수함에 유의하자.
Lemma \ref{flat-lemma-affine-locally-projective-pure}를 보라.
임의의 $i' \geq i$에 대하여 위 그림을 기저변환하여 $S_{i'}$ 위의 사상들
$(S'_{i'}, s'_{i'}) \to (S_{i'}, s_{i'})$와
$g_{i'} : X'_{i'} \to X_{i'}$를 갖는 그림을 얻을 수 있다. 또한 이를
기저변환하여 $S$ 위의 사상들 $(S', s') \to (S, s)$와
$g : X' \to X$를 갖는 그림을 얻을 수 있다.

\medskip\noindent
이제 순수성의 판정법을 사용할 수 있다. 에탈 사상
$X'_i \to X_i \times_{S_i} S'_i$의 상을
$W'_i \subset X_i \times_{S_i} S'_i$라 하자. 모든 $i' \geq i$에 대하여
마찬가지로 상 $W'_{i'} \subset X_{i'} \times_{S_{i'}} S'_{i'}$가 있고,
상 $W' \subset X \times_S S'$가 있다. 상을 취하는 것은 기저변환과
가환하므로 $W'_{i'} = W'_i \times_{S'_i} S'_{i'}$이고,
% P05-EN-CH38-0190: 정의된 W'_i 대신 미정의 W_i를 쓰는 원문을 그대로 보존했다.
$W' = W_i \times_{S'_i} S'$이다. $\mathcal{F}$가 $X_s$를 따라 순수하므로
Lemma \ref{flat-lemma-criterion}에 의해
\begin{equation}
\label{flat-equation-inclusion}
% P05-EN-CH38-0191: 기저변환된 공간에서 원래 f의 역상을 쓰는 원문을 그대로 보존했다.
f^{-1}(\Spec(\mathcal{O}_{S', s'})) \cap
\text{Ass}_{X \times_S S'/S'}(\mathcal{F} \times_S S') \subset W'
\end{equation}
이다. More on Morphisms, Lemma
\ref{more-morphisms-lemma-relative-assassin-constructible}에 의해
$$
E = \{t \in S' \mid \text{Ass}_{X_t}(\mathcal{F}_t) \subset W' \}
\quad\text{이고}\quad
E_{i'} = \{t \in S'_{i'} \mid
\text{Ass}_{X_t}(\mathcal{F}_{i', t}) \subset W'_{i'} \}
$$
는 각각 $S'$와 $S'_{i'}$의 국소 구성가능 부분집합이다.
More on Morphisms, Lemma
\ref{more-morphisms-lemma-base-change-assassin-in-U}에 의해 $E_{i'}$는
사상 $S'_{i'} \to S'_i$에 의한 $E_i$의 역상이고, $E$는 사상
$S' \to S'_i$에 의한 $E_i$의 역상이다. 따라서 Equation
(\ref{flat-equation-inclusion})은 $\Spec(\mathcal{O}_{S', s'})$가 $E_i$ 안으로
간다는 명제와 동치이다.
$\mathcal{O}_{S', s'} =
\colim_{i' \geq i} \mathcal{O}_{S'_{i'}, s'_{i'}}$
이므로 어떤 $i' \geq i$에 대하여
$\Spec(\mathcal{O}_{S'_{i'}, s'_{i'}})$가 $E_i$ 안으로 간다.
Limits, Lemma \ref{limits-lemma-limit-contained-in-constructible}을 보라.
이제 $S_{i'}$ 위의 상황에 Lemma \ref{flat-lemma-criterion}를 적용하면
$\mathcal{F}_{i'}$가 $(X_{i'})_{s_{i'}}$를 따라 순수하다고 결론짓는다.
\end{proof}

\begin{lemma}
\label{flat-lemma-flat-finite-presentation-purity-open}
$f : X \to S$를 유한표현 사상이라 하고, $\mathcal{F}$를 $S$ 위로
평탄한 유한표현 준연접 $\mathcal{O}_X$-가군이라 하자. 그러면 집합
$$
U = \{s \in S \mid \mathcal{F}\text{가 }X_s\text{를 따라 순수}\}
$$
는 $S$에서 열려 있다.
\end{lemma}

\begin{proof}
$s \in U$라 하자. Lemma
\ref{flat-lemma-finite-presentation-flat-along-fibre}를 사용하면 기본 에탈 근방
$(S', s') \to (S, s)$와 가환 그림
$$
\xymatrix{
X \ar[d] & X' \ar[l]^g \ar[d] \\
S & S' \ar[l]
}
$$
을 다음과 같이 찾을 수 있다. $g$는 에탈이고, $X_s \subset g(X')$이며,
스킴 $X'$, $S'$는 아핀이고, $\Gamma(X', g^*\mathcal{F})$는 사영
$\Gamma(S', \mathcal{O}_{S'})$-가군이다. $g^*\mathcal{F}$는 $S'$ 위에서
보편적으로 순수함에 유의하자. Lemma
\ref{flat-lemma-affine-locally-projective-pure}를 보라. 에탈 사상
$X' \to X \times_S S'$의 상을 $W' \subset X \times_S S'$라 하자.
% P05-EN-CH38-0192: 정의된 W' 대신 미정의 W를 쓰는 원문을 그대로 보존했다.
$W$는 $S'$ 위에서 열리고 준콤팩트임에 유의하자. 다음과 같이 두자.
$$
E = \{t \in S' \mid \text{Ass}_{X_t}(\mathcal{F}_t) \subset W' \}.
$$
More on Morphisms, Lemma
\ref{more-morphisms-lemma-relative-assassin-constructible}에 의해 $E$는
$S'$의 구성가능 부분집합이다. Lemma \ref{flat-lemma-criterion}에 의해
$\Spec(\mathcal{O}_{S', s'}) \subset E$이다. Morphisms, Lemma
\ref{morphisms-lemma-constructible-containing-open}에 의해 $E$는 $s'$의
열린 근방 $V'$를 포함한다. Lemma \ref{flat-lemma-criterion}를 한 번 더
적용하면 $S$에서 $V'$의 상에 속하는 임의의 점 $s_1$에 대하여 층
$\mathcal{F}$가 $X_{s_1}$를 따라 순수함을 안다. $S' \to S$는
에탈이므로 $S$에서 $V'$의 상은 열려 있고, 결론을 얻는다.
\end{proof}


\section{순수성의 사용법}
\label{flat-section-applications-purity}

\noindent
순수성을 사용하는 방법의 몇 가지 예를 든다. 첫 보조정리는 실제로
조금 더 약한 형태의 순수성을 사용한다.

\begin{lemma}
\label{flat-lemma-injectivity-map-source-flat-pure}
$f : X \to S$를 유한형 사상이라 하고, $\mathcal{F}$를 $X$ 위의
유한형 준연접 층이라 하자. $S$가 닫힌점이 $s$인 국소 스킴이라고
가정하자. $\mathcal{F}$가 $X_s$를 따라 순수하고 $\mathcal{F}$가
$S$ 위로 평탄하다고
가정하자.
% P05-EN-CH38-0193: 사상 도입 동사가 빠진 원문을 한국어로 자연스럽게 나타냈다.
$\varphi : \mathcal{F} \to \mathcal{G}$를 준연접
$\mathcal{O}_X$-가군들의 사상이라 하자. 다음은 서로 동치이다.
\begin{enumerate}
\item 모든 $x \in \text{Ass}_{X_s}(\mathcal{F}_s)$에 대하여 줄기 위의
사상 $\varphi_x$는 단사이다.
\item $\varphi$는 단사이다.
\end{enumerate}
\end{lemma}

\begin{proof}
$\mathcal{K} = \Ker(\varphi)$라 하자. 목표는 $\mathcal{K} = 0$을
증명하는 것이다. 이를 위해서는
$\text{WeakAss}_X(\mathcal{K}) = \emptyset$임을 증명하면 충분하다.
Divisors, Lemma \ref{divisors-lemma-weakly-ass-zero}를 보라.
Divisors, Lemma \ref{divisors-lemma-ses-weakly-ass}에 의해
$\text{WeakAss}_X(\mathcal{K}) \subset \text{WeakAss}_X(\mathcal{F})$이다.
$\mathcal{F}$가 평탄하므로 Lemma
\ref{flat-lemma-bourbaki-finite-type-general-base}에 의해
$\text{WeakAss}_X(\mathcal{F}) \subset \text{Ass}_{X/S}(\mathcal{F})$이다.
순수성에 의해 $\text{Ass}_{X/S}(\mathcal{F})$의 임의의 점 $x'$는
$X_s$의 한 점의 일반화이고, 따라서 Lemma
\ref{flat-lemma-associated-point-specializes}에 의해
% P05-EN-CH38-0194: 특수화 방향을 거꾸로 서술하는 원문을 그대로 보존했다.
어떤 점 $x \in \text{Ass}_{X_s}(\mathcal{F}_s)$의 특수화이다.
따라서 $\varphi_x$의 단사성은 $\varphi_{x'}$의 단사성을 함의하고,
그러므로 $\mathcal{K}_{x'} = 0$이다.
\end{proof}

\begin{proposition}
\label{flat-proposition-finite-presentation-flat-pure-is-projective}
$f : X \to S$를 유한표현 아핀 사상이라 하고, $\mathcal{F}$를
$S$ 위로 평탄한 유한표현 준연접 $\mathcal{O}_X$-가군이라 하자.
다음은 서로 동치이다.
\begin{enumerate}
\item $f_*\mathcal{F}$는 $S$ 위에서 국소 사영이다.
\item $\mathcal{F}$는 $S$에 대해 순수하다.
\end{enumerate}
특히 유한표현 환 준동형 $A \to B$와 $A$ 위로 평탄한 유한표현
$B$-가군 $N$이 주어졌을 때 다음이 성립한다. $N$이 $A$-가군으로서
사영일 필요충분조건은 $\Spec(B)$ 위의 $\widetilde{N}$이
$\Spec(A)$에 대해 순수한 것이다.
\end{proposition}

\begin{proof}
(1) $\Rightarrow$ (2)는 Lemma
\ref{flat-lemma-affine-locally-projective-pure}이다. $\mathcal{F}$가 $S$에 대해
순수하다고 가정하자. Lemma
\ref{flat-lemma-finite-type-flat-pure-along-fibre-is-universal}에 의해 이는
$\mathcal{F}$가 임의의 기저변환 뒤에도 순수함을 뜻함에 유의하자.
Descent, Lemma \ref{descent-lemma-locally-projective-descends}에 의해
$f_*\mathcal{F}$가 $S$ 위에서 fpqc 국소적으로 사영임을 증명하면 충분하다.
$s \in S$를 택하자. $f_*\mathcal{F}$를 $s$의 한 에탈 근방에 제한한
것이 국소 사영임을 증명하겠다. 실제로 Lemma
\ref{flat-lemma-finite-presentation-flat-along-fibre}에 의해 $S$를 $s$의
아핀 기본 에탈 근방으로 바꾼 뒤, 다음 그림이 존재한다고 가정해도 된다.
$$
\xymatrix{
X \ar[dr] & & X' \ar[ll]^g \ar[ld] \\
& S &
}
$$
여기서 두 스킴은 $S$ 위에서 아핀이고 유한표현이며, $g$는 에탈이고,
$X_s \subset g(X')$이며, $\Gamma(X', g^*\mathcal{F})$는 사영
$\Gamma(S, \mathcal{O}_S)$-가군이다. 이 경우 $g^*\mathcal{F}$는
$S$ 위에서 보편적으로 순수함에 유의하자. Lemma
\ref{flat-lemma-affine-locally-projective-pure}를 보라. 따라서 Lemma
\ref{flat-lemma-criterion}에 의해 열린집합 $g(X')$는
$\Spec(\mathcal{O}_{S, s})$ 위에 놓이는
$\text{Ass}_{X/S}(\mathcal{F})$의 점들을 포함한다. 다음과 같이 두자.
$$
E = \{t \in S \mid \text{Ass}_{X_t}(\mathcal{F}_t) \subset g(X') \}.
$$
More on Morphisms, Lemma
\ref{more-morphisms-lemma-relative-assassin-constructible}에 의해 $E$는
$S$의 구성가능 부분집합이다. 이미
$\Spec(\mathcal{O}_{S, s}) \subset E$임을 보았다. Morphisms, Lemma
\ref{morphisms-lemma-constructible-containing-open}에 의해 $E$는 $s$의
한 열린 근방을 포함한다. 따라서 $S$를 $s$의 아핀 근방으로 바꾼 뒤
$\text{Ass}_{X/S}(\mathcal{F}) \subset g(X')$라고 가정해도 된다.
Lemma \ref{flat-lemma-base-change-universally-flat}에 의해 이는
$$
\Gamma(X, \mathcal{F}) \longrightarrow \Gamma(X', g^*\mathcal{F})
$$
가 $\Gamma(S, \mathcal{O}_S)$-보편 단사라는 뜻이다. Algebra, Lemma
\ref{algebra-lemma-pure-submodule-ML}에 의해
% P05-EN-CH38-0195: Gamma 앞의 부정관사 오류를 한국어로 자연스럽게 나타냈다.
$\Gamma(X, \mathcal{F})$는 $\Gamma(S, \mathcal{O}_S)$-가군으로서
Mittag-Leffler이다. $\Gamma(X, \mathcal{F})$는
$\Gamma(S, \mathcal{O}_S)$-가군으로서 가산 생성이고 평탄하므로,
Algebra, Lemma \ref{algebra-lemma-countgen-projective}에 의해 사영이다.
\end{proof}

\noindent
이 명제를 사용하여 앞의 결과들 가운데 일부를 개선할 수 있다.
다음 보조정리는 Proposition
\ref{flat-proposition-finite-presentation-flat-at-point}의 개선이다.

\begin{lemma}
\label{flat-lemma-flat-finite-presentation-affine-neighbourhood-projective}
$f : X \to S$를 국소 유한표현 사상이라 하고, $\mathcal{F}$를 유한표현
준연접 $\mathcal{O}_X$-가군이라 하자. $s = f(x) \in S$인
$x \in X$를 택하자. $\mathcal{F}$가 $x$에서 $S$ 위로 평탄하면,
% P05-EN-CH38-0196: 조건절과 결론 사이의 구두점이 빠진 원문을 한국어로 자연스럽게 나타냈다.
아핀 기본 에탈 근방 $(S', s') \to (S, s)$와
$x' = (x, s')$를 포함하는 아핀 열린집합
$U' \subset X \times_S S'$가 존재하여
$\Gamma(U', \mathcal{F}|_{U'})$는 사영
$\Gamma(S', \mathcal{O}_{S'})$-가군이다.
\end{lemma}

\begin{proof}
증명 중에 $X$를 $x$의 열린 근방으로, $S$를 $s$의 기본 에탈 근방으로
바꾸어도 된다. 따라서 평탄성의 개방성에 의해(More on Morphisms,
Theorem \ref{more-morphisms-theorem-openness-flatness}를 보라)
$\mathcal{F}$가 $S$ 위로 평탄하다고 가정해도 된다. $S$와 $X$가
아핀이라고 가정해도 된다. $X$를 더 축소한 뒤
$\text{Ass}_{X_s}(\mathcal{F}_s)$의 임의의 점이 $x$의 일반화라고
가정해도 된다. 이 성질은 $(S, s)$를 기본 에탈 근방으로 바꾸어도
보존된다. 따라서 Lemma
\ref{flat-lemma-finite-presentation-flat-along-fibre}를 적용하여 다음 그림이
존재하는 상황에 이를 수 있다.
$$
\xymatrix{
X \ar[dr] & & X' \ar[ll]^g \ar[ld] \\
& S &
}
$$
여기서 두 스킴은 $S$ 위에서 아핀이고 유한표현이며, $g$는 에탈이고,
$X_s \subset g(X')$이며, $\Gamma(X', g^*\mathcal{F})$는 사영
$\Gamma(S, \mathcal{O}_S)$-가군이다. 이 경우 $g^*\mathcal{F}$는
$S$ 위에서 보편적으로 순수함에 유의하자. Lemma
\ref{flat-lemma-affine-locally-projective-pure}를 보라.

\medskip\noindent
$U \subset g(X')$를 $x$의 아핀 열린 근방이라 하자.
$\mathcal{F}|_U$가 $U_s$를 따라 순수하다고 주장한다. 이를 증명하면,
$S$를 축소한 뒤 $\mathcal{F}|_U$가 $S$에 대해 순수하게 되고(Lemma
\ref{flat-lemma-flat-finite-presentation-purity-open}을 보라), 이어서
Proposition \ref{flat-proposition-finite-presentation-flat-pure-is-projective}에
의해 사영성이 따르므로 보조정리가 증명된다. 주장을 증명하려면
$(S, s)$를 임의의 기본 에탈 근방으로 바꾼 뒤, 어떤
$s' \in S$, $s' \leadsto s$ 위에 놓이는
$\text{Ass}_{U/S}(\mathcal{F}|_U)$의 임의의 점 $\xi$가 $U_s$의 한 점으로
특수화함을 보여야 한다. $U \subset g(X')$이므로
% P05-EN-CH38-0197: xi' 앞의 부정관사 오류를 한국어로 자연스럽게 나타냈다.
$g(\xi') = \xi$인 $\xi' \in X'$를 찾을 수 있다.
$g^*\mathcal{F}$가 $S$ 위에서 순수하므로 Lemma
\ref{flat-lemma-associated-point-specializes}를 사용하면 어떤
$x' \in \text{Ass}_{X'_s}(g^*\mathcal{F}_s)$에 대하여 특수화
% P05-EN-CH38-0198: 올로 제한한 층에 g^*를 직접 적용하는 원문 표기를 그대로 보존했다.
$\xi' \leadsto x'$가 존재함을 안다. 그러면
$g(x') \in \text{Ass}_{X_s}(\mathcal{F}_s)$이다(예를 들어 뇌터 스킴들의
에탈 사상 $X'_s \to X_s$에 Lemma
\ref{flat-lemma-etale-weak-assassin-up-down}을 적용하라). 따라서 위에서 택한
$X$에 의해 $g(x') \leadsto x$이다! $x \in U$이므로
$g(x') \in U$라고 결론짓는다. 따라서 원하는 대로
$\xi = g(\xi') \leadsto g(x') \in U_s$이다.
\end{proof}

\noindent
다음 보조정리는 Lemma
\ref{flat-lemma-finite-type-flat-at-point-free-variant}의 개선이다.

\begin{lemma}
\label{flat-lemma-flat-finite-type-affine-neighbourhood-projective}
$f : X \to S$를 국소 유한형 사상이라 하고, $\mathcal{F}$를 유한형
준연접 $\mathcal{O}_X$-가군이라 하자. $s = f(x) \in S$인
$x \in X$를 택하자. $\mathcal{F}$가 $x$에서 $S$ 위로 평탄하면,
% P05-EN-CH38-0199: 조건절과 결론 사이의 구두점이 빠진 원문을 한국어로 자연스럽게 나타냈다.
아핀 기본 에탈 근방 $(S', s') \to (S, s)$와
$x' = (x, s')$를 포함하는 아핀 열린집합
$U' \subset X \times_S \Spec(\mathcal{O}_{S', s'})$가 존재하여
$\Gamma(U', \mathcal{F}|_{U'})$는 자유
$\mathcal{O}_{S', s'}$-가군이다.
\end{lemma}

\begin{proof}
문제는 $X$와 $S$ 위에서 자리스키 국소적이다. 따라서 $X$와 $S$가
아핀이라고 가정해도 된다. 그러면 $S$ 위의 닫힌 몰입
$i : X \to \mathbf{A}^n_S$를 찾을 수 있다. $\mathbf{A}^n_S$ 위의 층
$i_*\mathcal{F}$와 점 $i(x)$에 대하여 보조정리를 증명하면 충분함은
분명하다. 이로써 $X \to S$가 유한표현인 경우로 귀착된다. $S$를
$\Spec(\mathcal{O}_{S', s'})$로, $X$를
$X \times_S \Spec(\mathcal{O}_{S', s'})$의 열린집합으로 바꾼 뒤
$\mathcal{F}$가 유한표현이라고 가정해도 된다. Proposition
\ref{flat-proposition-finite-type-flat-at-point}를 보라. 이 경우 Lemma
\ref{flat-lemma-flat-finite-presentation-affine-neighbourhood-projective}와
Algebra, Theorem \ref{algebra-theorem-projective-free-over-local-ring}을
적용하면 결론을 얻는다.
\end{proof}

\begin{lemma}
\label{flat-lemma-flat-finite-type-local-colimit-free}
$A \to B$를 본질적으로 유한형인 국소환들의 국소 환 준동형이라 하자.
$N$을 $A$-가군으로서 평탄한 유한 $B$-가군이라 하자. $A$가
헨젤이면 $N$은 자유 $A$-가군 $F_i$들의 여과 여극한
$$
N = \colim_i F_i
$$
이다. 여기서 이 계의 모든 전이 사상 $u_i : F_i \to F_{i'}$는 단사 사상
$\overline{u}_i : F_i/\mathfrak m_AF_i \to F_{i'}/\mathfrak m_AF_{i'}$를
유도한다. 또한 $N$은 Mittag-Leffler $A$-가군이다.
\end{lemma}

\begin{proof}
유한형 사상 $X \to S = \Spec(A)$, $S$의 닫힌점 $s$ 위에 놓이는 점
$x \in X$, 그리고 $\mathcal{F}_x \cong N$이 $A$-가군의 동형이 되도록
하는 유한형 준연접 $\mathcal{O}_X$-가군 $\mathcal{F}$를 찾을 수 있다.
$X$를 축소한 뒤 $\text{Ass}_{X_s}(\mathcal{F}_s)$의 각 점이 $x$로
특수화한다고 가정해도 된다. Lemma
\ref{flat-lemma-flat-finite-type-affine-neighbourhood-projective}에 의해,
$\Gamma(U_i, \mathcal{F})$가 자유 $A$-가군 $F_i$가 되도록 하는 $x$의
아핀 열린 근방 $U_i \subset X$들의 근방 기저가 존재한다.
$U_{i'} \subset U_i$이면
$$
F_i/\mathfrak m_AF_i = \Gamma(U_{i, s}, \mathcal{F}_s)
\longrightarrow
\Gamma(U_{i', s}, \mathcal{F}_s) = F_{i'}/\mathfrak m_AF_{i'}
$$
는 단사임에 유의하자. 실제로 핵의 한 단면은 $x$와 만나지 않는
$X_s$의 닫힌 부분집합에 지지되어야 하는데, 이는 위에서 $X$를 택한
방식에 모순이다. 사상 $F_i \to F_{i'}$들은 $A$-보편 단사이므로
(Lemma \ref{flat-lemma-universally-injective-local}), Algebra, Lemma
\ref{algebra-lemma-colimit-universally-injective-ML}에 의해 $N$은
Mittag-Leffler이다.
\end{proof}

\noindent
처음부터 차례로 읽는 경우에는 다음 보조정리를 건너뛰어야 한다.

\begin{lemma}
\label{flat-lemma-flat-finite-type-local-valuation-ring-has-content}
$A \to B$를 본질적으로 유한형인 국소환들의 국소 환 준동형이라 하자.
$N$을 $A$-가군으로서 평탄한 유한 $B$-가군이라 하자. $A$가 값매김환이면
$N$의 임의의 원소는 내용 아이디얼 $I \subset A$를 갖는다(More on
Algebra, Definition \ref{more-algebra-definition-content-ideal}).
또한 $I$는 주 아이디얼이다.
\end{lemma}

\begin{proof}
마지막 명제는 More on Algebra, Lemma
\ref{more-algebra-lemma-content-finitely-generated}과 Algebra, Lemma
\ref{algebra-lemma-characterize-valuation-ring}에 의해 $I$가 유한 생성
아이디얼이라는 사실에서 따른다.

\medskip\noindent
이제 $I$의 존재를 증명하자. $A \subset A^h$를 헨젤화라 하자.
$B'$를 $B \otimes_A A^h$의 다음 극대 아이디얼에서의 국소화라 하자.
% P05-EN-CH38-0200: 텐서곱의 기저환 첨자가 빠진 원문을 그대로 보존했다.
$\mathfrak m_B \otimes A^h + B \otimes \mathfrak m_{A^h}$.
그러면 $B \to B'$는 평탄이고, 따라서 충실히 평탄이다.
$N' = N \otimes_B B'$라 하자. $x \in N$이라 하고 $x' \in N'$를
그 상이라 하자. 아이디얼 $I \subset A$에 대하여
$x \in IN \Leftrightarrow x' \in IN'$이라고 주장한다. 실제로
$N/IN \to N'/IN'$은 $B \to B'$와 $N/IN$의 텐서곱이고,
Algebra, Lemma
\ref{algebra-lemma-faithfully-flat-universally-injective}에 의해
$B \to B'$는 보편 단사이다. More on Algebra, Lemma
\ref{more-algebra-lemma-henselization-valuation-ring}과 Algebra, Lemma
\ref{algebra-lemma-ideals-valuation-ring}에 의해 사상 $A \to A^h$는
아이디얼들의 집합 사이에 포함 관계를 보존하는 전단사
$I \mapsto IA^h$를 정의한다. 따라서 $x$가 $A$에서 내용 아이디얼을
가질 필요충분조건은 $x'$가 $A^h$에서 내용 아이디얼을 갖는 것이다.
$x' \in N'$에 대한 명제는 Lemma
\ref{flat-lemma-flat-finite-type-local-colimit-free}와 Algebra, Lemma
\ref{algebra-lemma-minimal-contains}에서 따른다.
\end{proof}

\noindent
한 가지 응용은 다음과 같다.

\begin{lemma}
\label{flat-lemma-proper-flat-over-dvr-reduced-fibre}
$R$이 값매김환일 때 $X \to \Spec(R)$를 고유 평탄 사상이라 하자.
특별올이 축약이면,
% P05-EN-CH38-0201: 복합 주어와 단수 동사의 불일치를 한국어로 자연스럽게 나타냈다.
$X$와 $X \to \Spec(R)$의 모든 올은 축약이다.
\end{lemma}

\begin{proof}
특별올 $X_s$가 축약이라고 가정하자. 임의의 점 $x \in X$를 택하고
$\mathcal{O}_{X, x}$가 축약임을 보이자. 그러면 $X$가 축약임이 증명된다.
$x'$가 특별올에 속하는 특수화 $x \leadsto x'$를 택하자. 고유 사상은
닫힌 사상이므로 이러한 특수화가 존재한다. 국소환
$A = \mathcal{O}_{X, x'}$를 생각하자. $\mathcal{O}_{X, x}$는 $A$의
국소화이므로 $A$가 축약임을 보이면 충분하다. 영이 아닌
$a \in A$를 택하고 $I = (\pi) \subset R$를 그 내용 아이디얼이라 하자.
Lemma \ref{flat-lemma-flat-finite-type-local-valuation-ring-has-content}를 보라.
그러면 $a = \pi a'$이고, $\mathfrak m \subset R$이 극대 아이디얼일 때
$a'$은 $A/\mathfrak mA$의 영이 아닌 원소로 간다. $A$는 $R$ 위로
평탄하여 $\pi$가 영인자가 아니므로, $a$가 멱영이면 $a'$도 멱영이다.
그러나 $a'$은 축약환
$A/\mathfrak m A = \mathcal{O}_{X_s, x'}$의 영이 아닌 원소로 간다.
$A$가 축약이 아니면 이는 모순이며, 이것이 보이고자 한 바이다.

\medskip\noindent
물론 $X$가 축약이면 $R$ 위의 $X$의 일반올도 축약이다.
$\mathfrak p \subset R$이 소아이디얼이면 Algebra, Lemma
\ref{algebra-lemma-make-valuation-rings}에 의해 $R/\mathfrak p$는
값매김환이다. 따라서 $X$의 $R/\mathfrak p$로의 기저변환에 대하여
논증을 되풀이하면 $\mathfrak p$ 위의 올이 축약임을 증명할 수 있다.
\end{proof}






\section{평탄화 함자}
\label{flat-section-flattening-functors}

\noindent
$S$를 스킴이라 하자. 극한이 $T$인 아핀 스킴들의 모든 유향 역계
$\{T_i\}_{i \in I}$에 대하여 $F(T) = \colim_i F(T_i)$이면 함자
$F : (\Sch/S)^{opp} \to \textit{Sets}$를 극한 보존이라고 함을 상기하자.

\begin{situation}
\label{flat-situation-iso}
$f : X \to S$를 스킴의 사상이라 하고,
$u : \mathcal{F} \to \mathcal{G}$를 준연접 $\mathcal{O}_X$-가군들의
준동형이라 하자. $S$ 위의 임의의 스킴 $T$에 대하여 $u$의 $T$로의
기저변환을 $u_T : \mathcal{F}_T \to \mathcal{G}_T$로 나타내겠다.
다시 말해 $u_T$는 사영 사상 $X_T = X \times_S T \to X$에 의한
$u$의 당김이다. 이 상황에서 다음 함자를 생각할 수 있다.
\begin{equation}
\label{flat-equation-iso}
F_{iso} : (\Sch/S)^{opp} \longrightarrow \textit{Sets}, \quad
T \longrightarrow \left\{
\begin{matrix}
\{*\} & \text{다음이면} & u_T \text{는 동형이다}, \\
\emptyset & \text{그렇지 않으면.} &
\end{matrix}
\right.
\end{equation}
$u_T$가 각각 단사, 전사, 또는 영일 것을 요구하는 변형
$F_{inj}$, $F_{surj}$, $F_{zero}$도 있다.
\end{situation}

\begin{lemma}
\label{flat-lemma-iso-sheaf}
Situation \ref{flat-situation-iso}의 상황이라 하자.
\begin{enumerate}
\item 각 함자 $F_{iso}$, $F_{inj}$, $F_{surj}$, $F_{zero}$는 fpqc
위상에 대한 층 성질을 만족한다.
\item $f$가 준콤팩트이고 $\mathcal{G}$가 유한형이면 $F_{surj}$는
극한 보존이다.
% P05-EN-CH38-0202: F의 유한형 조건에서 동사가 빠진 원문을 한국어로 자연스럽게 나타냈다.
\item $f$가 준콤팩트이고 $\mathcal{F}$가 유한형이면 $F_{zero}$는
극한 보존이다.
\item $f$가 준콤팩트이고 $\mathcal{F}$가 유한형이며 $\mathcal{G}$가
유한표현이면 $F_{iso}$는 극한 보존이다.
\end{enumerate}
\end{lemma}

\begin{proof}
$\{T_i \to T\}_{i \in I}$를 $S$ 위의 스킴들의 fpqc 덮개라 하자.
$X_i = X_{T_i} = X \times_S T_i$와 $u_i = u_{T_i}$로 두자.
$\{X_i \to X_T\}_{i \in I}$는 $X_T$의 fpqc 덮개임에 유의하자.
Topologies, Lemma \ref{topologies-lemma-fpqc}를 보라. 특히 모든
$x \in X_T$에 대하여 $x$로 가는 어떤 $i \in I$와
$x_i \in X_i$가 존재한다. $\mathcal{O}_{X_T, x} \to
\mathcal{O}_{X_i, x_i}$는 평탄이고 따라서 충실히 평탄하므로(Algebra,
Lemma \ref{algebra-lemma-local-flat-ff}를 보라), $(u_i)_{x_i}$가 단사,
전사, 전단사, 또는 영일 필요충분조건은 $(u_T)_x$가 각각 단사, 전사,
전단사, 또는 영인 것이다. 이로부터 보조정리의 (1)을 얻는다.

\medskip\noindent
이제 (2)를 증명하자. $f$가 준콤팩트이고 $\mathcal{G}$가 유한형이라고
가정하자. $T = \lim_{i \in I} T_i$를 아핀 $S$-스킴들의 유향 극한이라
하고 $u_T$가 전사라고 가정하자.
$X_i = X_{T_i} = X \times_S T_i$와
$u_i = u_{T_i} : \mathcal{F}_i = \mathcal{F}_{T_i}
\to \mathcal{G}_i = \mathcal{G}_{T_i}$.
(2)를 증명하려면 어떤 $i$에 대하여 $u_i$가 전사임을
보여야 한다. $i_0 \in I$를 택하고 $I$를
$\{i \mid i \geq i_0\}$로 바꾸자. $f$가 준콤팩트이므로 스킴
$X_{i_0}$는 준콤팩트이다. 따라서 아핀 열린집합
$W_1, \ldots, W_m \subset X$와 아핀 열린 덮개
$X_{i_0} = U_{1, i_0} \cup \ldots \cup U_{m, i_0}$
를 다음과 같이 택할 수 있다. 각 $U_{j, i_0}$는 사영 사상
$X_{i_0} \to X$에 의해 $W_j$ 안으로 간다. 임의의 $i \in I$에 대하여
$U_{j, i}$를 $U_{j, i_0}$의 역상이라 하자. $U_j = \lim_i U_{j, i}$로
두면 $X_T = U_1 \cup \ldots \cup U_m$은 $X_T$의 아핀 열린 덮개이다.
이제 주어진 $j \in \{1, \ldots, m\}$마다 어떤
$i = i(j) \in I$에 대하여 $u_i|_{U_{j, i}}$가 전사임을 보이면 충분하다.
Properties, Lemma \ref{properties-lemma-finite-type-module}를 사용하면 이는
다음 대수 문제로 옮겨진다. $A$를 환이라 하고
$u : M \to N$을 $A$-가군 사상이라 하자.
$R = \colim_{i \in I} R_i$가 $A$-대수들의 유향 여극한이라고 하자.
$N$이 유한 $A$-가군이고
$u \otimes 1 : M \otimes_A R \to N \otimes_A R$이 전사이면, 어떤
$i$에 대하여 $u \otimes 1 : M \otimes_A R_i \to N \otimes_A R_i$가
전사이다. 이는 Algebra, Lemma
\ref{algebra-lemma-module-map-property-in-colimit}의 (2)이다.

\medskip\noindent
이제 (3)을 증명하자. (2)의 증명과 정확히 같은 논증으로 다음 대수
문제로 귀착된다. $A$를 환이라 하고 $u : M \to N$을 $A$-가군
사상이라 하자. $R = \colim_{i \in I} R_i$가 $A$-대수들의 유향
여극한이라고 하자. $M$이 유한 $A$-가군이고
$u \otimes 1 : M \otimes_A R \to N \otimes_A R$이 영이면, 어떤
$i$에 대하여 사상
$u \otimes 1 : M \otimes_A R_i \to N \otimes_A R_i$가 영이다.
이는 Algebra, Lemma \ref{algebra-lemma-module-map-property-in-colimit}의
(1)이다.

\medskip\noindent
이제 (4)를 증명하자.
% P05-EN-CH38-0203: 명제보다 강한 유한표현 가정을 두는 원문 증명을 그대로 보존했다.
$f$가 준콤팩트이고 $\mathcal{F}, \mathcal{G}$가 유한표현이라고
가정하자. 앞 문단과 정확히 같은 방식으로 논증하면(Properties, Lemma
\ref{properties-lemma-finite-presentation-module}도 추가로 사용한다),
% P05-EN-CH38-0204: (4)의 증명에서 (3)을 옮긴다고 하는 원문 번호를 그대로 보존했다.
(3)은 다음 대수 명제로 옮겨진다. $A$를 환이라 하고
$u : M \to N$을 $A$-가군 사상이라 하자.
$R = \colim_{i \in I} R_i$가 $A$-대수들의 유향 여극한이라고 하자.
$M$이 유한 $A$-가군이고 $N$이 유한표현 $A$-가군이며
$u \otimes 1 : M \otimes_A R \to N \otimes_A R$이 동형이라고
가정하자. 그러면 어떤 $i$에 대하여 사상
$u \otimes 1 : M \otimes_A R_i \to N \otimes_A R_i$는 동형이다.
이는 Algebra, Lemma \ref{algebra-lemma-module-map-property-in-colimit}의
(3)이다.
\end{proof}

\begin{situation}
\label{flat-situation-flat-at-point}
$(A, \mathfrak m_A)$를 국소환이라 하자.
$\mathcal{C}$를 다음과 같은 범주라 하자. 그 대상은 국소환인
$A$-대수 $A'$ 가운데 대수 구조 $A \to A'$가 국소환들의 국소
준동형인 것들이다. $\mathcal{C}$의 대상 $A', A''$ 사이의 사상은
$A$-대수들의 국소 준동형 $A' \to A''$이다.
$A \to B$를 국소환들의 국소환 사상이라 하고 $M$을 $B$-가군이라
하자. $A'$가 $\mathcal{C}$의 대상이면 $B' = B \otimes_A A'$로 두고,
$M' = M \otimes_A A'$를 $B'$-가군으로 둔다.
$A' \in \Ob(\mathcal{C})$가 주어졌을 때 다음 조건을 생각하자.
\begin{equation}
\label{flat-equation-flat-at-primes}
\forall \mathfrak q \in V(\mathfrak m_{A'}B' + \mathfrak m_B B')
\subset \Spec(B') :
M'_{\mathfrak q}\text{는 }A'\text{ 위에서 평탄이다}.
\end{equation}
More on Algebra, Equation (\ref{more-algebra-equation-flat-at-primes})과의
유사성에 유의하자. 특히 $A' \to A''$가 $\mathcal{C}$의 사상이고
(\ref{flat-equation-flat-at-primes})가 $A'$에 대하여 성립하면 $A''$에
대해서도 성립한다. More on Algebra, Lemma
\ref{more-algebra-lemma-base-change-flat-at-primes}를 보라. 따라서 다음
함자를 얻는다.
\begin{equation}
\label{flat-equation-flat-at-point}
F_{lf} : \mathcal{C} \longrightarrow \textit{Sets}, \quad
A' \longrightarrow \left\{
\begin{matrix}
\{*\} & \text{조건 }(\ref{flat-equation-flat-at-primes})\text{가 성립하면}, \\
\emptyset & \text{그렇지 않으면.} &
\end{matrix}
\right.
\end{equation}
\end{situation}

\begin{lemma}
\label{flat-lemma-flat-at-point}
Situation \ref{flat-situation-flat-at-point}의 상황이라 하자.
\begin{enumerate}
\item $A' \to A''$가 $\mathcal{C}$의 평탄 사상이면
$F_{lf}(A') = F_{lf}(A'')$이다.
\item $A \to B$가 본질적으로 유한표현이고 $M$이 유한표현
$B$-가군이면 $F_{lf}$는 극한 보존이다. 즉 $\{A_i\}_{i \in I}$가
$\mathcal{C}$의 대상들의 유향계를 이루면
$F_{lf}(\colim_i A_i) = \colim_i F_{lf}(A_i)$.
\end{enumerate}
\end{lemma}

\begin{proof}
(1)은 More on Algebra, Lemma
\ref{more-algebra-lemma-flat-descent-flat-at-primes}의 특수한 경우이다.
(2)는 More on Algebra, Lemma
\ref{more-algebra-lemma-limit-preserving-flat-at-primes}의 특수한 경우이다.
\end{proof}

\begin{lemma}
\label{flat-lemma-flat-at-point-finite}
Situation \ref{flat-situation-flat-at-point}의 상황이라 하자.
% P05-EN-CH38-0205: 동사와 관사가 손상된 원문 문장을 한국어로 자연스럽게 나타냈다.
$B \to C$를 국소 $A$-대수들의 국소 사상이라 하고 $N$을
$C$-가군이라 하자. 쌍 $(C, N)$에 결부된 함자를
$F'_{lf} : \mathcal{C} \to \textit{Sets}$라 하자. $M \cong N$이
$B$-가군으로서 성립하고 $B \to C$가 유한이면
$F_{lf} = F'_{lf}$.
\end{lemma}

\begin{proof}
$A'$를 $\mathcal{C}$의 대상이라 하자. Situation
\ref{flat-situation-flat-at-point}에서 $B'$, $M'$을 정의한 것과 마찬가지로
$C' = C \otimes_A A'$와 $N' = N \otimes_A A'$로 두자.
$M' \cong N'$이 $B'$-가군으로서 성립함에 유의하자. $B \to C$가
유한이라는 가정에는 두 결과가 있다. (a)
$\mathfrak m_C = \sqrt{\mathfrak m_B C}$이고, (b) $B' \to C'$가
유한이다. (a)로부터 다음을 얻는다.
$$
V(\mathfrak m_{A'}C' + \mathfrak m_C C')
=
\left(
\Spec(C') \to \Spec(B')
\right)^{-1}V(\mathfrak m_{A'}B' + \mathfrak m_B B').
$$
% P05-EN-CH38-0206: 점을 소집합으로 두는 원문의 유형 오류를 충실히 보존했다.
$\mathfrak q \subset V(\mathfrak m_{A'}B' + \mathfrak m_B B')$라고
가정하자. $M'_{\mathfrak q}$가 $A'$ 위에서 평탄일 필요충분조건은
$C'_{\mathfrak q}$-가군 $N'_{\mathfrak q}$가 $A'$ 위에서 평탄인
것이고(둘은 $A'$-가군으로서 동형이기 때문이다), 다시 이것의
필요충분조건은 $C'_{\mathfrak q}$의 모든 극대 아이디얼
$\mathfrak r$에 대하여 가군 $N'_{\mathfrak r}$가 $A'$ 위에서
평탄인 것이다(Algebra, Lemma
\ref{algebra-lemma-flat-localization}를 보라). (b)에 의해
$B'_{\mathfrak q} \to C'_{\mathfrak q}$가 유한이므로,
$C'_{\mathfrak q}$의 극대 아이디얼들은 $\mathfrak q$ 위에 놓이는
$C'$의 소아이디얼들과 정확히 대응한다(Algebra, Lemma
\ref{algebra-lemma-integral-going-up}을 보라). 위의 표시식에 의해 이
소아이디얼들은 모두 $V(\mathfrak m_{A'}C' + \mathfrak m_C C')$에
포함된다. 따라서 보조정리의 결론이 성립한다.
\end{proof}

\begin{lemma}
\label{flat-lemma-flat-at-point-go-up}
Situation \ref{flat-situation-flat-at-point}에서 $B \to C$가 국소환들의
평탄 국소 준동형이라고 가정하자. $N = M \otimes_B C$로 두고, 쌍
$(C, N)$에 결부된 함자를
$F'_{lf} : \mathcal{C} \to \textit{Sets}$라 하자. 그러면
$F_{lf} = F'_{lf}$이다.
\end{lemma}

\begin{proof}
$A'$를 $\mathcal{C}$의 대상이라 하자. Situation
\ref{flat-situation-flat-at-point}에서 $B'$, $M'$을 정의한 것과 마찬가지로
$C' = C \otimes_A A'$와
$N' = N \otimes_A A' = M' \otimes_{B'} C'$로 두자. 다음에 유의하자.
$$
V(\mathfrak m_{A'}B' + \mathfrak m_B B')
=
\Spec( \kappa(\mathfrak m_B) \otimes_A \kappa(\mathfrak m_{A'}) )
$$
그리고 $V(\mathfrak m_{A'}C' + \mathfrak m_C C')$에 대해서도
마찬가지이다. 환 사상
$$
\kappa(\mathfrak m_B) \otimes_A \kappa(\mathfrak m_{A'})
\longrightarrow
\kappa(\mathfrak m_C) \otimes_A \kappa(\mathfrak m_{A'})
$$
은 충실히 평탄이다. 따라서
$V(\mathfrak m_{A'}C' + \mathfrak m_C C') \to
V(\mathfrak m_{A'}B' + \mathfrak m_B B')$는 전사이다. 마지막으로
$\mathfrak r \in V(\mathfrak m_{A'}C' + \mathfrak m_C C')$가
$\mathfrak q \in V(\mathfrak m_{A'}B' + \mathfrak m_B B')$로 간다면,
$B' \to C'$가 평탄이므로 $M'_{\mathfrak q}$가 $A'$ 위에서 평탄일
필요충분조건은 $N'_{\mathfrak r}$가 $A'$ 위에서 평탄인 것이다.
Algebra, Lemma \ref{algebra-lemma-flatness-descends-more-general}을 보라.
이 관찰들로부터 보조정리가 형식적으로 따른다.
\end{proof}

\begin{situation}
\label{flat-situation-free-at-generic-points}
$f : X \to S$를 기하적으로 기약인 올을 갖는 매끄러운 사상이라 하자.
$\mathcal{F}$를 유한형 준연접 $\mathcal{O}_X$-가군이라 하자. $S$ 위의
임의의 스킴 $T$에 대하여 $\mathcal{F}$의 $T$로의 기저변환을
$\mathcal{F}_T$로 나타내겠다. 다시 말해 $\mathcal{F}_T$는 사영 사상
$X_T = X \times_S T \to X$에 의한 $\mathcal{F}$의 당김이다.
$X_T \to T$는 기하적으로 기약인 올을 갖는 매끄러운 사상이다.
Morphisms, Lemma \ref{morphisms-lemma-base-change-smooth}와 More on
Morphisms, Lemma
\ref{more-morphisms-lemma-base-change-fibres-geometrically-irreducible}를
보라. $p \geq 0$을 정수라 하자. 점 $t \in T$가 주어졌을 때 다음
조건을 생각하자.
\begin{equation}
\label{flat-equation-free-at-generic-point-fibre}
\mathcal{F}_T\text{는 }\xi_t\text{의 한 근방에서 계수 }p\text{인 자유 가군이다}
\end{equation}
여기서 $\xi_t$는 올 $X_t$의 일반점이다. 모든 $t \in T$에 대한 이
조건은 기저변환 아래에서 안정적이므로 다음 함자를 얻는다.
\begin{equation}
\label{flat-equation-free-at-generic-points}
H_p : (\Sch/S)^{opp} \longrightarrow \textit{Sets}, \quad
T \longrightarrow \left\{
\begin{matrix}
\{*\} & \mathcal{F}_T\text{가 모든 }\forall t\in T\text{에 대하여
(\ref{flat-equation-free-at-generic-point-fibre})를 만족하면}, \\
\emptyset & \text{그렇지 않으면.}
\end{matrix}
\right.
\end{equation}
\end{situation}

\begin{lemma}
\label{flat-lemma-free-at-generic-points}
Situation \ref{flat-situation-free-at-generic-points}의 상황이라 하자.
\begin{enumerate}
\item 함자 $H_p$는 fpqc 위상에 대한 층 성질을 만족한다.
% P05-EN-CH38-0207: H_p 앞의 관사가 빠진 원문을 한국어로 자연스럽게 나타냈다.
\item $\mathcal{F}$가 유한표현이면 함자 $H_p$는 극한 보존이다.
\end{enumerate}
\end{lemma}

\begin{proof}
$\{T_i \to T\}_{i \in I}$를 $S$ 위의 스킴들의 fpqc\footnote{유한표현
평탄 사상이 열린 사상이라는 사실을 쓰면 $H_p$가 fppf 위상에 대한
층임을 보이기는 매우 쉽다. 뒤에서 실제로 필요한 것은 이것뿐이다.
그러나 fpqc 위상에 대한 층 조건까지 만족한다는 것을 직접 증명하는
것도 제법 흥미롭다.} 덮개라 하자. $X_i = X_{T_i} = X \times_S T_i$로
두고, $X_i$로의 $\mathcal{F}$의 당김을 $\mathcal{F}_i$로 나타내자.
모든 $i$에 대하여 $\mathcal{F}_i$가
(\ref{flat-equation-free-at-generic-point-fibre})를 만족한다고 가정하자.
$t \in T$를 택하고 $\xi_t \in X_T$를 $X_t$의 일반점이라 하자.
% P05-EN-CH38-0208: 여기서 F_T 대신 F를 말하는 원문을 그대로 보존했다.
$\mathcal{F}$가 $\xi_t$의 한 근방에서 자유임을 보여야 한다.
% P05-EN-CH38-0209: 부정관사가 빠진 원문을 한국어로 자연스럽게 나타냈다.
어떤 $i \in I$에 대하여 $t$로 가는 $t_i \in T_i$를 찾을 수 있다.
$\xi_i \in X_i$를 $X_{t_i}$의 일반점이라 하자. 그러면 $\xi_i$는
$\xi_t$로 간다. $\mathcal{F}_i$가 $\xi_i$의 한 근방에서 계수 $p$인
자유 가군이라는 사실로부터 다음을 얻는다.
% P05-EN-CH38-0210: 정의된 xi_i 대신 x_i를 쓰는 두 원문 줄기 표기를 그대로 보존했다.
$(\mathcal{F}_i)_{x_i} \cong \mathcal{O}_{X_i, x_i}^{\oplus p}$
이는 다음을 함의한다.
$\mathcal{F}_{T, \xi_t} \cong \mathcal{O}_{X_T, \xi_t}^{\oplus p}$
$\mathcal{O}_{X_T, \xi_t} \to \mathcal{O}_{X_i, x_i}$가 평탄이기
때문이다. 예를 들어 Algebra, Lemma
\ref{algebra-lemma-finite-projective-descends}를 보라. 따라서 $X_T$에서
$\xi_t$의 아핀 근방 $U$와 전사 사상
$\mathcal{O}_U^{\oplus p} \to \mathcal{F}_U = \mathcal{F}_T|_U$
이 존재한다. Modules, Lemma
\ref{modules-lemma-finite-type-surjective-on-stalk}를 보라. $T$를
줄인 뒤 $U \to T$가 전사라고 가정해도 된다. 따라서 $U \to T$는
기하적으로 기약인 올을 갖는 아핀 스킴들의 매끄러운 사상이다. 더욱이
모든 $t' \in T$에 대하여 유도된 사상
$$
\alpha :
\mathcal{O}_{U, \xi_{t'}}^{\oplus p}
\longrightarrow
\mathcal{F}_{U, \xi_{t'}}
$$
은 동형이다(앞과 같은 논증으로 오른쪽 가군이 계수 $p$인 자유 가군이기
때문이다). Lemma \ref{flat-lemma-induction-step}로부터
$$
\Gamma(U, \mathcal{O}_U^{\oplus p})
\otimes_{\Gamma(T, \mathcal{O}_T)} \mathcal{O}_{T, t'}
\longrightarrow
\Gamma(U, \mathcal{F}_U)
\otimes_{\Gamma(T, \mathcal{O}_T)} \mathcal{O}_{T, t'}
$$
가 모든 $t' \in T$에 대하여 단사임을 얻는다. 따라서 전사 사상
$\alpha$는 동형이다. 이로써 (1)의 증명이 끝난다.

\medskip\noindent
$\mathcal{F}$가 유한표현이라고 가정하자. $T = \lim_{i \in I} T_i$를
아핀 $S$-스킴들의 유향 극한이라 하고 $\mathcal{F}_T$가
(\ref{flat-equation-free-at-generic-point-fibre})를 만족한다고 가정하자.
$X_i = X_{T_i} = X \times_S T_i$로 두고 $X_i$로의 $\mathcal{F}$의
당김을 $\mathcal{F}_i$로 나타내자. $U \subset X_T$를
$\mathcal{F}_T$가 $T$ 위에서 평탄인 점들로 이루어진 열린부분스킴이라
하자. More on Morphisms, Theorem
\ref{more-morphisms-theorem-openness-flatness}를 보라. 가정에 의해 모든
올의 모든 일반점은 $U$의 점이다. 즉 $U \to T$는 기하적으로 기약인
올을 갖는 매끄러운 전사 사상이다. $U$를 조금 줄여 $U$가
준콤팩트라고 가정해도 된다. Limits, Lemma
\ref{limits-lemma-descend-opens}를 사용하면, $X_T$에서의 역상이 $U$인
어떤 $i \in I$와 준콤팩트 열린집합 $U_i \subset X_i$를 찾을 수 있다.
$i$를 증가시킨 뒤 $\mathcal{F}_i|_{U_i}$가 $T_i$ 위에서 평탄이라고
가정해도 된다. Limits, Lemma
\ref{limits-lemma-descend-module-flat-finite-presentation}를 보라. 특히
$\mathcal{F}_i|_{U_i}$는 유한 국소 자유이므로 국소 상수 계수 함수
$\rho : U_i \to \{0, 1, 2, \ldots \}$를 정의한다. $(U_i)_p \subset U_i$를
$\rho$의 값이 $p$인 열린닫힌 부분집합이라 하자. $V_i \subset T_i$를
$(U_i)_p$의 상이라 하자. $V_i$는 열린집합이고 준콤팩트임에 유의하자.
가정에 의해 $T \to T_i$의 상은 $V_i$에 포함된다. 따라서 Limits,
Lemma \ref{limits-lemma-descend-opens}에 의해 $T_{i'} \to T_i$가
$V_i$를 통해 인수분해되는 어떤 $i' \geq i$가 존재한다. 그러면
$\mathcal{F}_{i'}$는 원하는 대로
(\ref{flat-equation-free-at-generic-point-fibre})를 만족한다. 몇몇 세부
사항은 생략한다.
\end{proof}

\begin{lemma}
\label{flat-lemma-pre-flat-dimension-n}
$f : X \to S$를 국소 유한형인 스킴 사상이라 하자. $\mathcal{F}$를
유한형 준연접 $\mathcal{O}_X$-가군이라 하자. $n \geq 0$이라 하자.
다음 조건들은 동치이다.
\begin{enumerate}
\item $s \in S$에 대하여 $\mathcal{F}$가 $S$ 위에서 평탄이지
않은 점들로 이루어진 닫힌 부분집합 $Z \subset X_s$(Lemma
\ref{flat-lemma-open-in-fibre-where-flat}을 보라)는 $\dim(Z) < n$을 만족하고,
\item $\mathcal{F}$가 $S$ 위의 $x$에서 평탄이지 않은 모든 $x \in X$에
대하여 $\text{trdeg}_{\kappa(f(x))}(\kappa(x)) < n$이다.
\end{enumerate}
이 조건들이 성립하면 임의의 기저변환 뒤에도 성립한다.
\end{lemma}

\begin{proof}
$x \in X$를 $s \in S$ 위의 점이라 하자. Varieties, Lemma
\ref{varieties-lemma-dimension-locally-algebraic}에 의해 $X_s$에서
$\{x\}$의 폐포의 차원은 $\text{trdeg}_{\kappa(s)}(\kappa(x))$이다.
반대로 $Z \subset X_s$가 차원 $d$인 닫힌 부분집합이면
$\text{trdeg}_{\kappa(s)}(\kappa(x)) = d$인 점 $x \in Z$가 존재한다
(같은 참조를 보라). 따라서 (1)과 (2)는 동치이다(각 올마다도
그러하다). 기저변환에 관한 명제는 Morphisms, Lemmas
\ref{morphisms-lemma-base-change-module-flat}와
\ref{morphisms-lemma-dimension-fibre-after-base-change}에서 따른다.
\end{proof}

\begin{definition}
\label{flat-definition-flat-dimension-n}
$f : X \to S$를 국소 유한형인 스킴 사상이라 하자. $\mathcal{F}$를
유한형 준연접 $\mathcal{O}_X$-가군이라 하자. $n \geq 0$이라 하자.
Lemma \ref{flat-lemma-pre-flat-dimension-n}의 동치인 조건들이 성립하면
{\it $\mathcal{F}$가 차원 $\geq n$에서 $S$ 위에 평탄하다}고 한다.
\end{definition}

\begin{situation}
\label{flat-situation-flat-dimension-n}
$f : X \to S$를 국소 유한형인 스킴 사상이라 하자. $\mathcal{F}$를
유한형 준연접 $\mathcal{O}_X$-가군이라 하자. $S$ 위의 임의의 스킴
$T$에 대하여 $\mathcal{F}$의 $T$로의 기저변환을 $\mathcal{F}_T$로
나타내겠다. 다시 말해 $\mathcal{F}_T$는 사영 사상
$X_T = X \times_S T \to X$에 의한 $\mathcal{F}$의 당김이다.
% P05-EN-CH38-0211: 국소 유한형 가정에서 기저변환을 유한형이라 하는 원문을 그대로 보존했다.
$X_T \to T$는 유한형이고 $\mathcal{F}_T$는 유한형
$\mathcal{O}_{X_T}$-가군이다(Morphisms, Lemma
\ref{morphisms-lemma-base-change-finite-type}와 Modules, Lemma
\ref{modules-lemma-pullback-finite-type}를 보라). $n \geq 0$이라 하자.
Definition \ref{flat-definition-flat-dimension-n}과 Lemma
\ref{flat-lemma-pre-flat-dimension-n}에 의해 다음 함자를 얻는다.
\begin{equation}
\label{flat-equation-flat-dimension-n}
F_n : (\Sch/S)^{opp} \longrightarrow \textit{Sets}, \quad
T \longrightarrow \left\{
\begin{matrix}
\{*\} & \mathcal{F}_T\text{가 }T\text{ 위에서 차원 }\dim \geq n\text{에 평탄하면}, \\
\emptyset & \text{그렇지 않으면.}
\end{matrix}
\right.
\end{equation}
\end{situation}

\begin{lemma}
\label{flat-lemma-flat-dimension-n}
Situation \ref{flat-situation-flat-dimension-n}의 상황이라 하자.
\begin{enumerate}
\item 함자 $F_n$은 fpqc 위상에 대한 층 성질을 만족한다.
\item $f$가 준콤팩트이고 국소 유한표현이며 $\mathcal{F}$가
유한표현이면 함자 $F_n$은 극한 보존이다.
\end{enumerate}
\end{lemma}

\begin{proof}
$\{T_i \to T\}_{i \in I}$를 $S$ 위의 스킴들의 fpqc 덮개라 하자.
$X_i = X_{T_i} = X \times_S T_i$로 두고 $X_i$로의 $\mathcal{F}$의
당김을 $\mathcal{F}_i$로 나타내자. 모든 $i$에 대하여
$\mathcal{F}_i$가 차원 $\geq n$에서 $T_i$ 위에 평탄이라고 가정하자.
$t \in T$라 하자. 지표 $i$와 $t$로 가는 점 $t_i \in T_i$를 택하자.
다음 카르테시안 그림을 생각하자.
$$
\xymatrix{
X_{\Spec(\mathcal{O}_{T, t})} \ar[d] &
X_{\Spec(\mathcal{O}_{T_i, t_i})} \ar[d] \ar[l] \\
\Spec(\mathcal{O}_{T, t}) &
\Spec(\mathcal{O}_{T_i, t_i}) \ar[l]
}
$$
아래 가로 사상이 평탄이므로 More on Morphisms, Lemma
\ref{more-morphisms-lemma-flat-locus-base-change}에 의해,
$\mathcal{F}_i$가 $T_i$ 위에서 평탄이지 않은 점들의 집합
$Z_i \subset X_{t_i}$와 $\mathcal{F}_T$가 $T$ 위에서 평탄이지 않은
점들의 집합 $Z \subset X_t$는 $Z_i = Z_{\kappa(t_i)}$라는 관계를
만족한다. 따라서 Morphisms, Lemma
\ref{morphisms-lemma-dimension-fibre-after-base-change}에 의해
$\mathcal{F}_T$는 차원 $\geq n$에서 $T$ 위에 평탄이다.

\medskip\noindent
$f$가 준콤팩트이고 국소 유한표현이며 $\mathcal{F}$가 유한표현이라고
가정하자. 이 문단에서는 먼저 (2)의 증명을 $f$가 유한표현인 경우로
귀착시킨다. $T = \lim_{i \in I} T_i$를 아핀 $S$-스킴들의 유향
극한이라 하고 $\mathcal{F}_T$가 차원 $\geq n$에서 평탄이라고
가정하자. $X_i = X_{T_i} = X \times_S T_i$로 두고 $X_i$로의
$\mathcal{F}$의 당김을 $\mathcal{F}_i$로 나타내자. 어떤 $i$에
대하여 $\mathcal{F}_i$가 차원 $\geq n$에서 평탄임을 보여야 한다.
$i_0 \in I$를 택하고 $I$를 $\{i \mid i \geq i_0\}$로 바꾸자.
$T_{i_0}$가 아핀이므로(따라서 준콤팩트이므로) 유한 개의 아핀
열린집합 $W_j \subset S$, $j = 1, \ldots, m$과 아핀 열린 덮개
% P05-EN-CH38-0212: covering의 철자가 손상된 원문을 한국어로 자연스럽게 나타냈다.
$T_{i_0} = \bigcup_{j = 1, \ldots, m} V_{j, i_0}$가 존재하여
$T_{i_0} \to S$는 $V_{j, i_0}$를 $W_j$ 안으로 보낸다.
$i \geq i_0$에 대하여 $T_i$에서의 $V_{j, i_0}$의 역상을
$V_{j, i}$로 나타내자. 각 $j$마다 어떤 $i$가 존재하여
% P05-EN-CH38-0213: 변하는 i 대신 고정된 i_0를 쓰는 원문 첨자를 그대로 보존했다.
$\mathcal{F}_{V_{j, i_0}}$가 차원 $\geq n$에서 평탄임을 보일 수
있다면 결론을 얻는다. 이와 같이 $S$가 아핀인 경우로 귀착된다.
이 경우 $X$는 준콤팩트이고 유한 아핀 열린 덮개
$X = W_1 \cup \ldots \cup W_m$를 택할 수 있다. 이때
$(X \to S, \mathcal{F})$에 대한 결과는
$(\coprod W_j, \coprod \mathcal{F}|_{W_j})$에 대한 결과와 동치이다.
따라서 $f$가 유한표현이라고 가정해도 된다.

\medskip\noindent
$f$와 $\mathcal{F}$가 유한표현이라고 가정하자. $U \subset X_T$를
$\mathcal{F}_T$가 $T$ 위에서 평탄인 점들로 이루어진 열린부분스킴이라
하자. More on Morphisms, Theorem
\ref{more-morphisms-theorem-openness-flatness}를 보라. 가정에 의해
$Z = X_T \setminus U$의 $T$ 위의 모든 올은 차원 $< n$이다. Limits,
Lemma \ref{limits-lemma-approximate-given-relative-dimension}에 의해, 모든
$t \in T$에 대하여 $\dim(Z'_t) < n$이고 $Z' \to X_T$가 유한표현인
닫힌부분스킴 $Z \subset Z' \subset X_T$를 찾을 수 있다. Limits,
Lemmas \ref{limits-lemma-descend-finite-presentation}와
\ref{limits-lemma-descend-closed-immersion-finite-presentation}에 의해,
$T$로의 기저변환이 $Z'$인 어떤 $i \in I$와 유한표현 닫힌부분스킴
$Z'_i \subset X_i$가 존재한다. Limits, Lemma
\ref{limits-lemma-limit-dimension}에 의해 $Z'_i \to T_i$의 모든 올의
차원이 $< n$이라고 가정해도 된다. Limits, Lemma
\ref{limits-lemma-descend-module-flat-finite-presentation}에 의해
% P05-EN-CH38-0214: 정의된 Z'_i 대신 T'_i를 쓰는 원문 여집합을 그대로 보존했다.
$\mathcal{F}_i|_{X_i \setminus T'_i}$가 $T_i$ 위에서 평탄이라고
가정해도 된다. 이는 $\mathcal{F}_i$가 차원 $\geq n$에서 평탄임을
함의한다. 여기서는 $Z' \to X_T$가 유한표현이고 따라서 여집합
$X_T \setminus Z'$가 준콤팩트라는 사실을 사용한다! 이로써 (2)가
증명되었고 보조정리의 증명이 끝난다.
\end{proof}

\begin{situation}
\label{flat-situation-flat}
$f : X \to S$를 스킴 사상이라 하고 $\mathcal{F}$를 준연접
$\mathcal{O}_X$-가군이라 하자. $S$ 위의 임의의 스킴 $T$에 대하여
$\mathcal{F}$의 $T$로의 기저변환을 $\mathcal{F}_T$로 나타내겠다.
다시 말해 $\mathcal{F}_T$는 사영 사상 $X_T = X \times_S T \to X$에
의한 $\mathcal{F}$의 당김이다. 평탄 가군의 기저변환은 평탄이므로
다음 함자를 얻는다.
\begin{equation}
\label{flat-equation-flat}
F_{flat} : (\Sch/S)^{opp} \longrightarrow \textit{Sets}, \quad
T \longrightarrow \left\{
\begin{matrix}
\{*\} & \mathcal{F}_T \text{가 }T\text{ 위에서 평탄이면}, \\
\emptyset & \text{그렇지 않으면.}
\end{matrix}
\right.
\end{equation}
\end{situation}

\begin{lemma}
\label{flat-lemma-flat}
Situation \ref{flat-situation-flat}의 상황이라 하자.
\begin{enumerate}
\item 함자 $F_{flat}$는 fpqc 위상에 대한 층 성질을 만족한다.
\item $f$가 준콤팩트이고 국소 유한표현이며 $\mathcal{F}$가
유한표현이면 함자 $F_{flat}$는 극한 보존이다.
\end{enumerate}
\end{lemma}

\begin{proof}
(1)은 다음 명제에서 따른다. $T' \to T$가 $S$ 위의 스킴들의 평탄
전사 사상이면 $\mathcal{F}_{T'}$가 $T'$ 위에서 평탄일 필요충분조건은
$\mathcal{F}_T$가 $T$ 위에서 평탄인 것이다. More on Morphisms, Lemma
\ref{more-morphisms-lemma-flat-locus-base-change}를 보라. (2)는 $X$와
$S$가 아핀인 경우로 귀착시킨 뒤 Limits, Lemma
\ref{limits-lemma-descend-module-flat-finite-presentation}에서 따른다
(Lemma \ref{flat-lemma-flat-dimension-n}의 증명과 비교하라).
\end{proof}




\section{평탄화 층화}
\label{flat-section-flattening}

\noindent
정의만 제시한다.
% P05-EN-CH38-0215: 주어와 술어를 부당하게 가르는 원문의 쉼표를 한국어에서 제거했다.
``일반 평탄성 층화''를 찾는 독자는 More on Morphisms, Section
\ref{more-morphisms-section-generic-flatness-stratification}를 참고하라.

\begin{definition}
\label{flat-definition-flattening}
$X \to S$를 스킴 사상이라 하고 $\mathcal{F}$를 준연접
$\mathcal{O}_X$-가군이라 하자. Situation \ref{flat-situation-flat}에서 정의한
함자 $F_{flat}$가 $S$ 위의 스킴 $S'$에 의해 표현가능이면
{\it $\mathcal{F}$의 보편 평탄화가 존재한다}고 한다.
$\mathcal{O}_X$의 보편 평탄화가 존재하면 {\it $X$의 보편 평탄화가
존재한다}고 한다.
\end{definition}

\noindent
$\mathcal{F}$의 보편 평탄화 $S'$\footnote{스킴 $S'$를 때때로
{\it 보편 평탄화자}라고 한다. \cite{GruRay}에서는
{\it platificateur universel}이라고 한다. 보편 평탄화의 존재를 More on
Algebra, Section \ref{more-algebra-section-blowup-flat}에서 논의한 유형의
결과와 혼동해서는 안 된다.}가 존재하면, 사상 $S' \to S$는 스킴들의
전사 단사 사상이고 $\mathcal{F}_{S'}$는 $S'$ 위에서 평탄이며, 사상
$T \to S$가 $S'$를 통해 인수분해될 필요충분조건은 $\mathcal{F}_T$가
$T$ 위에서 평탄인 것이다.

\begin{example}
\label{flat-example-no-universal-flattening}
$k$를 체라 하고 $X = S = \Spec(k[x, y])$라 하자.
$M = k[x, x^{-1}, y]/(y)$인 $\mathcal{F} = \widetilde{M}$이라 하자.
$k[x, y]$-대수 $A$에 대하여
$F_{flat}(A) = F_{flat}(\Spec(A))$로 두자. 모든 $n$에 대하여
$F_{flat}(k[x, y]/(x, y)^n) = \{*\}$이지만
$F_{flat}(k[[x, y]]) = \emptyset$이다. 이는 $F_{flat}$가 표현가능하지
않음을 뜻한다(대수공간에 의해서조차 그러하다. Formal Spaces, Lemma
\ref{formal-spaces-lemma-map-into-algebraic-space}를 보라). 따라서 이
경우에는 보편 평탄화가 존재하지 않는다.
\end{example}

\noindent
스킴 $S$의 (국소 유한이고 스킴론적인) {\it 층화}를 다음 자료로
정의한다(Topology, Remark
\ref{topology-remark-locally-finite-stratification}와 비교하라). 부분순서
집합 $I$로 지표화된 닫힌부분스킴 $Z_i \subset S$들이 주어져서
$S = \bigcup Z_i$가 집합론적으로 성립하고, $S$의 모든 점은 유한 개의
$Z_i$와만 만나는 근방을 가지며, 다음이 성립한다.
$$
Z_i \cap Z_j = \bigcup\nolimits_{k \leq i, j} Z_k.
$$
$S_i = Z_i \setminus \bigcup_{j < i} Z_j$로 두면 실제 층화는 국소 닫힌
부분스킴들로의 분해 $S = \coprod S_i$이다. 흔히 층 $S_i$만 제시하고
닫힌부분스킴 $Z_i$의 구성은 독자에게 맡긴다. 층화가 주어지면 단사 사상
$$
S' = \coprod\nolimits_{i \in I} S_i \longrightarrow S.
$$
을 얻는다. 이를 {\it 층화에 결부된 단사 사상}이라 부르겠다. 이
용어를 써서 평탄화 층화를 갖는다는 말의 뜻을 정의할 수 있다.

\begin{definition}
\label{flat-definition-flattening-stratification}
$X \to S$를 스킴 사상이라 하고 $\mathcal{F}$를 준연접
$\mathcal{O}_X$-가군이라 하자. Situation \ref{flat-situation-flat}에서 정의한
함자 $F_{flat}$가 국소 닫힌 부분스킴들에 의한 $S$의 층화에 결부된
단사 사상 $S' \to S$에 의해 표현가능이면 $\mathcal{F}$가
{\it 평탄화 층화}를 갖는다고 한다. $\mathcal{O}_X$가 평탄화 층화를
가지면 $X$가 {\it 평탄화 층화}를 갖는다고 한다.
\end{definition}

\noindent
평탄화 층화가 존재할 때는 층들을 지표화하는 집합과 그 부분순서를
이해하는 일이 흔히 중요하다. 이는 흔히 가군의 계수와 관련된다. 예를
들어 $X = S$이고 $\mathcal{F}$가 유한표현 $\mathcal{O}_S$-가군이면
평탄화 층화가 존재하며 $\mathcal{F}$의 피팅 아이디얼들에 의해
주어진다. Divisors, Lemma
\ref{divisors-lemma-finite-presentation-module}를 보라.




\section{아르틴 환 위의 평탄화 층화}
\label{flat-section-flattening-artinian}

\noindent
% P05-EN-CH38-0216: flattening의 철자가 손상된 원문을 의도대로 나타냈다.
기저 스킴이 아르틴 환의 스펙트럼이면 평탄화 층화가 존재한다.

\begin{lemma}
\label{flat-lemma-flattening-stratification-artinian}
$S$를 아르틴 환의 스펙트럼이라 하자.
% P05-EN-CH38-0217: 준연접 가군에 F라는 기호를 붙이지 않는 원문의 양화를 그대로 보존했다.
$S$ 위의 임의의 스킴 $X$와 임의의 준연접 $\mathcal{O}_X$-가군에
대하여 보편 평탄화가 존재한다. 실제로 보편 평탄화는 닫힌 몰입
$S' \to S$에 의해 주어지므로 $\mathcal{F}$에 대한 평탄화 층화이기도
하다.
\end{lemma}

\begin{proof}
아핀 열린 덮개 $X = \bigcup U_i$를 택하자. 그러면 $F_{flat}$는 각 쌍
$(U_i, \mathcal{F}|_{U_i})$에 결부된 함자들의 곱이다. 따라서 각
$(U_i, \mathcal{F}|_{U_i})$에 대하여 결과를 증명하면 충분하다. 아핀인
경우 보조정리는 More on Algebra, Lemma
\ref{more-algebra-lemma-flattening-artinian-universal-property}에서 즉시
따른다.
\end{proof}






\section{사상의 평탄화}
\label{flat-section-flattening-map}

\noindent
Theorem \ref{flat-theorem-flattening-map}은 뒤의 평탄화 명제들의 핵심이다.

\begin{lemma}
\label{flat-lemma-universally-separating}
$S$를 스킴이라 하자. $g : X' \to X$를 $S$ 위의 스킴들의 평탄
사상이라 하고 $X$가 $S$ 위에서 국소 유한형이라고 하자.
$\mathcal{F}$를 $S$ 위에서 평탄인 유한형 준연접
$\mathcal{O}_X$-가군이라 하자. 만일
$\text{Ass}_{X/S}(\mathcal{F}) \subset g(X')$이면 표준 사상
$$
\mathcal{F} \longrightarrow g_*g^*\mathcal{F}
$$
은 단사이고 임의의 기저변환 뒤에도 단사로 남는다.
\end{lemma}

\begin{proof}
마지막 주장은 임의의 사상 $T \to S$에 대하여
$\mathcal{F}_T \to (g_T)_*g_T^*\mathcal{F}_T$가 단사라는 뜻이다.
가정 $\text{Ass}_{X/S}(\mathcal{F}) \subset g(X')$는 기저변환 아래에서
보존된다. Divisors, Lemma
\ref{divisors-lemma-base-change-relative-assassin}과 Remark
\ref{divisors-remark-base-change-relative-assassin}을 보라. 평탄성과
유한형이라는 가정도 마찬가지이다. 따라서 표시된 화살표가 단사임을
증명하면 충분하다. $\mathcal{K} = \Ker(\mathcal{F} \to
g_*g^*\mathcal{F})$라 하자. $\mathcal{K} = 0$임을 증명하는 것이
목표이다. 이를 위해서는
$\text{WeakAss}_X(\mathcal{K}) = \emptyset$임을 증명하면 충분하다.
Divisors, Lemma \ref{divisors-lemma-weakly-ass-zero}를 보라. 또한
$\text{WeakAss}_X(\mathcal{K}) \subset \text{WeakAss}_X(\mathcal{F})$이다.
Divisors, Lemma \ref{divisors-lemma-ses-weakly-ass}를 보라.
$\mathcal{F}$가 평탄이므로 Lemma
\ref{flat-lemma-bourbaki-finite-type-general-base}에 의해
$\text{WeakAss}_X(\mathcal{F}) \subset \text{Ass}_{X/S}(\mathcal{F})$이다.
가정에 의해 $\text{Ass}_{X/S}(\mathcal{F})$의 임의의 점 $x$는 어떤
$x' \in X'$의 상이다. $g$가 평탄이므로 국소환 사상
$\mathcal{O}_{X, x} \to \mathcal{O}_{X', x'}$는 충실히 평탄이고,
따라서 사상
$$
\mathcal{F}_x
\longrightarrow
g^*\mathcal{F}_{x'} =
\mathcal{F}_x \otimes_{\mathcal{O}_{X, x}} \mathcal{O}_{X', x'}
$$
은 단사이다(Algebra, Lemma
\ref{algebra-lemma-faithfully-flat-universally-injective}를 보라). 이는
원하는 대로 $\mathcal{K}_x = 0$을 함의한다.
\end{proof}

\begin{lemma}
\label{flat-lemma-flattening-module-map}
$A$를 환이라 하고 $u : M \to N$을 $A$-가군들의 전사 사상이라 하자.
$M$이 $A$-가군으로서 사영이면, 임의의 환 사상
$\varphi : A \to B$에 대하여 다음 조건들이 동치가 되게 하는 아이디얼
$I \subset A$가 존재한다.
\begin{enumerate}
\item $u \otimes 1 : M \otimes_A B \to N \otimes_A B$는 동형이고,
\item $\varphi(I) = 0$.
\end{enumerate}
\end{lemma}

\begin{proof}
$M$이 사영이므로 $F = M \oplus C$가 자유 $A$-가군이 되게 하는 사영
$A$-가군 $C$를 찾을 수 있다. $u$를
$u \oplus 1 : F = M \oplus C \to N \oplus C$로 바꾸면 $M$이
자유라고 가정해도 됨을 알 수 있다. 이 경우 $M$의 어떤 고정된 기저에
관한 $\Ker(u)$의 모든 원소의 계수들로 생성되는 $A$의 아이디얼을
$I$라 하자. 이것이 성립하는 까닭은 $u$가 전사이고
$\otimes_A B$가 오른쪽 완전이므로 $\Ker(u \otimes 1)$이
$M \otimes_A B$ 안에서 $\Ker(u) \otimes_A B$의 상이기 때문이다.
\end{proof}

\begin{theorem}
\label{flat-theorem-flattening-map}
Situation \ref{flat-situation-iso}에서 다음을 가정하자.
\begin{enumerate}
\item $f$는 유한표현이고,
\item $\mathcal{F}$는 유한표현이며 $S$ 위에서 평탄이고 $S$에
상대적으로 순수하며,
\item $u$는 전사이다.
\end{enumerate}
그러면 $F_{iso}$는 닫힌 몰입 $Z \to S$에 의해 표현가능하다. 더욱이
$\mathcal{G}$가 유한표현이면 $Z \to S$도 유한표현이다.
\end{theorem}

\begin{proof}
$\mathcal{F}$가 $S$ 위에서 보편적으로 순수라는 사실을 더 언급하지
않고 사용하겠다. Lemma
\ref{flat-lemma-finite-type-flat-pure-along-fibre-is-universal}을 보라. Lemma
\ref{flat-lemma-iso-sheaf}와 Descent, Lemmas
\ref{descent-lemma-closed-immersion},
\ref{descent-lemma-descent-data-sheaves}에 의해 이 문제는 $S$ 위의 에탈
위상에 대하여 국소적이다. 따라서 $s \in S$가 주어졌을 때 정리가
성립하는 $(S, s)$의 에탈 근방이 존재함을 증명하면 충분하다.

\medskip\noindent
Lemma \ref{flat-lemma-finite-presentation-flat-along-fibre}를 사용하고 $S$를
$s$의 기본 에탈 근방으로 바꾼 뒤, 다음 가환 그림이 존재한다고
가정해도 된다.
$$
\xymatrix{
X \ar[dr] & & X' \ar[ll]^g \ar[ld] \\
& S &
}
$$
이는 $S$ 위의 유한표현 스킴들로 이루어지고, 여기서 $g$는 에탈이며
$X_s \subset g(X')$이고 스킴 $X'$와 $S$는 아핀이며,
% P05-EN-CH38-0218: 마지막 가군 서술에 동사가 빠진 원문을 한국어로 자연스럽게 나타냈다.
$\Gamma(X', g^*\mathcal{F})$는 사영
$\Gamma(S, \mathcal{O}_S)$-가군이다. $g^*\mathcal{F}$는 $S$ 위에서
보편적으로 순수임에 유의하자. Lemma
\ref{flat-lemma-affine-locally-projective-pure}를 보라. 따라서 Lemma
\ref{flat-lemma-criterion}에 의해 열린집합 $g(X')$는
$\Spec(\mathcal{O}_{S, s})$ 위에 놓이는
$\text{Ass}_{X/S}(\mathcal{F})$의 점들을 포함한다. 다음과 같이 두자.
$$
E = \{t \in S \mid \text{Ass}_{X_t}(\mathcal{F}_t) \subset g(X') \}.
$$
More on Morphisms, Lemma
\ref{more-morphisms-lemma-relative-assassin-constructible}에 의해 $E$는
$S$의 구성가능 부분집합이다. 앞에서
$\Spec(\mathcal{O}_{S, s}) \subset E$임을 보았다. Morphisms, Lemma
\ref{morphisms-lemma-constructible-containing-open}에 의해 $E$는 $s$의
열린 근방을 포함한다. 따라서 $S$를 $s$의 더 작은 아핀 근방으로
바꾼 뒤 $\text{Ass}_{X/S}(\mathcal{F}) \subset g(X')$라고 가정해도
된다.

\medskip\noindent
$u$가 전사라고 가정했으므로 $F_{iso} = F_{inj}$이다. Lemma
\ref{flat-lemma-universally-separating}에 의해
$u : \mathcal{F} \to \mathcal{G}$가 단사일 필요충분조건은
$g^*u : g^*\mathcal{F} \to g^*\mathcal{G}$가 단사인 것이며, 이는 임의의
기저변환 뒤에도 성립한다. 따라서 정리의 가정들에 더하여
$X \to S$가 아핀 스킴들의 사상이고 $\Gamma(X, \mathcal{F})$가 사영
$\Gamma(S, \mathcal{O}_S)$-가군인 경우로 귀착되었다. 이 경우는
Lemma \ref{flat-lemma-flattening-module-map}에서 즉시 따른다.

\medskip\noindent
$\mathcal{G}$가 유한표현일 때 $Z$가 유한표현임을 보이려면 Lemma
\ref{flat-lemma-iso-sheaf}의 (4)와 Limits, Remark
\ref{limits-remark-limit-preserving}을 함께 사용하면 된다.
\end{proof}

\begin{lemma}
\label{flat-lemma-Weil-restriction-closed-subschemes}
$f:X\to S$를 유한표현이고 평탄이며 순수인 스킴 사상이라 하자.
$Y$를 $X$의 닫힌부분스킴이라 하자. $F=f_*Y$를 $f$를 따른 $Y$의
베유 제한 함자라 하자. 이는 다음과 같이 정의된다.
$$
F : (\Sch/S)^{opp} \to \textit{Sets}, \quad
T \mapsto
\left\{
\begin{matrix}
\{*\} & \text{다음이면} & Y_T\to X_T \text{는 동형이다, }\\
\emptyset & \text{그렇지 않으면.} &
\end{matrix}
\right.
$$
그러면 $F$는 닫힌 몰입 $Z\to S$에 의해 표현가능하다. 더욱이
$Y\to S$가 유한표현이면 $Z\to S$도 유한표현이다.
\end{lemma}

\begin{proof}
$\mathcal{I}$를 $X$에서 $Y$를 정의하는 아이디얼 층이라 하고
$u:\mathcal{O}_X\to\mathcal{O}_X/\mathcal{I}$를 전사 사상이라 하자.
그러면 $S$-스킴 $T$에 대하여 닫힌 몰입 $Y_T\to X_T$가 동형일
필요충분조건은 $u_T$가 동형인 것이다. 따라서 결과는 Theorem
\ref{flat-theorem-flattening-map}에서 따른다.
\end{proof}



\section{국소 경우의 평탄화}
\label{flat-section-flattening-local}

\noindent
이 절에서는 앞의 내용을 적용하여 평탄화 층화의 윤곽을 얻기 시작한다.

\begin{theorem}
\label{flat-theorem-flattening-local}
Situation \ref{flat-situation-flat-at-point}에서 $A$가 헨젤 환이고 $B$가
$A$ 위에서 본질적으로 유한형이며 $M$이 유한 $B$-가군이라고
가정하자. 그러면 $A/I$가 범주 $\mathcal{C}$ 위의 함자 $F_{lf}$를
공표현하게 하는 아이디얼 $I \subset A$가 존재한다. 다시 말해
$B' = B \otimes_A A'$와 $M' = M \otimes_A A'$를 갖는 국소환들의
국소 준동형 $\varphi : A \to A'$가 주어졌을 때 다음 조건들은
동치이다.
\begin{enumerate}
\item $\forall \mathfrak q \in V(\mathfrak m_{A'}B' + \mathfrak m_B B')
\subset \Spec(B') :
M'_{\mathfrak q}\text{는 }A'\text{ 위에서 평탄이다}$이고,
\item $\varphi(I) = 0$.
\end{enumerate}
$B$가 $A$ 위에서 본질적으로 유한표현이고 $M$이 $B$ 위에서
유한표현이면 $I$는 유한생성 아이디얼이다.
\end{theorem}

\begin{proof}
$B = C_{\mathfrak q}$이고 $M = N_{\mathfrak q}$가 되도록 유한형 환 사상
$A \to C$, 유한 $C$-가군 $N$, 그리고 $C$의 소아이디얼
$\mathfrak q$를 택하자.
% P05-EN-CH38-0219: 닫는 따옴표가 빠진 원문의 약속을 한국어에서 완결했다.
이하에서 ``정리가 $(N/C/A, \mathfrak q)$에 대하여 성립한다''고 하면,
$B = C_{\mathfrak q}$이고 $M = N_{\mathfrak q}$인
$(A \to B, M)$에 대하여 성립한다는 뜻이다. Lemma
\ref{flat-lemma-flat-at-point-go-up}에 의해 $B$를 $B$ 위에서 평탄인
국소환으로 바꾸어도 함자 $F_{lf}$는 변하지 않는다. 따라서 $A$가
헨젤 환이므로 Lemma \ref{flat-lemma-existence-algebra}를 적용하여
$\mathfrak q$에서 $N/C/A$의 완전 데비사주가 존재한다고 가정해도 된다.

\medskip\noindent
$(A_i, B_i, M_i, \alpha_i, \mathfrak q_i)_{i = 1, \ldots, n}$을
$\mathfrak q$에서 $N/C/A$의 그러한 완전 데비사주라 하자. Definition
\ref{flat-definition-complete-devissage-at-x-algebra}에서와 같이
$\mathfrak q_i \subset B_i$ 위에 놓이는 유일한 소아이디얼을
$\mathfrak q'_i \subset A_i$라 하자. $C \to A_1$은 전사이고
$N \cong M_1$이 $C$-가군으로서 성립하므로, Lemma
\ref{flat-lemma-flat-at-point-finite}에 의해 정리가
$(M_1/A_1/A, \mathfrak q'_1)$에 대하여 성립함을 증명하면 충분하다.
% P05-EN-CH38-0220: 소아이디얼의 위아래 관계와 국소화 사상의 방향이 뒤집힌 원문을 그대로 보존했다.
$B_1 \to A_1$이 유한이고 $\mathfrak q_1$이 $\mathfrak q'_1$ 위에
놓이는 $B_1$의 유일한 소아이디얼이므로
$(A_1)_{\mathfrak q'_1} \to (B_1)_{\mathfrak q_1}$이 유한이다(Algebra,
Lemma \ref{algebra-lemma-unique-prime-over-localize-below} 또는 More on
Morphisms, Lemma
\ref{more-morphisms-lemma-finite-morphism-single-point-in-fibre}를 보라).
따라서 Lemma \ref{flat-lemma-flat-at-point-finite}에 의해 정리가
$(M_1/B_1/A, \mathfrak q_1)$에 대하여 성립함을 증명하면 충분하다.

\medskip\noindent
이제 데비사주의 길이 $n$에 대한 귀납법으로 정리가
$(M_2/B_2/A, \mathfrak q_2)$에 대하여 성립한다고 가정해도 된다.
($n = 1$이면 $M_2 = 0$이고 이는 $A$ 위에서 평탄이다.) 앞 문단의
마지막 두 단계를 거꾸로 적용하고,
% P05-EN-CH38-0221: 첫 여핵에 alpha_1 대신 alpha_2를 쓰는 원문을 그대로 보존했다.
$M_2 \cong \Coker(\alpha_2)$가 $B_1$-가군으로서 성립함을 사용하면,
정리가 $(\Coker(\alpha_1)/B_1/A, \mathfrak q_1)$에 대하여 성립함을
알 수 있다.

\medskip\noindent
$A'$를 $\mathcal{C}$의 대상이라 하자. 여기서 Lemma
\ref{flat-lemma-induction-step}를 사용하면, $\mathfrak m_{A'}$ 위에 놓이는
$B_1 \otimes_A A'$의 소아이디얼 $\mathfrak q'$에 대하여
$(M_1 \otimes_A A')_{\mathfrak q'}$가 $A'$ 위에서 평탄이면
$(\Coker(\alpha_1) \otimes_A A')_{\mathfrak q'}$도 $A'$ 위에서
평탄임을 알 수 있다. 따라서 $F_{lf}$는
$(B_1)_{\mathfrak q_1}$ 위의 가군
$\Coker(\alpha_1)_{\mathfrak q_1}$에 결부된 함자 $F'_{lf}$의
부분함자이다. 앞 문단에 의해 $F'_{lf}$는 어떤 아이디얼
$J \subset A$에 대한 $A/J$로 공표현된다. 따라서 $A$를 $A/J$로
바꾸어 $\Coker(\alpha_1)_{\mathfrak q_1}$이 $A$ 위에서 평탄이라고
가정해도 된다.

\medskip\noindent
% P05-EN-CH38-0222: 문법은 자연스럽게 나타내되 아직 정의되지 않은 N_1을 쓰는 원문 대상을 보존했다.
$\Coker(\alpha_1)$은 $\mathfrak q_1$에서 $N_1/B_1/A$의 완전 데비사주가
존재하는 $B_1$-가군이고 $\Coker(\alpha_1)_{\mathfrak q_1}$은 $A$
위에서 평탄이므로, Lemma
\ref{flat-lemma-complete-devissage-flat-finite-type-module}에 의해
$\Coker(\alpha_1)$은 $A$-가군으로서 자유이고 특히 $A$-가군으로서
평탄이다. 따라서 Lemma \ref{flat-lemma-induction-step}에 의해
% P05-EN-CH38-0223: alpha_1이어야 할 두 무첨자 alpha를 원문 그대로 보존했다.
$F_{lf}(A')$가 공집합이 아닐 필요충분조건은
$\alpha \otimes 1_{A'}$가 단사인 것이다. $N_1 = \Im(\alpha_1) \subset
M_1$로 두면 다음 완전열들을 얻는다.
$$
0 \to N_1 \to M_1 \to \Coker(\alpha_1) \to 0
\quad\text{이고}\quad
B_1^{\oplus r_1} \to N_1 \to 0
$$
$\Coker(\alpha_1)$의 평탄성은 첫째 열이 보편 완전임을 함의한다
(Algebra, Lemma \ref{algebra-lemma-flat-universally-injective}를 보라).
따라서 $\alpha \otimes 1_{A'}$가 단사일 필요충분조건은
$B_1^{\oplus r_1} \otimes_A A' \to N_1 \otimes_A A'$가 동형인 것이다.
마지막으로 Theorem \ref{flat-theorem-flattening-map}을 적용하면 이 함자가
어떤 아이디얼 $I$에 대한 $A/I$로 공표현됨을 알 수 있고, 따라서
$F_{lf}$도 $A/I$로 공표현된다는 결론을 얻는다.

\medskip\noindent
마지막 명제를 증명하기 위하여 $A \to B$가 본질적으로 유한표현이고
$M$이 $B$ 위에서 유한표현이라고 가정하자. $F_{lf}$가 $A/I$로
공표현되게 하는 아이디얼을 $I \subset A$라 하자. $I$에 포함된
유한생성 아이디얼 $I_\lambda$들을 움직여
$I = \bigcup I_\lambda$로 쓰자. $F_{lf}(A/I) = \{*\}$이므로 어떤
$\lambda$에 대하여 $F_{lf}(A/I_\lambda) = \{*\}$이다. Lemma
\ref{flat-lemma-flat-at-point}의 (2)를 보라. 이는 분명히
$I = I_\lambda$를 함의한다.
\end{proof}

\begin{remark}
\label{flat-remark-flattening-local-scheme-theoretic}
Theorem \ref{flat-theorem-flattening-local}을 스킴론적으로 다시 서술하면
다음과 같다. $(X, x) \to (S, s)$를 국소 유한형인 점을 갖춘 스킴들의
사상이라 하고 $\mathcal{F}$를 유한형 준연접 $\mathcal{O}_X$-가군이라
하자. $S$가 닫힌점 $s$를 갖는 헨젤 국소 스킴이라고 가정하자. 다음
성질을 갖는 닫힌부분스킴 $Z \subset S$가 존재한다. 점을 갖춘 스킴들의
임의의 사상 $(T, t) \to (S, s)$에 대하여 다음 조건들은 동치이다.
\begin{enumerate}
\item $x \in X_s$로 가는 올 $X_t$의 모든 점에서 $\mathcal{F}_T$는
$T$ 위에 평탄이고,
\item $\Spec(\mathcal{O}_{T, t}) \to S$는 $Z$를 통해 인수분해된다.
\end{enumerate}
더욱이 $X \to S$가 $x$에서 유한표현이고 $\mathcal{F}_x$가
$\mathcal{O}_{X, x}$ 위에서 유한표현이면 $Z \to S$는 유한표현이다.
\end{remark}

\noindent
이제 위의 결과로부터 매우 일반적인 몇몇 결과를 거의 공짜로 얻을 수
있다. 아마 가장 흥미로운 경우는 $E = X_s$일 때임에 유의하자!

\begin{lemma}
\label{flat-lemma-freebie}
$S$를 닫힌점 $s$를 갖는 헨젤 국소환의 스펙트럼이라 하자.
$X \to S$를 국소 유한형인 스킴 사상이라 하고 $\mathcal{F}$를 유한형
준연접 $\mathcal{O}_X$-가군이라 하자. $E \subset X_s$를 부분집합이라
하자. 다음 성질을 갖는 닫힌부분스킴 $Z \subset S$가 존재한다. 점을
갖춘 스킴들의 임의의 사상 $(T, t) \to (S, s)$에 대하여 다음 조건들은
동치이다.
\begin{enumerate}
\item $E \subset X_s$의 한 점으로 가는 올 $X_t$의 모든 점에서
$\mathcal{F}_T$는 $T$ 위에 평탄이고,
\item $\Spec(\mathcal{O}_{T, t}) \to S$는 $Z$를 통해 인수분해된다.
\end{enumerate}
더욱이 $X \to S$가 국소 유한표현이고 $\mathcal{F}$가 유한표현이며
$E \subset X_s$가 닫히고 준콤팩트이면 $Z \to S$는 유한표현이다.
\end{lemma}

\begin{proof}
$x \in X_s$에 대하여 Remark
\ref{flat-remark-flattening-local-scheme-theoretic}에서 찾은 닫힌부분스킴을
$Z_x \subset S$로 나타내자. 그러면
$Z = \bigcap_{x \in E} Z_x$가 조건을 만족함이 분명하다!

\medskip\noindent
마지막 명제를 증명하기 위하여 $X$가 국소 유한표현이고
$\mathcal{F}$가 유한표현이며
% P05-EN-CH38-0224: 가정의 E 대신 결론의 Z를 쓰는 원문을 그대로 보존했다.
$Z$가 닫히고 준콤팩트라고 가정하자. 먼저
$E \subset \bigcup W_j$가 되도록 유한 개의 아핀 열린집합
$W_j \subset X$를 택하자. 각 사상 $W_j \to S$, 층
% P05-EN-CH38-0225: 정의되지 않은 X_j에 제한하는 원문을 그대로 보존했다.
$\mathcal{F}|_{X_j}$, 그리고 닫힌 부분집합 $E \cap W_j$에 대하여
결과를 증명하면 충분함이 분명하다. 따라서 $X$가 아핀이라고 가정해도
된다. 이 경우 More on Algebra, Lemma
\ref{more-algebra-lemma-limit-preserving-flat-at-primes}는 (1)로 정의되는
함자가 ``극한 보존''임을 보여 준다. 따라서 Theorem
\ref{flat-theorem-flattening-local}의 증명 마지막 부분과 정확히 같은
방식으로 $Z \to S$가 유한표현임을 보일 수 있다.
\end{proof}

\begin{remark}
\label{flat-remark-flattening-complete-noetherian}
Lemma \ref{flat-lemma-freebie}의 증명을 그 기원까지 거슬러 올라가면 길고
구불구불한 길을 만나게 된다. 그러나 다음을 가정하면
\begin{enumerate}
\item $f$는 유한형이고,
\item $\mathcal{F}$는 유한형 $\mathcal{O}_X$-가군이며,
\item $E = X_s$이고,
\item $S$는 뇌터 완비 국소환의 스펙트럼이다.
\end{enumerate}
더 초등적인 대수에만 전적으로 의존하는 다음 증명이 있다. 먼저 유한
아핀 열린 덮개를 택하여 $X$가 아핀인 경우로 귀착시킨다. 이 경우
More on Algebra, Lemma
\ref{more-algebra-lemma-flattening-complete-local-universal-property}에 의해
$Z$가 존재한다. 이 증명의 핵심 단계는 절단
$\Spec(\mathcal{O}_{S, s}/\mathfrak m_s^n)$ 안에서 닫힌부분스킴 $Z$를
한 단계씩 구성하는 것이다. 이는 기저가 아르틴이면 평탄화 층화가
언제나 존재한다는 사실과
$\mathcal{O}_{S, s} = \lim \mathcal{O}_{S, s}/\mathfrak m_s^n$이라는
사실에 의존한다.
\end{remark}




\section{한 보조정리의 변형들}
\label{flat-section-variants-mod-injective}

\noindent
이 절에서는 Algebra, Lemmas \ref{algebra-lemma-mod-injective-general}과
\ref{algebra-lemma-mod-injective}의 변형들을 논한다. 가장 일반적인 판은
Proposition \ref{flat-proposition-finite-type-injective-into-flat-mod-m}이다.
이는 \cite[Lemma 4.2.2]{GruRay}로 진술되었지만 같은 문헌의 증명은
Lemma \ref{flat-lemma-upstairs-finite-type-injective-into-flat-mod-m}에 진술된
더 약한 결과만을 준다. Proposition
\ref{flat-proposition-finite-type-injective-into-flat-mod-m}의 정교한 증명은
Ofer Gabber에 의한 것이다. 지금은 이 명제를 적용할 곳이 없으므로,
독자에게 Lemma
\ref{flat-lemma-upstairs-finite-type-injective-into-flat-mod-m}의 증명을 읽은
뒤 다음 절로 건너뛰기를 권한다. 이 보조정리는 다음 절에서 Theorem
\ref{flat-theorem-check-flatness-at-associated-points}를 증명하는 데 쓰인다.

\begin{situation}
\label{flat-situation-mod-injective}
$\varphi : A \to B$를 본질적으로 유한형인 국소환들의 국소 준동형이라
하자. $M$을 평탄 $A$-가군, $N$을 유한 $B$-가군이라 하고,
$\overline{u} : N/\mathfrak m_AN \to M/\mathfrak m_AM$가 단사인
$A$-가군 사상 $u : N \to M$을 택하자.
\end{situation}

\noindent
이 상황에서 목표는 $u$가 $A$-보편 단사이고, $N$이 $B$ 위에서
유한표현이며, $N$이 $A$-가군으로서 평탄임을 보이는 것이다. 이것이
참이면 주어진 상황에서 {\it 보조정리가 성립한다}고 하겠다.

\begin{lemma}
\label{flat-lemma-Noetherian-finite-type-injective-into-flat-mod-m}
Situation \ref{flat-situation-mod-injective}에서 환 $A$가 뇌터이면 보조정리가
성립한다.
\end{lemma}

\begin{proof}
Algebra, Lemma \ref{algebra-lemma-mod-injective}를 적용하면 $u$가
단사이고
% P05-EN-CH38-0226: u:N->M과 맞지 않는 원문의 몫 N/u(M)을 그대로 보존했다.
$N/u(M)$이 $A$ 위에서 평탄임을 알 수 있다. 그러면 $u$는
$A$-보편 단사이고(Algebra, Lemma
\ref{algebra-lemma-flat-tor-zero}), $N$은 $A$-평탄이다(Algebra, Lemma
\ref{algebra-lemma-flat-ses}). 이 경우 $B$가 뇌터이므로 $N$은
유한표현이다.
\end{proof}

\begin{lemma}
\label{flat-lemma-reduce-finite-type-injective-into-flat-mod-m}
$A_0$를 국소환이라 하자. $A = A_0$이고, $B$가 $A$ 위의 다항식
대수의 국소화이며, $N$이 $B$ 위에서 유한표현인 모든 Situation
\ref{flat-situation-mod-injective}에 대하여 보조정리가 성립한다면,
$A = A_0$인 모든 Situation \ref{flat-situation-mod-injective}에 대하여
보조정리가 성립한다.
\end{lemma}

\begin{proof}
$A \to B$, $u : N \to M$을 Situation
\ref{flat-situation-mod-injective}에서와 같이 택하자. $C$가 $A$ 위의
다항식 대수를 한 소아이디얼에서 국소화한 것이 되도록 $B = C/I$로
쓰자. $N$이 $C$-가군으로서 유한표현임을 보일 수 있다면, 더욱더
$N$이 $B$-가군으로서 유한표현임을 얻는다(Algebra, Lemma
\ref{algebra-lemma-finitely-presented-over-subring}을 보라). 따라서
$B$가 다항식 대수의 국소화라고 가정해도 된다. 다음으로 어떤 부분가군
$K \subset B^{\oplus n}$에 대하여 $N = B^{\oplus n}/K$로 쓰자.
$B/\mathfrak m_AB$는 체 위에서 본질적으로 유한형이므로 뇌터이다.
따라서 $K' = \sum Bk_i$와 $N' = B^{\oplus n}/K'$에 대하여 표준 전사
$N' \to N$이 동형 $N'/\mathfrak m_AN' \cong N/\mathfrak m_AN$을
유도하게 하는 유한 개의 원소 $k_1, \ldots, k_s \in K$가 존재한다.
이제 합성 $u' : N' \to M$에 대하여 보조정리가 성립하면 $u'$는
단사이고, 따라서 $N' = N$과 $u' = u$이다. 그러므로 원래 상황에
대하여 보조정리가 성립한다.
\end{proof}

\begin{lemma}
\label{flat-lemma-henselian-finite-type-injective-into-flat-mod-m}
Situation \ref{flat-situation-mod-injective}에서 환 $A$가 헨젤이면 보조정리가
성립한다.
\end{lemma}

\begin{proof}
$B$가 $A$ 위에서 본질적으로 유한표현이고 $N$이 $B$ 위에서
유한표현인 경우를 증명하면 충분하다. Lemma
\ref{flat-lemma-reduce-finite-type-injective-into-flat-mod-m}를 보라. 잠시
$N$이 $A$ 위에서 평탄이라는 가정을 더하자. 그러면 $N$은 자유
$A$-가군 $F_i$들의 여과 여극한 $N = \colim_i F_i$이고, 전이 사상
$u_{ii'} : F_i \to F_{i'}$들은 $\mathfrak m_A$로 나눈 뒤 단사이다.
Lemma \ref{flat-lemma-flat-finite-type-local-colimit-free}를 보라. Lemma
\ref{flat-lemma-universally-injective-local}에 의해 각 합성
$u_i : F_i \to M$은 $A$-보편 단사이고, 따라서 원하는 대로
$u = \colim u_i$도 $A$-보편 단사이다.

\medskip\noindent
$A$를 헨젤 국소환, $B$를 $A$ 위에서 본질적으로 유한표현인 환,
$N$을 $B$ 위에서 유한표현인 가군이라 하자. Theorem
\ref{flat-theorem-flattening-local}에 의해 $N/IN$은 $A/I$ 위에서 평탄이고,
$I = 0$이 아니면 $N/I^2N$은 $A/I^2$ 위에서 평탄이지 않게 하는
유한생성 아이디얼 $I \subset A$가 존재한다. 앞 문단의 결과에 의해
$A/I$ 위의 $u \bmod I : N/IN \to M/IM$에 대하여 보조정리가 성립한다.
다음 가환 그림을 생각하자.
$$
\xymatrix{
0 \ar[r] &
M \otimes_A I/I^2 \ar[r] &
M/I^2M \ar[r] &
M/IM \ar[r] & 0 \\
&
N \otimes_A I/I^2 \ar[r] \ar[u]^u &
N/I^2N \ar[r] \ar[u]^u &
N/IN \ar[r] \ar[u]^u & 0
}
$$
$\otimes$의 오른쪽 완전성과 $M$이 $A$ 위에서 평탄이라는 사실에 의해
두 행은 완전이다. 왼쪽 세로 화살표는 사상
$N/IN \otimes_{A/I} I/I^2 \to M/IM \otimes_{A/I} I/I^2$이므로 단사이다.
그림 추적으로 왼쪽 아래 화살표가 단사임을 알 수 있다. 즉,
% P05-EN-CH38-0227: 아래 왼쪽 단사성을 나타내지 않는 원문의 Tor 식을 그대로 보존했다.
$\text{Tor}^1_{A/I^2}(I/I^2, M/I^2) = 0$이다. Algebra, Remark
\ref{algebra-remark-Tor-ring-mod-ideal}을 보라. 따라서 Algebra, Lemma
\ref{algebra-lemma-what-does-it-mean}에 의해 $N/I^2N$은 $A/I^2$ 위에서
평탄인데, 이는 $I = 0$이 아니면 모순이다.
\end{proof}

\noindent
다음 보조정리는 $M$이 $B$-가군 구조를 갖고 $u$가 $B$-선형인
Situation \ref{flat-situation-mod-injective}의 특수한 경우를 다룬다. 이는
실제로 가장 자주 쓰이는 경우이며 일반적인 경우보다 증명하기 훨씬 쉽다.

\begin{lemma}
\label{flat-lemma-upstairs-finite-type-injective-into-flat-mod-m}
$A \to B$를 본질적으로 유한형인 국소환들의 국소 준동형이라 하고
$u : N \to M$을 $B$-가군 사상이라 하자. $N$이 유한 $B$-가군이고
$M$이 $A$ 위에서 평탄이며
$\overline{u} : N/\mathfrak m_A N \to M/\mathfrak m_A M$가 단사이면,
$u$는 $A$-보편 단사이고 $N$은 $B$ 위에서 유한표현이며 $N$은 $A$ 위에서
평탄이다.
\end{lemma}

\begin{proof}
$A \to A^h$를 $A$의 헨젤화라 하자. $B'$를
$B \otimes_A A^h$를 극대 아이디얼
$\mathfrak m_B \otimes A^h + B \otimes \mathfrak m_{A^h}$에서 국소화한
것이라 하자. $B \to B'$는 평탄이고 따라서 충실히 평탄이므로(Algebra,
Lemma \ref{algebra-lemma-local-flat-ff}를 보라), $A \to B$를
$A^h \to B'$로, 가군 $M$을 $M \otimes_B B'$로, 가군 $N$을
$N \otimes_B B'$로, $u$를 $u \otimes \text{id}_{B'}$로 바꾸어도 된다.
Algebra, Lemmas \ref{algebra-lemma-descend-properties-modules}와
\ref{algebra-lemma-flatness-descends-more-general}을 보라. 따라서 $A$가
헨젤 국소환이라고 가정해도 된다. 이 경우 보조정리는 더 일반적인
Lemma \ref{flat-lemma-henselian-finite-type-injective-into-flat-mod-m}에서
따른다.
\end{proof}

\begin{lemma}
\label{flat-lemma-valuation-ring-finite-type-injective-into-flat-mod-m}
Situation \ref{flat-situation-mod-injective}에서 환 $A$가 값매김환이면
보조정리가 성립한다.
\end{lemma}

\begin{proof}
$A$-가군이 평탄일 필요충분조건은 꼬임이 없는 것임을 상기하자. More on
Algebra, Lemma \ref{more-algebra-lemma-valuation-ring-torsion-free-flat}을
보라. $T \subset N$을 $A$-꼬임이라 하자. 그러면 $u(T) = 0$이고
$N/T$는 $A$-평탄이다. 따라서 $N/T$는 $B$ 위에서 유한표현이다. More
on Algebra, Lemma
\ref{more-algebra-lemma-flat-finite-type-valuation-ring-finite-presentation}을
보라. 그러므로 $T$는 유한 $B$-가군이다. Algebra, Lemma
\ref{algebra-lemma-extension}을 보라. $N/T$가 $A$-평탄이므로
$T/\mathfrak m_A T \subset N/\mathfrak m_A N$이다. Algebra, Lemma
\ref{algebra-lemma-flat-tor-zero}을 보라. $\overline{u}$는 단사이지만
$u(T) = 0$이므로 $T/\mathfrak m_A T = 0$이라는 결론을 얻는다. 따라서
나카야마 보조정리에 의해 $T = 0$이다. Algebra, Lemma
\ref{algebra-lemma-NAK}를 보라. 이 시점에서 세 주장 가운데 둘($N$이
$A$-평탄이고 $B$ 위에서 유한표현이라는 것)을 증명했고, 남은 것은
$u$가 보편 단사임을 보이는 것이다.

\medskip\noindent
Algebra, Theorem \ref{algebra-theorem-universally-exact-criteria}에 의해
모든 유한표현 $A$-가군 $Q$에 대하여
$N \otimes_A Q \to M \otimes_A Q$가 단사임을 보이면 충분하다. More on
Algebra, Lemma \ref{more-algebra-lemma-generalized-valuation-ring-modules}에
의해 $a \in \mathfrak m_A$가 영이 아닌 $Q = A/(a)$인 경우로
가정해도 된다. 따라서 $N/aN \to M/aM$가 단사임을 보이면 충분하다.
$u(x) \in aM$인 $x \in N$을 택하자. Lemma
\ref{flat-lemma-flat-finite-type-local-valuation-ring-has-content}에 의해 $x$는
내용 아이디얼 $I \subset A$를 갖는다. $I$는 유한생성이고(More on
Algebra, Lemma \ref{more-algebra-lemma-content-finitely-generated}를 보라)
$A$는 값매김환이므로 어떤 $b$에 대하여 $I = (b)$이다(Algebra, Lemma
\ref{algebra-lemma-characterize-valuation-ring}을 보라). More on Algebra,
Lemma \ref{more-algebra-lemma-equal-content}에 의해 원소 $u(x)$도 내용
아이디얼 $I$를 갖는다. $u(x) \in aM$이므로 More on Algebra,
Definition \ref{more-algebra-definition-content-ideal}에 의해
$(b) \subset (a)$이다. $x \in bN$이므로 원하는 대로
$x \in aN$이라는 결론을 얻는다.
\end{proof}

\noindent
다음 상황을 생각하자.
\begin{equation}
\label{flat-equation-star}
\begin{matrix}
A \to B\text{는 유한표현이고, }S \subset B
\text{는 곱셈적 부분집합이며, }\\
N\text{은 유한표현 }S^{-1}B\text{-가군이다}
\end{matrix}
\end{equation}
이 상황에서 {\it 순수 확장모형}이란
$\Spec(S^{-1}B) \subset U$인 아핀 열린집합 $U \subset \Spec(B)$와
$N$을 확장하는 유한표현 $\mathcal{O}(U)$-가군 $N'$로서, $N'$이
$A$-사영이고 $N' \to N = S^{-1}N'$이 $A$-보편 단사인 것을 말한다.

\medskip\noindent
(\ref{flat-equation-star})에서 $A \to A_1$이 환 사상이면 다음과 같이
기저변환할 수 있다. $B_1 = B \otimes_A A_1$로 두고, $S_1 \subset B_1$을
$S$의 상이라 하며, $N_1 = N \otimes_A A_1$로 둔다. 이는
$S_1^{-1}B_1 = S^{-1}B \otimes_A A_1$이기 때문이다. 이하에서는 이를
더 언급하지 않고 사용하겠다.

\begin{lemma}
\label{flat-lemma-properties-pure-spreadout}
(\ref{flat-equation-star})에서 순수 확장모형이 존재하면 다음이 성립한다.
\begin{enumerate}
\item $N$의 원소들은 $A$에서 내용 아이디얼을 가지며,
\item $u : N \to M$이 평탄 $A$-가군 $M$으로 가는 사상이고 $A$의
모든 극대 아이디얼 $\mathfrak m$에 대하여
$N/\mathfrak m N \to M/\mathfrak m M$가 단사이면 $u$는
$A$-보편 단사이다.
\end{enumerate}
\end{lemma}

\begin{proof}
순수 확장모형의 정의에서와 같이 $U$, $N'$을 택하자. $N'$이
$A$-사영이므로 임의의 원소 $x' \in N'$은 $A$에서 내용 아이디얼을
갖는다(이는 직접 쉽게 알 수 있지만 More on Algebra, Lemma
\ref{more-algebra-lemma-content-exists-flat-Mittag-Leffler}와 Algebra,
Example \ref{algebra-example-ML}에서도 따른다). $N' \to N$은
$A$-보편 단사이므로 임의의 $x' \in N'$의 상 $x \in N$은 $A$에서
내용 아이디얼을 갖는다(이는 $x'$의 내용 아이디얼과 같다). 일반적인
$x \in N$에 대해서는 $s x$가 $N' \to N$의 상에 속하도록 하는
$s \in S$를 택하고, $x$와 $sx$가 같은 내용 아이디얼을 갖는다는
사실을 사용한다.

\medskip\noindent
$u : N \to M$을 (2)에서와 같이 택하자. $u$가 $A$-보편 단사임을
보이기 위해 $A$를 한 극대 아이디얼에서의 국소화로 바꾸어도 된다
(작은 세부 사항은 생략한다). $A$가 극대 아이디얼 $\mathfrak m$을
갖는 국소환이라고 가정하자. $s \in S$를 택하고 합성
$$
N' \to N \xrightarrow{1/s} N \xrightarrow{u} M
$$
을 생각하자. 이 사상들은 각각 $\mathfrak m$으로 나눈 뒤 단사이므로
Lemma \ref{flat-lemma-universally-injective-local}에 의해 그 합성은
$A$-보편 단사이다. $N = \colim_{s \in S} (1/s)N'$이므로
% P05-EN-CH38-0229: 정의되지 않은 inversally를 쓰는 원문을 의도된 보편 단사로 나타냈다.
보편 단사 사상들의 여극한인 $u$도 $A$-보편 단사라는 결론을 얻는다.
\end{proof}

\begin{lemma}
\label{flat-lemma-find-pure-spreadout}
(\ref{flat-equation-star})에서 모든 $\mathfrak p \in \Spec(A)$에 대하여
$A_\mathfrak p/I$ 위에서 순수 확장모형을 갖게 하는 유한생성 아이디얼
$I \subset \mathfrak pA_\mathfrak p$가 존재한다.
\end{lemma}

\begin{proof}
$A$를 $A_\mathfrak p$로 바꾸어도 된다. 따라서 $A$가 국소환이고
$\mathfrak p$가 $A$의 극대 아이디얼 $\mathfrak m$이라고 가정해도
된다. $S^{-1}B$ 위의 $N$의 표시에서 분모를 제거하면 어떤 유한표현
$B$-가군 $N'$에 대하여 $N = S^{-1}N'$으로 쓸 수 있다.
$B/\mathfrak m B$가 뇌터이므로
$N'/\mathfrak m N' \to N/\mathfrak m N$의 핵 $K$는 유한생성이다.
따라서 $s$가 $K$를 소멸시키는 $s \in S$를 택할 수 있다. $B$를
$B_s$로 바꾸어도 되는데, 이는 단지 $\Spec(B)$의 한 아핀
열린부분스킴으로 옮기는 것이기 때문이다. 그러면 $S$의 원소들이
$N'/\mathfrak m N'$ 위에서 단사로 작용함을 얻는다. 이제
$\mathbf{Z}$ 위에서 본질적으로 유한형인 국소 부분환 $A_0 \subset A$,
$B = A \otimes_{A_0} B_0$인 유한형 환 사상 $A_0 \to B_0$, 그리고
$N' = B \otimes_{B_0} N'_0 = A \otimes_{A_0} N'_0$인 유한
$B_0$-가군 $N'_0$을 택한다. $I = \mathfrak m_{A_0} A$가 조건을
만족한다고 주장한다. 실제로
$$
N'/IN' = N'_0/\mathfrak m_{A_0} N'_0 \otimes_{\kappa_{A_0}} A/I
$$
이고 이는 $A/I$ 위에서 자유이다. $S$의 원소에 의한 곱셈은 극대
아이디얼로 나눈 뒤 단사이므로, 예를 들어 Lemma
\ref{flat-lemma-invert-universally-injective}에 의해
$N'/IN' \to N/IN$은 보편 단사이다.
\end{proof}

\begin{lemma}
\label{flat-lemma-universally-injective-if-flat}
(\ref{flat-equation-star})에서 $N$이 $A$-평탄이고 $M$이 평탄 $A$-가군이며,
모든 $\mathfrak p \in \Spec(A)$에 대하여
$u \otimes \text{id}_{\kappa(\mathfrak p)}$가 단사인 $A$-가군 사상
$u : N \to M$이 주어졌다고 가정하자. 그러면 $u$는 $A$-보편
단사이다.
\end{lemma}

\begin{proof}
Algebra, Lemma \ref{algebra-lemma-check-universally-injective-into-flat}에
의해 모든 아이디얼 $I \subset A$에 대하여 $N/IN \to M/IM$가
단사임을 확인하면 충분하다. $A$를 $A/I$로 바꾸면 $u$가 단사임을
증명하면 충분함을 알 수 있다.

\medskip\noindent
$u$가 단사임을 증명하자. $x \in N$을 $u$의 핵에 속하는 영이 아닌
원소라 하자. 그러면 가군 $Ax$의 약연관 소아이디얼 $\mathfrak p$가
존재한다. Algebra, Lemma \ref{algebra-lemma-weakly-ass-zero}를 보라.
$A$를 $A_\mathfrak p$로 바꾸어 $A$가 국소환이라고 가정해도 되고,
소멸자의 근기가 $\mathfrak m_A$인 영이 아닌 원소 $y \in Ax$를 찾을
수 있다. Algebra, Lemma \ref{algebra-lemma-weakly-ass-local}을 보라.
따라서 $\text{Supp}(y) \subset \Spec(S^{-1}B)$는 공집합이 아니고
$\Spec(S^{-1}B) \to \Spec(A)$의 닫힌 올에 포함된다. $A/I$ 위에서
순수 확장모형을 갖도록 유한생성 아이디얼
$I \subset \mathfrak m_A$를 택하자. Lemma
\ref{flat-lemma-find-pure-spreadout}을 보라. 그러면 어떤 $n$에 대하여
$I^n y = 0$이다. 이제 평탄성에 의해
% P05-EN-CH38-0230: y가 N에 있지만 Ann_M을 쓰고 정의되지 않은 R 위에서 텐서하는 원문을 그대로 보존했다.
$y \in \text{Ann}_M(I^n) = \text{Ann}_A(I^n) \otimes_R N$이다. 따라서
원하는 모순을 얻기 위해서는 다음 사상이 단사임을 보이면 충분하다.
$$
\text{Ann}_A(I^n) \otimes_R N
\longrightarrow
\text{Ann}_A(I^n) \otimes_R M
$$
$N$과 $M$은 평탄이고 $\text{Ann}_A(I^n)$은 $I^n$에 의해 소멸되므로,
$I$에 의해 소멸되는 모든 $A$-가군 $Q$에 대하여
$Q \otimes_A N \to Q \otimes_A M$가 단사임을 보이면 충분하다. 이는
$I$를 택한 방식과 Lemma \ref{flat-lemma-properties-pure-spreadout}의
(2)에 의해 성립한다.
\end{proof}

\begin{lemma}
\label{flat-lemma-big-intersection-is-zero}
$A$를 체가 아닌 국소 정역이라 하고 $S$를 $A$의 유한생성
아이디얼들의 집합이라 하자. $S$가 곱에 대하여 닫혀 있고
$\bigcup_{I \in S} V(I)$가 $\Spec(A)$의 일반점의 여집합이라고
가정하자. 그러면 $\bigcap_{I \in S} I = (0)$이다.
\end{lemma}

\begin{proof}
$\mathfrak m_A \subset A$는 $\Spec(A)$의 일반점이 아니므로 적어도
하나의 $I \in S$에 대하여 $I \subset \mathfrak m_A$이다. 따라서
$\bigcap_{I \in S} I \subset \mathfrak m_A$이다. 영이 아닌
$f \in \mathfrak m_A$를 택하자. 그러면
$V(f) \subset \bigcup_{I \in S} V(I)$.
$V(f)$ 위의 구성가능 위상은 준콤팩트이므로(Topology, Lemma
\ref{topology-lemma-constructible-hausdorff-quasi-compact}와 Algebra,
Lemma \ref{algebra-lemma-spec-spectral}을 보라), 어떤 $I_j \in S$들에
대하여 $V(f) \subset V(I_1) \cup \ldots \cup V(I_n)$임을 얻는다.
$I_1 \ldots I_n \in S$이므로 어떤 $I$에 대하여
$V(f) \subset V(I)$이다. $I$가 유한생성이므로 이는 어떤 $m$에 대하여
$I^m \subset (f)$임을 함의하고, $S$가 곱에 대하여 닫혀 있으므로 어떤
$I \in S$에 대하여 $I \subset (f^2)$임을 알 수 있다. 그러면
$f \in I$일 수 없다.
\end{proof}

\begin{lemma}
\label{flat-lemma-closed-points-complement}
$A$를 국소환이라 하고 $I, J \subset A$를 아이디얼들이라 하자.
$J$가 유한생성이고 모든 $n \geq 1$에 대하여 $I \subset J^n$이면,
$V(I)$는 $\Spec(A) \setminus V(J)$의 닫힌점들을 포함한다.
\end{lemma}

\begin{proof}
$\mathfrak p \subset A$를 $\Spec(A) \setminus V(J)$의 닫힌점이라 하자.
$I \subset \mathfrak p$임을 보이고자 한다. 그렇지 않으면 어떤
$f \in I$의 $A/\mathfrak p$에서의 상이 영이 아니다.
$V(J) \cap \Spec(A/\mathfrak p)$는 일반점이 아닌 점들의 집합임에
유의하자. 따라서 아이디얼들의 모임 $J^nA/\mathfrak p$에 Lemma
\ref{flat-lemma-big-intersection-is-zero}을 적용하면 $f$의 상이
$A/\mathfrak p$에서 영이라는 결론을 얻는다.
\end{proof}

\begin{lemma}
\label{flat-lemma-make-smaller-flatness-ideal}
$A$를 국소환이라 하자. $I \subset A$를 아이디얼이라 하고,
$U \subset \Spec(A)$를 준콤팩트 열린집합이라 하자.
$M$을 $A$-가군이라 하자. 다음을 가정하자.
\begin{enumerate}
\item $M/IM$은 $A/I$ 위에서 평탄이다.
\item $M$은 $U$ 위에서 평탄이다.
\end{enumerate}
그러면 $I_2 = \Ker(I \to \Gamma(U, I/I^2))$라 둘 때
$M/I_2M$은 $A/I_2$ 위에서 평탄이다.
\end{lemma}

\begin{proof}
Algebra, Lemma \ref{algebra-lemma-what-does-it-mean-again}에 의해
$M \otimes_A I/I_2 \to IM/I_2M$이 단사임을 보이면 충분하다.
가정 (2)에 의해 이는 $U$ 위에서 성립한다. 따라서
$M \otimes_A I/I_2$가 $U$ 위의 그 단면들로 단사됨을 보이면 충분하다.
$M \otimes_A I/I_2 = M/IM \otimes_A I/I_2$이고 $M/IM$은 유한 자유
$A/I$-가군들의 여과 여극한이다(Algebra, Theorem
\ref{algebra-theorem-lazard}). 그러므로 $I/I_2$가 $U$ 위의 그 단면들로
단사됨을 보이면 충분한데, 이는 $I_2$의 구성에서 따른다.
\end{proof}

\begin{proposition}
\label{flat-proposition-finite-type-injective-into-flat-mod-m}
$A \to B$를 본질적으로 유한형인 국소환들의 국소환 준동형이라 하자.
$M$을 평탄 $A$-가군, $N$을 유한 $B$-가군이라 하고,
$\overline{u} : N/\mathfrak m_AN \to M/\mathfrak m_AM$가 단사인
$A$-가군 사상 $u : N \to M$이 주어졌다고 하자. 그러면 $u$는
$A$-보편 단사이고, $N$은 $B$ 위에서 유한표현이며, $N$은 $A$ 위에서
평탄이다.
\end{proposition}

\begin{proof}
Lemma \ref{flat-lemma-reduce-finite-type-injective-into-flat-mod-m}에 의해
$B$가 유한표현 $A$-대수 $B_0$의 국소화이고 $N$이 유한표현
$B_0$-가군 $M_0$의 국소화라고 가정해도 된다. More on Morphisms,
Lemma \ref{more-morphisms-lemma-generic-flatness-stratification}에 의해
$\Spec(A)$ 위의 $\Spec(B_0)$ 위에 $\widetilde{M_0}$의
``일반 평탄성 층화''가 존재한다. 이를 다시 $N$으로 옮기면
닫힌 부분스킴들의 열
$$
S = \Spec(A) \supset S_0 \supset S_1 \supset \ldots \supset S_t = \emptyset
$$
을 얻는다. 여기서 $S_i \subset S$는 $A$의 유한생성 아이디얼로
잘려 나오고, $\widetilde{N}$을
$\Spec(B) \times_S (S_i \setminus S_{i + 1})$로 당긴 것은
$S_i \setminus S_{i + 1}$ 위에서 평탄이다. 명제를 $t$에 관한 귀납법으로
증명하겠다(기저 경우 $t = 1$도 다른 단계들과 동시에 증명한다).
$\Spec(A/J_i)$를 $S_i \setminus S_{i + 1}$의 스킴론적 폐포라 하자.

\medskip\noindent
{\bf 주장 1.} $N/J_iN$은 $A/J_i$ 위에서 평탄이다. $i = t - 1$이면
이는 즉시 성립하고, $i > 0$이면 귀납가정에서 따른다. 따라서
$t > 1$, $S_{t - 1} \not = \emptyset$, $J_0 = 0$이라고 가정해도 되며,
$N$이 평탄임을 증명해야 한다. $J \subset A$를 $S_1$을 정의하는
아이디얼이라 하자. 다시 $t$에 관한 귀납법을 쓰면 $J$의 거듭제곱들로
나눈 뒤의 평탄성도 얻는다. $A^h$를 $A$의 헨젤화라 하고, $B'$을
$B \otimes_A A^h$의 극대 아이디얼
$\mathfrak m_B \otimes A^h + B \otimes \mathfrak m_{A^h}$에서의
국소화라 하자. 그러면 $B \to B'$은 충실히 평탄이다.
$N' = N \otimes_B B'$라 두자. $N'$이 $A^h$-평탄일 필요충분조건은
$N$이 $A$-평탄인 것이다. Theorem \ref{flat-theorem-flattening-local}에 의해
$N'/IN'$이 $A^h/I$ 위에서 평탄이게 하는 가장 작은 아이디얼
$I \subset A^h$가 존재하고, $I$는 유한생성이다. 위에서 본 바에 의해
모든 $n \geq 1$에 대하여 $I \subset J^nA^h$이다.
$S_i^h \subset \Spec(A^h)$를 $S_i \subset \Spec(A)$의 역상이라 하자.
Lemma \ref{flat-lemma-closed-points-complement}에 의해 $V(I)$는
$U = \Spec(A^h) - S_1^h$의 닫힌점들을 포함한다. 구성에 의해 $N'$은
$U$ 위에서 $A^h$-평탄이다. Lemma
\ref{flat-lemma-make-smaller-flatness-ideal}에 의해
% P05-EN-CH38-0231: I_2가 A^h의 아이디얼인데 A/I_2 위의 평탄성을 말하는 원문을 그대로 보존했다.
$I_2 = \Ker(I \to \Gamma(U, I/I^2))$라 둘 때 $N'/I_2N'$은
$A/I_2$ 위에서 평탄이다. 따라서 $I$의 최소성에 의해 $I = I_2$이다.
이는 $U$ 위에서 국소적으로 $I = I^2$, 즉 모든 $u \in U$에 대하여
$I\mathcal{O}_{U, u} = (0)$ 또는 $I\mathcal{O}_{U, u} = (1)$임을
뜻한다. $V(I)$가 $U$의 닫힌점들을 포함하므로 $U$ 위에서 $I = 0$이다.
이 문단의 처음에 $A$를 $A/J_0$로 바꾸었기 때문에
$U \subset \Spec(A^h)$는 스킴론적으로 조밀하고, 따라서 $I = 0$이다.
그러므로 $N'$은 $A^h$-평탄이고 주장 1이 성립한다.

\medskip\noindent
주장 1 앞에서 서술한 상황으로 돌아가자. $A^h$를 $A$의 헨젤화,
$B'$을 $B \otimes_A A^h$의 극대 아이디얼
$\mathfrak m_B \otimes A^h + B \otimes \mathfrak m_{A^h}$에서의 국소화,
$N' = N \otimes_B B'$라 하면, 이제 Theorem
\ref{flat-theorem-flattening-local}의 평탄화 아이디얼 $I \subset A^h$가
멱영임을 알 수 있다. $nil(A^h)$가 멱영원들의 아이디얼을 나타낸다면
$nil(A^h) = nil(A) A^h$
이다(More on Algebra, Lemma
\ref{more-algebra-lemma-henselization-nil}). 따라서 $N/I_0N$이 $A/I_0$
위에서 평탄이게 하는 유한생성 멱영 아이디얼 $I_0 \subset A$가 존재한다.

\medskip\noindent
{\bf 주장 2.} 모든 소아이디얼 $\mathfrak p \subset A$에 대하여 사상
$\kappa(\mathfrak p) \otimes_A N \to \kappa(\mathfrak p) \otimes_A M$
은 단사이다.
% P05-EN-CH38-0232: "bad it this is false"라는 원문의 철자 오류를 문맥에 맞게 번역했다.
이 명제가 거짓이면 $\mathfrak p$를 나쁘다고 하자. 나쁜
소아이디얼들의 공집합이 아닌 사슬 $C$가 주어졌다고 하고
$\mathfrak p^* = \bigcup_{\mathfrak p \in C} \mathfrak p$라 두자.
Lemma \ref{flat-lemma-find-pure-spreadout}에 의해 $V(\mathfrak a)$ 위에서
순수 확장모형이 존재하게 하는 유한생성 아이디얼
$\mathfrak a \subset \mathfrak p^*A_{\mathfrak p^*}$가 존재한다.
$\mathfrak p^*$가 좋다면 Lemma
\ref{flat-lemma-properties-pure-spreadout}에 의해 $V(\mathfrak a)$의 점들은
좋을 것이다. 그러나 $\mathfrak a$가 유한생성이고
% P05-EN-CH38-0233: 우변 합집합에 소아이디얼 인자가 빠진 원문의 수식을 그대로 보존했다.
$\mathfrak p^*A_{\mathfrak p^*} = \bigcup_{\mathfrak p \in C}A_{\mathfrak p^*}$
이므로, $V(\mathfrak a)$가 어떤 $\mathfrak p \in C$를 포함한다는 모순을
얻는다. 따라서 $\mathfrak p^*$는 나쁘다. 초른 보조정리에 의해 나쁜
소아이디얼이 존재하면 극대인 것이 존재하며, 이를 $\mathfrak p$라 하자.
다시 말해 모든 $\mathfrak p' \supset \mathfrak p$,
$\mathfrak p' \not = \mathfrak p$가 좋다고 가정해도 된다. 이 경우 모든
$f \in A$, $f \not \in \mathfrak p$에 대하여 사상
$u \otimes \text{id}_{A/(\mathfrak p + f)}$가 보편 단사임을 알 수 있다.
Lemma \ref{flat-lemma-universally-injective-if-flat}을 보라. 따라서 부분가군
$f(N/\mathfrak pN)$들이 정의하는 위상에 관하여 $N/\mathfrak p N$이
분리됨을 보이면 충분하다. $B \to B'$이 충실히 평탄이므로 가군
$N'/\mathfrak p N'$에 대하여 같은 명제를 증명하면 충분하다. Lemma
\ref{flat-lemma-flat-finite-type-local-colimit-free}와 More on Algebra, Lemma
\ref{more-algebra-lemma-content-exists-flat-Mittag-Leffler}에 의해
$N'/\mathfrak pN'$의 원소들은 $A^h/\mathfrak pA^h$에서 내용 아이디얼을
갖는다. 따라서
$\bigcap_{f \in A, f \not \in \mathfrak p} f(A^h/\mathfrak p A^h) = 0$.
임을 보이면 충분하다. 이어서 $\mathfrak p A^h$ 위에서 극소인 모든
소아이디얼 $\mathfrak q \subset A^h$에 대하여 $A^h/\mathfrak q A^h$에
관한 같은 명제를 보이면 충분하다. $A \to A^h$가 헨젤화이므로 모든
$\mathfrak q$는 $\mathfrak p$로 축약되고, 모든
$\mathfrak q' \supset \mathfrak q$, $\mathfrak q' \not = \mathfrak q$는
$\mathfrak p$를 진정으로 포함하는 소아이디얼 $\mathfrak p'$로 축약된다.
따라서 Lemma \ref{flat-lemma-big-intersection-is-zero}에서 이 교집합들의
소멸을 얻는다.

\medskip\noindent
이제 모든 것을 종합할 수 있다. 즉 주장 1과 주장 2 및 Lemma
\ref{flat-lemma-universally-injective-if-flat}에 의해
$N/I_0 N \to M/I_0M$은 $A/I_0$-보편 단사이다. 그러면 그림들
$$
\xymatrix{
N \otimes_A (I_0^n/I_0^{n + 1}) \ar[r] \ar[d] &
M \otimes_A (I_0^n/I_0^{n + 1}) \ar@{=}[d] \\
I_0^n N /I_0^{n + 1} N \ar[r] &
I_0^n M /I_0^{n + 1} M
}
$$
은 왼쪽 수직 화살표들이 단사임을 보여 준다. 따라서 Algebra, Lemma
\ref{algebra-lemma-what-does-it-mean-again}에 의해 $N$이 평탄임을 알 수
있다. 이와 비슷하게 $u$의 보편 단사성은 ($N$의 평탄성을 먼저
증명하지 않고도) $I_0$로 나눈 경우로 환원된다. 이로써 증명이 끝난다.
\end{proof}


\section{평탄 유한형 가군, 제III부}
\label{flat-section-finite-type-flat-III}

\noindent
다음 결과는 이 장의 주요 결과 가운데 하나이다.

\begin{theorem}
\label{flat-theorem-check-flatness-at-associated-points}
$f : X \to S$가 국소적으로 유한형이라고 하자.
$\mathcal{F}$를 유한형 준연접 $\mathcal{O}_X$-가군이라 하자.
$x \in X$의 상을 $s \in S$라 하자. 다음은 서로 동치이다.
\begin{enumerate}
\item $\mathcal{F}$는 $x$에서 $S$ 위로 평탄이다.
\item $x$로 특수화되는 모든
$x' \in \text{Ass}_{X_s}(\mathcal{F}_s)$에 대하여 $\mathcal{F}$는
$x'$에서 $S$ 위로 평탄이다.
\end{enumerate}
\end{theorem}

\begin{proof}
$x$로 특수화되는 모든 점에 대하여 $\mathcal{F}_{x'}$는
$\mathcal{F}_x$의 국소화이므로 (1)이 (2)를 함의함은 명백하다.
$A = \mathcal{O}_{S, s}$, $B = \mathcal{O}_{X, x}$,
$N = \mathcal{F}_x$라 두자. $\Sigma \subset B$를
$N/\mathfrak m_AN$ 위에서 영인자로 작용하지 않는 $B$의 원소들로
이루어진 곱셈적 부분집합이라 하자. 가정 (2)와 Lemma
\ref{flat-lemma-homothety-spectrum}에서의
% P05-EN-CH38-0234: 가군에 Spec을 적용한 원문의 수식을 그대로 보존했다.
$\Spec(\Sigma^{-1}N)$의 서술에 의해 $\Sigma^{-1}N$은 $A$-평탄이다.
한편 구성에 의해 사상 $N \to \Sigma^{-1}N$은 $\mathfrak m_A$로 나눈
뒤 단사이다. 따라서 Lemma
\ref{flat-lemma-upstairs-finite-type-injective-into-flat-mod-m}을 적용하면 된다.
\end{proof}

\noindent
이제 이를 직접 적용하여 다음의 유용한 결과들을 얻는다.

\begin{lemma}
\label{flat-lemma-check-along-closed-fibre}
$S$를 닫힌점 $s$를 갖는 국소 스킴이라 하자.
$f : X \to S$가 국소적으로 유한형이라고 하자.
$\mathcal{F}$를 유한형 준연접 $\mathcal{O}_X$-가군이라 하자.
다음을 가정하자.
\begin{enumerate}
\item $\text{Ass}_{X/S}(\mathcal{F})$의 모든 점은 닫힌 올 $X_s$의
한 점으로 특수화된다\footnote{예를 들어 $f$가 유한형이고
$\mathcal{F}$가 $X_s$를 따라 순수하거나, $f$가 고유이면 이는 성립한다.}.
\item $\mathcal{F}$는 $X_s$의 모든 점에서 $S$ 위로 평탄이다.
\end{enumerate}
그러면 $\mathcal{F}$는 $S$ 위에서 평탄이다.
\end{lemma}

\begin{proof}
Theorem \ref{flat-theorem-check-flatness-at-associated-points}에 의해
$S$ 위의 $\mathcal{F}$의 상대 연관점 집합에 속하는 점들에서만
평탄성을 확인하면 충분하다는 사실에서 즉시 따른다.
\end{proof}











\section{보편 평탄화}
\label{flat-section-flattening-final}

\noindent
$f : X \to S$가 스킴들의 고유 유한표현 사상이면 $f$의 보편 평탄화를
찾을 수 있다. 이 절에서는 이 사실과 그 몇 가지 변형을 논한다.

\begin{lemma}
\label{flat-lemma-free-at-generic-points-representable}
Situation \ref{flat-situation-free-at-generic-points}에서 다음이 성립한다.
각 $p \geq 0$에 대하여 함자 $H_p$
(\ref{flat-equation-free-at-generic-points})는 국소 닫힌 몰입
$S_p \to S$로 표현가능하다. $\mathcal{F}$가 유한표현이면
$S_p \to S$도 유한표현이다.
\end{lemma}

\begin{proof}
각 $S$에 대하여 모든 $p \geq 0$에 관한 명제를 동시에 증명하겠다.
Lemma \ref{flat-lemma-free-at-generic-points}에 의해 함자 $H_p$는 fppf 위상에
관한 층이다. 따라서 Descent, Lemma
\ref{descent-lemma-descent-data-sheaves}, More on Morphisms, Lemma
\ref{more-morphisms-lemma-separated-locally-quasi-finite-morphisms-fppf-descend}
및 Descent, Lemma
\ref{descent-lemma-descending-fppf-property-immersion}을 종합하면 이 문제는
$S$ 위의 에탈 위상에 관하여 국소적임을 알 수 있다. 특히 이 문제는
$S$ 위에서 자리스키 국소적이다.

\medskip\noindent
$s \in S$에 대하여 $\xi_s$가 올 $X_s$의 유일한 일반점을 나타내게 하자.
모든 $s \in S$에 대하여 $\mathcal{F}$의 제한 $\mathcal{F}_s$는
$\xi_s$의 어떤 근방에서 어떤 계수 $p(s) \geq 0$의 국소 자유 가군임에
유의하자. (이는 $X_s$가 기약이고 매끄러우므로 $X_s$ 위의
$\mathcal{F}_s$에 대한 일반 평탄성에서 따른다. Algebra, Lemma
\ref{algebra-lemma-generic-flatness-Noetherian}을 보라. 다만 이는
필요 이상으로 강한 결과이다.) 뒤에서 쓰기 위해 다음을 기록해 둔다.
$$
p(s) =
\dim_{\kappa(\xi_s)}(
\mathcal{F}_{\xi_s} \otimes_{\mathcal{O}_{X, \xi_s}} \kappa(\xi_s)
).
$$
특히 $H_{p(s)}(s)$는 공집합이 아니고, $q \not = p(s)$이면
$H_q(s)$는 공집합이다.

\medskip\noindent
$U \subset X$를 열린 부분스킴이라 하자. $f : X \to S$는 매끄러우므로
열린사상이다. (\ref{flat-equation-free-at-generic-points})에서 쌍
$(f|_U : U \to f(U), \mathcal{F}|_U)$에 대한 함자 $H_p$와 쌍
$(f|_{f^{-1}(f(U))}, \mathcal{F}|_{f^{-1}(f(U))})$
에 대한 함자 $H_p$가 같음은 즉시 따른다. 따라서 $f(U)$ 위에서
$S_p$의 존재를 증명할 때 언제나 $X$를 $U$로 바꾸어도 된다.

\medskip\noindent
$s \in S$를 택하자. $\mathcal{F}|_U$가 많아야 $p(s)$개의 원소로
생성되는 $\xi_s$의 아핀 열린근방 $U$가 존재한다. 위의 논증에 의해
% P05-EN-CH38-0235: "an neighbourhood"라는 원문의 관사 오류를 문맥에 맞게 번역했다.
$s$의 근방에서 $H_{p(s)}$에 관한 명제를 증명할 때 $\mathcal{F}$가
$p(s)$개의 원소로 생성된다고, 즉 전사
$$
u : \mathcal{O}_X^{\oplus p(s)} \longrightarrow \mathcal{F}
$$
가 존재한다고 가정해도 된다. 이 경우 $H_{p(s)}$가 사상 $u$에 대한
$F_{iso}$ (\ref{flat-equation-iso})와 같음은 명백하다(이는 Lemma
\ref{flat-lemma-injectivity-map-source-flat-pure}에서 즉시 따르지만, $S$와
$X$가 모두 아핀이 되도록 조금 더 축소한 뒤 Lemma
\ref{flat-lemma-induction-step-fp}을 적용해도 따른다). 따라서 Theorem
\ref{flat-theorem-flattening-map}을 적용하면 $H_{p(s)}$가 $s$의 근방에서
닫힌 몰입으로 표현가능함을 알 수 있다.

\medskip\noindent
이제 결과는 위의 내용에서 형식적으로 따른다. 실제로 위의 논증은
함수 $s \mapsto p(s)$가 $S$ 위에서 국소적으로 유계임을 보여 준다.
따라서 $p = \max_{s \in S} p(s)$에 관한 귀납법을 쓸 수 있다. 위에서
본 바에 의해 함자 $H_p$는 닫힌 몰입 $S_p \to S$로 표현가능하다.
$S$를 $S \setminus S_p$로 바꾸면 최댓값이 적어도 하나 감소하므로
귀납가정으로 끝난다.

\medskip\noindent
$\mathcal{F}$가 유한표현이라고 가정하자. Lemma
\ref{flat-lemma-free-at-generic-points}의 (2)와 Limits, Remark
\ref{limits-remark-limit-preserving}에 의해 $S_p \to S$는 국소적으로
유한표현이다. 이제 앞 문단의 귀납논증을 다시 쓰면 $S$가 아핀일 때
각 $S_p$가 준콤팩트임을 알 수 있다. 먼저
$p = \max_{s \in S} p(s)$이면 $S_p \subset S$는 닫혔으므로(위를 보라)
준콤팩트이다. 다음으로 $S_p \to S$가 유한표현 닫힌 몰입이므로
$U = S \setminus S_p$는 $S$의 준콤팩트 열린집합이다(예를 들어
Morphisms, Section \ref{morphisms-section-constructible}의 논의를 보라).
이어서 $S_{p - 1} \to U$는 유한표현 닫힌 몰입이므로 $S_{p - 1}$은
준콤팩트이고 $U' = S \setminus (S_p \cup S_{p - 1})$도 준콤팩트이다.
이와 같이 계속한다.
\end{proof}

\begin{lemma}
\label{flat-lemma-localize-flat-dimension-n}
Situation \ref{flat-situation-flat-dimension-n}에서 다음이 성립한다.
$h : X' \to X$를 에탈 사상이라 하자.
$\mathcal{F}' = h^*\mathcal{F}$ 및 $f' = f \circ h$라 두자.
$F_n'$을 $(f' : X' \to S, \mathcal{F}')$에 결부된
(\ref{flat-equation-flat-dimension-n})이라 하자. 그러면 $F_n$은 $F_n'$의
부분함자이고, $h(X') \supset \text{Ass}_{X/S}(\mathcal{F})$이면
$F_n = F'_n$이다.
\end{lemma}

\begin{proof}
$T \to S$를 임의의 사상이라 하자.
% P05-EN-CH38-0236: h_T를 정의되지 않은 g의 기저변환이라고 하는 원문을 그대로 보존했다.
그러면 $h_T : X'_T \to X_T$는 에탈 사상 $g$의 기저변환이므로 에탈이다.
$t \in T$에 대하여 $Z \subset X_t$가 $\mathcal{F}_T$가 $T$ 위에서
평탄이지 않은 점들의 집합을 나타내게 하고, 이와 비슷하게
$Z' \subset X'_t$가 $\mathcal{F}'_T$가 $T$ 위에서 평탄이지 않은 점들의
집합을 나타내게 하자. $\mathcal{F}'_T = h_T^*\mathcal{F}_T$이므로
$Z' = h_t^{-1}(Z)$이다. Morphisms, Lemma
\ref{morphisms-lemma-flat-permanence}을 보라. 따라서 $Z' \to Z$는 에탈
사상이므로 $\dim(Z') \leq \dim(Z)$이다(예를 들어 Descent, Lemma
\ref{descent-lemma-dimension-at-point-local}에 의하거나, 단지 에탈
사상이 상대차원 $0$인 매끄러운 사상이기 때문이다). 이는
$F_n \subset F_n'$을 함의한다.

\medskip\noindent
마지막으로 $h(X') \supset \text{Ass}_{X/S}(\mathcal{F})$이고,
$F_n'(T)$가 공집합이 아닌 사상 $T \to S$, 즉 $\mathcal{F}'_T$가
$T$ 위에서 차원 $\geq n$에서 평탄이게 하는 사상이 주어졌다고 하자.
점 $t \in T$를 택하고 $Z \subset X_t$와 $Z' \subset X'_t$를 위와 같이
두자. 모순을 얻기 위해 $\dim(Z) \geq n$이라고 가정하자. 차원
$\geq n$인 성분에 대응하는 일반점 $\xi \in Z$를 택하자.
$x \in \text{Ass}_{X_t}(\mathcal{F}_t)$를 $\xi$의 일반화라 하자.
그러면 Divisors, Lemma
\ref{divisors-lemma-base-change-relative-assassin}과 Remark
\ref{divisors-remark-base-change-relative-assassin}에 의해 $x$는
$\text{Ass}_{X/S}(\mathcal{F})$의 한 점으로 간다. 따라서 $x$는
$h_T$의 상에 속하여, 어떤 $x' \in X'_T$에 대하여 $x = h_T(x')$이다.
그러나 $x \leadsto \xi$이고 $\dim(Z') < n$이므로 $x' \not \in Z'$이다.
따라서 $\mathcal{F}'_T$는 $x'$에서 $T$ 위로 평탄이고, 이는
$\mathcal{F}_T$가 $x$에서 $T$ 위로 평탄임을 함의한다(Morphisms,
Lemma \ref{morphisms-lemma-flat-permanence}). 이것이 그러한 모든 $x$에
대하여 성립하므로 Theorem
\ref{flat-theorem-check-flatness-at-associated-points}에 의해
$\mathcal{F}_T$는 $\xi$에서 $T$ 위로 평탄이다. 이는 원하는 모순이다.
\end{proof}

\begin{lemma}
\label{flat-lemma-compare-H-F}
$X \to S$가 차원 $d$의 기하적으로 기약인 올들을 갖는 아핀 스킴들의
매끄러운 사상이고, $\mathcal{F}$가 유한표현 준연접
$\mathcal{O}_X$-가군이라고 가정하자. 그러면 어떤 $c \geq 0$에 대하여
$F_d = \coprod_{p = 0, \ldots, c} H_p$이다. 여기서 $F_d$는
(\ref{flat-equation-flat-dimension-n})에, $H_p$는
(\ref{flat-equation-free-at-generic-points})에 나오는 함자이다.
\end{lemma}

\begin{proof}
$X$가 아핀이고 $\mathcal{F}$가 유한표현 준연접이므로 $\mathcal{F}$는
어떤 $c \geq 0$개의 원소로 생성될 수 있다. 그러면 모든 점
$x \in X$에서 $\dim_{\kappa(x)}(\mathcal{F}_x \otimes \kappa(x))$는
$c$를 넘지 않는다. 특히 $p > c$이면 $H_p = \emptyset$이다. 또한 포함
$\coprod H_p \to F_d$가 분명히 존재함에 유의하자. 이제 이 보조정리의
내용은 다음과 같다. 기저변환 $\mathcal{F}_T$가 $T$ 위에서 차원
$\geq d$에서 평탄이고 $t \in T$이면, $\mathcal{F}_T$는 $X_t$의 유일한
일반점 $\xi$의 어떤 열린근방 $U \subset X_T$에서 어떤 계수 $r$의
자유 가군이다. 실제로 그러면 $H_r$은 $t$의 열린근방인 $U$의 상을
포함한다. $U$의 존재는 More on Morphisms, Lemma
\ref{more-morphisms-lemma-flat-and-free-at-point-fibre}에서 따른다.
\end{proof}

\begin{lemma}
\label{flat-lemma-flat-dimension-n-representable}
Situation \ref{flat-situation-flat-dimension-n}에서 다음이 성립한다.
% P05-EN-CH38-0237: 두 Let 절 사이의 접속사가 빠진 원문을 문맥에 맞게 번역했다.
$s \in S$ 및 $d \geq 0$을 택하자. 다음을 가정하자.
\begin{enumerate}
\item 어떤 점 $s \in S$ 위에서 $\mathcal{F}/X/S$의 완전 데비사주가
존재한다.
\item $X$는 $S$ 위에서 유한표현이다.
\item $\mathcal{F}$는 유한표현 $\mathcal{O}_X$-가군이다.
\item $\mathcal{F}$는 $S$ 위에서 차원 $\geq d + 1$에서 평탄이다.
\end{enumerate}
그러면 필요하다면 $S$를 $s$의 열린근방으로 바꾼 뒤 함자 $F_d$
(\ref{flat-equation-flat-dimension-n})는 유한표현 단사 사상
$Z_d \to S$로 표현가능하다.
\end{lemma}

\begin{proof}
먼저 $X$, $S$가 아핀 스킴들이고, $S$ 위의 아핀 스킴들의 범주에서
$F_d$가 유한표현 단사 사상 $Z_d \to S$로 표현가능함을 증명하면
충분함을 알아 두자. (물론 $Z_d$가 아핀일 것을 요구하지는 않는다.)
따라서 이 보조정리의 증명 전체에서 $S$ 위의 아핀 스킴들의 범주에서
작업한다.

\medskip\noindent
$s$ 위의 $\mathcal{F}/X/S$의 완전 데비사주를
$(Z_k, Y_k, i_k, \pi_k, \mathcal{G}_k, \alpha_k)_{k = 1, \ldots, n}$이라
하자. Definition \ref{flat-definition-complete-devissage}를 보라. 데비사주의
길이 $n$에 관한 귀납법을 쓰겠다. Definition
\ref{flat-definition-one-step-devissage}에 따라 $Y_k \to S$는 기하적으로
기약인 올들을 갖는 매끄러운 사상임을 상기하자. $d_k$를 $S$ 위의
$Y_k$의 상대차원이라 하자. 또한 $i_{k, *}\mathcal{G}_k =
\Coker(\alpha_k)$이고 $i_k$가 닫힌 몰입임을 상기하자. 위에서 참조한
정의들에 의해
$d_1 = \dim(\text{Supp}(\mathcal{F}_s))$이고
$$
d_k = \dim(\text{Supp}(\Coker(\alpha_{k - 1})_s))
= \dim(\text{Supp}(\mathcal{G}_{k, s}))
$$
이고, $k = 2, \ldots, n$에 대하여 뒤의 두 식이 성립한다. $\alpha_k$가
$(Y_k)_s$의 일반점에서 동형이므로
$d_1 > d_2 > \ldots > d_n \geq 0$이 따른다.

\medskip\noindent
$i_1$은 닫힌 몰입이고 $\mathcal{F} = i_{1, *}\mathcal{G}_1$임에
유의하자. 따라서 $T$가 아핀인 스킴 사상 $T \to S$가 주어지면
$\mathcal{F}_T = i_{1, T, *}\mathcal{G}_{1, T}$이고 $i_{1, T}$는
여전히 $T$ 위의 스킴들의 닫힌 몰입이다. 그러므로 $\mathcal{F}_T$가
$T$ 위에서 차원 $\geq d$에서 평탄일 필요충분조건은
$\mathcal{G}_{1, T}$가 $T$ 위에서 차원 $\geq d$에서 평탄인 것이다.
$\pi_1 : Z_1 \to Y_1$이 유한이므로 같은 방식으로
$\mathcal{G}_{1, T}$가 $T$ 위에서 차원 $\geq d$에서 평탄일
필요충분조건은 $\pi_{1, T, *}\mathcal{G}_{1, T}$가 $T$ 위에서 차원
$\geq d$에서 평탄인 것임을 알 수 있다. 같은 논증은
``차원 $\geq d + 1$에서 평탄''에도 적용되므로, $\mathcal{F}$에 대한
가정에 의해 특히 $\pi_{1, *}\mathcal{G}_1$은 $S$ 위에서 차원
$\geq d + 1$에서 평탄이다.

\medskip\noindent
$d_1 > d$라고 가정하자. 위의 논의에서 특히
$\pi_{1, *}\mathcal{G}_1$이 $(Y_1)_s$의 일반점에서 $S$ 위로 평탄임이
따른다. Lemma \ref{flat-lemma-induction-step-fp}에 의해 $S$를 $s$의 아핀
근방으로 바꾸고 $\alpha_1$이 $S$-보편 단사라고 가정해도 된다.
$\alpha_1$이 $S$-보편 단사이므로 $T$가 아핀인 임의의 사상
$T \to S$에 대하여 짧은 완전열
$$
0 \to \mathcal{O}_{Y_{1, T}}^{\oplus r_1}
\to \pi_{1, T, *}\mathcal{G}_{1, T} \to \Coker(\alpha_1)_T \to 0
$$
을 얻고 첫째 화살표는 여전히 $T$-보편 단사이다. 따라서
$(Y_1)_T$에서 $\pi_{1, T, *}\mathcal{G}_{1, T}$가 $T$ 위로 평탄인
점들의 집합은 $(Y_1)_T$에서
% P05-EN-CH38-0238: 이 기저변환 문맥에서 T 대신 S 위의 평탄성을 쓰는 원문을 그대로 보존했다.
$\Coker(\alpha_1)_T$가 $S$ 위로 평탄인 점들의 집합과 같다. 이렇게
문제는 길이 $n - 1$의 완전 데비사주를 갖는 층 $\Coker(\alpha_1)$에
관한 문제로 환원되며, 귀납법으로 끝난다.

\medskip\noindent
$d_1 < d$이면 $F_d$는 $S$로 표현되므로 끝난다.

\medskip\noindent
마지막 경우는 $d_1 = d$인 경우이다. 이는 Lemma
\ref{flat-lemma-compare-H-F}와 Lemma
\ref{flat-lemma-free-at-generic-points-representable}을 종합하면 따른다.
\end{proof}

\begin{theorem}
\label{flat-theorem-flat-dimension-n-representable}
Situation \ref{flat-situation-flat-dimension-n}에서 다음이 성립한다.
더 나아가 $f$가 유한표현이고, $\mathcal{F}$가 유한표현
$\mathcal{O}_X$-가군이며, $\mathcal{F}$가 $S$에 대하여 순수하다고
가정하자. 그러면 $F_n$은 유한표현 단사 사상
$Z_n \to S$로 표현가능하다.
\end{theorem}

\begin{proof}
Lemma \ref{flat-lemma-flat-dimension-n}에 의해 함자 $F_n$은 fppf 위상에 관한
층이다. 유한표현 단사 사상은 분리이고 준유한임에 유의하자(Morphisms,
Lemma
\ref{morphisms-lemma-monomorphism-loc-finite-type-loc-quasi-finite}).
따라서 Descent, Lemma \ref{descent-lemma-descent-data-sheaves}, More on
Morphisms, Lemma
\ref{more-morphisms-lemma-separated-locally-quasi-finite-morphisms-fppf-descend}
및 Descent, Lemmas \ref{descent-lemma-descending-property-monomorphism},
\ref{descent-lemma-descending-property-finite-presentation}을 종합하면
이 문제는 $S$ 위의 에탈 위상에 관하여 국소적임을 알 수 있다.

\medskip\noindent
특히 이 상황은 $S$ 위의 자리스키 위상에 관하여 국소적이므로 $S$가
아핀이라고 가정해도 된다. 이 경우 $f$의 올들의 차원은 위로 유계이므로
충분히 큰 $n$에 대하여 $F_n$이 표현가능함을 알 수 있다. 따라서
$n$에 관한 하강귀납법을 쓸 수 있다. $F_{n + 1}$이 유한표현 단사 사상
$Z_{n + 1} \to S$로 표현가능하다는 것을 안다고 하자. 기저변환
$X_{n + 1} = Z_{n + 1} \times_S X$와 $\mathcal{F}$를 $X_{n + 1}$로
당긴 $\mathcal{F}_{n + 1}$을 생각하자. 사상 $Z_{n + 1} \to S$는
유한표현 단사 사상이므로 준유한이다. 따라서 Lemma
\ref{flat-lemma-quasi-finite-base-change}에 의해 $\mathcal{F}_{n + 1}$은
$Z_{n + 1}$에 대하여 순수하다. $F_n$은 $F_{n + 1}$의 부분함자이므로
$F_n$에 관한 결과를 증명하려면 상황
$\mathcal{F}_{n + 1}/X_{n + 1}/Z_{n + 1}$에 대응하는 함자에 관한 결과를
증명하면 충분하다. 이렇게 하여
% P05-EN-CH38-0239: 표현 스킴 Z_{n+1} 대신 S_{n+1}을 쓰는 원문의 등식을 그대로 보존했다.
$S_{n + 1} = S$인 경우의 $F_n$에 관한 결과로 환원된다. 즉
$\mathcal{F}$가 $S$ 위에서 차원 $\geq n + 1$에서 평탄이라고 가정해도
된다.

\medskip\noindent
$n$을 고정하고 $\mathcal{F}$가 $S$ 위에서 차원 $\geq n + 1$에서
평탄이라고 가정하자. 증명을 끝내려면 $F_n$이 유한표현 단사 사상
$Z_n \to S$로 표현가능함을 보여야 한다. 이 문제는 $S$ 위의 에탈
위상에 관하여 국소적이므로, 모든 $s \in S$에 대하여 $S'$로 기저변환한
뒤 결과가 성립하게 하는 기본 에탈 근방
$(S', s') \to (S, s)$가 존재함을 보이면 충분하다. 따라서 Lemma
\ref{flat-lemma-existence-complete}에 의해 각 $j$마다 $s$ 위에서
$\mathcal{F}_j/Y_j/S$의 완전 데비사주가 존재하고
$X_s \subset \bigcup h_j(Y_j)$가 되게 하는 에탈 사상들
$h_j : Y_j \to X$, $j = 1, \ldots, m$이 존재한다고 가정해도 된다.
여기서 $\mathcal{F}_j$는 $\mathcal{F}$를 $Y_j$로 당긴 것이다.
% P05-EN-CH38-0240: "flat over in dimensions"라는 원문의 중복 전치사를 문맥에 맞게 번역했다.
Lemma \ref{flat-lemma-localize-flat-dimension-n}에 의해 층들 $\mathcal{F}_j$는
여전히 $S$ 위에서 차원 $\geq n + 1$에서 평탄임에 유의하자.
$W = \bigcup h_j(Y_j)$라 두면 이는 $X$의 준콤팩트 열린집합이다.
$\mathcal{F}$가 $X_s$를 따라 순수하므로
$$
E = \{t \in S \mid \text{Ass}_{X_t}(\mathcal{F}_t) \subset W \}.
$$
는 $s$의 모든 일반화를 포함한다. More on Morphisms, Lemma
\ref{more-morphisms-lemma-relative-assassin-constructible}에 의해 $E$는
$S$의 구성가능 부분집합이다. 앞에서
$\Spec(\mathcal{O}_{S, s}) \subset E$임을 보았다. Morphisms, Lemma
\ref{morphisms-lemma-constructible-containing-open}에 의해 $E$는 $s$의
열린근방을 포함한다. 따라서 $S$를 축소한 뒤 $E = S$라고 가정해도
된다. Lemma \ref{flat-lemma-localize-flat-dimension-n}에 의해
$X = \coprod Y_j$ 및 $\mathcal{F} = \coprod \mathcal{F}_j$에 결부된
함자 $F_n$에 대하여 보조정리를 증명하면 충분하다.
% P05-EN-CH38-0241: j번째 함자의 층을 F_i로 쓰는 원문의 첨자를 그대로 보존했다.
$F_{j, n}$이 $Y_j \to S$와 층 $\mathcal{F}_i$에 대한 함자를 나타내면
$F_n = \prod F_{j, n}$이다. 따라서 각 $F_{j, n}$이 어떤 유한표현
단사 사상 $Z_{j, n} \to S$로 표현가능함을 증명하면 충분하다. 실제로
$$
Z_n = Z_{1, n} \times_S \ldots \times_S Z_{m, n}
$$
이다. 이로써 정리는 Lemma
\ref{flat-lemma-flat-dimension-n-representable}에서 다룬 특수한 경우로
환원되었다.
\end{proof}

\noindent
다음 보조정리에서 이 정리가 보편 평탄화의 관점에서 무엇을 뜻하는지
명시한다.

\begin{lemma}
\label{flat-lemma-when-universal-flattening}
$f : X \to S$를 스킴들의 사상이라 하고, $\mathcal{F}$를 준연접
$\mathcal{O}_X$-가군이라 하자.
\begin{enumerate}
\item $f$가 유한표현이고, $\mathcal{F}$가 유한표현
$\mathcal{O}_X$-가군이며, $\mathcal{F}$가 $S$에 대하여 순수하면
$\mathcal{F}$의 보편 평탄화 $S' \to S$가 존재한다. 더 나아가
$S' \to S$는 유한표현 단사 사상이다.
\item $f$가 유한표현이고 $X$가 $S$에 대하여 순수하면 $X$의 보편
평탄화 $S' \to S$가 존재한다. 더 나아가 $S' \to S$는 유한표현 단사
사상이다.
\item $f$가 고유 유한표현이고 $\mathcal{F}$가 유한표현
$\mathcal{O}_X$-가군이면 $\mathcal{F}$의 보편 평탄화
$S' \to S$가 존재한다. 더 나아가 $S' \to S$는 유한표현 단사 사상이다.
\item $f$가 고유 유한표현이면 $X$의 보편 평탄화
$S' \to S$가 존재한다.
\end{enumerate}
\end{lemma}

\begin{proof}
이 명제들은 $F_0 = F_{flat}$에 Theorem
\ref{flat-theorem-flat-dimension-n-representable}을 적용하고, $f$가 고유이면
$\mathcal{F}$가 자동으로 기저 위에서 순수하다는 사실을 쓰면 즉시
따른다. Lemma \ref{flat-lemma-proper-pure}를 보라.
\end{proof}







\section{그로텐디크 존재 정리, IV}
\label{flat-section-existence}

\noindent
이 절에서는 Cohomology of Schemes, Sections
\ref{coherent-section-existence}, \ref{coherent-section-existence-proper},
\ref{coherent-section-existence-proper-support}의 논의를 이어간다. 다음
상황에서 작업하겠다.

\begin{situation}
\label{flat-situation-existence}
여기서는 전이사상이 전사이고 그 핵이 국소적으로 멱영인 환들의 역계
$(A_n)$이 주어져 있다. $A = \lim A_n$이라 두자. $A$ 위에서 분리이고
유한표현인 스킴 $X$가 주어져 있다.
$X_n = X \times_{\Spec(A)} \Spec(A_n)$이라 두고 이를 $X$의 닫힌
부분스킴으로 본다. 더 나아가 계 $(\mathcal{F}_n, \varphi_n)$이 주어져
있다고 가정한다. 여기서 $\mathcal{F}_n$은 $A_n$ 위에서 평탄이고 그
지지가 $A_n$ 위에서 고유인 유한표현 $\mathcal{O}_{X_n}$-가군이며,
$$
\varphi_n :
\mathcal{F}_n \otimes_{\mathcal{O}_{X_n}} \mathcal{O}_{X_{n - 1}}
\longrightarrow
\mathcal{F}_{n - 1}
$$
는 동형이다(이 표기는 Morphisms, Lemma
\ref{morphisms-lemma-i-star-equivalence}의 동치를 사용한다).
\end{situation}

\noindent
우리의 목표는 모든 $n$에 대하여
$\mathcal{F}_n = \mathcal{F} \otimes_{\mathcal{O}_X} \mathcal{O}_{X_n}$
이게 하는 $X$ 위의 준연접 층 $\mathcal{F}$를 찾을 수 있는지 알아보는
것이다.

\begin{lemma}
\label{flat-lemma-compute-what-it-should-be}
Situation \ref{flat-situation-existence}에서
$$
K = R\lim_{D_\QCoh(\mathcal{O}_X)}(\mathcal{F}_n) =
DQ_X(R\lim_{D(\mathcal{O}_X)}\mathcal{F}_n)
$$
를 생각하자. 그러면 $K$는 $D^b_{\QCoh}(\mathcal{O}_X)$에 속하고,
실제로 $K$의 영이 아닌 코호몰로지 층들은 차수 $\geq 0$에만 있다.
\end{lemma}

\begin{proof}
Derived Categories of Schemes, Example
\ref{perfect-example-inverse-limit-quasi-coherent}의 특수한 경우이다.
\end{proof}

\begin{lemma}
\label{flat-lemma-compute-against-perfect}
Situation \ref{flat-situation-existence}에서 $K$를 Lemma
\ref{flat-lemma-compute-what-it-should-be}와 같이 두자.
$D(\mathcal{O}_X)$의 임의의 완전 대상 $E$에 대하여 다음이 성립한다.
\begin{enumerate}
\item $M = R\Gamma(X, K \otimes^\mathbf{L} E)$은 $D(A)$의 완전 대상이고,
$D(A_n)$에 표준 동형
$R\Gamma(X_n, \mathcal{F}_n \otimes^\mathbf{L} E|_{X_n}) =
M \otimes_A^\mathbf{L} A_n$
이 존재한다.
\item $N = R\Hom_X(E, K)$은 $D(A)$의 완전 대상이고, $D(A_n)$에
표준 동형
$R\Hom_{X_n}(E|_{X_n}, \mathcal{F}_n) = N \otimes_A^\mathbf{L} A_n$
이 존재한다.
\end{enumerate}
두 명제 모두에서 $E|_{X_n}$은 $E$를 $X_n$으로 유도 당긴 것을
나타낸다.
\end{lemma}

\begin{proof}
(2)를 증명하자. $E_n = E|_{X_n}$ 및
$N_n = R\Hom_{X_n}(E_n, \mathcal{F}_n)$이라 쓰자.
$R\Hom_{X_n}(-, -)$는
$R\Gamma(X_n, R\SheafHom(-, -))$와 같음을 상기하자. Cohomology,
Section \ref{cohomology-section-global-RHom}을 보라. 따라서 Derived
Categories of Schemes, Lemma
\ref{perfect-lemma-base-change-RHom-perfect}에 의해 $N_n$은 $D(A_n)$의
완전 대상이고 그 구성은 기저변환과 가환한다. 그러므로 $\varphi_n$에서
오는 사상들
$N_n \otimes_{A_n}^\mathbf{L} A_{n - 1} \to N_{n - 1}$
은 동형이다. More on Algebra, Lemma
\ref{more-algebra-lemma-Rlim-perfect-gives-perfect}에 의해
$R\lim N_n$은 완전이고, 이를 다시 $A_n$으로 기저변환하면 $N_n$을
얻는다. 한편 삼각 범주들의 완전함자
$R\Hom_X(E, -) : D_\QCoh(\mathcal{O}_X) \to D(A)$
는 곱과 가환하고 따라서 유도 극한과도 가환한다. 그러므로
$$
R\Hom_X(E, K) =
R\lim R\Hom_X(E, \mathcal{F}_n) =
R\lim R\Hom_X(E_n, \mathcal{F}_n) =
R\lim N_n
$$
이다. 이로써 (2)를 증명했다. (1)은 Cohomology, Lemma
\ref{cohomology-lemma-dual-perfect-complex}을 사용하여 (2)로
옮겨 적으면 얻어진다.
\end{proof}

\begin{lemma}
\label{flat-lemma-relative-pseudo-coherence}
Situation \ref{flat-situation-existence}에서 $K$를 Lemma
\ref{flat-lemma-compute-what-it-should-be}와 같이 두자. 그러면 $K$는
$A$에 대하여 유사연접이다.
\end{lemma}

\begin{proof}
% P05-EN-CH38-0242: "Combinging"이라는 원문의 철자 오류를 문맥에 맞게 번역했다.
Lemma \ref{flat-lemma-compute-against-perfect}와 Derived Categories of Schemes,
Lemma \ref{perfect-lemma-perfect-enough}을 종합하면
$D(\mathcal{O}_X)$의 모든 유사연접 대상 $E$에 대하여
$R\Gamma(X, K \otimes^\mathbf{L} E)$가 $D(A)$에서 유사연접임을 알 수
있다. 따라서 이 보조정리는 More on Morphisms, Lemma
\ref{more-morphisms-lemma-characterize-pseudo-coherent}에서 따른다.
\end{proof}

\begin{lemma}
\label{flat-lemma-compute-over-affine}
Situation \ref{flat-situation-existence}에서 $K$를 Lemma
\ref{flat-lemma-compute-what-it-should-be}와 같이 두자. 임의의 준콤팩트
열린집합 $U \subset X$에 대하여
$$
R\Gamma(U, K) \otimes_A^\mathbf{L} A_n =
R\Gamma(U_n, \mathcal{F}_n)
$$
이 $D(A_n)$에서 성립한다. 여기서 $U_n = U \cap X_n$이다.
\end{lemma}

\begin{proof}
$n$을 고정하자. Derived Categories of Schemes, Lemma
\ref{perfect-lemma-computing-sections-as-colim}에 의해 다음을 만족하는
$X$ 위의 완전 복합체들의 계 $E_m$이 존재한다.
$R\Gamma(U, K) = \text{hocolim} R\Gamma(X, K \otimes^\mathbf{L} E_m)$.
실제로 이 공식은 $K$뿐 아니라 $D_\QCoh(\mathcal{O}_X)$의 모든 대상에
대하여 성립한다. 이를 $\mathcal{F}_n$에 적용하면
\begin{align*}
R\Gamma(U_n, \mathcal{F}_n)
& =
R\Gamma(U, \mathcal{F}_n) \\
& =
\text{hocolim}_m R\Gamma(X, \mathcal{F}_n \otimes^\mathbf{L} E_m) \\
& =
\text{hocolim}_m R\Gamma(X_n, \mathcal{F}_n \otimes^\mathbf{L} E_m|_{X_n})
\end{align*}
을 얻는다. Lemma \ref{flat-lemma-compute-against-perfect}와
$- \otimes_A^\mathbf{L} A_n$이 호모토피 여극한과 가환한다는 사실을
쓰면 결과를 얻는다.
\end{proof}

\begin{lemma}
\label{flat-lemma-finitely-presented}
Situation \ref{flat-situation-existence}에서 $K$를 Lemma
\ref{flat-lemma-compute-what-it-should-be}와 같이 두자.
% P05-EN-CH38-0243: "Denote X_0"에서 by가 빠진 원문을 문맥에 맞게 번역했다.
$\Spec(A)$의 닫힌 부분집합
$\Spec(A_1) = \Spec(A_2) = \ldots$ 위에 놓이는 점들로 이루어진 닫힌
부분집합을 $X_0 \subset X$라 하자. 다음을 만족하는, $X_0$을 포함하는
열린집합 $W \subset X$가 존재한다.
\begin{enumerate}
\item $i = 0$이 아니면 $H^i(K)|_W$는 영이다.
\item $\mathcal{F} = H^0(K)|_W$는 유한표현이다.
\item $\mathcal{F}_n = \mathcal{F} \otimes_{\mathcal{O}_X} \mathcal{O}_{X_n}$이다.
\end{enumerate}
\end{lemma}

\begin{proof}
$n \geq 1$을 고정하자. 구성에 의해 $D_\QCoh(\mathcal{O}_X)$에 표준
사상 $K \to \mathcal{F}_n$이 존재하고, 따라서 준연접 층들의 표준 사상
$H^0(K) \to \mathcal{F}_n$이 존재한다. 이것이 (3)의 의미를 설명한다.

\medskip\noindent
$x \in X_0$를 한 점이라 하자. (1), (2), (3)이 성립하게 하는 $x$의
열린근방 $W$를 찾겠다. $X_0$은 준콤팩트이므로 이로써 보조정리가
증명된다. $U \subset X$를 $x$의 아핀 열린근방이라 하고
$U = \Spec(B)$라 쓰자. $P$가 $A$ 위에서 매끄럽게 되는 전사
$P \to B$를 택하자. Lemma \ref{flat-lemma-relative-pseudo-coherence}와 상대
유사연접성의 정의에 의해 $Ri_*K$를 표현하는 유한 자유 $P$-가군들의
위로 유계인 복합체 $F^\bullet$가 존재한다. 여기서
$i : U \to \Spec(P)$는 이 표시에 의해 유도된 닫힌 몰입이다.
$M_n$을 $\mathcal{F}_n|_U$에 대응하는 $B$-가군이라 하자. Lemma
\ref{flat-lemma-compute-over-affine}에 의해
$$
H^i(F^\bullet \otimes_A A_n) =
\left\{
\begin{matrix}
0 & \text{이면} & i \not = 0 \\
M_n & \text{이면} & i = 0
\end{matrix}
\right.
$$
이다.
% P05-EN-CH38-0244: 닫힌 몰입과 복합체의 최대 차수에 같은 i를 쓰는 원문의 기호 재사용을 그대로 보존했다.
$F^i$가 영이 아닌 최대 첨자를 $i$라 하자. $i \leq 0$이면 (1), (2),
(3)이 성립한다. 그렇지 않으면 $i > 0$이고, 사상
$$
F^{i - 1} \to F^i
$$
% P05-EN-CH38-0246: "in the point x"라는 원문의 전치사 오류를 문맥에 맞게 번역했다.
의 $x$에서의 계수가 최대임을 알 수 있다. 따라서 $\Spec(P)$ 안에서
$x$의 한 열린근방 위에서도 그 계수가 최대이다. 그러므로 $P$를 한
주국소화로 바꾼 뒤 표시된 사상이 전사라고 가정해도 된다. $F^i$가
유한 자유이므로 분해 $F^{i - 1} = F' \oplus F^i$를 택할 수 있다.
그러면 $F^\bullet$을 복합체
$$
\ldots \to F^{i - 2} \to F' \to 0 \to \ldots
$$
로 바꿀 수 있고, $i$에 관한 귀납법으로 끝난다.
\end{proof}

\begin{lemma}
\label{flat-lemma-proper-support}
Situation \ref{flat-situation-existence}에서 $K$를 Lemma
\ref{flat-lemma-compute-what-it-should-be}와 같이 두고 $W \subset X$를
Lemma \ref{flat-lemma-finitely-presented}와 같이 두자.
$\mathcal{F} = H^0(K)|_W$라 두자. 그러면 필요하다면 열린집합 $W$를
축소한 뒤 $\mathcal{F}$의 지지는 $A$ 위에서 고유이다.
\end{lemma}

\begin{proof}
$n \geq 1$을 고정하고 $I_n = \Ker(A \to A_n)$이라 하자. More on
Algebra, Lemma \ref{more-algebra-lemma-limit-henselian}에 의해 쌍
$(A, I_n)$은 헨젤이다. $Z \subset W$를 $\mathcal{F}$의 지지라 하자.
$\mathcal{F}$가 유한표현이므로 이는 닫힌 부분집합이다. Lemma
\ref{flat-lemma-finitely-presented}의 (3)에 의해
$Z \times_{\Spec(A)} \Spec(A_n)$은 $\mathcal{F}_n$의 지지와 같고, 따라서
% P05-EN-CH38-0245: 정의된 I_n 대신 정의되지 않은 I를 쓰는 원문의 몫을 그대로 보존했다.
$\Spec(A/I)$ 위에서 고유이다. More on Morphisms, Lemma
\ref{more-morphisms-lemma-split-off-proper-part-henselian}에 의해
$Z = Z_1 \amalg Z_2$로 쓸 수 있다. 여기서 $Z_1, Z_2$는 $Z$ 안에서
열리고 닫혔으며, $Z_1$은 $A$ 위에서 고유이고,
$Z_1 \times_{\Spec(A)} \Spec(A/I_n)$은 $\mathcal{F}_n$의 지지와 같다.
다시 말해 $Z_2$는 $X_0$과 만나지 않는다. 따라서 $W$를
$W \setminus Z_2$로 바꾸면 보조정리를 얻는다.
\end{proof}

\begin{lemma}
\label{flat-lemma-monomorphism-isomorphism}
$A = \lim A_n$이 전이사상이 전사이고 그 핵이 국소적으로 멱영인
환들의 계의 극한이라고 하자. $S = \Spec(A)$라 두자.
$T \to S$를 국소적으로 유한형인 단사 사상이라 하자. 모든 $n$에
대하여 $\Spec(A_n) \to S$가 $T$를 통해 분해되면 $T = S$이다.
\end{lemma}

\begin{proof}
$S_n = \Spec(A_n)$이라 두자. $T_0 \subset T$를 분해사상들
$S_n \to T$의 공통 상이라 하자. 그러면 $T_0$은 준콤팩트이다.
$T' \subset T$를 $T_0$을 포함하는 준콤팩트 열린집합이라 하자.
그러면 $S_n \to T$는 $T'$을 통해 분해된다. $T' = S$임을 보일 수
있으면 $T' = T = S$이므로, $T$가 준콤팩트라고 가정해도 된다.

\medskip\noindent
$T$가 준콤팩트라고 가정하자. 이 경우 $T \to S$는 분리이고
준유한이다(Morphisms, Lemma
\ref{morphisms-lemma-monomorphism-loc-finite-type-loc-quasi-finite}).
자리스키 주정리를 More on Morphisms, Lemma
\ref{more-morphisms-lemma-quasi-finite-separated-pass-through-finite}의
형태로 사용하여 $W \to S$가 유한이고 $T \to W$가 열린 몰입인 분해
$T \to W \to S$를 택하자. $W = \Spec(B)$라 쓰자. (유일한) 분해사상들
$S_n \to T$를 $W$로 가는 사상으로 볼 수 있으므로
$$
A \longrightarrow B \longrightarrow \lim A_n = A
$$
를 얻는다. 오른쪽 화살표에서 오는 사상
$h : S = \Spec(A) \to \Spec(B) = W$를 생각하자. 그러면
$$
T \times_{W, h} S
$$
은 모든 $n$에 대하여 $S_n \to S$의 상을 포함하는 $S$의 열린
부분스킴이다. 증명을 끝내려면 어떤 $n \geq 1$에 대하여
$S_n \to S$의 상을 포함하는 임의의 열린집합 $U \subset S$가 $S$와
같음을 보이면 충분하다. 실제로 $(A, \Ker(A \to A_n))$은 헨젤 쌍이고
(More on Algebra, Lemma \ref{more-algebra-lemma-limit-henselian}),
따라서 $S$의 모든 닫힌점은 $S_n \to S$의 상에 포함되므로 이는
성립한다.
\end{proof}

\begin{theorem}[그로텐디크 존재 정리]
\label{flat-theorem-existence}
Situation \ref{flat-situation-existence}에서 다음을 만족하는, $A$ 위에서
평탄이고 그 지지가 $A$ 위에서 고유인 유한표현
$\mathcal{O}_X$-가군 $\mathcal{F}$가 존재한다.
$\mathcal{F}_n = \mathcal{F} \otimes_{\mathcal{O}_X} \mathcal{O}_{X_n}$
이 등식은 모든 $n$에 대하여 사상들 $\varphi_n$과 양립하게 성립한다.
\end{theorem}

\begin{proof}
Lemmas \ref{flat-lemma-compute-what-it-should-be},
\ref{flat-lemma-compute-against-perfect}, \ref{flat-lemma-relative-pseudo-coherence},
\ref{flat-lemma-compute-over-affine}, \ref{flat-lemma-finitely-presented},
\ref{flat-lemma-proper-support}을 적용하면 모든 $\Spec(A_n)$ 위에 놓이는
점들을 포함하는 열린 부분스킴 $W \subset X$와, 그 지지가 $A$ 위에서
고유이고 다음을 만족하는 유한표현 $\mathcal{O}_W$-가군
$\mathcal{F}$를 얻는다.
$\mathcal{F}_n = \mathcal{F} \otimes_{\mathcal{O}_W} \mathcal{O}_{X_n}$
이는 모든 $n \geq 1$에 대하여 성립한다. ($X_n \subset W$이므로 이
등식은 의미가 있다.) Lemma \ref{flat-lemma-proper-pure}에 의해
$\mathcal{F}$는 $\Spec(A)$에 대하여 보편적으로 순수하다. Theorem
\ref{flat-theorem-flat-dimension-n-representable}에 의해(설명은 Lemma
\ref{flat-lemma-when-universal-flattening}을 보라) $\mathcal{F}$의 보편
평탄화 $S' \to \Spec(A)$가 존재하고, 더 나아가 사상
$S' \to \Spec(A)$는 유한표현 단사 사상이다. $\mathcal{F}$를
$\Spec(A_n)$으로 기저변환한 것이 $\mathcal{F}_n$이므로 각 $n$에
대하여 $\Spec(A_n) \to \Spec(A)$는 $S'$을 통해 (유일하게) 분해된다.
Lemma \ref{flat-lemma-monomorphism-isomorphism}에 의해
$S' = \Spec(A)$이다. 이는 $\mathcal{F}$가 $A$ 위에서 평탄임을 뜻한다.
마지막으로 $\mathcal{F}$의 스킴론적 지지 $Z$가 $\Spec(A)$ 위에서
고유이므로 사상 $Z \to X$는 닫혔다. 따라서 직상
$(W \to X)_*\mathcal{F}$는 $W$ 위에 지지되고 원하는 모든 성질을
갖는다.
\end{proof}






\section{그로텐디크 존재 정리, V}
\label{flat-section-existence-derived}

\noindent
이 절에서는 준연접 가군에 대하여 Section \ref{flat-section-existence}에서
사용한 방법을 따라 유도 범주에서 그로텐디크 존재 정리의 유사형을
증명한다. 고전적인 경우는 Cohomology of Schemes, Sections
\ref{coherent-section-existence}, \ref{coherent-section-existence-proper},
\ref{coherent-section-existence-proper-support}에서 논한다. 다음
상황에서 작업하겠다.

\begin{situation}
\label{flat-situation-existence-derived}
여기서는 전이사상이 전사이고 그 핵이 국소적으로 멱영인 환들의 역계
$(A_n)$이 주어져 있다. $A = \lim A_n$이라 두자. $A$ 위에서 고유이고
평탄하며 유한표현인 스킴 $X$가 주어져 있다.
$X_n = X \times_{\Spec(A)} \Spec(A_n)$이라 두고 이를 $X$의 닫힌
부분스킴으로 본다. 더 나아가 계 $(K_n, \varphi_n)$이 주어졌다고
가정한다. 여기서 $K_n$은 $D(\mathcal{O}_{X_n})$의 유사연접 대상이고
$$
\varphi_n : K_n \longrightarrow K_{n - 1}
$$
는 $D(\mathcal{O}_{X_n})$의 사상으로서 $D(\mathcal{O}_{X_{n - 1}})$의
동형
$K_n \otimes_{\mathcal{O}_{X_n}}^\mathbf{L} \mathcal{O}_{X_{n - 1}}
\to K_{n - 1}$
을 유도한다.
\end{situation}

\noindent
더 정확히는
$\varphi_n : K_n \to Ri_{n - 1, *}K_{n - 1}$
이라고 써야 한다. 여기서 $i_{n - 1} : X_{n - 1} \to X_n$은 포함
사상이며, 이 표기에서 조건은 수반사상
$Li_{n - 1}^*K_n \to K_{n - 1}$이 동형이라는 것이다. 우리의 목표는
모든 $n$에 대하여
$K_n = K \otimes_{\mathcal{O}_X}^\mathbf{L} \mathcal{O}_{X_n}$이게 하는
유사연접 대상 $K \in D(\mathcal{O}_X)$를 찾는 것이다(같은 표기의
남용을 사용한다).

\begin{lemma}
\label{flat-lemma-compute-what-it-should-be-derived}
Situation \ref{flat-situation-existence-derived}에서
$$
K = R\lim_{D_\QCoh(\mathcal{O}_X)}(K_n) =
DQ_X(R\lim_{D(\mathcal{O}_X)} K_n)
$$
를 생각하자. 그러면 $K$는 $D^-_{\QCoh}(\mathcal{O}_X)$에 속한다.
\end{lemma}

\begin{proof}
$X$가 준콤팩트이고 준분리이므로 함자 $DQ_X$가 존재한다. Derived
Categories of Schemes, Lemma \ref{perfect-lemma-better-coherator}를 보라.
$DQ_X$는 우수반이므로 곱과 가환하고, 따라서 유도 극한과도 가환한다.
그러므로 보조정리에 적힌 등식이 성립한다.

\medskip\noindent
Derived Categories of Schemes, Lemma
\ref{perfect-lemma-boundedness-better-coherator}에 의해 함자 $DQ_X$의
코호몰로지 차원은 유계이다. 따라서
$R\lim K_n \in D^-(\mathcal{O}_X)$임을 보이면 충분하다. 이를 보이기
위해 $U \subset X$를 아핀 열린집합이라 하자. 그러면 Cohomology,
Lemma \ref{cohomology-lemma-RGamma-commutes-with-Rlim}에 의해 표준 완전열
$$
0 \to
R^1\lim H^{m - 1}(U, K_n) \to H^m(U, R\lim K_n) \to
\lim H^m(U, K_n) \to 0
$$
이 존재한다. $U$가 아핀이고 $K_n$이 유사연접이므로(Derived Categories
of Schemes, Lemma \ref{perfect-lemma-pseudo-coherent}에 의해 그
코호몰로지 층들이 준연접이다), Derived Categories of Schemes, Lemma
\ref{perfect-lemma-affine-compare-bounded}에 의해
$H^m(U, K_n) = H^m(K_n)(U)$이다. 따라서 $K_n$들이 $n$에 무관하게
위로 유계임을 보이면 충분하다.

\medskip\noindent
$K_n$은 유사연접이므로 $K_n \in D^-(\mathcal{O}_{X_n})$이다.
$H^{a_n}(K_n)$이 영이 아니게 하는 최대의 $a_n$을 택하자. 물론
$a_1 \leq a_2 \leq a_3 \leq \ldots$이다. $H^{a_n}(K_n)$은 유한표현
$\mathcal{O}_{X_n}$-가군임에 유의하자(Cohomology, Lemma
\ref{cohomology-lemma-finite-cohomology}). 다음 등식이 성립한다.
$H^{a_n}(K_{n - 1}) =
H^{a_n}(K_n) \otimes_{\mathcal{O}_{X_n}} \mathcal{O}_{X_{n - 1}}$.
$X_{n - 1} \to X_n$은 비후화이므로, 나카야마 보조정리(Algebra, Lemma
\ref{algebra-lemma-NAK})에 의해
$H^{a_n}(K_n) \otimes_{\mathcal{O}_{X_n}} \mathcal{O}_{X_{n - 1}}$
이 영이면 $H^{a_n}(K_n)$도 영이다. 따라서 모든 $n$에 대하여
$a_n = a_{n - 1}$이고 결론이 따른다.
\end{proof}

\begin{lemma}
\label{flat-lemma-compute-against-perfect-derived}
Situation \ref{flat-situation-existence-derived}에서 $K$를 Lemma
\ref{flat-lemma-compute-what-it-should-be-derived}와 같이 두자.
$D(\mathcal{O}_X)$의 임의의 완전 대상 $E$에 대하여
% P05-EN-CH38-0257: 유도 전역단면 대상을 cohomology라고 부르는 원문을 수학적 형에 맞게 나타냈다.
대상
$$
M = R\Gamma(X, K \otimes^\mathbf{L} E)
$$
는 $D(A)$의 유사연접 대상이고 표준 동형
$$
R\Gamma(X_n, K_n \otimes^\mathbf{L} E|_{X_n}) = M \otimes_A^\mathbf{L} A_n
$$
이 $D(A_n)$에 존재한다. 여기서 $E|_{X_n}$은 $E$를 $X_n$으로 유도
당긴 것을 나타낸다.
\end{lemma}

\begin{proof}
$E_n = E|_{X_n}$ 및
$M_n = R\Gamma(X_n, K_n \otimes^\mathbf{L} E|_{X_n})$이라 쓰자. Derived
Categories of Schemes, Lemma
\ref{perfect-lemma-flat-proper-pseudo-coherent-direct-image-general}에 의해
$M_n$은 $D(A_n)$의 유사연접 대상이고 그 구성은 기저변환과 가환한다.
따라서 $\varphi_n$에서 오는 사상들
$M_n \otimes_{A_n}^\mathbf{L} A_{n - 1} \to M_{n - 1}$
은 동형이다. More on Algebra, Lemma
\ref{more-algebra-lemma-Rlim-pseudo-coherent-gives-pseudo-coherent}에 의해
$R\lim M_n$은 유사연접이고, 이를 다시 $A_n$으로 기저변환하면 $M_n$을
얻는다. 한편 삼각 범주들의 완전함자
$R\Gamma(X, -) : D_\QCoh(\mathcal{O}_X) \to D(A)$
는 곱과 가환하고 따라서 유도 극한과도 가환한다. 그러므로
$$
R\Gamma(X, E \otimes^\mathbf{L} K) =
R\lim R\Gamma(X,  E \otimes^\mathbf{L} K_n) =
R\lim R\Gamma(X_n, E_n \otimes^\mathbf{L} K_n) =
R\lim M_n
$$
으로서 원하는 결론을 얻는다.
\end{proof}

\begin{lemma}
\label{flat-lemma-relative-pseudo-coherence-derived}
Situation \ref{flat-situation-existence-derived}에서 $K$를 Lemma
\ref{flat-lemma-compute-what-it-should-be-derived}와 같이 두자. 그러면 $K$는
$X$ 위에서 유사연접이다.
\end{lemma}

\begin{proof}
% P05-EN-CH38-0247: "Combinging"이라는 원문의 철자 오류를 문맥에 맞게 번역했다.
Lemma \ref{flat-lemma-compute-against-perfect-derived}와 Derived Categories of
Schemes, Lemma \ref{perfect-lemma-perfect-enough}을 종합하면
$D(\mathcal{O}_X)$의 모든 유사연접 대상 $E$에 대하여
$R\Gamma(X, K \otimes^\mathbf{L} E)$가 $D(A)$에서 유사연접임을 알 수
있다. 따라서 More on Morphisms, Lemma
\ref{more-morphisms-lemma-characterize-pseudo-coherent}에 의해 $K$는
$A$에 대하여 유사연접이다.
% P05-EN-CH38-0248: "is of flat and of finite presentation"이라는 원문의 중복 전치사를 문맥에 맞게 번역했다.
$X$가 $A$ 위에서 평탄하고 유한표현이므로 이는 $X$ 위에서 유사연접인
것과 같다. More on Morphisms, Lemma
\ref{more-morphisms-lemma-check-relative-pseudo-coherence-on-charts}를 보라.
\end{proof}

\begin{lemma}
\label{flat-lemma-compute-over-affine-derived}
Situation \ref{flat-situation-existence-derived}에서 $K$를 Lemma
\ref{flat-lemma-compute-what-it-should-be-derived}와 같이 두자. 임의의
준콤팩트 열린집합 $U \subset X$에 대하여
$$
R\Gamma(U, K) \otimes_A^\mathbf{L} A_n =
R\Gamma(U_n, K_n)
$$
이 $D(A_n)$에서 성립한다. 여기서 $U_n = U \cap X_n$이다.
\end{lemma}

\begin{proof}
$n$을 고정하자. Derived Categories of Schemes, Lemma
\ref{perfect-lemma-computing-sections-as-colim}에 의해 다음을 만족하는
$X$ 위의 완전 복합체들의 계 $E_m$이 존재한다.
$R\Gamma(U, K) = \text{hocolim} R\Gamma(X, K \otimes^\mathbf{L} E_m)$.
실제로 이 공식은 $K$뿐 아니라 $D_\QCoh(\mathcal{O}_X)$의 모든 대상에
대하여 성립한다. 이를 $K_n$에 적용하면
\begin{align*}
R\Gamma(U_n, K_n)
& =
R\Gamma(U, K_n) \\
& =
\text{hocolim}_m R\Gamma(X, K_n \otimes^\mathbf{L} E_m) \\
& =
\text{hocolim}_m R\Gamma(X_n, K_n \otimes^\mathbf{L} E_m|_{X_n})
\end{align*}
을 얻는다. Lemma \ref{flat-lemma-compute-against-perfect-derived}와
$- \otimes_A^\mathbf{L} A_n$이 호모토피 여극한과 가환한다는 사실을
쓰면 결과를 얻는다.
\end{proof}

\begin{theorem}[유도 그로텐디크 존재 정리]
\label{flat-theorem-existence-derived}
Situation \ref{flat-situation-existence-derived}에서 모든 $n$에 대하여
사상들 $\varphi_n$과 양립하게
$K_n = K \otimes_{\mathcal{O}_X}^\mathbf{L} \mathcal{O}_{X_n}$이게 하는
$D(\mathcal{O}_X)$의 유사연접 대상 $K$가 존재한다.
\end{theorem}

\begin{proof}
Lemmas \ref{flat-lemma-compute-what-it-should-be-derived},
\ref{flat-lemma-compute-against-perfect-derived},
\ref{flat-lemma-relative-pseudo-coherence-derived}를 적용하여
$D(\mathcal{O}_X)$의 유사연접 대상 $K$를 얻는다. Lemma
\ref{flat-lemma-compute-over-affine-derived}에서 아핀 열린집합들을 택하면
$K$를 $X_n$으로 제한한 것이 $K_n$임이 즉시 따른다.
\end{proof}

\begin{remark}
\label{flat-remark-correct-generality}
이 절의 결과는 일반화할 수 있다. $X \to \Spec(A)$가 분리이고
유한표현이며, $K_n$이 $A_n$ 위에서 고유인 $X_n$의 닫힌 부분집합에
지지되는 $A_n$에 대하여 유사연접인 대상이라고만 가정해도 아마
성립할 것이다. 결과는 $A$ 위에서 고유인 닫힌 부분집합에 지지되는,
$A$에 대하여 유사연접인 대상 $K$가 될 것이다. 이것이 필요해지면
여기에 정확한 명제를 정식화하고 증명할 것이다.
\end{remark}






\section{블로업과 평탄성}
\label{flat-section-blowup-flat}

\noindent
이 절에서는 “블로업 뒤에 강변환이 평탄해진다”는 형태의 결과에 관한
논의를 이어간다. More on Algebra, Section
\ref{more-algebra-section-blowup-flat}과 Divisors, Section
\ref{divisors-section-blowup-flat}을 보라. 이 절에서는 다음의
(어느 정도 표준적인) 표기를 사용한다. $X \to S$가 스킴들의 사상이고,
$\mathcal{F}$가 $X$ 위의 준연접 가군이며, $T \to S$가 스킴들의
사상이면, $\mathcal{F}_T$는 $\mathcal{F}$를 기저변환
$X_T = X \times_S T$로 당긴 것을 나타낸다.

\begin{remark}
\label{flat-remark-successive-blowups}
$S$를 준콤팩트 준분리 스킴이라 하고, $f : X \to S$를 스킴들의
사상이라 하자. $\mathcal{F}$를 $X$ 위의 준연접 가군이라 하고,
$U \subset S$를 준콤팩트 열린 부분스킴이라 하자. $U$-허용 블로업
$S' \to S$가 주어지면 $X'$은 $X$의 강변환을, $\mathcal{F}'$은
$\mathcal{F}$의 강변환을 나타내게 하고, 후자를 $X'$ 위의 준연접
가군으로 본다(Divisors, Lemma
\ref{divisors-lemma-strict-transform}을 사용한다). $P$를 $U$-허용
블로업에 대한 (위와 같은) 강변환 아래에서 안정적인
$\mathcal{F}/X/S$의 성질이라 하자. 이 절의 일반적인 문제는 다음과
같다. $\mathcal{F}/X/S$에 관한 보조 조건 아래에서 강변환
$\mathcal{F}'/X'/S'$가 $P$를 갖게 하는 $U$-허용 블로업
$S' \to S$가 존재함을 보여라.

\medskip\noindent
일반적인 전략은 $U$-허용 블로업들의 합성이 $U$-허용 블로업이라는
사실을 쓰는 것이다. Divisors, Lemma
\ref{divisors-lemma-composition-admissible-blowups}를 보라. 실제로는
더 정확한 Divisors, Lemma
\ref{divisors-lemma-composition-finite-type-blowups}을 Divisors, Lemma
\ref{divisors-lemma-strict-transform-composition-blowups}과 함께
사용한다. 그 결과 $U$-허용 블로업들의 열
$$
S = S_0 \leftarrow S_1 \leftarrow \ldots \leftarrow S_n
$$
을 찾으면 충분하다. 여기서 $\mathcal{F}_0 = \mathcal{F}$ 및
$X_0 = X$라 두고, $\mathcal{F}_i/X_i$는
$\mathcal{F}_{i - 1}/X_{i - 1}$의 강변환이라 두며, 마지막에 성질
$P$를 갖는 $\mathcal{F}_n/X_n/S_n$에 도달해야 한다.

\medskip\noindent
특히 $V(\mathcal{I}) = S \setminus U$인 유한형 준연접 아이디얼 층
$\mathcal{I} \subset \mathcal{O}_S$를 택하자. Properties, Lemma
\ref{properties-lemma-quasi-coherent-finite-type-ideals}를 보라.
$S' \to S$를 $\mathcal{I}$에 관한 블로업이라 하고 $E \subset S'$을
예외 인자라 하자(Divisors, Lemma
\ref{divisors-lemma-blowing-up-gives-effective-Cartier-divisor}). 그러면
문제는 다음 경우로 환원된다.
% P05-EN-CH38-0249: D가 S의 인자인데 지지를 X\U라고 하는 원문의 식을 그대로 보존했다.
그 지지가 $X \setminus U$인 유효 카르티에 인자 $D \subset S$가
존재하는 경우이다. 특히 $U$가 $S$ 안에서 스킴론적으로 조밀하다고
가정해도 된다(Divisors, Lemma
\ref{divisors-lemma-complement-effective-Cartier-divisor}).

\medskip\noindent
$P$가 $S$ 위에서 국소적이라고 가정하자. 즉
$S = \bigcup S_i$가 준콤팩트 열린집합들에 의한 유한 열린덮개이고
각 $\mathcal{F}_{S_i}/X_{S_i}/S_i$에 대하여 $P$가 성립하면
$\mathcal{F}/X/S$에 대하여도 $P$가 성립한다고 하자. 이 경우 위의
일반적인 문제도 $S$ 위에서 국소적이다. 즉 주어진 $s \in S$에 대하여
$\mathcal{F}_W/X_W/W$에 관한 문제가 풀리는 $s$의 준콤팩트 열린근방
$W$를 찾을 수 있으면 $\mathcal{F}/X/S$에 관한 문제도 풀린다. 이는
Divisors, Lemmas \ref{divisors-lemma-extend-admissible-blowups},
\ref{divisors-lemma-dominate-admissible-blowups}에서 따른다.
\end{remark}

\begin{lemma}
\label{flat-lemma-helper-blowup-affine-space}
$R$을 환이라 하고 $f \in R$라 하자. $r\geq 0$을 정수라 하자.
$R \to S$를 환 준동형이라 하고 $M$을 $S$-가군이라 하자. 다음을
가정하자.
\begin{enumerate}
\item $R \to S$는 유한표현이고 평탄이다.
\item 모든 올 환 $S \otimes_R \kappa(\mathfrak p)$는
% P05-EN-CH38-0250: 잉여류체 대신 R 위의 기하적 정역성을 말하는 원문의 기저를 그대로 보존했다.
$R$ 위에서 기하적으로 정역이다.
\item $M$은 유한 $S$-가군이다.
\item $M_f$는 유한표현 $S_f$-가군이다.
\item 모든 $\mathfrak p \in R$, $f \not \in \mathfrak p$에 대하여
$\mathfrak q = \mathfrak pS$라 두면 가군 $M_{\mathfrak q}$는
$S_\mathfrak q$ 위에서 계수 $r$의 자유 가군이다.
\end{enumerate}
그러면 $V(f) = V(I)$이고 다음 성질을 갖는 유한생성 아이디얼
$I \subset R$가 존재한다. 모든 $a \in I$에 대하여
$R' = R[\frac{I}{a}]$라 두면 몫
$$
M' = (M \otimes_R R')/a\text{-거듭제곱 꼬임}
$$
은 $S' = S \otimes_R R'$ 위에서 다음을 만족한다. 모든 소아이디얼
$\mathfrak p' \subset R'$에 대하여 $M'_g$가 $S'_g$ 위에서 계수 $r$의
자유 가군이 되게 하는 $g \in S'$, $g \not \in \mathfrak p'S'$가
존재한다.
\end{lemma}

\begin{proof}
이 보조정리는 More on Algebra, Lemma
\ref{more-algebra-lemma-blowup-module}의 일반화이다. 독자는 먼저 그
증명을 읽어 보기를 권한다.
% P05-EN-CH38-0251: 이 전사의 존재를 유한생성 가정 (3)이 아닌 (1)에 돌리는 원문을 그대로 보존했다.
(1)에 의해 가능한 전사 $S^{\oplus n} \to M$을 택하자. $f$를 가역화한
뒤 $S^{\oplus n}/K \to M$이 동형이게 하는 유한 부분가군
$K \subset \Ker(S^{\oplus n} \to M)$를 택하자. 이는 (4)에 의해
가능하다. $M_1 = S^{\oplus n}/K$라 두고 $M_1$에 대하여 보조정리를
증명할 수 있다고 하자. $I \subset R$가 대응하는 아이디얼이라고 하자.
그러면 $a \in I$에 대하여 사상
$$
M_1' = (M_1 \otimes_R R')/a\text{-거듭제곱 꼬임}
\longrightarrow
M' = (M \otimes_R R')/a\text{-거듭제곱 꼬임}
$$
은 전사이다. 또한 $R'_a = R_f$이므로 $R'$에서 $a$를 가역화한 뒤에는
동형이다. Algebra, Lemma \ref{algebra-lemma-blowup-in-principal}을 보라.
그런데 $a$는 $M'_1$ 위에서 영인자가 아니므로 표시된 사상은 동형이다.
따라서 $M$이 유한표현 $S$-가군인 경우에 보조정리를 증명하면 충분하다.

\medskip\noindent
(3)을 만족하는 유한표현 $S$-가군 $M$이 주어졌다고 가정하자. 그러면
$J = \text{Fit}_r(M) \subset S$는 유한생성 아이디얼이다. Lemma
\ref{flat-lemma-fibres-irreducible-flat-projective-nonnoetherian}에 의해
$S$를 자유 $R$-가군의 직합인자로 쓸 수 있다.
$\bigoplus_{\alpha \in A} R = S \oplus C$. 임의의 원소 $h \in S$를
위의 분해에서 $h = \sum a_\alpha$로 쓰면 $a_\alpha$들을 $h$의
계수들이라고 하자. $I' \subset R$를 $J$의 원소들의 계수들로 이루어진
아이디얼이라 하자. $S$의 한 원소에 의한 곱셈은 $R$-선형사상
$S \to S$를 정의하므로 $I'$은 $J$의 생성원들의 계수들로 생성된다.
즉 $I'$은 유한생성 아이디얼이다. $I = fI'$이 원하는 아이디얼이라고
주장한다.

\medskip\noindent
먼저 $V(f) = V(I)$임을 확인하자. 포함 $V(f) \subset V(I)$는 명백하다.
역으로 $f \not \in \mathfrak p$이면 성질 (5)와 More on Algebra, Lemma
\ref{more-algebra-lemma-fitting-ideal-generate-locally}에 의해
$\mathfrak q =  \mathfrak p S$는 $V(J)$의 원소가 아니다. 따라서
$S \otimes_R \kappa(\mathfrak p)$에서 영으로 가지 않는 $J$의 원소가
존재한다. 그러므로 $\mathfrak p$에 포함되지 않는 $I'$의 원소가
존재하고, 따라서 $\mathfrak p \not \in V(fI') = V(I)$이다.

\medskip\noindent
$a \in I$를 택하고 $R' = R[\frac{I}{a}]$라 두자. 어떤 $a' \in I'$에
대하여 $a = fa'$로 쓸 수 있다. Algebra, Lemmas
\ref{algebra-lemma-affine-blowup}, \ref{algebra-lemma-blowup-add-principal}에
의해 $I' R' = a'R'$이고 $a'$은 $R'$에서 영인자가 아니다.
% P05-EN-CH38-0252: R-대수의 기저변환에서 텐서 기저를 S로 쓰는 원문의 수식을 그대로 보존했다.
$S' = S \otimes_S R'$라 두자.
$JS' = \text{Fit}_r(M \otimes_S S')$의 모든 원소 $g$는 어떤
$c_\alpha \in I'R'$들에 대하여 $g = \sum_\alpha c_\alpha$로 쓸 수
있다. $I'R' = a'R'$이므로 어떤 $c'_\alpha \in R'$에 대하여
$c_\alpha = a'c'_\alpha$로 쓸 수 있고, $S'$에서
$g = (\sum c'_\alpha)a' = g' a'$이다. 더욱이 어떤 $\alpha$에 대하여
$a' = c_\alpha$가 되는 $g_0 \in J$가 존재한다. 이 원소에 대해서는
$S'$에서 $g_0 = g'_0 a'$이며, 여기서 $g'_0$은 $S'$의 단원이다.
$\mathfrak p' \subset R'$를 소아이디얼이라 하고
$\mathfrak q' = \mathfrak p'S'$라 두자. 위에서 본 바에 의해
$JS'_{\mathfrak q'}$는 영인자가 아닌 $a'$이 생성하는 주아이디얼이다.
More on Algebra, Lemma
\ref{more-algebra-lemma-principal-fitting-ideal}에 의해
$M'_{\mathfrak q'}$는 $r$개의 원소로 생성될 수 있다. $M'$은 유한이므로
대응하는 사상 $(S')^{\oplus r} \to M'$이 $g$를 가역화한 뒤 전사가
되게 하는 $m_1, \ldots, m_r \in M'$ 및
$g \in S'$, $g \not \in \mathfrak q'$가 존재한다.

\medskip\noindent
마지막으로 아이디얼 $J' = \text{Fit}_{r - 1}(M')$을 생각하자.
$J'S'_g$는 $m_1, \ldots, m_r$ 사이 관계들의 계수들로 생성됨에
유의하자(피팅 아이디얼과 기저변환의 양립성). 따라서 $J' = 0$임을
보이면 충분하다. More on Algebra, Lemma
\ref{more-algebra-lemma-fitting-ideal-finite-locally-free}를 보라.
$R'_a = R_f$(Algebra, Lemma
\ref{algebra-lemma-blowup-in-principal})이고 $M'_a = M_f$이므로, (5)에
의해 $\mathfrak p'' \subset R'_a$이고
$\mathfrak q'' = \mathfrak p''S'$인 임의의 소아이디얼
$\mathfrak q'' \subset S'$에 대하여 $J'_a$는
% P05-EN-CH38-0260: q''가 S'의 소아이디얼인데 S에서 국소화하는 원문의 식을 그대로 보존했다.
$S_{\mathfrak q''}$에서 영으로 간다. 또한
$S'_a \subset \prod_{\mathfrak q''\text{는 위와 같이}} S'_{\mathfrak q''}$
이다. 이는 Lemma \ref{flat-lemma-base-change-universally-flat}에 의해
$(S'_a)_{\mathfrak p''} \subset S'_{\mathfrak q''}$이기 때문이다.
따라서
% P05-EN-CH38-0253: S'의 아이디얼을 R'_a에 곱하는 원문의 식을 그대로 보존했다.
$J'R'_a = 0$임을 알 수 있다. $a$는 $R'$에서 영인자가 아니므로
$J' = 0$이라는 결론을 얻고 증명이 끝난다.
\end{proof}

\begin{lemma}
\label{flat-lemma-flatten-module-pre}
$S$를 준콤팩트 준분리 스킴이라 하고, $X \to S$를 스킴들의 사상이라
하자. $\mathcal{F}$를 $X$ 위의 준연접 가군이라 하고
$U \subset S$를 준콤팩트 열린집합이라 하자. 다음을 가정하자.
\begin{enumerate}
\item $X \to S$는 아핀이고 유한표현이며 평탄이고, 기하적으로 정역인
올들을 갖는다.
\item $\mathcal{F}$는 유한형 가군이다.
\item $\mathcal{F}_U$는 유한표현이다.
\item $\mathcal{F}$는 $U$의 점들 위에 놓이는 올들의 모든 일반점에서
$S$ 위로 평탄이다.
\end{enumerate}
그러면 다음을 만족하는 $U$-허용 블로업 $S' \to S$와 열린 부분스킴
$V \subset X_{S'}$이 존재한다. (a) $\mathcal{F}$의 강변환
$\mathcal{F}'$을 그 위로 제한한 것은 유한 국소 자유
$\mathcal{O}_V$-가군이고, (b) $V \to S'$은 전사이다.
\end{lemma}

\begin{proof}
보조정리의 가정들을 만족하는 $\mathcal{F}/X/S$와 $U \subset S$가
주어졌다고 하자. $P$가 “$\mathcal{F}$는 올들의 모든 일반점에서
$S$ 위로 평탄이다”라는 성질을 나타내게 하자. 강변환 아래에서 $P$가
보존됨은 명백하다. Divisors, Lemma
\ref{divisors-lemma-strict-transform-flat}과 Morphisms, Lemma
\ref{morphisms-lemma-base-change-module-flat}을 보라. $P$가 $S$ 위에서
국소적임도 명백하다. 따라서 Remark \ref{flat-remark-successive-blowups}의
모든 관찰이 이 보조정리가 제기한 문제에 적용된다.

\medskip\noindent
$u \in U$에 정수
$$
r(u) = \dim_{\kappa(\xi_u)}(\mathcal{F}_{\xi_u} \otimes \kappa(\xi_u))
$$
를 대응시키는 함수 $r : U \to \mathbf{Z}_{\geq 0}$을 생각하자.
여기서 $\xi_u$는 올 $X_u$의 일반점이다. More on Morphisms, Lemma
\ref{more-morphisms-lemma-flat-and-free-at-point-fibre}와 $X_S$의 열린집합의
$S$에서의 상이 열린집합이라는 사실에 의해 $r(u)$는 국소 상수이다.
따라서 $U = U_0 \amalg U_1 \amalg \ldots \amalg U_c$는 열리고 닫힌
부분스킴들의 유한 서로소 합이며, $r$은 $U_i$에서 값 $i$로 상수이다.
Divisors, Lemma
\ref{divisors-lemma-separate-disjoint-opens-by-blowing-up}에 의해 $S$를
두 스킴의 서로소 합으로 분해하는 $U$-허용 블로업을 찾을 수 있는데,
첫째 스킴은 $U_0$을 포함하고 둘째 스킴은
$U_1 \cup \ldots \cup U_c$를 포함한다. 이를 $c - 1$번 더 반복하면
$U_i \subset S_i$이고 $S$가 서로소 합
$S = S_0 \amalg S_1 \amalg \ldots \amalg S_c$으로 분해된다고
가정해도 된다. 따라서 위에서 정의한 함수 $r$이 어떤 값 $r$로
상수라고 가정해도 된다.

\medskip\noindent
Remark \ref{flat-remark-successive-blowups}에 의해 그 지지가
$S \setminus U$인 유효 카르티에 인자 $D \subset S$가 존재한다고
가정해도 된다. Remark \ref{flat-remark-successive-blowups}을 한 번 더
적용하고 Divisors, Lemma
\ref{divisors-lemma-characterize-effective-Cartier-divisor}를 종합하면,
어떤 영인자가 아닌 $f \in R$에 대하여
$S = \Spec(R)$ 및 $D = \Spec(R/(f))$라고 가정해도 됨을 알 수 있다.
이 경우는 Lemma \ref{flat-lemma-helper-blowup-affine-space}에서 다룬다.
\end{proof}

\begin{lemma}
\label{flat-lemma-trick-fitting-ideal}
$A \to C$를 계수 $d$의 유한 국소 자유 환 준동형이라 하자.
$C_h$가 $A$ 위에서 에탈이게 하는 원소 $h \in C$를 택하자.
$J \subset C$를 아이디얼이라 하자. $C/J$를 유한 $A$-가군으로 보아
$I = \text{Fit}_0(C/J)$라 두자. 그러면 어떤 아이디얼
$J' \subset C_h$에 대하여 $IC_h = JJ'$이다. $J$가 유한생성이면
$I$와 $J'$도 유한생성이다.
\end{lemma}

\begin{proof}
피팅 아이디얼의 기본 성질들을 사용하겠다. More on Algebra, Lemma
\ref{more-algebra-lemma-fitting-ideal-basics}를 보라. 그러면 $IC$는
$C/J \otimes_A C$의 피팅 아이디얼이다. $C \to C \otimes_A C$,
$c \mapsto 1 \otimes c$가 단면(곱셈 사상)을 가짐에 유의하자. 가정에
의해 $C \to C \otimes_A C$는 이 단면 아래에서 $\Spec(C_h)$의 상에
속하는 모든 소아이디얼에서 에탈이다. 따라서 곱셈 사상
$C \otimes_A C_h \to C_h$는 에탈이고, 특히 평탄이다. Algebra, Lemma
\ref{algebra-lemma-map-between-etale}을 보라. 그러므로
% P05-EN-CH38-0258: 존재하는 C_h-대수의 이름이 빠진 원문을 표시식에 따라 C'로 나타냈다.
$C_h$-대수가 존재하여 $C_h$-대수로서
$C \otimes_A C_h \cong C_h \oplus C'$이다. Algebra, Lemma
\ref{algebra-lemma-surjective-flat-finitely-presented}를 보라. 따라서
어떤 아이디얼 $I' \subset C'$에 대하여 $C_h$-가군으로서
$(C/J) \otimes_A C_h \cong (C_h/J_h) \oplus C'/I'$이다. 그러므로
$J' = \text{Fit}_0(C'/I')$라 두면 $IC_h = JJ'$이다. 여기서
% P05-EN-CH38-0259: 피팅 아이디얼에 들어가는 C'/I' 대신 C'/J'을 보는 원문의 몫을 그대로 보존했다.
$C'/J'$을 $C_h$-가군으로 본다.
\end{proof}

\begin{lemma}
\label{flat-lemma-push-ideal}
$A \to B$를 에탈 환 준동형이라 하고 $a \in A$를 영인자가 아닌
원소라 하자. $J \subset B$를 $V(J) \subset V(aB)$인 유한형
아이디얼이라 하자. 모든 $\mathfrak q \subset B$에 대하여 다음을
만족하는 $V(I) \subset V(a)$인 유한형 아이디얼 $I \subset A$와
$g \in B$, $g \not \in \mathfrak q$가 존재한다. 어떤 유한형 아이디얼
$J' \subset B_g$에 대하여 $IB_g = JJ'$이다.
\end{lemma}

\begin{proof}
$g \in B$, $g \not \in \mathfrak q$에서의 주국소화로 $B$를 바꾸어도
된다. 따라서 $B$가 표준 에탈이라고 가정해도 된다. Algebra,
Proposition \ref{algebra-proposition-etale-locally-standard}를 보라.
그러므로 $B$가 어떤 차수 $d$의 단항다항식 $f \in A[x]$에 대한
$C = A[x]/(f)$의 국소화라고 가정해도 된다. 어떤 $h \in C$에 대하여
$B = C_h$라고 하자. $B$ 위에서 $J$를 생성하는 원소들
$h_1, \ldots, h_n \in C$를 택하자. 조건 $V(J) \subset V(aB)$는 충분히
큰 어떤 $m$에 대하여 $B$에서 $a^m = \sum b_i h_i$임을 뜻한다.
$h_{n + 1} = a^m$이라 두자. Lemma
\ref{flat-lemma-trick-fitting-ideal}에서처럼
% P05-EN-CH38-0255: 정의한 h_{n+1} 대신 h_{r+1}을 두 번 쓰는 원문의 첨자를 그대로 보존했다.
$I = \text{Fit}_0(C/(h_1, \ldots, h_{r + 1}))$.
로 둔다. 가군 $C/(h_1, \ldots, h_{r + 1})$은 $a^m$에 의해
소멸되므로 $a^{dm} \in I$이고, 이는 $V(I) \subset V(a)$를 함의한다.
\end{proof}

\begin{lemma}
\label{flat-lemma-flatten-module-etale-localize}
$S$를 준콤팩트 준분리 스킴이라 하고 $X \to S$를 스킴들의 사상이라
하자. $\mathcal{F}$를 $X$ 위의 준연접 가군이라 하고
$U \subset S$를 준콤팩트 열린집합이라 하자. 다음 형태의 가환 그림들이
유한 개 존재한다고 가정하자.
$$
\xymatrix{
& X_i \ar[r]_{j_i} \ar[d] & X \ar[d] \\
S_i^* \ar[r] & S_i \ar[r]^{e_i} & S
}
$$
여기서
\begin{enumerate}
\item $e_i : S_i \to S$는 준콤팩트 에탈 사상이고
$S = \bigcup e_i(S_i)$이다.
\item $j_i : X_i \to X$는 에탈 사상이고
$X = \bigcup j_i(X_i)$이다.
\item $S^*_i \to S_i$는 $j_i^*\mathcal{F}$의 강변환
$\mathcal{F}_i^*$가 $S^*_i$ 위에서 평탄이게 하는
$e_i^{-1}(U)$-허용 블로업이다.
\end{enumerate}
그러면 $\mathcal{F}$의 강변환이 $S'$ 위에서 평탄이게 하는
$U$-허용 블로업 $S' \to S$가 존재한다.
\end{lemma}

\begin{proof}
이 보조정리의 가정들이 $U$-허용 블로업 아래에서 보존된다고 주장한다.
실제로 $b : S' \to S$가 준연접 아이디얼 층 $\mathcal{I}$에 관한
$U$-허용 블로업이라고 하자. 또한 $S^*_i \to S_i$가 준연접 아이디얼 층
$\mathcal{J}_i$에 관한 블로업이라고 하자. 그러면 사상들의 모임
$e'_i : S'_i = S_i \times_S S' \to S'$ 및
$j'_i : X_i' = X_i \times_S S' \to X \times_S S'$은 블로업
$S' \to S$에 대한 $\mathcal{F}$의 강변환 $\mathcal{F}'$에 관하여
조건 (1), (2), (3)을 만족한다. 먼저 $S_i'$은 $\mathcal{I}$의 당김에
관한 $S_i$의 블로업임에 유의하자. Divisors, Lemma
\ref{divisors-lemma-flat-base-change-blowing-up}을 보라. 다음으로
$\mathcal{J}_i$의 당김에 관한 $S_i'$의 블로업
$S_i^{\prime *} \to S_i'$을 생각하자. Divisors, Lemma
\ref{divisors-lemma-blowing-up-two-ideals}에 의해 가환 그림
$$
\xymatrix{
S_i^{\prime *} \ar[r] \ar[rd] \ar[d] & S'_i \ar[d] \\
S_i^* \ar[r] & S_i
}
$$
을 얻고, 위 그림의 모든 사상은 블로업이다. 따라서 Divisors, Lemmas
\ref{divisors-lemma-strict-transform-flat},
\ref{divisors-lemma-strict-transform-composition-blowups}에 의해
\begin{align*}
& \text{ }(j'_i)^*\mathcal{F}'\text{의 다음 사상 아래 강변환 }
S_i^{\prime *} \to S_i' \\
= &
\text{ }j_i^*\mathcal{F}\text{의 다음 사상 아래 강변환 }
S_i^{\prime *} \to S_i \\
= &
\text{ }\mathcal{F}_i'\text{의 다음 사상 아래 강변환 }
S_i^{\prime *} \to S_i' \\
= &
\text{ }\mathcal{F}_i^*\text{의 다음 사상을 통한 당김 }
X_i \times_{S_i} S_i^{\prime *} \to X_i
\end{align*}
임을 알 수 있다. 따라서 이는 $S_i^{\prime *}$ 위에서 평탄이다
(Morphisms, Lemma \ref{morphisms-lemma-base-change-module-flat}).
이제 보조정리와 같은 가정 아래에서 $\mathcal{F}$의 강변환이 기저
위에서 평탄이게 하는 $U$-허용 블로업을 찾는 문제에 Remark
\ref{flat-remark-successive-blowups}의 모든 관찰이 적용됨을 알 수 있다.
특히 $S \setminus U$가 유효 카르티에 인자 $D \subset S$의 지지라고
가정해도 된다. Remark \ref{flat-remark-successive-blowups}을 한 번 더
적용하고 Divisors, Lemma
\ref{divisors-lemma-characterize-effective-Cartier-divisor}를 종합하면
어떤 영인자가 아닌 $a \in A$에 대하여
$S = \Spec(A)$ 및 $D = \Spec(A/(a))$라고 가정해도 된다.

\medskip\noindent
$i$와 $s \in S_i$를 택하자. Lemma \ref{flat-lemma-push-ideal}에 의해
다음을 만족하는 열린근방 $s \in W_i \subset S_i$와 유한형 준연접
아이디얼 $\mathcal{I} \subset \mathcal{O}_S$를 찾을 수 있다. 어떤
유한형 준연접 아이디얼 $\mathcal{J}'_i \subset \mathcal{O}_{W_i}$에
대하여
$\mathcal{I} \cdot \mathcal{O}_{W_i} = \mathcal{J}_i \mathcal{J}'_i$이고
$V(\mathcal{I}) \subset V(a) = S \setminus U$이다. $S_i$는
준콤팩트이므로 $S_i$를 이러한 열린집합들의 유한 모임
$W_1, \ldots, W_n$으로 바꾸어도 된다. 따라서 각 $i$에 대하여 어떤
유한형 준연접 아이디얼
$\mathcal{J}'_i \subset \mathcal{O}_{S_i}$에 대해
$\mathcal{I}_i \cdot \mathcal{O}_{S_i} = \mathcal{J}_i \mathcal{J}'_i$인
준연접 아이디얼 층 $\mathcal{I}_i \subset \mathcal{O}_S$가 존재한다고
가정해도 된다. 증명의 첫 문단에서 논한 것처럼 곱
$\mathcal{I}_1 \ldots \mathcal{I}_n$에 관한 $S$의 블로업 $S'$을
생각하자(구성에 의해 이 블로업은 $U$-허용이다). $S' \to S$를
$S_i$로 기저변환한 것은
$$
\mathcal{J}_i \cdot
\mathcal{J}'_i \mathcal{I}_1 \ldots \hat{\mathcal{I}_i} \ldots \mathcal{I}_n
$$
에 관한 블로업이고, 이는 주어진 블로업 $S_i^* \to S_i$를 통해
분해된다(Divisors, Lemma
\ref{divisors-lemma-blowing-up-two-ideals}). 위 그림의 표기에서는
이것이 $S_i^{\prime *} = S_i'$을 뜻한다. 따라서 $S$를 $S'$으로 바꾼
뒤 $j_i^*\mathcal{F}$가 $S_i$ 위에서 평탄인 상황에 도달한다. 그러므로
$j_i^*\mathcal{F}$는 $S$ 위에서 평탄이다. Lemma
\ref{flat-lemma-etale-flat-up-down}을 보라. Morphisms, Lemma
\ref{morphisms-lemma-flat-permanence}에 의해 $\mathcal{F}$는 $S$ 위에서
평탄이다.
\end{proof}

\begin{theorem}
\label{flat-theorem-flatten-module}
$S$를 준콤팩트 준분리 스킴이라 하고 $X$를 $S$ 위의 스킴이라 하자.
$\mathcal{F}$를 $X$ 위의 준연접 가군이라 하고 $U \subset S$를
준콤팩트 열린집합이라 하자. 다음을 가정하자.
\begin{enumerate}
\item $X$는 준콤팩트이다.
\item $X$는 $S$ 위에서 국소적으로 유한표현이다.
\item $\mathcal{F}$는 유한형 가군이다.
\item $\mathcal{F}_U$는 유한표현이다.
\item $\mathcal{F}_U$는 $U$ 위에서 평탄이다.
\end{enumerate}
그러면 $\mathcal{F}$의 강변환 $\mathcal{F}'$이 유한표현
$\mathcal{O}_{X \times_S S'}$-가군이고 $S'$ 위에서 평탄이게 하는
$U$-허용 블로업 $S' \to S$가 존재한다.
\end{theorem}

\begin{proof}
먼저 강변환이 평탄이게 하는 $U$-허용 블로업을 찾을 수 있음을
증명한다. 이 문제는 원천과 목표 위에서 에탈 국소적이다. 정확한
명제는 Lemma \ref{flat-lemma-flatten-module-etale-localize}를 보라. 특히
$S = \Spec(R)$ 및 $X = \Spec(A)$가 아핀이라고 가정해도 된다.
$s \in S$에 대하여 $\mathcal{F}_s = \mathcal{F}|_{X_s}$라 쓰자
($\mathcal{F}$를 올로 당긴 것). $X \to S$가 유한형이므로
$d = \max_{s \in S} \dim(\text{Supp}(\mathcal{F}_s))$는 정수이다.
$d$에 관한 귀납법을 쓰겠다.

\medskip\noindent
$x \in X$를 $s \in S$ 위에 놓이고
$\dim_x(\text{Supp}(\mathcal{F}_s)) = d$인 $X$의 점이라 하자. Lemma
\ref{flat-lemma-elementary-devissage}를 적용하여
$g : X' \to X$, $e : S' \to S$, $i : Z' \to X'$, $\pi : Z' \to Y'$을
얻는다. $Y' \to S'$은 차원 $d$의 기하적으로 기약인 올들을 갖는
아핀 스킴들의 매끄러운 사상임에 유의하자. 문제는 에탈 국소적이므로
$g^*\mathcal{F}/X'/S'$에 대하여 정리를 증명하면 충분하다.
$i : Z' \to X'$가 유한표현 닫힌 몰입이고 강변환은 아핀 직상과
가환하므로(Divisors, Lemma
\ref{divisors-lemma-strict-transform-affine}) $\mathcal{G}$에 대한
평탄화 결과를 증명하면 충분하다. $\pi$는 유한이고 따라서 아핀이므로
$\pi_*\mathcal{G}/Y'/S'$에 대한 평탄화 결과를 증명하면 충분하다.
따라서 $X \to S$가 차원 $d$의 기하적으로 기약인 올들을 갖는 아핀
스킴들의 매끄러운 사상이라고 가정해도 된다.

\medskip\noindent
다음으로 Lemma \ref{flat-lemma-flatten-module-pre}와 같은 블로업을 적용한다.
그러면 $\mathcal{F}|_V$가 유한 국소 자유이고 $S$로 전사하는
열린집합 $V \subset X$가 존재하는 상황에 도달한다.
$\xi \in X$를 $X_s$의 일반점이라 하고
$r = \dim_{\kappa(\xi)} \mathcal{F}_\xi \otimes \kappa(\xi)$라 하자.
동형 $\kappa(\xi)^{\oplus r} \to
\mathcal{F}_\xi \otimes \kappa(\xi)$를 유도하는 사상
$\alpha : \mathcal{O}_X^{\oplus r} \to \mathcal{F}$를 택하자.
$\mathcal{F}$는 $V$ 위에서 국소 자유이므로 $\alpha$가 동형인
$\xi$의 열린근방 $W$를 찾을 수 있다. $W \to S$가 전사가 되도록
$S$를 $s$의 아핀 열린근방으로 축소하자. $\mathcal{F}$가 $A$-가군
$N$에 결부된 준연접 가군이라고 하자. $\mathcal{F}$는 올들의 모든
일반점에서(실제로 $W$의 모든 점에서) $S$ 위로 평탄이므로
$$
\alpha_\mathfrak p : A_\mathfrak p^{\oplus r} \to N_\mathfrak p
$$
는 $R$의 모든 소아이디얼 $\mathfrak p$에 대하여 보편 단사이다. Lemma
\ref{flat-lemma-induction-step}을 보라. 따라서 $\alpha$는 보편 단사이다.
Algebra, Lemma
\ref{algebra-lemma-universally-injective-check-stalks}를 보라.
$\mathcal{H} = \Coker(\alpha)$라 두자. Divisors, Lemma
\ref{divisors-lemma-strict-transform-universally-injective}에 의해
$U$-허용 블로업 $S' \to S$가 주어지면 $\mathcal{F}'$과
$\mathcal{H}'$의 강변환들은 완전열
$$
0 \to \mathcal{O}_{X \times_S S'}^{\oplus r} \to \mathcal{F}'
\to \mathcal{H}' \to 0
$$
을 이룬다. 따라서 Lemma \ref{flat-lemma-induction-step}은 한 점 $x'$에서
$\mathcal{F}'$이 평탄일 필요충분조건이 그 점에서 $\mathcal{H}'$이
평탄인 것임도 보여 준다. 특히 $\mathcal{H}_U$는 $U$ 위에서 평탄이고
$\mathcal{H}_U$는 유한표현 가군이다. $\mathcal{H}$에 귀납가정을 적용하면 강변환
$\mathcal{H}'$이 원하는 대로 평탄이게 하는 $U$-허용 블로업이
존재함을 알 수 있다.

\medskip\noindent
정리의 증명을 끝내려면 (필요하다면 또 다른 $U$-허용 블로업 뒤에)
$\mathcal{F}'$이 유한표현 가군임을 보여야 한다.
$U \subset S$가 스킴론적으로 조밀하다고 가정해도 되므로 이는 Lemma
\ref{flat-lemma-flat-finite-type-finitely-presented-over-dense-open}에서
따른다(Remark \ref{flat-remark-successive-blowups}의 셋째 문단을 보라).
이로써 정리의 증명이 끝난다.
\end{proof}





\section{응용}
\label{flat-section-applications}

\noindent
이 절에서는 위의 결과 몇 가지를 응용한다.

\begin{lemma}
\label{flat-lemma-flat-after-blowing-up}
$S$를 준콤팩트 준분리 스킴이라 하고 $X$를 $S$ 위의 스킴이라 하자.
$U \subset S$를 준콤팩트 열린집합이라 하자. 다음을 가정하자.
\begin{enumerate}
\item $X \to S$는 유한형이고 준분리이다.
\item $X_U \to U$는 평탄이고 국소적으로 유한표현이다.
\end{enumerate}
그러면 $X$의 강변환이 $S'$ 위에서 평탄이고 유한표현이게 하는
$U$-허용 블로업 $S' \to S$가 존재한다.
\end{lemma}

\begin{proof}
가정에 의해 $X \to S$는 준콤팩트이고 준분리이므로 블로업
$S' \to S$에 관한 $X$의 강변환도 준콤팩트이고 준분리이다. 따라서
보조정리를 증명하려면 강변환이 평탄이고 국소적으로 유한표현이게 하는
$U$-허용 블로업을 찾으면 충분하다. $X = W_1 \cup \ldots \cup W_n$을
유한 아핀 열린 덮개라 하자. $W_i$의 강변환이 평탄이고 국소적으로
유한표현이게 하는 $U$-허용 블로업 $S_i \to S$를 각각 찾을 수 있다면,
모든 $S_i \to S$를 지배하면서 원하는 성질을 갖는 $U$-허용 블로업
$S' \to S$가 존재한다(Divisors, Lemma
\ref{divisors-lemma-dominate-admissible-blowups} 및 Remark
\ref{flat-remark-successive-blowups}도 보라). 따라서 $X$가 아핀이라고
가정해도 된다.

\medskip\noindent
$X$가 아핀이라고 가정하자. Morphisms, Lemma
\ref{morphisms-lemma-quasi-affine-finite-type-over-S}에 의해 $S$ 위의
몰입 $j : X \to \mathbf{A}^n_S$를 택할 수 있다. $j$가 $S$ 위의 닫힌
몰입 $i : X \to V$를 유도하게 하는 준콤팩트 열린 부분스킴
$V \subset \mathbf{A}^n_S$를 택하자. Theorem
\ref{flat-theorem-flatten-module}를 $V \to S$와 준연접 가군
$i_*\mathcal{O}_X$에 적용한다. 그러면 $i_*\mathcal{O}_X$의 강변환이
$S'$ 위에서 평탄이고 $\mathcal{O}_{V \times_S S'}$-가군으로서
유한표현이게 하는 $U$-허용 블로업 $S' \to S$를 얻는다. $X'$을
$S' \to S$에 관한 $X$의 강변환이라 하고
$i' : X' \to V \times_S S'$을 유도된 사상이라 하자. 강변환을 취하는
것은 아핀 사상을 따른 직상과 가환하므로(Divisors, Lemma
\ref{divisors-lemma-strict-transform-affine}),
% P05-EN-CH38-0261
$i'_*\mathcal{O}_{X'}$는 $S$ 위에서 평탄이고
$\mathcal{O}_{V \times_S S'}$-가군으로서 유한표현이다. 이로부터
보조정리가 따른다.
\end{proof}

\begin{lemma}
\label{flat-lemma-finite-after-blowing-up}
$S$를 준콤팩트 준분리 스킴이라 하고 $X$를 $S$ 위의 스킴이라 하자.
$U \subset S$를 준콤팩트 열린집합이라 하자. 다음을 가정하자.
\begin{enumerate}
\item $X \to S$는 고유이다.
\item $X_U \to U$는 유한 국소 자유이다.
\end{enumerate}
그러면 $X$의 강변환이 $S'$ 위에서 유한 국소 자유이게 하는
$U$-허용 블로업 $S' \to S$가 존재한다.
\end{lemma}

\begin{proof}
Lemma \ref{flat-lemma-flat-after-blowing-up}에 의해 $X \to S$가 평탄이고
유한표현이라고 가정해도 된다. 필요하다면 $S$를 $U$-허용 블로업으로
바꾸어 $U \subset S$가 스킴론적으로 조밀하다고 가정해도 된다.
% P05-EN-CH38-0262
그러면 Lemma \ref{flat-lemma-proper-flat-finite-over-dense-open}에 의해
$f$는 유한이다. 따라서 Morphisms, Lemma
\ref{morphisms-lemma-finite-flat}에 의해 $f$는 유한 국소 자유이다.
\end{proof}

\begin{lemma}
\label{flat-lemma-zariski-after-blowup}
$\varphi : X \to S$를 유한형 분리 사상이라 하되 $S$가 준콤팩트이고
준분리라고 하자. $\varphi^{-1}U \to U$가 동형이게 하는 준콤팩트
열린집합 $U \subset S$가 주어졌다고 하자. 그러면 $X$의 강변환
$X'$이 $S'$의 열린 부분스킴과 동형이게 하는 $U$-허용 블로업
$S' \to S$가 존재한다.
\end{lemma}

\begin{proof}
Remark \ref{flat-remark-successive-blowups}의 논의를 적용할 수 있다. 따라서
먼저 $U$-허용 블로업 하나를 시행하여 여집합 $S \setminus U$가 유효
카르티에 인자 $D$의 지지라고 가정해도 된다. 특히 $U$는 $S$에서
스킴론적으로 조밀하다. 다음으로 또 다른 $U$-허용 블로업을 시행하여
$X \to S$가 평탄이고 유한표현인 상황에 이를 수 있다. Lemma
\ref{flat-lemma-flat-after-blowing-up}을 보라. 이 경우 결과는 Lemma
\ref{flat-lemma-zariski}에서 따른다.
\end{proof}

\noindent
다음 보조정리는 고유 변형을 블로업으로 지배할 수 있음을 말한다.

\begin{lemma}
\label{flat-lemma-dominate-modification-by-blowup}
$\varphi : X \to S$를 고유 사상이라 하되 $S$가 준콤팩트이고
준분리라고 하자. $\varphi^{-1}U \to U$가 동형이게 하는 준콤팩트
열린집합 $U \subset S$가 주어졌다고 하자. 그러면 $X$를 지배하는
$U$-허용 블로업 $S' \to S$, 즉 블로업 사상의 분해
$S' \to X \to S$가 존재하게 하는 블로업이 존재한다.
\end{lemma}

\begin{proof}
Remark \ref{flat-remark-successive-blowups}의 논의를 적용할 수 있다. 따라서
먼저 $U$-허용 블로업 하나를 시행하여 여집합 $S \setminus U$가 유효
카르티에 인자 $D$의 지지라고 가정해도 된다. 특히 $U$는 $S$에서
스킴론적으로 조밀하다. $X$의 강변환 $X'$이 $S'$의 열린 부분스킴이게
하는 또 다른 $U$-허용 블로업 $S' \to S$를 택하자. Lemma
\ref{flat-lemma-zariski-after-blowup}을 보라. $X' \to S'$는 고유이고
$U \subset S'$는 조밀하므로 $X' = S'$임을 알 수 있다. 몇 가지 세부
사항은 생략한다.
\end{proof}

\begin{lemma}
\label{flat-lemma-get-section-after-blowup}
$S$를 스킴이라 하고 $U \subset W \subset S$를 열린 부분스킴들이라 하자.
$f : X \to W$를 사상이라 하고 $s : U \to X$를
$f \circ s = \text{id}_U$를 만족하는 사상이라 하자. 다음을 가정하자.
\begin{enumerate}
\item $f$는 고유이다.
\item $S$는 준콤팩트이고 준분리이다.
\item $U$와 $W$는 준콤팩트이다.
\end{enumerate}
그러면 $U$-허용 블로업 $b : S' \to S$와 $s$의 연장인 사상
$s' : b^{-1}(W) \to X$가 존재하여
$f \circ s' = b|_{b^{-1}(W)}$를 만족한다.
\end{lemma}

\begin{proof}
$X$를 $s$의 스킴론적 상으로 바꾸어도 되므로 그렇게 하자. 그러면
$X \to W$는 $U$ 위에서 동형이다. Morphisms, Lemma
\ref{morphisms-lemma-scheme-theoretic-image-of-partial-section}을 보라.
Lemma \ref{flat-lemma-dominate-modification-by-blowup}에 의해 $U$-허용 블로업
$W' \to W$와 $s$의 연장 $W' \to X$가 존재한다. Divisors, Lemma
\ref{divisors-lemma-extend-admissible-blowups}을 적용하여 $W' \to W$를
$S$의 $U$-허용 블로업으로 연장하면 증명이 끝난다.
\end{proof}









\section{콤팩트화}
\label{flat-section-compactify}

\noindent
$S$를 준콤팩트 준분리 스킴이라 하자. $S$ 위의 스킴 $X$에 대하여,
$S$ 위에서 고유인 스킴 $\overline{X}$로 가는 준콤팩트 열린 몰입
$X \to \overline{X}$가 존재하면 {\it $S$ 위의 콤팩트화를 갖는다}
또는 {\it $S$ 위에서 콤팩트화 가능하다}고 하겠다. $X$가
$S$ 위의 콤팩트화를 가지면 $X \to S$는 분리이고 유한형이다. 그 역도
참이라는 것이 Nagata의 정리이다. \cite{Lutkebohmert},
\cite{Conrad-Nagata}, \cite{Nagata-1}, \cite{Nagata-2},
\cite{Nagata-3}, \cite{Nagata-4}를 보라. 다음 절에서 이 정리를
증명한다. Theorem \ref{flat-theorem-nagata}를 보라.

\medskip\noindent
$S$를 준콤팩트 준분리 스킴이라 하고 $X \to S$를 스킴들의 유한형
분리 사상이라 하자. {\it $S$ 위에서의 $X$의 콤팩트화들의 범주}를
다음과 같이 정의한다.
\begin{enumerate}
\item 대상은 $\overline{X} \to S$가 고유인 $S$ 위의 열린 몰입
$j : X \to \overline{X}$이다.
\item 사상 $(j' : X \to \overline{X}') \to (j : X \to \overline{X})$은
$f \circ j' = j$를 만족하는 $S$ 위의 스킴 사상
$f : \overline{X}' \to \overline{X}$이다.
\end{enumerate}
$j : X \to \overline{X}$가 콤팩트화이면 $j$는 준콤팩트 열린 몰입이다.
Schemes, Remark \ref{schemes-remark-quasi-compact-and-quasi-separated}를
보라.

\medskip\noindent
{\bf 주의.} 콤팩트화 $j : X \to \overline{X}$의 상이 조밀하다고
가정하지 {\it 않는다}. 따라서 $f : \overline{X}' \to \overline{X}$가
콤팩트화들의 사상이어도 $f^{-1}(j(X)) = j'(X)$가 성립하지 않을 수 있다.

\begin{lemma}
\label{flat-lemma-compactifications-cofiltered}
$S$를 준콤팩트 준분리 스킴이라 하고 $X$를 $S$ 위에서 콤팩트화
가능한 스킴이라 하자.
\begin{enumerate}
\item[(a)] $S$ 위에서의 $X$의 콤팩트화들의 범주는 쌍대여과 범주이다.
\item[(b)] $j(X)$가 $\overline{X}$에서 조밀하고 스킴론적으로 조밀한
콤팩트화 $j : X \to \overline{X}$들로 이루어진 충만 부분범주는
시초적이다(Categories, Definition \ref{categories-definition-initial}).
\item[(c)] $f : \overline{X}' \to \overline{X}$가 $X$의 콤팩트화들의
사상이고 $j'(X)$가 $\overline{X}'$에서 조밀하면
$f^{-1}(j(X)) = j'(X)$이다.
\end{enumerate}
\end{lemma}

\begin{proof}
(a)를 증명하려면 Categories, Definition
\ref{categories-definition-codirected}의 조건 (1), (2), (3)을 확인해야
한다. $X$가 콤팩트화 가능하다고 가정했으므로 조건 (1)이 성립한다.
$j_i : X \to \overline{X}_i$, $i = 1, 2$를 두 콤팩트화라 하자. 그러면
$(j_1, j_2) : X \to \overline{X}_1 \times_S \overline{X}_2$의 스킴론적 상
$\overline{X}$를 생각할 수 있다. 이는 두 $j_i$를 모두 지배하는 셋째
콤팩트화 $j : X \to \overline{X}$를 정한다.
$$
\xymatrix{
(X, \overline{X}_1) & (X, \overline{X}) \ar[l] \ar[r] & (X, \overline{X}_2)
}
$$
따라서 (2)가 성립한다. $f_1, f_2 : \overline{X}_1 \to \overline{X}_2$를
콤팩트화들 $j_i : X \to \overline{X}_i$, $i = 1, 2$ 사이의 두 사상이라
하자. $\overline{X} \subset \overline{X}_1$을 $f_1$과 $f_2$의 등화자라
하자. $\overline{X}_2 \to S$가 분리이므로 $\overline{X}$는
$\overline{X}_1$의 닫힌 부분스킴이고, 따라서 $S$ 위에서 고유이다.
또한 $f_1|_X = f_2|_X = \text{id}_X$이므로 열린 몰입
$X \to \overline{X}$를 얻는다. 닫힌 몰입
$\overline{X} \to \overline{X}_1$이 주는 사상
$(X \to \overline{X}) \to (j_1 : X \to \overline{X}_1)$은 $f_1$과
$f_2$를 등화한다. 이로써 조건 (3)이 증명된다.

\medskip\noindent
(b)의 증명. $j : X \to \overline{X}$를 콤팩트화라 하자.
$\overline{X}'$이 $\overline{X}$ 안에서 $X$의 스킴론적 폐포를 나타내면,
Morphisms, Lemma \ref{morphisms-lemma-quasi-compact-immersion}에 의해
$X$는 $\overline{X}'$에서 조밀하고 스킴론적으로 조밀하다. 이로써
Categories, Definition \ref{categories-definition-initial}의 첫째 조건이
증명된다. $X$의 콤팩트화들의 범주가 쌍대여과 범주임을 이미 보였으므로
Categories, Definition \ref{categories-definition-initial}의 둘째 조건은
첫째 조건에서 따른다(이 범주론 연습문제의 해답은 생략한다).

\medskip\noindent
(c)의 증명. $\overline{X}'$을 $j'(X)$의 스킴론적 폐포로 바꾸면
(이는 바탕 위상공간을 바꾸지 않는다), 결과는 Morphisms, Lemma
\ref{morphisms-lemma-scheme-theoretic-image-of-partial-section}에서 따른다.
\end{proof}

\noindent
$X$를 변화시키면서 모든 콤팩트화들의 범주도 생각할 수 있다. 이 범주를
내부에서 동형을 유도하는 사상들의 집합에 관하여 국소화하면 $S$ 위의
콤팩트화 가능한 스킴들의 범주와 동치임이 밝혀진다.

\begin{lemma}
\label{flat-lemma-compactifyable}
$S$를 준콤팩트 준분리 스킴이라 하자. $f : X \to Y$를 $S$ 위의 스킴
사상이라 하되 $Y$는 $S$ 위에서 분리이고 유한형이며 $X$는 $S$ 위에서
콤팩트화 가능하다고 하자. 그러면 $X$는 $Y$ 위의 콤팩트화를 갖는다.
\end{lemma}

\begin{proof}
$j : X \to \overline{X}$를 $S$ 위에서의 $X$의 콤팩트화라 하자.
$\overline{X}'$을 $(j, f) : X \to \overline{X} \times_S Y$의 스킴론적
상이라 하자. $\overline{X} \times_S Y \to Y$는
$\overline{X} \to S$의 기저변환이므로 고유이고, 따라서
$\overline{X}' \to Y$도 고유이다. 한편 $Y$는 $S$ 위에서 분리이므로
사상 $(1, f) : X \to X \times_S Y$는 닫힌 몰입이다(Schemes, Lemma
\ref{schemes-lemma-semi-diagonal}). 따라서 사영
$\overline{X} \times_S Y \to \overline{X}$의 ``부분 단면''
$s = (j, f)$에 Morphisms, Lemma
\ref{morphisms-lemma-scheme-theoretic-image-of-partial-section}을 적용하면
$X \to \overline{X}'$은 열린 몰입이다.
\end{proof}

\noindent
$S$를 준콤팩트 준분리 스킴이라 하자. {\it 콤팩트화들의 범주}를 다음
범주로 정의한다. 그 대상은 $\overline{X}$가 $S$ 위에서 고유인 스킴이고
$X \subset \overline{X}$가 준콤팩트 열린집합인 쌍
$(X, \overline{X})$이며, 그 사상은 $S$ 위의 스킴 사상들의 가환 그림
$$
\xymatrix{
X \ar[d] \ar[r]_f & Y \ar[d] \\
\overline{X} \ar[r]^{\overline{f}} & \overline{Y}
}
$$
이다.

\begin{lemma}
\label{flat-lemma-right-multiplicative-system}
$S$를 준콤팩트 준분리 스킴이라 하자. $u$가 동형인 사상
$(u, \overline{u}) : (X', \overline{X}') \to (X, \overline{X})$들의 모임은
콤팩트화들의 범주에서 화살표들의 오른쪽 곱셈계를 이룬다(Categories,
Definition \ref{categories-definition-multiplicative-system}).
\end{lemma}

\begin{proof}
공리 RMS1은 자명하게 확인된다. RMS2가 성립함을 확인하자. 그림
$$
\xymatrix{
& (X', \overline{X}') \ar[d]_{(u, \overline{u})} \\
(Y, \overline{Y}) \ar[r]^{(f, \overline{f})} & (X, \overline{X})
}
$$
이 주어졌고 $u : X' \to X$가 동형이라고 하자. 사영 사상
$v : Y' \to Y$(이는 동형이다)을 갖는 $Y' = Y \times_X X'$라 두자.
또한 사영 사상 $\overline{v} : \overline{Y}' \to \overline{Y}$를 갖는
$\overline{Y}' = \overline{Y} \times_{\overline{X}} \overline{X}'$라 두자.
% P05-EN-CH38-0263
$Y' \to \overline{Y}'$이 열린 몰입임은 명백하다. 그림
$$
\xymatrix{
(Y', \overline{Y}') \ar[r]_{(g, \overline{g})} \ar[d]_{(v, \overline{v})} &
(X', \overline{X}') \ar[d]_{(u, \overline{u})} \\
(Y, \overline{Y}) \ar[r]^{(f, \overline{f})} & (X, \overline{X})
}
$$
은 공리 RMS2가 성립함을 보여 준다.

\medskip\noindent
RMS3이 성립함을 확인하자. 콤팩트화들의 두 사상
$(f, \overline{f}), (g, \overline{g}) :
(X, \overline{X}) \to (Y, \overline{Y})$와 사상
$(v, \overline{v}) : (Y, \overline{Y}) \to (Y', \overline{Y}')$가
주어졌다고 하자. $v$는 동형이고
$(v, \overline{v}) \circ (f, \overline{f}) =
(v, \overline{v}) \circ (g, \overline{g})$라고 하자. 그러면 $f = g$이다.
따라서 $\overline{X}' \subset \overline{X}$을 $\overline{f}$와
$\overline{g}$의 등화자라 하면
$(u, \overline{u}) : (X, \overline{X}') \to (X, \overline{X})$은
콤팩트화들의 범주에서 사상이 되고, 원하는 대로
$(f, \overline{f}) \circ (u, \overline{u}) =
(g, \overline{g}) \circ (u, \overline{u})$를 만족한다.
\end{proof}

\begin{lemma}
\label{flat-lemma-invert-right-multiplicative-system}
$S$를 준콤팩트 준분리 스킴이라 하자. 함자
$(X, \overline{X}) \mapsto X$는 콤팩트화들의 범주를 Lemma
\ref{flat-lemma-right-multiplicative-system}의 오른쪽 곱셈계에 관하여
국소화한 범주(Categories, Lemma
\ref{categories-lemma-right-localization})와 $S$ 위의 콤팩트화 가능한
스킴들의 범주 사이의 동치를 정의한다.
\end{lemma}

\begin{proof}
$\mathcal{C}$를 콤팩트화들의 범주라 하고
$Q : \mathcal{C} \to \mathcal{C}'$를 Categories, Lemma
\ref{categories-lemma-properties-right-localization}의 국소화 함자라 하자.
$\mathcal{D}$를 $S$ 위의 콤팩트화 가능한 스킴들의 범주라 하자. 방금
인용한 보조정리와 곱셈계의 선택으로부터 함자
$\mathcal{C}' \to \mathcal{D}$를 얻음은 명백하다. 이 함자는 명백히
본질적으로 전사이다. $f : X \to Y$가 콤팩트화 가능한 스킴들의
사상이라 하자. $S$ 위에서 고유인 스킴으로 가는 열린 몰입
$Y \to \overline{Y}$를 택한 다음, $\overline{Y}$ 위에서 고유인 스킴
$\overline{X}$로 가는 매장 $X \to \overline{X}$를 택한다(이는
$X \to \overline{Y}$에 Lemma \ref{flat-lemma-compactifyable}을 적용하면
가능하다). 이로써 주어진 사상 $X \to Y$를 만들어 내는 콤팩트화들의
사상 $(X, \overline{X}) \to (Y, \overline{Y})$를 얻는다. 마지막으로
국소화된 범주에서 원천과 목표가 같은 두 사상이 주어졌다고 하자.
이를 각각
$$
a = ((f, \overline{f}) : (X', \overline{X}') \to (Y, \overline{Y}),
(u, \overline{u}) : (X', \overline{X}') \to (X, \overline{X}))
$$
그리고
$$
b = ((g, \overline{g}) : (X'', \overline{X}'') \to (Y, \overline{Y}),
(v, \overline{v}) : (X'', \overline{X}'') \to (X, \overline{X}))
$$
라 하자. 이들이 같은 $S$ 위의 사상 $X \to Y$를 만들어 낸다고, 즉
$f \circ u^{-1} = g \circ v^{-1}$라고 하자. Categories, Lemma
\ref{categories-lemma-morphisms-right-localization}에 의해
$(X', \overline{X}') = (X'', \overline{X}'')$이고
$(u, \overline{u}) = (v, \overline{v})$라고 가정해도 된다. 이 경우
$\overline{f}$와 $\overline{g}$의 등화자
$\overline{X}''' \subset \overline{X}'$을 생각할 수 있다. 사상
$(w, \overline{w}) : (X', \overline{X}''') \to (X', \overline{X}')$은
곱셈 부분집합에 속하고, $(w, \overline{w})$를 앞에 합성하면 국소화된
범주에서 $a = b$임을 알 수 있다.
\end{proof}






\section{Nagata 콤팩트화}
\label{flat-section-compactifications}

\noindent
이 절에서는 Section \ref{flat-section-compactify}에서 예고한 정리를 증명한다.

\begin{lemma}
\label{flat-lemma-check-separated}
$X \to S$를 스킴들의 사상이라 하자. $X = U \cup V$가 열린 덮개이고
$U \to S$와 $V \to S$가 분리이며 $U \cap V \to U \times_S V$가
닫혀 있으면 $X \to S$는 분리이다.
\end{lemma}

\begin{proof}
생략한다. 도움말: $U \times_S U$, $U \times_S V$, $V \times_S U$,
$V \times_S V$가 주는 $X \times_S X$의 열린 덮개를 사용하여
$\Delta : X \to X \times_S X$가 닫혀 있음을 확인하라.
\end{proof}

\begin{lemma}
\label{flat-lemma-separate-disjoint-locally-closed-by-blowing-up}
$X$를 준콤팩트 준분리 스킴이라 하고 $U \subset X$를 준콤팩트
열린집합이라 하자.
\begin{enumerate}
\item $Z_1, Z_2 \subset X$가 $Z_1 \cap Z_2 \cap U = \emptyset$을 만족하는
유한표현 닫힌 부분스킴들이면, $Z_1$과 $Z_2$의 강변환들이 서로소이게
하는 $U$-허용 블로업 $X' \to X$가 존재한다.
\item $T_1, T_2 \subset U$가 서로소인 구성가능 닫힌 부분집합들이면,
$T_1$과 $T_2$의 폐포들이 서로소이게 하는 $U$-허용 블로업
$X' \to X$가 존재한다.
\end{enumerate}
\end{lemma}

\begin{proof}
(1)의 증명. $Z_i \to X$가 유한표현이라는 가정은 $Z_i$의 준연접 아이디얼
층 $\mathcal{I}_i$가 유한형이라는 뜻이다. Morphisms, Lemma
\ref{morphisms-lemma-closed-immersion-finite-presentation}을 보라.
$Z \subset X$를 곱 $\mathcal{I}_1 \mathcal{I}_2$가 절단하는 닫힌
부분스킴이라 하자. $Z \cap U$는 $Z_1 \cap U$와 $Z_2 \cap U$의 서로소
합임에 유의하자. Divisors, Lemma
\ref{divisors-lemma-separate-disjoint-opens-by-blowing-up}에 의해 $Z_1$과
$Z_2$의 강변환들이 서로소이게 하는 $U \cap Z$-허용 블로업
$Z' \to Z$가 존재한다. $Y \subset Z$를 이 블로업의 중심이라 하자.
그러면 $Y \to X$는 $Y \to Z$와 $Z \to X$의 합성이므로 유한표현 닫힌
몰입이다(Divisors, Definition \ref{divisors-definition-admissible-blowup}
및 Morphisms, Lemma
\ref{morphisms-lemma-composition-finite-presentation}). 따라서 $Y$에서의
블로업 $X' \to X$는 $U$-허용 블로업이다. 강변환의 일반 성질에 의해
% P05-EN-CH38-0264
$X' \to X$에 관한 $Z_1, Z_2$의 강변환들은 $Z' \to Z$에 관한
$Z_1, Z_2$의 강변환들과 같다. Divisors, Lemma
\ref{divisors-lemma-strict-transform}을 보라. 이로써 (1)이 증명된다.

\medskip\noindent
(2)의 증명. Properties, Lemma
\ref{properties-lemma-quasi-coherent-finite-type-ideals}에 의해
$T_i = V(\mathcal{J}_i)$가 집합론적으로 성립하게 하는 유한형 준연접
아이디얼 층 $\mathcal{J}_i \subset \mathcal{O}_U$가 존재한다.
Properties, Lemma \ref{properties-lemma-extend}에 의해 $U$로의 제한이
$\mathcal{J}_i$인 유한형 준연접 아이디얼 층
$\mathcal{I}_i \subset \mathcal{O}_X$가 존재한다. 닫힌 부분스킴
$Z_i = V(\mathcal{I}_i)$에 (1)의 결과를 적용하면 된다.
\end{proof}

\begin{lemma}
\label{flat-lemma-blowup-iso-along}
$f : X \to Y$를 준콤팩트 준분리 스킴들의 고유 사상이라 하자.
$V \subset Y$를 준콤팩트 열린집합이라 하고 $U = f^{-1}(V)$라 하자.
$T \subset V$를 닫힌 부분집합이라 하되 $f|_U : U \to V$가 $V$ 안에서
$T$의 어떤 열린근방 위에서 동형이라고 하자. 그러면 $f$의 강변환
$f' : X' \to Y'$이 $Y'$ 안에서 $T$의 폐포의 어떤 열린근방 위에서
동형이게 하는 $V$-허용 블로업 $Y' \to Y$가 존재한다.
\end{lemma}

\begin{proof}
$T' \subset V$를 $f|_U$가 동형인 최대 열린집합의 여집합이라 하자.
그러면 $T', T$는 $V$에서 닫혀 있고 $T \cap T' = \emptyset$이다. $V$는
스펙트럼 위상공간이므로 $T \subset T_c$, $T' \subset T'_c$이고
$T_c \cap T'_c = \emptyset$인 구성가능 닫힌 부분집합 $T_c, T'_c$를
찾을 수 있다($T'$을 포함하고 $T$와 만나지 않는 $V$의 준콤팩트
열린집합 $W$를 택하여 $T_c = V \setminus W$라 둔 다음, $T_c$를
포함하고 $T'$과 만나지 않는 $V$의 준콤팩트 열린집합 $W'$을 택하여
$T'_c = V \setminus W'$라 둔다). Lemma
\ref{flat-lemma-separate-disjoint-locally-closed-by-blowing-up}에 의해 $Y$를
$V$-허용 블로업으로 바꾼 뒤 $T_c$와 $T'_c$의 $Y$ 안의 폐포들이
서로소라고 가정해도 된다. $Y_0 = Y \setminus \overline{T}'_c$,
$V_0 = V \setminus T'_c$, $U_0 = U \times_V V_0$,
$X_0 = X \times_Y Y_0$라 두자. $U_0 \to V_0$는 동형이므로 $X_0$의
강변환 $X'_0$이 $Y'_0$로 동형사상되게 하는 $V_0$-허용 블로업
$Y'_0 \to Y_0$를 찾을 수 있다. Lemma
\ref{flat-lemma-zariski-after-blowup}을 보라. Divisors, Lemma
\ref{divisors-lemma-extend-admissible-blowups}에 의해 $Y_0$로의 제한이
$Y'_0 \to Y_0$인 $V$-허용 블로업 $Y' \to Y$가 존재한다.
$f' : X' \to Y'$이 $f$의 강변환을 나타내면 $f'$은 $Y'_0$ 위에서
동형으로 제한되므로 원하는 결론이 참임을 알 수 있다.
\end{proof}

\begin{lemma}
\label{flat-lemma-find-common-blowups}
$S$를 준콤팩트 준분리 스킴이라 하자. $U \to X_1$과 $U \to X_2$를
$S$ 위의 스킴들의 열린 몰입이라 하고, $U$, $X_1$, $X_2$가 $S$ 위에서
유한형이고 분리라고 가정하자.
% P05-EN-CH38-0265
그러면 가환 그림
$$
\xymatrix{
X_1' \ar[d] \ar[r] & X & X_2' \ar[l] \ar[d] \\
X_1 & U \ar[l] \ar[lu] \ar[u] \ar[ru] \ar[r] & X_2
}
$$
이 존재한다. 이는 $S$ 위의 스킴들의 그림이며, $X_i' \to X_i$는
$U$-허용 블로업이고 $X_i' \to X$는 열린 몰입이며 $X$는 $S$ 위에서
분리이고 유한형이다.
\end{lemma}

\begin{proof}
증명 전체에서 모든 스킴은 $S$ 위에서 분리이고 유한형이라고 하겠다.
특히 이 스킴들은 준콤팩트이고 준분리이며 그 사이의 사상들은
준콤팩트이고 분리이다. Schemes, Sections
\ref{schemes-section-quasi-compact} 및
\ref{schemes-section-separation-axioms}를 보라. 다음 사실을 사용하겠다.
$U \to W$가 그러한 $S$ 위의 스킴들의 몰입이면 $W$ 안에서 $U$의
스킴론적 상 $Z$는 $W$의 닫힌 부분스킴이고, $U \to Z$는 열린 몰입이며,
$U \subset Z$는 스킴론적으로 조밀하고 $U \subset Z$는 위상적으로
조밀하다.
Morphisms, Lemma \ref{morphisms-lemma-quasi-compact-immersion}을 보라.

\medskip\noindent
$X_{12} \subset X_1 \times_S X_2$를 $U \to X_1 \times_S X_2$의
스킴론적 상이라 하자. Morphisms, Lemma
\ref{morphisms-lemma-scheme-theoretic-image-of-partial-section}에 의해
사영 $p_i : X_{12} \to X_i$는 동형 $p_i^{-1}(U) \to U$를 유도한다.
$X_{12}$의 강변환 $X_{12}^i$가 $X_i^i$의 열린 부분스킴과 동형이게 하는
$U$-허용 블로업 $X_i^i \to X_i$를 택하자. Lemma
\ref{flat-lemma-zariski-after-blowup}을 보라. 대응하는 유한형 준연접 아이디얼
층을 $\mathcal{I}_i \subset \mathcal{O}_{X_i}$라 하자.
$X_{12}^i \to X_{12}$는
$p_i^{-1}\mathcal{I}_i \mathcal{O}_{X_{12}}$에서의 블로업임을 상기하라.
Divisors, Lemma \ref{divisors-lemma-strict-transform}을 보라.
$X_{12}'$을 $p_1^{-1}\mathcal{I}_1 p_2^{-1}\mathcal{I}_2
\mathcal{O}_{X_{12}}$에서의 $X_{12}$의 블로업이라 하자. 이것이 뜻하는
바는 Divisors, Lemma \ref{divisors-lemma-blowing-up-two-ideals}를 보라.
특히 가환 그림
$$
\xymatrix{
X_{12}' \ar[d] \ar[r] & X_{12}^2 \ar[d] \\
X_{12}^1 \ar[r] & X_{12}
}
$$
을 얻으며, 모든 사상은 $U$-허용 블로업이다.
$X_{12}^i \subset X_i^i$가 열린집합이므로 $X_{12}' \to X_{12}^i$로
제한되는 $U$-허용 블로업 $X_i' \to X_i^i$를 택할 수 있다. Divisors,
Lemma \ref{divisors-lemma-extend-admissible-blowups}을 보라. 그러면
$X_{12}' \subset X_i'$은 열린 부분스킴이고 그림
$$
\xymatrix{
X_{12}' \ar[d] \ar[r] & X_i' \ar[d] \\
X_{12}^i \ar[r] & X_i^i
}
$$
은 가환하며 세로 화살표들은 블로업이고 가로 화살표들은 열린 몰입이다.
% P05-EN-CH38-0266
$X'_{12} \to X_1' \times_S X_2'$은 몰입이고 고유임에 유의하자
($X'_{12} \to X_{12}$는 고유이고 $X_{12} \to X_1 \times_S X_2$는
닫혀 있으며 $X_1' \times_S X_2' \to X_1 \times_S X_2$는 분리임을
사용하여 Morphisms, Lemma
\ref{morphisms-lemma-image-proper-scheme-closed}를 적용하라). 따라서
$X'_{12} \to  X_1' \times_S X_2'$은 닫힌 몰입이다. 이에 따라 공통 열린
부분스킴 $X_{12}'$을 따라 $X_1'$과 $X_2'$을 붙여 $X$를 정의하면
$X \to S$는 유한형이고 분리이다(Lemma \ref{flat-lemma-check-separated}).
$U$-허용 블로업들의 합성은 $U$-허용 블로업이므로(Divisors, Lemma
\ref{divisors-lemma-composition-admissible-blowups}) 보조정리가 증명된다.
\end{proof}

\begin{lemma}
\label{flat-lemma-replaced-by-strict-transform}
$X \to S$와 $Y \to S$를 스킴들의 사상이라 하고 $U \subset X$를 열린
부분스킴이라 하자. $V \to X \times_S Y$를 준콤팩트 사상이라 하되 첫째
사영과의 합성의 상이 $U$에 들어간다고 하자.
$Z \subset X \times_S Y$를 $V \to X \times_S Y$의 스킴론적 상이라 하고
$X' \to X$를 $U$-허용 블로업이라 하자. 그러면
$V \to X' \times_S Y$의 스킴론적 상은 이 블로업에 관한 $Z$의 강변환이다.
\end{lemma}

\begin{proof}
$Z' \to Z$를 강변환이라 하자. 사상 $Z' \to X'$은 닫힌 몰입인 사상
$Z' \to X' \times_S Y$를 유도한다(정의에 의해 $Z'$은
$X' \times_X Z$의 닫힌 부분스킴이다). 따라서 증명을 끝내려면
$V \to Z'$의 스킴론적 상 $Z''$이 $Z'$임을 보이면 충분하다.
$Z'' \subset Z'$은 $V \to Z'$가 $Z''$을 통하여 분해되는 닫힌
부분스킴임에 유의하자. $V \to X \times_S Y$와
$V \to X' \times_S Y$가 모두 준콤팩트이므로(뒤의 것은 Schemes, Lemma
\ref{schemes-lemma-quasi-compact-permanence}와 고유 사상의 기저변환
$X' \times_S Y \to X \times_S Y$가 분리라는 사실에서 따른다),
Morphisms, Lemma
\ref{morphisms-lemma-quasi-compact-scheme-theoretic-image}에 의해
$Z \cap (U \times_S Y) = Z'' \cap (U \times_S Y)$이다. 따라서 포함 사상
$Z'' \to Z'$은 $Z' \to Z$의 예외 인자 $E$ 밖에서 동형이다. 그러나
$Z'$의 구조층은 $E$ 위에 지지된 영이 아닌 단면을 갖지 않으므로
(강변환의 정의에 의해), 전사
$\mathcal{O}_{Z'} \to \mathcal{O}_{Z''}$은 동형이어야 한다.
\end{proof}

\begin{lemma}
\label{flat-lemma-compactification-dominates}
$S$를 준콤팩트 준분리 스킴이라 하고 $U$를 $S$ 위에서 유한형이고
분리인 스킴이라 하자. $V \subset U$를 준콤팩트 열린집합이라 하자.
$V$가 $S$ 위의 콤팩트화 $V \subset Y$를 가지면, $V \to U$가 고유 사상
$V' \to U$로 연장되게 하는 $V$-허용 블로업 $Y' \to Y$와 열린집합
$V \subset V' \subset Y'$이 존재한다.
\end{lemma}

\begin{proof}
``대각'' 사상 $V \to Y \times_S U$의 스킴론적 상
$Z \subset Y \times_S U$를 생각하자. $Y$를 $V$-허용 블로업으로
바꾸면 $Z$는 이 블로업에 관한 강변환으로 바뀐다. Lemma
\ref{flat-lemma-replaced-by-strict-transform}을 보라. 따라서 Lemma
\ref{flat-lemma-zariski-after-blowup}에 의해 $Z \to Y$가 열린 몰입이라고
가정해도 된다. $V' \subset Y$가 그 상을 나타내면, 사영
$Y \times_S U \to U$는 고유이고 $V' \cong Z$는 $Y \times_S U$의 닫힌
부분스킴이므로 유도된 사상 $V' \to U$가 고유임을 알 수 있다.
\end{proof}

\noindent
다음 보조정리는 뇌터인 경우에 대해서만 서술한다. 준콤팩트 준분리
스킴들에 대한 판도 참이지만 이 절의 주정리에서 자명하게 따를 것이다.

\begin{lemma}
\label{flat-lemma-two-compactifications}
$S$를 뇌터 스킴이라 하고 $U$를 $S$ 위에서 유한형이고 분리인 스킴이라
하자. $U = U_1 \cup U_2$를 열린집합들로 나타내되 $U_1$과 $U_2$가
$S$ 위의 콤팩트화들을 가지며 $U_1 \cap U_2$가 $U$에서 조밀하다고 하자.
그러면 $U$는 $S$ 위의 콤팩트화를 갖는다.
\end{lemma}

\begin{proof}
$i = 1, 2$에 대하여 콤팩트화 $U_i \subset X_i$를 택하자. $U_i$가
$X_i$에서 스킴론적으로 조밀하다고 가정해도 된다. 열린집합
$V_i \subset X_i$와 $\text{id} : U_i \to U_i$의 연장인 고유 사상
$\psi_i : V_i \to U$가 존재한다고 가정해도 된다. Lemma
\ref{flat-lemma-compactification-dominates}를 보라. 그림은 다음과 같다.
$$
\xymatrix{
U_i \ar[r] \ar[d] & V_i \ar[r] \ar[dl]^{\psi_i} & X_i \\
U
}
$$
$\{i, j\} = \{1, 2\}$이면
$Z_i = U \setminus U_j = U_i \setminus (U_1 \cap U_2)$
및
$Z_j = U \setminus U_i = U_j \setminus (U_1 \cap U_2)$.
라고 쓰자. 그러면
$$
U = U_1 \amalg Z_2 = Z_1 \amalg U_2 = Z_1 \amalg (U_1 \cap U_2) \amalg Z_2
$$
이다. $Z_{i, i} \subset V_i$를 $\psi_i$에 의한 $Z_i$의 역상이라 하자.
$\psi_i$는 $Z_i$의 어떤 열린근방 위에서 동형임에 유의하자.
$Z_{i, j} \subset V_i$를 $\psi_i$에 의한 $Z_j$의 역상이라 하자.
$\psi_i : Z_{i, j} \to Z_j$는 고유 사상임에 유의하자. $Z_i$와 $Z_j$는
$U$의 서로소인 닫힌 부분집합들이므로 $Z_{i, i}$와 $Z_{i, j}$는
$V_i$의 서로소인 닫힌 부분집합들이다.

\medskip\noindent
$\overline{Z}_{i, i}$와 $\overline{Z}_{i, j}$를 각각 $X_i$ 안에서
$Z_{i, i}$와 $Z_{i, j}$의 폐포라 하자. $X_i$를 $V_i$-허용 블로업으로
바꾼 뒤 $\overline{Z}_{i, i}$와 $\overline{Z}_{i, j}$가 서로소라고
가정해도 된다. Lemma
\ref{flat-lemma-separate-disjoint-locally-closed-by-blowing-up}을 보라. 이 성질이
$X_1$과 $X_2$ 모두에 대하여 성립한다고 가정한다. $X_i$를 또 다른
$V_i$-허용 블로업으로 바꾸어도 이 성질이 보존됨에 유의하자.

\medskip\noindent
$V_{12} = V_1 \times_U V_2$라 두자. 몰입
$V_{12} \to X_1 \times_S X_2$를 얻는데, 이는 닫힌 몰입
$V_{12} = V_1 \times_U V_2 \to V_1 \times_S V_2$(Schemes, Lemma
\ref{schemes-lemma-fibre-product-after-map})와 열린 몰입
$V_1 \times_S V_2 \to X_1 \times_S X_2$의 합성이다.
$X_{12} \subset X_1 \times_S X_2$를 $V_{12} \to X_1 \times_S X_2$의
스킴론적 상이라 하자. 사영 사상들
$$
p_1 : X_{12} \to X_1
\quad\text{이고}\quad
p_2 : X_{12} \to X_2
$$
은 $X_1$과 $X_2$가 $S$ 위에서 고유이므로 고유이다. $X_1$을
$V_1$-허용 블로업으로 바꾸면 $X_{12}$는 이 블로업에 관한 강변환으로
바뀐다. Lemma \ref{flat-lemma-replaced-by-strict-transform}을 보라.

\medskip\noindent
$\psi : V_{12} \to U$를 합성들
$\psi = \psi_1 \circ p_1|_{V_{12}} = \psi_2 \circ p_2|_{V_{12}}$라 하자.
닫힌 부분스킴
$$
Z_{12, 2} =
(p_1|_{V_{12}})^{-1}(Z_{1, 2}) =
(p_2|_{V_{12}})^{-1}(Z_{2, 2}) =
\psi^{-1}(Z_2) \subset V_{12}
$$
을 생각하자. $\psi_2 : V_2 \to U$는 $Z_2$의 어떤 열린근방 위에서
동형이고 $V_{12} = V_1 \times_U V_2$이므로 사상
$p_1|_{V_{12}} : V_{12} \to V_1$은 $Z_{1, 2}$의 어떤 열린근방 위에서
동형이다. Lemma \ref{flat-lemma-blowup-iso-along}에 의해 $p_1$의 강변환
$p'_1 : X'_{12} \to X'_1$이 $X'_1$ 안에서 $Z_{1, 2}$의 폐포의 어떤
열린근방 위에서 동형이게 하는 $V_1$-허용 블로업 $X_1' \to X_1$이
존재한다. $X_1$을 $X'_1$으로, $X_{12}$를 $X'_{12}$으로 바꾸고 나면
$p_1$이 $\overline{Z}_{1, 2}$의 어떤 열린근방 위에서 동형이라고
가정해도 된다.

\medskip\noindent
앞 문단의 축약은
$$
X_{12} \cap (\overline{Z}_{1, 2} \times_S \overline{Z}_{2, 1}) = \emptyset
$$
임을 말해 준다. 여기서 교차는 $X_1 \times_S X_2$ 안에서 취한다.
% P05-EN-CH38-0267
실제로 $X_{12}$ 안의 역상 $p_1^{-1}(\overline{Z}_{1, 2})$는
$\overline{Z}_{1, 2}$로 동형사상된다. 특히 $Z_{12, 2}$가
$p_1^{-1}(\overline{Z}_{1, 2})$에서 조밀함을 알 수 있다. 따라서 $p_2$는
$p_1^{-1}(\overline{Z}_{1, 2})$를 $\overline{Z}_{2, 2}$ 안으로 보낸다.
$\overline{Z}_{2, 2} \cap \overline{Z}_{2, 1} = \emptyset$이므로 결론을
얻는다.

\medskip\noindent
붙여서 얻는 스킴들
$$
W_i = U \coprod\nolimits_{U_i} (X_i \setminus \overline{Z}_{i, j}),
\quad i = 1, 2
$$
을 생각하자. Lemma \ref{flat-lemma-check-separated}를 적용하여 $W_i \to S$가
분리임을 보자. 먼저 $U \to S$와 $X_i \to S$는 분리이다. 몰입
$U_i \to U \times_S (X_i \setminus \overline{Z}_{i, j})$은 닫혀 있다.
실제로 $u_i \in U_i$이고 $u \in U \setminus U_i$인 임의의 특수화
$u_i \leadsto u$는 고유 사상 $\psi_i : V_i \to U$를 따라 $V_i$ 안의
특수화 $u_i \leadsto v_i$로 유일하게 올릴 수 있고, 이때 $v_i$는
$Z_{i, j}$에 속해야 한다. 따라서 몰입의 상은 닫혀 있고, 따라서 그
몰입은 닫힌 몰입이다.

\medskip\noindent
한편 분수체가 $K$인 $S$ 위의 임의의 값매김환 $A$와 $S$ 위의 임의의
사상 $\gamma : \Spec(K) \to (U_1 \cap U_2)$에 대하여, 어떤 $i$와
$\gamma$의 사상 $h_i : \Spec(A) \to W_i$로의 연장이 존재한다. 실제로
$i = 1, 2$ 각각에 대하여 $S$ 위에서 $X_i$의 고유성에 대한 값매김
판정법에 의해 $\gamma$를 연장하는 사상 $g_i : \Spec(A) \to X_i$가
존재한다. Morphisms, Lemma \ref{morphisms-lemma-characterize-proper}를
보라. 따라서 $i = 1, 2$ 모두에 대하여
$g_i(\mathfrak m_A) \in \overline{Z}_{i, j}$일 때만 문제가 생긴다.
이는 위에서 증명한 $X_{12}$와
$\overline{Z}_{1, 2} \times_S \overline{Z}_{2, 1}$의 교차가 공집합이라는
사실에 모순된다.
% P05-EN-CH38-0268

\medskip\noindent
Lemma \ref{flat-lemma-find-common-blowups}에서와 같은 그림
$$
\xymatrix{
W_1' \ar[d] \ar[r] & W & W_2' \ar[l] \ar[d] \\
W_1 & U \ar[l] \ar[lu] \ar[u] \ar[ru] \ar[r] & W_2
}
$$
을 생각하자. 앞 문단에 의해 모든 실선 그림
$$
\xymatrix{
\Spec(K) \ar[r]_\gamma  \ar[d] & W \ar[d] \\
\Spec(A) \ar@{..>}[ru] \ar[r] & S
}
$$
에서 $\Im(\gamma) \subset U_1 \cap U_2$이면 어떤 $i$와 $\gamma$의 연장
$h_i : \Spec(A) \to W_i$가 존재한다. $W'_i \to W_i$의 고유성에 대한
값매김 판정법을 사용하면 $h_i$를 $h'_i : \Spec(A) \to W'_i$로 올릴 수
있다. 따라서 그림의 점선 화살표가 존재한다. $W$는 $S$ 위에서
분리이므로 이 화살표는 유일하기도 하다. 이에 따라 Morphisms, Lemma
\ref{morphisms-lemma-refined-valuative-criterion-universally-closed}에 의해
$W \to S$는 보편적으로 닫힌 사상이다. $W \to S$는 이미 유한형이고
분리이므로 원하는 결론을 얻는다.
\end{proof}

\begin{theorem}
\label{flat-theorem-nagata}
\begin{reference}
\cite{Lutkebohmert}, \cite{Conrad-Nagata}, \cite{Nagata-1},
\cite{Nagata-2}, \cite{Nagata-3}, \cite{Nagata-4}
를 보라.
\end{reference}
$S$를 준콤팩트 준분리 스킴이라 하고 $X \to S$를 유한형 분리 사상이라
하자. 그러면 $X$는 $S$ 위의 콤팩트화를 갖는다.
\end{theorem}

\begin{proof}
먼저 뇌터인 경우로 축약한다. 독자에게 이 문단을 건너뛸 것을 강하게
권한다. $X' \to S$가 유한표현이고 분리인 닫힌 몰입 $X \to X'$가
존재한다. Limits, Proposition
\ref{limits-proposition-separated-closed-in-finite-presentation}을 보라.
$S$ 위에서 $X'$의 콤팩트화를 찾으면, 그 안에서 $X$의 스킴론적 상을
취하여 $S$ 위에서의 $X$의 콤팩트화를 얻는다. 따라서 $X \to S$가
분리이고 유한표현이라고 가정해도 된다. $S = \lim S_i$를 아핀 전이
사상들을 갖는 뇌터 스킴들의 계의 유향 극한으로 쓸 수 있다. Limits,
Proposition \ref{limits-proposition-approximate}를 보라. $S$로 기저변환하면
$X \to S$가 되는 유한표현 사상 $X_i \to S_i$와 $i$를 택할 수 있다.
Limits, Lemma \ref{limits-lemma-descend-finite-presentation}를 보라. $i$를
늘린 뒤 $X_i \to S_i$가 분리라고 가정해도 된다. Limits, Lemma
\ref{limits-lemma-descend-separated-finite-presentation}을 보라. $S_i$ 위에서
$X_i$의 콤팩트화를 찾을 수 있으면 이를 $S$로 기저변환한 것은 $S$
위에서의 $X$의 콤팩트화이다. 이로써 다음 문단에서 다룰 경우로 축약된다.

\medskip\noindent
$S$가 뇌터라고 가정하자. $U_1 \cap \ldots \cap U_n$이 $X$에서 조밀하게
하는 유한 아핀 열린 덮개 $X = \bigcup_{i = 1, \ldots, n} U_i$를 택할 수
있다. 이는 Properties, Lemma
\ref{properties-lemma-point-and-maximal-points-affine} 및 $X$가 유한 개의
기약 성분을 갖는 준콤팩트 스킴이라는 사실에서 따른다. 각 $i$에 대하여
Morphisms, Lemma \ref{morphisms-lemma-quasi-affine-finite-type-over-S}에 의해
$n_i \geq 0$과 몰입 $U_i \to \mathbf{A}^{n_i}_S$를 택할 수 있다. 따라서
$\mathbf{P}^{n_i}_S$ 안에서 스킴론적 상을 취하면
$i = 1, \ldots, n$에 대하여 $U_i$는 $S$ 위의 콤팩트화를 갖는다. Lemma
\ref{flat-lemma-two-compactifications}를 $(n - 1)$번 적용하면 정리가 참이라는
결론을 얻는다.
\end{proof}






\section{h 위상}
\label{flat-section-h}

\noindent
느슨하게 말하면, 여기서 h 층은 Zariski 위상에 대한 층이면서 유한표현
전사 고유 사상들에 대한 층 성질을 만족하는 것이다. Lemma
\ref{flat-lemma-characterize-sheaf-h}를 보라. 다만 스킴들의 범주 위의 h 위상의
정의가 문헌에 따라 다름은 지적할 만하다.

\medskip\noindent
Voevodsky는 처음에 h 덮개를 합동으로 보편 잠입적인 유한형 사상들의
유한 모임으로 정의했다(Morphisms, Definition
\ref{morphisms-definition-submersive}). \cite[Definition 3.1.2]{Voevodsky}를
보라. 이 정의는 바탕 스킴들의 범주를 고정된 뇌터 기저 스킴 위의 모든
유한형 스킴들로 제한할 때 가장 잘 작동한다. 이 상황에서 Voevodsky는
h 덮개와 ph 덮개를 관련짓는다. ph 위상은 Zariski 덮개들과 고유 전사
사상들에 의해 생성된다. 자세한 내용은 Topologies, Section
\ref{topologies-section-ph}를 보라.

\medskip\noindent
Topologies, Section \ref{topologies-section-V}에서는 V 위상을 연구한다.
$Y$의 점들의 모든 특수화가 $X$의 점들의 어떤 특수화의 상이고 임의의
기저변환 뒤에도 마찬가지이면, 준콤팩트 사상 $X \to Y$는 V 덮개를
정의한다(Topologies, Lemma \ref{topologies-lemma-refine-qcqs-V}). 이 경우
$X \to Y$는 보편 잠입적이다(Topologies, Lemma
\ref{topologies-lemma-V-covering-universally-submersive}). 뇌터가 아닌
스킴들을 다룰 때 V 덮개의 개념은 ``고정된 목표를 갖고 합동으로 보편
잠입적인 사상들의 족''을 훌륭하게 대체한다.

\medskip\noindent
먼저 국소적으로 유한표현인 사상들의 (무한일 수도 있는) 족에 대하여
ph 덮개와 V 덮개의 동치를 증명하겠다. 그런 다음 Stacks project에서
이 족들을 h 덮개의 개념으로 사용한다. 뇌터 스킴들과 유한 족의 경우 이
덮개들은 Voevodsky의 정의에 나오는 것들과 일치한다. Lemma
\ref{flat-lemma-Noetherian-h-covering}을 보라. 고정된 준콤팩트 준분리 스킴
$S$ 위의 유한표현 스킴들의 범주에서 이 덮개들은
\cite[Definition 2.7]{Witt-Grass}의 것과 같은 위상을 정한다.

\begin{lemma}
\label{flat-lemma-equivalence-h-v-locally-finite-presentation}
$\{f_i : X_i \to X\}_{i \in I}$를 고정된 목표를 갖는 스킴 사상들의
족이라 하고 모든 $i$에 대하여 $f_i$가 국소적으로 유한표현이라고 하자.
다음은 동치이다.
\begin{enumerate}
\item $\{X_i \to X\}$는 ph 덮개이다.
\item $\{X_i \to X\}$는 V 덮개이다.
\end{enumerate}
\end{lemma}

\begin{proof}
$U \subset X$를 아핀 열린집합이라 하자. Topologies, Definitions
\ref{topologies-definition-ph-covering} 및
\ref{topologies-definition-V-covering}을 보면 기저변환
$\{X_i \times_X U \to U\}$가 표준 ph 덮개로 세분될 필요충분조건이 표준
V 덮개로 세분될 수 있는 것임을 보이면 충분하다. 따라서 $X$가
아핀이라고 가정해도 되고, $\{X_i \to X\}$가 표준 ph 덮개로 세분될
필요충분조건이 표준 V 덮개로 세분되는 것임을 보여야 한다. 표준 ph
덮개는 표준 V 덮개이므로 다른 함의만 증명하면 충분하다. Topologies,
Lemma \ref{topologies-lemma-standard-ph-standard-V}를 보라.

\medskip\noindent
$X$가 아핀이고 $\{f_i : X_i \to X\}_{i \in I}$가 표준 V 덮개
$\{g_j : Y_j \to X\}_{j = 1, \ldots, m}$로 세분될 수 있다고 하자.
각 $j$에 대하여 $g_j = f_{i_j} \circ h_j$를 만족하는 $i_j$와 사상
$h_j : Y_j \to X_{i_j}$를 택하자. $Y_j$는 아핀이므로 준콤팩트이다.
따라서 각 $j$에 대하여 $\Im(h_j) \subset \bigcup U_{j, k}$가 되게 하는
유한 개의 아핀 열린집합 $U_{j, k} \subset X_{i_j}$를 찾을 수 있다.
그러면 $\{U_{j, k} \to X\}_{j, k}$는 $\{X_i \to X\}$를 세분하며 표준
V 덮개이다(아핀들의 사상들의 유한 족이고 $\{Y_j \to X\}$의 대응하는
성질로부터 값매김환에 대한 올림 성질을 물려받기 때문이다). 이로써
다음 문단에서 다룰 경우로 축약된다.

\medskip\noindent
$\{f_i : X_i \to X\}_{i = 1, \ldots, n}$가 표준 V 덮개이고 $f_i$가
유한표현이라고 하자. $\{X_i \to X\}$가 표준 ph 덮개로 세분될 수 있음을
보여야 한다. 아핀들의 유한표현 사상
$$
X = S \supset S_0 \supset S_1 \supset \ldots \supset S_t = \emptyset
$$
와 같은 일반 평탄성 층화를
$$
\coprod\nolimits_{i = 1, \ldots, n} f_i :
\coprod\nolimits_{i = 1, \ldots, n} X_i
\longrightarrow
X
$$
에 대하여 More on Morphisms, Lemma
\ref{more-morphisms-lemma-generic-flatness-stratification-scheme}에서와 같이
택하자. 앞으로는 이 층화의 모든 성질을 더 언급하지 않고 사용하겠다.
구성에 의해 각 $f_i$를 $U_k = S_k \setminus S_{k + 1}$로 기저변환한
사상은 평탄이다. $Y_k$를 $S_k$ 안에서 $U_k$의 스킴론적 폐포라 하자.
$U_k \to S_k$는 준콤팩트 열린 몰입이므로(Properties, Lemma
\ref{properties-lemma-quasi-coherent-finite-type-ideals}를 보라),
$U_k \subset Y_k$는 준콤팩트하고 조밀하며 스킴론적으로 조밀한 열린
몰입이다. Morphisms, Lemma
\ref{morphisms-lemma-quasi-compact-scheme-theoretic-image}를 보라. 사상
$\coprod_{k = 0, \ldots, t - 1} Y_k \to X$는 유한 전사이므로
$\{Y_k \to X\}$는 표준 ph 덮개이고, 따라서 표준 V 덮개이다(위를 보라).
표준 V 덮개의 추이성(Topologies, Lemma
\ref{topologies-lemma-composition-standard-V})에 의해 덮개
$\{X_i \to X\}$를 각 $Y_k$로 당긴 것이 표준 V 덮개로 세분될 수 있음을
보이면 충분하다. 이로써 다음 문단에서 설명할 경우로 축약된다.

\medskip\noindent
$\{f_i : X_i \to X\}_{i = 1, \ldots, n}$가 표준 V 덮개이고 $f_i$가
유한표현이며, $X_i \times_X U \to U$가 평탄이게 하는 조밀한 준콤팩트
열린집합 $U \subset X$가 존재한다고 하자. Theorem
\ref{flat-theorem-flatten-module}에 의해 $f_i$의 강변환
$f'_i : X'_i \to X'$이 평탄이게 하는 $U$-허용 블로업 $X' \to X$가
존재한다. $U \subset X$가 조밀하고 $U$가 $X'$의 열린집합과
동일시되므로 사영적(따라서 닫힌) 사상 $X' \to X$는 전사임에 유의하자.
필요하다면 $X'$을 또 다른 $U$-허용 블로업으로 바꾸어
$U \subset X'$가 스킴론적으로 조밀하다고도 가정해도 된다(Remark
\ref{flat-remark-successive-blowups}를 보라). 따라서 모든 점 $x \in X'$에
대하여 값매김환 $V$와 사상 $g : \Spec(V) \to X'$가 존재하여
$\Spec(V)$의 일반점은 $U$ 안으로, $\Spec(V)$의 닫힌점은 $x$로 간다.
Morphisms,
Lemma \ref{morphisms-lemma-reach-points-scheme-theoretic-image}를 보라.
$\{X_i \to X\}$는 표준 V 덮개이므로 값매김환의 확대 $V \subset W$,
지표 $i$, 사상 $\Spec(W) \to X_i$를 택하여 그림
$$
\xymatrix{
\Spec(W) \ar[d] \ar[rr] & & X_i \ar[d] \\
\Spec(V) \ar[r] & X' \ar[r] & X
}
$$
이 가환하게 할 수 있다. $X'_i \subset X' \times_X X_i$는 열린집합
$U \times_X X_i$를 포함하는 닫힌 부분스킴이고, $\Spec(W)$는 정역
스킴이며, 유도된 사상 $h : \Spec(W) \to X' \times_X X_i$는
$\Spec(W)$의 일반점을 $U \times_X X_i$ 안으로 보내므로, $h$는 닫힌
부분스킴 $X'_i \subset X' \times_X X_i$를 통하여 분해된다. 따라서
$\{f'_i : X'_i \to X'\}$는 V 덮개이다. 특히 $\coprod f'_i$는 전사이고,
특히 $\{X'_i \to X'\}$는 fppf 덮개이다. fppf 덮개는 ph 덮개이므로
(More on Morphisms, Lemma \ref{more-morphisms-lemma-fppf-ph}),
% P05-EN-CH38-0269
$\{X'_i \to X\}$를 세분하는 표준 ph 덮개 $\{Y_j \to X'\}$를 찾을 수
있다. 이 덮개가 고유 전사 사상 $Y \to X'$와 유한 아핀 열린 덮개
$Y = \bigcup Y_j$에 의해 주어진다고 하자. 그러면 합성 $Y \to X$는
고유 전사이고, 따라서 $\{Y_j \to X\}$는 표준 ph 덮개이다. 이로써
증명이 끝난다.
\end{proof}

\noindent
이제 우리의 정의를 제시한다.

\begin{definition}
\label{flat-definition-h-covering}
$T$를 스킴이라 하자. {\it $T$의 h 덮개}란 각 $f_i$가 국소적으로
유한표현이고 Lemma
\ref{flat-lemma-equivalence-h-v-locally-finite-presentation}의 동치 조건 중
하나를 만족하는 사상들의 족 $\{f_i : T_i \to T\}_{i \in I}$이다.
\end{definition}

\noindent
뇌터 스킴들에 대하여 이는 ph 덮개와 같고(이를 아래의 Lemma
\ref{flat-lemma-Noetherian-h-ph}에 기록한다), Voevodsky의 개념을 회복한다.

\begin{lemma}
\label{flat-lemma-Noetherian-h-covering}
$X$를 뇌터 스킴이라 하고 $\{X_i \to X\}_{i \in I}$를 유한형 사상들의
유한 족이라 하자. 다음은 동치이다.
\begin{enumerate}
\item $\coprod_{i \in I} X_i \to X$는 보편 잠입적이다(Morphisms,
Definition \ref{morphisms-definition-submersive}).
\item $\{X_i \to X\}_{i \in I}$는 h 덮개이다.
\end{enumerate}
\end{lemma}

\begin{proof}
(2) $\Rightarrow$ (1)은 더 일반적인 Topologies, Lemma
\ref{topologies-lemma-V-covering-universally-submersive}와 h 덮개의 정의에서
따른다. $\coprod X_i \to X$가 보편 잠입적이라고 하자.
$\{X_i \to X\}$가 ph 덮개로 세분될 수 있음을 보이겠다. 그러면
Topologies, Lemma \ref{topologies-lemma-refine-by-ph}와 h 덮개의 정의에
의해 충분하다. 논증은 Lemma
\ref{flat-lemma-equivalence-h-v-locally-finite-presentation}의 증명에서 사용한
것과 같다.

\medskip\noindent
유한표현 사상
$$
X = S \supset S_0 \supset S_1 \supset \ldots \supset S_t = \emptyset
$$
과 같은 일반 평탄성 층화를
$$
\coprod\nolimits_{i = 1, \ldots, n} f_i :
\coprod\nolimits_{i = 1, \ldots, n} X_i
\longrightarrow
X
$$
에 대하여 More on Morphisms, Lemma
\ref{more-morphisms-lemma-generic-flatness-stratification-scheme}에서와 같이
택하자. 앞으로는 이 층화의 모든 성질을 더 언급하지 않고 사용하겠다.
구성에 의해 각 $f_i$를 $U_k = S_k \setminus S_{k + 1}$로 기저변환한
사상은 평탄이다. $Y_k$를 $S_k$ 안에서 $U_k$의 스킴론적 폐포라 하자.
$U_k \to S_k$는 준콤팩트 열린 몰입이므로(이 문단의 모든 스킴은
뇌터이다), $U_k \subset Y_k$는 준콤팩트하고 조밀하며 스킴론적으로
조밀한 열린 몰입이다. Morphisms, Lemma
\ref{morphisms-lemma-quasi-compact-scheme-theoretic-image}를 보라. 사상
$\coprod_{k = 0, \ldots, t - 1} Y_k \to X$는 유한 전사이므로
$\{Y_k \to X\}$는 ph 덮개이다. ph 덮개의 추이성(Topologies, Lemma
\ref{topologies-lemma-ph})에 의해 덮개 $\{X_i \to X\}$를 각 $Y_k$로
당긴 것이 ph 덮개로 세분될 수 있음을 보이면 충분하다. 이로써 다음
문단에서 설명할 경우로 축약된다.

\medskip\noindent
$\coprod X_i \to X$가 보편 잠입적이고, 모든 $i$에 대하여
$X_i \times_X U \to U$가 평탄이게 하는 조밀한 열린집합
$U \subset X$가 존재한다고 하자. Theorem \ref{flat-theorem-flatten-module}에
의해 모든 $i$에 대하여 $f_i$의 강변환 $f'_i : X'_i \to X'$이 평탄이게
하는 $U$-허용 블로업 $X' \to X$가 존재한다. $U \subset X$가 조밀하고
$U$가 $X'$의 열린집합과 동일시되므로 사영적(따라서 닫힌) 사상
$X' \to X$는 전사임에 유의하자. 필요하다면 $X'$을 또 다른 $U$-허용
블로업으로 바꾸어 $U \subset X'$가 조밀하다고도 가정해도 된다(Remark
\ref{flat-remark-successive-blowups}를 보라). 따라서 모든 점 $x \in X'$에
대하여 이산 값매김환 $A$와 사상 $g : \Spec(A) \to X'$가 존재하여
$\Spec(A)$의 일반점은 $U$ 안으로 가고 $\Spec(A)$의 닫힌점은 $x$로
간다. Limits, Lemma \ref{limits-lemma-reach-point-closure-Noetherian}를 보라.
다음과 같이 두자.
$$
W = \Spec(A) \times_X \coprod X_i = \coprod \Spec(A) \times_X X_i
$$
$\coprod X_i \to X$가 보편 잠입적이므로 $W$ 안에 특수화
$w' \leadsto w$가 존재하여 $w'$은 $\Spec(A)$의 일반점으로 가고 $w$는
$\Spec(A)$의 닫힌점으로 간다(그렇지 않으면 $W \to \Spec(A)$의 닫힌
올은 일반화 아래 안정적이어서 열린집합이 되며, 이는
$W \to \Spec(A)$가 잠입적이라는 사실에 모순된다). 어떤 지표에 대하여
$w' \in \Spec(A) \times_X X_i$라고 하자. 그러면 물론
$w \in \Spec(A) \times_X X_i$이기도 하다. $x'_i \leadsto x_i$를
$X' \times_X X_i$ 안에서 $w' \leadsto w$의 상이라 하자.
$x'_i \in X'_i$이고 $X'_i \subset X' \times_X X_i$는 닫힌 부분스킴이므로
$x_i \in X'_i$임을 알 수 있다. $x_i$는 $x \in X'$로 가므로
$\coprod X'_i \to X'$는 전사이다! 특히 $\{X'_i \to X'\}$는 fppf
덮개이다. 그런데 fppf 덮개는 ph 덮개이다(More on Morphisms, Lemma
\ref{more-morphisms-lemma-fppf-ph}). $X' \to X$는 고유 전사이므로
$\{X'_i \to X\}$가 ph 덮개라는 결론을 얻고 증명이 끝난다.
\end{proof}

\begin{lemma}
\label{flat-lemma-Noetherian-h-ph}
$X$를 국소 뇌터 스킴이라 하자. $X$를 목표로 하는 사상들의 족
$\{f_i : X_i \to X\}_{i \in I}$가 h 덮개일 필요충분조건은 ph 덮개인
것이다.
\end{lemma}

\begin{proof}
Definition \ref{flat-definition-h-covering}에 의해 h 덮개는 ph 덮개이다.
역으로 $\{f_i : X_i \to X\}$가 ph 덮개이면 사상들 $f_i$는
국소적으로 유한형이다(Topologies, Definition
\ref{topologies-definition-ph-covering}). $X$가 국소 뇌터이므로 각
$f_i$는 국소적으로 유한표현이고, 따라서 정의에 의해 h 덮개를 얻는다.
\end{proof}

\noindent
다음 보조정리와 \cite[Theorem 8.4]{rydh_descent}는 우리의 정의가 David
Rydh의 논문 \cite{rydh_descent}에 나오는 정의와 일치함(또는 적어도
밀접하게 관련됨)을 보여 준다. 간단히 하기 위하여 아핀 기저로 제한한다.
% P05-EN-CH38-0270

\begin{lemma}
\label{flat-lemma-approximate-h-cover}
$X$를 아핀 스킴이라 하고 $\{X_i \to X\}_{i \in I}$를 h 덮개라 하자.
그러면 유한표현인(!) 고유 전사 사상
$$
Y \longrightarrow X
$$
과 유한 아핀 열린 덮개 $Y = \bigcup_{j = 1, \ldots, m} Y_j$가 존재하여
$\{Y_j \to X\}_{j = 1, \ldots, m}$가 $\{X_i \to X\}_{i \in I}$를 세분한다.
\end{lemma}

\begin{proof}
가정에 의해 고유 전사 사상 $Y \to X$와 유한 아핀 열린 덮개
$Y = \bigcup_{j = 1, \ldots, m} Y_j$가 존재하여
$\{Y_j \to X\}_{j = 1, \ldots, m}$가 $\{X_i \to X\}_{i \in I}$를
세분한다. 이는 각 $j$에 대하여 지표 $i_j \in I$와 $X$ 위의 사상
$h_j : Y_j \to X_{i_j}$가 존재한다는 뜻이다. Definition
\ref{flat-definition-h-covering} 및 Topologies, Definition
\ref{topologies-definition-ph-covering}을 보라. 문제는 $Y \to X$가
유한표현인지 알지 못한다는 것이다. Limits, Lemma
\ref{limits-lemma-proper-limit-of-proper-finite-presentation}에 의해
$$
Y = \lim Y_\lambda
$$
를 $X$ 위에서 고유이고 유한표현인 스킴들 $Y_\lambda$의 유향 극한으로
쓸 수 있으며, 사상들 $Y \to Y_\lambda$와 전이 사상들은 닫힌 몰입이다.
% P05-EN-CH38-0271
각 $Y_\lambda \to X$가 전사임에 유의하자. Limits, Lemma
\ref{limits-lemma-descend-opens}에 의해 어떤 $\lambda$와
$j = 1, \ldots, m$에 대한 준콤팩트 열린집합
$Y_{\lambda, j} \subset Y_\lambda$를 찾을 수 있는데, 이들은
$Y_\lambda$를 덮고 $Y$로 제한하면 $Y_j$가 된다. 그러면
$Y_j = \lim Y_{\lambda, j}$이다. $\lambda$를 늘린 뒤 모든 $j$에 대하여
$Y_{\lambda, j}$가 아핀이라고 가정해도 된다. Limits, Lemma
\ref{limits-lemma-limit-affine}를 보라. 마지막으로 $X_i \to X$가
국소적으로 유한표현이므로 국소적으로 유한표현인 사상들의 함자적
특징화(Limits, Proposition
\ref{limits-proposition-characterize-locally-finite-presentation})를
사용하여 다음 성질의 $\lambda$를 찾을 수 있다. 각 $j$에 대하여 위에서
택한 사상 $h_j$가 되도록 $Y_j$로 제한되는 사상
$h_{\lambda, j} : Y_{\lambda, j} \to X_{i_j}$가 존재한다. 따라서
$\{Y_{\lambda, j} \to X\}$는 $\{X_i \to X\}$를 세분하며 증명이 끝난다.
\end{proof}

\noindent
h 덮개의 일반 이론 전개로 돌아가자.

\begin{lemma}
\label{flat-lemma-zariski-h}
fppf 덮개는 h 덮개이다. 따라서 syntomic, 매끄러운, 에탈, Zariski
덮개들도 h 덮개이다.
\end{lemma}

\begin{proof}
fppf 덮개의 사상들은 국소적으로 유한표현이어야 하고 fppf 덮개들은 ph
덮개들이므로 참이다. More on Morphisms, Lemma
\ref{more-morphisms-lemma-fppf-ph}를 보라.
% P05-EN-CH38-0272
둘째 명제는 첫째 명제와 Topologies, Lemma
\ref{topologies-lemma-zariski-etale-smooth-syntomic-fppf}에서 따른다.
\end{proof}

\begin{lemma}
\label{flat-lemma-surjective-proper-finite-presentation-h}
$f : Y \to X$를 유한표현인 스킴들의 고유 전사 사상이라 하자. 그러면
$\{Y \to X\}$는 h 덮개이다.
\end{lemma}

\begin{proof}
Topologies, Lemmas
\ref{topologies-lemma-zariski-etale-smooth-syntomic-fppf-fpqc-ph-V} 및
\ref{topologies-lemma-surjective-proper-ph}를 결합하라.
\end{proof}

\begin{lemma}
\label{flat-lemma-refine-by-h}
$T$를 스킴이라 하자. $\{f_i : T_i \to T\}_{i \in I}$를 사상들의 족이라
하되 모든 $i$에 대하여 $f_i$가 국소적으로 유한표현이라고 하자. 다음은
동치이다.
\begin{enumerate}
\item $\{T_i \to T\}_{i \in I}$는 h 덮개이다.
\item $\{T_i \to T\}_{i \in I}$를 세분하는 h 덮개가 존재한다.
\item $\{\coprod_{i \in I} T_i \to T\}$는 h 덮개이다.
\end{enumerate}
\end{lemma}

\begin{proof}
이는 ph 덮개들에 대한 유사한 명제(Topologies, Lemma
\ref{topologies-lemma-refine-by-ph}) 또는 V 덮개들에 대한 유사한
명제(Topologies, Lemma \ref{topologies-lemma-refine-by-V})에서 따른다.
\end{proof}

\noindent
다음으로 h 덮개의 우리 개념이 Sites, Definition
\ref{sites-definition-site}의 조건들을 만족함을 보인다.

\begin{lemma}
\label{flat-lemma-h}
$T$를 스킴이라 하자.
\begin{enumerate}
\item $T' \to T$가 동형이면 $\{T' \to T\}$는 $T$의 h 덮개이다.
\item $\{T_i \to T\}_{i\in I}$가 h 덮개이고 각 $i$에 대하여 h 덮개
$\{T_{ij} \to T_i\}_{j\in J_i}$가 주어지면
$\{T_{ij} \to T\}_{i \in I, j\in J_i}$는 h 덮개이다.
\item $\{T_i \to T\}_{i\in I}$가 h 덮개이고 $T' \to T$가 스킴들의
사상이면 $\{T' \times_T T_i \to T'\}_{i\in I}$는 h 덮개이다.
\end{enumerate}
\end{lemma}

\begin{proof}
ph 덮개 또는 V 덮개에 대한 대응하는 명제(Topologies, Lemma
\ref{topologies-lemma-ph} 또는 \ref{topologies-lemma-V})와 국소적으로
유한표현인 사상들의 부류가 기저변환과 합성 아래 보존된다는 사실에서
곧바로 따른다.
\end{proof}

\noindent
다음으로 Stacks project에서 사용할 큰 h 사이트들을 정의한다. 계속하기
전에 Topologies, Section \ref{topologies-section-procedure}의 일반적인
논의를 읽는 것이 도움이 된다.

\begin{definition}
\label{flat-definition-big-h-site}
{\it 큰 h 사이트}란 다음과 같이 구성된 Sites, Definition
\ref{sites-definition-site}에서와 같은 임의의 사이트 $\Sch_h$이다.
\begin{enumerate}
\item 스킴들의 임의의 집합 $S_0$와 이 스킴들 사이의 h 덮개들의 임의의
집합 $\text{Cov}_0$를 택한다.
\item 바탕 범주로서 집합 $S_0$에서 시작하여 Sets, Lemma
\ref{sets-lemma-construct-category}에서와 같이 구성된 임의의 범주
$\Sch_\alpha$를 취한다.
\item 범주 $\Sch_\alpha$, h 덮개들의 부류, 위에서 택한 집합
$\text{Cov}_0$에서 시작하여 Sets, Lemma
\ref{sets-lemma-coverings-site}에서와 같은 덮개들의 임의의 집합을 택한다.
\end{enumerate}
\end{definition}

\noindent
큰 사이트 정의의 동기와 설명은 Topologies, Definition
\ref{topologies-definition-big-zariski-site} 뒤의 비고들을 보라.

\begin{definition}
\label{flat-definition-standard-h}
$T$를 아핀 스킴이라 하자. $T$의 {\it 표준 h 덮개}란 각 $T_i$가
아핀이고 $f_i$가 유한표현이며 다음 두 동치 조건 중 하나를 만족하는
족 $\{f_i : T_i \to T\}_{i = 1, \ldots, n}$이다. (1)
$\{T_i \to T\}$는 표준 ph 덮개로 세분될 수 있다. (2)
$\{T_i \to T\}$는 V 덮개이다.
\end{definition}

\noindent
조건들의 동치는 Lemma
\ref{flat-lemma-equivalence-h-v-locally-finite-presentation}, Topologies,
Definition \ref{topologies-definition-ph-covering}, Lemma
\ref{topologies-lemma-refine-by-ph}에서 따른다.

\medskip\noindent
스킴 $S$의 큰 h 사이트를 도입하기 전에, 큰 h 사이트 $\Sch_h$ 위의
위상은 어떤 의미에서 모든 스킴들의 범주 위의 h 위상에서 유도된다는
점을 지적하자.

\begin{lemma}
\label{flat-lemma-h-induced}
$\Sch_h$를 Definition \ref{flat-definition-big-h-site}에서와 같은 큰 h
사이트라 하자. $T \in \Ob(\Sch_h)$라 하고
$\{T_i \to T\}_{i \in I}$를 $T$의 임의의 h 덮개라 하자.
\begin{enumerate}
\item $\{T_i \to T\}_{i \in I}$를 세분하는, 사이트 $\Sch_h$ 안에서의
$T$의 덮개 $\{U_j \to T\}_{j \in J}$가 존재한다.
\item $\{T_i \to T\}_{i \in I}$가 표준 h 덮개이면 이는 자명하게
$\Sch_h$의 덮개와 동치이다.
\item $\{T_i \to T\}_{i \in I}$가 Zariski 덮개이면 이는 자명하게
$\Sch_h$의 덮개와 동치이다.
\end{enumerate}
\end{lemma}

\begin{proof}
생략한다. 도움말: 이는 Topologies, Lemma
\ref{topologies-lemma-ph-induced}의 증명과 정확히 같다.
\end{proof}

\begin{definition}
\label{flat-definition-big-small-h}
$S$를 스킴이라 하고 $\Sch_h$를 $S$를 포함하는 큰 h 사이트라 하자.
\begin{enumerate}
\item $(\Sch/S)_h$로 나타내는 {\it $S$의 큰 h 사이트}는 Sites, Section
\ref{sites-section-localize}에서 도입한 사이트 $\Sch_h/S$이다.
\item $(\textit{Aff}/S)_h$로 나타내는 {\it $S$의 큰 아핀 h 사이트}는
대상들이 아핀 $U/S$인 $(\Sch/S)_h$의 충만 부분범주이다.
$(\textit{Aff}/S)_h$의 덮개란 표준 h 덮개인 $(\Sch/S)_h$의 임의의 덮개
$\{U_i \to U\}$이다.
\end{enumerate}
\end{definition}

\noindent
큰 아핀 h 사이트가 사이트임을 명시적으로 서술한다.

\begin{lemma}
\label{flat-lemma-verify-site-h}
$S$를 스킴이라 하고 $\Sch_h$를 $S$를 포함하는 큰 h 사이트라 하자.
그러면 $(\textit{Aff}/S)_h$는 사이트이다.
\end{lemma}

\begin{proof}
Topologies, Lemma \ref{topologies-lemma-verify-site-etale}의 증명과 같이
논증하면, 표준 h 덮개들의 모임이 Sites, Definition
\ref{sites-definition-site}의 성질 (1), (2), (3)을 만족함을 보이면
충분하다. 예를 들어 표준 h 덮개 $\{T_i \to T\}_{i\in I}$와 각 $i$에
대한 표준 h 덮개 $\{T_{ij} \to T_i\}_{j \in J_i}$가 주어지면,
$\{T_{ij} \to T\}_{i \in I, j\in J_i}$는 h 덮개이고(Lemma
\ref{flat-lemma-h}), $\bigcup_{i\in I} J_i$는 유한하며 각 $T_{ij}$는
아핀이다. 따라서 $\{T_{ij} \to T\}_{i \in I, j\in J_i}$는 표준 h
덮개이다. 그러므로 이는 명백하다.
\end{proof}

\begin{lemma}
\label{flat-lemma-fibre-products-h}
$S$를 스킴이라 하고 $\Sch_h$를 $S$를 포함하는 큰 h 사이트라 하자.
사이트들 $\Sch_h$, $(\Sch/S)_h$, $(\textit{Aff}/S)_h$의 바탕 범주들은
올곱들을 갖는다. 각 경우 모든 스킴들의 범주 $\Sch$로 가는 자명한
함자는 올곱을 취하는 것과 가환한다. 범주 $(\Sch/S)_h$는 끝 대상
$S/S$를 갖는다.
\end{lemma}

\begin{proof}
$\Sch_h$에 대해서는 구성에 의해 참이다. Sets, Lemma
\ref{sets-lemma-what-is-in-it}를 보라. $U, V, W \in \Ob(\Sch_h)$인 스킴
사상들 $U \to S$, $V \to U$, $W \to U$가 주어졌다고 하자. $\Sch_h$
안의 올곱 $V \times_U W$는 $\Sch$ 안의 올곱이고, $S$ 위의 모든 스킴들의
범주에서 $U/S$ 위의 $V/S$와 $W/S$의 올곱이며, 따라서
$(\Sch/S)_h$ 안에서도 올곱이다. 이로써 $(\Sch/S)_h$에 대한 결과가
증명된다. $U, V, W$가 아핀이면 $V \times_U W$도 아핀이므로
$(\textit{Aff}/S)_h$에 대한 결과도 성립한다.
\end{proof}

\noindent
다음으로 큰 아핀 사이트가 큰 사이트와 같은 토포스를 정의함을 확인한다.

\begin{lemma}
\label{flat-lemma-affine-big-site-h}
$S$를 스킴이라 하고 $\Sch_h$를 $S$를 포함하는 큰 h 사이트라 하자.
함자 $(\textit{Aff}/S)_h \to (\Sch/S)_h$는 쌍대연속이고 토포스들의 동치
$\Sh((\textit{Aff}/S)_h)$에서 $\Sh((\Sch/S)_h)$로 가는 동치를 유도한다.
\end{lemma}

\begin{proof}
특별한 쌍대연속 함자의 개념은 Sites, Definition
\ref{sites-definition-special-cocontinuous-functor}에서 도입한다. 따라서
Sites, Lemma \ref{sites-lemma-equivalence}의 가정 (1)--(5)를 확인해야
한다. 포함 함자를 $u : (\textit{Aff}/S)_h \to (\Sch/S)_h$라 하자.
$T$가 아핀일 때 $T/S$의 모든 h 덮개는, 예를 들어 Lemma
\ref{flat-lemma-approximate-h-cover}에 의해 표준 h 덮개로 세분될 수 있으므로
쌍대연속이다. 따라서 (1)이 성립한다. 표준 h 덮개는 h 덮개이므로
$u$가 연속임을 곧바로 알 수 있다. 따라서 (2)가 성립한다. (3)과 (4)는
$u$가 충만충실이라는 사실에서 곧바로 따른다. 마지막으로 모든 스킴은
아핀 열린 덮개(이는 h 덮개이다)를 가지므로 조건 (5)가 따른다.
\end{proof}

\begin{lemma}
\label{flat-lemma-characterize-sheaf-h}
$\mathcal{F}$를 $(\Sch/S)_h$ 위의 전층이라 하자. 그러면
$\mathcal{F}$가 층일 필요충분조건은 다음과 같다.
\begin{enumerate}
\item $\mathcal{F}$는 Zariski 덮개들에 대한 층 조건을 만족한다.
\item $f : V \to U$가 고유이고 전사이며 유한표현이면
$\mathcal{F}(U)$는 두 사상
$\mathcal{F}(V) \to \mathcal{F}(V \times_U V)$의 등화자로 전단사된다.
\end{enumerate}
더욱이 (1)이 성립할 때 성질 (2)는 다음 성질과 동치이다.
\begin{enumerate}
\item[(2')] $U$가 아핀인 (2)에서와 같은 $\{V \to U\}$에 대한 층 성질.
\end{enumerate}
\end{lemma}

\begin{proof}
(1)과 (2)가 성립하면 $\mathcal{F}$가 층임을 보이겠다.
% P05-EN-CH38-0273
$\{T_i \to T\}$를 $(\Sch/S)_h$ 안의 덮개라 하자. 이 덮개에 대한 층
조건을 확인하겠다. $s_i \in \mathcal{F}(T_i)$를
$T_i \times_T T_{i'}$ 위에서 같은 단면으로 제한되는 단면들이라 하자.
$T_i$ 위에서 $s_i$로 제한되는 유일한 단면 $s \in \mathcal{F}(T)$가
존재함을 보이겠다. $T = \bigcup U_j$를 아핀 열린 덮개라 하자. 성질
(1)에 의해 $U_j \cap U_{j'}$에서 일치하는 단면들
$s_j \in \mathcal{F}(U_j)$를 만들면 $s$를 만들기에 충분하다. 덮개들
$\{T_i \times_T U_j \to U_j\}$를 생각하자. 그러면
$s_{ji} = s_i|_{T_i \times_T U_j}$는
$(T_i \times_T U_j) \times_{U_j} (T_{i'} \times_T U_j)$ 위에서 일치하는
단면들이다. 유한표현 고유 전사 사상 $V_j \to U_j$와 유한 아핀 열린
덮개 $V_j = \bigcup V_{jk}$를 택하여 $\{V_{jk} \to U_j\}$가
$\{T_i \times_T U_j \to U_j\}$를 세분하게 하자. Lemma
\ref{flat-lemma-approximate-h-cover}를 보라. $s_{jk} \in \mathcal{F}(V_{jk})$가
암시된 사상들에 의한 $s_{ji}$의 $V_{jk}$로의 당김을 나타내면,
$s_{jk}$들이 붙어서 단면 $s'_j \in \mathcal{F}(V_j)$를 이룸을 알 수
있다. 겹침에서의 일치를 다시 사용하면 $s'_j$가 두 사상
$\mathcal{F}(V_j) \to \mathcal{F}(V_j \times_{U_j} V_j)$의 등화자에
속함을 알 수 있다. 따라서 (2)에 의해 $s'_j$는 유일한 단면
$s_j \in \mathcal{F}(U_j)$에서 온다. 이 단면들 $s_j$가 원하는 모든
성질을 갖는다는 확인은 생략한다.

\medskip\noindent
(1)이 성립할 때 (2)와 (2')가 동치임의 증명. $V \to U$가 고유이고
전사이며 유한표현인 $(\Sch/S)_h$의 사상이라고 하자. 아핀 열린 덮개
$U = \bigcup U_i$를 택하고 $V_i = V \times_U U_i$라 두자. (2')에 의해
$\mathcal{F}(U_i) \to \mathcal{F}(V_i)$가 단사이고 (1)에 의해
$\mathcal{F}(U) \to \prod \mathcal{F}(U_i)$가 단사임을 아니까
$\mathcal{F}(U) \to \mathcal{F}(V)$가 단사임을 알 수 있다. 마지막으로
두 사상 $\mathcal{F}(V) \to \mathcal{F}(V \times_U V)$의 등화자에 속하는
$t \in \mathcal{F}(V)$가 주어졌다고 하자. 그러면 모든 $i$에 대하여
$t|_{V_i}$는 두 사상
$\mathcal{F}(V_i) \to \mathcal{F}(V_i \times_{U_i} V_i)$의 등화자에
속한다. 따라서 (2')에 의해 모든 $i$에 대하여 $t|_{V_i}$로 가는 유일한
단면 $s_i \in \mathcal{F}(U_i)$를 얻는다. 모든 $i, j$에 대하여
$s_i|_{U_i \cap U_j} = s_j|_{U_i \cap U_j}$임의 확인은 생략한다. 이는
방금 보인 유일성을 사용한다. 덮개 $U = \bigcup U_i$에 대한 층 성질에
의해 단면 $s \in \mathcal{F}(U)$를 얻는다. $s$가
$\mathcal{F}(V)$에서 $t$로 간다는 증명은 생략한다.
\end{proof}

\noindent
다음으로 이 사이트들에 결부된 토포스들 사이의 몇 가지 관계를 확립한다.

\begin{lemma}
\label{flat-lemma-morphism-big-h}
$\Sch_h$를 큰 h 사이트라 하고 $f : T \to S$를 $\Sch_h$ 안의 사상이라
하자. 함자
$$
u : (\Sch/T)_h \longrightarrow (\Sch/S)_h,
\quad
V/T \longmapsto V/S
$$
는 쌍대연속이고 연속인 오른쪽 수반함자
$$
v : (\Sch/S)_h \longrightarrow (\Sch/T)_h,
\quad
(U \to S) \longmapsto (U \times_S T \to T).
$$
를 갖는다. 이들은 같은 토포스 사상
$$
f_{big} :
\Sh((\Sch/T)_h)
\longrightarrow
\Sh((\Sch/S)_h)
$$
을 유도한다. 또한
$f_{big}^{-1}(\mathcal{G})(U/T) = \mathcal{G}(U/S)$이고,
$f_{big, *}(\mathcal{F})(U/S) = \mathcal{F}(U \times_S T/T)$이다.
$f_{big}^{-1}$은 올곱 및 등화자와 가환하는 왼쪽 수반함자 $f_{big!}$을
갖는다.
\end{lemma}

\begin{proof}
함자 $u$는 쌍대연속이고 연속이며 올곱 및 등화자와 가환한다. 따라서
Sites, Lemmas \ref{sites-lemma-when-shriek} 및
\ref{sites-lemma-preserve-equalizers}를 적용하여 $f_{big}^{-1}$의 공식과
$f_{big!}$의 존재를 얻는다. 더욱이 $U/T$와 $V/S$가 주어졌을 때 원하는
대로 $\Mor_S(u(U), V) = \Mor_T(U, V \times_S T)$이므로 함자 $v$는 오른쪽
수반함자이다. 따라서 Sites, Lemmas
\ref{sites-lemma-have-functor-other-way} 및
\ref{sites-lemma-have-functor-other-way-morphism}을 적용하여
$f_{big, *}$의 공식을 얻을 수 있다.
\end{proof}

\begin{lemma}
\label{flat-lemma-composition-h}
$X$, $Y$, $Y$가 $(\Sch/S)_h$ 안의 스킴들이고 사상
$f : X \to Y$, $g : Y \to Z$가 주어지면
% P05-EN-CH38-0274
$g_{big} \circ f_{big} = (g \circ f)_{big}$.
\end{lemma}

\begin{proof}
이는 Lemma \ref{flat-lemma-morphism-big-h}에서 얻은 큰 사이트들 위 함자들의
직상과 역상의 간단한 기술에서 따른다.
\end{proof}










\section{h 위상에 관한 추가 결과}
\label{flat-section-h-more}

\noindent
이 절에서는 h 위상에 관한 결과 몇 가지를 더 증명한다. 먼저 몇 가지
반례를 살펴보자.

\begin{example}
\label{flat-example-structure-sheaf-not-h-sheaf}
``구조층'' $\mathcal{O}$는 h 위상에서 층이 아니다. 예를 들어 유한표현
전사 닫힌 몰입 $X \to Y$를 생각하자. 그러면, 예컨대 Lemma
\ref{flat-lemma-surjective-proper-finite-presentation-h}에 의해
$\{X \to Y\}$는 h 덮개이다. 또한 $X \times_Y X = X$임에 유의하자.
% P05-EN-CH38-0275
따라서 $\mathcal{O}$가 h 위상에서 층이라면
$\mathcal{O}_Y(Y) \to \mathcal{O}_X(X)$는 전단사일 것이다. $X$, $Y$가
아핀이고 사상 $X \to Y$가 동형이 아니기만 하면 이는 성립하지 않는다.
\end{example}

\begin{example}
\label{flat-example-representable-sheaf-not-h-sheaf}
임의의 사이트 $(\Sch/S)_h$에서 위상은 준표준이 아니다. 다시 말해
표현가능 층들은 층이 아니다.
% P05-EN-CH38-0276
실제로 $(\Sch/S)_h$에서 $\mathcal{O}(X) = \Mor_S(X, \mathbf{A}^1_S)$이므로
``구조층'' $\mathcal{O}$는 표현가능하고, Example
\ref{flat-example-structure-sheaf-not-h-sheaf}에서 $\mathcal{O}$가 층이 아님을
보았다.
\end{example}

\begin{lemma}
\label{flat-lemma-limit-h-topology}
$T$를 아핀 스킴들의 유향 역계의 극한
$T = \lim_{i \in I} T_i$로 쓰이는 아핀 스킴이라 하자.
\begin{enumerate}
\item $\mathcal{V} = \{V_j \to T\}_{j = 1, \ldots, m}$를 $T$의 표준 h
덮개라 하자. Definition \ref{flat-definition-standard-h}를 보라. 그러면 지표
$i$와 표준 h 덮개
$\mathcal{V}_i = \{V_{i, j} \to T_i\}_{j = 1, \ldots, m}$가 존재하여
이를 $T$로 기저변환한 $T \times_{T_i} \mathcal{V}_i$는
$\mathcal{V}$와 동형이다.
\item $\mathcal{V}_i$, $\mathcal{V}'_i$를 $T_i$의 두 표준 h 덮개라 하자.
$f : T \times_{T_i} \mathcal{V}_i \to T \times_{T_i} \mathcal{V}'_i$가
$T$의 덮개들의 사상이면 지표 $i' \geq i$와 사상
% P05-EN-CH38-0277
$f_{i'} : T_{i'} \times_{T_i} \mathcal{V} \to
T_{i'} \times_{T_i} \mathcal{V}'_i$
가 존재하여 이를 $T$로 기저변환하면 $f$가 된다.
\item
% P05-EN-CH38-0278
$f, g : \mathcal{V} \to \mathcal{V}'_i$
가 $T_i$의 표준 h 덮개들의 사상들이고 이들을 $T$로 기저변환한
$f_T, g_T$가 같으면, $f_{T_{i'}} = g_{T_{i'}}$이게 하는 지표
$i' \geq i$가 존재한다.
\end{enumerate}
다시 말해 $T$의 표준 h 덮개들의 범주는 $T_i$의 표준 h 덮개들의
범주들을 $I$ 위에서 취한 여극한이다.
\end{lemma}

\begin{proof}
Limits, Lemma \ref{limits-lemma-descend-finite-presentation}에 의해 $T$
위의 유한표현 스킴들의 범주는 $T_i$ 위의 유한표현 스킴들의 범주들을
$I$ 위에서 취한 여극한이다.
% P05-EN-CH38-0279
Limits, Lemma \ref{limits-lemma-descend-affine-finite-presentation}에 의해
$T$ 위에서 아핀이고 유한표현인 스킴들의 범주에 대해서도 마찬가지이다.
보조정리의 증명을 끝내려면 다음을 보이면 충분하다.
$\{V_{j, i} \to T_i\}_{j = 1, \ldots, m}$가 $V_{j, i}$들이 아핀인
유한표현 사상들의 유한 족이고, 이를 기저변환한 족
$\{T \times_{T_i} V_{j, i} \to T\}$가 h 덮개이면, 어떤 $i' \geq i$에
대하여 족 $\{T_{i'} \times_{T_i} V_{j, i} \to T_{i'}\}$가 h 덮개이다.
이를 보기 위하여 Lemma \ref{flat-lemma-approximate-h-cover}를 사용해 유한표현
고유 전사 사상 $Y \to T$와 유한 아핀 열린 덮개
$Y = \bigcup_{k = 1, \ldots, n} Y_k$를 택하여
$\{Y_k \to T\}_{k = 1, \ldots, n}$가
$\{T \times_{T_i} V_{j, i} \to T\}$를 세분하게 하자. 위의 논증과 Limits,
Lemmas \ref{limits-lemma-eventually-proper},
\ref{limits-lemma-descend-surjective}, \ref{limits-lemma-descend-opens}를
사용하면 $i' \geq i$, 유한표현 고유 전사 사상
$Y_{i'} \to T_{i'}$, 아핀 열린 덮개
$Y_{i'} = \bigcup_{k = 1, \ldots, n} Y_{i', k}$를 찾을 수 있으며,
더욱이
% P05-EN-CH38-0280
$\{Y_{i', k} \to Y_{i'}\}$는
$\{T_{i'} \times_{T_i} V_{j, i} \to T_{i'}\}$를 세분한다. 따라서 마지막에
언급한 족은 h 덮개이고 증명이 끝난다.
\end{proof}

\begin{lemma}
\label{flat-lemma-extend-sheaf-h}
$S$를 큰 사이트 $\Sch_h$에 포함된 스킴이라 하자.
$F : (\Sch/S)_h^{opp} \to \textit{Sets}$를
$\mathcal{C} = (\Sch/S)_h$일 때 Topologies, Lemma
\ref{topologies-lemma-extend}의 성질 (b)를 만족하는 h 층이라 하자. 그러면
$F$를 $S$ 위의 모든 스킴들의 범주로 연장한 $F'$은 모든 h 덮개에 대한
층 조건을 만족하고 극한을 보존한다(Limits, Remark
\ref{limits-remark-limit-preserving}).
\end{lemma}

\begin{proof}
Lemmas \ref{flat-lemma-limit-h-topology} 및 \ref{flat-lemma-h-induced}를 사용하여
Topologies, Lemmas \ref{topologies-lemma-extend-sheaf-general} 및
\ref{topologies-lemma-extend-sheaf}의 증명에 주어진 논증으로 증명된다.
세부 사항은 생략한다.
\end{proof}








\section{블로업 사각형과 ph 위상}
\label{flat-section-blow-up-ph}

\noindent
$X$를 스킴이라 하자. $Z \subset X$를 포함 사상이 유한표현인 닫힌
부분스킴, 즉 $Z$에 대응하는 준연접 아이디얼 층이 유한형인 닫힌
부분스킴이라 하자. $b : X' \to X$를 $Z$에서의 $X$의 블로업이라 하고
$E = b^{-1}(Z)$를 예외 인자라 하자. Divisors, Section
\ref{divisors-section-blowing-up}을 보라. 이 상황과 이 절에서는
\begin{equation}
\label{flat-equation-blow-up-square}
\vcenter{
\xymatrix{
E \ar[d] \ar[r] & X' \ar[d]^b \\
Z \ar[r] & X
}
}
\end{equation}
을 {\it 블로업 사각형}이라고 하겠다.

\begin{lemma}
\label{flat-lemma-blow-up-square-ph}
$\mathcal{F}$를 사이트 $(\Sch/S)_{ph}$ 위의 층이라 하자. Topologies,
Definition \ref{topologies-definition-big-small-ph}를 보라. 그러면 범주
$(\Sch/S)_{ph}$ 안의 모든 블로업 사각형
(\ref{flat-equation-blow-up-square})에 대하여 그림
$$
\xymatrix{
\mathcal{F}(E) & \mathcal{F}(X') \ar[l] \\
\mathcal{F}(Z) \ar[u] & \mathcal{F}(X) \ar[u] \ar[l]
}
$$
은 집합들의 범주에서 카르테시안이다.
\end{lemma}

\begin{proof}
$Z \amalg X' \to X$는 고유 전사 사상이므로
$\{Z \amalg X' \to X\}$는 ph 덮개이다(Topologies, Lemma
\ref{topologies-lemma-surjective-proper-ph}). 또한
$$
(Z \amalg X') \times_X (Z \amalg X') =
Z \amalg E \amalg E \amalg X' \times_X X'
$$
이다. $\mathcal{F}$는 Zariski 층이므로 $\mathcal{F}$는 서로소 합을
곱으로 보냄을 알 수 있다. 따라서 덮개 $\{Z \amalg X' \to X\}$에 대한 층 조건은
$\mathcal{F}(X) \to \mathcal{F}(Z) \times \mathcal{F}(X')$가 단사이고,
그 상이 (a) $t|_E = s'|_E$이고 (b) $s'$가 두 사상
$\mathcal{F}(X') \to \mathcal{F}(X' \times_X X')$의 등화자에 속하는 쌍
$(t, s')$들의 집합이라는 뜻이다. 다음으로 자명한 사상
$$
E \times_Z E \amalg X' \longrightarrow X' \times_X X'
$$
은 고유 전사 사상임에 유의하자. 실제로 $b$는 동형
$X' \setminus E \to X \setminus Z$를 유도한다. 따라서
$\mathcal{F}(X' \times_X X') \to
\mathcal{F}(E \times_Z E) \times \mathcal{F}(X')$가 단사라는 결론을
얻는다. 이에 따라 (a) $\Rightarrow$ (b)이고, 이는 보조정리가 참이라는
뜻이다.
\end{proof}

\begin{lemma}
\label{flat-lemma-thickening-ph}
$\mathcal{F}$를 Topologies, Definition
\ref{topologies-definition-big-small-ph}에서와 같은 사이트
$(\Sch/S)_{ph}$ 위의 층이라 하자. $X \to X'$를 비후화인
$(\Sch/S)_{ph}$의 사상이라 하자. 그러면
$\mathcal{F}(X') \to \mathcal{F}(X)$는 전단사이다.
\end{lemma}

\begin{proof}
$X \to X'$는 고유 전사 사상이고 $X \times_{X'} X = X$임에 유의하자.
% P05-EN-CH38-0281
ph 덮개 $\{X \to X'\}$에 대한 층 성질(Topologies, Lemma
\ref{topologies-lemma-surjective-proper-ph})에서 결론이 따른다.
\end{proof}








\section{거의 블로업 사각형과 h 위상}
\label{flat-section-blow-up-h}

\noindent
블로업 사각형 (\ref{flat-equation-blow-up-square})을 생각하자. 사상
$b : X' \to X$는 사영적이지만(Divisors, Lemma
\ref{divisors-lemma-blowing-up-projective}), 일반적으로 $b$가 유한표현임을
보장하는 간단한 방법은 없다. h 덮개는 유한표현 사상들을 사용하여
구성되므로 변형된 개념이 필요하다. 즉, 스킴들의 가환 그림
\begin{equation}
\label{flat-equation-almost-blow-up-square}
\vcenter{
\xymatrix{
E \ar[d] \ar[r] & X' \ar[d]^b \\
Z \ar[r] & X
}
}
\end{equation}
이 다음 조건들을 만족하면 {\it 거의 블로업 사각형}이라고 하겠다.
\begin{enumerate}
\item $Z \to X$는 유한표현 닫힌 몰입이다.
\item $E = b^{-1}(Z)$는 $X'$의 국소 주 아이디얼 닫힌 부분스킴이다.
\item $b$는 고유이고 유한표현이다.
\item $E$ 위에 지지된 $\mathcal{O}_{X'}$의 단면들의 준연접 아이디얼이
절단하는 닫힌 부분스킴 $X'' \subset X'$은 $Z$에서의 $X$의
블로업이다(Properties, Lemma
\ref{properties-lemma-sections-supported-on-closed-subset}).
\end{enumerate}
그러면 사상 $b$는 동형 $X' \setminus E \to X \setminus Z$를 유도한다.
거의 블로업 사각형의 매우 간단한 예들은 Examples
\ref{flat-example-one-generator} 및 \ref{flat-example-two-generators}를 보라.

\medskip\noindent
블로업의 기저변환은 보통 블로업이 아니지만, 거의 블로업은 기저변환과
호환된다.

\begin{lemma}
\label{flat-lemma-base-change-almost-blow-up}
거의 블로업 사각형 (\ref{flat-equation-almost-blow-up-square})을 생각하자.
$Y \to X$를 임의의 사상이라 하자. 그러면 기저변환
$$
\xymatrix{
Y \times_X E \ar[d] \ar[r] & Y \times_X X' \ar[d] \\
Y \times_X Z \ar[r] & Y
}
$$
도 거의 블로업 사각형이다.
\end{lemma}

\begin{proof}
Morphisms, Lemmas \ref{morphisms-lemma-base-change-proper} 및
\ref{morphisms-lemma-base-change-finite-presentation}에 의해 사상
$Y \times_X X' \to Y$는 고유이고 유한표현이다. 사상
$Y \times_X Z \to Y$는 유한표현 닫힌 몰입이다(Morphisms, Lemma
\ref{morphisms-lemma-base-change-closed-immersion}). $Y \times_X X'$ 안에서
$Y \times_X Z$의 역상은 $Y \times_X X'$ 안에서 $E$의 역상과 같으므로
국소 주 아이디얼이다(Divisors, Lemma
\ref{divisors-lemma-pullback-locally-principal}). 각각 $E$ 및
$Y \times_X E$ 위에 지지된 $\mathcal{O}_{X'}$ 및
% P05-EN-CH38-0282
$\mathcal{O}_{Y \times_Y X'}$의 단면들의 준연접 아이디얼에 대응하는
닫힌 부분스킴들을 각각 $X'' \subset X'$ 및
$Y'' \subset Y \times_X X'$이라 하자. 명백히
$Y'' \subset Y \times_X X''$은 $Y \times_X (E \cap X'')$ 위에 지지된
% P05-EN-CH38-0283
$\mathcal{O}_{Y \times_Y X''}$의 단면들의 준연접 아이디얼에 대응하는
닫힌 부분스킴이다. 따라서 $Y''$은 블로업 $X'' \to X$에 관한 $Y$의
강변환이다. Divisors, Definition
\ref{divisors-definition-strict-transform}을 보라. 따라서 Divisors, Lemma
\ref{divisors-lemma-strict-transform}에 의해 $Y''$은 $Y$ 위에서
$Y \times_X Z$의 블로업임을 알 수 있다.
% P05-EN-CH38-0284
\end{proof}

\noindent
거의 블로업 사각형을 축소할 수 있다.

\begin{lemma}
\label{flat-lemma-shrink-almost-blow-up}
거의 블로업 사각형 (\ref{flat-equation-almost-blow-up-square})을 생각하자.
$W \to X'$를 유한표현 닫힌 몰입이라 하자. 다음은 동치이다.
\begin{enumerate}
\item $X' \setminus E$는 스킴론적으로 $W$에 포함된다.
\item $Z$에서의 $X$의 블로업 $X''$은 스킴론적으로 $W$에 포함된다.
\item 그림
$$
\xymatrix{
E \cap W \ar[d] \ar[r] & W \ar[d] \\
Z \ar[r] & X
}
$$
은 거의 블로업 사각형이다. 여기서 $E \cap W$는 스킴론적 교차이다.
\end{enumerate}
\end{lemma}

\begin{proof}
(1)을 가정하자. 그러면 전사 $\mathcal{O}_{X'} \to \mathcal{O}_W$는
열린집합 $X' \subset E$ 위에서 동형이다.
% P05-EN-CH38-0285
$X'' \subset X'$의 아이디얼 층은 $E$ 위에 지지된
$\mathcal{O}_{X'}$의 단면들이므로(거의 블로업 사각형의 우리의 정의에
의해) (2)가 참이라는 결론을 얻는다. (2)가 참이면 (3)이 성립한다.
(3)이 성립하면 $X'' \cap (X' \setminus E)$가 $X \setminus Z$와
동형이고, 다시 이는 $X' \setminus E$와 동형이므로 (1)이 성립한다.
\end{proof}

\noindent
실제 블로업은 주어진 임의의 거의 블로업의 축소들의 극한이다.

\begin{lemma}
\label{flat-lemma-blow-up-limit-almost-blow-up}
$X$가 준콤팩트이고 준분리인 거의 블로업 사각형
(\ref{flat-equation-almost-blow-up-square})을 생각하자. 그러면 $Z$에서의
$X$의 블로업 $X''$은
$$
X'' = \lim X'_i
$$
로 쓸 수 있다. 여기서 극한은 Lemma \ref{flat-lemma-shrink-almost-blow-up}의
동치 조건들을 만족하는 유한표현 닫힌 부분스킴
$X'_i \subset X'$들의 유향계에 관하여 취한다.
\end{lemma}

\begin{proof}
$\mathcal{I} \subset \mathcal{O}_{X'}$를 $X''$에 대응하는 준연접 아이디얼
층이라 하자. Properties, Lemma
\ref{properties-lemma-quasi-coherent-colimit-finite-type}에 의해
$\mathcal{I}$를 유한형 준연접 부분가군들의 여과 여극한
$\mathcal{I} = \colim \mathcal{I}_i$로 쓸 수 있다. 이 가군들과 닫힌
부분스킴들 $X'_i$가 $1$-대-$1$로 대응하므로 증명이 끝난다.
\end{proof}

\noindent
거의 블로업 사각형은 존재한다.

\begin{lemma}
\label{flat-lemma-almost-blow-up-square}
$X$를 준콤팩트 준분리 스킴이라 하고 $Z \subset X$를 유한형 준연접
아이디얼 층이 절단하는 닫힌 부분스킴이라 하자. 그러면
(\ref{flat-equation-almost-blow-up-square})에서와 같은 거의 블로업 사각형이
존재한다.
\end{lemma}

\begin{proof}
$X = \lim X_i$를 아핀 전이 사상들을 갖는 뇌터 스킴들의 역계의 유향
극한으로 쓸 수 있다. Limits, Proposition
\ref{limits-proposition-approximate}를 보라. $X$로 기저변환하면 닫힌 몰입
$Z \to X$가 되는 닫힌 몰입 $Z_i \to X_i$와 지표 $i$를 찾을 수 있다.
Limits, Lemmas \ref{limits-lemma-descend-finite-presentation} 및
\ref{limits-lemma-descend-closed-immersion-finite-presentation}을 보라.
$b_i : X'_i \to X_i$를 중심이 $Z_i$인 블로업이라 하자. 이는 블로업
사각형
$$
\xymatrix{
E_i \ar[r] \ar[d] & X'_i \ar[d]^{b_i} \\
Z_i \ar[r] & X_i
}
$$
을 만들어 낸다. 모든 사상은 뇌터 스킴들의 유한형 사상이므로
유한표현이다. 따라서 이는 거의 블로업 사각형이다. Lemma
\ref{flat-lemma-base-change-almost-blow-up}에 의해 이 그림을 $X$로
기저변환하면 원하는 거의 블로업 사각형을 얻는다.
\end{proof}

\noindent
거의 블로업 사각형은 Lemma \ref{flat-lemma-shrink-almost-blow-up}에서와 같은
축소를 제외하고 유일하다.

\begin{lemma}
\label{flat-lemma-almost-blow-up-unique}
$X$를 준콤팩트 준분리 스킴이라 하고 $Z \subset X$를 유한형 준연접
아이디얼 층이 절단하는 닫힌 부분스킴이라 하자. 거의 블로업 사각형들
(\ref{flat-equation-almost-blow-up-square})
$$
\xymatrix{
E_k \ar[r] \ar[d] & X_k' \ar[d] \\
Z \ar[r] & X
}
$$
이 $k = 1, 2$에 대하여 주어졌다고 하자. 그러면 거의 블로업 사각형
$$
\xymatrix{
E \ar[r] \ar[d] & X' \ar[d] \\
Z \ar[r] & X
}
$$
과 $E = i_k^{-1}(E_k)$를 만족하는 $X$ 위의 닫힌 몰입들
$i_k : X' \to X'_k$가 존재한다.
\end{lemma}

\begin{proof}
$X'' \to X$를 $X$에서의 $Z$의 블로업이라 하자.
% P05-EN-CH38-0286
$X''$을 $X'_1$과 $X'_2$ 모두의 닫힌 부분스킴으로 본다. Lemma
\ref{flat-lemma-blow-up-limit-almost-blow-up}에서와 같이
$X'' = \lim X'_{1, i}$로 쓰자. Limits, Proposition
\ref{limits-proposition-characterize-locally-finite-presentation}에 의해
포함들 $X'' \subset X'_{1, i}$ 및 $X'' \subset X'_2$와 일치하는 사상
$h : X'_{1, i} \to X'_2$와 $i$가 존재한다. Limits, Lemma
\ref{limits-lemma-eventually-closed-immersion}에 의해 어떤 $i' \geq i$에
대하여 $h$를 $X'_{1, i'}$로 제한한 사상은 닫힌 몰입이다. 이로써 증명이
끝난다.
\end{proof}

\noindent
유리한 상황에서 블로업에 대한 우리의 평탄화 기법은 거의 블로업으로
이어진다.

\begin{lemma}
\label{flat-lemma-flat-after-almost-blowing-up}
$Y$를 준콤팩트 준분리 스킴이라 하고 $X$를 $Y$ 위의 유한표현 스킴이라
하자. $V \subset Y$를 $X_V \to V$가 평탄이게 하는 준콤팩트 열린집합이라
하자. 그러면 가환 그림
$$
\xymatrix{
E \ar[ddd] \ar[rd] & & & D \ar[lll] \ar[ddd] \ar[ld] \\
& Y' \ar[d] & X' \ar[l] \ar[d] \\
& Y & X \ar[l] \\
Z \ar[ru] & & & T \ar[lll] \ar[lu]
}
$$
이 존재하여, 오른쪽 및 왼쪽 사각형은 거의 블로업 사각형이고 아래쪽 및
위쪽 사각형은 카르테시안이며, $Z \cap V = \emptyset$이고,
$X' \to Y'$은 평탄이다(그리고 유한표현이다).
\end{lemma}

\begin{proof}
$Y$가 뇌터 스킴이면 이 보조정리는 Lemma
\ref{flat-lemma-flat-after-blowing-up}에서 곧바로 따른다. 이 경우 블로업
사각형은 거의 블로업 사각형이기 때문이다(강변환들이 블로업이라는
사실도 사용한다). 일반적인 경우는 절대 뇌터 근사로 뇌터인 경우에
축약된다.

\medskip\noindent
$Y = \lim Y_i$를 아핀 전이 사상들을 갖는 뇌터 스킴들의 역계의 유향
극한으로 쓸 수 있다. Limits, Proposition
\ref{limits-proposition-approximate}를 보라. $Y$로 기저변환하면
$X \to Y$가 되는 유한표현 사상 $X_i \to Y_i$와 지표 $i$를 찾을 수 있다.
Limits, Lemmas \ref{limits-lemma-descend-finite-presentation}를 보라.
% P05-EN-CH38-0287
$i$를 늘린 뒤 $V$가 열린 부분스킴 $V_i \subset Y_i$의 역상이라고
가정해도 된다. Limits, Lemma \ref{limits-lemma-descend-opens}를 보라.
마지막으로 $i$를 늘린 뒤 $X_{i, V_i} \to V_i$가 평탄이라고 가정해도
된다. Limits, Lemma \ref{limits-lemma-descend-flat-finite-presentation}을
보라. 뇌터인 경우에 의해 $X_i \to Y_i \supset V_i$에 대하여 보조정리와
같은 그림을 구성할 수 있다. 이 그림을 $Y \to Y_i$에 의해 기저변환하면
해를 얻는다. 기저변환이 사상들의 성질을 보존한다는 사실(Morphisms,
Lemmas \ref{morphisms-lemma-base-change-proper},
\ref{morphisms-lemma-base-change-finite-presentation},
\ref{morphisms-lemma-base-change-closed-immersion},
\ref{morphisms-lemma-base-change-flat})과 거의 블로업 사각형의 기저변환이
거의 블로업 사각형이라는 사실(Lemma
\ref{flat-lemma-base-change-almost-blow-up})을 사용하라.
\end{proof}

\begin{lemma}
\label{flat-lemma-blow-up-square-h}
$\mathcal{F}$를 Definition \ref{flat-definition-big-small-h}에서 구성한 사이트
$(\Sch/S)_h$ 중 하나 위의 층이라 하자. 그러면 범주 $(\Sch/S)_h$ 안의
모든 거의 블로업 사각형 (\ref{flat-equation-almost-blow-up-square})에 대하여
그림
$$
\xymatrix{
\mathcal{F}(E) & \mathcal{F}(X') \ar[l] \\
\mathcal{F}(Z) \ar[u] & \mathcal{F}(X) \ar[u] \ar[l]
}
$$
은 집합들의 범주에서 카르테시안이다.
\end{lemma}

\begin{proof}
$Z \amalg X' \to X$는 유한표현 고유 전사 사상이므로
$\{Z \amalg X' \to X\}$는 h 덮개이다(Lemma
\ref{flat-lemma-surjective-proper-finite-presentation-h}). 또한
$$
(Z \amalg X') \times_X (Z \amalg X') =
Z \amalg E \amalg E \amalg X' \times_X X'
$$
이다. $\mathcal{F}$는 Zariski 층이므로 $\mathcal{F}$는 서로소 합을 곱으로
보냄을 알 수 있다. 따라서 덮개 $\{Z \amalg X' \to X\}$에 대한 층
조건은 $\mathcal{F}(X) \to \mathcal{F}(Z) \times \mathcal{F}(X')$가
단사이고, 그 상이 (a) $t|_E = s'|_E$이고 (b) $s'$가 두 사상
$\mathcal{F}(X') \to \mathcal{F}(X' \times_X X')$의 등화자에 속하는 쌍
$(t, s')$들의 집합이라는 뜻이다. 다음으로 자명한 사상
$$
E \times_Z E \amalg X' \longrightarrow X' \times_X X'
$$
은 유한표현 고유 전사 사상임에 유의하자. 실제로 $b$는 동형
$X' \setminus E \to X \setminus Z$를 유도한다. 따라서
$\mathcal{F}(X' \times_X X') \to
\mathcal{F}(E \times_Z E) \times \mathcal{F}(X')$가 단사라는 결론을
얻는다. 이에 따라 (a) $\Rightarrow$ (b)이고, 이는 보조정리가 참이라는
뜻이다.
\end{proof}

\begin{lemma}
\label{flat-lemma-thickening-h}
$\mathcal{F}$를 Definition \ref{flat-definition-big-small-h}에서 구성한 사이트
$(\Sch/S)_h$ 중 하나 위의 층이라 하자. $X \to X'$를 비후화이면서
유한표현인 $(\Sch/S)_h$의 사상이라 하자. 그러면
$\mathcal{F}(X') \to \mathcal{F}(X)$는 전단사이다.
\end{lemma}

\begin{proof}
첫째 증명. $X \to X'$는 유한표현 고유 전사 사상이고
$X \times_{X'} X = X$임에 유의하자. h 덮개 $\{X \to X'\}$에 대한 층
성질(Lemma \ref{flat-lemma-surjective-proper-finite-presentation-h})에서 결론이
따른다.

\medskip\noindent
둘째 증명(우스꽝스러운 증명). $X$에서의 $X'$의 블로업은 공스킴이다.
그 이유는 $a$가 $A$의 멱영 원소이면 아핀 블로업 대수
$A[\frac{I}{a}]$(Algebra, Section \ref{algebra-section-blow-up})가 영이기
때문이다. 세부 사항은 생략한다. 따라서 다음 꼴의 거의 블로업 사각형을
얻는다.
$$
\xymatrix{
\emptyset \ar[r] \ar[d] & \emptyset \ar[d] \\
X \ar[r] & X'
}
$$
$\mathcal{F}$는 층이므로 $\mathcal{F}(\emptyset)$은 한원소집합이다.
Lemma \ref{flat-lemma-blow-up-square-h}를 적용하면 결론을 얻는다.
\end{proof}

\begin{proposition}
\label{flat-proposition-check-h}
$\mathcal{F}$를 Definition \ref{flat-definition-big-small-h}에서 구성한 사이트
$(\Sch/S)_h$ 중 하나 위의 전층이라 하자. 그러면 $\mathcal{F}$가 층일
필요충분조건은 다음 조건들이 성립하는 것이다.
\begin{enumerate}
\item $\mathcal{F}$는 Zariski 위상에 대한 층이다.
\item $Y$가 아핀이고 $f$가 전사, 평탄, 고유, 유한표현인
$(\Sch/S)_h$의 사상 $f : X \to Y$가 주어지면 $\mathcal{F}(Y)$는 두 사상
$\mathcal{F}(X) \to \mathcal{F}(X \times_Y X)$의 등화자이다.
\item 범주 $(\Sch/S)_h$에서 $X$가 아핀인 거의 블로업 사각형
(\ref{flat-equation-almost-blow-up-square})이 주어지면 그림
$$
\xymatrix{
\mathcal{F}(E) & \mathcal{F}(X') \ar[l] \\
\mathcal{F}(Z) \ar[u] & \mathcal{F}(X) \ar[u] \ar[l]
}
$$
은 집합들의 범주에서 카르테시안이다.
\end{enumerate}
\end{proposition}

\begin{proof}
$\mathcal{F}$가 층이라고 하자. Zariski 덮개는 h 덮개이므로 조건 (1)이
성립한다. Lemma \ref{flat-lemma-zariski-h}를 보라. (2)에서와 같은 $f$에
대하여 $\{X \to Y\}$는 fppf 덮개이고(이는 명백하다), 따라서 h
덮개이므로 조건 (2)가 성립한다. Lemma \ref{flat-lemma-zariski-h}를 보라.
조건 (3)은 Lemma \ref{flat-lemma-blow-up-square-h}에서 따른다.

\medskip\noindent
역으로 $\mathcal{F}$가 (1), (2), (3)을 만족한다고 하자. Lemma
\ref{flat-lemma-characterize-sheaf-h}를 적용하여 $\mathcal{F}$가 층임을
증명하겠다. $Y$가 아핀인 $(\Sch/S)_h$의 유한표현 고유 전사 사상
$f : X \to Y$를 생각하자. $\mathcal{F}(Y)$가 두 사상
$\mathcal{F}(X) \to \mathcal{F}(X \times_Y X)$의 등화자임을 보이면
충분하다.

\medskip\noindent
먼저 $f : X \to Y$가 또한 닫힌 몰입이라고(다시 말해 $f$가
비후화라고) 하자. 그러면 $X$에서의 $Y$의 블로업은 공스킴이고, 이는
꼭짓점들이 $\emptyset, \emptyset, X, Y$인 거의 블로업 사각형을 만든다
(Lemma \ref{flat-lemma-thickening-h}의 둘째 증명과 비교하라).
% P05-EN-CH38-0288
따라서 조건 (3)에 의해
$$
\xymatrix{
\mathcal{F}(\emptyset) & \mathcal{F}(\emptyset) \ar[l] \\
\mathcal{F}(X) \ar[u] & \mathcal{F}(Y) \ar[u] \ar[l]
}
$$
이 집합들의 범주에서 카르테시안임을 알 수 있다. $\mathcal{F}$는 Zariski
위상에 대한 층이므로 $\mathcal{F}(\emptyset)$은 한원소집합이다. 따라서
$\mathcal{F}(X) = \mathcal{F}(Y)$임을 알 수 있다.

\medskip\noindent
막간 A: $T \to T'$를 비후화이고 유한표현인 $(\Sch/S)_h$의 사상이라
하자. 그러면 $\mathcal{F}(T') \to \mathcal{F}(T)$는 전단사이다. 실제로
아핀 열린 덮개 $T' = \bigcup T'_i$를 택하고
$T_i = T \times_{T'} T'_i$를 대응하는 $T$의 아핀 열린집합들이라 하자.
그러면 앞 문단의 결과에 의해 모든 $i$에 대하여
$\mathcal{F}(T'_i) \to \mathcal{F}(T_i)$가 전단사이다. Zariski 층 성질을
사용하면 $\mathcal{F}(T') \to \mathcal{F}(T)$가 단사임을 알 수 있다.
논증을 반복하면 이것이 전단사임을 알 수 있다. 사소한 세부 사항은
생략한다.

\medskip\noindent
막간 B: 범주 $(\Sch/S)_h$ 안의 거의 블로업 사각형
(\ref{flat-equation-almost-blow-up-square})을 생각하자. 그림
$$
\xymatrix{
\mathcal{F}(E) & \mathcal{F}(X') \ar[l] \\
\mathcal{F}(Z) \ar[u] & \mathcal{F}(X) \ar[u] \ar[l]
}
$$
이 집합들의 범주에서 카르테시안이라고 주장한다. $X$의 아핀 열린
덮개를 택하고 막간 A에서와 같이 논증하면 이는 조건 (3)에서 따른다.
세부 사항은 생략한다.

\medskip\noindent
다음으로 $Y$가 아핀인 $(\Sch/S)_h$의 유한표현 고유 전사 사상
$f : X \to Y$를 택하자. $f : X \to Y$에 대하여 More on Morphisms,
Lemma \ref{more-morphisms-lemma-generic-flatness-stratification-scheme}에서와
같은 일반 평탄성 층화
$$
Y \supset Y_0 \supset Y_1 \supset \ldots \supset Y_t = \emptyset
$$
를 택하자. 앞으로는 이 층화의 모든 성질을 더 언급하지 않고
사용하겠다. $X_0 = X \times_Y Y_0$라 두자. 막간 B에 의해
$\mathcal{F}(Y_0) = \mathcal{F}(Y)$,
$\mathcal{F}(X_0) = \mathcal{F}(X)$,
$\mathcal{F}(X_0 \times_{Y_0} X_0) = \mathcal{F}(X \times_Y X)$이다.

\medskip\noindent
$t$에 관한 귀납법으로 결과를 증명하겠다. $t = 1$이면
$X_0 \to Y_0$는 전사, 고유, 평탄, 유한표현이므로 성질 (2)에 의해
결과가 성립한다. $t > 1$이면 $Y$를 $Y_0$으로, $X$를 $X_0$으로
바꾸어(위를 보라) $Y = Y_0$이라고 가정해도 된다.

\medskip\noindent
준콤팩트 열린 부분스킴 $V = Y \setminus Y_1 = Y_0 \setminus Y_1$을
생각하자. $f : X \to Y \supset V$에 대하여 Lemma
\ref{flat-lemma-flat-after-almost-blowing-up}에서와 같은 그림
$$
\xymatrix{
E \ar[ddd] \ar[rd] & & & D \ar[lll] \ar[ddd] \ar[ld] \\
& Y' \ar[d] & X' \ar[l] \ar[d] \\
& Y & X \ar[l] \\
Z \ar[ru] & & & T \ar[lll] \ar[lu]
}
$$
을 택하자. 그러면 $f' : X' \to Y'$은 평탄이고 유한표현이다. 또한
$f'$은 고유이다(이를 보려면 Morphisms, Lemmas
\ref{morphisms-lemma-composition-proper} 및
\ref{morphisms-lemma-image-proper-scheme-closed}를 사용하라). 따라서 상
$W = f'(X') \subset Y'$은 $Y'$의 열린(Morphisms, Lemma
\ref{morphisms-lemma-fppf-open}) 닫힌 부분스킴이다. $Y' \setminus E$가
$W$에 포함됨에 유의하자. Lemma \ref{flat-lemma-shrink-almost-blow-up}에 의해
이는 위 그림에서 $Y'$을 $W$로 바꾸어도 된다는 뜻이다. 다시 말해
$f'$이 전사라고 가정해도 되며 그렇게 하자. 이제
$$
\vcenter{
\xymatrix{
\mathcal{F}(E) & \mathcal{F}(Y') \ar[l] \\
\mathcal{F}(Z) \ar[u] & \mathcal{F}(Y) \ar[u] \ar[l]
}
}
\quad\text{이고}\quad
\vcenter{
\xymatrix{
\mathcal{F}(D) & \mathcal{F}(X') \ar[l] \\
\mathcal{F}(T) \ar[u] & \mathcal{F}(X) \ar[u] \ar[l]
}
}
$$
는 막간 B에 의해 카르테시안이다. $Z \cap Y_1 \to Z$는 유한표현
비후화임에 유의하자($Z$는 $V$와 서로소인 $Y$의 닫힌 부분스킴으로서
집합론적으로 $Y_1$에 포함되기 때문이다). 따라서 $f$를 $T$로 제한한
$T = Z \times_Y X \to Z$에 대하여 위와 같은 여과
$$
Z \supset Z \cap Y_1 \supset Z \cap Y_2 \subset \ldots \subset
Z \cap Y_t = \emptyset
$$
를 얻는다.
% P05-EN-CH38-0289
따라서 귀납가정에 의해 $\mathcal{F}(Z) \to \mathcal{F}(T)$는 집합들의
단사 사상이고 그 상은 두 사상
$\mathcal{F}(T) \to \mathcal{F}(T \times_Z T)$의 등화자이다.

\medskip\noindent
$s \in \mathcal{F}(X)$가 두 사상
$\mathcal{F}(X) \to \mathcal{F}(X \times_Y X)$의 등화자에 속한다고 하자.
위의 결과에 의해 제한 $s|_T$는 유일한 원소
$t \in \mathcal{F}(Z)$에서 오고, 마찬가지로 제한 $s|_{X'}$는 유일한
원소 $t' \in \mathcal{F}(Y')$에서 온다. 위의 거대한 가환 그림의
화살표들에 대응하는 $\mathcal{F}$의 제한 사상들을 따라 단면들을
추적하면, $t$와 $t'$가 $\mathcal{F}(E)$의 같은 원소로 제한됨을 알 수
있다. 실제로 이들은 $\mathcal{F}(D)$의 같은 원소로 제한되고 (2)가
성립한다. 여기서 $D \to E$가 $X' \to Y'$의 제한으로서 전사, 평탄,
고유, 유한표현이라는 사실을 사용한다. 따라서 위에 표시한 두 카르테시안
사각형 중 첫째 것에 의해 $Z$와 $Y'$에서 각각 $t$와 $t'$로 제한되는
유일한 단면 $u \in \mathcal{F}(Y)$를 얻는다. $u$가 $X$에서 $s$로
제한됨을 보려면 둘째 그림을 사용하라.
% P05-EN-CH38-0290
\end{proof}

\begin{example}
\label{flat-example-one-generator}
$A$를 환이라 하자. $f \in A$를 원소라 하자. $J \subset A$를
$f$의 어떤 거듭제곱에 의해 소멸되는 유한생성 아이디얼이라 하자. 그러면
$$
\xymatrix{
E = \Spec(A/fA + J) \ar[r] \ar[d] & \Spec(A/J) = X' \ar[d] \\
Z = \Spec(A/fA) \ar[r] & \Spec(A) = X
}
$$
은 거의 블로업 사각형이다.
\end{example}

\begin{example}
\label{flat-example-two-generators}
$A$를 환이라 하고 $f_1, f_2 \in A$를 원소라 하자.
$$
\xymatrix{
E = \text{Proj}(A/(f_1, f_2)[T_0, T_1]) \ar[r] \ar[d] &
\text{Proj}(A[T_0, T_1]/(f_2 T_0 - f_1 T_1) = X' \ar[d] \\
Z = \Spec(A/(f_1, f_2)) \ar[r] & \Spec(A) = X
}
$$
은 거의 블로업 사각형이다.
% P05-EN-CH38-0291
\end{example}

\begin{lemma}
\label{flat-lemma-refine-check-h}
정의 \ref{flat-definition-big-small-h}에서 구성한 사이트 $(\Sch/S)_h$ 중 하나 위의
전층을 $\mathcal{F}$라 하자. 그러면 $\mathcal{F}$가 층일 필요충분조건은 다음
조건들이 성립하는 것이다.
\begin{enumerate}
\item $\mathcal{F}$는 자리스키 위상에 대한 층이다.
\item $Y$가 아핀이고 $f$가 전사, 평탄, 고유, 유한표현인 $(\Sch/S)_h$의
사상 $f : X \to Y$가 주어지면, $\mathcal{F}(Y)$는 두 사상
$\mathcal{F}(X) \to \mathcal{F}(X \times_Y X)$의 등화자이다.
\item $\mathcal{F}$는 범주 $(\Sch/S)_h$에서 예
\ref{flat-example-one-generator}과 같은 거의 블로업 사각형을 집합들의 카르테시안
그림으로 보낸다.
\item $\mathcal{F}$는 범주 $(\Sch/S)_h$에서 예
\ref{flat-example-two-generators}와 같은 거의 블로업 사각형을 집합들의 카르테시안
그림으로 보낸다.
\end{enumerate}
\end{lemma}

\begin{proof}
명제 \ref{flat-proposition-check-h}에 의해 다음을 보이면 충분하다. 범주
$(\Sch/S)_h$에서 $X$가 아핀인 거의 블로업 사각형
(\ref{flat-equation-almost-blow-up-square})이 주어지면 그림
$$
\xymatrix{
\mathcal{F}(E) & \mathcal{F}(X') \ar[l] \\
\mathcal{F}(Z) \ar[u] & \mathcal{F}(X) \ar[u] \ar[l]
}
$$
은 집합의 범주에서 카르테시안이다. 증명의 대략적인 착상은 이 사상을 국소적으로
가정 (3)과 (4)를 적용할 수 있는 다른 거의 블로업 사각형들로 우세화하는 것이다.

\medskip\noindent
범주 $(\Sch/S)_h$의 거의 블로업 사각형
(\ref{flat-equation-almost-blow-up-square}), 열린 덮개 $X = \bigcup U_i$, 그리고
다음 그림들이 카르테시안이 되게 하는 열린 덮개
$U_i \cap U_j = \bigcup U_{ijk}$가 있다고 하자.
$$
\vcenter{
\xymatrix{
\mathcal{F}(E \cap b^{-1}(U_i)) & \mathcal{F}(b^{-1}(U_i)) \ar[l] \\
\mathcal{F}(Z \cap U_i) \ar[u] & \mathcal{F}(U_i) \ar[u] \ar[l]
}
}
\quad\text{이고}\quad
\vcenter{
\xymatrix{
\mathcal{F}(E \cap b^{-1}(U_{ijk})) &
\mathcal{F}(b^{-1}(U_{ijk})) \ar[l] \\
\mathcal{F}(Z \cap U_{ijk}) \ar[u] &
\mathcal{F}(U_{ijk}) \ar[u] \ar[l]
}
}
$$
그러면 다음 그림도 카르테시안이다.
$$
\xymatrix{
\mathcal{F}(E) & \mathcal{F}(X') \ar[l] \\
\mathcal{F}(Z) \ar[u] & \mathcal{F}(X) \ar[u] \ar[l]
}
$$
이는 $\mathcal{F}$가 자리스키 위상에서 층이라는 사실로부터 따른다.

\medskip\noindent
특히 $b : X' \to X$가 닫힌 몰입이고 $Z$가 국소 주 아이디얼 닫힌 부분스킴인
블로업 사각형 (\ref{flat-equation-almost-blow-up-square})이 있으면
$\mathcal{F}(X) = \mathcal{F}(X') \times_{\mathcal{F}(E)} \mathcal{F}(Z)$임을
알 수 있다. 실제로 $X$ 위 아핀 국소적으로 (3)과 같은 거의 블로업 사각형을
얻는다.
% P05-EN-CH38-0292

\medskip\noindent
$Z \subset X$, $E_k \subset X'_k \to X$, $E \subset X' \to X$, 그리고
$i_k : X' \to X'_k$가 보조정리 \ref{flat-lemma-almost-blow-up-unique}의 명제와
같다고 하자. 그러면
$$
\xymatrix{
E \ar[d] \ar[r] & X' \ar[d] \\
E_k \ar[r] & X'_k
}
$$
은 앞 단락에서 논한 종류의 거의 블로업 사각형이다. 따라서 앞 단락의 결과에
의해 $k = 1, 2$에 대하여
$$
\mathcal{F}(X'_k) = \mathcal{F}(X') \times_{\mathcal{F}(E)} \mathcal{F}(E_k)
$$
이고, 그러므로
$$
\mathcal{F}(X) \longrightarrow
\mathcal{F}(X'_k) \times_{\mathcal{F}(E_k)} \mathcal{F}(Z)
$$
가 $k = 1$에 대하여 전단사일 필요충분조건은 $k = 2$에 대하여 전단사인 것이다.
따라서 $X$가 준콤팩트하고 준분리인 유한표현 닫힌 몰입 $Z \to X$가 주어지면,
$\mathcal{F}(X) = \mathcal{F}(X') \times_{\mathcal{F}(E)} \mathcal{F}(Z)$
의 성립 여부는 선택한 거의 블로업 사각형
(\ref{flat-equation-almost-blow-up-square})의 선택과 무관하다. (더욱이 보조정리
\ref{flat-lemma-almost-blow-up-square}에 의해 $Z \subset X$에 대한 거의 블로업
사각형이 실제로 존재한다.)

\medskip\noindent
끝으로 $(\Sch/S)_h$의 아핀 대상 $X$와 유한표현 닫힌 몰입 $Z \to X$를
생각하자. $X$ 안에서 $Z$를 정의하는 아이디얼의 생성원 수 $r$에 관한 귀납법으로
쌍 $(X, Z)$에 대해 원하는 성질을 증명하겠다. 생성원 수가 $\leq 2$이면 거의
블로업 사각형으로 예 \ref{flat-example-two-generators}의 것을 선택할 수 있으므로
가정 (4)에 의해 결론을 얻는다.

\medskip\noindent
귀납 단계. $r > 2$이고 $X = \Spec(A)$,
$Z = \Spec(A/(f_1, \ldots, f_r))$라고 하자. 쌍 $(X, Z)$에 대한 블로업
사각형 (\ref{flat-equation-almost-blow-up-square})을 선택하자.
$Z_1 = \Spec(A/(f_1, f_2))$라 놓고
$$
\xymatrix{
E_1 \ar[d] \ar[r] & Y \ar[d] \\
Z_1 \ar[r] & X
}
$$
를 예 \ref{flat-example-two-generators}에서 구성한 거의 블로업 사각형이라 하자.
보조정리 \ref{flat-lemma-base-change-almost-blow-up}에 의해 기저변환
$$
(I)
\vcenter{
\xymatrix{
Y \times_X E \ar[r] \ar[d] & Y \times_X X' \ar[d] \\
Y \times_X Z \ar[r] & Y
}
}
\quad\text{이고}\quad
(II)
\vcenter{
\xymatrix{
E \ar[r] \ar[d] & Z_1 \times_X X' \ar[d] \\
Z \ar[r] & Z_1
}
}
$$
은 거의 블로업 사각형들이다. $Z_1$ 안의 $Z$의 아이디얼은 $r - 2$개의
원소로 생성된다. $Y \times_X Z$의 아이디얼은 $f_1, \ldots, f_r$를 $Y$로
당긴 원소들로 생성된다. $Y$ 위 국소적으로 $f_1, f_2$가 생성하는 아이디얼은
하나의 원소로 생성될 수 있으므로, $Y \times_X Z$는 $Y$ 위 아핀 국소적으로
많아야 $r - 1$개의 원소에 의해 잘려 나온다. 귀납 가정과 위의 논의에 의해
$$
\mathcal{F}(Y) =
\mathcal{F}(Y \times_X X') \times_{\mathcal{F}(Y \times_X E)}
\mathcal{F}(Y \times_X Z)
$$
이고
$$
\mathcal{F}(Z_1) =
\mathcal{F}(Z_1 \times_X X') \times_{\mathcal{F}(E)}
\mathcal{F}(Z)
$$
가정 (4)에 의해
$$
\mathcal{F}(X) =
\mathcal{F}(Y) \times_{\mathcal{F}(E_1)} \mathcal{F}(Z_1)
$$
를 얻는다. 이제 $s' \in \mathcal{F}(X')$와 $t \in \mathcal{F}(Z)$가
$\mathcal{F}(E)$에서 같은 제한을 갖는 쌍 $(s', t)$를 이룬다고 하자. 그러면
$(s'|{Z_1 \times_X X'}, t)$는 유일한 원소
$t_1 \in \mathcal{F}(Z_1)$의 상이다.
% P05-EN-CH38-0293
마찬가지로 $(s'|_{Y \times_X X'}, t|_{Y \times_X Z})$는 유일한 원소
$s_Y \in \mathcal{F}(Y)$의 상이다. $s_Y$와 $t_1$이
$\mathcal{F}(E_1)$의 같은 원소로 제한됨을 주장한다. 이는 거의 블로업 사각형
$$
\xymatrix{
E_1 \times_X E \ar[r] \ar[d] & E_1 \times_X X' \ar[d] \\
E_1 \times_X Z \ar[r] & E_1
}
$$
이 $E_1 \to Y$에 의한 거의 블로업 사각형 (I)의 기저변환인 동시에
$E_1 \to Z_1$에 의한 거의 블로업 사각형 (II)의 기저변환이고, $s_Y$와
$t_1$을 구성하는 데 사용한 단면의 쌍들이 서로 맞기 때문이다. 따라서 셋째
올곱 등식에 의해 $\mathcal{F}(Y)$에서 $s_Y$로, $\mathcal{F}(Z)$에서
$t_1$로 가는 유일한 $s \in \mathcal{F}(X)$가 있음을 알 수 있다.
% P05-EN-CH38-0294
$s$가 $\mathcal{F}(X')$에서 $s'$로, $\mathcal{F}(Z)$에서 $t$로 감의 확인은
생략한다. 도움말: 방금 구성한 $s$의 유일성을 사용하고 아핀 국소적으로 작업하라.
\end{proof}

\begin{lemma}
\label{flat-lemma-refine-check-h-stack}
$p : \mathcal{S} \to (\Sch/S)_h$를 그루포이드로 섬유화된 범주라 하자.
그러면 $\mathcal{S}$가 그루포이드 스택일 필요충분조건은 다음 조건들이
성립하는 것이다.
\begin{enumerate}
\item $\mathcal{S}$는 자리스키 위상에 대한 그루포이드 스택이다.
\item $Y$가 아핀이고 $f$가 전사, 평탄, 고유, 유한표현인 $(\Sch/S)_h$의
사상 $f : X \to Y$가 주어지면
$$
\mathcal{S}_Y \longrightarrow
\mathcal{S}_X \times_{\mathcal{S}_{X \times_Y X}} \mathcal{S}_X
$$
는 범주의 동치이다.
\item 범주 $(\Sch/S)_h$에서 예 \ref{flat-example-one-generator} 또는
\ref{flat-example-two-generators}와 같은 거의 블로업 사각형에 대하여 함자
$$
\mathcal{S}_X \longrightarrow
\mathcal{S}_Z \times_{\mathcal{S}_E} \mathcal{S}_{X'}
$$
는 범주의 동치이다.
\end{enumerate}
\end{lemma}

\begin{proof}
이 보조정리는 보조정리 \ref{flat-lemma-refine-check-h}와 그루포이드 스택에 대한
우리의 정의의 형식적 귀결이다. 예를 들어 (1), (2), (3)을 가정하자.
$\mathcal{S}$가 스택임을 보이려면 사상과 대상에 대한 강하를 증명해야 한다.
스택, 정의 \ref{stacks-definition-stack-in-groupoids}를 보라.

\medskip\noindent
$x, y$가 $(\Sch/S)_h$의 대상 $U$ 위에 놓인 $\mathcal{S}$의 대상들이면,
우리의 가정에 의해 $\mathit{Isom}(x, y)$는 보조정리
\ref{flat-lemma-refine-check-h}의 (1), (2), (3), (4)를 만족하는 $(\Sch/U)_h$
위의 전층이고 따라서 층이다. 몇몇 세부사항은 생략한다.

\medskip\noindent
$\{U_i \to U\}_{i \in I}$를 $(\Sch/S)_h$의 덮개라 하자.
$(x_i, \varphi_{ij})$를 족 $\{U_i \to U\}_{i \in I}$에 대한
$\mathcal{S}$의 강하 데이터라 하자. 스택, 정의
\ref{stacks-definition-descent-data}를 보라. $(\Sch/U)_h$의 $V/U$에 다음과 같은
쌍 $(y, \psi_i)$들의 집합을 대응시키는 규칙 $F$를 생각하자. 여기서 $y$는
$\mathcal{S}_V$의 대상이고,
$\psi_i : y|_{U_i \times_U V} \to x_i|_{U_i \times_U V}$는
$U_i \times_U V$ 위의 $\mathcal{S}$의 사상이며 다음을 만족한다.
$$
\varphi_{ij}|_{U_i \times_U U_j \times_U V}
\circ
\psi_i|_{U_i \times_U U_j \times_U V}
=
\psi_j|_{U_i \times_U U_j \times_U V}
$$
이는 동형을 제외하고 성립한다. 사상에 대한 강하는 이미 확보했으므로
$F(V/U)$가 공집합이거나 한원소 집합임이 분명하다. 한편
$F(U_{i_0}/U)$는 $(x_{i_0}, \varphi_{i_0i})$를 포함하므로 공집합이 아니다.
목표는 $F(U/U)$가 공집합이 아님을 증명하는 것이므로 $F$가 $(\Sch/U)_h$
위의 층임을 보이면 충분하다. 이를 위해 보조정리
\ref{flat-lemma-refine-check-h}의 판정법을 사용할 수 있다. 그런데 우리의 가정
(1), (2), (3)은 (여기서는 생략하는 몇몇 가환 그림을 그려 보면) 보조정리
\ref{flat-lemma-refine-check-h}의 성질 (1), (2), (3), (4)가 $F$에 대하여
성립함을 함의한다.

\medskip\noindent
$\mathcal{S}$가 그루포이드 스택이면 (1), (2), (3)이 성립한다는 확인은
생략한다.
\end{proof}





\section{절대 약한 정규화와 h 덮개}
\label{flat-section-h-and-weak-normalization}

\noindent
이 절에서는 절 \ref{flat-section-blow-up-h}에서 얻은 판정법을 사용하여 몇 가지
h 층을 제시하고, 구조층의 h 층화와 절대 약한 정규화를 연관 짓는다. 이를 위해
다음 초등적 보조정리가 필요하다.

\begin{lemma}
\label{flat-lemma-funny-blow-up}
$Z, X, X', E$를 예 \ref{flat-example-two-generators}과 같은 거의 블로업
사각형이라 하자. 그러면 $p > 0$에 대하여
$H^p(X', \mathcal{O}_{X'}) = 0$이고,
$\Gamma(X, \mathcal{O}_X) \to \Gamma(X', \mathcal{O}_{X'})$는 핵이
제곱 영 아이디얼인 환의 전사 사상이다.
\end{lemma}

\begin{proof}
먼저 $A = \mathbf{Z}[f_1, f_2]$가 다항식환이라고 가정하자. 이 경우 우리의
거의 블로업 사각형은 닫힌 부분스킴 $Z$에서 $X = \Spec(A)$를 블로업한 것이며,
실제로 $X' \subset \mathbf{P}^1_X$는 $\mathcal{O}_{\mathbf{P}^1_X}(1)$의
전역단면 $f_2T_0 - f_1 T_1$에 의해 잘려 나오는 유효 카르티에 인자이다.
따라서 분해
$$
0 \to
\mathcal{O}_{\mathbf{P}^1_X}(-1) \to
\mathcal{O}_{\mathbf{P}^1_X} \to
\mathcal{O}_{X'} \to
0
$$
를 얻는다. 스킴의 코호몰로지, 절
\ref{coherent-section-cohomology-projective-space}에 주어진 코호몰로지의
기술을 사용하면, 이 경우
$\Gamma(X, \mathcal{O}_X) \to \Gamma(X', \mathcal{O}_{X'})$가 동형이고
$H^1(X', \mathcal{O}_{X'}) = 0$임이 따른다.

\medskip\noindent
다음으로 예 \ref{flat-example-two-generators}과 같은 모든 그림은 앞 단락의
그림을 환 사상 $\mathbf{Z}[f_1, f_2] \to A$로 기저변환한 것임을 관찰하자.
그러므로 사상론 보충, 보조정리
\ref{more-morphisms-lemma-check-h1-fibre-zero},
\ref{more-morphisms-lemma-h1-fibre-zero}, 및
\ref{more-morphisms-lemma-h1-fibre-zero-check-h0-kappa}
에 의해 일반적으로 $H^1(X', \mathcal{O}_{X'})$가 영이고 사상
$H^0(X, \mathcal{O}_X) \to H^0(X', \mathcal{O}_{X'})$가 전사임을 얻는다.
% P05-EN-CH38-0295

\medskip\noindent
다음으로 일반적인 경우의 핵을 조사하자. $a \in A$가 영으로 가면 아핀
차트들에서 살펴보아
$$
a = (f_1x - f_2)(a_0 + a_1x + \ldots + a_rx^r)\text{는 }A[x]\text{에서 성립}
$$
가 어떤 $r \geq 0$ 및 $a_0, \ldots, a_r \in A$에 대하여 성립하고, 마찬가지로
$$
a = (f_1 - f_2y)(b_0 + b_1y + \ldots + b_s y^s)\text{는 }A[y]\text{에서 성립}
$$
가 어떤 $s \geq 0$ 및 $b_0, \ldots, b_s \in A$에 대하여 성립함을 알 수 있다.
이는
$$
a = f_2 a_0,\ f_1 a_0 = f_2 a_1,\ \ldots,\ f_1 a_r = 0,
\ a = f_1 b_0,\ f_2 b_0 = f_1 b_1,\ \ldots,\ f_2 b_s = 0
$$
을 뜻한다. $(a', r', a'_i, s', b'_j)$가 이와 같은 둘째 계라면
$$
aa' = f_1f_2a_0b'_0 = f_1f_2a_1b'_1 = f_1f_2a_2b'_2 = \ldots = 0
$$
을 얻으며, 이것이 원하는 결론이다.
\end{proof}

\noindent
$\mathbf{F}_p$-대수 $A$에 대하여 $\colim_F A$를 계
$$
A \xrightarrow{F} A \xrightarrow{F} A \xrightarrow{F} \ldots
$$
의 여극한으로 정의한다. 여기서 $F : A \to A$, $a \mapsto a^p$는 Frobenius
자기준동형이다.

\begin{lemma}
\label{flat-lemma-h-sheaf-colim-F}
$p$를 소수라 하고 $S$를 $\mathbf{F}_p$ 위의 스킴이라 하자.
$(\Sch/S)_h$를 정의 \ref{flat-definition-big-small-h}와 같은 사이트라 하자.
$(\Sch/S)_h$의 임의의 준콤팩트 준분리 대상 $X$에 대하여
$$
\mathcal{F}(X) = \colim_F \Gamma(X, \mathcal{O}_X)
$$
를 만족하는 $(\Sch/S)_h$ 위의 유일한 층 $\mathcal{F}$가 존재한다.
\end{lemma}

\begin{proof}
함자
$$
X \longrightarrow \colim_F \Gamma(X, \mathcal{O}_X)
$$
의 자리스키 층화를 $\mathcal{F}$라 쓰자. 층, 보조정리
\ref{sheaves-lemma-directed-colimits-sections}와 $\mathcal{O}$가 자리스키 위상에
대한 층이라는 사실에 의해, 준콤팩트 준분리 스킴 $X$에 대하여
$\mathcal{F}(X) = \colim_F \Gamma(X, \mathcal{O}_X)$이다.
% P05-EN-CH38-0296
% P05-EN-CH38-0297
따라서 $\mathcal{F}$가 h 층임을 보이면 충분하다. 이를 증명하기 위해 보조정리
\ref{flat-lemma-refine-check-h}의 조건 (1), (2), (3), (4)를 확인한다. (거의
불필요한) 자리스키 층화를 수행했으므로 조건 (1)이 성립한다. $\mathcal{O}$가
fppf 층이고(강하, 보조정리 \ref{descent-lemma-sheaf-condition-holds}),
$A$가 $\mathbf{F}_p$-대수의 두 사상 $B \to C$의 등화자이면
$\colim_F A$가 두 사상 $\colim_F B \to \colim_F C$의 등화자이므로 조건
(2)가 성립한다.

\medskip\noindent
조건 (3)을 확인하자. $A, f, J$가 예 \ref{flat-example-one-generator}과 같다고
하자. 다음을 보여야 한다.
$$
\colim_F A = \colim_F A/J \times_{\colim_F A/fA + J} \colim_F A/fA
$$
이는 다음 대수 문제로 환원된다. $a', a'' \in A$가
$F^n(a' - a'') \in fA + J$를 만족한다고 하자. $a - F^m(a') \in J$ 및
$a - F^m(a'') \in fA$를 만족하는 $a \in A$와 $m \geq 0$을 찾고, 쌍
$(a, m)$이 $(a, m) \mapsto (F(a), m + 1)$ 꼴의 치환까지 유일하게 결정됨을
보여라. 이를 위해 $h \in A$와 $g \in J$를 써서 단지
$F^n(a' - a'') = f h + g$라 쓰고,
$a = F^n(a') - g = F^n(a'') + fh$ 및 $m = n$으로 놓으면 된다.
유일성을 보이기 위해 $(a_1, m_1)$이 둘째 해라고 하자. 위와 같은 꼴의
치환을 사용하여 $m = m_1$이라 가정할 수 있다. 그러면
$a - a_1 \in J$이고 $a - a_1 \in fA$임을 알 수 있다. $J$는 $f$의 어떤
거듭제곱에 의해 소멸되므로 $a - a_1$은 멱영원이다. 따라서 충분히 큰 어떤
$k$에 대하여 $F^k(a - a_1)$은 영이다. 그러므로 치환을 더 수행한 뒤
$a = a_1$을 얻는다.

\medskip\noindent
조건 (4)를 확인하자. $X, X', Z, E$가 예
\ref{flat-example-two-generators}과 같다고 하자. 보조정리
\ref{flat-lemma-funny-blow-up}에 의해
$$
\mathcal{F}(X) = \colim_F \Gamma(X, \mathcal{O}_X)
\longrightarrow
\colim_F \Gamma(X', \mathcal{O}_{X'}) = \mathcal{F}(X')
$$
가 전단사임을 알 수 있다. 이 경우 $E = \mathbf{P}^1_Z$이므로
$\mathcal{F}(Z) \to \mathcal{F}(E)$도 전단사임을 알 수 있다. 따라서 이
경우에도 결론이 성립한다.
\end{proof}

\noindent
$p$를 소수라 하자. $\mathbf{F}_p$-대수 $A$에 대하여 $\lim_F A$를 역계
$$
\ldots \xrightarrow{F} A \xrightarrow{F} A \xrightarrow{F} A
$$
의 극한으로 정의한다. 여기서 $F : A \to A$, $a \mapsto a^p$는 Frobenius
자기준동형이다.

\begin{lemma}
\label{flat-lemma-h-sheaf-lim-F}
$p$를 소수라 하고 $S$를 $\mathbf{F}_p$ 위의 스킴이라 하자.
$(\Sch/S)_h$를 정의 \ref{flat-definition-big-small-h}와 같은 사이트라 하자.
규칙
$$
\mathcal{F}(X) = \lim_F \Gamma(X, \mathcal{O}_X)
$$
는 $(\Sch/S)_h$ 위의 층을 정의한다.
\end{lemma}

\begin{proof}
$\mathcal{F}$가 층임을 증명하기 위해 보조정리
\ref{flat-lemma-refine-check-h}의 조건 (1), (2), (3), (4)를 확인하자. 층들의
극한은 층이고 $\mathcal{O}$는 자리스키 층이므로 조건 (1)이 성립한다.
$\mathcal{O}$가 fppf 층이고(강하, 보조정리
\ref{descent-lemma-sheaf-condition-holds}), $A$가 $\mathbf{F}_p$-대수의 두
사상 $B \to C$의 등화자이면 $\lim_F A$가 두 사상
$\lim_F B \to \lim_F C$의 등화자이므로 조건 (2)가 성립한다.

\medskip\noindent
조건 (3)을 확인하자. $A, f, J$가 예 \ref{flat-example-one-generator}과 같다고
하자. 다음 사상이 전단사임을 보여야 한다.
\begin{align*}
\lim_F A
& \to
\lim_F A/J \times_{\lim_F A/fA + J} \lim_F A/fA \\
& =
\lim_F (A/J \times_{A/fA + J} A/fA) \\
& =
\lim_F A/(fA \cap J)
\end{align*}
$J$가 $f$의 어떤 거듭제곱에 의해 소멸되므로
$\mathfrak a = fA \cap J$가 멱영 아이디얼임을, 즉
$\mathfrak a^n = 0$인 어떤 $n$이 존재함을 알 수 있다. 이 경우
$\lim_F A \to \lim_F A/\mathfrak a$가 전단사임은 곧바로 확인된다.

\medskip\noindent
조건 (4)를 확인하자. $X, X', Z, E$가 예
\ref{flat-example-two-generators}과 같다고 하자. 보조정리
\ref{flat-lemma-funny-blow-up}와 위와 같은 논증에 의해
$$
\mathcal{F}(X) = \lim_F \Gamma(X, \mathcal{O}_X)
\longrightarrow
\lim_F \Gamma(X', \mathcal{O}_{X'}) = \mathcal{F}(X')
$$
가 전단사임을 알 수 있다. 이 경우 $E = \mathbf{P}^1_Z$이므로
$\mathcal{F}(Z) \to \mathcal{F}(E)$도 전단사임을 알 수 있다. 따라서 이
경우에도 결론이 성립한다.
\end{proof}

\noindent
다음 보조정리에서는 스킴 $X$의 절대 약한 정규화 $X^{awn}$을 사용한다.
사상론, 절 \ref{morphisms-section-seminormalization}를 보라.

\begin{lemma}
\label{flat-lemma-weak-normalization-ph-sheaf}
$(\Sch/S)_{ph}$를 위상론, 정의
\ref{topologies-definition-big-small-ph}와 같은 사이트라 하자. 규칙
$$
X \longmapsto \Gamma(X^{awn}, \mathcal{O}_{X^{awn}})
$$
는 $(\Sch/S)_{ph}$ 위의 층이다.
\end{lemma}

\begin{proof}
$\mathcal{F}$가 층임을 증명하기 위해 위상론, 보조정리
\ref{topologies-lemma-characterize-sheaf}의 조건 (1)과 (2)를 확인하자.
$X^{awn}$의 형성은 열린 덮개와 가환하므로 조건 (1)이 성립한다. 사상론,
보조정리 \ref{morphisms-lemma-seminormalization}와 그 증명을 보라.

\medskip\noindent
$\pi : Y \to X$를 전사 고유 사상이라 하자. 두 사상
$$
\Gamma(Y^{awn}, \mathcal{O}_{Y^{awn}}) \to
\Gamma((Y \times_X Y)^{awn}, \mathcal{O}_{(Y \times_X Y)^{awn}})
$$
의 등화자가 $\Gamma(X^{awn}, \mathcal{O}_{X^{awn}})$와 같음을 보여야 한다.
$f$를 이 등화자의 원소라 하자. 이제 사상
$$
f : Y^{awn} \longrightarrow \mathbf{A}^1_X
$$
을 생각하자. $Y^{awn} \to X$는 보편 폐이므로 $f$의 스킴론적 상 $Z$는
$X$ 위에서 고유한 $\mathbf{A}^1_X$의 닫힌 부분스킴이고
$f : Y^{awn} \to Z$는 전사이다. 사상론, 보조정리
\ref{morphisms-lemma-scheme-theoretic-image-is-proper}를 보라. 따라서
$Z \to X$는 유한이고(사상론, 보조정리
\ref{morphisms-lemma-finite-proper}) 전사이다.

\medskip\noindent
$k$를 체라 하고 $z_1, z_2 : \Spec(k) \to Z$를 $Z \to X$에 의해
등화되는 두 사상이라 하자. $z_1 = z_2$임을 주장한다. 상
$\lambda_i = z_i^*f \in k$들이 일치함을 보이면 충분하다($Z$의 구조층은
$X$의 구조층 위에서 $f$로 생성되기 때문이다). 이를 보이기 위해 체 확대
$K/k$와 사상 $y_1, y_2 : \Spec(K) \to Y^{awn}$을
$z_i \circ (\Spec(K) \to \Spec(k)) = f \circ y_i$가 되도록 선택하자.
사상 $Y^{awn} \to Z$가 전사이므로 이는 가능하다. 매우 큰 기수의 대수적으로
닫힌 확대 $\Omega/k$를 선택하자. 임의의 $k$-대수 사상
$\sigma_i : K \to \Omega$에 대하여
$$
\Spec(\Omega)
\xrightarrow{\sigma_1, \sigma_2}
\Spec(K \otimes_k K)
\xrightarrow{y_1, y_2}
Y^{awn} \times_X Y^{awn}
$$
를 얻는다. 표준 사상
$(Y \times_X Y)^{awn} \to Y^{awn} \times_X Y^{awn}$이 보편
위상동형사상이고 $\Omega$가 대수적으로 닫혀 있으므로, 위 합성을 유일하게
사상 $\Spec(\Omega) \to (Y \times_X Y)^{awn}$으로 들어 올릴 수 있다.
$f$가 위의 등화자에 속하므로 이는
$\sigma_1(\lambda_1) = \sigma_2(\lambda_2)$임을 증명한다. 체 확대에 관한
쉬운 보조정리에 의해 이로부터 $\lambda_1 = \lambda_2$가 따른다. 세부사항은
생략한다.

\medskip\noindent
$Z \to X$가 보편 단사라고 결론 내린다. 즉 $Z \to X$는 점들 위에서
단사이고 순수 비분리 잉여체 확대들을 유도한다(사상론, 보조정리
\ref{morphisms-lemma-universally-injective}). 종합하면 $Z \to X$는 보편
위상동형사상이다. 사상론, 보조정리
\ref{morphisms-lemma-universal-homeomorphism}를 보라.
% P05-EN-CH38-0298

\medskip\noindent
$g : X^{awn} \to Z$를 $X^{awn}$의 보편성으로부터 얻은 사상이라 하자.
그러면 $Y^{awn} \to X^{awn} \to Z$와 $f : Y^{awn} \to Z$는 $X$ 위의 두
사상이다. $Y^{awn} \to Y$의 보편성에 의해 이에 대응하는 $Y$ 위의 두 사상
$Y^{awn} \to Y \times_X Z$는 같아야 한다. 이는 $\mathbf{A}^1_X$로 가는
사상으로서 $g \circ \pi^{wan} = f$임을 함의하며, 따라서
$g \in \Gamma(X^{awn}, \mathcal{O}_{X^{awn}})$가 우리가 찾던 원소라고
결론 내린다.
% P05-EN-CH38-0299
\end{proof}

\begin{lemma}
\label{flat-lemma-weak-normalization-h-sheaf}
$S$를 스킴이라 하자. 정의 \ref{flat-definition-big-small-h}와 같은 사이트
$(\Sch/S)_h$를 선택하자. 규칙
$$
X \longmapsto \Gamma(X^{awn}, \mathcal{O}_{X^{awn}})
$$
는 $(\Sch/S)_h$ 위의 ``구조층'' $\mathcal{O}$의 층화이다. ph 위상에
대해서도 마찬가지이다.
\end{lemma}

\begin{proof}
보조정리 \ref{flat-lemma-weak-normalization-ph-sheaf}에서 보조정리의 규칙
$\mathcal{F}$가 ph 위상의 층, 따라서 더욱이 h 위상의 층을 정의함을 보았다.
환의 전층들의 표준 사상 $\mathcal{O} \to \mathcal{F}$가 있음은 분명하다.
증명을 끝내려면 다음을 보이면 충분하다.
\begin{enumerate}
\item $f \in \mathcal{O}(X)$가 $\mathcal{F}(X)$에서 영으로 가면,
$f|_{X_i} = 0$이 되게 하는 h 덮개 $\{X_i \to X\}$가 존재한다.
\item $f \in \mathcal{F}(X)$가 주어지면, $f|_{X_i}$가 어떤
$f_i \in \mathcal{O}(X_i)$의 상이 되게 하는 h 덮개 $\{X_i \to X\}$가
존재한다.
\end{enumerate}
$f$가 (1)과 같다고 하자. 그러면 $f|_{X^{awn}} = 0$이다. 이는 $f$가
국소적으로 멱영임을 뜻한다. 따라서 $X' \subset X$가 $f$에 의해 잘려 나오는
닫힌 부분스킴이면 $X' \to X$는 유한표현인 전사 닫힌 몰입이다. 그러므로
$\{X' \to X\}$가 원하는 h 덮개이다. $f$가 (2)와 같다고 하자. $X$를 어떤
아핀 열린 덮개의 원소들로 치환하여 $X = \Spec(A)$가 아핀이라고 가정할 수
있다. 그러면 $f \in A^{awn}$이다. 사상론, 보조정리
\ref{morphisms-lemma-seminormalization-ring}을 보라. 사상론, 보조정리
\ref{morphisms-lemma-universal-homeo-limit}에 의해 유한표현 환 사상 $A \to B$를
$\Spec(B) \to \Spec(A)$가 보편 위상동형사상이고, $f$가 표준 사상
$B \to A^{awn}$ 아래에서 어떤 원소 $b \in B$의 상이 되도록 찾을 수 있다.
그러면 $\{\Spec(B) \to \Spec(A)\}$는 h 덮개이고 결론을 얻는다. ph 위상에
관한 명제도 같은 방식으로 따른다(또는 h 위상에 관한 명제로부터 유도할 수
있다).
\end{proof}

\noindent
$p$를 소수라 하자. $\mathbf{F}_p$-대수 $A$에서 사상
$F : A \to A$, $x \mapsto x^p$가 $A$의 자기동형이면 그 대수를
{\it 완전}이라고 한다.

\begin{lemma}
\label{flat-lemma-perfect-weankly-normal}
$p$를 소수라 하자. $\mathbf{F}_p$-대수 $A$가 절대 약정규일 필요충분조건은
완전인 것이다.
% P05-EN-CH38-0300
\end{lemma}

\begin{proof}
사상론, 정의 \ref{morphisms-definition-seminormal-ring}의 조건 (2)(b)로부터
$A$가 절대 약정규이면 완전이라는 사실이 곧바로 따른다.

\medskip\noindent
$A$가 완전이라고 가정하자. $x, y \in A$이고 $x^3 = y^2$라고 하자.
$p > 3$이면 어떤 $n, m > 0$에 대하여 $p = 2n + 3m$이라 쓸 수 있다.
$a^p = x$이고 $b^p = y$인 $a, b \in A$를 선택하자. $c = a^n b^m$으로
놓으면
$$
c^{2p} = x^{2n} y^{2m} = x^{2n + 3m} = x^p
$$
이고 따라서 $c^2 = x$이다. 마찬가지로 $c^3 = y$이다. $p = 2$이면
$x = a^2$이라 써서 $a^6 = y^2$을 얻으며, 이는 $a^3 = y$를 함의한다.
$p = 3$이면 $y = a^3$이라 써서 $x^3 = a^6$을 얻으며, 이는
$x = a^2$을 함의한다.

\medskip\noindent
$\ell$이 어떤 소수이고 $x, y \in A$가
$\ell^\ell x = y^\ell$을 만족한다고 하자. $\ell \not = p$이면
$a = y/\ell$은 $a^\ell = x$와 $\ell a = y$를 만족한다.
$\ell = p$이면 $y = 0$이고 어떤 $a$에 대하여 $x = a^p$이다.
\end{proof}

\begin{lemma}
\label{flat-lemma-char-p}
$p$를 소수라 하자.
\begin{enumerate}
\item $A$가 $\mathbf{F}_p$-대수이면 $\colim_F A = A^{awn}$이다.
\item $S$가 $\mathbf{F}_p$ 위의 스킴이면 $\mathcal{O}$의 h 층화는 준콤팩트
준분리 $X$를 $\colim_F \Gamma(X, \mathcal{O}_X)$로 보낸다.
\end{enumerate}
\end{lemma}

\begin{proof}
(1)의 증명. 대수학, 보조정리 \ref{algebra-lemma-p-ring-map}에 의해
$A \to \colim_F A$는 스펙트럼 위에서 보편 위상동형사상을 유도한다.
따라서 $B = \colim_F A$가 절대 약정규임을 보이면 충분하다. 사상론, 보조정리
\ref{morphisms-lemma-seminormalization-ring}을 보라. 환 사상 $F : B \to B$는
자기동형임에 유의하자. 달리 말해 $B$는 완전환이다. 그러므로 보조정리
\ref{flat-lemma-perfect-weankly-normal}을 적용할 수 있다.

\medskip\noindent
(2)의 증명. 아핀 국소적으로 살펴보면 이는 (1)과 보조정리
\ref{flat-lemma-h-sheaf-colim-F} 및
\ref{flat-lemma-weak-normalization-h-sheaf}에서 따른다.
\end{proof}






\section{양의 표수에서 벡터다발의 강하}
\label{flat-section-descent-vectorbundles-h}

\noindent
이 절의 참고문헌은 \cite{Witt-Grass}이다.

\medskip\noindent
스킴 $S$에 대하여 유한 국소 자유 $\mathcal{O}_S$-가군들의 범주를
$\textit{Vect}(S)$로 쓰자. $p$를 소수라 하고 $S$를 $\mathbf{F}_p$ 위의
준콤팩트 준분리 스킴이라 하자. 이 절에서는 범주
$$
\colim_F \textit{Vect}(S) =
\colim
\left(
\textit{Vect}(S) \xrightarrow{F^*}
\textit{Vect}(S) \xrightarrow{F^*}
\textit{Vect}(S) \xrightarrow{F^*} \ldots
\right)
$$
를 다룬다. 여기서 $F : S \to S$는 절대 Frobenius 사상이다. 구체적으로 이
범주의 대상은 쌍 $(\mathcal{E}, n)$이다. 여기서 $\mathcal{E}$는 유한 국소
자유 $\mathcal{O}_S$-가군이고 $n \geq 0$은 정수이다. 사상은 다음과 같이
정한다.
$$
\Hom_{\colim_F \textit{Vect}(S)}((\mathcal{E}, n), (\mathcal{G}, m)) =
\colim_N \Hom_S(F^{N - n, *}\mathcal{E}, F^{N - m, *}\mathcal{G})
$$
여기서 $F : S \to S$는 $S$의 절대 Frobenius 사상이다. 따라서 대상
$(\mathcal{E}, n)$은 대상 $(F^*\mathcal{E}, n + 1)$과 동형이다.

\begin{lemma}
\label{flat-lemma-colim-F-Vect}
$p$를 소수라 하고 $S$를 $\mathbf{F}_p$ 위의 준콤팩트 준분리 스킴이라 하자.
범주 $\colim_F \textit{Vect}(S)$는 $S$ 위의 환의 층
$\colim_F \mathcal{O}_S$ 위의 유한 국소 자유 가군들의 범주와 동치이다.
\end{lemma}

\begin{proof}
생략한다.
\end{proof}

\begin{lemma}
\label{flat-lemma-vector-bundle-I}
$p$를 소수라 하자. 예 \ref{flat-example-one-generator}과 같은 표수 $p$의 거의
블로업 사각형 $X, X', Z, E$를 생각하자. 그러면 함자
$$
\colim_F \textit{Vect}(X)
\longrightarrow
\colim_F \textit{Vect}(Z)
\times_{\colim_F \textit{Vect}(E)}
\colim_F \textit{Vect}(X')
$$
는 동치이다.
\end{lemma}

\begin{proof}
$A, f, J$가 예 \ref{flat-example-one-generator}과 같다고 하자. 우리의 모든
스킴이 아핀이고 벡터다발의 범주에 내부 Hom이 있으므로, 다음을 보이면 함자의
충만충실성이 따른다.
$$
\colim P \otimes_{A, F^N} A =
\colim P \otimes_{A, F^N} A/J
\times_{\colim P \otimes_{A, F^N} A/fA + J}
\colim P \otimes_{A, F^N} A/fA
$$
이는 유한 사영 $A$-가군 $P$에 대한 등식이다. $P$를 유한 자유 가군의
직접가산항으로 쓰면, 이는 $P$가 유한 자유인 경우에서 따른다. 이 경우는
곧바로 $P = A$인 경우로 환원된다. $P = A$인 경우는 보조정리
\ref{flat-lemma-h-sheaf-colim-F}에서 따른다(실제로 이 보조정리의 증명에서 이
경우를 직접 증명했다).

\medskip\noindent
본질적 전사성. 여기서는 다음 대수 문제를 얻는다. $P_1$이 유한 사영
$A/J$-가군이고 $P_2$가 유한 사영 $A/fA$-가군이며
$$
\varphi :
P_1 \otimes_{A/J} A/fA + J
\longrightarrow
P_2 \otimes_{A/fA} A/fA + J
$$
가 동형이라고 하자. 목표는 어떤 $N$, 유한 사영 $A$-가군 $P$, 동형
$\varphi_1 : P \otimes_A A/J \to P_1 \otimes_{A/J, F^N} A/J$, 그리고
$\varphi$와 자명한 방식으로 양립하는 동형
$\varphi_2 : P \otimes_A A/fA \to P_2 \otimes_{A/fA, F^N} A/fA$가
존재함을 보이는 것이다. 이는 다음과 같이 알 수 있다. 먼저
$$
A/(J \cap fA) = A/J \times_{A/fA + J} A/fA
$$
임을 관찰하자. 그러므로 대수학 보충, 보조정리
\ref{more-algebra-lemma-finitely-presented-module-over-fibre-product}에 의해
$A/(J \cap fA)$ 위의 유한 사영 가군 $P'$가 존재하며, 여기에는
$\varphi$와 양립하는 동형
$\varphi'_1 : P' \otimes_A A/J \to P_1$ 및
% P05-EN-CH38-0303
$\varphi_2 : P' \otimes_A A/fA \to P_2$가 주어진다. $J$가 유한생성
아이디얼이고 $f$-멱 꼬임이므로 $J \cap fA$가 멱영 아이디얼임을 알 수 있다.
따라서 어떤 $N$에 대하여 $F^N$의 분해
$$
A \xrightarrow{\alpha} A/(J \cap fA) \xrightarrow{\beta} A
$$
가 존재한다. $P = P' \otimes_{A/(J \cap fA), \beta} A$로 놓으면 결론을
얻는다.
\end{proof}

\begin{lemma}
\label{flat-lemma-vector-bundle-II}
$p$를 소수라 하자. 예 \ref{flat-example-two-generators}과 같은 표수 $p$의 거의
블로업 사각형 $X, X', Z, E$를 생각하자. 그러면 함자
$$
G :
\colim_F \textit{Vect}(X)
\longrightarrow
\colim_F \textit{Vect}(Z)
\times_{\colim_F \textit{Vect}(E)}
\colim_F \textit{Vect}(X')
$$
는 동치이다.
\end{lemma}

\begin{proof}
충만충실성. $(\mathcal{E}, n)$과 $(\mathcal{F}, m)$을
$\colim_F \textit{Vect}(X)$의 대상들이라 하자. 우변의 사상
$(a, b) : G(\mathcal{E}, n) \to G(\mathcal{F}, m)$을 택하자.
$N \gg 0$을 선택하여 $a$를 사상
$a : F^{N - n, *}\mathcal{E}|_Z \to F^{N - m, *}\mathcal{F}|_Z$로,
$b$를 $E$ 위에서 앞의 사상과 일치하는 사상
$b : F^{N - n, *}\mathcal{E}|_{X'} \to F^{N - m, *}\mathcal{F}|_{X'}$로
생각할 수 있다. $\mathcal{E}|_{X_i}$와 $\mathcal{F}|_{X_i}$가 유한 자유
$\mathcal{O}_{X_i}$-가군이 되게 하는 유한 아핀 열린 덮개
$X = X_1 \cup \ldots \cup X_n$을 선택하자. 각 $i$에 대하여 기저변환
$$
\xymatrix{
E_i \ar[r] \ar[d] & X'_i \ar[d] \\
Z_i \ar[r] & X_i
}
$$
은 예 \ref{flat-example-two-generators}과 같은 또 하나의 거의 블로업 사각형이다.
이 사각형들에 대하여 보조정리 \ref{flat-lemma-h-sheaf-colim-F}에 의해(그
보조정리의 증명을 보라)
$$
\colim_F H^0(X_i, \mathcal{O}_{X_i}) =
\colim_F H^0(Z_i, \mathcal{O}_{Z_i})
\times_{\colim_F H^0(E_i, \mathcal{O}_{E_i})}
\colim_F H^0(X'_i, \mathcal{O}_{X'_i})
$$
임을 안다. 따라서 $N$을 늘린 뒤 사상 $a|_{Z_i}$와 $b|_{X'_i}$가 사상
$c_i : F^{N - n, *}\mathcal{E}|_{X_i} \to F^{N - m, *}\mathcal{F}|_{X_i}$에서
온다고 가정할 수 있다. 필요하면 $N$을 다시 늘려 $c_i$와 $c_j$가
$X_i \cap X_j$에서 일치한다고 가정할 수 있다. 그러므로 이 사상들은 접합되어
좌변에서 원하는 사상 $(\mathcal{E}, n) \to (\mathcal{F}, m)$을 준다.

\medskip\noindent
본질적 전사성. $(\mathcal{F}, \mathcal{G}, \varphi)$를 유한 국소 자유
$\mathcal{O}_Z$-가군 $\mathcal{F}$, 유한 국소 자유 $\mathcal{O}_{X'}$-가군
$\mathcal{G}$, 그리고 동형 $\varphi : \mathcal{F}|_E \to \mathcal{G}|_E$로
이루어진 삼중항이라 하자. 이 삼중항을 어떤 Frobenius 거듭제곱 밑당김으로
치환한 뒤 유한 국소 자유 $\mathcal{O}_X$-가군에서 옴을 보여야 한다.

\medskip\noindent
뇌터 환원. 독자는 이 단락을 건너뛰기를 권한다. $X = \Spec(A)$,
$Z = \Spec(A/(f_1, f_2))$,
$X' = \text{Proj}(A[T_0, T_1]/(f_2T_0 - f_1T_1))$, 그리고
$E = \mathbf{P}^1_Z$임을 상기하자. 극한, 보조정리
\ref{limits-lemma-descend-invertible-modules}에 의해 $f_1$과 $f_2$를 포함하는
유한생성 $\mathbf{F}_p$-부분대수 $A_0 \subset A$를 삼중항
$(\mathcal{F}, \mathcal{G}, \varphi)$가
$X_0 = \Spec(A_0)$ 및 $Z_0 = \Spec(A_0/(f_1, f_2))$,
$X_0' = \text{Proj}(A_0[T_0, T_1]/(f_2T_0 - f_1T_1))$, 그리고
$E_0 = \mathbf{P}^1_{Z_0}$로 강하하도록 찾을 수 있다. 따라서 우리의
스킴들이 뇌터라고 가정할 수 있다.

\medskip\noindent
$X$가 뇌터라고 가정하자. $\mathcal{F}|_{Z \cap X_i}$가 자유가 되게 하는 유한
아핀 열린 덮개 $X = X_1 \cup \ldots \cup X_n$을 선택할 수 있다. 자리스키
위상에서 $\colim_F \textit{Vect}(X)$의 대상들을 접합할 수 있고(보조정리
\ref{flat-lemma-colim-F-Vect}), $X_i$와 $X_i \cap X_j$ 위의 충만충실성을 이미
알고 있으므로(증명의 첫 단락을 보라), 각 $X_i$ 위의 존재를 증명하면 충분하다.
이는 다음 단락에서 논하는 경우로 우리를 환원한다.

\medskip\noindent
$X$가 뇌터이고 $\mathcal{F} = \mathcal{O}_Z^{\oplus r}$이라고 가정하자.
$\varphi$를 사용하여 동형
$\mathcal{O}_E^{\oplus r} \to \mathcal{G}|_E$를 얻는다.
$I = (f_1, f_2) \subset A$라 하자. $\mathcal{I} \subset \mathcal{O}_{X'}$를
$E$의 아이디얼층이라 하자. 이는 $f_1$과 $f_2$에 의해 전역적으로 생성된다.
임의의 $n$에 대하여 전사
$$
(\mathcal{I}^n/\mathcal{I}^{n + 1})^{\oplus r} =
\mathcal{I}^n/\mathcal{I}^{n + 1} \otimes_{\mathcal{O}_E}
\mathcal{G}|_E \longrightarrow
\mathcal{I}^n\mathcal{G}/\mathcal{I}^{n + 1}\mathcal{G}
$$
가 존재한다. 그러므로 이 가군의 첫째 코호몰로지군은 영이다. 여기서는
$E = \mathbf{P}^1_Z$이므로 그 구조층, 그리고 실제로 전역적으로 생성되는 모든
준연접 가군의 $H^1$이 영이라는 사실을 사용한다. 사상론 보충, 보조정리
\ref{more-morphisms-lemma-globally-generated-vanishing}와 비교하라. 이제 짧은
완전열
$$
0 \to  \mathcal{I}^n\mathcal{G}/\mathcal{I}^{n + 1}\mathcal{G} \to
\mathcal{G}/\mathcal{I}^{n + 1}\mathcal{G} \to
\mathcal{G}/\mathcal{I}^n\mathcal{G} \to 0
$$
과 귀납법을 사용하면
$$
\lim H^0(X', \mathcal{G}/\mathcal{I}^n\mathcal{G})
\to
H^0(E, \mathcal{G}|_E) = H^0(E, \mathcal{O}_E^{\oplus r}) =
% P05-EN-CH38-0304
A/I^{\oplus r}
$$
가 전사임을 알 수 있다. 형식함수 정리(스킴의 코호몰로지, 정리
\ref{coherent-theorem-formal-functions})에 의해 이는
$$
H^0(X', \mathcal{G}) \to
H^0(E, \mathcal{G}|_E) = H^0(E, \mathcal{O}_E^{\oplus r}) =
A/I^{\oplus r}
$$
가 전사임을 함의한다. 따라서 주어진 $\mathcal{G}|_E$의 자명화와 양립하는 사상
$\alpha : \mathcal{O}_{X'}^{\oplus r} \to \mathcal{G}$를 선택할 수 있다.
그러면 $\alpha$는 $X'$ 안의 $E$의 어떤 열린근방 위에서 동형이다. 따라서
$Z$의 모든 점은 문제를 풀 수 있는 아핀 열린근방을 갖는다.
$X' \setminus E \to X \setminus Z$가 동형이므로 $Z$에 속하지 않는 $X$의
점들에 대해서도 마찬가지이다. 그러므로 또 하나의 자리스키 접합 논증으로
증명이 끝난다.
\end{proof}

\begin{proposition}
\label{flat-proposition-h-descent-vector-bundles-p}
$p$를 소수라 하고 $S$를 표수 $p$의 스킴이라 하자. 그러면 그루포이드로
섬유화된 범주
$$
p : \mathcal{S} \longrightarrow (\Sch/S)_h
$$
로서, $U$ 위의 섬유범주가 $U$ 위의 유한 국소 자유
$\colim_F \mathcal{O}_U$-가군들의 범주인 것은 그루포이드 스택이다. 더욱이
$U$가 준콤팩트하고 준분리이면 $\mathcal{S}_U$는
$\colim_F \textit{Vect}(U)$이다.
\end{proposition}

\begin{proof}
마지막 주장은 보조정리 \ref{flat-lemma-colim-F-Vect}의 내용이다. 명제를 증명하기
위해 보조정리 \ref{flat-lemma-refine-check-h-stack}의 조건 (1), (2), (3)을
확인하겠다.

\medskip\noindent
정의에 의해 자리스키 위상에 대한 접합이 있으므로 조건 (1)이 성립한다.

\medskip\noindent
조건 (2)를 보기 위해 $Y$가 아핀이고 $f : X \to Y$가 $S$ 위의 전사, 평탄,
고유, 유한표현 사상이라고 하자. $Y, X, X \times_Y X$가 준콤팩트하고
준분리이므로 명제의 명제문에 주어진 섬유범주의 기술을 사용할 수 있다. 그러면
$$
\textit{Vect}(Y)
\longrightarrow
\textit{Vect}(X)
\times_{\textit{Vect}(X \times_Y X)}
\textit{Vect}(X)
$$
가 동치임을 보이면 충분함이 분명하다(그러면 여극한에 대해서도 같은 결과가
따른다). 이는 유한 국소 자유 가군의 fppf 강하에서 곧바로 따른다. 강하, 명제
\ref{descent-proposition-fpqc-descent-quasi-coherent} 및 보조정리
\ref{descent-lemma-finite-locally-free-descends}를 보라.

\medskip\noindent
조건 (3)은 보조정리 \ref{flat-lemma-vector-bundle-I} 및
\ref{flat-lemma-vector-bundle-II}의 내용이다.
\end{proof}

\begin{lemma}
\label{flat-lemma-trivial-fibres-dvr}
$S$가 이산 값매김환의 스펙트럼이고 $f : X \to S$가 기하적으로 연결된
올들을 갖는 고유 사상이라 하자. 일반점을 $\eta \in S$로 쓰고, $S$의 어떤
균등화원의 $n$제곱에 의해 잘려 나오는 닫힌 부분스킴을 $X_n \subset X$로
쓰자. 그러면 다음이 성립하는 정수 $n$이 존재한다. 유한 국소 자유
$\mathcal{O}_X$-가군 $\mathcal{E}$가 $\mathcal{E}|_{X_\eta}$와
$\mathcal{E}|_{X_n}$이 자유라는 성질을 가지면, 그 가군은 자유이다.
\end{lemma}

\begin{proof}
먼저 $X \to S$가 단면을 갖는 경우로 환원하자. $S = \Spec(A)$라 쓰자.
$X_\eta$의 폐점 $\xi$를 선택하자. $B$의 분수체가 $\kappa(\xi)$가 되게 하는
이산 값매김환의 확대 $A \subset B$를 선택하자. Krull--Akizuki 정리(대수학,
보조정리 \ref{algebra-lemma-integral-closure-Dedekind})와 $\kappa(\xi)$가 $A$의
분수체의 유한 확대라는 사실에 의해 이는 가능하다. 고유성의 값매김 판정법
(사상론, 보조정리 \ref{morphisms-lemma-characterize-proper})에 의해 단면
$\sigma : \Spec(B) \to X_B$를 유도하는 $B$-값 점
$\tau : \Spec(B) \to X$를 얻는다. 유한 국소 자유 $\mathcal{O}_X$-가군
$\mathcal{E}$에 대하여 기저변환 $X_B$로의 밑당김을 $\mathcal{E}_B$라 하자.
평탄 기저변환(스킴의 코호몰로지, 보조정리
\ref{coherent-lemma-flat-base-change-cohomology})에 의해
$H^0(X_B, \mathcal{E}_B) = H^0(X, \mathcal{E}) \otimes_A B$임을 알 수 있다.
따라서 $\mathcal{E}_B$가 계수 $r$의 자유 가군이면 $H^0(X, \mathcal{E})$의
단면들은 자유 $B$-가군
$\tau^*\mathcal{E} = \sigma^*\mathcal{E}_B$를 생성한다. 특히
$\tau^*\mathcal{E}$를 생성하는 $\mathcal{E}$의 전역단면
$s_1, \ldots, s_r$을 $r$개 찾을 수 있다. 그러면
$$
s_1, \ldots, s_r :
\mathcal{O}_X^{\oplus r}
\longrightarrow
\mathcal{E}
$$
는 계수 $r$의 유한 국소 자유 $\mathcal{O}_X$-가군들의 사상이고, $X_B$로의
밑당김은 자유 $\mathcal{O}_{X_B}$-가군들의 사상으로서 각 올의 한 점에서
동형으로 제한된다. 행렬식을 취하면 각 올의 한 점에서 가역인 함수
% P05-EN-CH38-0306
$g \in \Gamma(X_\eta, \mathcal{O}_{X_B})$
를 얻는다. 올들이 고유이고 연결되어 있으므로 $g$는 가역이어야 한다(세부사항은
생략한다. 도움말: 대수다양체, 보조정리
\ref{varieties-lemma-proper-geometrically-reduced-global-sections}를 사용하라).
따라서 기저변환 $X_B \to \Spec(B)$에 대하여 보조정리를 증명하면 충분하다.

\medskip\noindent
단면 $\sigma : S \to X$가 있다고 가정하자. $\mathcal{E}$를 일반올과
$X_n$ 위에서 자유라고 가정한 유한 국소 자유 $\mathcal{O}_X$-가군이라 하자
($n$은 나중에 선택한다). 동형
$\sigma^*\mathcal{E} = \mathcal{O}_S^{\oplus r}$을 선택하자. $D(A)$의 사상
$$
K = R\Gamma(X, \mathcal{E}) \longrightarrow
R\Gamma(S, \sigma^*\mathcal{E}) = A^{\oplus r}
$$
을 생각하자. 위와 같이 논하면, 유도된 사상
$H^0(K) = H^0(X, \mathcal{E}) \to A^{\oplus r}$이 전사일 때 (그리고 그때에만)
$\mathcal{E}$가 자유임을 알 수 있다.

\medskip\noindent
$L = R\Gamma(X, \mathcal{O}_X^{\oplus r})$로 놓고, 이에 대응하는 사상
$L \to A^{\oplus r}$가 원하는 성질을 가짐을 관찰하자. 평탄 기저변환과
$\mathcal{E}$가 일반올에서 자유라는 가정에 의해
$K \otimes_A Q(A) \cong L \otimes_A Q(A)$임을 관찰하자.
$\pi \in A$를 균등화원이라 하자. 다음에 유의하라.
$$
K \otimes_A^\mathbf{L} A/\pi^m A =
R\Gamma(X, \mathcal{E} \xrightarrow{\pi^m} \mathcal{E})
$$
그리고 $L$에 대해서도 마찬가지이다. 특수올 위에 지지되는 단면들의 연접
부분층을 $\mathcal{E}_{tors} \subset \mathcal{E}$로 쓰고, 다른
$\mathcal{O}_X$-가군들에 대해서도 같은 표기를 쓰자. 다음 사상이 단사가
되게 하는 $k > 0$을 선택하자.
$(\mathcal{O}_X)_{tors} \to \mathcal{O}_X/\pi^k \mathcal{O}_X$
(스킴의 코호몰로지, 보조정리 \ref{coherent-lemma-Artin-Rees}).
$\mathcal{E}$가 국소 자유이므로
$\mathcal{E}_{tors} \subset \mathcal{E}/\pi^k\mathcal{E}$임을 알 수 있다.
그러면 $n \geq m + k$에 대하여 $D(\mathcal{O}_X)$에서 다음 동형들이 있다.
\begin{align*}
(\mathcal{E} \xrightarrow{\pi^m} \mathcal{E})
& \cong
(\mathcal{E}/\pi^k\mathcal{E} \xrightarrow{\pi^m}
\mathcal{E}/\pi^{k + m}\mathcal{E}) \\
& \cong
(\mathcal{O}_X^{\oplus r}/\pi^k\mathcal{O}_X^{\oplus r} \xrightarrow{\pi^m}
\mathcal{O}_X^{\oplus r}/\pi^{k + m}\mathcal{O}_X^{\oplus r}) \\
& \cong
(\mathcal{O}_X^{\oplus r} \xrightarrow{\pi^m} \mathcal{O}_X^{\oplus r})
\end{align*}
이는 $D(A)$에서 동형
$$
K \otimes_A^\mathbf{L} A/\pi^m A \cong L \otimes_A^\mathbf{L} A/\pi^m A
$$
을 결정한다(이는 $n \geq m + k$일 때 성립한다). 이 동형들은 $\sigma$에 의한
밑당김과 양립함에 유의하자. 따라서 특히
$K \otimes_A^\mathbf{L} A/\pi^m A \to (A/\pi^m A)^{\oplus r}$
가 차수 $0$의 코호몰로지 가군 위에서 전사를 정의한다고 결론 내린다($L$에
대하여 이것이 참이기 때문이다). $A$가 이산 값매김환이므로 $D(A)$에서
$$
K \cong \bigoplus H^i(K)[-i]
\quad\text{이고}\quad
L \cong
\bigoplus H^i(L)[-i]
$$
이다. 대수학 보충, 예
\ref{more-algebra-example-finite-injective-finite-global-dimension}를 보라.
스킴의 코호몰로지, 보조정리
\ref{coherent-lemma-proper-over-affine-cohomology-finite}에 의해 코호몰로지군
$H^i(K) = H^i(X, \mathcal{E})$와
$H^i(L) = H^i(X, \mathcal{O}_X)^{\oplus r}$는 유한 $A$-가군이다.
대수학 보충, 보조정리
\ref{more-algebra-lemma-generalized-valuation-ring-modules}에 의해 이 가군들은
순환 가군들의 직합이다. 위에서 $H^i(K)$와 $H^i(L)$의 자유 부분의 계수
$\beta_i$가 같음을 보았다. 다음을 관찰하자.
$$
H^i(L \otimes_A^\mathbf{L} A/\pi^m A) =
H^i(L)/\pi^m H^i(L) \oplus H^{i + 1}(L)[\pi^m]
$$
그리고 $K$에 대해서도 마찬가지이다. 어떤 $i$에 대하여 $A/\pi^eA$가
$H^i(X, \mathcal{O}_X)$, 또는 동치로 $H^i(L)$의 직접가산항으로 나타나는
가장 큰 정수를 $e$라 하자. 그러면 $m = e + 1$로 취할 때
$H^i(L \otimes_A^\mathbf{L} A/\pi^m A)$는 $A/\pi^m A$의 $\beta_i$개 복사본과,
각각 $\pi^e$에 의해 소멸되는 몇몇 다른 순환 가군들의 직합임을 알 수 있다.
$K$에 대한 같은 논증과 동형
$K \otimes_A^\mathbf{L} A/\pi^m A \cong L \otimes_A^\mathbf{L} A/\pi^m A$
에 의해 $H^i(K)$에는 $l > e$인 $A/\pi^l A$ 꼴의 순환 직접가산항이 있을 수
없음이 따른다. ($K$가 $D(A)$의 대상으로서 $L$과 동형이라는 것도 따르지만,
이는 필요하지 않다.) 이제 사상
$$
H^0(K \otimes^\mathbf{L}_A A/\pi^{e + 1} A) =
H^0(K)/\pi^{e + 1}H^0(K) \oplus H^1(K)[\pi^{e + 1}]
\longrightarrow
(A/\pi^{e + 1} A)^{\oplus r}
$$
가 전사일 수 있는 유일한 경우는 첫째 직접가산항 위에서 전사인 경우이다.
이것이 우리가 보이려던 것이다. (정확히 말해 단면 $\sigma$가 있다면,
보조정리의 명제문에 있는 정수 $n$은 $k + e + 1$이어야 한다. 여기서 $k$와
$e$는 위와 같고 $X$에만 의존한다.)
\end{proof}

\begin{lemma}
\label{flat-lemma-trivial-over-dvrs}
$f : X \to S$를 스킴의 사상이라 하고 $\mathcal{E}$를 유한 국소 자유
$\mathcal{O}_X$-가군이라 하자. 다음을 가정하자.
\begin{enumerate}
\item $f$는 평탄이고 고유이며 $\mathcal{O}_S = f_*\mathcal{O}_X$이다.
\item $S$는 정규 뇌터 스킴이다.
\item 모든 여차원 $1$인 점 $s \in S$에 대하여 $\mathcal{E}$를
$X \times_S \Spec(\mathcal{O}_{S, s})$로 당긴 것은 자유이다.
\end{enumerate}
그러면 $\mathcal{E}$는 어떤 유한 국소 자유 $\mathcal{O}_S$-가군의 밑당김과
동형이다.
\end{lemma}

\begin{proof}
표준 사상
$$
\Phi : f^*f_*\mathcal{E} \longrightarrow \mathcal{E}
$$
이 동형임을 증명하겠다. 평탄 기저변환(스킴의 코호몰로지, 보조정리
\ref{coherent-lemma-flat-base-change-cohomology})과 가정 (1), (3)에 의해,
모든 여차원 $1$인 점 $s \in S$에 대하여 이를
$X \times_S \Spec(\mathcal{O}_{S, s})$로 당긴 것이 동형임을 알 수 있다.
인자, 보조정리 \ref{divisors-lemma-check-isomorphism-via-depth-and-ass}에 의해
여차원 $\geq 2$인 점 $s \in S$로 가는 임의의 점 $x \in X$에 대하여
$\text{depth}((f^*f_*\mathcal{E})_x) \geq 2$임을 증명하면 충분하다. $f$가
평탄이고
$(f^*f_*\mathcal{E})_x = (f_*\mathcal{E})_s \otimes_{\mathcal{O}_{S, s}}
\mathcal{O}_{X, x}$이므로 $\text{depth}((f_*\mathcal{E})_s) \geq 2$임을
증명하면 충분하다. 대수학, 보조정리
\ref{algebra-lemma-apply-grothendieck}을 보라. $S$가 정규 뇌터 스킴이고
$\dim(\mathcal{O}_{S, s}) \geq 2$이므로
$\text{depth}(\mathcal{O}_{S, s}) \geq 2$이다. 성질론, 보조정리
\ref{properties-lemma-criterion-normal}을 보라. 따라서 인자, 보조정리
\ref{divisors-lemma-depth-pushforward}로부터 원하는 결과를 얻는다.
\end{proof}

\noindent
위 결과들을 사용하여 다음 놀라운 명제를 증명할 수 있다.

\begin{theorem}
\label{flat-theorem-pullback-trivial-fibres}
$p$를 소수라 하고 $Y$를 $\mathbf{F}_p$ 위의 준콤팩트 준분리 스킴이라 하자.
$f : X \to Y$를 기하적으로 연결된 올들을 갖는 고유 전사 유한표현 사상이라
하자. 그러면 함자
$$
\colim_F \textit{Vect}(Y) \longrightarrow \colim_F \textit{Vect}(X)
$$
는 충만충실하며 그 본질적 상은 다음과 같이 기술된다. $\mathcal{E}$를 유한
국소 자유 $\mathcal{O}_X$-가군이라 하자. 모든 $y \in Y$에 대하여 다음을
만족하는 정수 $n_y, r_y \geq 0$이 존재한다고 가정하자.
$$
F^{n_y, *}\mathcal{E}|_{X_{y, red}}
\cong
\mathcal{O}_{X_{y, red}}^{\oplus r_y}
$$
그러면 어떤 $n \geq 0$에 대하여 $n$번째 Frobenius 거듭제곱 밑당김
$F^{n, *}\mathcal{E}$는 어떤 유한 국소 자유 $\mathcal{O}_Y$-가군의
밑당김이다.
\end{theorem}

\begin{proof}
충만충실성의 증명. $Y$ 위의 벡터다발들은 국소적으로 자명하므로 이는 사상
$$
\colim_F \Gamma(Y, \mathcal{O}_Y)
\longrightarrow
\colim_F \Gamma(X, \mathcal{O}_X)
$$
이 전단사라는 명제로 환원된다. $\{X \to Y\}$는 h 덮개이므로, 두 사상
$$
\colim_F \Gamma(X, \mathcal{O}_X)
\longrightarrow
\colim_F \Gamma(X \times_Y X, \mathcal{O}_{X \times_Y X})
$$
이 같음을 보이면 보조정리 \ref{flat-lemma-h-sheaf-colim-F}로부터 원하는 결론이
따른다. $g \in \Gamma(X, \mathcal{O}_X)$라 하고, $g$를
$X \times_Y X$로 당기는 두 밑당김을 $g_1$과 $g_2$로 쓰자.
$X_{y, red}$는 기하적으로 연결되어 있으므로
$H^0(X_{y, red}, \mathcal{O}_{X_{y, red}})$가 $\kappa(y)$의 순수 비분리
확대임을 알 수 있다. 대수다양체, 보조정리
\ref{varieties-lemma-proper-geometrically-reduced-global-sections}를 보라.
따라서 어떤 $p$-거듭제곱 $q$에 대하여(이는 $y$에 의존할 수 있다)
$g^q|_{X_{y, red}}$는 $\kappa(y)$의 원소에서 온다. 그러므로
$g_1^q$와 $g_2^q$는 $(X \times_Y X)_y = X_y \times_y X_y$의 모든 점에서
잉여체의 같은 원소로 간다. 따라서 $g_1 - g_2$는
$(X \times_Y X)_{red}$에서 영으로 제한된다. 그러므로 원하는 대로
$p$-거듭제곱으로 택할 수 있는 어떤 $n$에 대하여 $(g_1 - g_2)^n = 0$이다.

\medskip\noindent
본질적 상의 기술. $\mathcal{E}$가 정리의 명제문과 같다고 하자. 먼저 뇌터인
경우로 환원하자.

\medskip\noindent
$y \in Y$를 점이라 하고 이를 잉여체의 스펙트럼에서 $Y$로 가는 사상
$y \to Y$로 보자. $y \to Y$를 $Y_i$가 아핀인 유한표현 사상
$Y_i \to Y$들의 여과 극한으로 쓸 수 있다. (이를 직접 증명하는 것이 가장
좋지만, 극한, 보조정리 \ref{limits-lemma-relative-approximation} 및
\ref{limits-lemma-limit-affine}에서도 형식적으로 따른다.) 각 $i$에 대하여
$Z_i = Y_i \times_Y X$로 놓자. 그러면 $X_y = \lim Z_i$이고
$X_{y, red} = \lim Z_{i, red}$이다. 극한, 보조정리
\ref{limits-lemma-descend-modules-finite-presentation}에 의해
$F^{n_y, *}\mathcal{E}|_{Z_{i, red}} \cong
\mathcal{O}_{Z_{i, red}}^{\oplus r_y}$인 $i$를 찾을 수 있다.
이 $i$를 고정하자. 여기서 $Z_{i, j} \to Z_i$는 유한표현 비후화인
$Z_{i, red} = \lim Z_{i, j}$를 얻는다(극한, 보조정리
\ref{limits-lemma-closed-is-limit-closed-and-finite-presentation}). 앞과 같은
보조정리를 사용하여
$F^{n_y, *}\mathcal{E}|_{Z_{i, j}} \cong \mathcal{O}_{Z_{i, j}}^{\oplus r_y}$.
인 $j$를 찾을 수 있다. 각 $y \in Y$에 대하여 상이 $y$를 포함하는 유한표현
사상 $Y_y \to Y$와 다음을 만족하는 비후화
$Z_y \to Y_y \times_Y X$가 존재한다고 결론 내린다.
$F^{n_y, *}\mathcal{E}|_{Z_y} \cong \mathcal{O}_{Z_y}^{\oplus r_y}$.
$Y_y \to Y$의 상은 구성가능하다(사상론, 보조정리
\ref{morphisms-lemma-chevalley}). $Y$는 구성가능 위상에서 준콤팩트이므로(위상,
보조정리 \ref{topology-lemma-constructible-hausdorff-quasi-compact} 및 성질론,
보조정리 \ref{properties-lemma-quasi-compact-quasi-separated-spectral}), 유한표현
사상들이 유한 개
$$
Y_1 \to Y,\ Y_2 \to Y,\ \ldots,\ Y_N \to Y
$$
존재하여 집합론적으로 $Y = \bigcup \Im(Y_a \to Y)$이고, 각
$a \in \{1, \ldots, N\}$에 대하여 정수 $n_a, r_a \geq 0$과 유한표현 비후화
$Z_a \subset Y_a \times_Y X$가 존재하며
$F^{n_a, *}\mathcal{E}|_{Z_a} \cong \mathcal{O}_{Z_a}^{\oplus r_a}$.

\medskip\noindent
이와 같이 정식화하면 조건은 절대 뇌터 근사로 강하한다. 독자는 이 단락을
건너뛰기를 강력히 권한다. 먼저 $Y = \lim_{i \in I} Y_i$를 아핀 전이 사상을
갖는 $\mathbf{F}_p$ 위 유한형 스킴들의 공여과 극한으로 쓰자(극한, 보조정리
\ref{limits-lemma-relative-approximation}). 다음으로 $Y$로의 기저변환이
$f : X \to Y$를 복원하는 고유 사상 $f_i : X_i \to Y_i$들이 있다고 가정할
수 있다. 극한, 보조정리 \ref{limits-lemma-descend-finite-presentation}을 보라.
$i$를 늘린 뒤, $X$로 당긴 것이 $\mathcal{E}$와 동형인 유한 국소 자유
$\mathcal{O}_{X_i}$-가군 $\mathcal{E}_i$가 존재한다고 가정할 수 있다.
극한, 보조정리 \ref{limits-lemma-descend-invertible-modules}를 보라.
$0 \in I$를 선택하고 $f_0$의 기하학적 올들이 연결되는 구성가능 부분집합을
$E \subset Y_0$로 쓰자. 사상론 보충, 보조정리
\ref{more-morphisms-lemma-nr-geom-connected-components-constructible}를 보라.
그러면 $Y \to Y_0$는 $E$로 간다. 사상론 보충, 보조정리
\ref{more-morphisms-lemma-base-change-fibres-geometrically-connected}를 보라.
따라서 $i \gg 0$에 대하여 $Y_i \to Y_0$는 $E$로 간다. 극한, 보조정리
\ref{limits-lemma-limit-contained-in-constructible}를 보라. 그러므로
$i \gg 0$에 대하여 $f_i$의 올들이 기하적으로 연결됨을 알 수 있다. 극한,
보조정리 \ref{limits-lemma-descend-finite-presentation}에 의해 충분히 큰 $i$에
대하여 $Y$로의 기저변환이 $Y_a \to Y$를 복원하는 유한형 사상
$Y_{i, a} \to Y_i$, $a  \in \{1, \ldots, N\}$을 찾을 수 있다. 필요하면
$i$를 늘려, $Y_a \times_Y X$로의 기저변환이 $Z_a$를 복원하는 비후화
$Z_{i, a} \subset Y_{i, a} \times_{Y_i} X_i$를 찾을 수 있다(앞과 같은
참고문헌을 극한, 보조정리
\ref{limits-lemma-descend-closed-immersion-finite-presentation} 및
\ref{limits-lemma-descend-surjective}와 함께 사용하라). $Z_a = \lim Z_{i, a}$이므로
$i$를 늘린 뒤
$F^{n_a, *}\mathcal{E}_i|_{Z_{i, a}} \cong
\mathcal{O}_{Z_{i, a}}^{\oplus r_a}$라고 가정할 수 있다. 극한, 보조정리
\ref{limits-lemma-descend-modules-finite-presentation}을 보라. 끝으로 $i$를 한 번
더 늘려 극한, 보조정리 \ref{limits-lemma-descend-surjective}에 의해
$\coprod Y_{i, a} \to Y_i$가 전사라고 가정할 수 있다. 이제 모든 가정이
$X_i \to Y_i$와 $\mathcal{E}_i$에 대하여 성립하며, $X_i \to Y_i$와
$\mathcal{E}_i$에 대하여 결과를 증명하면 충분함을 알 수 있다.

\medskip\noindent
$Y$가 $\mathbf{F}_p$ 위 유한형이라고 가정하자. 결과를 증명하기 위해
$\dim(Y)$에 관한 귀납법을 사용하겠다. $\mathcal{E}$가 결정하는
$\colim_F \textit{Vect}(X)$의 대상으로 당겨지는
$\colim_F \textit{Vect}(Y)$의 대상을 찾으려 한다. 이미 증명한 충만충실성과
명제 \ref{flat-proposition-h-descent-vector-bundles-p}에 의해, $Y$를 어떤 h 덮개의
원소들로 치환하고 $X$를 이에 대응하는 기저변환으로 치환한 뒤
$\mathcal{E}$의 강하를 구성하면 충분하다. 이는 $Y$의 차원을 늘리지 않는 한
$Y$를 $Y$ 위 고유 전사인 스킴으로 치환할 수 있다는 뜻이다.
$T \subset T'$이 $\mathbf{F}_p$ 위 유한형 스킴들의 비후화이면
$\colim_F \textit{Vect}(T) = \colim_F \textit{Vect}(T')$
이다. 실제로 $\{T \to T'\}$는 $T \times_{T'} T = T$를 만족하는 h 덮개이다.
$T' \to T$가 $\mathbf{F}_p$ 위 유한형 스킴들의 보편 위상동형사상이면
$\colim_F \textit{Vect}(T) = \colim_F \textit{Vect}(T')$
이다. 실제로 $\{T \to T'\}$는 대각사상
$T \subset T \times_{T'} T$이 비후화인 h 덮개이다.
% P05-EN-CH38-0311

\medskip\noindent
위의 일반적인 논평을 사용하여 $X$를 그 축소로 치환하며, $X$가 축소되었다고
가정할 수 있다. Stein 분해 $X \to Y' \to Y$를 생각하자. 사상론 보충, 정리
\ref{more-morphisms-theorem-stein-factorization-Noetherian}를 보라. 그러면
$Y' \to Y$는 $\mathbf{F}_p$ 위 유한형 스킴들의 보편 위상동형사상이다.
위 결과에 의해 $Y$를 $Y'$로 치환할 수 있다. 따라서
$f_*\mathcal{O}_X = \mathcal{O}_Y$이고 $Y$가 축소되었다고 가정할 수 있다.
이는 다음 단락에서 논하는 경우로 우리를 환원한다.

\medskip\noindent
$Y$가 축소되었고 $Y$의 어떤 조밀한 열린 부분스킴 위에서
$f_*\mathcal{O}_X = \mathcal{O}_Y$라고 가정하자. 그러면 어떤 조밀한 열린
부분스킴 $V \subset Y$ 위에서 $X \to Y$는 평탄하다. 사상론, 명제
\ref{morphisms-proposition-generic-flatness-reduced}를 보라. 보조정리
\ref{flat-lemma-flat-after-blowing-up}에 의해 $X$의 강변환 $X'$가 $Y'$ 위에서
평탄해지게 하는 $V$-허용 블로업 $Y' \to Y$가 존재한다. $Y$와 $Y'$가 공통의
조밀한 열린 부분스킴을 가지므로 $\dim(Y') = \dim(Y)$임을 관찰하자. 사상론
보충, 보조정리
\ref{more-morphisms-lemma-proper-flat-nr-geom-conn-comps-lower-semicontinuous}와
$V \subset Y'$가 조밀하다는 사실에 의해 $f' : X' \to Y'$의 모든 올은
기하적으로 연결되어 있다. 여전히
$(f'_*\mathcal{O}_{X'})|_V = \mathcal{O}_V$이다. 다음과 같이 쓰자.
$$
Y' \times_Y X = X' \cup E \times_Y X
$$
여기서 $E \subset Y'$는 블로업의 예외 인자이다. 위의 일반적인 논평에 의해
$Y' \times_Y X \to Y'$와 $\mathcal{E}$의 $Y' \times_Y X$로의 제한에 대하여
존재를 증명하면 충분하다. 그 대상을 당기면 $\mathcal{E}$의 $X'$로의 제한이
되는(여극한 범주의 대상으로 본다) 어떤 대상 $\xi'$을
$\colim_F \textit{Vect}(Y')$에서 찾았다고 하자. $\dim(Y)$에 관한
귀납법으로 $E \times_Y X$로 당기면 $\mathcal{E}$의 그곳으로의 제한이 되는
대상 $\xi''$을 $\colim_F \textit{Vect}(E)$에서 찾을 수 있다. 그러면
충만충실성에 의해 $\mathcal{E}$의 $E \times_{Y'} X'$로의 제한과 주어진
동일시들과 양립하는 유일한 동형 $\xi'|_E \to \xi''$이 결정된다. 이제
$$
\{E \times_Y X \to Y' \times_Y X, X' \to Y' \times_Y X\}
$$
는 다음을 만족하는 닫힌 몰입 한 쌍으로 주어지는 h 덮개이다.
$$
(E \times_Y X) \times_{(Y' \times_Y X)} X' = E \times_{Y'} X'
$$
따라서 $\xi'$은 $Y' \times_Y X$로 당겼을 때 $\mathcal{E}$의 그곳으로의 제한이
된다고 결론 내린다. 그러므로 $\xi'$을 찾으면 충분하며, 다음 단락에서 논하는
경우로 환원된다.

\medskip\noindent
$Y$가 축소되었고 $f$가 평탄이며 $Y$의 어떤 조밀한 열린 부분스킴 위에서
$f_*\mathcal{O}_X = \mathcal{O}_Y$라고 가정하자. 이 경우 정규화
$Y^\nu \to Y$를 생각한다(사상론, 절
\ref{morphisms-section-normalization}). 이는 조밀한 열린집합 위에서 동형인
유한 전사 사상이다(사상론, 보조정리
\ref{morphisms-lemma-nagata-normalization} 및
\ref{morphisms-lemma-ubiquity-nagata}). 따라서 우리의 일반적인 논평에 의해
$Y$를 $Y^\nu$로, $X$를 $Y^\nu \times_Y X$로 치환할 수 있다. 이 치환 뒤에는
$\mathcal{O}_Y = f_*\mathcal{O}_X$임을 알 수 있다(이 경우 Stein 분해가
동형이어야 하기 때문이다. 작은 세부사항은 생략한다).

\medskip\noindent
$Y$가 정규 뇌터 스킴이고 $f$가 평탄이며
$f_*\mathcal{O}_X = \mathcal{O}_Y$라고 가정하자. $\mathcal{E}$를 적당한
Frobenius 거듭제곱 밑당김으로 치환한 뒤, $Y$의 기약 성분들의 일반점에서
$f$의 스킴론적 올 위의 $\mathcal{E}$가 자명하다고 가정할 수 있다
($\colim_F \textit{Vect}(-)$가 보편 위상동형사상 위에서 동치이기 때문이다.
위를 보라). 위의 논증과 마찬가지로(뇌터인 경우로 환원하는 논증에서),
$\mathcal{E}|_{f^{-1}(V)}$가 자유가 되게 하는 조밀한 열린 부분스킴
$V \subset Y$가 존재한다고 결론 내린다. 집합론적으로 $Y = V \amalg Z$가
되게 하는 닫힌 부분스킴 $Z \subset Y$를 택하자.
$z_1, \ldots, z_t \in Z$를 $Z$의 여차원 $1$인 기약 성분들의 일반점들이라
하자. 그러면 $A_i = \mathcal{O}_{Y, z_i}$는 이산 값매김환이다.
$n_i$를 $A_i$ 위의 스킴 $X_{A_i}$에 대하여 보조정리
\ref{flat-lemma-trivial-fibres-dvr}에서 얻은 정수라 하자. $\mathcal{E}$를 적당한
Frobenius 거듭제곱 밑당김으로 치환한 뒤,
$X_{A_i/\mathfrak m_i^{n_i}}$ 위에서 $\mathcal{E}$가 자유라고 가정할 수 있다
(다시 말해 여극한 범주가 보편 위상동형사상 아래 불변이기 때문이다. 위를
보라). 그러면 보조정리 \ref{flat-lemma-trivial-fibres-dvr}에 의해
$\mathcal{E}$는 $X_{A_i}$ 위에서 자유이다. 따라서 끝으로 보조정리
\ref{flat-lemma-trivial-over-dvrs}를 적용하여 결론을 얻는다.
\end{proof}










\section{복합체의 블로업}
\label{flat-section-blowup-complexes}

\noindent
이 절에서는 블로업한 뒤 유도범주의 완전 대상들에 대한 표준형을 찾는다.

\begin{lemma}
\label{flat-lemma-fitting-ideals-complex}
$X$를 스킴이라 하고 $E \in D(\mathcal{O}_X)$를 유사연접이라 하자. 모든
$p, k \in \mathbf{Z}$에 대하여 다음 성질을 갖는 유한형 준연접 아이디얼층
$\text{Fit}_{p, k}(E) \subset \mathcal{O}_X$가 존재한다. $U \subset X$가
열린집합이고 $E|_U$가 다음과 동형이면
$$
\ldots \to
\mathcal{O}_U^{\oplus n_{b - 2}}
\xrightarrow{d_{b - 2}}
\mathcal{O}_U^{\oplus n_{b - 1}}
\xrightarrow{d_{b - 1}}
\mathcal{O}_U^{\oplus n_b} \to 0 \to \ldots
$$
제한 $\text{Fit}_{p, k}(E)|_U$는 $d_p$의 행렬의 다음 크기 소행렬식들로
생성된다.
$$
- k + n_{p + 1} - n_{p + 2} + \ldots + (-1)^{b - p + 1} n_b
$$
규약: $r \leq 0$이면 $r \times r$ 소행렬식들로 생성되는 아이디얼은
$\mathcal{O}_U$이고, $r > \min(n_p, n_{p + 1})$이면 $r \times r$
소행렬식들로 생성되는 아이디얼은 영이다.
\end{lemma}

\begin{proof}
$X$ 위 국소적으로 $E$는 보조정리의 명제문과 같은 꼴을 갖는다. 대수학 보충,
절 \ref{more-algebra-section-pseudo-coherent}, 코호몰로지, 절
\ref{cohomology-section-pseudo-coherent}, 그리고 스킴의 유도범주, 절
\ref{perfect-section-spell-out}를 보라. 따라서 소행렬식들의 아이디얼이 선택한
표현자와 무관함을 증명하면 충분하다. 이를 위해 국소환에서 확인하면 충분하다.
국소환 $(R, \mathfrak m, \kappa)$ 위에서 위로 유계인 복합체
$$
F^\bullet :
\ldots \to
R^{\oplus n_{b - 2}}
\xrightarrow{d_{b - 2}}
R^{\oplus n_{b - 1}}
\xrightarrow{d_{b - 1}}
R^{\oplus n_b} \to 0 \to \ldots
$$
를 생각하자. $d_p$의 행렬에서 크기가
$k - n_{p + 1} + n_{p + 2} - \ldots + (-1)^{b - p} n_b$인 소행렬식들로
생성되는 아이디얼을
% P05-EN-CH38-0312
$\text{Fit}_{k, p}(F^\bullet) \subset R$로 쓰자. $F^\bullet$의 어떤 미분의
어떤 행렬 성분이 가역이라고 하자. $d_i$가 가역 행렬 성분을 갖게 하는 가장
큰 정수 $i$를 선택하자. 대수학, 보조정리
\ref{algebra-lemma-add-trivial-complex}에 의해 복합체 $F^\bullet$은 차수 $i$와
$i + 1$에서 영이 아닌 항들을 갖는 자명한 복합체
$\ldots \to 0 \to R \to R \to 0 \to \ldots$와 복합체 $(F')^\bullet$의
직합과 동형이다. 다음의 확인은 독자에게 맡긴다.
$\text{Fit}_{p, k}(F^\bullet) = \text{Fit}_{p, k}((F')^\bullet)$;
여기서 소행렬식들의 크기에 대한 공식을 사용한다. $(F')^\bullet$의 또 다른
미분이 가역 행렬 성분을 가지면 이를 다시 반복하는 식으로 진행한다. 이와 같이
계속하면 마침내 모든 미분의 행렬 성분이 $\mathfrak m$에 속하는 복합체
$(F^\infty)^\bullet$에 도달한다. 여기서는 무한히 많은 단계를 수행해야 할 수
있지만, 임의의 절단점에 대하여 이 단계들 중 유한 개만이 차수가 $\geq $
절단점인 복합체에 영향을 준다. 따라서 ``극한'' $(F^\infty)^\bullet$은 잘 정의된
유한 자유 가군들의 위로 유계인 복합체이고, 우리가 시작한 복합체로 가는
준동형
$(F^\infty)^\bullet \to F^\bullet$을 갖추며,
$\text{Fit}_{p, k}(F^\bullet) = \text{Fit}_{p, k}((F^\infty)^\bullet)$.
대수학 보충, 보조정리
\ref{more-algebra-lemma-lift-pseudo-coherent-from-residue-field}에 의해 복합체
$(F^\infty)^\bullet$은 동형까지 유일하므로 증명이 끝난다.
\end{proof}

\begin{lemma}
\label{flat-lemma-blowup-complex}
$X$를 스킴이라 하고 $E \in D(\mathcal{O}_X)$를 완전이라고 하자. 모든
$i \in \mathbf{Z}$에 대하여 $H^i(E|_U)$가 일정한 계수 $r_i$의 유한 국소
자유가 되게 하는 스킴론적으로 조밀한 열린 부분스킴 $U \subset X$를 택하자.
그러면 모든 $i \in \mathbf{Z}$에 대하여 $H^i(Lb^*E)$가 Tor 차원
$\leq 1$인 완전 $\mathcal{O}_{X'}$-가군이 되게 하는 $U$-허용 블로업
$b : X' \to X$가 존재한다.
\end{lemma}

\begin{proof}
블로업을 아핀 국소적으로 구성하고 연구하겠다. 즉 $V \subset X$가 아핀 열린
부분스킴이고 $E|_V$가 복합체
$$
\mathcal{O}_V^{\oplus n_a} \xrightarrow{d_a}
\ldots \xrightarrow{d_{b - 1}} \mathcal{O}_V^{\oplus n_b}
$$
로 표현된다고 하자.
$k_ i = r_{i + 1} - r_{i + 2} + \ldots + (-1)^{b - i + 1}r_b$로 놓자.
생략하는 계산에 의해 $U \cap V$ 위에서 $d_i$의 계수는
$$
\rho_i = - k_i + n_{i + 1} - n_{i + 2} + \ldots + (-1)^{b - i + 1}n_b
$$
이다. 여기서 그 의미는 $d_i$의 여핵이 계수 $n_{i + 1} - \rho_i$의 유한
국소 자유라는 것이다. $\mathcal{I}_i \subset \mathcal{O}_V$를 $d_i$의
행렬에서 크기 $\rho_i \times \rho_i$인 소행렬식들로 생성되는 아이디얼이라
하자.

\medskip\noindent
한편 보조정리 \ref{flat-lemma-fitting-ideals-complex}와 비교하면 아이디얼
$\mathcal{I}_i$는 $E|_V$를 표현하는 복합체의 선택과 무관함을 보인 전역
아이디얼 $\text{Fit}_{i, k_i}(E)$에 대응함을 알 수 있다. 다른 한편
$\mathcal{I}_i$는 $\Coker(d_i)$의 $(n_{i + 1} - \rho_i)$번째 피팅
아이디얼이다. 이를 기억해 두라.

\medskip\noindent
$b : X' \to X$를 아이디얼들 $\text{Fit}_{i, k_i}(E)$의 곱에서의 블로업이라
하자. $X$ 위 국소적으로 이 아이디얼들 중 거의 모두가 단위 아이디얼과 같으므로
이는 의미가 있다(위를 보라). 이 블로업은 아이디얼
$\text{Fit}_{i, k_i}(E)$에서의 블로업 $b_i : X'_i \to X$들을 우세화한다.
인자, 보조정리 \ref{divisors-lemma-blowing-up-two-ideals}를 보라. 인자,
보조정리 \ref{divisors-lemma-blowup-fitting-ideal}에 의해 각 $b_i$는
$U$-허용 블로업이다. 따라서 $b$도 $U$-허용 블로업이다(아주 작은 세부사항은
생략한다. 인자, 보조정리
\ref{divisors-lemma-dominate-admissible-blowups}의 증명과 비교하라). 끝으로
$U$는 여전히 $X'$의 스킴론적으로 조밀한 열린 부분스킴이다. 따라서 $X$를
$X'$로 치환하면 다음 단락에서 논하는 상황에 이른다.

\medskip\noindent
모든 $i$에 대하여 $\text{Fit}_{i, k_i}(E)$가 가역 아이디얼이라고 가정하자.
위와 같이 $E|_V$를 표현하는 유한 자유 가군들의 복합체와 아핀 열린집합 $V$를
선택하자. 인자, 보조정리 \ref{divisors-lemma-blowup-fitting-ideal}에 의해
$\Coker(d_i)$는 Tor 차원 $\leq 1$이다. 따라서 $\Im(d_i)$는 유한 국소 자유
가군에서 Tor 차원 $\leq 1$인 유한표현 가군으로 가는 사상의 핵으로서 유한
국소 자유이다. 따라서 $\Ker(d_i)$도 유한 국소 자유이다(같은 논증). 그러므로
짧은 완전열
$$
0 \to \Im(d_{i - 1}) \to \Ker(d_i) \to H^i(E)|_V \to 0
$$
은 우리가 원하는 것을 보이며 증명이 끝난다.
\end{proof}

\begin{lemma}
\label{flat-lemma-blowup-complex-integral}
$X$를 정역 스킴이라 하고 $E \in D(\mathcal{O}_X)$를 완전이라고 하자.
그러면 모든 $i \in \mathbf{Z}$에 대하여 $H^i(E|_U)$가 일정한 계수 $r_i$의
유한 국소 자유가 되게 하는 공집합이 아닌 열린집합 $U \subset X$가 존재하고,
모든 $i \in \mathbf{Z}$에 대하여 $H^i(Lb^*E)$가 Tor 차원 $\leq 1$인
완전 $\mathcal{O}_{X'}$-가군이 되게 하는 $U$-허용 블로업
$b : X' \to X$가 존재한다.
\end{lemma}

\begin{proof}
독자는 $U$의 존재에 대한 자신의 증명을 찾기를 강력히 권한다.
$\eta \in X$를 일반점이라 하자. $E$를 $\eta$로 제한한 것은
$D(\kappa(\eta))$에서 미분들이 영인 유한차원 벡터공간들의 유한 복합체
$V^\bullet$과 동형이다. $r_i = \dim_{\kappa(\eta)} V^i$로 놓자. 그러면 항들이
$\mathcal{O}_X^{\oplus r_i}$이고 미분들이 영인 복합체로 표현되는
$D(\mathcal{O}_X)$의 완전 대상 $E'$은 $\eta$로 당긴 뒤 $E$와 동형이 된다.
따라서 스킴의 유도범주, 보조정리
\ref{perfect-lemma-descend-relatively-perfect}에 의해 $E|_U$와 $E'|_U$가 동형이
되게 하는 $\eta$의 열린근방 $U$가 존재한다. 이는 첫째 주장을 증명한다.
공집합이 아닌 모든 열린집합은 정역 스킴 $X$에서 스킴론적으로 조밀하므로,
둘째 주장은 첫째 주장과 보조정리 \ref{flat-lemma-blowup-complex}에서 따른다.
\end{proof}

\begin{remark}
\label{flat-remark-when-you-have-a-complex}
$X$를 스킴이라 하자. $E \in D(\mathcal{O}_X)$를 모든 $i \in \mathbf{Z}$에
대하여 $H^i(E)$가 Tor 차원 $\leq 1$인 완전 $\mathcal{O}_X$-가군이 되는
완전 대상이라 하자. 이 성질은 때때로 $E$에 관한 문제를 $H^i(E)$에 관한
문제로 환원하게 해 준다. 예를 들어
$$
\mathcal{E}^a \xrightarrow{d^a} \ldots
\xrightarrow{d^{b - 2}} \mathcal{E}^{b - 1}
\xrightarrow{d^{b - 1}} \mathcal{E}^b
$$
가 $E$를 표현하는 유한 국소 자유 $\mathcal{O}_X$-가군들의 유계 복합체라고
하자. 그러면 모든 $i$에 대하여 $\Im(d^i)$와 $\Ker(d^i)$는 유한 국소 자유
$\mathcal{O}_X$-가군이다. 실제로 $i$보다 큰 모든 지표에 대하여 이를 안다고
귀납적으로 가정하자. 먼저 짧은 완전열
$$
0 \to \Im(d^i) \to \Ker(d^{i + 1}) \to H^{i + 1}(E) \to 0
$$
과 $H^{i + 1}(E)$가 Tor 차원 $\leq 1$인 완전 가군이라는 가정을 사용하여
$\Im(d^i)$가 유한 국소 자유라고 결론 내릴 수 있다. 같은 논증을 짧은 완전열
$$
0 \to \Ker(d^i) \to \mathcal{E}^i \to \Im(d^i) \to 0
$$
에 다시 적용하면 $\Ker(d^i)$가 유한 국소 자유임을 얻는다. 따라서 특별 삼각형
$$
\tau_{\leq k - 1}E \to \tau_{\leq k}E \to H^k(E)[-k] \to
(\tau_{\leq k - 1}E)[1]
$$
은 유한 국소 자유 가군들의 유계 복합체들로 이루어진 다음 짧은 완전열들로
표현된다.
$$
\begin{matrix}
& &
& &
& &
0 \\
& &
& &
& &
\downarrow \\
\mathcal{E}^a & \to &
\ldots & \to &
\mathcal{E}^{k - 2} & \to &
\Ker(d^{k - 1}) \\
\downarrow & &
& &
\downarrow & &
\downarrow \\
\mathcal{E}^a & \to &
\ldots & \to &
\mathcal{E}^{k - 2} & \to &
\mathcal{E}^{k - 1} & \to &
\Ker(d^k) \\
& &
& &
& &
\downarrow & &
\downarrow \\
& &
& &
& &
\Im(d^{k - 1}) & \to &
\Ker(d^k) \\
& &
& &
& &
\downarrow \\
& &
& &
& &
0
\end{matrix}
$$
여기서 각 행이 복합체이며 ``자명한'' 영들은 그림에서 생략했다.
\end{remark}





\section{완전 가군의 블로업}
\label{flat-section-blowup-perfect-modules}

\noindent
이 절에서는 블로업한 뒤 Tor 차원 $\leq 1$인 완전 가군의 표준형을 찾으려
한다. 우리는 부분적으로만 성공한다.

\begin{lemma}
\label{flat-lemma-blowup-pd1-derived}
$X$를 스킴이라 하고 $\mathcal{F}$를 Tor 차원 $\leq 1$인 완전
$\mathcal{O}_X$-가군이라 하자. 임의의 블로업 $b : X' \to X$에 대하여
$Lb^*\mathcal{F} = b^*\mathcal{F}$이고 $b^*\mathcal{F}$는 Tor 차원
$\leq 1$인 완전 $\mathcal{O}_X$-가군이다.
% P05-EN-CH38-0313
\end{lemma}

\begin{proof}
$X = \Spec(A)$가 아핀이고 $\mathcal{F}$에 대응하는 $A$-가군 $M$이 표시
$$
0 \to A^{\oplus m} \to A^{\oplus n} \to M \to 0
$$
를 갖는다고 가정할 수 있다. $I \subset A$가 아이디얼이고 $a \in I$라고
하자. 아핀 블로업 대수 $A[\frac{I}{a}]$는 $A_a$의 부분환임을 상기하자.
국소화는 완전하므로 $A_a^{\oplus m} \to A_a^{\oplus n}$이 단사임을 알 수
있다. 따라서 $A[\frac{I}{a}]^{\oplus m} \to A[\frac{I}{a}]^{\oplus n}$도
단사이다. 이는 보조정리를 증명한다.
\end{proof}

\begin{lemma}
\label{flat-lemma-blowup-pd1}
$X$를 스킴이라 하고 $\mathcal{F}$를 Tor 차원 $\leq 1$인 완전
$\mathcal{O}_X$-가군이라 하자. $\mathcal{F}|_U$가 일정한 계수 $r$의 유한
국소 자유가 되게 하는 스킴론적으로 조밀한 열린집합 $U \subset X$를 택하자.
그러면 표준적인 짧은 완전열
$$
0 \to \mathcal{K} \to b^*\mathcal{F} \to \mathcal{Q} \to 0
$$
이 존재하게 하는 $U$-허용 블로업 $b : X' \to X$가 존재한다. 여기서
$\mathcal{Q}$는 계수 $r$의 유한 국소 자유이고, $\mathcal{K}$는 $U$로의
제한이 영인 Tor 차원 $\leq 1$의 완전 $\mathcal{O}_X$-가군이다.
% P05-EN-CH38-0314
\end{lemma}

\begin{proof}
인자, 보조정리 \ref{divisors-lemma-blowup-fitting-ideal}와 보조정리
\ref{flat-lemma-blowup-pd1-derived}를 결합하라.
\end{proof}

\begin{lemma}
\label{flat-lemma-canonical-blowup-torsion-pd1}
$X$를 스킴이라 하고 $\mathcal{F}$를 Tor 차원 $\leq 1$인 완전
$\mathcal{O}_X$-가군이라 하자. $\mathcal{F}|_U = 0$이 되게 하는 열린집합
$U \subset X$를 택하자. 그러면 $U$-허용 블로업
$$
b : X' \to X
$$
이 존재하여 $\mathcal{F}' = b^*\mathcal{F}$에는 두 표준적인 국소 유한 여과
$$
0 = F^0 \subset F^1 \subset F^2 \subset \ldots \subset \mathcal{F}'
\quad\text{이고}\quad
\mathcal{F}' = F_1 \supset F_2 \supset F_3 \supset \ldots \supset 0
$$
가 갖추어져 있고, 각 $n \geq 1$에 대하여 다음 성질을 갖는 유효 카르티에 인자
$D_n \subset X'$가 존재한다.
$$
F^i/F^{i - 1}
\quad\text{이고}\quad
F_i/F_{i + 1}
$$
는 $D_i$ 위에서 계수 $i$의 유한 국소 자유이다.
% P05-EN-CH38-0315
\end{lemma}

\begin{proof}
표시
$$
0 \to \mathcal{O}_V^{\oplus n} \xrightarrow{A} \mathcal{O}_V^{\oplus n} \to
\mathcal{F} \to 0
$$
가 어떤 $n$ 및 어떤 행렬 $A$에 대하여 존재하게 하는 아핀 열린집합
$V \subset X$를 선택하자. 우리가 블로업할 아이디얼은 $k \geq 0$에 대한 피팅
아이디얼들 $\text{Fit}_k(\mathcal{F})$의 곱이다. 위의 아핀 상황에서
$k > n$이면 $\text{Fit}_k(\mathcal{F})|_V = \mathcal{O}_V$임을 알 수 있으므로
이는 의미가 있다. 이것이 $U$-허용 블로업임은 분명하다. 인자, 보조정리
\ref{divisors-lemma-blowing-up-two-ideals}에 의해 $X'$ 위에서 아이디얼들
$\text{Fit}_k(\mathcal{F})$가 가역임을 알 수 있다. 따라서 다음 단락에서
논하는 경우로 환원된다.

\medskip\noindent
$k \geq 0$에 대하여 $\text{Fit}_k(\mathcal{F})$가 가역 아이디얼이라고
가정하자. $E_k \subset X$가 $k \geq 0$에 대하여
$\text{Fit}_k(\mathcal{F})$에 의해 정의되는 유효 카르티에 인자이면,
보조정리의 명제문에 있는 유효 카르티에 인자 $D_k$들은 다음을 만족한다.
$$
E_k = D_{k + 1} + 2 D_{k + 2} + 3 D_{k + 3} + \ldots
$$
모음 $D_k$가 국소 유한일 것이므로 이는 의미가 있다. 더욱이 이는 유효
카르티에 인자들 $D_k$를 유일하게 결정하므로 $D_k$를 국소적으로 구성하면
충분하다.

\medskip\noindent
위와 같은 $\mathcal{F}|_V$의 표시와 아핀 열린집합 $V \subset X$를 선택하자.
표시에 있는 정수 $n$에 관한 귀납법으로 인자와 여과를 구성하겠다. $k > n$에
대하여 $D_k|_V = \emptyset$으로 놓고 $D_n|V = E_{n - 1}|_V$로 놓는다.
% P05-EN-CH38-0316
$V$를 축소한 뒤 $\text{Fit}_{n - 1}(\mathcal{F})|_V$가 하나의 영인자가 아닌
원소 $f \in \Gamma(V, \mathcal{O}_V)$로 생성된다고 가정할 수 있다.
$\text{Fit}_{n - 1}(\mathcal{F})|_V$는 $A$의 성분들로 생성되는 아이디얼이므로
$A = fA'$을 만족하는 행렬 $A'$가 $\Gamma(V, \mathcal{O}_V)$에 있음을 알 수
있다. 짧은 완전열
$$
0 \to \mathcal{O}_V^{\oplus n} \xrightarrow{A'} \mathcal{O}_V^{\oplus n} \to
\mathcal{F}' \to 0
$$
로 $V$ 위의 $\mathcal{F}'$을 정의하자. $A'$의 성분들이
$\Gamma(V, \mathcal{O}_V)$의 단위 아이디얼을 생성하므로 $V$ 위 국소적으로
$\mathcal{F}'$은 $n$이 $1$만큼 감소된 표시를 가짐을 알 수 있다. 대수학,
보조정리 \ref{algebra-lemma-add-trivial-complex}을 보라. 또한
$f^{n - k}\text{Fit}_k(\mathcal{F}') = \text{Fit}_k(\mathcal{F})|_V$
가 $k = 0, \ldots, n$에 대하여 성립함에 유의하자. 따라서 모든 $k$에 대하여
$\text{Fit}_k(\mathcal{F}')$은 가역 아이디얼이다. 귀납법에 의해
$\mathcal{F}'$이 보조정리의 명제문과 같은 두 표준적 여과를 갖게 하는 유효
카르티에 인자 $D'_k \subset V$들이 존재한다고 결론 내린다. 이제
$k = 1, \ldots, n - 1$에 대하여 $D_k|_V = D'_k$로 놓는다. 위에 표시한 유효
카르티에 인자들 사이의 등식들이 이 선택에 대하여 성립함을 관찰하자. 끝으로
여과의 구성에 이른다. 즉 짧은 완전열
$$
0 \to
\mathcal{O}_{D_n \cap V}^{\oplus n} \to
\mathcal{F} \to \mathcal{F}' \to 0
\quad\text{이고}\quad
0 \to
\mathcal{F}' \to \mathcal{F} \to
\mathcal{O}_{D_n \cap V}^{\oplus n} \to 0
$$
이 있으며, 이는 $A$의 두 분해 $A = A'f = f A'$에서 온다. 이 완전열들은
표준적이다. 첫째 완전열에서 부분가군은
$\Ker(f : \mathcal{F} \to \mathcal{F})$이고, 둘째 완전열에서 몫가군은
$\Coker(f : \mathcal{F} \to \mathcal{F})$이기 때문이다.
\end{proof}

\begin{lemma}
\label{flat-lemma-blowup-map-pd1}
$X$를 스킴이라 하고 $\varphi : \mathcal{F} \to \mathcal{G}$를 Tor 차원
$\leq 1$인 완전 $\mathcal{O}_X$-가군들의 준동형이라 하자. $U \subset X$를
$\mathcal{F}|_U = 0$이고 $\mathcal{G}|_U = 0$인 스킴론적으로 조밀한
열린집합이라 하자. 그러면 $b^*\varphi$의 핵, 상, 여핵이 Tor 차원
$\leq 1$인 완전 $\mathcal{O}_{X'}$-가군들이 되게 하는 $U$-허용 블로업
$b : X' \to X$가 존재한다.
\end{lemma}

\begin{proof}
가정들에 의해 $D(\mathcal{O}_X)$의 대상
$(\mathcal{F} \to \mathcal{G})$는 완전이다. 따라서 보조정리
\ref{flat-lemma-blowup-complex} 및 \ref{flat-lemma-blowup-pd1-derived}에 의해(밑당김 뒤
복합체의 꼴을 보기 위해 후자를 사용한다) 여핵과 핵에 대하여 작동하는
$U$-허용 블로업을 얻는다. 상은 여핵으로 가는 사상의 핵이므로 역시 Tor 차원
$\leq 1$인 완전 가군이다.
\end{proof}











\section{Berthelot와 Ogus가 도입한 연산자}
\label{flat-section-eta}

\noindent
먼저 코호몰로지, 절 \ref{cohomology-section-eta}를 읽으라.

\medskip\noindent
$X$를 스킴이라 하고 $D \subset X$를 유효 카르티에 인자라 하자.
$\mathcal{I} = \mathcal{I}_D \subset \mathcal{O}_X$를 $D$의 아이디얼층이라
하자. 인자, 절 \ref{divisors-section-effective-Cartier-invertible}을 보라.
코호몰로지, 절 \ref{cohomology-section-eta}의 논의를 $X$와 $\mathcal{I}$에
적용할 수 있음은 분명하다.

\begin{lemma}
\label{flat-lemma-eta-stalks}
$X$를 스킴이라 하고 $D \subset X$를 아이디얼층
$\mathcal{I} \subset \mathcal{O}_X$를 갖는 유효 카르티에 인자라 하자. 모든
$i$에 대하여 $\mathcal{F}^i$가 $\mathcal{I}$-꼬임없음이 되게 하는 준연접
$\mathcal{O}_X$-가군들의 복합체를 $\mathcal{F}^\bullet$이라 하자. 그러면
$\eta_\mathcal{I}\mathcal{F}^\bullet$은 준연접 $\mathcal{O}_X$-가군들의
복합체이다. 더욱이 $U = \Spec(A) \subset X$가 아핀 열린집합이고
$D \cap U = V(f)$이면 $\eta_f(\mathcal{F}^\bullet(U))$는
$(\eta_\mathcal{I}\mathcal{F}^\bullet)(U)$와 표준적으로 동형이다.
\end{lemma}

\begin{proof}
생략한다.
\end{proof}

\begin{lemma}
\label{flat-lemma-Leta}
$X$를 스킴이라 하고 $D \subset X$를 아이디얼층
$\mathcal{I} \subset \mathcal{O}_X$를 갖는 유효 카르티에 인자라 하자.
코호몰로지, 보조정리 \ref{cohomology-lemma-Leta}의 함자
$L\eta_\mathcal{I} : D(\mathcal{O}_X) \to D(\mathcal{O}_X)$는
$D_\QCoh(\mathcal{O}_X)$를 그 자신으로 보낸다. 더욱이
$X = \Spec(A)$가 아핀이고 $D = V(f)$이면, 대수학 보충, 보조정리
\ref{more-algebra-lemma-Leta}에서 정의한 $D(A)$ 위의 함자 $L\eta_f$와
$D_\QCoh(\mathcal{O}_X)$ 위의 함자 $L\eta_\mathcal{I}$는 스킴의 유도범주,
보조정리 \ref{perfect-lemma-affine-compare-bounded}의 동치를 통해 대응한다.
\end{lemma}

\begin{proof}
생략한다.
\end{proof}





\section{복합체의 블로업 II}
\label{flat-section-blowup-complexes-II}

\noindent
이 절의 내용은 절 \ref{flat-section-blowup-complexes-III}에서 Macpherson의 그래프
구성의 한 변형을 구성하는 데 사용된다.
 
\begin{situation}
\label{flat-situation-complex-and-divisor}
여기서 $X$는 스킴이고, $D \subset X$는 아이디얼층
$\mathcal{I} \subset \mathcal{O}_X$를 갖는 유효 카르티에 인자이며, $M$은
$D(\mathcal{O}_X)$의 완전 대상이다.
\end{situation}

\noindent
$(X, D, M)$을 상황 \ref{flat-situation-complex-and-divisor}와 같은 삼중항이라 하자.
다음을 만족하는 아핀 열린집합 $U = \Spec(A) \subset X$를 생각하자.
\begin{enumerate}
\item 어떤 영인자가 아닌 원소 $f \in A$에 대하여 $D \cap U = V(f)$이다.
\item $M|_U$를 표현하는 유한 자유 $A$-가군들의 유계 복합체 $M^\bullet$이
존재한다(스킴의 유도범주, 보조정리
\ref{perfect-lemma-affine-compare-bounded}의 동치를 사용한다).
\end{enumerate}
$(U, A, f, M^\bullet)$을 $(X, D, M)$의 {\it 아핀 차트}라고 하겠다. 대수학
보충, 절 \ref{more-algebra-section-perfect-eta}에서 정의한 아이디얼
$I_i(M^\bullet, f) \subset A$를 생각하자. 모든 $x \in D$에 대하여
$x \in U$이고 모든 $i \in \mathbf{Z}$에 대하여 $I_i(M^\bullet, f)$가 주
아이디얼이 되게 하는 아핀 차트 $(U, A, f, M^\bullet)$이 존재하면
$(X, S, M)$을 {\it 좋은 삼중항}이라고 하자.
% P05-EN-CH38-0318

\begin{lemma}
\label{flat-lemma-pullback-triple-ideals-good}
상황 \ref{flat-situation-complex-and-divisor}에서 $h : Y \to X$를 $D$의 당김
$E = h^{-1}D$가 정의되는 스킴 사상이라고 하자(인자, 정의
\ref{divisors-definition-pullback-effective-Cartier-divisor}).
$(U, A, f, M^\bullet)$을 $(X, D, M)$의 아핀 차트라고 하자.
$V = \Spec(B) \subset Y$를 $h(V) \subset U$인 아핀 열린집합이라고 하자.
$f \in A$의 상을 $g \in B$로 나타내자. 그러면
\begin{enumerate}
\item $(V, B, g, M^\bullet \otimes_A B)$은 $(Y, E, Lh^*M)$의 아핀 차트이고,
\item $B$에서 $I_i(M^\bullet, f)B = I_i(M^\bullet \otimes_A B, g)$이며,
\item $(X, D, M)$이 좋은 삼중항이면 $(Y, E, Lh^*M)$도 좋은 삼중항이다.
\end{enumerate}
\end{lemma}

\begin{proof}
첫째 명제는 다음 관찰들에서 따른다. $g$는 $B$의 영인자가 아닌 원소이고
$E \cap V \subset V$를 정의한다. 또한 $M^\bullet \otimes_A B$는
$M^\bullet \otimes_A^\mathbf{L} B$를 표현하므로 스킴의 유도범주, 보조정리
\ref{perfect-lemma-quasi-coherence-pullback}에 따라 $M$의 $V$로의 당김을
표현한다. (2)는 (1)과 대수학 보충, 보조정리
\ref{more-algebra-lemma-eta-base-change-pre}에서 따른다. 이를 대수학 보충,
보조정리 \ref{more-algebra-lemma-eta-base-change-pre}와 결합하면 보조정리의
둘째 명제를 얻는다.
% P05-EN-CH38-0319
\end{proof}

\begin{lemma}
\label{flat-lemma-good-triple}
$X, D, \mathcal{I}, M$을 상황 \ref{flat-situation-complex-and-divisor}에서와
같이 두자. $(X, D, M)$이 좋은 삼중항이면 $L\eta_\mathcal{I}M$은
$D(\mathcal{O}_X)$의 완전 대상이다.
\end{lemma}

\begin{proof}
이는 대수학 보충, 보조정리
\ref{more-algebra-lemma-eta-locally-free}를 옮긴 것이다.
이 옮김에는 보조정리 \ref{flat-lemma-Leta}를 사용한다.
\end{proof}

\begin{lemma}
\label{flat-lemma-good-triple-bdd-loc-free}
$X, D, \mathcal{I}, M$을 상황 \ref{flat-situation-complex-and-divisor}에서와
같이 두고 $(X, D, M)$이 좋은 삼중항이라고 가정하자. $M$을 표현하는
유한 국소 자유 $\mathcal{O}_X$-가군들의 국소 유계 복합체
$\mathcal{M}^\bullet$이 존재하면, $L\eta_\mathcal{I}M$을 표현하는 유한
국소 자유 $\mathcal{O}_{X'}$-가군들의 국소 유계 복합체
$\mathcal{Q}^\bullet$이 존재한다.
% P05-EN-CH38-0320
\end{lemma}

\begin{proof}
코호몰로지, 보조정리 \ref{cohomology-lemma-Leta}에 따라 복합체
$\mathcal{Q}^\bullet = \eta_\mathcal{I}\mathcal{M}^\bullet$은
$L\eta_\mathcal{I}M$을 표현한다. 이 복합체가 국소 유계이고 유한 국소
자유 가군들로 이루어짐을 확인하기 위해 아핀 국소적으로 작업해도 된다.
그러면 유계성은 명백하다. 아이디얼 $I_i(M^\bullet, f)$들이 주
아이디얼이고 모든 $i$에 대하여 $\mathcal{M}^i|_U$가 유한 자유인
$(X, D, M)$의 아핀 차트 $(U, A, f, M^\bullet)$을 택하자.
$(X, D, M)$이 좋은 삼중항이라는 가정 때문에 이렇게 할 수 있다.
$N^i = \Gamma(U, \mathcal{M}^i|_U)$라고 쓰면, 유한 자유 $A$-가군들의
유계 복합체 $N^\bullet$을 얻으며 이는 $D(A)$에서 복합체 $M^\bullet$과
같은 대상을 표현한다(스킴의 유도범주, 보조정리
\ref{perfect-lemma-affine-compare-bounded}). 그러면 대수학 보충, 보조정리
\ref{more-algebra-lemma-ideal-well-defined}에 따라 모든 $i$에 대하여
$I_i(N^\bullet, f)$는 주 아이디얼이다. 따라서 복합체
$\eta_fN^\bullet$은 유한 국소 자유 $A$-가군들의 유계 복합체이다.
보조정리 \ref{flat-lemma-eta-stalks}에 따라 $\mathcal{Q}^i|_U$는
$\eta_fN^i$에 대응하는 준연접 $\mathcal{O}_U$-가군이므로 결론을 얻는다.
\end{proof}

\begin{lemma}
\label{flat-lemma-complex-and-divisor-eta-pull}
상황 \ref{flat-situation-complex-and-divisor}에서 $h : Y \to X$를 당김
$E = h^{-1}D$가 정의되는 스킴 사상이라고 하자. $(X, D, M)$이 좋은
삼중항이면
$$
Lh^*(L\eta_\mathcal{I}M) = L\eta_\mathcal{J}(Lh^*M)
$$
$D(\mathcal{O}_Y)$에서 위 등식이 성립한다. 여기서 $\mathcal{J}$는
$E$의 아이디얼 층이다.
\end{lemma}

\begin{proof}
이는 대수학 보충, 보조정리
\ref{more-algebra-lemma-eta-base-change}를 옮긴 것이다. 이 옮김에는
보조정리 \ref{flat-lemma-eta-stalks}와 \ref{flat-lemma-Leta}를 사용한다.
\end{proof}

\begin{lemma}
\label{flat-lemma-complex-and-divisor-blowup-pre}
상황 \ref{flat-situation-complex-and-divisor}에서 다음을 만족하는 유일한 사상
$b : X' \to X$가 존재한다.
\begin{enumerate}
\item 당김 $D' = b^{-1}D$가 정의되고, $M' = Lb^*M$이라 할 때
$(X', D', M')$은 좋은 삼중항이다.
\item 당김 $E = h^{-1}D$가 정의되고 $(Y, E, Lh^*M)$이 좋은 삼중항인
임의의 스킴 사상 $h : Y \to X$에 대하여 $h$는 $b$를 통한 유일한
분해를 갖는다.
\end{enumerate}
더욱이 임의의 아핀 차트 $(U, A, f, M^\bullet)$에 대하여 제한
$b^{-1}(U) \to U$는 아이디얼 $I_i(M^\bullet, f)$들의 곱에 대한
블로업이고, 임의의 준콤팩트 열린집합 $W \subset X$에 대하여 제한
$b|_{b^{-1}(W)} : b^{-1}(W) \to W$는 $W \setminus D$-허용 블로업이다.
\end{lemma}

\begin{proof}
증명은 각 아핀 차트 $(U, A, f, M^\bullet)$에서 곱 아이디얼
$I_i(M^\bullet, f)$들에 대하여 $X$를 국소적으로 블로업하는 것이다.
대수학 보충, 절 \ref{more-algebra-section-perfect-eta}의 처음 몇
보조정리는 이것이 잘 정의됨을 보인다. 그러면 보편 성질 (2)는 블로업의
보편 성질에서 따른다. 자세한 내용은 다음과 같다.

\medskip\noindent
$U, A, f, M^\bullet$을 $(X, D, M)$의 아핀 차트라고 하자. 유한 개를
제외한 모든 아이디얼 $I_i(M^\bullet, f)$가 $A$와 같으므로
$$
I = \prod\nolimits_i I_i(M^\bullet, f)
$$
를 생각할 수 있으며 이는 $A$의 유한 생성 아이디얼이다.
$$
b_U : U' \to U
$$
를 $I$에 대한 $U$의 블로업으로 나타내자. 그러면 인자, 보조정리
\ref{divisors-lemma-blow-up-pullback-effective-Cartier}에 따라
$b_U^{-1}(U \cap D)$가 정의된다. $f^{r_i} \in I_i(M^\bullet, f)$임을
상기하면 $b_U$는 $(U \setminus D)$-허용 블로업이다. 인자, 보조정리
\ref{divisors-lemma-blowing-up-two-ideals}에 따라 각 $i$에 대하여 사상
$b_U$는 $U' \to U'_i \to U$로 분해된다. 여기서 $U'_i \to U$는
$I_i(M^\bullet, f)$에 대한 블로업이고 $U' \to U'_i$는 또 다른
블로업이다. 따라서 $I_i(M^\bullet, f)$의 $U'$로의 당김
$I_i(M^\bullet, f)\mathcal{O}_{U'}$는 가역 아이디얼 층이다. 인자,
보조정리 \ref{divisors-lemma-blow-up-pullback-effective-Cartier}와
\ref{divisors-lemma-blowing-up-gives-effective-Cartier-divisor}를 보라.
따라서 $(U', b^{-1}D, Lb^*M|_U)$는 좋은 삼중항이다. 당김 아래에서의
아이디얼 $I_i(-,-)$의 거동에 대해서는 보조정리
\ref{flat-lemma-pullback-triple-ideals-good}를 보라.
% P05-EN-CH38-0321
마지막으로 $b_U : U' \to U$가 보조정리의 명제 (2)에 언급된 보편
성질을 갖는다고 주장한다. 즉, $h : Y \to U$가 당김
$E = h^{-1}(D \cap U)$가 정의되고 $(Y, E, Lh^*M)$이 좋은 삼중항인
스킴 사상이라고 하자. 그러면 $Y$는 모든 $i$에 대하여
$I_i(N^\bullet, g)$가 가역 아이디얼인 아핀 차트
$(V, B, g, N^\bullet)$들로 덮인다. $g$와 $f$의 $B$에서의 상은 모두
유효 카르티에 인자 $E \cap V$를 잘라 내므로 단위원만큼 다르다.
따라서 대수학 보충, 보조정리
\ref{more-algebra-lemma-eta-change-unit}에 의해 $g$가 $f$의 상이라고
가정해도 된다.
그러면 대수학 보충, 보조정리
\ref{more-algebra-lemma-ideal-well-defined}에 따라 $B$-가군으로서
$I_i(N^\bullet, g)$는 $I_i(M^\bullet \otimes_A B, g)$와 동형이다.
따라서 $I_i(M^\bullet \otimes_A B, g) = I_i(M^\bullet, f)B$
(보조정리 \ref{flat-lemma-pullback-triple-ideals-good})는 가역 $B$-가군이다.
따라서 아이디얼 $IB$는 가역이고, 이에 따라 $I\mathcal{O}_Y$도
가역이다. 그러므로 인자, 보조정리
\ref{divisors-lemma-universal-property-blowing-up}에 의해 $h$의 $b_U$를
통한 유일한 분해를 얻는다.

\medskip\noindent
$(X, D, M)$의 아핀 차트 $(U, A, f, M^\bullet)$이 존재하는 아핀
열린집합 $U \subset X$들의 집합을 $\mathcal{B}$라고 하자. 그러면
$\mathcal{B}$는 $X$의 위상의 기저이다. 자세한 내용은 생략한다.
$U \in \mathcal{B}$에 대하여 위에서 구성한 사상 $b_U : U' \to U$가
있으며 이는 $U$ 위에서 보편 성질을 만족한다.
$U_1 \subset U_2 \subset X$가 모두 $\mathcal{B}$에 속하면
$b_{U_1} : U'_1 \to U_1$은 표준적으로
$$
b_{U_2}|_{b_{U_2}^{-1}(U_1)} : b_{U_2}^{-1}(U_1) \longrightarrow U_1
$$
와 동형이다. 이는 보편 성질에서 따른다. 달리 말하면 $U_1$ 위의 동형
$U'_1 \to b_{U_2}^{-1}(U_1)$을 얻는다. 이 동형들은 (다시 보편 성질에
따라) 코사이클 조건을 만족하므로 구성, 보조정리
\ref{constructions-lemma-relative-glueing}에 의해 각 $U$가
$\mathcal{B}$에 속할 때 그 위에서의 제한이 $U' \to U$와 동형인 사상
$b : X' \to X$를 얻는다. 이 성질들은 국소적으로 확인할 수 있으므로
사상 $b : X' \to X$는 보조정리 명제의 성질 (1)과 (2)를 만족한다.
자세한 내용은 생략한다.

\medskip\noindent
보조정리의 마지막 주장을 증명하는 일이 남았다. $W \subset X$를
준콤팩트 열린집합이라고 하자. 각 $1 \leq t \leq T$에 대하여 아핀 차트
$(U_t, A_t, f_t, M_t^\bullet)$이 존재하도록 유한 덮개
$W = U_1 \cup \ldots \cup U_T$를 택하자. 아래에서는 임의의 아핀
열린집합 $V = \Spec(B) \subset U_t \cap U_{t'}$에 대하여 다음을
사용한다. (a) $f_t$와 $f_{t'}$의 $B$에서의 상은 단위원만큼 다르고,
(b) 복합체 $M_t^\bullet \otimes_A B$와 $M_{t'} \otimes_A B$는
$D(B)$에서 동형인 대상들을 정의한다.
% P05-EN-CH38-0322
$i \in \mathbf{Z}$에 대하여
$$
N_i = \max\nolimits_{t = 1, \ldots, T}
\left(\sum\nolimits_{j \geq i} (-1)^{j - i} rk(M_t^j)\right)
$$
라고 두자. 그러면
$N_t - \sum\nolimits_{j \geq i} (-1)^{j - i} rk(M_t^j) \geq 0$이고,
다음 아이디얼들을 생각할 수 있다.
% P05-EN-CH38-0323
$$
I_{t, i} =
f_t^{N_i - \sum\nolimits_{j \geq i} (-1)^{j - i} rk(M_t^j)}
I_i(M_t^\bullet, f_t)
\subset A_t
$$
대수학 보충, 보조정리 \ref{more-algebra-lemma-eta-change-unit}와
\ref{more-algebra-lemma-ideal-well-defined}에 따라 아이디얼
$I_{t, i}$들은 이어붙어 준연접 유한형 아이디얼
$\mathcal{I}_i \subset \mathcal{O}_W$가 된다. 더욱이 유한 개를 제외한
모든 이 아이디얼은 $\mathcal{O}_W$와 같다. 위에서 구성한 사상
$X' \to X$의 제한은 명백히 아이디얼 $\mathcal{I}_i$들의 곱에 대한
$W$의 블로업이다. 이로써 증명이 끝난다.
\end{proof}

\begin{lemma}
\label{flat-lemma-complex-and-divisor-blowup}
상황 \ref{flat-situation-complex-and-divisor}에서 $b : X' \to X$를 보조정리
\ref{flat-lemma-complex-and-divisor-blowup-pre}의 사상이라고 하자. 아이디얼 층
$\mathcal{I}' \subset \mathcal{O}_{X'}$을 갖는 유효 카르티에 인자
$D' = b^{-1}D$를 생각하자. 그러면 $Q = L\eta_{\mathcal{I}'}Lb^*M$은
$D(\mathcal{O}_{X'})$의 완전 대상이다.
\end{lemma}

\begin{proof}
보조정리 \ref{flat-lemma-complex-and-divisor-blowup-pre}와
\ref{flat-lemma-good-triple}에서 따른다.
\end{proof}

\begin{lemma}
\label{flat-lemma-complex-and-divisor-blowup-base-change}
상황 \ref{flat-situation-complex-and-divisor}에서 $h : Y \to X$를 당김
$E = h^{-1}D$가 정의되는 스킴 사상이라고 하자. 각각 $D \subset X$와
$M$, 그리고 $E \subset Y$와 $Lh^*M$에 대하여 보조정리
\ref{flat-lemma-complex-and-divisor-blowup-pre}에서 구성한 사상을 각각
$b : X' \to X$, $c : Y' \to Y$라고 하자. 그러면 $Y'$은
$b : X' \to X$에 관한 $Y$의 강변환이고(정확한 정식화는 증명을 보라),
$$
L\eta_{\mathcal{J}'}L(h \circ c)^*M = L(Y' \to X')^*Q
$$
이다. 여기서 $Q = L\eta_{\mathcal{I}'}Lb^*M$는 보조정리
\ref{flat-lemma-complex-and-divisor-blowup}에서와 같다. 특히
$(Y, E, Lh^*M)$이 좋은 삼중항이고 $k : Y \to X'$가
$h = b \circ k$를 만족하는 유일한 사상이면
$L\eta_\mathcal{J}Lh^*M = Lk^*Q$이다.
\end{lemma}

\begin{proof}
$E' = c^{-1}E$라고 쓰자. 그러면 $(Y', E', L(h \circ c)^*M)$은 좋은
삼중항이다. 따라서 보조정리
\ref{flat-lemma-complex-and-divisor-blowup-pre}의 보편 성질에 의해 유일한 사상
$$
h' : Y' \longrightarrow X'
$$
이 존재하여 $b \circ h' = h \circ c$를 만족한다. 특히 사상
$(h', c) : Y' \to X' \times_X Y$가 존재한다. $W \subset X$가
$b^{-1}(W) \to W$가 블로업인 준콤팩트 열린집합이면, 이 사상이
$Y'|_W$를 $b^{-1}(W) \to W$에 관한 $Y_W$의 강변환과 식별한다고
주장한다. 이 주장을 확인하는 것은 $W$ 위의 국소 문제이므로 아핀 차트
위에서 증명해도 된다. 다음 문단에서 이를 수행한다.

\medskip\noindent
$(U, A, f, M^\bullet)$을 $(X, D, M)$의 아핀 차트라고 하자. 보조정리
\ref{flat-lemma-complex-and-divisor-blowup}의 증명에서 $b : X' \to X$의
$U$로의 제한은 아이디얼 $I_i(M^\bullet, f)$들의 곱에 대한
$U = \Spec(A)$의 블로업임을 상기하자. 이제
$V = \Spec(B) \subset Y$가 $h(V) \subset U$인 임의의 아핀
열린집합이면, $f$의 상을 $g \in B$라 할 때
$(V, B, g, M^\bullet \otimes_A B)$는 $(Y, E, Lh^*M)$의 아핀 차트이다.
보조정리 \ref{flat-lemma-pullback-triple-ideals-good}를 보라. 따라서
$c : Y' \to Y$의 $V$로의 제한은 아이디얼
$I_i(M^\bullet, f)B$들의 곱에 대한 블로업이다. 즉, $h^{-1}(U)$ 위의
사상 $c : Y' \to Y$는 아이디얼
$\prod I_i(M^\bullet, f) \mathcal{O}_{h^{-1}(U)}$에 대한
$h^{-1}(U)$의 블로업이다. 강변환에 대해서도 마찬가지이므로 강변환에
관한 우리의 주장이 참이다.

\medskip\noindent
이제 등식
$L\eta_{\mathcal{J}'}L(h \circ c)^*M = L(Y' \to X')^*Q$
은 보조정리 \ref{flat-lemma-complex-and-divisor-eta-pull}에서 따른다.
마지막 명제는 이것의 특수한 경우, 즉 $c = \text{id}_Y$이고 $k = h'$인
경우이다.
\end{proof}

\begin{lemma}
\label{flat-lemma-complex-and-divisor-blowup-good}
상황 \ref{flat-situation-complex-and-divisor}에서 $W \subset X$를 $M$의
코호몰로지 층들이 국소 자유인 최대 열린 부분스킴이라고 하자. 그러면
보조정리 \ref{flat-lemma-complex-and-divisor-blowup-pre}의 사상
$b : X' \to X$는 $W$ 위에서 동형이다.
\end{lemma}

\begin{proof}
$U \subset W$인 임의의 아핀 차트 $(U, A, f, M^\bullet)$에 대하여
대수학 보충, 보조정리
\ref{more-algebra-lemma-eta-cohomology-free}에 따라
$I_i(M^\bullet, f)$들이 국소적으로 $f$의 거듭제곱으로 생성되기
때문이다. $f$는 영인자가 아니므로 블로업 $b^{-1}(U) \to U$는
동형이다.
\end{proof}

\begin{lemma}
\label{flat-lemma-complex-and-divisor-blowup-good-T}
$X, D, \mathcal{I}, M$을 상황 \ref{flat-situation-complex-and-divisor}에서와
같이 두자. $(X, D, M)$이 좋은 삼중항이면 다음 성질들을 갖는 유한 표시
닫힌 몰입
$$
i : T \longrightarrow D
$$
이 존재한다.
\begin{enumerate}
\item $W \subset X$가 $M$의 코호몰로지 층들이 국소 자유인 최대 열린
부분스킴일 때, $T$는 $D \cap W$를 스킴론적으로 포함한다.
\item $Li^*L\eta_\mathcal{I}M$의 코호몰로지 층들은 국소 자유이다.
\item 상이 $x = i(t) \in W$인 임의의 점 $t \in T$에 대하여
$\mathcal{O}_{X, x}$ 위의 $H^i(M)_x$의 계수와
$\mathcal{O}_{T, t}$ 위의 $H^i(Li^*L\eta_\mathcal{I}M)_t$의 계수는
같다.
\end{enumerate}
\end{lemma}

\begin{proof}
$(U, A, f, M^\bullet)$을 모든 $i \in \mathbf{Z}$에 대하여
$I_i(M^\bullet, f)$가 주 아이디얼인 $(X, D, M)$의 아핀 차트라고 하자.
그러면 $T \cap U \subset D \cap U$를 다음 아이디얼로 정의되는 닫힌
부분스킴으로 정의한다.
$$
J(M^\bullet, f) = \sum J_i(M^\bullet, f) \subset A/fA
$$
이 아이디얼은 대수학 보충, 보조정리
\ref{more-algebra-lemma-eta-vanishing-beta-plus-pre}와
\ref{more-algebra-lemma-eta-vanishing-beta-plus}에서 연구했다. 둘째
보조정리의 표기로 $T \cap U \to D \cap U$는 거기서 연구한 환 준동형
$A/fA \to C$로 주어짐을 알 수 있다. $(X, D, M)$이 좋은 삼중항이므로
$X$를 이러한 꼴의 아핀 차트들로 덮을 수 있고, 두 보조정리 중 첫째에
따라 이 구성은 이어붙는다. 따라서 닫힌 부분스킴 $T \subset D$를 얻으며,
위와 같은 좋은 아핀 차트 위에서는 아이디얼 $J(M^\bullet, f)$로
주어진다. 그러면 성질 (1)과 (2)는 둘째 보조정리에서 따른다. 자세한
내용은 생략한다. 독자를 돕기 위해 다음을 관찰하자. 대수학 보충,
보조정리 \ref{more-algebra-lemma-eta-locally-free}에 따라
$\eta_fM^\bullet$은 국소 자유 가군들의 복합체이므로
$Li^*L\eta_\mathcal{I}M|_{T \cap U}$는 $C$-가군들의 복합체
$\eta_fM^\bullet \otimes_A C$로 표현된다. 계수에 관한 명제 (3)은
코호몰로지, 보조정리
\ref{cohomology-lemma-eta-cohomology-locally-free}에서 따른다.
\end{proof}

\begin{lemma}
\label{flat-lemma-complex-and-divisor-blowup-T}
상황 \ref{flat-situation-complex-and-divisor}에서 $b : X' \to X$와 $D'$를
보조정리 \ref{flat-lemma-complex-and-divisor-blowup-pre}에서와 같이 두고,
$Q = L\eta_{\mathcal{I}'}Lb^*M$를 보조정리
\ref{flat-lemma-complex-and-divisor-blowup}에서와 같이 두자. $W \subset X$를
$M$이 국소 자유 코호몰로지 가군들을 갖는 최대 열린집합이라고 하자.
그러면 다음을 만족하는 유한 표시 닫힌 몰입 $i : T \to D'$가 존재한다.
\begin{enumerate}
\item $D' \cap b^{-1}(W) \subset T$가 스킴론적으로 성립한다.
\item $Li^*Q$는 국소 자유 코호몰로지 층들을 갖는다.
\item $w \in W$로 가는 $t \in T$에 대하여 $\mathcal{O}_{T, t}$ 위의
$H^i(Li^*Q)_t$의 계수는 $\mathcal{O}_{X, x}$ 위의 $H^i(M)_x$의
계수와 같다.
% P05-EN-CH38-0324
\end{enumerate}
\end{lemma}

\begin{proof}
보조정리 \ref{flat-lemma-complex-and-divisor-blowup-good}에 따르면 $b$는
$W$ 위에서 동형이다. 따라서 $b^{-1}(W) \subset X'$는 $Lb^*M$이 국소
자유 코호몰로지 층들을 갖는 최대 열린집합 $W' \subset X'$에 포함된다.
이제 보조정리의 실제 명제들은 $(X', D', Lb^*M)$에 적용한 보조정리
\ref{flat-lemma-complex-and-divisor-blowup-good-T}와 명제에 언급된 다른
보조정리들에서 즉시 따른다.
\end{proof}

\begin{lemma}
\label{flat-lemma-complex-and-divisor-blowup-T-ranks}
상황 \ref{flat-situation-complex-and-divisor}에서 $b : X' \to X$,
$D' \subset X'$, $Q$를 보조정리
\ref{flat-lemma-complex-and-divisor-blowup}에서와 같이 두자.
$\rho = (\rho_i)_{i \in \mathbf{Z}}$를 정수들의 족이라고 하자.
$W(\rho) \subset X$를 모든 $i$에 대하여 $H^i(M)$이 계수 $\rho_i$인
국소 자유 층인 최대 열린 부분스킴이라고 하자. $i : T \to D'$를
보조정리 \ref{flat-lemma-complex-and-divisor-blowup-T}에서와 같이 두자.
그러면 $D' \cap b^{-1}(W(\rho))$를 스킴론적으로 포함하고 모든 $i$에
대하여 $H^i(Li^*Q|_{T(\rho)})$가 계수 $\rho_i$인 국소 자유 층이 되는
열리고 닫힌 부분스킴 $T(\rho) \subset T$가 존재한다.
\end{lemma}

\begin{proof}
$T(\rho) \subset T$를 모든 $i$에 대하여 $H^i(Li^*Q)$가 계수
$\rho_i$를 갖는 열리고 닫힌 부분스킴이라고 하자. 그러면 명제는
보조정리 \ref{flat-lemma-complex-and-divisor-blowup-T}의 코호몰로지 가군의
계수에 관한 주장으로부터 즉시 따른다.
\end{proof}

\begin{lemma}
\label{flat-lemma-complex-and-divisor-blowup-T-represented-bdd-loc-free}
상황 \ref{flat-situation-complex-and-divisor}에서 $b : X' \to X$,
$D' \subset X'$, $Q$를 보조정리
\ref{flat-lemma-complex-and-divisor-blowup}에서와 같이 두자. $M$을 표현하는
유한 국소 자유 $\mathcal{O}_X$-가군들의 국소 유계 복합체
$\mathcal{M}^\bullet$이 존재하면, $Q$를 표현하는 유한 국소 자유
$\mathcal{O}_{X'}$-가군들의 국소 유계 복합체 $\mathcal{Q}^\bullet$이
존재한다.
\end{lemma}

\begin{proof}
$\mathcal{I}'$가 유효 카르티에 인자 $D'$의 아이디얼 층일 때
$Q = L\eta_{\mathcal{I}'}Lb^*M$임을 상기하자. 유한 국소 자유
$\mathcal{O}_{X'}$-가군들의 국소 유계 복합체
$(\mathcal{M}')^\bullet = b^*\mathcal{M}^\bullet$은 $Lb^*M$을 표현한다.
따라서 보조정리는 보조정리
\ref{flat-lemma-good-triple-bdd-loc-free}에서 따른다.
\end{proof}

\begin{lemma}
\label{flat-lemma-complex-and-divisor-derived}
$X$를 스킴, $D \subset X$를 유효 카르티에 인자라고 하자.
$M \in D(\mathcal{O}_X)$을 완전 대상이라고 하자. $W \subset X$를
코호몰로지 층 $H^i(M)$들이 국소 자유인 최대 열린집합이라고 하자.
그러면 다음 성질들을 갖는 고유 사상 $b : X' \longrightarrow X$와
$D(\mathcal{O}_{X'})$의 대상 $Q$가 존재한다.
\begin{enumerate}
\item $b : X' \to X$는 $X \setminus D$ 위에서 동형이다.
\item $b : X' \to X$는 $W$ 위에서 동형이다.
\item $D' = b^{-1}D$는 유효 카르티에 인자이다.
\item $\mathcal{I}'$가 $D'$의 아이디얼 층일 때
$Q = L\eta_{\mathcal{I}'}Lb^*M$이다.
\item $Q$는 $D(\mathcal{O}_{X'})$의 완전 대상이다.
\item 다음을 만족하는 유한 표시 닫힌 몰입 $i : T \to D'$가 존재한다.
\begin{enumerate}
\item $D' \cap b^{-1}(W) \subset T$가 스킴론적으로 성립한다.
\item $Li^*Q$는 유한 국소 자유 코호몰로지 층들을 갖는다.
\item 상이 $w \in W$인 $t \in T$에 대하여 $\mathcal{O}_{T, t}$ 위의
$H^i(Li^*Q)_t$의 계수는 $\mathcal{O}_{X, x}$ 위의 $H^i(M)_x$의
계수와 같다.
\end{enumerate}
\item $(X, D, M)$의 임의의 아핀 차트 $(U, A, f, M^\bullet)$에 대하여
$b$의 $U$로의 제한은 아이디얼 $I = \prod I_i(M^\bullet, f)$에 대한
$U = \Spec(A)$의 블로업이다.
\item $I_i(N^\bullet, g)$가 주 아이디얼인
$(X', D', Lb^*N)$의 임의의 아핀 차트 $(V, B, g, N^\bullet)$에 대하여
다음이 성립한다.
% P05-EN-CH38-0325
\begin{enumerate}
\item $Q|_V$는 $\eta_gN^\bullet$에 대응한다.
\item $T \subset V \cap D'$는 아이디얼
$J(N^\bullet, g) = \sum J_i(N^\bullet, g) \subset B/gB$
에 대응하며, 이 아이디얼은 대수학 보충, 보조정리
\ref{more-algebra-lemma-eta-vanishing-beta-plus}에서 연구했다.
% P05-EN-CH38-0326
\end{enumerate}
\item $M$이 유한 국소 자유 $\mathcal{O}_X$-가군들의 국소 유계
복합체로 표현될 수 있으면, $Q$는 유한 국소 자유
$\mathcal{O}_{X'}$-가군들의 유계 복합체로 표현될 수 있다.
% P05-EN-CH38-0327
\end{enumerate}
\end{lemma}

\begin{proof}
이 명제는 보조정리 \ref{flat-lemma-pullback-triple-ideals-good},
\ref{flat-lemma-good-triple},
\ref{flat-lemma-complex-and-divisor-eta-pull},
\ref{flat-lemma-complex-and-divisor-blowup-pre},
\ref{flat-lemma-complex-and-divisor-blowup},
\ref{flat-lemma-complex-and-divisor-blowup-base-change},
\ref{flat-lemma-complex-and-divisor-blowup-good},
\ref{flat-lemma-complex-and-divisor-blowup-good-T},
\ref{flat-lemma-complex-and-divisor-blowup-T}, 그리고
\ref{flat-lemma-complex-and-divisor-blowup-T-represented-bdd-loc-free}에서
얻은 정보를 모은 것이다.
\end{proof}










\section{복합체의 블로업 III}
\label{flat-section-blowup-complexes-III}

\noindent
이 절에서는 \cite[Section 18.1]{F}에 주어진 맥퍼슨의 그래프 구성의
``대수 버전''을 제시한다.

\medskip\noindent
$X$를 스킴이라고 하자. $E$를 $D(\mathcal{O}_X)$의 완전 대상이라고
하자. $U \subset X$를 $E|_U$가 국소 자유 코호몰로지 층들을 갖는 최대
열린 부분스킴이라고 하자.

\medskip\noindent
다음 가환 도식을 생각하자.
$$
\xymatrix{
\mathbf{A}^1_X \ar[r] \ar[rd] &
\mathbf{P}^1_X \ar[d]^p &
(\mathbf{P}^1_X)_\infty \ar[l] \ar[ld] \\
& X \ar@/_1em/[ur]_\infty
}
$$
여기서 $\mathbf{A}^1 = D_+(T_0)$은 $\mathbf{P}^1$의 첫째 표준 아핀
열린집합이고 $\infty = V_+(T_0)$은 그 여집합인 유효 카르티에 인자이며,
위 도식은 이 스킴들을 $X$로 당긴 것임을 상기하자.
$\infty : X \to (\mathbf{P}^1_X)_\infty$는 동형임을 관찰하자. 그러면
$$
(\mathbf{P}^1_X, (\mathbf{P}^1_X)_\infty, Lp^*E)
$$
은 절 \ref{flat-section-blowup-complexes-II}의 상황
\ref{flat-situation-complex-and-divisor}에서와 같은 삼중항이다.
$$
b : W \longrightarrow \mathbf{P}^1_X\quad\text{이고}\quad
W_\infty = b^{-1}((\mathbf{P}^1_X)_\infty)
$$
를 이 삼중항에서 출발하여 보조정리
\ref{flat-lemma-complex-and-divisor-blowup-pre}에서 구성한 블로업과 유효
카르티에 인자라고 하자. 또한
$$
Q = L\eta_{\mathcal{I}} Lb^*M = L\eta_\mathcal{I} L(p \circ b)^*E
$$
를 보조정리 \ref{flat-lemma-complex-and-divisor-blowup}에서 생각한
$D(\mathcal{O}_W)$의 완전 대상으로 나타내자. 여기서
$\mathcal{I} \subset \mathcal{O}_W$는 $W_\infty$의 아이디얼 층이다.
% P05-EN-CH38-0328

\begin{lemma}
\label{flat-lemma-graph-construction}
위 구성은 다음 성질들을 갖는다.
\begin{enumerate}
\item $b$는 $\mathbf{P}^1_U \cup \mathbf{A}^1_X$ 위에서 동형이다.
\item $Q$의 $\mathbf{A}^1_X$로의 제한은 $E$의 당김과 같다.
\item $(W_\infty \to X)^{-1}U \subset T$가 스킴론적으로 성립하고
$Li^*Q$가 국소 자유 코호몰로지 층들을 갖도록 하는 유한 표시 닫힌 몰입
$i : T \to W_\infty$가 존재한다.
\item 상이 $u \in U$인 $t \in T$에 대하여 $\mathcal{O}_{T, t}$ 위의
$H^i(Li^*Q)_t$의 계수는 $\mathcal{O}_{U, u}$ 위의 $H^i(M)_u$의
계수와 같다.
% P05-EN-CH38-0329
\item $E$가 유한 국소 자유 $\mathcal{O}_X$-가군들의 국소 유계
복합체로 표현될 수 있으면 $Q$는 유한 국소 자유
$\mathcal{O}_W$-가군들의 국소 유계 복합체로 표현될 수 있다.
\end{enumerate}
\end{lemma}

\begin{proof}
이는 절 \ref{flat-section-blowup-complexes-II}의 결과들에서 즉시 따른다.
필요한 모든 것을 모은 명제는 보조정리
\ref{flat-lemma-complex-and-divisor-derived}를 보라.
\end{proof}
